\documentclass[UTF8,fontset=none,11pt,oneside,openany]{ctexbook}

\usepackage{iftex}
\usepackage{fontspec}
\ifXeTeX
  \usepackage{xeCJK}
\fi
\defaultfontfeatures{Ligatures=TeX}
\setmainfont{texgyretermes-regular.otf}[
  BoldFont=texgyretermes-bold.otf,
  ItalicFont=texgyretermes-italic.otf,
  BoldItalicFont=texgyretermes-bolditalic.otf
]
\setsansfont{texgyreheros-regular.otf}
\setmonofont{lmmono10-regular.otf}[ItalicFont=lmmono10-italic.otf]
\setCJKmainfont{FandolSong-Regular.otf}[
  BoldFont=FandolSong-Bold.otf,
  ItalicFont=FandolKai-Regular.otf,
  BoldItalicFont=FandolSong-Bold.otf
]
\setCJKsansfont{FandolHei-Regular.otf}
\setCJKmonofont{FandolFang-Regular.otf}
\ifXeTeX
  \xeCJKsetup{PunctStyle=kaiming,CheckSingle=true}
  \XeTeXlinebreaklocale "zh"
  \XeTeXlinebreakskip=0pt plus 1pt
\fi

\usepackage{amsmath}
\ifLuaTeX
  \usepackage{unicode-math}
  \setmathfont{texgyretermes-math.otf}
  \providecommand{\square}{\mdlgwhtsquare}
\else
  \usepackage{amssymb,mathrsfs}
\fi
\usepackage{graphicx}
\usepackage{tikz-cd}
\ifLuaTeX
  \usepackage{luatex85}
\fi
\usepackage[all]{xy}
% Canonical EGA shorthand used by the locked indexes and errata.
% Keep these as aliases rather than rewriting mathematical source payloads.
\providecommand{\sh}{\mathcal}
\providecommand{\vphi}{\phi}
\providecommand{\kres}{k}
\providecommand{\supertilde}{{\,\widetilde{\,}\,}}
\providecommand{\Hom}{\operatorname{Hom}}
\providecommand{\Proj}{\operatorname{Proj}}
\providecommand{\Spec}{\operatorname{Spec}}
\providecommand{\Hyp}{\operatorname{Hyp}}
\providecommand{\isoto}{\simeq}
\ifLuaTeX
  % Xy-pic derives its PDF stroke width from \fontdimen8\textfont3.  Under
  % unicode-math/TeX Gyre Termes Math that value is 7.2489pt rather than the
  % legacy 0.43799pt, producing unusably heavy diagram arrows.  Retain the
  % established Xy-pic geometry while keeping OpenType math and ToUnicode.
  \makeatletter
  \AtBeginDocument{%
    \edef\xP@preclw{0.43799pt}%
    \edef\xP@lw{\xP@dim\xP@preclw}%
  }
  \makeatother
\fi
\usepackage[shortlabels]{enumitem}
\usepackage{longtable}
\usepackage{marginnote}
\ifLuaTeX\else\usepackage{microtype}\fi
\usepackage{fancyhdr}
\usepackage[
  a4paper,
  inner=26mm,
  outer=24mm,
  top=24mm,
  bottom=24mm,
  headheight=14pt,
  headsep=8mm,
  footskip=12mm
]{geometry}
\usepackage[
  unicode=true,
  colorlinks=true,
  linkcolor=black,
  citecolor=black,
  urlcolor=blue,
  hypertexnames=false
]{hyperref}

\hypersetup{
  pdftitle={代数几何学原理 II：几类态射的整体性质（大陆简体中文独立工作版）},
  pdfauthor={Alexander Grothendieck；Jean Dieudonné（协作编写）},
  pdfsubject={依据 NUMDAM 法文印本与锁定的法文外交式 TeX 制作的大陆简体中文版本},
  pdfkeywords={代数几何；概形；态射；EGA II；简体中文}
}

\AtBeginDocument{\fontsize{10.95pt}{16.5pt}\selectfont}
\setlength{\parindent}{2em}
\setlength{\parskip}{0pt}
\setlength{\emergencystretch}{2em}
\flushbottom

\pagestyle{fancy}
\fancyhf{}
\fancyhead[L]{\small 代数几何学原理 II}
\fancyhead[R]{\small 几类态射的整体性质}
\fancyfoot[C]{\thepage}
\renewcommand{\headrulewidth}{0.4pt}
\renewcommand{\footrulewidth}{0pt}
\fancypagestyle{plain}{%
  \fancyhf{}%
  \fancyhead[L]{\small 代数几何学原理 II}%
  \fancyhead[R]{\small 几类态射的整体性质}%
  \fancyfoot[C]{\thepage}%
  \renewcommand{\headrulewidth}{0.4pt}%
  \renewcommand{\footrulewidth}{0pt}%
}

\renewcommand{\contentsname}{目录}
\renewcommand{\bibname}{参考文献}
\renewcommand{\figurename}{图}
\renewcommand{\tablename}{表}
\renewcommand{\thesection}{\arabic{section}}
\renewcommand{\thesubsection}{\thesection.\arabic{subsection}}

\newenvironment{env}[1][]{\par\medskip\noindent\textbf{(#1)}\ }{\par}
\newenvironment{definition}[1][]{\par\medskip\noindent\textbf{定义 (#1).}\ }{\par}
\newenvironment{lemma}[1][]{\par\medskip\noindent\textbf{引理 (#1).}\ }{\par}
\newenvironment{theorem}[1][]{\par\medskip\noindent\textbf{定理 (#1).}\ }{\par}
\newenvironment{proposition}[1][]{\par\medskip\noindent\textbf{命题 (#1).}\ }{\par}
\newenvironment{corollary}[1][]{\par\medskip\noindent\textbf{推论 (#1).}\ }{\par}
\newenvironment{remark}[1][]{\par\medskip\noindent\textbf{注 (#1).}\ }{\par}
\newenvironment{namedtheorem}[2][]{\par\medskip\noindent\textbf{#2 (#1).}\ }{\par}
\newenvironment{notation}[1][]{\par\medskip\noindent\textbf{记号 (#1).}\ }{\par}
\newenvironment{example}[1][]{\par\medskip\noindent\textbf{例 (#1).}\ }{\par}
\newenvironment{proof}{\par\medskip\noindent\textit{证明.}\ }{\hfill$\square$\par}

\newcommand{\oldpage}[2][]{%
  \hypertarget{ega-source-page-#1-#2}{}%
  \marginnote{\scriptsize\textbf{#1}~|~#2}\ignorespaces}

% Cross-volume references preserve the exact source key without creating a dead link
% in the standalone EGA II reader. Internal EGA II references use \hyperref directly.
\newcommand{\egaxref}[2]{\textup{#2}}
\newcommand{\egaxcite}[2][]{%
  \textup{[#2\if\relax\detokenize{#1}\relax\else, #1\fi]}}
\makeatletter
\newcommand{\egaref}[2]{%
  \@ifundefined{r@#1}{\textup{#2}}{\hyperref[#1]{#2}}}
\makeatother

\let\EGAOriginalLabel\label
\renewcommand{\label}[1]{\phantomsection\EGAOriginalLabel{#1}}

\begin{document}
\pagestyle{fancy}
\pagenumbering{arabic}
\setcounter{page}{5}
% Locked French authority: front-fr.tex, 1423 B,
% SHA-256 F199DE0C2C8C49E2BDAEF48D764086952806E77D8066795B97B3595F7EA03198.
% Authority binding: NUMDAM leaf 2 / printed p.5.

\clearpage
\thispagestyle{plain}
\markboth{第二章}{第二章}
\oldpage[II]{5}
\phantomsection
\label{chapter:II-zh}

\begin{center}
{\large 第二章\par}
\vspace{1.5em}
{\Large\bfseries 几类态射的初等整体研究\par}
\vspace{1.8em}
{\bfseries 内容提要\par}
\end{center}

\medskip
\begin{tabular}{@{}r@{\ }p{0.82\textwidth}@{}}
\S~1. & 仿射态射。\\
\S~2. & 齐次素谱。\\
\S~3. & 分次代数层的齐次谱。\\
\S~4. & 泛射影丛；丰沛层。\\
\S~5. & 拟仿射态射；拟射影态射；紧合态射；射影态射。\\
\S~6. & 整型态射与有限态射。\\
\S~7. & 赋值判别法。\\
\S~8. & 概形的暴涨；投影锥；泛射影闭包。
\end{tabular}

\medskip
本章研究上述各类态射时，并不大量使用上同调方法；第三章将运用这些方法作更深入的研究，
其中主要使用第二章的 \S\S~2、4 和~5。初读时可以略去 \S~8：它补充
\S\S~1 至~3 所建立的形式体系，其内容归结为这一形式体系的若干简单应用；与本章其他结果相比，
后文使用这些补充结果的频率也较低。
% Locked French authority: ega2-1-fr.tex, 820505 B,
% SHA-256 84EDBE3E83530AF2959B441796337C9DC21EAFCA6A13114A26778760FBF437AC.
% Sequential source scope: lines 1-89, through 1.1.3 / printed pp.5-6.

\setcounter{section}{0}
\section{仿射态射}
\label{section:II.1-zh}

本节的大部分结果，是第一章 \S~1 中相应结果的“整体”版本；因此，它们在本质上并非新的，
只是为后文提供一种方便的语言。

\setcounter{subsection}{0}
\subsection{\texorpdfstring{$S$-预概形与 $\mathcal{O}_S$-代数层}{S-预概形与结构代数层}}
\label{subsection:II.1.1-zh}

\begin{env}[1.1.1]
\label{II.1.1.1-zh}
设 $S$ 是一个预概形，$X$ 是一个 $S$-预概形，$f:X\to S$ 是它的结构态射。
已知（\egaxref{0.4.2.4-fr}{0, 4.2.4}），直像
$f_*(\mathcal{O}_X)$ 是一个 $\mathcal{O}_S$-代数层；在不致混淆时，我们把它
\oldpage[II]{6}
记作 $\mathcal{A}(X)$。若 $U$ 是 $S$ 的一个开集，则
\[
  \mathcal{A}(f^{-1}(U))=\mathcal{A}(X)|U.
\]
同样，对任一 $\mathcal{O}_X$-模层 $\mathcal{F}$（相应地，对任一
$\mathcal{O}_X$-代数层 $\mathcal{B}$），我们把直像 $f_*(\mathcal{F})$
（相应地，$f_*(\mathcal{B})$）记作 $\mathcal{A}(\mathcal{F})$
（相应地，$\mathcal{A}(\mathcal{B})$）；它是一个 $\mathcal{A}(X)$-模层
（相应地，一个 $\mathcal{A}(X)$-代数层），而不只是一个 $\mathcal{O}_S$-模层
（相应地，一个 $\mathcal{O}_S$-代数层）。
\end{env}

\begin{env}[1.1.2]
\label{II.1.1.2-zh}
设 $Y$ 是第二个 $S$-预概形，$g:Y\to S$ 是它的结构态射，而
$h:X\to Y$ 是一个 $S$-态射；于是有交换图
\[
  \xymatrix{
    X\ar[rr]^h\ar[rd]_f && Y\ar[ld]^g\\
    & S
  }
  \tag{1.1.2.1}\label{II.1.1.2.1-zh}
\]

按定义，$h=(\psi,\theta)$，其中
$\theta:\mathcal{O}_Y\to h_*(\mathcal{O}_X)=\psi_*(\mathcal{O}_X)$
是一个环层同态；由此（\egaxref{0.4.2.2-fr}{0, 4.2.2}）得到一个
$\mathcal{O}_S$-代数层同态
\[
  g_*(\theta):g_*(\mathcal{O}_Y)\to
  g_*(h_*(\mathcal{O}_X))=f_*(\mathcal{O}_X),
\]
换言之，是一个 $\mathcal{O}_S$-代数层同态
$\mathcal{A}(Y)\to\mathcal{A}(X)$，我们也把它记作 $\mathcal{A}(h)$。
若 $h':Y\to Z$ 是第二个 $S$-态射，则显然有
$\mathcal{A}(h'\circ h)=\mathcal{A}(h)\circ\mathcal{A}(h')$。
这样，我们就在 $S$-预概形范畴上定义了一个取值于
$\mathcal{O}_S$-代数层范畴的反变函子 $X\mapsto\mathcal{A}(X)$。

现在设 $\mathcal{F}$ 是一个 $\mathcal{O}_X$-模层，$\mathcal{G}$ 是一个
$\mathcal{O}_Y$-模层，而 $u:\mathcal{G}\to\mathcal{F}$ 是一个
$h$-态射，也就是说（\egaxref{0.4.4.1-fr}{0, 4.4.1}），它是一个
$\mathcal{O}_Y$-模层同态 $\mathcal{G}\to h_*(\mathcal{F})$。于是
\[
  g_*(u):g_*(\mathcal{G})\to g_*(h_*(\mathcal{F}))=f_*(\mathcal{F})
\]
是一个 $\mathcal{O}_S$-模层同态
$\mathcal{A}(\mathcal{G})\to\mathcal{A}(\mathcal{F})$，我们把它记作
$\mathcal{A}(u)$；此外，二元组
$(\mathcal{A}(h),\mathcal{A}(u))$ 构成从 $\mathcal{A}(Y)$-模层
$\mathcal{A}(\mathcal{G})$ 到 $\mathcal{A}(X)$-模层
$\mathcal{A}(\mathcal{F})$ 的一个\emph{双同态}。
\end{env}

\begin{env}[1.1.3]
\label{II.1.1.3-zh}
固定预概形 $S$。可以把二元组 $(X,\mathcal{F})$ 看成一个\emph{范畴}的对象，
其中 $X$ 是一个 $S$-预概形，$\mathcal{F}$ 是一个 $\mathcal{O}_X$-模层；
从 $(X,\mathcal{F})$ 到 $(Y,\mathcal{G})$ 的一个\emph{态射}定义为二元组
$(h,u)$，其中 $h:X\to Y$ 是一个 $S$-态射，而
$u:\mathcal{G}\to\mathcal{F}$ 是一个 $h$-态射。于是可以说，赋值
$(X,\mathcal{F})\mapsto(\mathcal{A}(X),\mathcal{A}(\mathcal{F}))$
是一个\emph{反变函子}；其值域范畴的对象，是由一个 $\mathcal{O}_S$-代数层
及其上的一个模层组成的二元组，而态射则是双同态。
\end{env}
% Locked French authority: ega2-1-fr.tex, 820505 B,
% SHA-256 84EDBE3E83530AF2959B441796337C9DC21EAFCA6A13114A26778760FBF437AC.
% Sequential source scope: lines 91-232, subsection 1.2 / printed pp.6-8.

\subsection{相对仿射的预概形}
\label{subsection:II.1.2-zh}

\begin{definition}[1.2.1]
\label{II.1.2.1-zh}
设 $X$ 是一个 $S$-预概形，$f:X\to S$ 是它的结构态射。若存在 $S$ 的一个仿射开集覆盖
$(S_\alpha)$，使得对每个 $\alpha$，$X$ 在开集 $f^{-1}(S_\alpha)$ 上诱导的预概形都是仿射的，
则称 $X$ \emph{相对于 $S$ 仿射}。
\end{definition}

\begin{example}[1.2.2]
\label{II.1.2.2-zh}
$S$ 的每个闭子概形都是一个相对于 $S$ 仿射的 $S$-预概形
（\egaxref{I.4.2.3-fr}{I, 4.2.3} 及 \egaxref{I.4.2.4-fr}{4.2.4}）。
\end{example}

\begin{remark}[1.2.3]
\label{II.1.2.3-zh}
相对于 $S$ 仿射的预概形 $X$ 不一定是仿射概形，例 $X=S$
（\egaref{II.1.2.2-zh}{1.2.2}）即说明了这一点。另一方面，若仿射概形 $X$ 是一个
$S$-预概形，则 $X$ 不一定相对于
\oldpage[II]{7}
$S$ 仿射（见 \egaref{II.1.3.3-zh}{1.3.3}）。不过要回顾，若 $S$ 是一个\emph{概形}，
则作为仿射概形的每个 $S$-预概形都相对于 $S$ 仿射
（\egaxref{I.5.5.10-fr}{I, 5.5.10}）。
\end{remark}

\begin{proposition}[1.2.4]
\label{II.1.2.4-zh}
每个相对于 $S$ 仿射的 $S$-预概形都相对于 $S$ 分离（换言之，它是一个 $S$-概形）。
\end{proposition}

这立即由 \egaxref{I.5.5.5-fr}{I, 5.5.5} 和
\egaxref{I.5.5.8-fr}{I, 5.5.8} 得出。

\begin{proposition}[1.2.5]
\label{II.1.2.5-zh}
设 $X$ 是一个相对于 $S$ 仿射的 $S$-概形，$f:X\to S$ 是结构态射。对每个开集
$U\subset S$，$f^{-1}(U)$ 都相对于 $U$ 仿射。
\end{proposition}

根据定义 \egaref{II.1.2.1-zh}{1.2.1}，可以归约到
$S=\operatorname{Spec}(A)$ 与 $X=\operatorname{Spec}(B)$ 均为仿射概形的情形；此时
$f=({}^a\varphi,\widetilde{\varphi})$，其中 $\varphi:A\to B$ 是一个同态。
由于 $g\in A$ 时的 $D(g)$ 构成 $S$ 的一个拓扑基，只需考察 $U=D(g)$ 的情形；而这时已知
$f^{-1}(U)=D(\varphi(g))$（\egaxref{I.1.2.2.2-fr}{I, 1.2.2.2}），命题得证。

\begin{proposition}[1.2.6]
\label{II.1.2.6-zh}
设 $X$ 是一个相对于 $S$ 仿射的 $S$-概形，$f:X\to S$ 是结构态射。对每个拟凝聚
$\mathcal{O}_X$-模层 $\mathcal{F}$，$f_*(\mathcal{F})$ 都是一个拟凝聚
$\mathcal{O}_S$-模层。
\end{proposition}

结合 \egaref{II.1.2.4-zh}{1.2.4}，这由
\egaxref{I.9.2.2-fr}{I, 9.2.2, a)} 得出。

特别地，$\mathcal{O}_S$-代数层
$\mathcal{A}(X)=f_*(\mathcal{O}_X)$ 是\emph{拟凝聚的}。

\begin{proposition}[1.2.7]
\label{II.1.2.7-zh}
设 $X$ 是一个相对于 $S$ 仿射的 $S$-概形。对每个 $S$-预概形 $Y$，
$h\mapsto\mathcal{A}(h)$ 从集合 $\operatorname{Hom}_S(Y,X)$ 到集合
$\operatorname{Hom}(\mathcal{A}(X),\mathcal{A}(Y))$ 的映射
（\egaref{II.1.1.2-zh}{1.1.2}）是双射。
\end{proposition}

设 $f:X\to S$ 与 $g:Y\to S$ 是结构态射。先假设
$S=\operatorname{Spec}(A)$ 与 $X=\operatorname{Spec}(B)$ 均为仿射概形；需要证明，对每个
$\mathcal{O}_S$-代数层同态
$\omega:f_*(\mathcal{O}_X)\to g_*(\mathcal{O}_Y)$，存在唯一的 $S$-态射
$h:Y\to X$，使得 $\mathcal{A}(h)=\omega$。按定义，对每个开集 $U\subset S$，
$\omega$ 给出一个 $\Gamma(U,\mathcal{O}_S)$-代数同态
\[
  \omega_U=\Gamma(U,\omega):
  \Gamma(f^{-1}(U),\mathcal{O}_X)\to
  \Gamma(g^{-1}(U),\mathcal{O}_Y).
\]
特别地，取 $U=S$，便得到一个 $\Gamma(S,\mathcal{O}_S)$-代数同态
$\varphi:\Gamma(X,\mathcal{O}_X)\to\Gamma(Y,\mathcal{O}_Y)$；由于 $X$ 是仿射的，
它对应于一个唯一确定的 $S$-态射 $h:Y\to X$
（\egaxref{I.2.2.4-fr}{I, 2.2.4}）。尚需证明 $\mathcal{A}(h)=\omega$；换言之，
对 $S$ 的某个拓扑基中的每个开集 $U$，$\omega_U$ 都等于代数同态 $\varphi_U$，其中
$\varphi_U$ 对应于 $h$ 的限制 $g^{-1}(U)\to f^{-1}(U)$。

只需考察 $U=D(\lambda)$ 且 $\lambda\in A$ 的情形\footnote{法文印本在此写作
$\lambda\in S$。但本段已设 $S=\operatorname{Spec}(A)$，而且下文令
$\mu=\rho(\lambda)$、$\rho:A\to B$，故类型上必须有 $\lambda\in A$。本版把这一可逆的读者层校正
明确标出，并未改动法文外交式来源。}；若
$f=({}^a\rho,\widetilde{\rho})$，其中 $\rho:A\to B$ 是环同态，则
$f^{-1}(U)=D(\mu)$，其中 $\mu=\rho(\lambda)$，而
$\Gamma(f^{-1}(U),\mathcal{O}_X)$ 是分式环 $B_\mu$。图
\[
  \xymatrix{
    B\ar[r]^\varphi\ar[d] & \Gamma(Y,\mathcal{O}_Y)\ar[d]\\
    B_\mu\ar[r]_{\varphi_U} & \Gamma(g^{-1}(U),\mathcal{O}_Y)
  }
\]
是交换的；把 $\varphi_U$ 换成 $\omega_U$ 后所得的类似图也交换。因此，由分式环的泛性质
（\egaxref{0.1.2.4-fr}{0, 1.2.4}），得到 $\varphi_U=\omega_U$。

现在考察一般情形。取 $S$ 的一个仿射开集覆盖 $(S_\alpha)$，
\oldpage[II]{8}
使每个
$f^{-1}(S_\alpha)$ 都是仿射的。于是，任一 $\mathcal{O}_S$-代数层同态
$\omega:\mathcal{A}(X)\to\mathcal{A}(Y)$ 通过限制给出一族
\[
  \omega_\alpha:
  \mathcal{A}(f^{-1}(S_\alpha))\to
  \mathcal{A}(g^{-1}(S_\alpha))
\]
$\mathcal{O}_{S_\alpha}$-代数层同态；由前述结果，它们给出一族 $S_\alpha$-态射
$h_\alpha:g^{-1}(S_\alpha)\to f^{-1}(S_\alpha)$。现在只需说明：对
$S_\alpha\cap S_\beta$ 的某个拓扑基中的每个仿射开集 $U$，$h_\alpha$ 与
$h_\beta$ 在 $g^{-1}(U)$ 上的限制相同。这是显然的，因为按上述结果，这两个限制都对应于
$\omega$ 的限制 $\mathcal{A}(X)|U\to\mathcal{A}(Y)|U$。

\begin{corollary}[1.2.8]
\label{II.1.2.8-zh}
设 $X$、$Y$ 是两个相对于 $S$ 仿射的 $S$-概形。一个 $S$-态射 $h:Y\to X$ 是同构，
当且仅当 $\mathcal{A}(h):\mathcal{A}(X)\to\mathcal{A}(Y)$ 是同构。
\end{corollary}

这立即由 \egaref{II.1.2.7-zh}{1.2.7} 以及 $\mathcal{A}(X)$ 的函子性得出。
% Locked French authority: ega2-1-fr.tex, 820505 B,
% SHA-256 84EDBE3E83530AF2959B441796337C9DC21EAFCA6A13114A26778760FBF437AC.
% Sequential source scope: lines 234-369, subsection 1.3 / printed pp.8-9.

\subsection{与 \texorpdfstring{$\mathcal{O}_S$}{O-S}-代数层相伴的相对仿射预概形}
\label{subsection:II.1.3-zh}

\begin{proposition}[1.3.1]
\label{II.1.3.1-zh}
设 $S$ 是一个预概形。对每个拟凝聚 $\mathcal{O}_S$-代数层 $\mathcal{B}$，存在一个相对于
$S$ 仿射的预概形 $X$，它在唯一的 $S$-同构意义下确定，并满足
$\mathcal{A}(X)=\mathcal{B}$。
\end{proposition}

唯一性由 \egaref{II.1.2.8-zh}{1.2.8} 得出；下面证明 $X$ 的存在性。对每个仿射开集
$U\subset S$，令 $X_U$ 为预概形 $\operatorname{Spec}(\Gamma(U,\mathcal{B}))$；由于
$\Gamma(U,\mathcal{B})$ 是一个 $\Gamma(U,\mathcal{O}_S)$-代数，$X_U$ 是一个
$S$-预概形（\egaxref{I.1.6.1-fr}{I, 1.6.1}）。此外，由于 $\mathcal{B}$ 是拟凝聚的，
$\mathcal{O}_S$-代数层 $\mathcal{A}(X_U)$ 与 $\mathcal{B}|U$ 典范同构
（\egaxref{I.1.3.7-fr}{I, 1.3.7}、\egaxref{I.1.3.13-fr}{1.3.13} 及
\egaxref{I.1.6.3-fr}{1.6.3}）。设 $V$ 是 $S$ 的第二个仿射开集，并以 $f_U$ 表示结构态射
$X_U\to S$；令 $X_{U,V}$ 为 $X_U$ 在 $f_U^{-1}(U\cap V)$ 上诱导的预概形。
$X_{U,V}$ 与 $X_{V,U}$ 都相对于 $U\cap V$ 仿射
（\egaref{II.1.2.5-zh}{1.2.5}），而按定义，$\mathcal{A}(X_{U,V})$ 与
$\mathcal{A}(X_{V,U})$ 都同 $\mathcal{B}|(U\cap V)$ 典范同构。因此
（\egaref{II.1.2.8-zh}{1.2.8}），存在典范的 $S$-同构
$\theta_{U,V}:X_{V,U}\to X_{U,V}$。

再设 $W$ 是 $S$ 的第三个仿射开集，并令 $\theta'_{U,V}$、$\theta'_{V,W}$、
$\theta'_{U,W}$ 分别为 $\theta_{U,V}$、$\theta_{V,W}$、$\theta_{U,W}$ 在相应结构态射下
$U\cap V\cap W$ 的逆像上的限制；这些逆像分别位于 $X_V$、$X_W$、$X_W$ 中。于是
$\theta'_{U,V}\circ\theta'_{V,W}=\theta'_{U,W}$。因此，存在一个预概形 $X$、一个由仿射开集
$(T_U)$ 构成的 $X$ 的覆盖，以及对每个 $U$ 的同构 $\varphi_U:X_U\to T_U$，使得
$\varphi_U$ 把 $f_U^{-1}(U\cap V)$ 映到 $T_U\cap T_V$，并且
$\theta_{U,V}=\varphi_U^{-1}\circ\varphi_V$
（\egaxref{I.2.3.1-fr}{I, 2.3.1}）。态射
$g_U=f_U\circ\varphi_U^{-1}$ 使 $T_U$ 成为一个 $S$-预概形；$g_U$ 与 $g_V$ 在
$T_U\cap T_V$ 上一致，故 $X$ 是一个 $S$-预概形。此外，由定义显然可见，$X$ 相对于
$S$ 仿射且 $\mathcal{A}(T_U)=\mathcal{B}|U$，因此
$\mathcal{A}(X)=\mathcal{B}$。

称这样定义的 $S$-概形 $X$ \emph{与 $\mathcal{O}_S$-代数层 $\mathcal{B}$ 相伴}，
或称它为 $\mathcal{B}$ 的\emph{谱}，并记作 $\operatorname{Spec}(\mathcal{B})$。

\begin{corollary}[1.3.2]
\label{II.1.3.2-zh}
设 $X$ 是一个相对于 $S$ 仿射的预概形，$f:X\to S$ 是结构态射。对每个仿射开集
$U\subset S$，在 $f^{-1}(U)$ 上诱导的预概形都是一个仿射概形，其环为
$\Gamma(U,\mathcal{A}(X))$。
\end{corollary}

\oldpage[II]{9}
由 \egaref{II.1.2.6-zh}{1.2.6} 和 \egaref{II.1.3.1-zh}{1.3.1}，可以假设
$X$ 与某个 $\mathcal{O}_S$-代数层相伴；于是本推论由
\egaref{II.1.3.1-zh}{1.3.1} 中所述的 $X$ 的构造得出。

\begin{example}[1.3.3]
\label{II.1.3.3-zh}
设 $S$ 是域 $K$ 上把点 $0$ 加倍所得的仿射平面
（\egaxref{I.5.5.11-fr}{I, 5.5.11}）。沿用同一处
（\egaxref{I.5.5.11-fr}{I, 5.5.11}）的记号，$S$ 是两个仿射开集
$Y_1$、$Y_2$ 的并。若 $f$ 是开浸入 $Y_1\to S$，则
$f^{-1}(Y_2)=Y_1\cap Y_2$ 不是 $Y_1$ 中的仿射开集（同上）。这就给出了一个不相对于
$S$ 仿射的仿射概形的例子。
\end{example}

\begin{corollary}[1.3.4]
\label{II.1.3.4-zh}
设 $S$ 是一个仿射概形。一个 $S$-预概形 $X$ 相对于 $S$ 仿射，当且仅当 $X$ 是仿射概形。
\end{corollary}

\begin{corollary}[1.3.5]
\label{II.1.3.5-zh}
设 $X$ 是一个相对于预概形 $S$ 仿射的预概形，$Y$ 是一个 $X$-预概形。则 $Y$ 相对于
$S$ 仿射，当且仅当 $Y$ 相对于 $X$ 仿射。
\end{corollary}

立即可归约到 $S$ 是仿射概形的情形；此时 $X$ 也同样是仿射概形
（\egaref{II.1.3.4-zh}{1.3.4}）。命题中的两个条件此时都表示 $Y$ 是仿射概形
（\egaref{II.1.3.4-zh}{1.3.4}）。

\begin{env}[1.3.6]
\label{II.1.3.6-zh}
设 $X$ 是一个相对于 $S$ 仿射的预概形。根据 \egaref{II.1.3.5-zh}{1.3.5}，为了定义一个
\emph{相对于 $X$ 仿射}的预概形 $Y$，等价地可以给出一个\emph{相对于 $S$ 仿射}的预概形
$Y$ 以及一个 $S$-态射 $g:Y\to X$；换言之
（\egaref{II.1.3.1-zh}{1.3.1} 及 \egaref{II.1.2.7-zh}{1.2.7}），这等价于给出
一个拟凝聚 $\mathcal{O}_S$-代数层 $\mathcal{B}$，以及一个 $\mathcal{O}_S$-代数层同态
$\mathcal{A}(X)\to\mathcal{B}$（也可以把它看作在 $\mathcal{B}$ 上定义了一个
$\mathcal{A}(X)$-代数层结构）。若 $f:X\to S$ 是结构态射，则此时有
$\mathcal{B}=f_*(g_*(\mathcal{O}_Y))$。
\end{env}

\begin{corollary}[1.3.7]
\label{II.1.3.7-zh}
设 $X$ 是一个相对于 $S$ 仿射的预概形。则 $X$ 在 $S$ 上为有限型，当且仅当拟凝聚
$\mathcal{O}_S$-代数层 $\mathcal{A}(X)$ 为有限型
（\egaxref{I.9.6.2-fr}{I, 9.6.2}）。
\end{corollary}

按定义（\egaxref{I.9.6.2-fr}{I, 9.6.2}），可以归约到 $S$ 仿射的情形；此时 $X$ 是仿射概形
（\egaref{II.1.3.4-zh}{1.3.4}）。若 $S=\operatorname{Spec}(A)$、
$X=\operatorname{Spec}(B)$，则 $\mathcal{A}(X)$ 是 $\mathcal{O}_S$-代数层
$\widetilde{B}$；由于 $\Gamma(U,\widetilde{B})=B$，本推论由
\egaxref{I.9.6.2-fr}{I, 9.6.2} 和 \egaxref{I.6.3.3-fr}{I, 6.3.3} 得出。

\begin{corollary}[1.3.8]
\label{II.1.3.8-zh}
设 $X$ 是一个相对于 $S$ 仿射的预概形。则 $X$ 是约化的，当且仅当拟凝聚
$\mathcal{O}_S$-代数层 $\mathcal{A}(X)$ 是约化的
（\egaxref{0.4.1.4-fr}{0, 4.1.4}）。
\end{corollary}

事实上，这个问题显然在 $S$ 上是局部的；根据 \egaref{II.1.3.2-zh}{1.3.2}，本推论由
\egaxref{I.5.1.1-fr}{I, 5.1.1} 和 \egaxref{I.5.1.4-fr}{I, 5.1.4} 得出。
% Locked French authority: ega2-1-fr.tex, 820505 B,
% SHA-256 84EDBE3E83530AF2959B441796337C9DC21EAFCA6A13114A26778760FBF437AC.
% Sequential source scope: lines 371-622, subsection 1.4 / printed pp.9-12.

\subsection{相对于 \texorpdfstring{$S$}{S} 仿射的预概形上的拟凝聚层}
\label{subsection:II.1.4-zh}

\begin{proposition}[1.4.1]
\label{II.1.4.1-zh}
设 $X$ 是一个相对于 $S$ 仿射的预概形，$Y$ 是一个 $S$-预概形，
$\mathcal{F}$（相应地，$\mathcal{G}$）是一个拟凝聚 $\mathcal{O}_X$-模层
（相应地，$\mathcal{O}_Y$-模层）。则映射
$(h,u)\mapsto(\mathcal{A}(h),\mathcal{A}(u))$ 从态射
$(Y,\mathcal{G})\to(X,\mathcal{F})$ 的集合到双同态
$(\mathcal{A}(X),\mathcal{A}(\mathcal{F}))\to
(\mathcal{A}(Y),\mathcal{A}(\mathcal{G}))$ 的集合是双射
（\egaref{II.1.1.2-zh}{1.1.2} 及
\egaref{II.1.1.3-zh}{1.1.3}）。
\end{proposition}

证明完全沿用 \egaref{II.1.2.7-zh}{1.2.7} 的步骤，只需使用
\egaxref{I.2.2.5-fr}{I, 2.2.5} 和 \egaxref{I.2.2.4-fr}{I, 2.2.4}；
细节留给读者。

\begin{corollary}[1.4.2]
\label{II.1.4.2-zh}
除 \egaref{II.1.4.1-zh}{1.4.1} 的假设外，再假设 $Y$ 相对于 $S$ 仿射。
则 $(h,u)$ 是同构，当且仅当 $(\mathcal{A}(h),\mathcal{A}(u))$ 是双同构。
\end{corollary}

\oldpage[II]{10}
\begin{proposition}[1.4.3]
\label{II.1.4.3-zh}
对任意一个由拟凝聚 $\mathcal{O}_S$-代数层 $\mathcal{B}$ 和拟凝聚
$\mathcal{B}$-模层 $\mathcal{M}$ 组成的偶对 $(\mathcal{B},\mathcal{M})$
（$\mathcal{M}$ 作为 $\mathcal{O}_S$-模层或作为 $\mathcal{B}$-模层是拟凝聚的，
这两者等价；见 \egaxref{I.9.6.1-fr}{I, 9.6.1}），存在一个由相对于 $S$ 仿射的
预概形 $X$ 和拟凝聚 $\mathcal{O}_X$-模层 $\mathcal{F}$ 组成的偶对
$(X,\mathcal{F})$，使得 $\mathcal{A}(X)=\mathcal{B}$ 且
$\mathcal{A}(\mathcal{F})=\mathcal{M}$；而且该偶对在唯一同构意义下确定。
\end{proposition}

唯一性由 \egaref{II.1.4.1-zh}{1.4.1} 和 \egaref{II.1.4.2-zh}{1.4.2}
得出；存在性的证明与 \egaref{II.1.3.1-zh}{1.3.1} 相同，细节仍留给读者。
把 $\mathcal{O}_X$-模层 $\mathcal{F}$ 记作 $\widetilde{\mathcal{M}}$，并称它
\emph{与}拟凝聚 $\mathcal{B}$-模层 $\mathcal{M}$\emph{相伴}。对每个仿射开集
$U\subset S$，$\widetilde{\mathcal{M}}|p^{-1}(U)$（其中 $p$ 是结构态射
$X\to S$）典范同构于 $(\Gamma(U,\mathcal{M}))^{\sim}$。

\begin{corollary}[1.4.4]
\label{II.1.4.4-zh}
在拟凝聚 $\mathcal{B}$-模层的范畴中，$\mathcal{M}\mapsto
\widetilde{\mathcal{M}}$ 是协变加性正合函子，并且与归纳极限和直和可交换。
\end{corollary}

事实上，立即可归约到 $S$ 仿射的情形，此时本推论由
\egaxref{I.1.3.5-fr}{I, 1.3.5}、\egaxref{I.1.3.9-fr}{I, 1.3.9} 和
\egaxref{I.1.3.11-fr}{I, 1.3.11} 得出。

\begin{corollary}[1.4.5]
\label{II.1.4.5-zh}
在 \egaref{II.1.4.3-zh}{1.4.3} 的假设下，$\widetilde{\mathcal{M}}$ 是有限型
$\mathcal{O}_X$-模层，当且仅当 $\mathcal{M}$ 是有限型 $\mathcal{B}$-模层。
\end{corollary}

立即可归约到 $S=\operatorname{Spec}(A)$ 是仿射概形的情形。此时
$\mathcal{B}=\widetilde{B}$，其中 $B$ 是一个有限型 $A$-代数
（\egaxref{I.9.6.2-fr}{I, 9.6.2}），且 $\mathcal{M}=\widetilde{M}$，
其中 $M$ 是一个 $B$-模（\egaxref{I.1.3.13-fr}{I, 1.3.13}）。
\emph{在预概形 $X$ 上}，$\mathcal{O}_X$ 与环 $B$ 相伴，而
$\widetilde{\mathcal{M}}$ 与 $B$-模 $M$ 相伴。因此，
$\widetilde{\mathcal{M}}$ 是有限型的，当且仅当 $M$ 是有限型的
（\egaxref{I.1.3.13-fr}{I, 1.3.13}），这就证明了断言。

\begin{proposition}[1.4.6]
\label{II.1.4.6-zh}
设 $Y$ 是一个相对于 $S$ 仿射的预概形，$X$、$X'$ 是两个相对于 $Y$ 仿射的
预概形（因而也相对于 $S$ 仿射；见 \egaref{II.1.3.5-zh}{1.3.5}）。令
$\mathcal{B}=\mathcal{A}(Y)$、$\mathcal{A}=\mathcal{A}(X)$、
$\mathcal{A}'=\mathcal{A}(X')$。
则 $X\times_Y X'$ 相对于 $Y$ 仿射（因而也相对于 $S$ 仿射），并且
$\mathcal{A}(X\times_Y X')$ 典范同构于
$\mathcal{A}\otimes_{\mathcal{B}}\mathcal{A}'$。
\end{proposition}

事实上（\egaxref{I.9.6.1-fr}{I, 9.6.1}），
$\mathcal{A}\otimes_{\mathcal{B}}\mathcal{A}'$ 是一个拟凝聚
$\mathcal{B}$-代数层，因而也是一个拟凝聚 $\mathcal{O}_S$-代数层
（\egaxref{I.9.6.1-fr}{I, 9.6.1}）。令
$Z=\operatorname{Spec}(\mathcal{A}\otimes_{\mathcal{B}}\mathcal{A}')$
（\egaref{II.1.3.1-zh}{1.3.1}）。典范的 $\mathcal{B}$-同态
$\mathcal{A}\to\mathcal{A}\otimes_{\mathcal{B}}\mathcal{A}'$ 和
$\mathcal{A}'\to\mathcal{A}\otimes_{\mathcal{B}}\mathcal{A}'$ 分别对应于
（\egaref{II.1.2.7-zh}{1.2.7}）$Y$-态射 $p:Z\to X$ 和 $p':Z\to X'$。
为证明三元组 $(Z,p,p')$ 是积 $X\times_Y X'$，只需考虑 $S$ 是环 $C$ 的仿射概形
这一情形（\egaxref{I.3.2.6.4-fr}{I, 3.2.6.4}）。这时 $Y$、$X$、$X'$ 都是
仿射概形（\egaref{II.1.3.4-zh}{1.3.4}），其环 $B$、$A$、$A'$ 是 $C$-代数，
并满足 $\mathcal{B}=\widetilde{B}$、$\mathcal{A}=\widetilde{A}$、
$\mathcal{A}'=\widetilde{A'}$。已知（\egaxref{I.1.3.13-fr}{I, 1.3.13}），
$\mathcal{A}\otimes_{\mathcal{B}}\mathcal{A}'$ 典范同构于
$\mathcal{O}_S$-代数层 $\widetilde{A\otimes_B A'}$；所以环 $A(Z)$ 同构于
$A\otimes_B A'$，而态射 $p$、$p'$ 分别对应于典范同态
$A\to A\otimes_B A'$、$A'\to A\otimes_B A'$。本命题遂由
\egaxref{I.3.2.2-fr}{I, 3.2.2} 得出。

\begin{corollary}[1.4.7]
\label{II.1.4.7-zh}
设 $\mathcal{F}$（相应地，$\mathcal{F}'$）是拟凝聚 $\mathcal{O}_X$-模层
（相应地，$\mathcal{O}_{X'}$-模层）。则
$\mathcal{A}(\mathcal{F}\otimes_Y\mathcal{F}')$ 典范同构于
$\mathcal{A}(\mathcal{F})\otimes_{\mathcal{A}(Y)}
\mathcal{A}(\mathcal{F}')$。
\end{corollary}

已知 $\mathcal{F}\otimes_Y\mathcal{F}'$ 在 $X\times_Y X'$ 上是拟凝聚的
（\egaxref{I.9.1.2-fr}{I, 9.1.2}）。设 $g:Y\to S$、$f:X\to Y$、
$f':X'\to Y$ 为结构态射，从而结构态射
\oldpage[II]{11}
$h:Z\to S$ 同时等于 $g\circ f\circ p$ 和 $g\circ f'\circ p'$。
按如下方式定义一个典范同态
\[
  \mathcal{A}(\mathcal{F})\otimes_{\mathcal{A}(Y)}\mathcal{A}(\mathcal{F}')
  \longrightarrow \mathcal{A}(\mathcal{F}\otimes_Y\mathcal{F}'):
\]
对每个开集 $U\subset S$，存在典范同态
\[
  \begin{aligned}
  \Gamma(f^{-1}(g^{-1}(U)),\mathcal{F})
  &\longrightarrow \Gamma(h^{-1}(U),p^*(\mathcal{F})),\quad\text{以及}\\
  \Gamma(f^{\prime-1}(g^{-1}(U)),\mathcal{F}')
  &\longrightarrow \Gamma(h^{-1}(U),p^{\prime*}(\mathcal{F}')).
  \end{aligned}
\]
由 \egaxref{0.4.4.3-fr}{0, 4.4.3} 可从中导出典范同态
\[
  \begin{aligned}
  &\Gamma(f^{-1}(g^{-1}(U)),\mathcal{F})
    \otimes_{\Gamma(g^{-1}(U),\mathcal{O}_Y)}
    \Gamma(f^{\prime-1}(g^{-1}(U)),\mathcal{F}')\longrightarrow\\
  &\hspace{17em}
    \Gamma(h^{-1}(U),p^*(\mathcal{F}))
    \otimes_{\Gamma(h^{-1}(U),\mathcal{O}_Z)}
    \Gamma(h^{-1}(U),p^{\prime*}(\mathcal{F}')).
  \end{aligned}
\]

为证明所得的确是 $\mathcal{A}(Z)$-模层的同构，只需考虑 $S$（从而 $X$、$X'$、
$Y$、$X\times_Y X'$）都是仿射概形的情形。采用
\egaref{II.1.4.6-zh}{1.4.6} 的记号，令
$\mathcal{F}=\widetilde{M}$、$\mathcal{F}'=\widetilde{M'}$，其中 $M$
（相应地，$M'$）是 $A$-模（相应地，$A'$-模）。此时
$\mathcal{F}\otimes_Y\mathcal{F}'$ 同构于 $X\times_Y X'$ 上与
$(A\otimes_B A')$-模 $M\otimes_B M'$ 相伴的层
（\egaxref{I.9.1.3-fr}{I, 9.1.3}）。本推论于是由 $\mathcal{O}_S$-模层
$\widetilde{M\otimes_B M'}$ 与
$\widetilde{M}\otimes_{\widetilde{B}}\widetilde{M'}$ 的典范同构得出；
这里 $M$、$M'$ 和 $B$ 均视为 $C$-模
（\egaxref{I.1.3.13-fr}{I, 1.3.13} 及
\egaxref{I.1.6.3-fr}{I, 1.6.3}）。

特别地，在 \egaref{II.1.4.7-zh}{1.4.7} 中取 $X=Y$ 且
$\mathcal{F}'=\mathcal{O}_{X'}$，可见 $\mathcal{A}'$-模层
$\mathcal{A}(f^{\prime*}(\mathcal{F}))$ 同构于
$\mathcal{A}(\mathcal{F})\otimes_{\mathcal{B}}\mathcal{A}'$。

\begin{env}[1.4.8]
\label{II.1.4.8-zh}
特别地，当 $X=X'=Y$（$X$ 相对于 $S$ 仿射）时，若 $\mathcal{F}$、
$\mathcal{G}$ 是两个拟凝聚 $\mathcal{O}_X$-模层，则在函子性典范同构意义下有
\[
  \mathcal{A}(\mathcal{F}\otimes_{\mathcal{O}_X}\mathcal{G})
  =\mathcal{A}(\mathcal{F})\otimes_{\mathcal{A}(X)}\mathcal{A}(\mathcal{G})
  \tag{1.4.8.1}\label{II.1.4.8.1-zh}
\]
。再若 $\mathcal{F}$ 有有限表示，则由
\egaxref{I.1.6.3-fr}{I, 1.6.3} 和 \egaxref{I.1.3.12-fr}{I, 1.3.12}
可得典范同构
\[
  \mathcal{A}(\mathcal{H}om_X(\mathcal{F},\mathcal{G}))
  =\mathcal{H}om_{\mathcal{A}(X)}
    (\mathcal{A}(\mathcal{F}),\mathcal{A}(\mathcal{G}))
  \tag{1.4.8.2}\label{II.1.4.8.2-zh}
\]
。
\end{env}

\begin{remark}[1.4.9]
\label{II.1.4.9-zh}
若 $X$、$X'$ 是两个相对于 $S$ 仿射的预概形，则和 $X\sqcup X'$ 也相对于
$S$ 仿射，因为两个仿射概形的和仍是仿射概形。
\end{remark}

\begin{proposition}[1.4.10]
\label{II.1.4.10-zh}
设 $S$ 是一个预概形，$\mathcal{B}$ 是一个拟凝聚 $\mathcal{O}_S$-代数层，
$X=\operatorname{Spec}(\mathcal{B})$。对 $\mathcal{B}$ 的每个拟凝聚理想层
$\mathcal{J}$，$\widetilde{\mathcal{J}}$ 是 $\mathcal{O}_X$ 的拟凝聚理想层，
而由 $\widetilde{\mathcal{J}}$ 定义的 $X$ 的闭子预概形 $Y$ 典范同构于
$\operatorname{Spec}(\mathcal{B}/\mathcal{J})$。
\end{proposition}

事实上，由 \egaxref{I.4.2.3-fr}{I, 4.2.3} 立即可知 $Y$ 相对于 $S$ 仿射；
根据 \egaref{II.1.3.1-zh}{1.3.1}，因此可归约到 $S$ 仿射的情形，本命题遂由
\egaxref{I.4.1.2-fr}{I, 4.1.2} 立即得出。

还可把 \egaref{II.1.4.10-zh}{1.4.10} 的结果表述为：若
$h:\mathcal{B}\to\mathcal{B}'$ 是拟凝聚 $\mathcal{O}_S$-代数层的
\emph{满}同态，则 ${}^a h:\operatorname{Spec}(\mathcal{B}')\to
\operatorname{Spec}(\mathcal{B})$ 是一个\emph{闭浸入}。

\oldpage[II]{12}
\begin{proposition}[1.4.11]
\label{II.1.4.11-zh}
设 $S$ 是一个预概形，$\mathcal{B}$ 是一个拟凝聚 $\mathcal{O}_S$-代数层，
$X=\operatorname{Spec}(\mathcal{B})$。对 $\mathcal{O}_S$ 的每个拟凝聚理想层
$\mathcal{K}$，若 $f:X\to S$ 是结构态射，则在典范同构意义下有
$f^*(\mathcal{K})\mathcal{O}_X=\widetilde{\mathcal{K}\mathcal{B}}$。
\end{proposition}

这个问题在 $S$ 上是局部的，故可归约到 $S=\operatorname{Spec}(A)$ 为仿射概形的
情形；此时本命题正是 \egaxref{I.1.6.9-fr}{I, 1.6.9}。
% Locked French authority: ega2-1-fr.tex, 820505 B,
% SHA-256 84EDBE3E83530AF2959B441796337C9DC21EAFCA6A13114A26778760FBF437AC.
% Sequential source scope: lines 624-800, subsection 1.5 / printed pp.12-14.

\subsection{基预概形的换基}
\label{II.1.5-zh}

\begin{proposition}[1.5.1]
\label{II.1.5.1-zh}
设 $X$ 是一个相对于 $S$ 仿射的预概形。对基预概形的任意换基
$g:S'\to S$，$X'=X_{(S')}=X\times_S S'$ 相对于 $S'$ 仿射。
\end{proposition}

若 $f'$ 是投影 $X'\to S'$，只需证明：对 $S'$ 的每个仿射开集 $U'$，只要
$g(U')$ 包含于 $S$ 的某个仿射开集 $U$，$f^{\prime-1}(U')$ 就是仿射开集
（\egaref{II.1.2.1-zh}{1.2.1}）。因此可限于 $S$ 和 $S'$ 均仿射的情形，并且只需
证明此时 $X'$ 是仿射概形（\egaref{II.1.3.4-zh}{1.3.4}）。但这时 $X$ 也是仿射概形
（\egaref{II.1.3.4-zh}{1.3.4}）；若 $A$、$A'$、$B$ 分别为 $S$、$S'$、$X$
的环，则已知 $X'$ 是环为 $A'\otimes_A B$ 的仿射概形
（\egaxref{I.3.2.2-fr}{I, 3.2.2}）。

\begin{corollary}[1.5.2]
\label{II.1.5.2-zh}
在 \egaref{II.1.5.1-zh}{1.5.1} 的假设下，设 $f:X\to S$ 是结构态射，
$f':X'\to S'$、$g':X'\to X$ 是投影，从而下图交换：
\[
  \xymatrix{
    X\ar[d] & X'\ar[l]_{g'}\ar[d]^{f'}\\
    S & S'\ar[l]_g
  }
\]
对每个拟凝聚 $\mathcal{O}_X$-模层 $\mathcal{F}$，存在
$\mathcal{O}_{S'}$-模层的典范同构
\[
  u:g^*(f_*(\mathcal{F}))\xrightarrow{\sim}
  f'_*(g^{\prime*}(\mathcal{F})).
  \tag{1.5.2.1}\label{II.1.5.2.1-zh}
\]
特别地，存在 $g^*(\mathcal{A}(X))$ 到 $\mathcal{A}(X')$ 的典范同构。
\end{corollary}

为定义 $u$，只需定义同态
\[
  v:f_*(\mathcal{F})\longrightarrow
  g_*(f'_*(g^{\prime*}(\mathcal{F})))
  =f_*(g'_*(g^{\prime*}(\mathcal{F})))
\]
并取 $u=v^\sharp$（\egaxref{0.4.4.3-fr}{0, 4.4.3}）。取
$v=f_*(\rho)$，其中 $\rho$ 是典范同态
$\mathcal{F}\to g'_*(g^{\prime*}(\mathcal{F}))$
（\egaxref{0.4.4.3-fr}{0, 4.4.3}）。为证明 $u$ 是同构，只需考虑 $S$、$S'$，
从而 $X$、$X'$ 都仿射的情形。沿用 \egaref{II.1.5.1-zh}{1.5.1} 的记号，此时
$\mathcal{F}=\widetilde{M}$，其中 $M$ 是一个 $B$-模。立即可见，
$g^*(f_*(\mathcal{F}))$ 与 $f'_*(g^{\prime*}(\mathcal{F}))$ 都等于
$S'$ 上与 $A'$-模 $A'\otimes_A M$ 相伴的 $\mathcal{O}_{S'}$-模层
（这里把 $M$ 视为 $A$-模），而 $u$ 是与恒等同态相伴的同态
（\egaxref{I.1.6.3-fr}{I, 1.6.3}、\egaxref{I.1.6.5-fr}{1.6.5} 及
\egaxref{I.1.6.7-fr}{1.6.7}）。

\begin{remark}[1.5.3]
\label{II.1.5.3-zh}
不要以为当不再假设 $X$ 相对于 $S$ 仿射时，\egaref{II.1.5.2-zh}{1.5.2}
仍然成立；即使取 $S'=\operatorname{Spec}(k(s))$（$s\in S$），并令
$S'\to S$ 为典范态射（\egaxref{I.2.4.5-fr}{I, 2.4.5}），结论也未必成立。
在这种情形，$X'$ 正是\emph{纤维} $f^{-1}(s)$
（\egaxref{I.3.6.2-fr}{I, 3.6.2}）。换言之，当 $X$ 不相对于 $S$ 仿射时，
\oldpage[II]{13}
“拟凝聚层的直像”运算与“取纤维”运算并不交换。不过，我们将在第三章
（\egaxref{III.4.2.4-fr}{III, 4.2.4}）看到一个具有“渐近”性质的同类结果：
当 $f$ 是紧合态射（\egaref{II.5.4-zh}{5.4}）、$S$ 是诺特概形时，该结果对
$X$ 上的\emph{凝聚}层成立。
\end{remark}

\begin{corollary}[1.5.4]
\label{II.1.5.4-zh}
对每个相对于 $S$ 仿射的预概形 $X$ 和每个 $s\in S$，纤维 $f^{-1}(s)$
（其中 $f$ 表示结构态射 $X\to S$）是仿射概形。
\end{corollary}

只需对 $S'=\operatorname{Spec}(k(s))$ 应用
\egaref{II.1.5.1-zh}{1.5.1}，再使用 \egaref{II.1.3.4-zh}{1.3.4}。

\begin{corollary}[1.5.5]
\label{II.1.5.5-zh}
设 $X$ 是一个 $S$-预概形，$S'$ 是一个相对于 $S$ 仿射的预概形。则
$X'=X_{(S')}$ 是一个相对于 $X$ 仿射的预概形。此外，若 $f:X\to S$ 是结构态射，
则存在 $\mathcal{O}_X$-代数层的典范同构
$\mathcal{A}(X')\xrightarrow{\sim}f^*(\mathcal{A}(S'))$；而对每个拟凝聚
$\mathcal{A}(S')$-模层 $\mathcal{M}$，存在典范双同构
$f^*(\mathcal{M})\xrightarrow{\sim}
\mathcal{A}(f^{\prime*}(\widetilde{\mathcal{M}}))$，其中
$f'=f_{(S')}$ 表示结构态射 $X'\to S'$。
\end{corollary}

只需在 \egaref{II.1.5.1-zh}{1.5.1} 和 \egaref{II.1.5.2-zh}{1.5.2}
中交换 $X$ 与 $S'$ 的角色。

\begin{env}[1.5.6]
\label{II.1.5.6-zh}
现在设 $S$、$S'$ 是两个预概形，$q:S'\to S$ 是一个态射，
$\mathcal{B}$（相应地，$\mathcal{B}'$）是一个拟凝聚
$\mathcal{O}_S$-代数层（相应地，$\mathcal{O}_{S'}$-代数层），
$u:\mathcal{B}\to\mathcal{B}'$ 是一个 $q$-态射（也就是说，是
$\mathcal{O}_S$-代数层的同态
$\mathcal{B}\to q_*(\mathcal{B}')$）。若
$X=\operatorname{Spec}(\mathcal{B})$、$X'=\operatorname{Spec}(\mathcal{B}')$，
则由此典范地导出态射
\[
  v=\operatorname{Spec}(u):X'\to X
\]
，使得下图交换：
\[
  \xymatrix{
    X'\ar[r]^v\ar[d] & X\ar[d]\\
    S'\ar[r]_q & S
  }
  \tag{1.5.6.1}\label{II.1.5.6.1-zh}
\]
其中竖直箭头是结构态射。事实上，给出 $u$ 等价于给出拟凝聚
$\mathcal{O}_{S'}$-代数层的同态
$u^\sharp:q^*(\mathcal{B})\to\mathcal{B}'$
（\egaxref{0.4.4.3-fr}{0, 4.4.3}）；因此它典范地定义一个 $S'$-态射
\[
  w:\operatorname{Spec}(\mathcal{B}')\to
  \operatorname{Spec}(q^*(\mathcal{B}))
\]
，使得 $\mathcal{A}(w)=u^\sharp$
（\egaref{II.1.2.7-zh}{1.2.7}）。另一方面，由
\egaref{II.1.5.2-zh}{1.5.2} 可知
$\operatorname{Spec}(q^*(\mathcal{B}))$ 典范同构于 $X\times_S S'$。
态射 $v$ 是 $w$ 与第一投影的复合
$X'\xrightarrow{w}X\times_S S'\xrightarrow{p_1}X$，而
\egaref{II.1.5.6.1-zh}{1.5.6.1} 的交换性由定义得出。设 $U$
（相应地，$U'$）是 $S$（相应地，$S'$）的仿射开集，且
$q(U')\subset U$；令 $A=\Gamma(U,\mathcal{O}_S)$、
$A'=\Gamma(U',\mathcal{O}_{S'})$ 为相应的环，令
$B=\Gamma(U,\mathcal{B})$、$B'=\Gamma(U',\mathcal{B}')$。
$u$ 限制所得的 $(q|U')$-态射
$\mathcal{B}|U\to\mathcal{B}'|U'$ 对应于代数的一个双同态
$B\to B'$。若 $V$、$V'$ 分别是 $U$、$U'$ 在结构态射下于 $X$、$X'$
中的逆像，则 $v$ 的限制 $V'\to V$ 对应于上述双同态
（\egaxref{I.1.7.3-fr}{I, 1.7.3}）。
\end{env}

\begin{env}[1.5.7]
\label{II.1.5.7-zh}
在 \egaref{II.1.5.6-zh}{1.5.6} 的相同假设下，设 $\mathcal{M}$ 是一个拟凝聚
$\mathcal{B}$-模层。则存在 $\mathcal{O}_{X'}$-模层的典范同构
\[
  v^*(\widetilde{\mathcal{M}})\xrightarrow{\sim}
  \bigl(q^*(\mathcal{M})\otimes_{q^*(\mathcal{B})}\mathcal{B}'\bigr)^{\sim}
  \tag{1.5.7.1}\label{II.1.5.7.1-zh}
\]

\oldpage[II]{14}
事实上，典范同构 \egaref{II.1.5.2.1-zh}{1.5.2.1} 给出了
$p_1^*(\widetilde{\mathcal{M}})$ 到
$\operatorname{Spec}(q^*(\mathcal{B}))$ 上与
$q^*(\mathcal{B})$-模层 $q^*(\mathcal{M})$ 相伴之层的典范同构；随后只需应用
\egaref{II.1.4.7-zh}{1.4.7}。
\end{env}
% Locked French authority: ega2-1-fr.tex, 820505 B,
% SHA-256 84EDBE3E83530AF2959B441796337C9DC21EAFCA6A13114A26778760FBF437AC.
% Sequential source scope: lines 802-858, subsection 1.6 / printed p.14.

\subsection{仿射态射}
\label{II.1.6-zh}

\begin{env}[1.6.1]
\label{II.1.6.1-zh}
若预概形的态射 $f:X\to Y$ 使 $X$ 成为一个相对于 $Y$ 仿射的预概形，就称
$f$ 是\emph{仿射的}。用这种语言，预概形相对于另一个预概形仿射的性质可表述如下。
\end{env}

\begin{proposition}[1.6.2]
\label{II.1.6.2-zh}
\begin{enumerate}
  \item[\rm(i)] 闭浸入是仿射的。
  \item[\rm(ii)] 两个仿射态射的复合是仿射的。
  \item[\rm(iii)] 若 $f:X\to Y$ 是仿射 $S$-态射，则对基的每个换基
    $S'\to S$，$f_{(S')}:X_{(S')}\to Y_{(S')}$ 是仿射的。
  \item[\rm(iv)] 若 $f:X\to Y$ 和 $f':X'\to Y'$ 是两个仿射 $S$-态射，则
    \[
      f\times_S f':X\times_S X'\to Y\times_S Y'
    \]
    是仿射的。
  \item[\rm(v)] 若 $f:X\to Y$ 和 $g:Y\to Z$ 是两个态射，$g\circ f$ 仿射且
    $g$ 可分离，则 $f$ 仿射。
  \item[\rm(vi)] 若 $f$ 仿射，则 $f_{\mathrm{red}}$ 也仿射。
\end{enumerate}
\end{proposition}

根据 \egaxref{I.5.5.12-fr}{I, 5.5.12}，只需证明 (i)、(ii)、(iii)。
其中 (i) 正是例 \egaref{II.1.2.2-zh}{1.2.2}，(ii) 正是
\egaref{II.1.3.5-zh}{1.3.5}；最后，(iii) 由
\egaref{II.1.5.1-zh}{1.5.1} 得出，因为 $X_{(S')}$ 同构于积
$X\times_Y Y_{(S')}$（\egaxref{I.3.3.11-fr}{I, 3.3.11}）。

\begin{corollary}[1.6.3]
\label{II.1.6.3-zh}
若 $X$ 是仿射概形、$Y$ 是概形，则每个态射 $X\to Y$ 都是仿射的。
\end{corollary}

\begin{proposition}[1.6.4]
\label{II.1.6.4-zh}
设 $Y$ 是一个局部诺特预概形，$f:X\to Y$ 是一个有限型态射。则 $f$ 仿射，
当且仅当 $f_{\mathrm{red}}$ 仿射。
\end{proposition}

由 \egaref{II.1.6.2-zh}{1.6.2, (vi)}，只需证明条件的充分性。只需证明：
若 $Y$ 是仿射诺特概形，则 $X$ 是仿射的。此时 $Y_{\mathrm{red}}$ 是仿射的，
由假设 $X_{\mathrm{red}}$ 也仿射。又因为 $X$ 是诺特的，结论遂由
\egaxref{I.6.1.7-fr}{I, 6.1.7} 得出。
% Locked French authority: ega2-1-fr.tex, 820505 B,
% SHA-256 84EDBE3E83530AF2959B441796337C9DC21EAFCA6A13114A26778760FBF437AC.
% Sequential source scope: lines 860-1313, subsection 1.7 / printed pp.14-19.

\subsection{与模层相伴的向量丛}
\label{II.1.7-zh}

\begin{env}[1.7.1]
\label{II.1.7.1-zh}
设 $A$ 是一个环，$E$ 是一个 $A$-模。回顾一下，$E$ 上的\emph{对称代数}
$\mathbf{S}(E)$（或 $\mathbf{S}_A(E)$）是张量代数 $\mathbf{T}(E)$ 对由所有元素
$x\otimes y-y\otimes x$（$x,y$ 遍历 $E$）生成的双边理想的商代数。
代数 $\mathbf{S}(E)$ 由下述泛性质刻画：设 $\sigma:E\to\mathbf{S}(E)$ 是典范映射
（即 $E\to\mathbf{T}(E)$ 与典范映射
$\mathbf{T}(E)\to\mathbf{S}(E)$ 的复合）。若 $B$ 是一个\emph{交换}
$A$-代数，则每个 $A$-线性映射 $E\to B$ 都唯一分解为
$E\xrightarrow{\sigma}\mathbf{S}(E)\xrightarrow{g}B$，其中 $g$ 是一个
$A$-\emph{代数}同态。由这一刻画立即得出：对任意两个 $A$-模 $E$、$F$，在典范同构
意义下有
\[
  \mathbf{S}(E\oplus F)=\mathbf{S}(E)\otimes\mathbf{S}(F)
\]
。

\oldpage[II]{15}
此外，$\mathbf{S}(E)$ 关于 $E$ 是从 $A$-模范畴到交换 $A$-代数范畴的协变函子。
上述刻画还表明，若 $E=\varinjlim E_\lambda$，则在典范同构意义下
$\mathbf{S}(E)=\varinjlim\mathbf{S}(E_\lambda)$。在不致混淆时，通常把乘积
$\sigma(x_1)\sigma(x_2)\cdots\sigma(x_n)$（$x_i\in E$）简写为
$x_1x_2\cdots x_n$。代数 $\mathbf{S}(E)$ 是\emph{分次的}，其中
$\mathbf{S}_n(E)$ 是 $E$ 的 $n$ 个元素之乘积的所有线性组合构成的集合
（$n\geq 0$）。代数 $\mathbf{S}(A)$ 典范同构于单个不定元的多项式代数
$A[T]$，而 $\mathbf{S}(A^n)$ 典范同构于 $n$ 个不定元的多项式代数
$A[T_1,\ldots,T_n]$。
\end{env}

\begin{env}[1.7.2]
\label{II.1.7.2-zh}
设 $\varphi:A\to B$ 是一个环同态。若 $F$ 是一个 $B$-模，则典范映射
$F\to\mathbf{S}(F)$ 给出典范映射
$F_{[\varphi]}\to\mathbf{S}(F)_{[\varphi]}$，因而它分解为
$F_{[\varphi]}\to\mathbf{S}(F_{[\varphi]})\to
\mathbf{S}(F)_{[\varphi]}$。典范同态
$\mathbf{S}(F_{[\varphi]})\to\mathbf{S}(F)_{[\varphi]}$ 是满的，但未必是双射。
若 $E$ 是一个 $A$-模，则每个双同态 $E\to F$（即每个 $A$-同态
$E\to F_{[\varphi]}$）都典范地给出 $A$-代数同态
$\mathbf{S}(E)\to\mathbf{S}(F_{[\varphi]})\to
\mathbf{S}(F)_{[\varphi]}$，也就是代数的双同态
$\mathbf{S}(E)\to\mathbf{S}(F)$。

仍沿用这些记号。对每个 $A$-模 $E$，$\mathbf{S}(E\otimes_A B)$ 典范同构于代数
$\mathbf{S}(E)\otimes_A B$；这由 $\mathbf{S}(E)$ 的泛性质
（\egaref{II.1.7.1-zh}{1.7.1}）立即得出。
\end{env}

\begin{env}[1.7.3]
\label{II.1.7.3-zh}
设 $R$ 是环 $A$ 的一个乘法子集。把 \egaref{II.1.7.2-zh}{1.7.2} 应用于
$B=R^{-1}A$，并注意 $R^{-1}E=E\otimes_A R^{-1}A$，可见在典范同构意义下
$\mathbf{S}(R^{-1}E)=R^{-1}\mathbf{S}(E)$。此外，若 $R'\supset R$ 是 $A$
的另一个乘法子集，则下图交换：
\[
  \xymatrix{
    R^{-1}E\ar[r]\ar[d] & {R'}^{-1}E\ar[d]\\
    \mathbf{S}(R^{-1}E)\ar[r] & \mathbf{S}({R'}^{-1}E)
  }
\]
\end{env}

\begin{env}[1.7.4]
\label{II.1.7.4-zh}
现在设 $(S,\mathcal{A})$ 是一个赋环空间，$\mathcal{E}$ 是 $S$ 上的
$\mathcal{A}$-模层。对每个开集 $U\subset S$，令它对应于
$\Gamma(U,\mathcal{A})$-模 $\mathbf{S}(\Gamma(U,\mathcal{E}))$。
由 $\mathbf{S}(E)$ 的函子性（\egaref{II.1.7.2-zh}{1.7.2}），这定义了一个
代数预层。称其相伴层 $\mathbf{S}(\mathcal{E})$（或
$\mathbf{S}_{\mathcal{A}}(\mathcal{E})$）为 $\mathcal{A}$-模层
$\mathcal{E}$ 的\emph{对称 $\mathcal{A}$-代数层}。由
\egaref{II.1.7.1-zh}{1.7.1} 立即可知，$\mathbf{S}(\mathcal{E})$ 是一个泛问题的解：
若 $\mathcal{B}$ 是一个 $\mathcal{A}$-代数层，则每个 $\mathcal{A}$-模层同态
$\mathcal{E}\to\mathcal{B}$ 都唯一分解为
$\mathcal{E}\to\mathbf{S}(\mathcal{E})\to\mathcal{B}$，其中第二个箭头是
$\mathcal{A}$-代数层同态。因此，$\mathcal{A}$-模层同态
$\mathcal{E}\to\mathcal{B}$ 与 $\mathcal{A}$-代数层同态
$\mathbf{S}(\mathcal{E})\to\mathcal{B}$ 一一对应。特别地，每个
$\mathcal{A}$-模层同态 $u:\mathcal{E}\to\mathcal{F}$ 都定义
$\mathcal{A}$-代数层同态
$\mathbf{S}(u):\mathbf{S}(\mathcal{E})\to\mathbf{S}(\mathcal{F})$，
所以 $\mathbf{S}(\mathcal{E})$ 关于 $\mathcal{E}$ 是协变函子。

\oldpage[II]{16}
由 \egaref{II.1.7.2-zh}{1.7.2} 以及 $\mathbf{S}$ 与归纳极限可交换，
对每个点 $s\in S$ 有
$(\mathbf{S}(\mathcal{E}))_s=\mathbf{S}(\mathcal{E}_s)$。
若 $\mathcal{E}$、$\mathcal{F}$ 是两个 $\mathcal{A}$-模层，则
$\mathbf{S}(\mathcal{E}\oplus\mathcal{F})$ 典范同构于
$\mathbf{S}(\mathcal{E})\otimes_{\mathcal{A}}\mathbf{S}(\mathcal{F})$；
考察相应的预层即可看出这一点。

还应注意，$\mathbf{S}(\mathcal{E})$ 是一个分次 $\mathcal{A}$-代数层，即各
$\mathbf{S}_n(\mathcal{E})$ 的无限直和；其中 $\mathcal{A}$-模层
$\mathbf{S}_n(\mathcal{E})$ 是预层
$U\mapsto\mathbf{S}_n(\Gamma(U,\mathcal{E}))$ 的相伴层。特别取
$\mathcal{E}=\mathcal{A}$，可见
$\mathbf{S}_{\mathcal{A}}(\mathcal{A})$ 同构于
$\mathcal{A}[T]=\mathcal{A}\otimes_{\mathbf{Z}}\mathbf{Z}[T]$；
这里 $T$ 是不定元，$\mathbf{Z}$ 视为常值层。
\end{env}

\begin{env}[1.7.5]
\label{II.1.7.5-zh}
设 $(T,\mathcal{B})$ 是第二个赋环空间，$f:(S,\mathcal{A})\to
(T,\mathcal{B})$ 是一个态射。若 $\mathcal{F}$ 是一个 $\mathcal{B}$-模层，
则 $\mathbf{S}(f^*(\mathcal{F}))$ 典范同构于
$f^*(\mathbf{S}(\mathcal{F}))$。事实上，若 $f=(\psi,\theta)$，则根据定义
（\egaxref{0.4.3.1-fr}{0, 4.3.1}）
\[
  \mathbf{S}(f^*(\mathcal{F}))
  =\mathbf{S}(\psi^*(\mathcal{F})
    \otimes_{\psi^*(\mathcal{B})}\mathcal{A})
  =\mathbf{S}(\psi^*(\mathcal{F}))
    \otimes_{\psi^*(\mathcal{B})}\mathcal{A}
\]
。
这里使用了 \egaref{II.1.7.2-zh}{1.7.2}。对 $S$ 的每个开集 $U$ 和
$\mathbf{S}(\psi^*(\mathcal{F}))$ 在 $U$ 上的每个截面 $h$，在每个点
$s\in U$ 的某个邻域 $V$ 中，$h$ 都与
$\mathbf{S}(\Gamma(V,\psi^*(\mathcal{F})))$ 的一个元素相同。回顾
$\psi^*(\mathcal{F})$ 的定义（\egaxref{0.3.7.1-fr}{0, 3.7.1}），并注意对任意模
$E$，$\mathbf{S}(E)$ 的每个元素都是 $E$ 的元素之有限多个乘积的线性组合，可知在
$T$ 中存在 $\psi(s)$ 的一个邻域 $W$、$\mathbf{S}(\mathcal{F})$ 在 $W$ 上的一个
截面 $h'$，并且存在 $s$ 的一个邻域
$V'\subset V\cap\psi^{-1}(W)$，使 $h$ 在 $V'$ 上等于
$t\mapsto h'(\psi(t))$。断言由此得证。
\end{env}

\begin{proposition}[1.7.6]
\label{II.1.7.6-zh}
设 $A$ 是一个环，$S=\operatorname{Spec}(A)$ 是它的素谱，
$\mathcal{E}=\widetilde{M}$ 是与 $A$-模 $M$ 相伴的 $\mathcal{O}_S$-模层。
则 $\mathcal{O}_S$-代数层 $\mathbf{S}(\mathcal{E})$ 与 $A$-代数
$\mathbf{S}(M)$ 相伴。
\end{proposition}

\begin{proof}
事实上，对每个 $f\in A$，有
$\mathbf{S}(M_f)=(\mathbf{S}(M))_f$
（\egaref{II.1.7.3-zh}{1.7.3}），故本命题由定义
（\egaxref{I.1.3.4-fr}{I, 1.3.4}）得出。
\end{proof}

\begin{corollary}[1.7.7]
\label{II.1.7.7-zh}
设 $S$ 是一个预概形，$\mathcal{E}$ 是一个拟凝聚 $\mathcal{O}_S$-模层。
则 $\mathcal{O}_S$-代数层 $\mathbf{S}(\mathcal{E})$ 是拟凝聚的。若
$\mathcal{E}$ 还是有限型的，则每个 $\mathcal{O}_S$-模层
$\mathbf{S}_n(\mathcal{E})$ 都是有限型的。
\end{corollary}

\begin{proof}
第一个断言是 \egaref{II.1.7.6-zh}{1.7.6} 与
\egaxref{I.1.4.1-fr}{I, 1.4.1} 的直接推论。第二个断言来自如下事实：若 $E$
是有限型 $A$-模，则 $\mathbf{S}_n(E)$ 是有限型 $A$-模；再应用
\egaxref{I.1.3.13-fr}{I, 1.3.13} 即可。
\end{proof}

\begin{definition}[1.7.8]
\label{II.1.7.8-zh}
设 $\mathcal{E}$ 是一个拟凝聚 $\mathcal{O}_S$-模层。称拟凝聚
$\mathcal{O}_S$-代数层 $\mathbf{S}(\mathcal{E})$ 的谱
（\egaref{II.1.3.1-zh}{1.3.1}）为\emph{由 $\mathcal{E}$ 定义的 $S$ 上的向量丛}，
并记作 $\mathbf{V}(\mathcal{E})$。
\end{definition}

根据 \egaref{II.1.2.7-zh}{1.2.7}，对每个 $S$-预概形 $X$，$S$-态射
$X\to\mathbf{V}(\mathcal{E})$ 与 $\mathcal{O}_S$-代数层同态
$\mathbf{S}(\mathcal{E})\to\mathcal{A}(X)$ 典范地一一对应，因而也与
$\mathcal{O}_S$-模层同态
$\mathcal{E}\to\mathcal{A}(X)=f_*(\mathcal{O}_X)$ 典范地一一对应；这里 $f$
是结构态射 $X\to S$。特别地，有下述结论。

\begin{env}[1.7.9]
\label{II.1.7.9-zh}
取 $X$ 为 $S$ 在开集 $U\subset S$ 上诱导的子预概形。则 $S$-态射
$U\to\mathbf{V}(\mathcal{E})$ 正是 $\mathbf{V}(\mathcal{E})$ 在开集
$p^{-1}(U)$ 上诱导的 $U$-预概形的 $U$-截面
（\egaxref{I.2.5.5-fr}{I, 2.5.5}）；这里 $p$ 是结构态射
$\mathbf{V}(\mathcal{E})\to S$。由刚才所见，这些 $U$-截面与
$\mathcal{O}_S$-模层同态
$\mathcal{E}\to j_*(\mathcal{O}_S|U)$ 一一对应，其中 $j:U\to S$ 是典范注入；
等价地（\egaxref{0.4.4.3-fr}{0, 4.4.3}），它们与

\oldpage[II]{17}
$\mathcal{O}_S|U$-模层同态
$j^*(\mathcal{E})=\mathcal{E}|U\to\mathcal{O}_S|U$ 一一对应。此外，显然，
一个 $S$-态射 $U\to\mathbf{V}(\mathcal{E})$ 限制到开集
$U'\subset U$，对应于相应同态
$\mathcal{E}|U\to\mathcal{O}_S|U$ 限制到 $U'$。由此得出，
$\mathbf{V}(\mathcal{E})$ 的\emph{$S$-截面芽层}典范同构于
$\mathcal{E}$ 的\emph{对偶 $\check{\mathcal{E}}$}。

特别取 $X=U=S$，则零同态 $\mathcal{E}\to\mathcal{O}_S$ 对应于
$\mathbf{V}(\mathcal{E})$ 的一个典范 $S$-截面，称为\emph{零 $S$-截面}
（参见 \egaref{II.8.3.3-zh}{8.3.3}）。
\end{env}

\begin{env}[1.7.10]
\label{II.1.7.10-zh}
现在取 $X$ 为域 $K$ 的谱 $\{\xi\}$。结构态射 $f:X\to S$ 对应于单同态
$k(s)\to K$，其中 $s=f(\xi)$
（\egaxref{I.2.4.6-fr}{I, 2.4.6}）。因此，$S$-态射
$\{\xi\}\to\mathbf{V}(\mathcal{E})$ 正是
\emph{$\mathbf{V}(\mathcal{E})$ 取值于 $k(s)$ 的扩域 $K$ 的几何点}
（\egaxref{I.3.4.5-fr}{I, 3.4.5}），这些点都位于 $p^{-1}(s)$ 的点上。
这一点集也可称为\emph{$\mathbf{V}(\mathcal{E})$ 在点 $s$ 上的 $K$-有理几何纤维}。
根据 \egaref{II.1.7.8-zh}{1.7.8}，它同构于 $\mathcal{O}_S$-模层同态
$\mathcal{E}\to f_*(\mathcal{O}_X)$ 的集合；等价地
（\egaxref{0.4.4.3-fr}{0, 4.4.3}），它同构于 $\mathcal{O}_X$-模层同态
$f^*(\mathcal{E})\to\mathcal{O}_X=K$ 的集合。但按定义
（\egaxref{0.4.3.1-fr}{0, 4.3.1}），令
$\mathcal{E}^s=\mathcal{E}_s/\mathfrak{m}_s\mathcal{E}_s$，就有
$f^*(\mathcal{E})=\mathcal{E}_s\otimes_{\mathcal{O}_s}K
=\mathcal{E}^s\otimes_{k(s)}K$。因此，
$\mathbf{V}(\mathcal{E})$ 在 $s$ 上的 $K$-有理几何纤维同构于
\emph{$K$-向量空间 $\mathcal{E}^s\otimes_{k(s)}K$ 的对偶}。
若 $\mathcal{E}^s$ 或 $K$ 在 $k(s)$ 上是有限维的，该对偶还同构于
$(\mathcal{E}^s)^\vee\otimes_{k(s)}K$，其中
$(\mathcal{E}^s)^\vee$ 表示 $k(s)$-向量空间 $\mathcal{E}^s$ 的对偶。
\end{env}

\begin{proposition}[1.7.11]
\label{II.1.7.11-zh}
\begin{enumerate}
  \item[\rm(i)] $\mathbf{V}(\mathcal{E})$ 关于 $\mathcal{E}$ 是从拟凝聚
    $\mathcal{O}_S$-模层范畴到仿射 $S$-概形范畴的反变函子。
  \item[\rm(ii)] 若 $\mathcal{E}$ 是有限型 $\mathcal{O}_S$-模层，则
    $\mathbf{V}(\mathcal{E})$ 在 $S$ 上是有限型的。
  \item[\rm(iii)] 若 $\mathcal{E}$、$\mathcal{F}$ 是两个拟凝聚
    $\mathcal{O}_S$-模层，则 $\mathbf{V}(\mathcal{E}\oplus\mathcal{F})$
    典范同构于
    $\mathbf{V}(\mathcal{E})\times_S\mathbf{V}(\mathcal{F})$。
  \item[\rm(iv)] 设 $g:S'\to S$ 是一个态射。对每个拟凝聚
    $\mathcal{O}_S$-模层 $\mathcal{E}$，$\mathbf{V}(g^*(\mathcal{E}))$
    典范同构于
    $\mathbf{V}(\mathcal{E})_{(S')}
    =\mathbf{V}(\mathcal{E})\times_S S'$。
  \item[\rm(v)] 拟凝聚 $\mathcal{O}_S$-模层的满同态
    $\mathcal{E}\to\mathcal{F}$ 对应于闭浸入
    $\mathbf{V}(\mathcal{F})\to\mathbf{V}(\mathcal{E})$。
\end{enumerate}
\end{proposition}

\begin{proof}
考虑到每个 $\mathcal{O}_S$-模层同态
$\mathcal{E}\to\mathcal{F}$ 都典范地定义 $\mathcal{O}_S$-代数层同态
$\mathbf{S}(\mathcal{E})\to\mathbf{S}(\mathcal{F})$，(i) 是
\egaref{II.1.2.7-zh}{1.2.7} 的直接推论。(ii) 立即由定义
（\egaxref{I.6.3.1-fr}{I, 6.3.1}）以及如下事实得出：若 $E$ 是有限型 $A$-模，
则 $\mathbf{S}(E)$ 是有限型 $A$-代数。为证明 (iii)，只需从典范同构
$\mathbf{S}(\mathcal{E}\oplus\mathcal{F})
\simeq\mathbf{S}(\mathcal{E})
\otimes_{\mathcal{O}_S}\mathbf{S}(\mathcal{F})$
（\egaref{II.1.7.4-zh}{1.7.4}）出发，再应用
\egaref{II.1.4.6-zh}{1.4.6}。同样，为证明 (iv)，只需从典范同构
$\mathbf{S}(g^*(\mathcal{E}))
\simeq g^*(\mathbf{S}(\mathcal{E}))$
（\egaref{II.1.7.5-zh}{1.7.5}）出发，再应用
\egaref{II.1.5.2-zh}{1.5.2}。最后，为证明 (v)，只需注意：若同态
$\mathcal{E}\to\mathcal{F}$ 是满的，则相应的 $\mathcal{O}_S$-代数层同态
$\mathbf{S}(\mathcal{E})\to\mathbf{S}(\mathcal{F})$ 也是满的；结论遂由
\egaref{II.1.4.10-zh}{1.4.10} 得出。
\end{proof}

\begin{env}[1.7.12]
\label{II.1.7.12-zh}
特别取 $\mathcal{E}=\mathcal{O}_S$。预概形 $\mathbf{V}(\mathcal{O}_S)$ 是
$\mathcal{O}_S$-代数层 $\mathbf{S}(\mathcal{O}_S)$ 的谱，因而是仿射
$S$-概形；该代数层同构于
$\mathcal{O}_S[T]=\mathcal{O}_S\otimes_{\mathbf{Z}}\mathbf{Z}[T]$。

\oldpage[II]{18}
这里 $T$ 是不定元。当 $S=\operatorname{Spec}(\mathbf{Z})$ 时，由
\egaref{II.1.7.6-zh}{1.7.6} 立即可知这一点；一般情形则通过考虑结构态射
$S\to\operatorname{Spec}(\mathbf{Z})$ 并使用
\egaref{II.1.7.11-zh}{1.7.11, (iv)} 得出。鉴于这一结果，仍记
$\mathbf{V}(\mathcal{O}_S)=S[T]$，于是有公式
\[
  S[T]=S\otimes_{\mathbf{Z}}\mathbf{Z}[T].
  \tag{1.7.12.1}\label{II.1.7.12.1-zh}
\]

把 $S[T]$ 的 $S$-截面芽层同 $\mathcal{O}_S$ 等同起来这一事实，已在
\egaxref{I.3.3.15-fr}{I, 3.3.15} 中出现；在这里，它又在更一般的语境中作为
\egaref{II.1.7.9-zh}{1.7.9} 的特例得到。
\end{env}

\begin{env}[1.7.13]
\label{II.1.7.13-zh}
对每个 $S$-预概形 $X$，由 \egaref{II.1.7.8-zh}{1.7.8} 已知
$\operatorname{Hom}_S(X,S[T])$ 典范同构于
$\operatorname{Hom}_{\mathcal{O}_S}(\mathcal{O}_S,\mathcal{A}(X))$，
后者又典范同构于 $\Gamma(S,\mathcal{A}(X))$，因而带有环结构。此外，对每个
$S$-态射 $h:X\to Y$，都有保持环结构的态射
$\Gamma(\mathcal{A}(h)):\Gamma(S,\mathcal{A}(Y))
\to\Gamma(S,\mathcal{A}(X))$
（\egaref{II.1.1.2-zh}{1.1.2}）。赋予
$\operatorname{Hom}_S(X,S[T])$ 这样定义的环结构后，可见
$\operatorname{Hom}_S(X,S[T])$ 是关于 $X$ 的\emph{反变函子}，从
$S$-预概形范畴取值到环范畴。

另一方面，$\operatorname{Hom}_S(X,\mathbf{V}(\mathcal{E}))$ 同样同构于
$\operatorname{Hom}_{\mathcal{O}_S}(\mathcal{E},\mathcal{A}(X))$；
这里把 $\mathcal{A}(X)$ 视为一个\emph{$\mathcal{O}_S$-模层}。因此，可以典范地
赋予它在环 $\operatorname{Hom}_S(X,S[T])$ 上的\emph{模}结构。与上面相同可见，偶对
\[
  \bigl(\operatorname{Hom}_S(X,S[T]),
  \operatorname{Hom}_S(X,\mathbf{V}(\mathcal{E}))\bigr)
\]
是关于 $X$ 的反变函子，取值于如下范畴：对象是由环 $A$ 与 $A$-模 $M$ 组成的偶对
$(A,M)$，态射是双同态。

我们把这些事实解释为：$S[T]$ 是一个\emph{$S$-环概形}，而
$\mathbf{V}(\mathcal{E})$ 是环概形 $S[T]$ 上的一个\emph{$S$-模概形}
（参见第 0 章第 8 节）。
\end{env}

\begin{env}[1.7.14]
\label{II.1.7.14-zh}
下面将看到，$S$-概形 $\mathbf{V}(\mathcal{E})$ 上定义的 $S$-模概形结构，能够在
唯一同构意义下重构 $\mathcal{O}_S$-模层 $\mathcal{E}$。为此，我们将证明
$\mathcal{E}$ 典范同构于
$\mathbf{S}(\mathcal{E})=\mathcal{A}(\mathbf{V}(\mathcal{E}))$ 的一个
$\mathcal{O}_S$-子模层，并且该子模层可由这一结构定义。事实上，由
\egaref{II.1.7.4-zh}{1.7.4}，$\mathcal{O}_S$-\emph{代数层}同态的集合
$\operatorname{Hom}_{\mathcal{O}_S}(\mathbf{S}(\mathcal{E}),
\mathcal{A}(X))$ 典范同构于 $\mathcal{O}_S$-\emph{模层}同态的集合
$\operatorname{Hom}_{\mathcal{O}_S}(\mathcal{E},\mathcal{A}(X))$。
若 $h$、$h'$ 是后一集合中的两个元素，$s_i$（$1\leq i\leq n$）是
$\mathcal{E}$ 在开集 $U\subset S$ 上的截面，$t$ 是 $\mathcal{A}(X)$ 在 $U$
上的截面，则按定义
\[
  (h+h')(s_1s_2\cdots s_n)
  =\prod_{i=1}^n\bigl(h(s_i)+h'(s_i)\bigr)
\]
并且
\[
  (t.h)(s_1s_2\cdots s_n)=t^n\prod_{i=1}^n h(s_i).
\]

现在，若 $z$ 是 $\mathbf{S}(\mathcal{E})$ 在 $U$ 上的一个截面，则
$h\mapsto h(z)$ 是一个以
$\operatorname{Hom}_S(X,\mathbf{V}(\mathcal{E}))$ 为定义域、以
$\Gamma(U,\mathcal{A}(X))$ 为值域的映射；这里定义域典范同构于
$\operatorname{Hom}_{\mathcal{O}_S}
(\mathbf{S}(\mathcal{E}),\mathcal{A}(X))$。我们将

\oldpage[II]{19}
证明：$\mathcal{E}$ 等同于 $\mathbf{S}(\mathcal{E})$ 中具有下述性质的子模层：
\emph{对每个开集 $U\subset S$、该 $\mathcal{O}_S$-子模层在 $U$ 上的每个截面
$z$，以及每个 $S$-预概形 $X$，映射 $h\mapsto h(z)$ 从
$\operatorname{Hom}_{\mathcal{O}_S|U}(\mathbf{S}(\mathcal{E})|U,
\mathcal{A}(X)|U)$ 到 $\Gamma(U,\mathcal{A}(X))$ 是
$\Gamma(U,\mathcal{A}(X))$-模同态}。

$\mathcal{E}$ 具有这一性质是显然的。为证明逆命题，只需证明：当
$S=\operatorname{Spec}(A)$、$\mathcal{E}=\widetilde{M}$ 时，若
$\mathbf{S}(\mathcal{E})$ 在 $S$ 上的一个截面 $z$ 对 $U=S$ 具有上述性质，
则 $z$ 必为 $\mathcal{E}$ 的截面。此时
$z=\sum_{n=0}^{\infty}z_n$，其中 $z_n\in\mathbf{S}_n(M)$；需要证明
$n\ne1$ 时 $z_n=0$。令 $B=\mathbf{S}(M)$，并取
$X=\operatorname{Spec}(B[T])$，其中 $T$ 是不定元。集合
$\operatorname{Hom}_{\mathcal{O}_S}(\mathbf{S}(\mathcal{E}),
\mathcal{A}(X))$ 同构于环同态 $h:B\to B[T]$ 的集合
（\egaxref{I.1.3.13-fr}{I, 1.3.13}）。由上面的讨论，
$(T.h)(z)=\sum_{n=0}^{\infty}T^nh(z_n)$。关于 $z$ 的假设意味着：对每个同态
$h$，都有
$\sum_{n=0}^{\infty}T^nh(z_n)
=T.\sum_{n=0}^{\infty}h(z_n)$。特别取 $h$ 为典范注入，得到
$\sum_{n=0}^{\infty}T^nz_n
=T.\sum_{n=0}^{\infty}z_n$，这正蕴含 $n\ne1$ 时 $z_n=0$。
\end{env}

\begin{proposition}[1.7.15]
\label{II.1.7.15-zh}
设 $Y$ 是一个底拓扑空间为诺特空间的预概形，或者是一个拟紧概形。每个在 $Y$ 上有限型的
仿射 $Y$-概形 $X$，都 $Y$-同构于某个形如 $\mathbf{V}(\mathcal{E})$ 的
$Y$-概形的闭 $Y$-子概形，其中 $\mathcal{E}$ 是一个有限型拟凝聚
$\mathcal{O}_Y$-模层。
\end{proposition}

\begin{proof}
事实上，拟凝聚 $\mathcal{O}_Y$-代数层 $\mathcal{A}(X)$ 是有限型的
（\egaref{II.1.3.7-zh}{1.3.7}）。所作假设蕴含：$\mathcal{A}(X)$ 由一个有限型
拟凝聚 $\mathcal{O}_Y$-子模层 $\mathcal{E}$ 生成
（\egaxref{I.9.6.5-fr}{I, 9.6.5}）。按定义，这意味着典范延拓注入
$\mathcal{E}\to\mathcal{A}(X)$ 所得的典范同态
$\mathbf{S}(\mathcal{E})\to\mathcal{A}(X)$ 是\emph{满的}。结论遂由
\egaref{II.1.4.10-zh}{1.4.10} 得出。
\end{proof}
% Locked French authority: ega2-1-fr.tex, 820505 B,
% SHA-256 84EDBE3E83530AF2959B441796337C9DC21EAFCA6A13114A26778760FBF437AC.
% Sequential source scope: lines 1315-1649, section II.2 and subsection 2.1 /
% printed pp.19-23.

\section{齐次素谱}
\label{section:II.2-zh}

\subsection{分次环与分次模概述}
\label{subsection:II.2.1-zh}

\begin{notation}[2.1.1]
\label{II.2.1.1-zh}
给定一个按非负次数\emph{分次}的环 $S$，用 $S_n$ 表示 $S$ 中所有 $n$ 次齐次元素
组成的部分（$n\geq0$），用 $S_+$ 表示所有 $n>0$ 的 $S_n$ 的直和。于是
$1\in S_0$，$S_0$ 是 $S$ 的子环，$S_+$ 是 $S$ 的分次理想，并且 $S$ 是
$S_0$ 与 $S_+$ 的直和。若 $M$ 是 $S$ 上的一个\emph{分次模}
（次数可以为正或负），同样用 $M_n$ 表示 $M$ 中所有 $n$ 次齐次元素组成的
$S_0$-模（$n\in\mathbf{Z}$）。

对每个整数 $d>0$，用 $S^{(d)}$ 表示各 $S_{nd}$ 的直和。把 $S_{nd}$ 的元素看作
$n$ 次齐次元素，这些 $S_{nd}$ 就在 $S^{(d)}$ 上定义了分次环结构。

对每个满足 $0\leq k\leq d-1$ 的整数 $k$，用 $M^{(d,k)}$ 表示

\oldpage[II]{20}
各 $M_{nd+k}$（$n\in\mathbf{Z}$）的直和。把 $M_{nd+k}$ 的元素看作 $n$ 次
齐次元素，$M^{(d,k)}$ 就成为分次 $S^{(d)}$-模。把 $M^{(d,0)}$ 简记为
$M^{(d)}$。

沿用上述记号。对每个整数 $n$（正或负），用 $M(n)$ 表示满足
$(M(n))_k=M_{n+k}$（$k\in\mathbf{Z}$）的分次 $S$-模。特别地，
$S(n)$ 是满足 $(S(n))_k=S_{n+k}$ 的分次 $S$-模；约定 $n<0$ 时 $S_n=0$。
若一个分次 $S$-模 $M$ 作为\emph{分次模}同构于若干个形如 $S(n)$ 的模的直和，
就称 $M$ 是\emph{自由的}。因为 $S(n)$ 是由 $S$ 的元素 $1$（视为 $-n$ 次元素）
生成的单生成 $S$-模，这也等价于说 $M$ 有一个由\emph{齐次}元素组成的 $S$-基。

若存在正合序列 $P\to Q\to M\to0$，其中 $P$、$Q$ 是有限多个形如 $S(n)$ 的模
的直和，且各同态均为 $0$ 次同态（参见 \egaref{II.2.1.2-zh}{2.1.2}），就称
分次 $S$-模 $M$ \emph{有有限表示}。
\end{notation}

\begin{env}[2.1.2]
\label{II.2.1.2-zh}
设 $M$、$N$ 是两个分次 $S$-模。如下在 $M\otimes_S N$ 上定义分次 $S$-模结构。
先在张量积 $M\otimes_{\mathbf{Z}}N$ 上定义分次 $\mathbf{Z}$-模结构
（这里 $\mathbf{Z}$ 的分次为
$\mathbf{Z}_0=\mathbf{Z}$，而 $n\ne0$ 时 $\mathbf{Z}_n=0$），令
\[
  (M\otimes_{\mathbf{Z}}N)_q
  =\bigoplus_{m+n=q}M_m\otimes_{\mathbf{Z}}N_n.
\]
因为 $M$、$N$ 分别是各 $M_m$、各 $N_n$ 的直和，已知
$M\otimes_{\mathbf{Z}}N$ 可典范地等同于所有
$M_m\otimes_{\mathbf{Z}}N_n$ 的直和。现在，
$M\otimes_S N=(M\otimes_{\mathbf{Z}}N)/P$，其中 $P$ 是
$M\otimes_{\mathbf{Z}}N$ 中由所有元素
$(xs)\otimes y-x\otimes(sy)$（$x\in M$、$y\in N$、$s\in S$）生成的
$\mathbf{Z}$-子模。显然，$P$ 是 $M\otimes_{\mathbf{Z}}N$ 的分次
$\mathbf{Z}$-子模；取商后立即得到 $M\otimes_S N$ 上的分次 $S$-模结构。

对两个分次 $S$-模 $M$、$N$，回顾：若 $S$-模同态 $u:M\to N$ 对每个
$j\in\mathbf{Z}$ 都满足 $u(M_j)\subset N_{j+k}$，就称它是\emph{$k$ 次的}。
若 $H_n$ 表示从 $M$ 到 $N$ 的所有 $n$ 次同态组成的集合，则用
$\operatorname{Hom}_S(M,N)$ 表示所有 $H_n$（$n\in\mathbf{Z}$）在
从 $M$ 到 $N$ 的全部 $S$-模同态所成的 $S$-模 $H$ 中的直和。一般而言，
$\operatorname{Hom}_S(M,N)$ 不等于 $H$。不过，当 $M$ 是\emph{有限型}时，
$H=\operatorname{Hom}_S(M,N)$。事实上，此时可设 $M$ 由有限多个齐次元素
$x_i$（$1\leq i\leq n$）生成。每个同态 $u\in H$ 唯一写成
$\sum_{k\in\mathbf{Z}}u_k$，其中对每个 $k$，$u_k(x_i)$ 等于 $u(x_i)$ 的
$k+\deg(x_i)$ 次齐次分量（$1\leq i\leq n$）。这蕴含除有限多个指标外
$u_k=0$；而按定义 $u_k\in H_k$，断言得证。

称 $\operatorname{Hom}_S(M,N)$ 中的 $0$ 次元素为\emph{分次 $S$-模同态}。
显然 $S_mH_n\subset H_{m+n}$，所以各 $H_n$ 在
$\operatorname{Hom}_S(M,N)$ 上定义了分次 $S$-模结构。

由这些定义立即得到：对两个分次 $S$-模 $M$、$N$，
\[
  \label{II.2.1.2.1-zh}
  M(m)\otimes_S N(n)=(M\otimes_S N)(m+n),
  \tag{2.1.2.1}
\]
\[
  \label{II.2.1.2.2-zh}
  \operatorname{Hom}_S(M(m),N(n))
  =(\operatorname{Hom}_S(M,N))(n-m).
  \tag{2.1.2.2}
\]

\oldpage[II]{21}
设 $S$、$S'$ 是两个分次环。所谓\emph{分次环同态}
$\varphi:S\to S'$，是指对每个 $n\in\mathbf{Z}$ 都满足
$\varphi(S_n)\subset S'_n$ 的环同态；换言之，它必须是分次
$\mathbf{Z}$-模的\emph{$0$ 次}同态。给出这样的同态就在 $S'$ 上定义了分次
$S$-模结构；连同原有的分次环结构，此时称 $S'$ 是一个\emph{分次 $S$-代数}。

若 $M$ 是一个分次 $S$-模，则分次 $S$-模的张量积 $M\otimes_S S'$ 自然带有
分次 $S'$-模结构，其分次即上面定义的分次。
\end{env}

\begin{lemma}[2.1.3]
\label{II.2.1.3-zh}
设 $S$ 是一个按非负次数分次的环。设 $E$ 是由 $S_+$ 的齐次元素组成的子集。
则 $E$ 作为 $S$-模生成 $S_+$，当且仅当 $E$ 作为 $S_0$-代数生成 $S$。
\end{lemma}

\begin{proof}
条件的充分性显然。下面证明必要性。用 $E_n$（相应地，$E^n$）表示 $E$ 中次数等于
$n$（相应地，不超过 $n$）的元素组成的集合。只需对 $n>0$ 作归纳，证明 $S_n$
是由所有次数为 $n$ 且为 $E^n$ 中元素之乘积的元素生成的 $S_0$-模。
由假设，$n=1$ 时这是显然的；假设还给出
$S_n=\sum_{p=0}^{n-1}S_pE_{n-p}$，归纳论证由此立即完成。
\end{proof}

\begin{corollary}[2.1.4]
\label{II.2.1.4-zh}
$S_+$ 是有限型理想，当且仅当 $S$ 是有限型 $S_0$-代数。
\end{corollary}

\begin{proof}
事实上，总可以假设 $S_0$-代数 $S$（相应地，$S$-理想 $S_+$）的一组有限生成元
由齐次元素组成：只需把所取的每个生成元换成它的各个齐次分量。
\end{proof}

\begin{corollary}[2.1.5]
\label{II.2.1.5-zh}
$S$ 是诺特环，当且仅当 $S_0$ 是诺特环且 $S$ 是有限型 $S_0$-代数。
\end{corollary}

\begin{proof}
条件的充分性显然。必要性来自 $S_0\simeq S/S_+$，并且 $S_+$ 必须是有限型理想
（\egaref{II.2.1.4-zh}{2.1.4}）。
\end{proof}

\begin{lemma}[2.1.6]
\label{II.2.1.6-zh}
设 $S$ 是一个按非负次数分次的环，并且是有限型 $S_0$-代数。设 $M$ 是一个有限型
分次 $S$-模。则：
\begin{enumerate}
  \item[\rm(i)] 各 $M_n$ 都是有限型 $S_0$-模，并且存在整数 $n_0$，使
    $n\leq n_0$ 时 $M_n=0$。
  \item[\rm(ii)] 存在整数 $n_1$ 和整数 $h>0$，使得对每个整数 $n\geq n_1$，
    $M_{n+h}=S_hM_n$。
  \item[\rm(iii)] 对每对满足 $d>0$、$0\leq k\leq d-1$ 的整数 $(d,k)$，
    $M^{(d,k)}$ 都是有限型 $S^{(d)}$-模。
  \item[\rm(iv)] 对每个整数 $d>0$，$S^{(d)}$ 是有限型 $S_0$-代数。
  \item[\rm(v)] 存在整数 $h>0$，使得对每个 $m>0$ 都有
    $S_{mh}=(S_h)^m$。
  \item[\rm(vi)] 对每个整数 $n>0$，存在整数 $m_0$，使得当 $m\geq m_0$ 时
    $S_m\subset S_+^n$。
\end{enumerate}
\end{lemma}

\begin{proof}
可以假设 $S$ 作为 $S_0$-代数由次数为 $h_i$ 的齐次元素 $f_i$
（$1\leq i\leq r$）生成，而 $M$ 作为 $S$-模由次数为 $k_j$ 的齐次元素 $x_j$
（$1\leq j\leq s$）生成。显然，$M_n$ 由所有形如

\oldpage[II]{22}
$f_1^{\alpha_1}\cdots f_r^{\alpha_r}x_j$ 的元素以 $S_0$ 中的系数作线性组合而成，
其中各 $\alpha_i$ 是非负整数且满足
$k_j+\sum_i\alpha_i h_i=n$。因为各 $h_i>0$，对每个 $j$，满足这个等式的
$(\alpha_i)$ 只有有限组。这证明了 (i) 的第一个断言；第二个断言显然。

另一方面，令 $h$ 为各 $h_i$ 的最小公倍数，并令
$g_i=f_i^{h/h_i}$（$1\leq i\leq r$），于是所有 $g_i$ 都是 $h$ 次的。
令 $z_\mu$ 遍历 $M$ 中所有形如
$f_1^{\alpha_1}\cdots f_r^{\alpha_r}x_j$ 且对
$1\leq i\leq r$ 满足 $0\leq\alpha_i<h/h_i$ 的元素。这些元素只有有限多个；
令 $n_1$ 为它们次数的最大值。显然，当 $n\geq n_1$ 时，$M_{n+h}$ 的每个元素
都是各 $z_\mu$ 的线性组合，其系数是关于各 $g_i$ 的正次数单项式；所以
$M_{n+h}=S_hM_n$，这证明了 (ii)。

用同样的方法可见，对每个 $d>0$，$M^{(d,k)}$ 的每个元素都是形如
$g^df_1^{\alpha_1}\cdots f_r^{\alpha_r}x_j$ 的元素以 $S_0$ 中的系数作线性组合，
其中 $0\leq\alpha_i<d$，而 $g$ 是 $S$ 的齐次元素。这证明了 (iii)。
在 (iii) 中取 $M=S_+$，并注意
$(S_+)^{(d)}=(S^{(d)})_+$，再应用
\egaref{II.2.1.3-zh}{2.1.3}，即得 (iv)。在 (ii) 中取 $M=S$，即得 (v)。
最后，给定 $n$，满足 $\alpha_i\geq0$ 且
$\sum_i\alpha_i<n$ 的 $(\alpha_i)$ 只有有限组。若 $m_0$ 是这些组中
$\sum_i\alpha_i h_i$ 的最大值，则 $m>m_0$ 时
$S_m\subset S_+^n$，这证明了 (vi)。
\end{proof}
% Continuation of locked French source scope: ega2-1-fr.tex lines 1526-1649,
% subsection 2.1 / printed pp.22-23.

\begin{corollary}[2.1.7]
\label{II.2.1.7-zh}
若 $S$ 是诺特环，则对每个整数 $d>0$，$S^{(d)}$ 也是诺特环。
\end{corollary}

\begin{proof}
这由 \egaref{II.2.1.5-zh}{2.1.5} 和
\egaref{II.2.1.6-zh}{2.1.6, (iv)} 得出。
\end{proof}

\begin{env}[2.1.8]
\label{II.2.1.8-zh}
设 $\mathfrak{p}$ 是分次环 $S$ 的一个分次素理想；因此
$\mathfrak{p}$ 是各子群 $\mathfrak{p}_n=\mathfrak{p}\cap S_n$ 的直和。
假设 $\mathfrak{p}$ 不包含 $S_+$。若 $f\in S_+$ 且
$f\notin\mathfrak{p}$，则关系 $f^nx\in\mathfrak{p}$ 等价于
$x\in\mathfrak{p}$。特别地，若 $f\in S_d$（$d>0$），则对每个
$x\in S_{m-nd}$，关系 $f^nx\in\mathfrak{p}_m$ 等价于
$x\in\mathfrak{p}_{m-nd}$。
\end{env}

\begin{proposition}[2.1.9]
\label{II.2.1.9-zh}
设 $n_0>0$ 是一个整数，并且对每个 $n\geq n_0$，设
$\mathfrak{p}_n$ 是 $S_n$ 的一个子群。存在一个不包含 $S_+$ 的分次素理想
$\mathfrak{p}$，使得对每个 $n\geq n_0$ 都有
$\mathfrak{p}\cap S_n=\mathfrak{p}_n$，当且仅当满足下列条件：
\begin{enumerate}
  \item[$1^\circ$] 对每个 $m\geq0$ 和每个 $n\geq n_0$，
    $S_m\mathfrak{p}_n\subset\mathfrak{p}_{m+n}$。
  \item[$2^\circ$] 若 $m\geq n_0$、$n\geq n_0$、
    $f\in S_m$、$g\in S_n$，则
    $fg\in\mathfrak{p}_{m+n}$ 蕴含
    $f\in\mathfrak{p}_m$ 或 $g\in\mathfrak{p}_n$。
  \item[$3^\circ$] 至少对一个 $n\geq n_0$，有
    $\mathfrak{p}_n\neq S_n$。
\end{enumerate}
而且，此时分次素理想 $\mathfrak{p}$ 是唯一的。
\end{proposition}

\begin{proof}
条件 $1^\circ$、$2^\circ$ 的必要性显然。此外，若
$\mathfrak{p}\not\supset S_+$，则至少存在一个 $k>0$，使
$\mathfrak{p}\cap S_k\neq S_k$。取
$f\in S_k\setminus\mathfrak{p}$。根据
\egaref{II.2.1.8-zh}{2.1.8}，关系
$\mathfrak{p}\cap S_n=S_n$ 蕴含
$\mathfrak{p}\cap S_{n-mk}=S_{n-mk}$。因此，若从某个 $n$ 起都有
$\mathfrak{p}\cap S_n=S_n$，就会得到
$\mathfrak{p}\supset S_+$，与假设矛盾。这证明了 $3^\circ$ 的必要性。

反过来，假设 $1^\circ$、$2^\circ$、$3^\circ$ 均成立。注意：若
$d\geq n_0$ 且 $f\in S_d\setminus\mathfrak{p}_d$，则在
$\mathfrak{p}$ 存在时，对 $m<n_0$，$\mathfrak{p}_m$ 必等于所有满足下述条件的
$x\in S_m$ 所成的集合：除有限多个 $r$ 外，都有
$f^rx\in\mathfrak{p}_{m+rd}$。这已经证明了 $\mathfrak{p}$ 若存在便是唯一的。
还需证明：若用上述条件定义 $m<n_0$ 时的 $\mathfrak{p}_m$，则
$\mathfrak{p}=\sum_{n=0}^{\infty}\mathfrak{p}_n$ 是素理想。首先，由
$2^\circ$，对 $m\geq n_0$，$\mathfrak{p}_m$ 也等于所有满足下述条件的
$x\in S_m$ 所成的集合：除有限多个 $r$ 外，都有
$f^rx\in\mathfrak{p}_{m+rd}$。

\oldpage[II]{23}
现在，若 $g\in S_m$、$x\in\mathfrak{p}_n$，则除有限多个 $r$ 外，
$f^rgx\in\mathfrak{p}_{m+n+rd}$，所以
$gx\in\mathfrak{p}_{m+n}$。这证明了 $\mathfrak{p}$ 是 $S$ 的理想。
为证明该理想是素的，也就是由各子群
$S_n/\mathfrak{p}_n$ 分次的环 $S/\mathfrak{p}$ 是整环，只需考察
$S/\mathfrak{p}$ 中两个元素的最高次分量，并证明：若
$x\in S_m$、$y\in S_n$ 满足
$x\notin\mathfrak{p}_m$、$y\notin\mathfrak{p}_n$，则
$xy\notin\mathfrak{p}_{m+n}$。否则，当 $r$ 充分大时会有
$f^{2r}xy\in\mathfrak{p}_{m+n+2rd}$。但对每个 $r>0$ 都有
$f^ry\notin\mathfrak{p}_{n+rd}$；于是由 $2^\circ$ 可知，除有限多个 $r$ 外，
$f^rx\in\mathfrak{p}_{m+rd}$，从而将得到
$x\in\mathfrak{p}_m$，与假设矛盾。
\end{proof}

\begin{env}[2.1.10]
\label{II.2.1.10-zh}
若 $S_+$ 的一个子集 $\mathfrak{J}$ 是 $S$ 的理想，就称它是
\emph{$S_+$ 的理想}。若 $\mathfrak{J}$ 是 $S_+$ 与 $S$ 的某个
\emph{不包含 $S_+$} 的分次素理想之交，就称 $\mathfrak{J}$ 是
\emph{$S_+$ 的分次素理想}；根据
\egaref{II.2.1.9-zh}{2.1.9}，这个素理想还是唯一的。

若 $\mathfrak{J}$ 是 $S_+$ 的理想，则
\emph{$\mathfrak{J}$ 在 $S_+$ 中的根}是 $S_+$ 中某次幂属于
$\mathfrak{J}$ 的所有元素组成的集合；换言之，它是
$r_+(\mathfrak{J})=r(\mathfrak{J})\cap S_+$。特别地，$0$ 在 $S_+$ 中的根
也称为 $S_+$ 的\emph{幂零根}，记作 $\mathfrak{N}_+$；它正是 $S_+$ 中所有
幂零元素组成的集合。若 $\mathfrak{J}$ 是 $S_+$ 的\emph{分次}理想，则它的根
$r_+(\mathfrak{J})$ 是分次理想。事实上，取商环
$S/\mathfrak{J}$ 可归约到 $\mathfrak{J}=0$ 的情形；此时只需说明：若
$x=x_h+x_{h+1}+\cdots+x_k$ 是幂零元，则每个
$x_i\in S_i$（$1\leq h\leq i\leq k$）也是幂零元。可假设 $x_k\neq0$；
$x^n$ 的最高次分量是 $x_k^n$，所以 $x_k$ 是幂零元，再对 $k$ 作归纳即可。
若分次环 $S$ 满足 $\mathfrak{N}_+=0$，也就是 $S_+$ 不含非零幂零元，就称
$S$ 是\emph{本质约化的}。
\end{env}

\begin{env}[2.1.11]
\label{II.2.1.11-zh}
注意，在分次环 $S$ 中，若元素 $x$ 是零因子，则它的最高次分量也是零因子。
若环 $S$ 的环 $S_+$（\emph{没有单位元}）非零且不含零因子，就称 $S$ 是
\emph{本质整的}。为此，只需 $S_+$ 中每个非零齐次元素都不是该环的零因子。
显然，若 $\mathfrak{p}$ 是 $S_+$ 的分次素理想，则
$S/\mathfrak{p}$ 是本质整的。

设 $S$ 是一个本质整的分次环，且 $x_0\in S_0$。若存在\emph{一个}
$S_+$ 中的非零齐次元素 $f$，使 $x_0f=0$，则
$x_0S_+=0$。事实上，对每个 $g\in S_+$，
$(x_0g)f=(x_0f)g=0$，而假设蕴含 $x_0g=0$。因此，$S$ 是整环，当且仅当
$S_0$ 是整环且 $S_+$ 在 $S_0$ 中的零化子等于 $0$。
\end{env}
% Locked French authority: ega2-1-fr.tex, 820505 B,
% SHA-256 84EDBE3E83530AF2959B441796337C9DC21EAFCA6A13114A26778760FBF437AC.
% Sequential source scope: lines 1651-1803, subsection 2.2 / printed pp.23-25.

\subsection{分次环的分式环}
\label{subsection:II.2.2-zh}

\begin{env}[2.2.1]
\label{II.2.2.1-zh}
设 $S$ 是一个按非负次数分次的环，$f$ 是 $S$ 中一个次数为 $d>0$ 的
\emph{齐次}元素。分式环 $S'=S_f$ 可按如下方式分次：令 $S'_n$ 为所有
$x/f^k$ 组成的集合，其中 $x\in S_{n+kd}$ 且 $k\geq0$（注意，这里的 $n$
可以取任意负值）。我们用 $S_{(f)}$ 表示 $S'$ 中由\emph{次数为 $0$} 的元素
组成的子环 $S'_0=(S_f)_0$。

若 $f\in S_d$，则 $S_f$ 中的单项式 $(f/1)^h$（$h$ 为正整数或负整数）
构成环 $S_{(f)}$ 上的一个\emph{自由系}，而它们的线性组合所成的集合恰为
\oldpage[II]{24}
环 $(S^{(d)})_f$；因此，该环\emph{同构于
$S_{(f)}[T,T^{-1}]=S_{(f)}\otimes_{\mathbf{Z}}\mathbf{Z}[T,T^{-1}]$}
（其中 $T$ 是不定元）。事实上，若有关系
$\sum_{h=-a}^{b}z_h(f/1)^h=0$，其中 $z_h=x_h/f^m$ 且所有 $x_h$ 都属于
$S_{md}$，则按定义，这一关系等价于存在 $k\geq a$，使得
$\sum_{h=-a}^{b}f^{h+k}x_h=0$。由于这个和式中各项的次数互不相同，
对每个 $h$ 都有 $f^{h+k}x_h=0$，从而对每个 $h$ 都有 $z_h=0$。

若 $M$ 是一个分次 $S$-模，则 $M'=M_f$ 是一个分次 $S_f$-模，其中
$M'_n$ 是所有 $z/f^k$ 组成的集合，$z\in M_{n+kd}$ 且 $k\geq0$。
用 $M_{(f)}$ 表示 $M'$ 中所有次数为 $0$ 的齐次元素组成的集合。显然，
$M_{(f)}$ 是一个 $S_{(f)}$-模，并且
$(M^{(d)})_f=M_{(f)}\otimes_{S_{(f)}}(S^{(d)})_f$。
\end{env}

\begin{lemma}[2.2.2]
\label{II.2.2.2-zh}
设 $d$、$e$ 是两个整数 $>0$，$f\in S_d$，$g\in S_e$。存在典范环同构
\[
  S_{(fg)}\xrightarrow{\sim}(S_{(f)})_{g^d/f^e}\text{；}
\]
若通过这个同构把两个环典范地等同起来，则存在典范模同构
\[
  M_{(fg)}\xrightarrow{\sim}(M_{(f)})_{g^d/f^e}\text{。}
\]
\end{lemma}

\begin{proof}
事实上，$fg$ 整除 $f^eg^d$，而后一个元素又整除 $(fg)^{de}$，所以分次环
$S_{fg}$ 与 $S_{f^eg^d}$ 典范地等同。另一方面，$S_{f^eg^d}$ 也与
$(S_{f^e})_{g^d/1}$ 典范地等同（\egaxref{0.1.4.6-fr}{0, I.4.6}）；
由于 $f^e/1$ 在 $S_{f^e}$ 中可逆，$S_{f^eg^d}$ 还与
$(S_{f^e})_{g^d/f^e}$ 典范地等同。元素 $g^d/f^e$ 在 $S_{f^e}$ 中的
次数为 $0$；因此，$(S_{f^e})_{g^d/f^e}$ 中次数为 $0$ 的元素组成的子环
就是 $(S_{(f^e)})_{g^d/f^e}$。又显然有 $S_{(f^e)}=S_{(f)}$，这就证明了
引理的第一部分；第二部分同理可证。
\end{proof}

\begin{env}[2.2.3]
\label{II.2.2.3-zh}
在 \egaref{II.2.2.2-zh}{2.2.2} 的假设下，显然，典范同态
$S_f\to S_{fg}$（\egaxref{0.1.4.1-fr}{0, I.4.1}）把 $x/f^k$ 变为
$g^kx/(fg)^k$，其次数为 $0$；因而限制后得到一个\emph{典范同态
$S_{(f)}\to S_{(fg)}$}，使图
\[
  \xymatrix{
    & S_{(f)}\ar[dl]\ar[dr]\\
    S_{(fg)}\ar[rr]^-{\sim} && (S_{(f)})_{g^d/f^e}
  }
\]
交换。同样可以定义典范同态 $M_{(f)}\to M_{(fg)}$。
\end{env}

\begin{lemma}[2.2.4]
\label{II.2.2.4-zh}
若 $f$、$g$ 是 $S_+$ 的两个齐次元素，则环 $S_{(fg)}$ 由 $S_{(f)}$ 与
$S_{(g)}$ 的典范像之并生成。
\end{lemma}

\begin{proof}
由 \egaref{II.2.2.2-zh}{2.2.2}，只需说明
$1/(g^d/f^e)=f^{d+e}/(fg)^d$ 属于 $S_{(g)}$ 在 $S_{(fg)}$ 中的典范像；
这由定义显然成立。
\end{proof}

\begin{proposition}[2.2.5]
\label{II.2.2.5-zh}
设 $d$ 是整数 $>0$，$f\in S_d$。则存在典范环同构
$S_{(f)}\xrightarrow{\sim}S^{(d)}/(f-1)S^{(d)}$。若通过这个同构把两个环
等同起来，则存在典范模同构
$M_{(f)}\xrightarrow{\sim}M^{(d)}/(f-1)M^{(d)}$。
\end{proposition}

\begin{proof}
第一个同构的定义如下：把 $x/f^n$（其中 $x\in S_{nd}$）对应到元素
$\overline{x}$，即 $x$ 模 $(f-1)S^{(d)}$ 的剩余类。这个映射是良定义的，
因为对每个 $x\in S^{(d)}$ 都有同余式
$f^hx\equiv x\pmod{(f-1)S^{(d)}}$；所以，若对某个 $h>0$ 有 $f^hx=0$，
\oldpage[II]{25}
则 $\overline{x}=0$。另一方面，若 $x\in S_{nd}$ 且
$x=(f-1)y$，其中
$y=y_{hd}+y_{(h+1)d}+\cdots+y_{kd}$、$y_{jd}\in S_{jd}$、$y_{hd}\neq0$，
则必有 $h=n$、$x=-y_{hd}$，并且
$y_{(j+1)d}=fy_{jd}$（$h\leq j\leq k-1$）、$fy_{kd}=0$；最终得到
$f^{k-n+1}x=0$。因此，可以把 $x\in S_{nd}$ 的每个模
$(f-1)S^{(d)}$ 剩余类 $\overline{x}$ 对应到 $S_{(f)}$ 中的元素 $x/f^n$，
因为上面的说明表明这个映射是良定义的。显然，这样定义的两个映射都是
环同态，并且互为逆映射。对 $M$ 完全同样处理。
\end{proof}

\begin{corollary}[2.2.6]
\label{II.2.2.6-zh}
若 $S$ 是诺特环，则对每个次数 $>0$ 的齐次元素 $f$，$S_{(f)}$ 也是
诺特环。
\end{corollary}

\begin{proof}
这立即由 \egaref{II.2.1.7-zh}{2.1.7} 和
\egaref{II.2.2.5-zh}{2.2.5} 得出。
\end{proof}

\begin{env}[2.2.7]
\label{II.2.2.7-zh}
设 $T$ 是 $S_+$ 中一个由\emph{齐次}元素组成的乘法子集；于是
$T_0=T\cup\{1\}$ 是 $S$ 的乘法子集。由于 $T_0$ 的元素都是齐次的，
环 $T_0^{-1}S$ 仍以显然的方式分次。用 $S_{(T)}$ 表示 $T_0^{-1}S$ 中
由次数为 $0$ 的元素组成的子环，也就是由形如 $x/h$ 的元素组成的子环，
其中 $h\in T$，而 $x$ 是与 $h$ 同次的齐次元素。已知
（\egaxref{0.1.4.5-fr}{0, 1.4.5}），$T_0^{-1}S$ 典范地等同于环 $S_f$
（$f$ 遍历 $T$）关于典范同态 $S_f\to S_{fg}$ 的归纳极限；由于这个等同
保持次数，它把 $S_{(T)}$ 与各 $S_{(f)}$（$f\in T$）的\emph{归纳极限}
等同起来。对每个分次 $S$-模 $M$，同样定义 $S_{(T)}$ 上的模 $M_{(T)}$：
它由 $T_0^{-1}M$ 中次数为 $0$ 的元素组成；并且可见，这个模是各
$M_{(f)}$（$f\in T$）的归纳极限。

若 $\mathfrak{p}$ 是 $S_+$ 的一个分次素理想，则分别用
$S_{(\mathfrak{p})}$ 与 $M_{(\mathfrak{p})}$ 表示环 $S_{(T)}$ 与模
$M_{(T)}$，其中 $T$ 是 $S_+$ 中所有不属于 $\mathfrak{p}$ 的\emph{齐次}
元素组成的集合。
\end{env}
% Locked French authority: ega2-1-fr.tex, 820505 B,
% SHA-256 84EDBE3E83530AF2959B441796337C9DC21EAFCA6A13114A26778760FBF437AC.
% Sequential source scope: lines 1805-1981, subsection 2.3 through proposition 2.3.6 /
% printed pp.25-27.

\subsection{分次环的齐次素谱}
\label{subsection:II.2.3-zh}

\begin{env}[2.3.1]
\label{II.2.3.1-zh}
给定一个按非负次数分次的环 $S$，称其所有分次素理想
（\egaref{II.2.1.10-zh}{2.1.10}）所成的集合为 $S$ 的\emph{齐次素谱}，
并记作 $\operatorname{Proj}(S)$；这里所说的分次素理想是 $S_+$ 的分次素理想，
或等价地，是 $S$ 中\emph{不包含} $S_+$ 的分次素理想。下面将在底集
$\operatorname{Proj}(S)$ 上定义一个\emph{概形}结构。
\end{env}

\begin{env}[2.3.2]
\label{II.2.3.2-zh}
对 $S$ 的任一子集 $E$，令 $V_+(E)$ 为 $S$ 中所有包含 $E$ 而不包含
$S_+$ 的分次素理想所成的集合；因此，它就是 $\operatorname{Spec}(S)$ 的子集
$V(E)\cap\operatorname{Proj}(S)$。由
\egaxref{I.1.1.2-fr}{I, 1.1.2} 得：
\[
  \label{II.2.3.2.1-zh}
  V_+(0)=\operatorname{Proj}(S),\quad
  V_+(S)=V_+(S_+)=\varnothing
  \tag{2.3.2.1}
\]
\[
  \label{II.2.3.2.2-zh}
  V_+\left(\bigcup_\lambda E_\lambda\right)
  =\bigcap_\lambda V_+(E_\lambda)
  \tag{2.3.2.2}
\]
\[
  \label{II.2.3.2.3-zh}
  V_+(EE')=V_+(E)\cup V_+(E')
  \tag{2.3.2.3}
\]
把 $E$ 换成由 $E$ 生成的分次理想并不改变 $V_+(E)$。此外，若
$\mathfrak{J}$ 是 $S$ 的分次理想，则
\[
  \label{II.2.3.2.4-zh}
  V_+(\mathfrak{J})
  =V_+\left(\bigcup_{q\geq n}(\mathfrak{J}\cap S_q)\right)
  \tag{2.3.2.4}
\]
\oldpage[II]{26}
对每个 $n>0$ 都成立。事实上，若
$\mathfrak{p}\in\operatorname{Proj}(S)$ 包含 $\mathfrak{J}$ 中次数
$\geq n$ 的齐次元素，则按假设，存在某个不属于 $\mathfrak{p}$ 的齐次元素
$f\in S_d$。对每个 $m\geq0$ 及每个
$x\in S_m\cap\mathfrak{J}$，除有限多个 $r$ 外，都有
$f^rx\in\mathfrak{J}\cap S_{m+rd}$，从而
$f^rx\in\mathfrak{p}\cap S_{m+rd}$；于是
$x\in\mathfrak{p}\cap S_m$（\egaref{II.2.1.9-zh}{2.1.9}）。

最后，对 $S$ 的任一分次理想 $\mathfrak{J}$，都有
\[
  \label{II.2.3.2.5-zh}
  V_+(\mathfrak{J})=V_+(r_+(\mathfrak{J})).
  \tag{2.3.2.5}
\]
\end{env}

\begin{env}[2.3.3]
\label{II.2.3.3-zh}
按定义，$V_+(E)$ 是 $X=\operatorname{Proj}(S)$ 在
$\operatorname{Spec}(S)$ 的谱拓扑所诱导的拓扑中的闭子集；这个拓扑仍称为
$X$ 上的\emph{谱拓扑}。对每个 $f\in S$，令
\[
  \label{II.2.3.3.1-zh}
  D_+(f)=D(f)\cap\operatorname{Proj}(S)
  =\operatorname{Proj}(S)\setminus V_+(f)
  \tag{2.3.3.1}
\]
于是，对 $S$ 的两个元素 $f$、$g$，有
（\egaxref{I.1.1.9.1-fr}{I, 1.1.9.1}）
\[
  \label{II.2.3.3.2-zh}
  D_+(fg)=D_+(f)\cap D_+(g).
  \tag{2.3.3.2}
\]
\end{env}

\begin{proposition}[2.3.4]
\label{II.2.3.4-zh}
当 $f$ 遍历 $S_+$ 的所有齐次元素时，各 $D_+(f)$ 构成
$X=\operatorname{Proj}(S)$ 的一个拓扑基。
\end{proposition}

\begin{proof}
事实上，由 \egaref{II.2.3.2.2-zh}{2.3.2.2} 和
\egaref{II.2.3.2.4-zh}{2.3.2.4}，$X$ 的每个闭子集都是形如
$V_+(f)$ 的集合之交，其中 $f$ 是次数 $>0$ 的齐次元素。
\end{proof}

\begin{env}[2.3.5]
\label{II.2.3.5-zh}
设 $f$ 是 $S_+$ 中次数为 $d>0$ 的\emph{齐次}元素。对 $S$ 的任一不包含
$f$ 的分次素理想 $\mathfrak{p}$，已知所有满足
$x\in\mathfrak{p}$、$n\geq0$ 的 $x/f^n$ 组成分式环 $S_f$ 的一个素理想
（\egaxref{0.1.2.6-fr}{0, 1.2.6}）。它与 $S_{(f)}$ 的交因而是该环的一个
素理想，记作 $\psi_f(\mathfrak{p})$；它由所有满足
$n\geq0$、$x\in\mathfrak{p}\cap S_{nd}$ 的 $x/f^n$ 组成。这样就定义了映射
\[
  \psi_f:D_+(f)\longrightarrow\operatorname{Spec}(S_{(f)})\text{；}
\]
此外，若 $g\in S_e$ 是 $S_+$ 的另一个齐次元素，则有交换图
\[
  \label{II.2.3.5.1-zh}
  \xymatrix{
    D_+(f)\ar[r]^{\psi_f}
    & \operatorname{Spec}(S_{(f)})\\
    D_+(fg)\ar[u]\ar[r]_{\psi_{fg}}
    & \operatorname{Spec}(S_{(fg)})\ar[u]
  }
  \tag{2.3.5.1}
\]
其中左边的竖箭头是包含映射，右边的竖箭头是由典范同态
$\omega=\omega_{fg,f}:S_{(f)}\to S_{(fg)}$ 导出的映射
${}^a\omega_{fg,f}$（\egaxref{I.1.2.1-fr}{I, 1.2.1}）。事实上，若
$x/f^n\in\omega^{-1}(\psi_{fg}(\mathfrak{p}))$，其中
$fg\notin\mathfrak{p}$，则按定义
$g^nx/(fg)^n\in\psi_{fg}(\mathfrak{p})$，所以
$g^nx\in\mathfrak{p}$，进而 $x\in\mathfrak{p}$；逆命题显然。
\end{env}

\begin{proposition}[2.3.6]
\label{II.2.3.6-zh}
映射 $\psi_f$ 是从 $D_+(f)$ 到 $\operatorname{Spec}(S_{(f)})$ 的同胚。
\end{proposition}

\begin{proof}
首先，$\psi_f$ 连续。事实上，若 $h\in S_{nd}$ 且
$h/f^n\in\psi_f(\mathfrak{p})$，则按定义 $h\in\mathfrak{p}$，反之亦然；
所以
\[
  \psi_f^{-1}(D(h/f^n))=D_+(hf)
\]
而断言由公式 \egaref{II.2.3.3.2-zh}{2.3.3.2} 得出。此外，当 $h$ 遍历
各 $S_{nd}$ 时，各 $D_+(hf)$ 由
\egaref{II.2.3.4-zh}{2.3.4} 和公式
\egaref{II.2.3.3.2-zh}{2.3.3.2} 构成 $D_+(f)$ 的一个拓扑基。因此，
\oldpage[II]{27}
结合 $D_+(f)$ 与 $\operatorname{Spec}(S_{(f)})$ 中都成立的 $(T_0)$ 公理，
上面的论证说明 $\psi_f$ 是单射，并且逆映射
$\psi_f(D_+(f))\longrightarrow D_+(f)$ 连续。最后，为证明 $\psi_f$ 满射，
设 $\mathfrak{q}_0$ 是 $S_{(f)}$ 的一个素理想；对每个 $n>0$，令
$\mathfrak{p}_n$ 为所有满足 $x^d/f^n\in\mathfrak{q}_0$ 的
$x\in S_n$ 组成的集合。各 $\mathfrak{p}_n$ 满足
\egaref{II.2.1.9-zh}{2.1.9} 的条件：事实上，若 $x\in S_n$、$y\in S_n$
满足 $x^d/f^n\in\mathfrak{q}_0$ 和
$y^d/f^n\in\mathfrak{q}_0$，则
$(x+y)^{2d}/f^{2n}\in\mathfrak{q}_0$，从而
$(x+y)^d/f^n\in\mathfrak{q}_0$，因为 $\mathfrak{q}_0$ 是素理想。
这证明各 $\mathfrak{p}_n$ 是 $S_n$ 的子群；考虑到 $\mathfrak{q}_0$ 是
素理想，\egaref{II.2.1.9-zh}{2.1.9} 的其余条件可立即验证。若
$\mathfrak{p}$ 是这样定义的 $S$ 的分次素理想，则确有
$\psi_f(\mathfrak{p})=\mathfrak{q}_0$；因为对 $x\in S_{nd}$，关系
$x/f^n\in\mathfrak{q}_0$ 与 $x^d/f^{nd}\in\mathfrak{q}_0$ 等价，
这是由于 $\mathfrak{q}_0$ 是素理想。
\end{proof}
% Locked French authority: ega2-1-fr.tex, 820505 B,
% SHA-256 84EDBE3E83530AF2959B441796337C9DC21EAFCA6A13114A26778760FBF437AC.
% Sequential source scope: lines 1983-2139, subsection 2.3 from corollary 2.3.7 /
% through 2.3.17 / printed pp.27-28.

\begin{corollary}[2.3.7]
\label{II.2.3.7-zh}
$D_+(f)=\varnothing$ 当且仅当 $f$ 是幂零元。
\end{corollary}

\begin{proof}
事实上，$\operatorname{Spec}(S_{(f)})=\varnothing$ 当且仅当
$S_{(f)}=0$，也就是 $1=0$ 在 $S_f$ 中成立；按定义，这意味着 $f$ 是幂零元。
\end{proof}

\begin{corollary}[2.3.8]
\label{II.2.3.8-zh}
设 $E$ 是 $S_+$ 的一个子集。下列条件等价：
\begin{enumerate}
  \item[\rm a)] $V_+(E)=X=\operatorname{Proj}(S)$。
  \item[\rm b)] $E$ 的每个元素都是幂零元。
  \item[\rm c)] $E$ 的每个元素的各个齐次分量都是幂零元。
\end{enumerate}
\end{corollary}

\begin{proof}
显然，c) 蕴含 b)，而 b) 蕴含 a)。若 $\mathfrak{J}$ 是 $S$ 中由 $E$
生成的分次理想，则条件 a) 等价于 $V_+(\mathfrak{J})=X$；特别地，a) 蕴含：
其中每个齐次元素 $f\in\mathfrak{J}$ 都满足 $V_+(f)=X$，因而由
\egaref{II.2.3.7-zh}{2.3.7}，$f$ 是幂零元。
\end{proof}

\begin{corollary}[2.3.9]
\label{II.2.3.9-zh}
若 $\mathfrak{J}$ 是 $S_+$ 的分次理想，则 $r_+(\mathfrak{J})$ 等于所有
包含 $\mathfrak{J}$ 的 $S_+$ 分次素理想之交。
\end{corollary}

\begin{proof}
考察分次环 $S/\mathfrak{J}$，可以归约到 $\mathfrak{J}=0$ 的情形。需要证明：
若 $f\in S_+$ 不是幂零元，则存在一个不包含 $f$ 的 $S$ 的分次素理想。
然而，$f$ 的齐次分量中至少有一个不是幂零元，因而可以假设 $f$ 是齐次的；
此时断言由 \egaref{II.2.3.7-zh}{2.3.7} 得出。
\end{proof}

\begin{env}[2.3.10]
\label{II.2.3.10-zh}
对 $X=\operatorname{Proj}(S)$ 的任一子集 $Y$，令 $\mathfrak{j}_+(Y)$ 为所有
满足 $Y\subset V_+(f)$ 的 $f\in S_+$ 组成的集合；这也等价于说
$\mathfrak{j}_+(Y)=\mathfrak{j}(Y)\cap S_+$。因此，$\mathfrak{j}_+(Y)$ 是
$S_+$ 的一个理想，并且等于它在 $S_+$ 中的根。
\end{env}

\begin{proposition}[2.3.11]
\label{II.2.3.11-zh}
\begin{enumerate}
  \item[\rm(i)] 对 $S_+$ 的任一子集 $E$，$\mathfrak{j}_+(V_+(E))$ 等于
    $S_+$ 中由 $E$ 生成的分次理想在 $S_+$ 中的根。
  \item[\rm(ii)] 对 $X$ 的任一子集 $Y$，
    $V_+(\mathfrak{j}_+(Y))=\overline{Y}$，即 $Y$ 在 $X$ 中的闭包。
\end{enumerate}
\end{proposition}

\begin{proof}
\begin{enumerate}
  \item[\rm(i)] 若 $\mathfrak{J}$ 是 $S_+$ 中由 $E$ 生成的分次理想，则
    $V_+(E)=V_+(\mathfrak{J})$，于是断言由
    \egaref{II.2.3.9-zh}{2.3.9} 得出。
  \item[\rm(ii)] 由于
    $V_+(\mathfrak{J})=\bigcap_{f\in\mathfrak{J}}V_+(f)$，关系
    $Y\subset V_+(\mathfrak{J})$ 蕴含：对每个
    $f\in\mathfrak{J}$，都有 $Y\subset V_+(f)$；因而
    $\mathfrak{j}_+(Y)\supset\mathfrak{J}$，从而
    $V_+(\mathfrak{j}_+(Y))\subset V_+(\mathfrak{J})$。按闭集的定义，
    这就证明了 (ii)。
\end{enumerate}
\end{proof}

\begin{corollary}[2.3.12]
\label{II.2.3.12-zh}
$X=\operatorname{Proj}(S)$ 的闭子集 $Y$ 与 $S_+$ 中等于自身在 $S_+$ 中之根的
分次理想，通过反序映射
$Y\longmapsto\mathfrak{j}_+(Y)$、
$\mathfrak{J}\longmapsto V_+(\mathfrak{J})$ 一一对应。两个闭子集
$X$ 的两个闭子集之并 $Y_1\cup Y_2$ 对应于
\oldpage[II]{28}
$\mathfrak{j}_+(Y_1)\cap\mathfrak{j}_+(Y_2)$；任意一族闭子集
$(Y_\lambda)$ 的交对应于各 $\mathfrak{j}_+(Y_\lambda)$ 之和在 $S_+$ 中的根。
\end{corollary}

\begin{corollary}[2.3.13]
\label{II.2.3.13-zh}
设 $\mathfrak{J}$ 是 $S_+$ 中的分次理想。$V_+(\mathfrak{J})=\varnothing$
当且仅当 $S_+$ 的每个元素都有某个幂属于 $\mathfrak{J}$。
\end{corollary}

最后一个推论也可等价地表述为下列形式之一：

\begin{corollary}[2.3.14]
\label{II.2.3.14-zh}
设 $(f_\alpha)$ 是 $S_+$ 的一族齐次元素。各 $D_+(f_\alpha)$ 构成
$X=\operatorname{Proj}(S)$ 的一个覆盖，当且仅当 $S_+$ 的每个元素都有某个幂
属于由各 $f_\alpha$ 生成的理想。
\end{corollary}

\begin{corollary}[2.3.15]
\label{II.2.3.15-zh}
设 $(f_\alpha)$ 是 $S_+$ 的一族齐次元素，$f$ 是 $S_+$ 的一个元素。
下列关系等价：
\begin{enumerate}
  \item[\rm a)] $D_+(f)\subset\bigcup_\alpha D_+(f_\alpha)$；
  \item[\rm b)] $V_+(f)\supset\bigcap_\alpha V_+(f_\alpha)$；
  \item[\rm c)] $f$ 的某个幂属于由各 $f_\alpha$ 生成的理想。
\end{enumerate}
\end{corollary}

\begin{corollary}[2.3.16]
\label{II.2.3.16-zh}
$X=\operatorname{Proj}(S)$ 为空，当且仅当 $S_+$ 的每个元素都是幂零元。
\end{corollary}

\begin{corollary}[2.3.17]
\label{II.2.3.17-zh}
在 \egaref{II.2.3.12-zh}{2.3.12} 所述的一一对应中，$X$ 的不可约闭子集
对应于 $S_+$ 中的分次素理想。
\end{corollary}

\begin{proof}
事实上，若 $Y=Y_1\cup Y_2$，其中 $Y_1$、$Y_2$ 是闭集且都不等于 $Y$，则
\[
  \mathfrak{j}_+(Y)
  =\mathfrak{j}_+(Y_1)\cap\mathfrak{j}_+(Y_2)
\]
而理想 $\mathfrak{j}_+(Y_1)$、$\mathfrak{j}_+(Y_2)$ 都不等于
$\mathfrak{j}_+(Y)$，所以 $\mathfrak{j}_+(Y)$ 不是素理想。反过来，若
$\mathfrak{J}$ 是 $S_+$ 中一个非素的分次理想，则存在 $S_+$ 的两个元素
$f$、$g$，使得 $f\notin\mathfrak{J}$、$g\notin\mathfrak{J}$、
$fg\in\mathfrak{J}$。于是
$V_+(f)\not\supset V_+(\mathfrak{J})$、
$V_+(g)\not\supset V_+(\mathfrak{J})$，但由
\egaref{II.2.3.2.3-zh}{2.3.2.3}，
$V_+(\mathfrak{J})\subset V_+(f)\cup V_+(g)$。因此，
$V_+(\mathfrak{J})$ 是两个闭集
$V_+(f)\cap V_+(\mathfrak{J})$ 与
$V_+(g)\cap V_+(\mathfrak{J})$ 的并，而这两个闭集都不等于
$V_+(\mathfrak{J})$。
\end{proof}
% Locked French authority: ega2-1-fr.tex, 820505 B,
% SHA-256 84EDBE3E83530AF2959B441796337C9DC21EAFCA6A13114A26778760FBF437AC.
% Sequential source scope: lines 2141-2259, subsection II.2.4 through 2.4.4 /
% printed pp.28-29.

\subsection{\texorpdfstring{$\operatorname{Proj}(S)$}{Proj(S)} 上的概形结构}
\label{subsection:II.2.4-zh}

\begin{env}[2.4.1]
\label{II.2.4.1-zh}
设 $f$、$g$ 是 $S_+$ 的两个齐次元素；考察仿射概形
$Y_f=\operatorname{Spec}(S_{(f)})$、
$Y_g=\operatorname{Spec}(S_{(g)})$ 和
$Y_{fg}=\operatorname{Spec}(S_{(fg)})$。由
\egaref{II.2.2.2-zh}{2.2.2}，从 $Y_{fg}$ 到 $Y_f$ 的态射
$w_{fg,f}=({}^a\omega_{fg,f},\widetilde{\omega}_{fg,f})$，它对应于典范同态
$\omega_{fg,f}:S_{(f)}\to S_{(fg)}$，是一个\emph{开浸入}
（\egaref{I.1.3.6-zh}{I, 1.3.6}）。借助
$\psi_f:D_+(f)\to Y_f$ 的逆同胚（\egaref{II.2.3.6-zh}{2.3.6}），可以把
$Y_f$ 的仿射概形结构搬到 $D_+(f)$ 上；由图表
\egaref{II.2.3.5.1-zh}{2.3.5.1} 的交换性，仿射概形 $D_+(fg)$ 因而等同于
仿射概形 $D_+(f)$ 的底空间的开子集 $D_+(fg)$ 上所诱导的概形。于是，考虑到
\egaref{II.2.3.4-zh}{2.3.4}，显然 $X=\operatorname{Proj}(S)$ 带有唯一的
\emph{预概形}结构，并且它在每个 $D_+(f)$ 上的限制就是刚才定义的仿射概形。
此外：
\end{env}

\begin{proposition}[2.4.2]
\label{II.2.4.2-zh}
预概形 $\operatorname{Proj}(S)$ 是一个概形。
\end{proposition}

\begin{proof}
只需证明（\egaref{I.5.5.6-zh}{I, 5.5.6}）：对 $S_+$ 中任意两个齐次元素
$f$、$g$，$D_+(f)\cap D_+(g)=D_+(fg)$ 都是仿射的，而且它的环由
$D_+(f)$ 和 $D_+(g)$ 的环的典范像生成。第一点由定义显然成立，第二点来自
\egaref{II.2.2.4-zh}{2.2.4}。
\end{proof}
\oldpage[II]{29}
以后把齐次素谱 $\operatorname{Proj}(S)$ 称为\emph{概形}时，所指的总是刚才定义的
结构。

\begin{example}[2.4.3]
\label{II.2.4.3-zh}
取 $S=K[T_1,T_2]$，其中 $K$ 是一个域，$T_1$、$T_2$ 是两个不定元，而 $S$
按总次数分次。由 \egaref{II.2.3.14-zh}{2.3.14}，
$\operatorname{Proj}(S)$ 是 $D_+(T_1)$ 与 $D_+(T_2)$ 的并；立即可见，这两个
仿射概形都与 $K[T]$ 典范同构，而 $\operatorname{Proj}(S)$ 由
\egaref{I.2.3.2-zh}{I, 2.3.2} 所述的这两个仿射概形的粘合得到（参见
\egaref{II.7.4.14-zh}{7.4.14}）。
\end{example}

\begin{proposition}[2.4.4]
\label{II.2.4.4-zh}
设 $S$ 是一个按非负次数分次的环，并设 $X$ 是概形
$\operatorname{Proj}(S)$。
\begin{enumerate}
  \item[\rm(i)] 若 $\mathfrak{N}_+$ 是 $S_+$ 的幂零根
    （\egaref{II.2.1.10-zh}{2.1.10}），则概形 $X_{\mathrm{red}}$ 与
    $\operatorname{Proj}(S/\mathfrak{N}_+)$ 典范同构；特别地，若 $S$ 本质约化，
    则 $\operatorname{Proj}(S)$ 是约化的。
  \item[\rm(ii)] 假设 $S$ 本质约化。则 $X$ 为整概形的充分必要条件是 $S$
    本质整。
\end{enumerate}
\end{proposition}

\begin{proof}
\begin{enumerate}
  \item[\rm(i)] 设 $\overline{S}$ 是分次环 $S/\mathfrak{N}_+$，并用
    $x\mapsto\overline{x}$ 表示 $0$ 次典范同态 $S\to\overline{S}$。对每个
    $f\in S_d$（$d>0$），典范同态 $S_f\to\overline{S}_{\overline{f}}$
    （\egaref{0.1.5.1-zh}{0, 1.5.1}）是满的且为 $0$ 次，因而限制后给出满同态
    $S_{(f)}\to\overline{S}_{(\overline{f})}$。若假设
    $f\notin\mathfrak{N}_+$，则立即验证出
    $\overline{S}_{(\overline{f})}$ 是约化的，而且前一同态的核是
    $S_{(f)}$ 的幂零根；换言之，
    $\overline{S}_{(\overline{f})}=(S_{(f)})_{\mathrm{red}}$。因此，这一同态对应于
    一个闭浸入 $D_+(\overline{f})\to D_+(f)$，它把
    $D_+(\overline{f})$ 等同于 $(D_+(f))_{\mathrm{red}}$
    （\egaref{I.5.1.2-zh}{I, 5.1.2}），并且特别是这两个仿射概形的底空间之间的
    一个同胚。此外，若 $g\notin\mathfrak{N}_+$ 是 $S_+$ 的另一个齐次元素，则图表
    \[
      \xymatrix{
        S_{(f)}\ar[r]\ar[d]
        & \overline{S}_{(\overline{f})}\ar[d]\\
        S_{(fg)}\ar[r]
        & \overline{S}_{(\overline{fg})}
      }
    \]
    交换；又因为当 $f$ 遍历所有满足次数 $>0$ 且
    $f\notin\mathfrak{N}_+$ 的齐次元素时，各 $D_+(f)$ 构成
    $X=\operatorname{Proj}(S)$ 的一个覆盖
    （\egaref{II.2.3.7-zh}{2.3.7}），所以各态射
    $D_+(\overline{f})\to D_+(f)$ 是某个闭浸入
    $\operatorname{Proj}(\overline{S})\to\operatorname{Proj}(S)$ 的限制，而这个闭浸入
    是底空间之间的同胚；结论由此得出
    （\egaref{I.5.1.2-zh}{I, 5.1.2}）。
  \item[\rm(ii)] 假设 $S$ 本质整；换言之，$(0)$ 是 $S_+$ 中一个不同于 $S_+$
    的分次素理想。则由 (i)，$X$ 是约化的；由
    \egaref{II.2.3.17-zh}{2.3.17}，它又是不可约的。反过来，假设 $S$ 本质约化且
    $X$ 是整概形。于是，对 $S_+$ 中任一非零齐次元素 $f\neq0$，有
    $D_+(f)\neq\varnothing$（\egaref{II.2.3.7-zh}{2.3.7}）。$X$ 不可约这一假设
    蕴含：对 $S_+$ 中任意两个齐次元素 $f$、$g$，只要二者均满足 $\neq0$，就有
    $D_+(f)\cap D_+(g)\neq\varnothing$；因而由
    \egaref{II.2.3.3.2-zh}{2.3.3.2}，$fg\neq0$。所以 $S_+$ 没有零因子，
    由此得证。
\end{enumerate}
\end{proof}
% Locked French authority: ega2-1-fr.tex, 820505 B,
% SHA-256 84EDBE3E83530AF2959B441796337C9DC21EAFCA6A13114A26778760FBF437AC.
% Sequential source scope: lines 2261-2360, subsection II.2.4 remainder /
% printed pp.29-30.

\begin{env}[2.4.5]
\label{II.2.4.5-zh}
给定一个交换环 $A$。回顾：称一个分次环 $S$ 是一个\emph{分次 $A$-代数}，
如果它带有一个 $A$-代数结构，使得每个子群 $S_n$ 都是一个 $A$-模。为此，只需
$S_0$ 是一个
\oldpage[II]{30}
$A$-代数；换言之，只须在 $S_0$ 上定义一个 $A$-代数结构，并对
$\alpha\in A$、$x\in S_n$ 令 $\alpha\cdot x=(\alpha\cdot1)x$，就定义了
$S$ 上的分次 $A$-代数结构。
\end{env}

\begin{proposition}[2.4.6]
\label{II.2.4.6-zh}
假设 $S$ 是一个分次 $A$-代数。则在 $X=\operatorname{Proj}(S)$ 上，结构层
$\mathcal{O}_X$ 是一个 $A$-代数层（这里把 $A$ 看作 $X$ 上的常值层）；换言之，
$X$ 是 $\operatorname{Spec}(A)$ 上的一个概形。
\end{proposition}

\begin{proof}
只须注意：对 $S_+$ 中任一齐次元素 $f$，$S_{(f)}$ 是 $A$ 上的一个代数，而且对
$S_+$ 中任意齐次元素 $f$、$g$，图表
\[
  \xymatrix{
    S_{(f)}\ar[rr] && S_{(fg)}\\
    & A\ar[ul]\ar[ur] &
  }
\]
交换。
\end{proof}

\begin{proposition}[2.4.7]
\label{II.2.4.7-zh}
设 $S$ 是一个按非负次数分次的环。
\begin{enumerate}
  \item[\rm(i)] 对每个整数 $d>0$，存在概形 $\operatorname{Proj}(S)$ 到概形
    $\operatorname{Proj}(S^{(d)})$ 的一个典范同构。
  \item[\rm(ii)] 设 $S'$ 是满足 $S'_0=\mathbf{Z}$、
    $S'_n=S_n$（视为 $\mathbf{Z}$-模，$n>0$）的分次环。存在概形
    $\operatorname{Proj}(S)$ 到概形 $\operatorname{Proj}(S')$ 的一个典范同构。
\end{enumerate}
\end{proposition}

\begin{proof}
\begin{enumerate}
  \item[\rm(i)] 先证明映射
    $\mathfrak{p}\mapsto\mathfrak{p}\cap S^{(d)}$ 是集合
    $\operatorname{Proj}(S)$ 到 $\operatorname{Proj}(S^{(d)})$ 的双射。事实上，
    设给定一个分次素理想
    $\mathfrak{p}'\in\operatorname{Proj}(S^{(d)})$，并令
    $\mathfrak{p}_{nd}=\mathfrak{p}'\cap S_{nd}$（$n\geq0$）。对每个不是 $d$ 的
    倍数的 $n>0$，把 $\mathfrak{p}_n$ 定义为所有满足
    $x^d\in\mathfrak{p}_{nd}$ 的 $x\in S_n$ 组成的集合。若
    $x\in\mathfrak{p}_n$、$y\in\mathfrak{p}_n$，则
    $(x+y)^{2d}\in\mathfrak{p}_{2nd}$；因为 $\mathfrak{p}'$ 是素理想，所以
    $(x+y)^d\in\mathfrak{p}_{nd}$。立即可见，对所有 $n\geq0$ 如此定义的
    $\mathfrak{p}_n$ 满足 \egaref{II.2.1.9-zh}{2.1.9} 的条件，因而存在唯一的
    素理想 $\mathfrak{p}\in\operatorname{Proj}(S)$，使得
    $\mathfrak{p}\cap S^{(d)}=\mathfrak{p}'$。因为对 $S_+$ 中每个齐次元素 $f$，
    都有 $V_+(f)=V_+(f^d)$（\egaref{II.2.3.2.3-zh}{2.3.2.3}），所以前一双射是
    拓扑空间之间的同胚。最后，沿用同样的记号，$S_{(f)}$ 和
    $S^{(d)}_{(f^d)}$ 典范等同（\egaref{II.2.2.2-zh}{2.2.2}），所以
    $\operatorname{Proj}(S)$ 和 $\operatorname{Proj}(S^{(d)})$ 作为概形典范等同。
  \item[\rm(ii)] 若对每个 $\mathfrak{p}\in\operatorname{Proj}(S)$，让它对应于
    唯一的素理想 $\mathfrak{p}'\in\operatorname{Proj}(S')$，使得对每个 $n>0$
    都有 $\mathfrak{p}'\cap S_n=\mathfrak{p}\cap S_n$
    （\egaref{II.2.1.9-zh}{2.1.9}），则显然定义了底空间之间的典范同胚
    $\operatorname{Proj}(S)\xrightarrow{\sim}\operatorname{Proj}(S')$；事实上，当
    $f$ 是 $S_+$ 的齐次元素时，$V_+(f)$ 对 $S$ 和 $S'$ 是同一个集合。又因为
    $S_{(f)}=S'_{(f)}$，所以 $\operatorname{Proj}(S)$ 和
    $\operatorname{Proj}(S')$ 作为概形彼此等同。
\end{enumerate}
\end{proof}

\begin{corollary}[2.4.8]
\label{II.2.4.8-zh}
若 $S$ 是一个分次 $A$-代数，$S'_A$ 是满足 $(S'_A)_0=A$、
$(S'_A)_n=S_n$（$n>0$）的分次 $A$-代数，则存在
$\operatorname{Proj}(S)$ 到 $\operatorname{Proj}(S'_A)$ 的一个典范同构。
\end{corollary}

\begin{proof}
事实上，沿用 \egaref{II.2.4.7-zh}{2.4.7, (ii)} 的记号，这两个概形都与
$\operatorname{Proj}(S')$ 典范同构。
\end{proof}
% Locked French authority: ega2-1-fr.tex, 820505 B,
% SHA-256 84EDBE3E83530AF2959B441796337C9DC21EAFCA6A13114A26778760FBF437AC.
% Sequential source scope: lines 2362-2544, subsection II.2.5 through 2.5.9 /
% printed pp.30-32.

\subsection{分次模的伴随层}
\label{subsection:II.2.5-zh}

\begin{env}[2.5.1]
\label{II.2.5.1-zh}
设 $M$ 是一个分次 $S$-模。对 $S_+$ 中任一齐次元素 $f$，$M_{(f)}$ 是一个
$S_{(f)}$-模，因而在仿射概形 $\operatorname{Spec}(S_{(f)})$ 上有与之相伴的
拟凝聚层 $\widetilde{M_{(f)}}$；这里把该仿射概形与 $D_+(f)$ 等同
（\egaref{I.1.3.4-zh}{I, 1.3.4}）。
\end{env}
\oldpage[II]{31}

\begin{proposition}[2.5.2]
\label{II.2.5.2-zh}
在 $X=\operatorname{Proj}(S)$ 上，存在唯一的拟凝聚 $\mathcal{O}_X$-模
$\widetilde{M}$，使得对 $S_+$ 中任一齐次元素 $f$，都有
$\Gamma(D_+(f),\widetilde{M})=M_{(f)}$；而当 $f$、$g$ 是 $S_+$ 中的齐次元素时，
限制同态
$\Gamma(D_+(f),\widetilde{M})\to\Gamma(D_+(fg),\widetilde{M})$
就是典范同态 $M_{(f)}\to M_{(fg)}$
（\egaref{II.2.2.3-zh}{2.2.3}）。
\end{proposition}

\begin{proof}
设 $f\in S_d$、$g\in S_e$。由 \egaref{II.2.2.2-zh}{2.2.2}，
$D_+(fg)$ 可与 $(S_{(f)})_{g^d/f^e}$ 的素谱等同，所以 $D_+(f)$ 上的层
$\widetilde{M_{(f)}}$ 限制到 $D_+(fg)$ 后，典范地等同于模
$(M_{(f)})_{g^d/f^e}$ 的伴随层
（\egaref{I.1.3.6-zh}{I, 1.3.6}），因而也等同于
$\widetilde{M_{(fg)}}$（\egaref{II.2.2.2-zh}{2.2.2}）。由此存在典范同构
\[
  \theta_{g,f}:\widetilde{M_{(f)}}|D_+(fg)
  \xrightarrow{\sim}\widetilde{M_{(g)}}|D_+(fg),
\]
并且若 $h$ 是 $S_+$ 中第三个齐次元素，则在 $D_+(fgh)$ 中有
$\theta_{f,h}=\theta_{f,g}\circ\theta_{g,h}$。于是由
\egaref{0.3.3.1-zh}{0, 3.3.1}，在 $X$ 上存在一个拟凝聚
$\mathcal{O}_X$-模 $\mathcal{F}$，并且对 $S_+$ 中每个齐次元素 $f$，存在从
$\mathcal{F}|D_+(f)$ 到 $\widetilde{M_{(f)}}$ 的同构 $\eta_f$，使
$\theta_{g,f}=\eta_g\circ\eta_f^{-1}$。再考察 $X$ 的由各 $D_+(f)$ 组成的拓扑基上以
$D_+(f)\to M_{(f)}$ 定义的预层，其中以典范同态
$M_{(f)}\to M_{(fg)}$ 作为限制同态，并记其层化为 $\mathcal{G}$。考虑到
\egaref{I.1.3.7-zh}{I, 1.3.7}，以上论证表明 $\mathcal{F}$ 与
$\mathcal{G}$ 同构。把 $\mathcal{G}$ 记作 $\widetilde{M}$，它确实满足命题中的
条件。特别地，$\widetilde{S}=\mathcal{O}_X$。
\end{proof}

\begin{definition}[2.5.3]
\label{II.2.5.3-zh}
称 \egaref{II.2.5.2-zh}{2.5.2} 中定义的拟凝聚
$\mathcal{O}_X$-模 $\widetilde{M}$ 与分次 $S$-模 $M$ \emph{相伴}。
\end{definition}

回顾：若把分次模同态限制为\emph{$0$ 次同态}，分次 $S$-模就构成一个范畴。
在这一约定下：

\begin{proposition}[2.5.4]
\label{II.2.5.4-zh}
函子 $M\mapsto\widetilde{M}$ 是从分次 $S$-模范畴到拟凝聚
$\mathcal{O}_X$-模范畴的正合加性协变函子，并且与归纳极限和直和交换。
\end{proposition}

\begin{proof}
这些性质都是局部的，所以只须在
$\widetilde{M}|D_+(f)=\widetilde{M_{(f)}}$ 上验证。事实上，函子
$M\mapsto M_f$、分次 $S_f$-模范畴中的函子 $N\mapsto N_0$，以及
$S_{(f)}$-模范畴中的伴随层函子 $P\mapsto\widetilde{P}$，都具有上述正合性，
且都与归纳极限和直和交换
（\egaref{I.1.3.5-zh}{I, 1.3.5} 与
\egaref{I.1.3.9-zh}{I, 1.3.9}）。命题由此得证。
\end{proof}

若 $u:M\to N$ 是一个 $0$ 次同态，则把与之对应的同态
$\widetilde{M}\to\widetilde{N}$ 记作 $\widetilde{u}$。由
\egaref{II.2.5.4-zh}{2.5.4} 立即推出：
\egaref{I.1.3.9-zh}{I, 1.3.9} 与
\egaref{I.1.3.10-zh}{I, 1.3.10} 的结果，对分次 $S$-模及其 $0$ 次同态仍然
成立（其中 $\widetilde{M}$ 取本节赋予的意义），因为那些证明纯属形式推演。

\begin{proposition}[2.5.5]
\label{II.2.5.5-zh}
对每个 $\mathfrak{p}\in X=\operatorname{Proj}(S)$，有
$\widetilde{M}_{\mathfrak{p}}=M_{(\mathfrak{p})}$。
\end{proposition}

\begin{proof}
按定义，
$\widetilde{M}_{\mathfrak{p}}=
\varinjlim\Gamma(D_+(f),\widetilde{M})$，其中 $f$ 遍历 $S_+$ 中所有满足
$f\notin\mathfrak{p}$ 的齐次元素。因为
$\Gamma(D_+(f),\widetilde{M})=M_{(f)}$，命题由
$M_{(\mathfrak{p})}$ 的定义得出
（\egaref{II.2.2.7-zh}{2.2.7}）。
\end{proof}
\oldpage[II]{32}

特别地，\emph{局部环 $(\mathcal{O}_X)_{\mathfrak{p}}$} 正是环
$S_{(\mathfrak{p})}$；它由分式 $x/f$ 组成，其中 $f$ 是 $S_+$ 中不属于
$\mathfrak{p}$ 的齐次元素，而齐次元素 $x$ 与 $f$ 具有相同次数。

更特别地，若 $S$ \emph{本质整}，从而
$\operatorname{Proj}(S)=X$ 是\emph{整概形}
（\egaref{II.2.4.4-zh}{2.4.4}），且 $\xi=(0)$ 是 $X$ 的泛点，则
\emph{有理函数域} $R(X)=\mathcal{O}_\xi$ 由元素 $fg^{-1}$ 组成，其中
$f$、$g$ 是 $S_+$ 中次数相同的齐次元素，且 $g\neq0$。

\begin{proposition}[2.5.6]
\label{II.2.5.6-zh}
若对每个 $z\in M$ 和 $S_+$ 中每个齐次元素 $f$，都存在 $f$ 的某个幂零化
$z$，则 $\widetilde{M}=0$。当 $S_0$-代数 $S$ 由次数为 $1$ 的齐次元素集
$S_1$ 生成时，这一充分条件也是必要的。
\end{proposition}

\begin{proof}
条件 $\widetilde{M}=0$ 等价于：对 $S_+$ 中每个齐次元素 $f$，都有
$M_{(f)}=0$。另一方面，若 $f\in S_d$，则 $M_{(f)}=0$ 意味着：对每个次数
为 $d$ 的倍数的齐次元素 $z\in M$，都存在一个幂 $f^n$ 使
$f^nz=0$。因此，若对每个 $f\in S_1$ 都有 $M_{(f)}=0$，命题中的条件便对
每个 $f\in S_1$ 成立；而当 $S_1$ 生成 $S$ 时，它自然也对 $S_+$ 中每个
齐次元素 $f$ 成立，因为 $S_+$ 中每个齐次元素都是 $S_1$ 中元素之积的线性组合。
\end{proof}

\begin{proposition}[2.5.7]
\label{II.2.5.7-zh}
设 $d$ 是一个整数（$>0$），$f\in S_d$。则对每个 $n\in\mathbf{Z}$，
$(\mathcal{O}_X|D_+(f))$-模
$\widetilde{S(nd)}|D_+(f)$ 与 $\mathcal{O}_X|D_+(f)$ 典范同构。
\end{proposition}

\begin{proof}
$S_f$ 中可逆元素 $(f/1)^n$ 的乘法给出从
$S_{(f)}=(S_f)_0$ 到
$(S_f)_{nd}=(S_f(nd))_0=(S(nd)_f)_0=S(nd)_{(f)}$ 的双射。换言之，
$S_{(f)}$-模 $S_{(f)}$ 与 $S(nd)_{(f)}$ 典范同构；命题由此得证。
\end{proof}

\begin{corollary}[2.5.8]
\label{II.2.5.8-zh}
在开集
\[
  U=\bigcup_{f\in S_d}D_+(f)
\]
上，$\mathcal{O}_X$-模 $\widetilde{S(nd)}$ 的限制是一个可逆的
$(\mathcal{O}_X|U)$-模
（\egaref{0.5.4.1-zh}{0, 5.4.1}）。
\end{corollary}

\begin{corollary}[2.5.9]
\label{II.2.5.9-zh}
若 $S$ 的理想 $S_+$ 由次数为 $1$ 的齐次元素集 $S_1$ 生成，则对每个
$n\in\mathbf{Z}$，$\mathcal{O}_X$-模 $\widetilde{S(n)}$ 都可逆。
\end{corollary}

\begin{proof}
由假设和 \egaref{II.2.3.14-zh}{2.3.14}，
$X=\bigcup_{f\in S_1}D_+(f)$；只须对 $U=X$ 应用
\egaref{II.2.5.8-zh}{2.5.8}。
\end{proof}
% Locked French authority: ega2-1-fr.tex, 820505 B,
% SHA-256 84EDBE3E83530AF2959B441796337C9DC21EAFCA6A13114A26778760FBF437AC.
% Sequential source scope: lines 2546-2694, subsection II.2.5.10-2.5.12 /
% printed pp.32-34.

\begin{env}[2.5.10]
\label{II.2.5.10-zh}
在本节后文中，对每个 $n\in\mathbf{Z}$，都令
\[
\label{II.2.5.10.1-zh}
  \mathcal{O}_X(n)=\widetilde{S(n)}
\tag{2.5.10.1}
\]
并且对 $X$ 的每个开子集 $U$、每个 $(\mathcal{O}_X|U)$-模
$\mathcal{F}$ 以及每个 $n\in\mathbf{Z}$，令
\[
\label{II.2.5.10.2-zh}
  \mathcal{F}(n)=
  \mathcal{F}\otimes_{\mathcal{O}_X|U}(\mathcal{O}_X(n)|U)
\tag{2.5.10.2}
\]
若理想 $S_+$ 由 $S_1$ 生成，则对每个 $n\in\mathbf{Z}$，函子
$\mathcal{F}(n)$ 关于 $\mathcal{F}$ 是\emph{正合的}，因为
$\mathcal{O}_X(n)$ 此时是一个可逆 $\mathcal{O}_X$-模。
\end{env}

\begin{env}[2.5.11]
\label{II.2.5.11-zh}
设 $M$、$N$ 是两个分次 $S$-模。对每个 $f\in S_d$（$d>0$），定义函子性的
典范 $S_{(f)}$-模同态
\[
\label{II.2.5.11.1-zh}
  \lambda_f:M_{(f)}\otimes_{S_{(f)}}N_{(f)}
  \longrightarrow(M\otimes_S N)_{(f)}
\tag{2.5.11.1}
\]
\oldpage[II]{33}
如下：把由嵌入 $M_{(f)}\to M_f$、$N_{(f)}\to N_f$ 与
$S_{(f)}\to S_f$ 所得的同态
$M_{(f)}\otimes_{S_{(f)}}N_{(f)}\to M_f\otimes_{S_f}N_f$，同典范同构
$M_f\otimes_{S_f}N_f\xrightarrow{\sim}(M\otimes_S N)_f$
（\egaref{0.1.3.4-zh}{0, 1.3.4}）复合；再注意到，按两个分次模之张量积的
定义，后一同构保持次数。于是，对
$x\in M_{md}$、$y\in N_{nd}$（$m\geq0$、$n\geq0$），有
\[
  \lambda_f((x/f^m)\otimes(y/f^n))=(x\otimes y)/f^{m+n}.
\]

由定义立即可知：若 $g\in S_e$（$e>0$），则图表
\[
  \xymatrix{
    M_{(f)}\otimes_{S_{(f)}}N_{(f)}\ar[r]^{\lambda_f}\ar[d]
    & (M\otimes_S N)_{(f)}\ar[d]\\
    M_{(fg)}\otimes_{S_{(fg)}}N_{(fg)}\ar[r]_{\lambda_{fg}}
    & (M\otimes_S N)_{(fg)}
  }
\]
交换；其中右侧竖直箭头是典范同态，左侧竖直箭头由各典范同态诱导。因此，各
$\lambda$ 诱导出函子性的典范 $\mathcal{O}_X$-模同态
\[
\label{II.2.5.11.2-zh}
  \lambda:\widetilde{M}\otimes_{\mathcal{O}_X}\widetilde{N}
  \longrightarrow\widetilde{M\otimes_S N}.
\tag{2.5.11.2}
\]

特别考察 $S$ 的两个分次理想 $\mathfrak{J}$、$\mathfrak{K}$。由于
$\widetilde{\mathfrak{J}}$、$\widetilde{\mathfrak{K}}$ 是
$\mathcal{O}_X$ 的理想层，存在典范同态
$\widetilde{\mathfrak{J}}\otimes_{\mathcal{O}_X}
\widetilde{\mathfrak{K}}\to\mathcal{O}_X$；此时图表
\[
\label{II.2.5.11.3-zh}
  \xymatrix{
    \widetilde{\mathfrak{J}}\otimes_{\mathcal{O}_X}
    \widetilde{\mathfrak{K}}\ar[rr]^\lambda\ar[dr]
    && \widetilde{\mathfrak{J}\otimes_S\mathfrak{K}}\ar[dl]\\
    & \mathcal{O}_X &
  }
\tag{2.5.11.3}
\]
交换。事实上，只须在每个开集 $D_+(f)$ 上验证，其中 $f$ 是 $S_+$ 中的齐次
元素；这立即来自 $\lambda_f$ 的定义
\egaref{II.2.5.11.1-zh}{2.5.11.1} 与
\egaref{I.1.3.13-zh}{I, 1.3.13}。

最后注意：若 $M$、$N$、$P$ 是三个分次 $S$-模，则图表
\[
\label{II.2.5.11.4-zh}
  \xymatrix{
    \widetilde{M}\otimes_{\mathcal{O}_X}\widetilde{N}
    \otimes_{\mathcal{O}_X}\widetilde{P}
    \ar[r]^{\lambda\otimes1}\ar[d]_{1\otimes\lambda}
    & \widetilde{M\otimes_S N}\otimes_{\mathcal{O}_X}\widetilde{P}
    \ar[d]^\lambda\\
    \widetilde{M}\otimes_{\mathcal{O}_X}\widetilde{N\otimes_S P}
    \ar[r]_\lambda
    & \widetilde{M\otimes_S N\otimes_S P}
  }
\tag{2.5.11.4}
\]
\oldpage[II]{34}
交换。仍然只须在每个 $D_+(f)$ 上验证，而这立即由定义和
\egaref{I.1.3.13-zh}{I, 1.3.13} 得出。
\end{env}

\begin{env}[2.5.12]
\label{II.2.5.12-zh}
在 \egaref{II.2.5.11-zh}{2.5.11} 的假设下，定义函子性的典范
$S_{(f)}$-模同态
\[
\label{II.2.5.12.1-zh}
  \mu_f:\bigl(\operatorname{Hom}_S(M,N)\bigr)_{(f)}
  \longrightarrow
  \operatorname{Hom}_{S_{(f)}}(M_{(f)},N_{(f)})
\tag{2.5.12.1}
\]
如下：若 $u$ 是一个 $nd$ 次同态，则把 $u/f^n$ 对应到同态
$M_{(f)}\to N_{(f)}$，后者把
$x/f^m$（$x\in M_{md}$）变为 $u(x)/f^{m+n}$。对
$g\in S_e$（$e>0$），仍有交换图表
\[
  \xymatrix{
    \bigl(\operatorname{Hom}_S(M,N)\bigr)_{(f)}
    \ar[r]^{\mu_f}\ar[d]
    & \operatorname{Hom}_{S_{(f)}}(M_{(f)},N_{(f)})\ar[d]\\
    \bigl(\operatorname{Hom}_S(M,N)\bigr)_{(fg)}
    \ar[r]_{\mu_{fg}}
    & \operatorname{Hom}_{S_{(fg)}}(M_{(fg)},N_{(fg)})
  }
\]
其中左侧竖直箭头是典范同态，右侧竖直箭头由各典范同态诱导。考虑到
\egaref{I.1.3.8-zh}{I, 1.3.8}，由此又得到：各 $\mu_f$ 定义函子性的典范
$\mathcal{O}_X$-模同态
\[
\label{II.2.5.12.2-zh}
  \mu:\widetilde{\operatorname{Hom}_S(M,N)}
  \longrightarrow
  \mathcal{H}om_{\mathcal{O}_X}(\widetilde{M},\widetilde{N}).
\tag{2.5.12.2}
\]
\end{env}
% Locked French authority: ega2-1-fr.tex, 820505 B,
% SHA-256 84EDBE3E83530AF2959B441796337C9DC21EAFCA6A13114A26778760FBF437AC.
% Sequential source scope: lines 2696-2791, subsection II.2.5.13 and proof /
% printed pp.34-35.

\begin{proposition}[2.5.13]
\label{II.2.5.13-zh}
假设理想 $S_+$ 由 $S_1$ 生成。则同态 $\lambda$
（\egaref{II.2.5.11.2-zh}{2.5.11.2}）是同构；若分次 $S$-模 $M$ 有限表示
（\egaref{II.2.1.1-zh}{2.1.1}），则同态 $\mu$
（\egaref{II.2.5.12.2-zh}{2.5.12.2}）也是同构。
\end{proposition}

\begin{proof}
由 \egaref{II.2.3.14-zh}{2.3.14}，$X$ 是所有 $D_+(f)$（$f\in S_1$）
的并，所以在上述假设下，只须证明：当 $f$ 是一个\emph{次数为 $1$}的齐次元素时，
$\lambda_f$ 与 $\mu_f$ 是同构。此时可以定义一个 $\mathbf{Z}$-双线性映射
\[
  M_m\times N_n\longrightarrow
  M_{(f)}\otimes_{S_{(f)}}N_{(f)}
\]
把 $(x,y)$ 对应到元素 $(x/f^m)\otimes(y/f^n)$；当 $m<0$ 时，记号
$x/f^m$ 表示 $f^{-m}x/1$。这些映射定义一个 $\mathbf{Z}$-线性映射
\[
  M\otimes_{\mathbf{Z}}N\longrightarrow
  M_{(f)}\otimes_{S_{(f)}}N_{(f)}.
\]
若 $s\in S_q$，则此映射把 $(sx)\otimes y$ 变为
\[
  (s/f^q)\bigl((x/f^m)\otimes(y/f^n)\bigr)
\]
（其中 $x\in M_m$、$y\in N_n$）。因此得到关于典范同态
$S\to S_{(f)}$ 的模双同态
\[
  \gamma_f:M\otimes_S N\longrightarrow
  M_{(f)}\otimes_{S_{(f)}}N_{(f)},
\]
这里该典范同态把 $s\in S_q$ 变为 $s/f^q$。再设
$M\otimes_S N$ 的某个元素
$\sum_i(x_i\otimes y_i)$ 满足
\[
  f^r\sum_i(x_i\otimes y_i)=0,
\]
其中 $x_i$、$y_i$ 是次数分别为 $m_i$、$n_i$ 的齐次元素；换言之，
$\sum_i(f^rx_i\otimes y_i)=0$。由
\egaref{0.1.3.4-zh}{0, 1.3.4} 得
\[
  \sum_i(f^rx_i/f^{m_i+r})\otimes(y_i/f^{n_i})=0,
\]
即 $\gamma_f(\sum_i(x_i\otimes y_i))=0$。所以 $\gamma_f$ 分解为
\[
  M\otimes_S N\longrightarrow(M\otimes_S N)_f
  \xrightarrow{\gamma'_f}
  M_{(f)}\otimes_{S_{(f)}}N_{(f)}.
\]
若 $\lambda'_f$ 是 $\gamma'_f$ 在
\oldpage[II]{35}
$(M\otimes_S N)_{(f)}$ 上的限制，则立即验证出 $\lambda_f$ 与
$\lambda'_f$ 是互逆的 $S_{(f)}$-同态；命题第一部分得证。

为证明第二部分，设 $M$ 是某个分次 $S$-模同态 $P\to Q$ 的余核，其中
$P$、$Q$ 都是有限多个形如 $S(n)$ 的模的直和。利用
$\operatorname{Hom}_S(L,N)$ 关于 $L$ 的左正合性，以及 $M_{(f)}$ 关于
$M$ 的正合性，立即归结为证明 $M=S(n)$ 时 $\mu_f$ 是同构。
对 $N$ 中任一齐次元素 $z$，令 $u_z$ 为从 $S(n)$ 到 $N$ 且满足
$u_z(1)=z$ 的同态；立即可见，$\eta:z\mapsto u_z$ 是从
$N(-n)$ 到 $\operatorname{Hom}_S(S(n),N)$ 的 $0$ 次同构。与之对应有同构
\[
  \eta_f:\bigl(N(-n)\bigr)_{(f)}
  \longrightarrow
  \bigl(\operatorname{Hom}_S(S(n),N)\bigr)_{(f)}.
\]
另一方面，令 $\eta'_f$ 为同构
\[
  N_{(f)}\longrightarrow
  \operatorname{Hom}_{S_{(f)}}\bigl(S(n)_{(f)},N_{(f)}\bigr),
\]
它把每个 $z'\in N_{(f)}$ 对应到同态 $v_{z'}$，其中
$v_{z'}(s/f^k)=sz'/f^{n+k}$
（$s\in S_{n+k}=(S(n))_k$）。容易看出，复合映射
\[
  \bigl(N(-n)\bigr)_{(f)}
  \xrightarrow{\eta_f}
  \bigl(\operatorname{Hom}_S(S(n),N)\bigr)_{(f)}
  \xrightarrow{\mu_f}
  \operatorname{Hom}_{S_{(f)}}\bigl(S(n)_{(f)},N_{(f)}\bigr)
  \xrightarrow{{\eta'_f}^{-1}}N_{(f)}
\]
就是从 $\bigl(N(-n)\bigr)_{(f)}$ 到 $N_{(f)}$ 的同构
$z/f^h\mapsto z/f^{h-n}$；故 $\mu_f$ 是同构。
\end{proof}
% Locked French authority: ega2-1-fr.tex, 820505 B,
% SHA-256 84EDBE3E83530AF2959B441796337C9DC21EAFCA6A13114A26778760FBF437AC.
% Sequential source scope: lines 2793-2907, subsection II.2.5 remainder /
% printed pp.35-36.

若理想 $S_+$ 由 $S_1$ 生成，则由
\egaref{II.2.5.13-zh}{2.5.13}，对 $S$ 的任一分次理想 $\mathfrak{J}$ 和任一
分次 $S$-模 $M$，在相差一个典范同构的意义下有
\[
\label{II.2.5.13.1-zh}
  \widetilde{\mathfrak{J}}\cdot\widetilde{M}
  =\widetilde{\mathfrak{J}\cdot M}
\tag{2.5.13.1}
\]
事实上，这来自图表
\[
  \xymatrix{
    \widetilde{\mathfrak{J}}\otimes_{\mathcal{O}_X}\widetilde{M}
    \ar[rr]^\lambda\ar[dr]
    && \widetilde{\mathfrak{J}\otimes_S M}\ar[dl]\\
    & \widetilde{M} &
  }
\]
的交换性；其验证与 \egaref{II.2.5.11.3-zh}{2.5.11.3} 相同。

\begin{corollary}[2.5.14]
\label{II.2.5.14-zh}
假设 $S$ 由 $S_1$ 生成。则对 $\mathbf{Z}$ 中任意 $m$、$n$，在相差一个典范同构的
意义下有
\[
\label{II.2.5.14.1-zh}
  \mathcal{O}_X(m)\otimes_{\mathcal{O}_X}\mathcal{O}_X(n)
  =\mathcal{O}_X(m+n)
\tag{2.5.14.1}
\]
\[
\label{II.2.5.14.2-zh}
  \mathcal{O}_X(n)=\bigl(\mathcal{O}_X(1)\bigr)^{\otimes n}
\tag{2.5.14.2}
\]
\end{corollary}

\begin{proof}
第一式来自 \egaref{II.2.5.13-zh}{2.5.13}，以及 $0$ 次典范同构
$S(m)\otimes_S S(n)\xrightarrow{\sim}S(m+n)$ 的存在。它把元素
$1\otimes1$ 对应到
$1\in(S(m+n))_{-m-n}$；这里第一个因子 $1$ 属于 $(S(m))_{-m}$，第二个属于
$(S(n))_{-n}$。于是，只须对 $n=-1$ 证明第二式；再由
\egaref{II.2.5.13-zh}{2.5.13}，这归结为证明
$\operatorname{Hom}_S(S(1),S)$ 与 $S(-1)$ 典范同构。回到定义
\egaref{II.2.1.2-zh}{2.1.2}，并注意到 $S(1)$ 是一个单生成 $S$-模，便立即
得到这一点。
\end{proof}

\oldpage[II]{36}
\begin{corollary}[2.5.15]
\label{II.2.5.15-zh}
假设 $S$ 由 $S_1$ 生成。则对任一分次 $S$-模 $M$ 及任一
$n\in\mathbf{Z}$，在相差一个典范同构的意义下有
\[
\label{II.2.5.15.1-zh}
  \widetilde{M(n)}=\widetilde{M}(n)
\tag{2.5.15.1}
\]
\end{corollary}

\begin{proof}
这来自定义 \egaref{II.2.5.10.2-zh}{2.5.10.2}、
\egaref{II.2.5.10.1-zh}{2.5.10.1}，命题
\egaref{II.2.5.13-zh}{2.5.13}，以及 $0$ 次典范同构
$M(n)\xrightarrow{\sim}M\otimes_S S(n)$ 的存在；它把每个
$z\in(M(n))_h=M_{n+h}$ 对应到
$z\otimes1\in M_{n+h}\otimes(S(n))_{-n}\subset(M\otimes_S S(n))_h$。
\end{proof}

\begin{env}[2.5.16]
\label{II.2.5.16-zh}
以 $S'$ 表示满足 $S'_0=\mathbf{Z}$、$S'_n=S_n$（$n>0$）的分次环。若
$f\in S_d$（$d>0$），则对每个 $n\in\mathbf{Z}$，有
$(S(n))_{(f)}=(S'(n))_{(f)}$：事实上，$(S'(n))_{(f)}$ 的元素形如
$x/f^k$，其中 $x\in S'_{n+kd}$（$k>0$），而总可以把 $k$ 取得使
$n+kd\neq0$。由于 $X=\operatorname{Proj}(S)$ 与
$X'=\operatorname{Proj}(S')$ 典范等同
（\egaref{II.2.4.7-zh}{2.4.7, (ii)}），可见对每个 $n\in\mathbf{Z}$，
$\mathcal{O}_X(n)$ 与 $\mathcal{O}_{X'}(n)$ 在上述等同下互为像。

另一方面，注意到对每个 $d>0$、每个 $n\in\mathbf{Z}$，都有
\[
  (S^{(d)}(n))_h=S_{(n+h)d}=(S(nd))_{hd}.
\]
故对 $f\in S_d$，有
$(S^{(d)}(n))_{(f)}=(S(nd))_{(f)}$。由
\egaref{II.2.4.7-zh}{2.4.7, (i)}，概形
$X=\operatorname{Proj}(S)$ 与
$X^{(d)}=\operatorname{Proj}(S^{(d)})$ 典范等同。以上说明：若
$S_0$-代数 $S^{(d)}$ 由 $S_d$ 生成，则对每个 $n\in\mathbf{Z}$，
$\mathcal{O}_X(nd)$ 与 $\mathcal{O}_{X^{(d)}}(n)$ 在这一等同下互为像。
\end{env}

\begin{proposition}[2.5.17]
\label{II.2.5.17-zh}
设 $d$ 是一个整数（$>0$），且
$U=\bigcup_{f\in S_d}D_+(f)$。则对每个整数 $n$，典范同态
$\mathcal{O}_X(nd)\otimes_{\mathcal{O}_X}\mathcal{O}_X(-nd)
\to\mathcal{O}_X$ 在 $U$ 上的限制是同构。
\end{proposition}

\begin{proof}
由 \egaref{II.2.5.16-zh}{2.5.16}，可以归结到 $d=1$ 的情形；结论于是来自
\egaref{II.2.5.13-zh}{2.5.13} 的证明。
\end{proof}
% Locked French authority: ega2-1-fr.tex, 820505 B,
% SHA-256 84EDBE3E83530AF2959B441796337C9DC21EAFCA6A13114A26778760FBF437AC.
% Sequential source scope: lines 2909-3000, subsection II.2.6 through 2.6.2 /
% printed pp.36-37.

\subsection{\texorpdfstring{$\operatorname{Proj}(S)$ 上一个层所对应的分次 $S$-模}
  {Proj(S) 上一个层所对应的分次 S-模}}
\label{subsection:II.2.6-zh}

\emph{本节始终假设 $S$ 的理想 $S_+$ 由次数为 $1$ 的齐次元素所成的集合
$S_1$ 生成。}

\begin{env}[2.6.1]
\label{II.2.6.1-zh}
此时，$\mathcal{O}_X$-模 $\mathcal{O}_X(1)$ 是\emph{可逆的}
（\egaref{II.2.5.9-zh}{2.5.9}）；于是对任一 $\mathcal{O}_X$-模
$\mathcal{F}$，置（\egaxref{0.5.4.6-zh}{0, 5.4.6}）
\[
\label{II.2.6.1.1-zh}
  \Gamma_*(\mathcal{F})
  =\Gamma_*(\mathcal{O}_X(1),\mathcal{F})
  =\bigoplus_{n\in\mathbf{Z}}\Gamma(X,\mathcal{F}(n))
\tag{2.6.1.1}
\]
这里已利用 \egaref{II.2.5.14.2-zh}{2.5.14.2}。回顾
\egaxref{0.5.4.6-zh}{0, 5.4.6}，$\Gamma_*(\mathcal{O}_X)$ 带有一个
\emph{分次环}结构，而 $\Gamma_*(\mathcal{F})$ 带有一个
$\Gamma_*(\mathcal{O}_X)$ 上的\emph{分次模}结构。

由于 $\mathcal{O}_X(n)$ 是局部自由的，$\Gamma_*(\mathcal{F})$ 关于
$\mathcal{F}$ 是一个协变加性\emph{左正合}函子；特别地，若
$\mathcal{J}$ 是 $\mathcal{O}_X$ 的一个理想层，则
$\Gamma_*(\mathcal{J})$ 典范地等同于 $\Gamma_*(\mathcal{O}_X)$ 的一个
\emph{分次理想}。
\end{env}

\begin{env}[2.6.2]
\label{II.2.6.2-zh}
设 $M$ 是一个分次 $S$-模；对任一 $f\in S_d$（$d>0$），映射
$x\mapsto x/1$ 是阿贝尔群同态 $M_0\to M_{(f)}$。又因 $M_{(f)}$ 典范地
\oldpage[II]{37}
等同于 $\Gamma(D_+(f),\widetilde{M})$，故由此得到阿贝尔群同态
\[
  \alpha_0^f:M_0\longrightarrow\Gamma(D_+(f),\widetilde{M}).
\]
显然，对任一 $g\in S_e$（$e>0$），图表
\[
  \xymatrix{
    & \Gamma(D_+(f),\widetilde{M})\ar[dd]\\
    M_0\ar[ur]^{\alpha_0^f}\ar[dr]_{\alpha_0^{fg}}\\
    & \Gamma(D_+(fg),\widetilde{M})
  }
\]
是交换的；这意味着，对任一 $x\in M_0$，$\widetilde{M}$ 的截面
$\alpha_0^f(x)$ 与 $\alpha_0^g(x)$ 在
$D_+(f)\cap D_+(g)$ 上相同。因而存在唯一的截面
$\alpha_0(x)\in\Gamma(X,\widetilde{M})$，它在每个 $D_+(f)$ 上的限制都是
$\alpha_0^f(x)$。这样便定义了一个阿贝尔群同态（这里并未使用 $S$ 由
$S_1$ 生成这一假设）
\[
\label{II.2.6.2.1-zh}
  \alpha_0:M_0\longrightarrow\Gamma(X,\widetilde{M}).
\tag{2.6.2.1}
\]

把这一结果应用于分次 $S$-模 $M(n)$（$n\in\mathbf{Z}$ 任意），便对每个
$n\in\mathbf{Z}$ 得到阿贝尔群同态
\[
\label{II.2.6.2.2-zh}
  \alpha_n:M_n=(M(n))_0\longrightarrow\Gamma(X,\widetilde{M}(n))
\tag{2.6.2.2}
\]
（这里已利用 \egaref{II.2.5.15-zh}{2.5.15}）；从而得到分次阿贝尔群的一个
函子性同态（次数为 $0$）
\[
\label{II.2.6.2.3-zh}
  \alpha:M\longrightarrow\Gamma_*(\widetilde{M})
\tag{2.6.2.3}
\]
（也记作 $\alpha_M$），它在每个 $M_n$ 上都与 $\alpha_n$ 相同。

特别取 $M=S$。由 $\Gamma_*(\mathcal{O}_X)$ 中乘法的定义
（\egaxref{0.5.4.6-zh}{0, 5.4.6}），立即验证出
$\alpha:S\to\Gamma_*(\mathcal{O}_X)$ 是分次环同态，并且对任一分次
$S$-模 $M$，\egaref{II.2.6.2.3-zh}{2.6.2.3} 是分次模的一个双同态。
\end{env}
% Locked French authority: ega2-1-fr.tex, 820505 B,
% SHA-256 84EDBE3E83530AF2959B441796337C9DC21EAFCA6A13114A26778760FBF437AC.
% Sequential source scope: lines 3002-3109, subsection II.2.6.3 through 2.6.5 /
% printed pp.37-38.

\begin{proposition}[2.6.3]
\label{II.2.6.3-zh}
对任一 $f\in S_d$（$d>0$），$D_+(f)$ 恰好等于所有使
$\mathcal{O}_X(d)$ 的截面 $\alpha_d(f)$ 不消失的
$\mathfrak{p}\in X$ 所成的集合（\egaxref{0.5.5.2-zh}{0, 5.5.2}）。
\end{proposition}

\begin{proof}
按假设 $X=\bigcup_{g\in S_1}D_+(g)$，所以只须验证：对任一
$g\in S_1$，使 $\alpha_d(f)$ 不消失的所有
$\mathfrak{p}\in D_+(g)$ 所成的集合恰好是 $D_+(fg)$。事实上，$\alpha_d(f)$ 在
$D_+(g)$ 上的限制按定义是与 $(S(d))_{(g)}$ 中元素 $f/1$ 对应的截面；在
典范同构 $(S(d))_{(g)}\xrightarrow{\sim}S_{(g)}$
（\egaref{II.2.5.7-zh}{2.5.7}）之下，$\mathcal{O}_X(d)$ 在 $D_+(g)$ 上的
这个截面等同于 $\mathcal{O}_X$ 在 $D_+(g)$ 上与 $S_{(g)}$ 中元素
$f/g^d$ 对应的截面。说这个截面在
$\mathfrak{p}\in D_+(g)$ 处消失，就是说 $f/g^d\in\mathfrak{q}$，其中
$\mathfrak{q}$ 是 $S_{(g)}$ 中与 $\mathfrak{p}$ 对应的素理想
（\egaref{II.2.3.6-zh}{2.3.6}）；按定义，这就是说
$f\in\mathfrak{p}$，命题得证。
\end{proof}

\begin{env}[2.6.4]
\label{II.2.6.4-zh}
现在设 $\mathcal{F}$ 是一个 $\mathcal{O}_X$-模，并置
$M=\Gamma_*(\mathcal{F})$；由于存在分次环同态
$\alpha:S\to\Gamma_*(\mathcal{O}_X)$，可以把 $M$ 看成一个分次
$S$-模。对任一 $f\in S_d$（$d>0$），由
\egaref{II.2.6.3-zh}{2.6.3}，$\mathcal{O}_X(d)$ 的截面
$\alpha_d(f)$ 在 $D_+(f)$ 上的限制是可逆的；因而对任一 $n>0$，
$\mathcal{O}_X(nd)$ 的截面 $\alpha_{nd}(f^n)$\footnote{法文印本在本段三个幂次项中
均把 alpha 的下标误排为 d；本译依齐次次数分别采用 nd、kd、nd。
参见源文勘误记录 EGA2-ZH-DEF-0002 及守护者处置 EGA2-CANON-DISP-0002。}在 $D_+(f)$ 上的限制也是
可逆的。于是设
$z\in M_{nd}=\Gamma(X,\mathcal{F}(nd))$（$n>0$）；若存在整数
$k\geq0$，使 $f^kz$ 在 $D_+(f)$ 上的限制，即
% EGA2-CANON-DISP-0002 confirms the printed index error and permits the
% explicit reversible nd/kd/nd target readings; French source remains unchanged.
\oldpage[II]{38}
截面
$(z|D_+(f))(\alpha_{kd}(f^k)|D_+(f))$（它是
$\mathcal{F}((n+k)d)$ 的截面）为零，则由前面的说明，也有
$z|D_+(f)=0$。这表明，可以定义一个 $S_{(f)}$-同态
$\beta_f:M_{(f)}\to\Gamma(D_+(f),\mathcal{F})$，把元素 $z/f^n$ 对应到
$\mathcal{F}$ 在 $D_+(f)$ 上的截面
$(z|D_+(f))(\alpha_{nd}(f^n)|D_+(f))^{-1}$。此外立即验证出图表
\[
\label{II.2.6.4.1-zh}
  \xymatrix{
    M_{(f)}\ar[r]^{\beta_f}\ar[d]
    & \Gamma(D_+(f),\mathcal{F})\ar[d]\\
    M_{(fg)}\ar[r]_{\beta_{fg}}
    & \Gamma(D_+(fg),\mathcal{F})
  }
\tag{2.6.4.1}
\]
对 $g\in S_e$（$e>0$）是交换的。回顾 $M_{(f)}$ 典范地等同于
$\Gamma(D_+(f),\widetilde{M})$，并且所有 $D_+(f)$ 构成 $X$ 的一个拓扑基
（\egaref{II.2.3.4-zh}{2.3.4}），可见这些 $\beta_f$ 来自唯一的典范
$\mathcal{O}_X$-模同态
\[
\label{II.2.6.4.2-zh}
  \beta:\widetilde{\Gamma_*(\mathcal{F})}
  \longrightarrow\mathcal{F}
\tag{2.6.4.2}
\]
（也记作 $\beta_{\mathcal{F}}$），它显然是函子性的。
\end{env}

\begin{proposition}[2.6.5]
\label{II.2.6.5-zh}
设 $M$ 是一个分次 $S$-模，$\mathcal{F}$ 是一个 $\mathcal{O}_X$-模；则复合
同态
\[
\label{II.2.6.5.1-zh}
  \widetilde{M}\xrightarrow{\widetilde{\alpha}}
  \widetilde{\Gamma_*(\widetilde{M})}
  \xrightarrow{\beta}\widetilde{M}
\tag{2.6.5.1}
\]
\[
\label{II.2.6.5.2-zh}
  \Gamma_*(\mathcal{F})\xrightarrow{\alpha}
  \Gamma_*(\widetilde{\Gamma_*(\mathcal{F})})
  \xrightarrow{\Gamma_*(\beta)}\Gamma_*(\mathcal{F})
\tag{2.6.5.2}
\]
都是恒等同构。
\end{proposition}

\begin{proof}
对 \egaref{II.2.6.5.1-zh}{2.6.5.1} 的验证是局部的：在开集
$D_+(f)$ 中，它立即来自定义以及如下事实：应用于拟凝聚层的 $\beta$
由其在 $D_+(f)$ 上截面的作用所确定
（\egaxref{I.1.3.8-zh}{I, 1.3.8}）。对
\egaref{II.2.6.5.2-zh}{2.6.5.2} 的验证按次数分别进行：置
$M=\Gamma_*(\mathcal{F})$，则
$M_n=\Gamma(X,\mathcal{F}(n))$，且
$(\Gamma_*(\widetilde{M}))_n=\Gamma(X,\widetilde{M}(n))
=\Gamma(X,\widetilde{M(n)})$。又若 $f\in S_1$、$z\in M_n$，则
$\alpha_n^f(z)$ 是 $(M(n))_{(f)}$ 中的元素 $z/1$，它等于
$(f/1)^n(z/f^n)$；经 $\beta_f$，它所对应的是 $D_+(f)$ 上的截面
\[
  ((\alpha_1(f))^n|D_+(f))
  ((z|D_+(f))((\alpha_1(f))^n|D_+(f))^{-1}),
\]
也就是 $z$ 在 $D_+(f)$ 上的限制；这便验证了
\egaref{II.2.6.5.2-zh}{2.6.5.2}。
\end{proof}
% Locked French authority: ega2-1-fr.tex, 820505 B,
% SHA-256 84EDBE3E83530AF2959B441796337C9DC21EAFCA6A13114A26778760FBF437AC.
% Sequential source scope: lines 3111-3211, subsection II.2.7 through 2.7.3 /
% printed pp.38-40.

\subsection{有限性条件}
\label{subsection:II.2.7-zh}

\begin{proposition}[2.7.1]
\label{II.2.7.1-zh}
\begin{enumerate}
  \item[\rm(i)] 若 $S$ 是一个诺特分次环，则
    $X=\operatorname{Proj}(S)$ 是一个诺特概形。
  \item[\rm(ii)] 若 $S$ 是一个有限型分次 $A$-代数，则
    $X=\operatorname{Proj}(S)$ 是 $Y=\operatorname{Spec}(A)$ 上的一个
    有限型概形。
\end{enumerate}
\end{proposition}

\oldpage[II]{39}
\begin{proof}
\begin{enumerate}
  \item[\rm(i)] 若 $S$ 是诺特的，则理想 $S_+$ 有一组有限的齐次生成元
    $f_i$（$1\leq i\leq p$）。因此由
    \egaref{II.2.3.14-zh}{2.3.14}，底空间 $X$ 是所有
    $D_+(f_i)=\operatorname{Spec}(S_{(f_i)})$ 的并；问题归结为证明每个
    $S_{(f_i)}$ 都是诺特的，而这来自
    \egaref{II.2.2.6-zh}{2.2.6}。
  \item[\rm(ii)] 此假设蕴涵 $S_0$ 是有限型 $A$-代数，且 $S$ 是有限型
    $S_0$-代数，因而 $S_+$ 是有限型理想
    （\egaref{II.2.1.4-zh}{2.1.4}）。与 (i) 一样，问题归结为证明：对
    $f\in S_d$，$S_{(f)}$ 是有限型 $A$-代数。由
    \egaref{II.2.2.5-zh}{2.2.5}，只须证明 $S^{(d)}$ 是有限型
    $A$-代数，而这来自 \egaref{II.2.1.6-zh}{2.1.6}。
\end{enumerate}
\end{proof}

\begin{env}[2.7.2]
\label{II.2.7.2-zh}
下文将考虑分次 $S$-模 $M$ 的如下有限性条件：
\begin{enumerate}
  \item[\rm(TF)] 存在整数 $n$，使子模
    $\bigoplus_{k\geq n}M_k$ 是有限型 $S$-模。
  \item[\rm(TN)] 存在整数 $n$，使 $M_k=0$ 对 $k\geq n$ 成立。
\end{enumerate}

若 $M$ 满足 (TN)，则对 $S_+$ 中任一齐次元素 $f$，都有
$M_{(f)}=0$，从而 $\widetilde{M}=0$。

设 $M$、$N$ 是两个分次 $S$-模。称次数为 $0$ 的同态
$u:M\to N$ 是 (TN)-\emph{单射的}（相应地，(TN)-\emph{满射的}、
(TN)-\emph{双射的}），如果存在整数 $n$，使
$u_k:M_k\to N_k$ 对 $k\geq n$ 是单射的（相应地，满射的、双射的）。
所以，说 $u$ 是 (TN)-单射的（相应地，(TN)-满射的），等价于说
$\operatorname{Ker}u$（相应地，$\operatorname{Coker}u$）满足 (TN)。由
\egaref{II.2.5.4-zh}{2.5.4}，若 $u$ 是 (TN)-单射的（相应地，
(TN)-满射的、(TN)-双射的），则 $\widetilde{u}$ 是单射的（相应地，
满射的、双射的）；当 $u$ 是 (TN)-双射的时，也称 $u$ 是一个
(TN)-\emph{同构}。
\end{env}

\begin{proposition}[2.7.3]
\label{II.2.7.3-zh}
设 $S$ 是一个分次环，且理想 $S_+$ 为有限型；设 $M$ 是一个分次 $S$-模。
\begin{enumerate}
  \item[\rm(i)] 若 $M$ 满足条件 (TF)，则 $\mathcal{O}_X$-模
    $\widetilde{M}$ 为有限型。
  \item[\rm(ii)] 假设 $M$ 满足 (TF)。则 $\widetilde{M}=0$ 当且仅当
    $M$ 满足 (TN)。
\end{enumerate}
\end{proposition}

\begin{proof}
刚已看到，条件 (TN) 蕴涵 $\widetilde{M}=0$。若 $M$ 满足 (TF)，则按假设
为有限型的分次子模 $M'=\bigoplus_{k\geq n}M_k$ 使 $M/M'$ 满足
(TN)；从而 $\widetilde{M/M'}=0$。函子 $M\mapsto\widetilde{M}$ 的正合性
（\egaref{II.2.5.4-zh}{2.5.4}）蕴涵
$\widetilde{M}=\widetilde{M'}$。因此，为证明 $\widetilde{M}$ 为有限型，
可归结到 $M$ 为有限型的情形。这个问题是局部的，故只须证明
$M_{(f)}$ 是有限型 $S_{(f)}$-模
（\egaxref{I.1.3.9-zh}{I, 1.3.9}）；而 $M^{(d)}$ 是有限型
$S^{(d)}$-模（\egaref{II.2.1.6-zh}{2.1.6, (iii)}），断言于是来自
\egaref{II.2.2.5-zh}{2.2.5}。

现在假设 $M$ 满足 (TF) 且 $\widetilde{M}=0$。沿用上述记号，有
$\widetilde{M'}=0$，而 $M'$ 的条件 (TN) 等价于 $M$ 的条件 (TN)。因此，
为证明 $\widetilde{M}=0$ 蕴涵 $M$ 满足 (TN)，仍可只考虑 $M$ 由有限多个
齐次元素 $x_i$（$1\leq i\leq p$）生成的情形。另一方面，设
$(f_j)_{1\leq j\leq q}$ 是理想 $S_+$ 的一组齐次生成元。按假设，对每个
$j$ 都有 $M_{(f_j)}=0$，所以存在整数 $n$，使
$f_j^nx_i=0$ 对任意 $i$、$j$ 成立。设 $n_j$ 是 $f_j$ 的次数，并令
$m$ 为 $\sum_jr_jn_j$ 在所有满足 $\sum_jr_j\leq nq$ 的有限整数族
$(r_j)$ 上所取的最大值；于是显然

\oldpage[II]{40}
当 $k>m$ 时，对任一 $i$ 都有 $S_kx_i=0$。若 $h$ 是各 $x_i$ 次数中的
最大者，便推出 $M_k=0$ 对 $k>h+m$ 成立，证明完毕。
\end{proof}
% Locked French authority: ega2-1-fr.tex, 820505 B,
% SHA-256 84EDBE3E83530AF2959B441796337C9DC21EAFCA6A13114A26778760FBF437AC.
% Sequential source scope: lines 3213-3290, subsection II.2.7.4 through 2.7.8 /
% printed pp.40-41.

\begin{corollary}[2.7.4]
\label{II.2.7.4-zh}
设 $S$ 是一个分次环，且理想 $S_+$ 为有限型。则
$X=\operatorname{Proj}(S)=\varnothing$，当且仅当存在 $n$，使
$S_k=0$ 对 $k\geq n$ 成立。
\end{corollary}

\begin{proof}
事实上，条件 $X=\varnothing$ 等价于
$\mathcal{O}_X=\widetilde{S}=0$，而 $S$ 是单生成 $S$-模。
\end{proof}

\begin{theorem}[2.7.5]
\label{II.2.7.5-zh}
假设理想 $S_+$ 由有限多个次数为 $1$ 的齐次元素生成；设
$X=\operatorname{Proj}(S)$。则对任一拟凝聚 $\mathcal{O}_X$-模
$\mathcal{F}$，典范同态
$\beta:\widetilde{\Gamma_*(\mathcal{F})}
\to\mathcal{F}$（\egaref{II.2.6.4-zh}{2.6.4}）是同构。
\end{theorem}

\begin{proof}
事实上，若 $S_+$ 由有限多个元素 $f_i\in S_1$ 生成，则 $X$ 是各子空间
$\operatorname{Spec}(S_{(f_i)})$ 的并
（\egaref{II.2.3.6-zh}{2.3.6}）；这些子空间拟紧，所以它拟紧。此外，
$X$ 是概形（\egaref{II.2.4.2-zh}{2.4.2}）。由
\egaxref{I.9.3.2-zh}{I, 9.3.2}、
\egaref{II.2.5.14.2-zh}{2.5.14.2} 和
\egaref{II.2.6.3-zh}{2.6.3}，对任一 $f\in S_d$（$d>0$），有典范同构
$(\Gamma_*(\mathcal{F}))_{(\alpha_d(f))}\xrightarrow{\sim}
\Gamma(D_+(f),\mathcal{F})$。另一方面，按定义，
$(\Gamma_*(\mathcal{F}))_{(\alpha_d(f))}$（这里把
$\Gamma_*(\mathcal{F})$ 看作 $\Gamma_*(\mathcal{O}_X)$-模）正是
$(\Gamma_*(\mathcal{F}))_{(f)}$（这里把 $\Gamma_*(\mathcal{F})$ 看作
$S$-模）。回到上述典范同构的定义
（\egaxref{I.9.3.1-zh}{I, 9.3.1}），便看出它与同态 $\beta_f$ 相同，
定理得证。
\end{proof}

\begin{remark}[2.7.6]
\label{II.2.7.6-zh}
若再假设分次环 $S$ 是\emph{诺特的}，则只要假设理想 $S_+$ 由所有次数为
$1$ 的齐次元素组成的集合 $S_1$ 生成，
\egaref{II.2.7.5-zh}{2.7.5} 的条件就\emph{自动}满足。
\end{remark}

\begin{corollary}[2.7.7]
\label{II.2.7.7-zh}
在 \egaref{II.2.7.5-zh}{2.7.5} 的假设下，任一拟凝聚
$\mathcal{O}_X$-模 $\mathcal{F}$ 都同构于形如 $\widetilde{M}$ 的
$\mathcal{O}_X$-模，其中 $M$ 是一个分次 $S$-模。
\end{corollary}

\begin{corollary}[2.7.8]
\label{II.2.7.8-zh}
在 \egaref{II.2.7.5-zh}{2.7.5} 的假设下，任一有限型拟凝聚
$\mathcal{O}_X$-模 $\mathcal{F}$ 都同构于形如 $\widetilde{N}$ 的
$\mathcal{O}_X$-模，其中 $N$ 是一个有限型分次 $S$-模。
\end{corollary}

\begin{proof}
可以假设 $\mathcal{F}=\widetilde{M}$，其中 $M$ 是一个分次 $S$-模
（\egaref{II.2.7.7-zh}{2.7.7}）。设
$(f_\lambda)_{\lambda\in L}$ 是 $M$ 的一组齐次生成元；对 $L$ 的任一有限
子集 $H$，设 $M_H$ 是由所有满足 $\lambda\in H$ 的 $f_\lambda$ 生成的
$M$ 的分次子模。显然，$M$ 是其各子模 $M_H$ 的归纳极限，因而
$\mathcal{F}$ 是其各 $\mathcal{O}_X$-子模 $\widetilde{M_H}$ 的归纳极限
（\egaref{II.2.5.4-zh}{2.5.4}）。但由于 $\mathcal{F}$ 为有限型且底空间
$X$ 拟紧，由 \egaxref{0.5.2.3-zh}{0, 5.2.3}，存在 $L$ 的某个有限子集
$H$，使 $\mathcal{F}=\widetilde{M_H}$。
\end{proof}
% Locked French authority: ega2-1-fr.tex, 820505 B,
% SHA-256 84EDBE3E83530AF2959B441796337C9DC21EAFCA6A13114A26778760FBF437AC.
% Sequential source scope: lines 3291-3377, subsection II.2.7.9 through 2.7.11 /
% printed p.41.

\begin{corollary}[2.7.9]
\label{II.2.7.9-zh}
在 \egaref{II.2.7.5-zh}{2.7.5} 的假设下，设 $\mathcal{F}$ 是一个有限型
拟凝聚 $\mathcal{O}_X$-模。则存在整数 $n_0$，使得对任一
$n\geq n_0$，$\mathcal{F}(n)$ 同构于形如 $\mathcal{O}_X^k$ 的
$\mathcal{O}_X$-模的一个商（$k>0$ 依赖于 $n$），从而由它在 $X$ 上的
有限多个截面生成（\egaxref{0.5.1.1-zh}{0, 5.1.1}）。
\end{corollary}

\begin{proof}
由 \egaref{II.2.7.8-zh}{2.7.8}，可以假设
$\mathcal{F}=\widetilde{M}$，其中 $M$ 是有限多个形如 $S(m_i)$ 的
$S$-模的直和的商；再由 \egaref{II.2.5.4-zh}{2.5.4}，可只考虑
$M=S(m)$ 的情形。于是
$\mathcal{F}(n)=\widetilde{S(m+n)}
=\mathcal{O}_X(m+n)$（\egaref{II.2.5.15-zh}{2.5.15}）。所以只须证明

\oldpage[II]{41}
\begin{lemma}[2.7.9.1]
\label{II.2.7.9.1-zh}
在 \egaref{II.2.7.5-zh}{2.7.5} 的假设下，对任一 $n\geq0$，存在一个
整数 $k$（依赖于 $n$）和一个满同态
$\mathcal{O}_X^k\to\mathcal{O}_X(n)$。
\end{lemma}

事实上，只须由 \egaref{II.2.7.2-zh}{2.7.2} 证明：对适当的 $k$，存在
次数为 $0$ 的 (TN)-满射同态 $u$，从直积分次 $S$-模 $S^k$ 到
$S(n)$。但 $(S(n))_0=S_n$，且按假设，对任一 $h>0$ 都有
$S_h=S_1^h$，所以 $SS_n=S_n+S_{n+1}+\cdots$。由于 $S_n$ 是有限型
$S_0$-模（\egaref{II.2.1.5-zh}{2.1.5} 和
\egaref{II.2.1.6-zh}{2.1.6, (i)}），取这个模的一组生成元
$(a_i)_{1\leq i\leq k}$；同态 $u$ 把 $S^k$ 的典范基的第 $i$ 个元素
对应到 $a_i$（$1\leq i\leq k$）。\footnote{法文印本此处把对应方向写反；
由已声明的同态类型和生成元的角色，方向必须从典范基元素指向生成元。本译依
EGA2-CANON-DISP-0003 作可逆校正；法文外交转写保持原样。}
% EGA2-CANON-DISP-0003 / EGA2DEF68: guardian-confirmed correction with
% explicit reversible Chinese provenance; French diplomatic source unchanged.
此时 $\operatorname{Coker}u$ 等同于
$(S(n))_{-n}+\cdots+(S(n))_{-1}$，所以 $u$ 正是所求同态。
\end{proof}

\begin{corollary}[2.7.10]
\label{II.2.7.10-zh}
在 \egaref{II.2.7.5-zh}{2.7.5} 的假设下，设 $\mathcal{F}$ 是一个有限型
拟凝聚 $\mathcal{O}_X$-模。则存在整数 $n_0$，使得对任一
$n\geq n_0$，$\mathcal{F}$ 同构于形如 $(\mathcal{O}_X(-n))^k$ 的
$\mathcal{O}_X$-模的一个商（$k$ 依赖于 $n$）。
\end{corollary}

\begin{proposition}[2.7.11]
\label{II.2.7.11-zh}
假设 \egaref{II.2.7.5-zh}{2.7.5} 的各项假设成立，并设 $M$ 是一个分次
$S$-模。则：
\begin{enumerate}
  \item[\rm(i)] 典范同态
    $\widetilde{\alpha}:\widetilde{M}\to
    \widetilde{\Gamma_*(\widetilde{M})}$ 是同构。
  \item[\rm(ii)] 设 $\mathcal{G}$ 是 $\widetilde{M}$ 的一个拟凝聚
    $\mathcal{O}_X$-子模，并设 $N$ 是 $M$ 的分次 $S$-子模，即
    $\Gamma_*(\mathcal{G})$ 在 $\alpha$ 下的逆像。则
    $\widetilde{N}=\mathcal{G}$（由
    \egaref{II.2.5.4-zh}{2.5.4}，这里把 $\widetilde{N}$ 等同于
    $\widetilde{M}$ 的一个 $\mathcal{O}_X$-子模）。
\end{enumerate}
\end{proposition}

\begin{proof}
由 \egaref{II.2.7.5-zh}{2.7.5}，
$\beta:\widetilde{\Gamma_*(\widetilde{M})}
\to\widetilde{M}$ 是同构；再由
\egaref{II.2.6.5.1-zh}{2.6.5.1}，$\widetilde{\alpha}$ 是其逆同构，
从而得到 (i)。设 $P$ 是 $\Gamma_*(\widetilde{M})$ 的分次子模
$\alpha(M)$。因为 $M\mapsto\widetilde{M}$ 是正合函子
（\egaref{II.2.5.4-zh}{2.5.4}），$\widetilde{M}$ 在
$\widetilde{\alpha}$ 下的像等于 $\widetilde{P}$；所以由 (i)，
$\widetilde{P}=\widetilde{\Gamma_*(\widetilde{M})}$。置
$Q=\Gamma_*(\mathcal{G})\cap P$；由以上结论及
\egaref{II.2.5.4-zh}{2.5.4}，有
$\widetilde{Q}=\widetilde{\Gamma_*(\mathcal{G})}$，故由
\egaref{II.2.7.5-zh}{2.7.5}，$\beta$ 在 $\widetilde{Q}$ 上的限制是从这个
$\mathcal{O}_X$-模到 $\mathcal{G}$ 的一个\emph{同构}。但由 $N$ 的定义
及 \egaref{II.2.5.4-zh}{2.5.4}，同构 $\widetilde{\alpha}$ 在
$\widetilde{N}$ 上的限制是从 $\widetilde{N}$ 到 $\widetilde{Q}$ 的
同构；结论于是来自 \egaref{II.2.6.5.1-zh}{2.6.5.1}。
\end{proof}
% Locked French authority: ega2-1-fr.tex, 820505 B,
% SHA-256 84EDBE3E83530AF2959B441796337C9DC21EAFCA6A13114A26778760FBF437AC.
% Sequential source scope: lines 3379-3435, subsection II.2.8 through 2.8.1 /
% printed pp.41-42.

\subsection{函子性}
\label{subsection:II.2.8-zh}

\begin{env}[2.8.1]
\label{II.2.8.1-zh}
设 $S$、$S'$ 是两个按非负次数分次的环，
$\varphi:S'\to S$ 是分次环同态。以 $G(\varphi)$ 表示
$X=\operatorname{Proj}(S)$ 中开子集，即 $V_+(\varphi(S'_+))$ 的补集；
等价地，它是所有 $D_+(\varphi(f'))$ 的并，其中 $f'$ 遍历 $S'_+$ 的齐次
元素。从 $\operatorname{Spec}(S)$ 到 $\operatorname{Spec}(S')$ 的连续映射
${}^a\varphi$（\egaxref{I.1.2.1-zh}{I, 1.2.1}）在 $G(\varphi)$ 上的限制
因而是从 $G(\varphi)$ 到 $\operatorname{Proj}(S')$ 的连续映射；仍滥用语言
把它记作 ${}^a\varphi$。若 $f'\in S'_+$ 是齐次的，则
\[
\label{II.2.8.1.1-zh}
  {}^a\varphi^{-1}(D_+(f'))=D_+(\varphi(f'))
\tag{2.8.1.1}
\]
这是因为 ${}^a\varphi$ 把 $G(\varphi)$ 映入
$\operatorname{Proj}(S')$，并且由
\egaxref{I.1.2.2.2-zh}{I, 1.2.2.2}。另一方面，同态 $\varphi$ 典范地定义
了一个分次环同态 $S'_{f'}\to S_f$（仍沿用相同记号）；限制到次数为
$0$ 的元素便得到

\oldpage[II]{42}
一个同态 $S'_{(f')}\to S_{(f)}$，记作 $\varphi_{(f)}$。与之对应
（\egaxref{I.1.6.1-zh}{I, 1.6.1}）有仿射概形之间的态射
$({}^a\varphi_{(f)},\widetilde{\varphi}_{(f)}):
\operatorname{Spec}(S_{(f)})\to\operatorname{Spec}(S'_{(f')})$。若把
$\operatorname{Spec}(S_{(f)})$ 典范地等同于
$\operatorname{Proj}(S)$ 在 $D_+(f)$ 上所诱导的概形
（\egaref{II.2.3.6-zh}{2.3.6}），就定义了一个态射
$\Phi_f:D_+(f)\to D_+(f')$，而 ${}^a\varphi_{(f)}$ 等同于
${}^a\varphi$ 在 $D_+(f)$ 上的限制。另一方面，若 $g'$ 是 $S'_+$ 的另一个
齐次元素，且 $g=\varphi(g')$，则立即看出图表
\[
  \xymatrix{
    D_+(f)\ar[r]^{\Phi_f} & D_+(f')\\
    D_+(fg)\ar[u]\ar[r]_{\Phi_{fg}} & D_+(f'g')\ar[u]
  }
\]
是交换的，这是因为图表
\[
  \xymatrix{
    S'_{(f')}\ar[r]^{\varphi_{(f)}}\ar[d]_{\omega_{f'g',f'}}
    & S_{(f)}\ar[d]^{\omega_{fg,f}}\\
    S'_{(f'g')}\ar[r]_{\varphi_{(fg)}} & S_{(fg)}
  }
\]
交换。结合 $G(\varphi)$ 的定义与
\egaref{II.2.3.3.2-zh}{2.3.3.2}，由此可见：
\end{env}
% Locked French authority: ega2-1-fr.tex, 820505 B,
% SHA-256 84EDBE3E83530AF2959B441796337C9DC21EAFCA6A13114A26778760FBF437AC.
% Sequential source scope: lines 3437-3492, subsection II.2.8.2 through 2.8.3 /
% printed pp.42-43.

\begin{proposition}[2.8.2]
\label{II.2.8.2-zh}
给定一个分次环同态 $\varphi:S'\to S$，存在唯一一个从诱导预概形
$G(\varphi)$ 到 $\operatorname{Proj}(S')$ 的态射
$({}^a\varphi,\widetilde{\varphi})$；称它\emph{与 $\varphi$ 相伴}，并记作
$\operatorname{Proj}(\varphi)$。它具有如下性质：对任一齐次元素
$f'\in S'_+$，此态射在 $D_+(\varphi(f'))$ 上的限制，与由对应于
$\varphi$ 的同态 $S'_{(f')}\to S_{(\varphi(f'))}$ 所相伴的态射相同。
\end{proposition}

\begin{proof}
注意到，沿用上述记号，若 $f'\in S'_d$，则图表
\[
\label{II.2.8.2.1-zh}
  \xymatrix{
    S'_{(f')}\ar[r]^{\varphi_{(f)}}\ar[d]_{\sim}
    & S_{(f)}\ar[d]^{\sim}\\
    {S'}^{(d)}/(f'-1){S'}^{(d)}\ar[r]
    & S^{(d)}/(f-1)S^{(d)}
  }
\tag{2.8.2.1}
\]
是交换的（竖直箭头是 \egaref{II.2.2.5-zh}{2.2.5} 的同构）。
\end{proof}

\begin{corollary}[2.8.3]
\label{II.2.8.3-zh}
\begin{enumerate}
  \item[\rm(i)] 态射 $\operatorname{Proj}(\varphi)$ 是仿射的。
  \item[\rm(ii)] 若 $\operatorname{Ker}(\varphi)$ 是幂零的（特别地，若
    $\varphi$ 是单射），则态射 $\operatorname{Proj}(\varphi)$ 是支配的。
\end{enumerate}
\end{corollary}

\begin{proof}
(i) 立即来自 \egaref{II.2.8.2-zh}{2.8.2} 与
\egaref{II.2.8.1.1-zh}{2.8.1.1}。另一方面，若
$\operatorname{Ker}(\varphi)$ 是幂零的，则对 $S'_+$ 中任一齐次元素
$f'$，立即验证出 $\operatorname{Ker}(\varphi_f)$（其中
$f=\varphi(f')$）是幂零的，因而
$\operatorname{Ker}(\varphi_{(f)})$ 也是幂零的；结论于是来自
\egaref{II.2.8.2-zh}{2.8.2} 与
\egaxref{I.1.2.7-zh}{I, 1.2.7}。
\end{proof}

\oldpage[II]{43}
注意，一般而言，存在从 $\operatorname{Proj}(S)$ 到
$\operatorname{Proj}(S')$ 的非仿射态射，因而它们并不来自分次环同态
$S'\to S$。一个例子是当 $A$ 为域时的结构态射
$\operatorname{Proj}(S)\to\operatorname{Spec}(A)$（这里把
$\operatorname{Spec}(A)$ 与 $\operatorname{Proj}(A[T])$ 等同，参见
\egaref{II.3.1.7-zh}{3.1.7}）；这确实来自
\egaxref{I.2.3.2-zh}{I, 2.3.2}。
% Locked French authority: ega2-1-fr.tex, 820505 B,
% SHA-256 84EDBE3E83530AF2959B441796337C9DC21EAFCA6A13114A26778760FBF437AC.
% Sequential source scope: lines 3494-3597, subsection II.2.8.4 through 2.8.7 /
% printed pp.43-44.

\begin{env}[2.8.4]
\label{II.2.8.4-zh}
设 $S''$ 是第三个正次数分次环，$\varphi':S''\to S'$ 是一个分次环同态，
并令 $\varphi''=\varphi\circ\varphi'$。由
\egaref{II.2.8.1.1-zh}{2.8.1.1} 以及公式
${}^a\varphi''=({}^a\varphi')\circ({}^a\varphi)$，立即验证出
$G(\varphi'')\subset G(\varphi)$；若 $\Phi$、$\Phi'$、$\Phi''$ 分别是与
$\varphi$、$\varphi'$、$\varphi''$ 相伴的态射，则
$\Phi''=\Phi'\circ(\Phi|G(\varphi''))$。
\end{env}

\begin{env}[2.8.5]
\label{II.2.8.5-zh}
设 $S$（相应地 $S'$）是分次 $A$-代数（相应地分次 $A'$-代数），并设
$\psi:A'\to A$ 是一个使图表
\[
\label{II.2.8.5.1-zh}
  \xymatrix{
    A'\ar[r]^{\psi}\ar[d] & A\ar[d]\\
    S'\ar[r]_{\varphi} & S
  }
\tag{2.8.5.1}
\]
交换的环同态。此时可分别把 $G(\varphi)$ 与 $\operatorname{Proj}(S')$ 看作
$\operatorname{Spec}(A)$ 与 $\operatorname{Spec}(A')$ 上的概形；若
$\Phi$ 与 $\Psi$ 分别是与 $\varphi$ 和 $\psi$ 相伴的态射，则图表
\[
  \xymatrix{
    G(\varphi)\ar[r]^{\Phi}\ar[d] & \operatorname{Proj}(S')\ar[d]\\
    \operatorname{Spec}(A)\ar[r]_{\Psi} & \operatorname{Spec}(A')
  }
\]
交换。事实上，只需对 $\Phi$ 在 $D_+(f)$ 上的限制证明这一点，其中
$f=\varphi(f')$，而 $f'$ 是 $S'_+$ 中的齐次元素；这时结论来自图表
\[
  \xymatrix{
    A'\ar[r]^{\psi}\ar[d] & A\ar[d]\\
    S'_{(f')}\ar[r]_{\varphi_{(f)}} & S_{(f)}
  }
\]
的交换性。
\end{env}

\begin{env}[2.8.6]
\label{II.2.8.6-zh}
现设 $M$ 是一个分次 $S$-模，并考虑显然也分次的 $S'$-模
$M_{[\varphi]}$。设 $f'$ 是 $S'_+$ 中的齐次元素，并令
$f=\varphi(f')$；由 \egaxref{0.1.5.2-zh}{0, 1.5.2}，存在典范同构
$(M_{[\varphi]})_{f'}\xrightarrow{\sim}(M_f)_{[\varphi_f]}$，且此同构显然
保持次数，因而存在典范同构
$(M_{[\varphi]})_{(f')}\xrightarrow{\sim}
(M_{(f)})_{[\varphi_{(f)}]}$。与这个同构典范地对应着层的同构
$\widetilde{M_{[\varphi]}}|D_+(f')
\xrightarrow{\sim}(\Phi_f)_*(\widetilde{M}|D_+(f))$
（\egaref{II.2.5.2-zh}{2.5.2} 与
\egaxref{I.1.6.3-zh}{I, 1.6.3}）。此外，

\oldpage[II]{44}
若 $g'$ 是 $S'_+$ 中第二个齐次元素，且 $g=\varphi(g')$，则图表
\[
  \xymatrix{
    (M_{[\varphi]})_{(f')}\ar[r]^{\sim}\ar[d]
    & (M_{(f)})_{[\varphi_{(f)}]}\ar[d]\\
    (M_{[\varphi]})_{(f'g')}\ar[r]^{\sim}
    & (M_{(fg)})_{[\varphi_{(fg)}]}
  }
\]
交换。由此立即得出，同构
\[
  \widetilde{M_{[\varphi]}}|D_+(f'g')
  \xrightarrow{\sim}(\Phi_{fg})_*(\widetilde{M}|D_+(fg))
\]
是同构
$\widetilde{M_{[\varphi]}}|D_+(f')
\xrightarrow{\sim}(\Phi_f)_*(\widetilde{M}|D_+(f))$
在 $D_+(f'g')$ 上的限制。由于 $\Phi_f$ 是态射 $\Phi$ 在 $D_+(f)$ 上的
限制，结合 \egaref{II.2.8.1.1-zh}{2.8.1.1}，并记
$X'=\operatorname{Proj}(S')$，便得到：
\end{env}

\begin{proposition}[2.8.7]
\label{II.2.8.7-zh}
存在一个函子性的典范同构，把 $\mathcal{O}_{X'}$-模
$\widetilde{M_{[\varphi]}}$ 同构到 $\mathcal{O}_{X'}$-模
$\Phi_*(\widetilde{M}|G(\varphi))$ 上。
\end{proposition}

由此立即得到一个函子性的典范映射：它把从某个分次 $S'$-模到分次
$S$-模 $M$ 的 $\varphi$-态射 $M'\to M$ 全体所成集合，映到 $\Phi$-态射
$\widetilde{M'}\to\widetilde{M}|G(\varphi)$ 全体所成集合。
沿用 \egaref{II.2.8.4-zh}{2.8.4} 的记号，若 $M''$ 是一个分次
$S''$-模，则一个 $\varphi$-态射 $M'\to M$ 与一个 $\varphi'$-态射
$M''\to M'$ 的复合，典范地对应于
$\widetilde{M'}|G(\varphi')\to\widetilde{M}|G(\varphi'')$ 与
$\widetilde{M''}\to\widetilde{M'}|G(\varphi')$ 的复合。
% Locked French authority: ega2-1-fr.tex, 820505 B,
% SHA-256 84EDBE3E83530AF2959B441796337C9DC21EAFCA6A13114A26778760FBF437AC.
% Sequential source scope: lines 3599-3695, subsection II.2.8.8 through 2.8.9 /
% printed pp.44-45.

\begin{proposition}[2.8.8]
\label{II.2.8.8-zh}
在 \egaref{II.2.8.1-zh}{2.8.1} 的假设下，设 $M'$ 是一个分次
$S'$-模。存在一个函子性的典范同态 $\nu$，从
$(\mathcal{O}_X|G(\varphi))$-模 $\Phi^*(\widetilde{M'})$ 到
$(\mathcal{O}_X|G(\varphi))$-模
$\widetilde{M'\otimes_{S'}S}|G(\varphi)$。若理想 $S'_+$ 由
$S'_1$ 生成，则 $\nu$ 是同构。
\end{proposition}

\begin{proof}
事实上，对 $f'\in S'_d$（$d>0$），令 $f=\varphi(f')$，定义函子性的
典范 $S_{(f)}$-模同态
\[
\label{II.2.8.8.1-zh}
  \nu_f:M'_{(f')}\otimes_{S'_{(f')}}S_{(f)}
  \longrightarrow(M'\otimes_{S'}S)_{(f)}
\tag{2.8.8.1}
\]
如下：将同态
$M'_{(f')}\otimes_{S'_{(f')}}S_{(f)}\to
M'_{f'}\otimes_{S'_{f'}}S_f$ 与典范同构
$M'_{f'}\otimes_{S'_{f'}}S_f\xrightarrow{\sim}(M'\otimes_{S'}S)_f$
复合（\egaxref{0.1.5.4-zh}{0, 1.5.4}），并注意后一个同构保持次数。
立即验证出各 $\nu_f$ 与从 $D_+(f)$ 到 $D_+(fg)$ 的限制算子相容
（其中 $g'\in S'_+$ 是另一个元素，且 $g=\varphi(g')$）；结合
\egaxref{I.1.6.5-zh}{I, 1.6.5}，由此定义同态
\[
  \nu:\Phi^*(\widetilde{M'})
  \longrightarrow\widetilde{M'\otimes_{S'}S}|G(\varphi).
\]
为证明第二个断言，只需证明对每个 $f'\in S'_1$，$\nu_f$ 都是同构，
因为这时 $G(\varphi)$ 是各 $D_+(\varphi(f'))$ 的并。先定义一个
$\mathbf{Z}$-双线性映射
$M'_m\times S_n\to M'_{(f')}\otimes_{S'_{(f')}}S_{(f)}$，把
$(x',s)$ 对应到元素 $(x'/{f'}^m)\otimes(s/f^n)$（约定
$x'/{f'}^m$

\oldpage[II]{45}
当 $m<0$ 时指 ${f'}^{-m}x'/1$）。如同
\egaref{II.2.5.13-zh}{2.5.13} 的证明一样，可见这个映射诱导出模的
双同态
\[
  \eta_f:M'\otimes_{S'}S
  \longrightarrow M'_{(f')}\otimes_{S'_{(f')}}S_{(f)}.
\]
此外，若对 $r>0$ 有 $f^r\sum_i(x'_i\otimes s_i)=0$，也可写成
$\sum_i({f'}^rx'_i\otimes s_i)=0$；于是由
\egaxref{0.1.5.4-zh}{0, 1.5.4}，有
$\sum_i({f'}^rx'_i/{f'}^{m_i+r})\otimes(s_i/f^{n_i})=0$，也就是
$\eta_f(\sum_i x'_i\otimes s_i)=0$。这证明 $\eta_f$ 分解为
\[
  M'\otimes_{S'}S\longrightarrow(M'\otimes_{S'}S)_f
  \xrightarrow{\eta'_f}M'_{(f')}\otimes_{S'_{(f')}}S_{(f)}~;
\]
最后验证 $\eta'_f$ 与 $\nu_f$ 互为逆同构。

特别地，由 \egaref{II.2.1.2.1-zh}{2.1.2.1}，对每个
$n\in\mathbf{Z}$ 都有典范同态
\[
\label{II.2.8.8.2-zh}
  \Phi^*(\mathcal{O}_{X'}(n))\xrightarrow{\sim}
  \mathcal{O}_X(n)|G(\varphi)
\tag{2.8.8.2}
\]
。
\end{proof}

\begin{env}[2.8.9]
\label{II.2.8.9-zh}
设 $A$、$A'$ 是两个环，$\psi:A'\to A$ 是一个环同态，它定义态射
$\Psi:\operatorname{Spec}(A)\to\operatorname{Spec}(A')$。设 $S'$ 是一个
正次数分次 $A'$-代数，并令 $S=S'\otimes_{A'}A$；显然，它是一个以
$S'_n\otimes_{A'}A$ 为各次齐次部分的分次 $A$-代数。映射
$\varphi:s'\mapsto s'\otimes1$ 是一个使图表
\egaref{II.2.8.5.1-zh}{2.8.5.1} 交换的分次环同态。这里，$S_+$ 是由
$\varphi(S'_+)$ 生成的 $A$-模，故
$G(\varphi)=\operatorname{Proj}(S)=X$。于是，令
$X'=\operatorname{Proj}(S')$，得到交换图表
\[
\label{II.2.8.9.1-zh}
  \xymatrix{
    X\ar[r]^{\Phi}\ar[d]_p & X'\ar[d]\\
    Y\ar[r]_{\Psi} & Y'
  }
\tag{2.8.9.1}
\]

现设 $M'$ 是一个分次 $S'$-模，并令
$M=M'\otimes_{A'}A=M'\otimes_{S'}S$。在这些条件下：
\end{env}
% Locked French authority: ega2-1-fr.tex, 820505 B,
% SHA-256 84EDBE3E83530AF2959B441796337C9DC21EAFCA6A13114A26778760FBF437AC.
% Sequential source scope: lines 3697-3816, subsection II.2.8.10 through 2.8.13 /
% printed pp.45-46.

\begin{proposition}[2.8.10]
\label{II.2.8.10-zh}
图表 \egaref{II.2.8.9.1-zh}{2.8.9.1} 把概形 $X$ 等同于纤维积
$X'\times_{Y'}Y$；此外，典范同态
$\nu:\Phi^*(\widetilde{M'})\to\widetilde{M}$
（\egaref{II.2.8.8-zh}{2.8.8}）是同构。
\end{proposition}

\begin{proof}
为证明第一个断言，只需证明：对 $S'_+$ 中每个齐次元素 $f'$，令
$f=\varphi(f')$，则 $\Phi$ 与 $p$ 在 $D_+(f)$ 上的限制把这个概形等同于
纤维积 $D_+(f')\times_{Y'}Y$（\egaxref{I.3.2.6.2-zh}{I, 3.2.6.2}）。
换言之，只需（\egaxref{I.3.2.2-zh}{I, 3.2.2}）证明
$S_{(f)}$ 典范地等同于 $S'_{(f')}\otimes_{A'}A$；这立即来自保持次数的
典范同构 $S_f\xrightarrow{\sim}S'_{f'}\otimes_{A'}A$ 的存在
（\egaxref{0.1.5.4-zh}{0, 1.5.4}）。第二个断言来自如下事实：由上述结果，
$M'_{(f')}\otimes_{S'_{(f')}}S_{(f)}$ 等同于
$M'_{(f')}\otimes_{A'}A$，而后者又等同于 $M_{(f)}$，因为 $M_f$ 由一个
保持次数的同构典范地等同于 $M'_{f'}\otimes_{A'}A$。
\end{proof}

\begin{corollary}[2.8.11]
\label{II.2.8.11-zh}
对每个整数 $n\in\mathbf{Z}$，$\widetilde{M(n)}$ 等同于
$\Phi^*(\widetilde{M'(n)})=\widetilde{M'(n)}\otimes_{Y'}\mathcal{O}_Y$；
特别地，
$\mathcal{O}_X(n)=\Phi^*(\mathcal{O}_{X'}(n))
=\mathcal{O}_{X'}(n)\otimes_{Y'}\mathcal{O}_Y$。
\end{corollary}

\begin{proof}
这来自 \egaref{II.2.8.10-zh}{2.8.10} 与
\egaref{II.2.5.15-zh}{2.5.15}。
\end{proof}

\oldpage[II]{46}

\begin{env}[2.8.12]
\label{II.2.8.12-zh}
在 \egaref{II.2.8.9-zh}{2.8.9} 的假设下，对
$f'\in S'_d$（$d>0$）及 $f=\varphi(f')$，图表
\[
\label{II.2.8.12.1-zh}
  \xymatrix{
    M'_{(f')}\ar[r]^{\sim}\ar[d]
    & M'^{(d)}/(f'-1)M'^{(d)}\ar[d]\\
    M_{(f)}\ar[r]^{\sim}
    & M^{(d)}/(f-1)M^{(d)}
  }
\tag{2.8.12.1}
\]
（比较 \egaref{II.2.2.5-zh}{2.2.5}）交换。
\end{env}

\begin{env}[2.8.13]
\label{II.2.8.13-zh}
仍沿用 \egaref{II.2.8.9-zh}{2.8.9} 的记号和假设，设
$\mathcal{F}'$ 是一个 $\mathcal{O}_{X'}$-模；若令
$\mathcal{F}=\Phi^*(\mathcal{F}')$，则对每个 $n\in\mathbf{Z}$，由
\egaref{II.2.8.11-zh}{2.8.11} 与
\egaxref{0.4.3.3-zh}{0, 4.3.3} 有
$\mathcal{F}(n)=\Phi^*(\mathcal{F}'(n))$。于是由
\egaxref{0.3.7.1-zh}{0, 3.7.1}，得到典范同态
\[
  \Gamma(\rho):\Gamma(X',\mathcal{F}'(n))
  \longrightarrow\Gamma(X,\mathcal{F}(n))
\]
，它给出分次模的典范双同态
\[
  \Gamma_\bullet(\mathcal{F}')\longrightarrow
  \Gamma_\bullet(\mathcal{F}).
\]

假设理想 $S_+$ 由 $S_1$ 生成，并且
$\mathcal{F}'=\widetilde{M'}$，从而 $\mathcal{F}=\widetilde{M}$，其中
$M=M'\otimes_{A'}A$。若 $f'$ 是 $S'_+$ 中的齐次元素，且
$f=\varphi(f')$，则已知 $M_{(f)}=M'_{(f')}\otimes_{A'}A$，所以图表
\[
  \xymatrix{
    M'_0\ar[r]\ar[d]
    & M'_{(f')}\ar[d]
    & =\Gamma(D_+(f'),\widetilde{M'})\\
    M_0\ar[r]
    & M_{(f)}
    & =\Gamma(D_+(f),\widetilde{M})
  }
\]
交换。由这一点以及同态 $\alpha$ 的定义
（\egaref{II.2.6.2-zh}{2.6.2}），立即得出图表
\[
\label{II.2.8.13.1-zh}
  \xymatrix{
    M'\ar[r]^{\alpha_{M'}}\ar[d]
    & \Gamma_\bullet(\widetilde{M'})\ar[d]\\
    M\ar[r]_{\alpha_M}
    & \Gamma_\bullet(\widetilde{M})
  }
\tag{2.8.13.1}
\]
交换。同样，图表
\[
\label{II.2.8.13.2-zh}
  \xymatrix{
    \widetilde{\Gamma_\bullet(\mathcal{F}')}\ar[r]^{\beta_{\mathcal{F}'}}
      \ar[d]
    & \mathcal{F}'\ar[d]\\
    \widetilde{\Gamma_\bullet(\mathcal{F})}\ar[r]_{\beta_{\mathcal{F}}}
    & \mathcal{F}
  }
\tag{2.8.13.2}
\]
交换（其中右边的竖直箭头是典范 $\Phi$-态射
$\mathcal{F}'\to\Phi^*(\mathcal{F}')=\mathcal{F}$）。
\end{env}
% Locked French authority: ega2-1-fr.tex, 820505 B,
% SHA-256 84EDBE3E83530AF2959B441796337C9DC21EAFCA6A13114A26778760FBF437AC.
% Sequential source scope: lines 3818-3906, subsection II.2.8.14 through 2.8.15 /
% printed pp.47-48.

\oldpage[II]{47}

\begin{env}[2.8.14]
\label{II.2.8.14-zh}
仍保持 \egaref{II.2.8.9-zh}{2.8.9} 的假设与记号，设 $N'$ 是第二个
分次 $S'$-模，并令 $N=N'\otimes_{A'}A$。显然，典范双同态
$M'\to M$、$N'\to N$ 给出双同态
$M'\otimes_{S'}N'\to M\otimes_S N$（它相对于典范环同态
$S'\to S$），从而也给出次数为 $0$ 的 $S$-同态
\[
  (M'\otimes_{S'}N')\otimes_{A'}A\longrightarrow M\otimes_S N,
\]
结合 \egaref{II.2.8.10-zh}{2.8.10}，它对应于 $\mathcal{O}_X$-模同态
\[
\label{II.2.8.14.1-zh}
  \Phi^*(\widetilde{M'\otimes_{S'}N'})
  \longrightarrow\widetilde{M\otimes_S N}.
\tag{2.8.14.1}
\]

此外，立即验证出图表
\[
\label{II.2.8.14.2-zh}
  \xymatrix{
    \Phi^*(\widetilde{M'}\otimes_{\mathcal{O}_{X'}}\widetilde{N'})
      \ar[r]^{\sim}\ar[d]_{\Phi^*(\lambda)}
    & \widetilde{M}\otimes_{\mathcal{O}_X}\widetilde{N}
      \ar[d]^{\lambda}
    & =\Phi^*(\widetilde{M'})\otimes_{\mathcal{O}_X}
      \Phi^*(\widetilde{N'})\\
    \Phi^*(\widetilde{M'\otimes_{S'}N'})\ar[r]
    & \widetilde{M\otimes_S N}
  }
\tag{2.8.14.2}
\]
交换；第一行是 \egaxref{0.4.3.3-zh}{0, 4.3.3} 的典范同构。若理想
$S'_+$ 由 $S'_1$ 生成，则显然 $S_+$ 由 $S_1$ 生成，而且
\egaref{II.2.8.14.2-zh}{2.8.14.2} 的两条竖直箭头都是同构
（\egaref{II.2.5.13-zh}{2.5.13}）；因此
\egaref{II.2.8.14.1-zh}{2.8.14.1} 也是同构。

同样有典范双同态
$\operatorname{Hom}_{S'}(M',N')\to\operatorname{Hom}_S(M,N)$：把次数为
$k$ 的同态 $u'$ 对应到仍为次数 $k$ 的同态 $u'\otimes1$。由此又得到
次数为 $0$ 的 $S$-同态
\[
  (\operatorname{Hom}_{S'}(M',N'))\otimes_{A'}A
  \longrightarrow\operatorname{Hom}_S(M,N)
\]
，从而得到 $\mathcal{O}_X$-模同态
\[
\label{II.2.8.14.3-zh}
  \Phi^*(\widetilde{\operatorname{Hom}_{S'}(M',N')})
  \longrightarrow\widetilde{\operatorname{Hom}_S(M,N)}.
\tag{2.8.14.3}
\]

此外，图表
\[
  \xymatrix{
    \Phi^*(\widetilde{\operatorname{Hom}_{S'}(M',N')})\ar[r]
      \ar[d]_{\Phi^*(\mu)}
    & \widetilde{\operatorname{Hom}_S(M,N)}\ar[d]^{\mu}\\
    \Phi^*(\mathcal{H}om_{\mathcal{O}_{X'}}
      (\widetilde{M'},\widetilde{N'}))\ar[r]
    & \mathcal{H}om_{\mathcal{O}_X}(\widetilde{M},\widetilde{N})
  }
\]
交换（第二条水平箭头是 \egaxref{0.4.4.6-zh}{0, 4.4.6} 的典范同态）。
\end{env}

\oldpage[II]{48}

\begin{env}[2.8.15]
\label{II.2.8.15-zh}
记号和假设仍与 \egaref{II.2.8.1-zh}{2.8.1} 相同。由
\egaref{II.2.4.7-zh}{2.4.7}，若先把 $S_0$ 与 $S'_0$ 换成
$\mathbf{Z}$，并把 $\varphi_0$ 换成恒等映射，再分别把 $S$ 与 $S'$ 换成
$S^{(d)}$ 与 ${S'}^{(d)}$（$d>0$），并把 $\varphi$ 换成它在
$S^{(d)}$ 上的限制 $\varphi^{(d)}$，则在同构意义下并不改变态射
$\Phi$。
\end{env}
% Locked French authority: ega2-1-fr.tex, 820505 B,
% SHA-256 84EDBE3E83530AF2959B441796337C9DC21EAFCA6A13114A26778760FBF437AC.
% Sequential source scope: lines 3908-3974, subsection II.2.9 through 2.9.2 /
% printed p.48.

\subsection{概形 \texorpdfstring{$\operatorname{Proj}(S)$}{Proj(S)} 的闭子概形}
\label{subsection:II.2.9-zh}

\begin{env}[2.9.1]
\label{II.2.9.1-zh}
若 $\varphi:S\to S'$ 是一个分次环同态，则称 $\varphi$ 是 (TN)-
\emph{满射的}（相应地，(TN)-\emph{单射的}、(TN)-\emph{双射的}），
如果存在一个整数 $n$，使得当 $k\geq n$ 时，
$\varphi_k:S_k\to S'_k$ 是满射（相应地，单射、双射）。由
\egaref{II.2.8.15-zh}{2.8.15}，对 $\Phi$ 的研究可归结到 $\varphi$
为\emph{满射}（相应地，\emph{单射}、\emph{双射}）的情形。也把
$\varphi$ 是 (TN)-双射的说法表述为：它是一个 (TN)-\emph{同构}。
\end{env}

\begin{proposition}[2.9.2]
\label{II.2.9.2-zh}
设 $S$ 是一个正次数分次环，并设 $X=\operatorname{Proj}(S)$。
\begin{enumerate}
  \item[\rm(i)] 若 $\varphi:S\to S'$ 是一个分次环的 (TN)-满同态，则
    对应态射 $\Phi$（\egaref{II.2.8.1-zh}{2.8.1}）在整个
    $\operatorname{Proj}(S')$ 上都有定义，并且是从
    $\operatorname{Proj}(S')$ 到 $X$ 的闭浸入。若 $\mathfrak{J}$ 是
    $\varphi$ 的核，则与 $\Phi$ 相伴的 $X$ 的闭子概形由
    $\mathcal{O}_X$ 的拟凝聚理想层 $\widetilde{\mathfrak{J}}$ 定义。
  \item[\rm(ii)] 再假设理想 $S_+$ 由有限多个次数为 $1$ 的齐次元素
    生成。设 $X'$ 是 $X$ 的一个闭子概形，由 $\mathcal{O}_X$ 的一个
    拟凝聚理想层 $\mathcal{J}$ 定义。设 $\mathfrak{J}$ 是 $S$ 的分次
    理想，它是 $\Gamma_\bullet(\mathcal{J})$ 在典范同态
    $\alpha:S\to\Gamma_\bullet(\mathcal{O}_X)$
    （\egaref{II.2.6.2-zh}{2.6.2}）下的逆像，并令
    $S'=S/\mathfrak{J}$。那么，$X'$ 就是与分次环典范同态
    $S\to S'$ 所对应的闭浸入 $\operatorname{Proj}(S')\to X$ 相伴的
    子概形。
\end{enumerate}
\end{proposition}

\begin{proof}
\begin{enumerate}
  \item[\rm(i)] 可以假设 $\varphi$ 是满射
    （\egaref{II.2.9.1-zh}{2.9.1}）。按照假设，$\varphi(S_+)$ 生成
    $S'_+$，故 $G(\varphi)=\operatorname{Proj}(S')$。另一方面，第二个
    断言可在 $X$ 上局部验证；因此，设 $f$ 是 $S_+$ 中的一个齐次元素，
    并令 $f'=\varphi(f)$。由于 $\varphi$ 是分次环的满同态，立即验证出
    $\varphi_{(f')}:S_{(f)}\to S'_{(f')}$ 是满射，且其核为
    $\mathfrak{J}_{(f)}$；这便完成了 (i) 的证明
    （\egaxref{I.4.2.3-zh}{I, 4.2.3}）。
  \item[\rm(ii)] 由 (i)，只需验证从典范单射
    $j:\mathfrak{J}\to S$ 得到的同态
    $\widetilde{j}:\widetilde{\mathfrak{J}}\to\mathcal{O}_X$，把
    $\widetilde{\mathfrak{J}}$ 同构到 $\mathcal{J}$ 上；这来自
    \egaref{II.2.7.11-zh}{2.7.11}。
\end{enumerate}
\end{proof}

注意，$\mathfrak{J}$ 是满足
$\widetilde{j}(\widetilde{\mathfrak{J}'})=\mathcal{J}$ 的 $S$ 的分次
理想 $\mathfrak{J}'$ 中\emph{最大的}一个；事实上，直接返回定义
（\egaref{II.2.6.2-zh}{2.6.2}）即可验证，上述关系蕴含
$\alpha(\mathfrak{J}')\subset\Gamma_\bullet(\mathcal{J})$。
% Locked French authority: ega2-1-fr.tex, 820505 B,
% SHA-256 84EDBE3E83530AF2959B441796337C9DC21EAFCA6A13114A26778760FBF437AC.
% Sequential source scope: lines 3976-4043, subsection II.2.9.3 through the
% section II.3 heading / printed pp.48-49.

\begin{corollary}[2.9.3]
\label{II.2.9.3-zh}
假设 \egaref{II.2.9.2-zh}{2.9.2, (i)} 的假设成立，并且理想 $S_+$ 由
$S_1$ 生成；那么，对每个 $n\in\mathbf{Z}$，
$\Phi^*(\widetilde{S(n)})$ 典范地同构于 $\widetilde{S'(n)}$，从而对每个
$\mathcal{O}_X$-模 $\mathcal{F}$，$\Phi^*(\mathcal{F}(n))$ 典范地同构于
$\Phi^*(\mathcal{F})(n)$。
\end{corollary}

\begin{proof}
结合 \egaref{II.2.5.10.2-zh}{2.5.10.2}，这是
\egaref{II.2.8.8-zh}{2.8.8} 的一个特殊情形。
\end{proof}

\begin{corollary}[2.9.4]
\label{II.2.9.4-zh}
假设 \egaref{II.2.9.2-zh}{2.9.2, (ii)} 的假设成立。$X$ 的闭子预概形
$X'$ 为整概形的充分必要条件是分次理想 $\mathfrak{J}$ 在 $S$ 中是
素理想。
\end{corollary}

\oldpage[II]{49}

\begin{proof}
由于 $X'$ 同构于 $\operatorname{Proj}(S/\mathfrak{J})$，由
\egaref{II.2.4.4-zh}{2.4.4}，这个条件是充分的。为证明必要性，考虑正合列
$0\to\mathcal{J}\to\mathcal{O}_X\to\mathcal{O}_X/\mathcal{J}$；它给出
正合列
\[
  0\longrightarrow\Gamma_\bullet(\mathcal{J})
  \longrightarrow\Gamma_\bullet(\mathcal{O}_X)
  \longrightarrow\Gamma_\bullet(\mathcal{O}_X/\mathcal{J}).
\]
只需证明：若 $f\in S_m$、$g\in S_n$，且 $\alpha_{n+m}(fg)$ 在
$\Gamma_\bullet(\mathcal{O}_X/\mathcal{J})$ 中的像为零，则
$\alpha_m(f)$、$\alpha_n(g)$ 的像之一为零。按照定义，这两个像分别是
整概形 $X'$ 上的可逆
$(\mathcal{O}_X/\mathcal{J})$-模
$\mathcal{L}=(\mathcal{O}_X/\mathcal{J})(m)$ 与
$\mathcal{L}'=(\mathcal{O}_X/\mathcal{J})(n)$ 的截面。假设蕴含这两个
截面之积在 $\mathcal{L}\otimes\mathcal{L}'$ 中为零
（\egaref{II.2.9.3-zh}{2.9.3} 与
\egaref{II.2.5.14.1-zh}{2.5.14.1}），故由
\egaxref{I.7.4.4-zh}{I, 7.4.4}，其中一个截面为零。
\end{proof}

\begin{corollary}[2.9.5]
\label{II.2.9.5-zh}
设 $A$ 是一个环，$M$ 是一个 $A$-模，$S$ 是一个由次数为 $1$ 的齐次
元素全体 $S_1$ 生成的分次 $A$-代数，$u:M\to S_1$ 是一个满的
$A$-模同态，$\overline{u}:\mathbf{S}(M)\to S$ 是把 $M$ 的对称代数
$\mathbf{S}(M)$ 映到 $S$ 中并延拓 $u$ 的同态（即 $A$-代数同态）。
此时，与 $\overline{u}$ 对应的态射是从 $\operatorname{Proj}(S)$ 到
$\operatorname{Proj}(\mathbf{S}(M))$ 的闭浸入。
\end{corollary}

\begin{proof}
事实上，按照假设 $\overline{u}$ 是满射，只需应用
\egaref{II.2.9.2-zh}{2.9.2}。
\end{proof}

\section{分次代数层的齐次谱}
\label{section:II.3-zh}
% Locked French authority: ega2-1-fr.tex, 820505 B,
% SHA-256 84EDBE3E83530AF2959B441796337C9DC21EAFCA6A13114A26778760FBF437AC.
% Sequential source scope: lines 4045-4114, subsection II.3.1 through 3.1.1 /
% printed pp.49-50.

\subsection{拟凝聚分次 \texorpdfstring{$\mathcal{O}_Y$}{O-Y}-代数层的齐次谱}
\label{subsection:II.3.1-zh}
\label{II.3.1-zh}

\begin{env}[3.1.1]
\label{II.3.1.1-zh}
设 $Y$ 是一个预概形，$\mathcal{S}$ 是一个分次 $\mathcal{O}_Y$-代数层，
$\mathcal{M}$ 是一个分次 $\mathcal{S}$-模层。若 $\mathcal{S}$ 是
\emph{拟凝聚的}，则它的每个齐次分量 $\mathcal{S}_n$ 都是拟凝聚
$\mathcal{O}_Y$-模层，因为它是 $\mathcal{S}$ 在 $\mathcal{S}$ 的一个
自同态下的像
（\egaxref{I.1.3.8-zh}{I, 1.3.8} 与
\egaxref{I.1.3.9-zh}{I, 1.3.9}）。同样，若把 $\mathcal{M}$ 视为
$\mathcal{O}_Y$-模层时它是拟凝聚的，则它的各齐次分量
$\mathcal{M}_n$ 也拟凝聚，反之亦然。若 $d$ 是一个 $>0$ 的整数，记
$\mathcal{S}^{(d)}$ 为各 $\mathcal{O}_Y$-模层 $\mathcal{S}_{nd}$
（$n\in\mathbf{Z}$）的直和；若 $\mathcal{S}$ 拟凝聚，则这个直和也
拟凝聚（\egaxref{I.1.3.9-zh}{I, 1.3.9}）。对满足
$0\leq k\leq d-1$ 的每个整数 $k$，记 $\mathcal{M}^{(d,k)}$（当
$k=0$ 时也记作 $\mathcal{M}^{(d)}$）为各 $\mathcal{M}_{nd+k}$
（$n\in\mathbf{Z}$）的直和；它是一个分次 $\mathcal{S}^{(d)}$-模层，
并且当 $\mathcal{S}$ 与 $\mathcal{M}$ 拟凝聚时也是拟凝聚的
（\egaxref{I.9.6.1-zh}{I, 9.6.1}）。记 $\mathcal{M}(n)$ 为满足
$(\mathcal{M}(n))_k=\mathcal{M}_{n+k}$（对每个 $k\in\mathbf{Z}$）的
分次 $\mathcal{S}$-模层；若 $\mathcal{S}$ 与 $\mathcal{M}$ 拟凝聚，
则 $\mathcal{M}(n)$ 是拟凝聚分次 $\mathcal{S}$-模层
（\egaxref{I.9.6.1-zh}{I, 9.6.1}）。

称 $\mathcal{M}$ 是\emph{有限型}分次 $\mathcal{S}$-模层（相应地，
称它有一个\emph{有限表示}），如果对每个 $y\in Y$，存在 $y$ 的一个
开邻域 $U$ 和一些整数 $n_i$（相应地，一些整数 $m_i$ 与 $n_j$），
使得存在次数为 $0$ 的满同态
\[
  \bigoplus_{i=1}^r(\mathcal{S}(n_i)|U)\longrightarrow\mathcal{M}|U
\]
（相应地，$\mathcal{M}|U$ 同构于某个次数为 $0$ 的同态
\[
  \bigoplus_{i=1}^r(\mathcal{S}(m_i)|U)
  \longrightarrow\bigoplus_{j=1}^s(\mathcal{S}(n_j)|U)).
\]
的余核）。

设 $U$ 是 $Y$ 的一个仿射开集，其环为
$A=\Gamma(U,\mathcal{O}_Y)$。按照假设，分次
$(\mathcal{O}_Y|U)$-代数层 $\mathcal{S}|U$ 同构于
$\widetilde{S}$，其中 $S=\Gamma(U,\mathcal{S})$ 是一个分次
$A$-代数

\oldpage[II]{50}

（\egaxref{I.1.4.3-zh}{I, 1.4.3}）；令
$X_U=\operatorname{Proj}(\Gamma(U,\mathcal{S}))$。设 $U'\subset U$ 是
$Y$ 的第二个仿射开集，$A'$ 是它的环，$j:U'\to U$ 是对应于限制同态
$A\to A'$ 的典范单射。于是 $\mathcal{S}|U'=j^*(\mathcal{S}|U)$，从而
$S'=\Gamma(U',\mathcal{S})$ 典范地等同于 $S\otimes_A A'$
（\egaxref{I.1.6.4-zh}{I, 1.6.4}）。由
\egaref{II.2.8.10-zh}{2.8.10}，$X_{U'}$ 典范地等同于
$X_U\times_U U'$，因而也等同于 $f_U^{-1}(U')$；这里 $f_U$ 表示结构
态射 $X_U\to U$（\egaxref{I.4.4.1-zh}{I, 4.4.1}）。把由此定义的典范
同构 $f_U^{-1}(U')\xrightarrow{\sim}X_{U'}$ 记作 $\sigma_{U',U}$，并把
$\sigma_{U',U}^{-1}$ 与典范单射 $f_U^{-1}(U')\to X_U$ 的复合所得开
浸入 $X_{U'}\to X_U$ 记作 $\rho_{U',U}$。显然，若
$U''\subset U'$ 是 $Y$ 的第三个仿射开集，则
$\rho_{U'',U}=\rho_{U',U}\circ\rho_{U'',U'}$。
\end{env}
% Locked French authority: ega2-1-fr.tex, 820505 B,
% SHA-256 84EDBE3E83530AF2959B441796337C9DC21EAFCA6A13114A26778760FBF437AC.
% Sequential source scope: lines 4116-4179, subsection II.3.1.2 through 3.1.3 /
% printed p.50.

\begin{proposition}[3.1.2]
\label{II.3.1.2-zh}
设 $Y$ 是一个预概形。对每个正次数拟凝聚分次
$\mathcal{O}_Y$-代数层 $\mathcal{S}$，存在一个 $Y$ 上的预概形 $X$，
且在 $Y$-同构意义下唯一，具有如下性质：若 $f$ 是结构态射
$X\to Y$，则对 $Y$ 的每个仿射开集 $U$，存在一个从诱导预概形
$f^{-1}(U)$ 到 $X_U=\operatorname{Proj}(\Gamma(U,\mathcal{S}))$ 的
\emph{同构} $\eta_U$，使得当 $V$ 是包含于 $U$ 的 $Y$ 的第二个仿射
开集时，图表
\[
\label{II.3.1.2.1-zh}
  \xymatrix{
    f^{-1}(V)\ar[r]^{\eta_V}\ar[d]
    & X_V\ar[d]^{\rho_{V,U}}\\
    f^{-1}(U)\ar[r]_{\eta_U}
    & X_U
  }
\tag{3.1.2.1}
\]
交换。
\end{proposition}

\begin{proof}
对 $Y$ 的两个仿射开集 $U$、$V$，令 $X_{U,V}$ 为 $X_U$ 在
$f_U^{-1}(U\cap V)$ 上诱导的预概形；我们将定义一个 $Y$-同构
$\theta_{U,V}:X_{V,U}\xrightarrow{\sim}X_{U,V}$。为此，考虑一个仿射
开集 $W\subset U\cap V$：复合同构
\[
  f_U^{-1}(W)\xrightarrow{\sigma_{W,U}}X_W
  \xrightarrow{\sigma_{W,V}^{-1}}f_V^{-1}(W),
\]
得到一个同构 $\tau_W$。立即验证出，若 $W'\subset W$ 是一个仿射
开集，则 $\tau_{W'}$ 是 $\tau_W$ 在 $f_U^{-1}(W')$ 上的限制；故各
$\tau_W$ 确为一个 $Y$-同构 $\theta_{V,U}$ 的限制。此外，若
$U$、$V$、$W$ 是 $Y$ 的三个仿射开集，而 $\theta'_{U,V}$、
$\theta'_{V,W}$、$\theta'_{U,W}$ 分别是 $\theta_{U,V}$、
$\theta_{V,W}$、$\theta_{U,W}$ 在 $U\cap V\cap W$ 于 $X_V$、$X_W$、
$X_W$ 中的逆像上的限制，则由上述定义有
$\theta'_{U,V}\circ\theta'_{V,W}=\theta'_{U,W}$。因此，满足所述性质的
$X$ 的存在性来自 \egaxref{I.2.3.1-zh}{I, 2.3.1}；结合
\egaref{II.3.1.2.1-zh}{3.1.2.1}，它在 $Y$-同构意义下的唯一性是显然的。
\end{proof}

\begin{env}[3.1.3]
\label{II.3.1.3-zh}
称 \egaref{II.3.1.2-zh}{3.1.2} 中定义的预概形 $X$ 是拟凝聚分次
$\mathcal{O}_Y$-代数层 $\mathcal{S}$ 的\emph{齐次谱}，并记作
$\operatorname{Proj}(\mathcal{S})$。显然，
$\operatorname{Proj}(\mathcal{S})$ 在 $Y$ 上是\emph{分离的}
（\egaref{II.2.4.2-zh}{2.4.2} 与
\egaxref{I.5.5.5-zh}{I, 5.5.5}）；若 $\mathcal{S}$ 是\emph{有限型}
$\mathcal{O}_Y$-代数层（\egaxref{I.9.6.2-zh}{I, 9.6.2}），则
$\operatorname{Proj}(\mathcal{S})$ 在 $Y$ 上是\emph{有限型的}
（\egaref{II.2.7.1-zh}{2.7.1, (ii)} 与
\egaxref{I.6.3.1-zh}{I, 6.3.1}）。

若 $f$ 是结构态射 $X\to Y$，则显然对由 $Y$ 在 $Y$ 的开集 $U$ 上诱导的每个
预概形，$f^{-1}(U)$ 都等同于齐次谱
$\operatorname{Proj}(\mathcal{S}|U)$。
\end{env}
% Locked French authority: ega2-1-fr.tex, 820505 B,
% SHA-256 84EDBE3E83530AF2959B441796337C9DC21EAFCA6A13114A26778760FBF437AC.
% Sequential source scope: lines 4181-4284, subsection II.3.1.4 through the
% paragraph after 3.1.7 / printed pp.50-51.

\begin{proposition}[3.1.4]
\label{II.3.1.4-zh}
设 $f\in\Gamma(Y,\mathcal{S}_d)$，其中 $d>0$。在
$X=\operatorname{Proj}(\mathcal{S})$ 的底空间中存在一个开子集 $X_f$，
具有如下性质：对 $Y$ 的每个仿射开集 $U$，

\oldpage[II]{51}

在等同于 $X_U=\operatorname{Proj}(\Gamma(U,\mathcal{S}))$ 的
$\varphi^{-1}(U)$ 中有
$X_f\cap\varphi^{-1}(U)=D_+(f|U)$（这里 $\varphi$ 表示结构态射
$X\to Y$）。此外，在 $X_f$ 上诱导的预概形在 $Y$ 上是仿射的，并且
典范地同构于
$\operatorname{Spec}(\mathcal{S}^{(d)}/(f-1)\mathcal{S}^{(d)})$
（\egaref{II.1.3.1-zh}{1.3.1}）。
\end{proposition}

\begin{proof}
有 $f|U\in\Gamma(U,\mathcal{S}_d)=(\Gamma(U,\mathcal{S}))_d$。若
$U$、$U'$ 是 $Y$ 的两个仿射开集，且 $U'\subset U$，则 $f|U'$ 是
$f|U$ 在限制同态
\[
  \Gamma(U,\mathcal{S})\longrightarrow\Gamma(U',\mathcal{S})
\]
下的像。因此，沿用 \egaref{II.3.1.1-zh}{3.1.1} 的记号，
$D_+(f|U')$ 等于 $X_{U'}$ 在逆像
$\rho_{U',U}^{-1}(D_+(f|U))$ 上诱导的预概形
（\egaref{II.2.8.1-zh}{2.8.1}）；这就证明了第一个断言。此外，
$X_U$ 在 $D_+(f|U)$ 上诱导的预概形典范地等同于
$\operatorname{Spec}((\Gamma(U,\mathcal{S}))_{(f|U)})$，而这些等同与
限制同态相容（\egaref{II.2.8.1-zh}{2.8.1}）；第二个断言于是来自
\egaref{II.2.2.5-zh}{2.2.5} 以及图表
\egaref{II.2.8.2.1-zh}{2.8.2.1} 的交换性。
\end{proof}

仍称 $X_f$（把它看作底空间 $X$ 中的开集）为 $f$ 
\emph{不消失}的点 $x\in X$ 全体所成集合。

\begin{corollary}[3.1.5]
\label{II.3.1.5-zh}
若 $f\in\Gamma(Y,\mathcal{S}_d)$、$g\in\Gamma(Y,\mathcal{S}_e)$，则
\[
\label{II.3.1.5.1-zh}
  X_{fg}=X_f\cap X_g.
\tag{3.1.5.1}
\]
\end{corollary}

\begin{proof}
事实上，只需把等式两边分别与集合 $\varphi^{-1}(U)$ 相交，其中 $U$ 是
$Y$ 中的一个仿射开集，然后应用公式
\egaref{II.2.3.3.2-zh}{2.3.3.2}。
\end{proof}

\begin{corollary}[3.1.6]
\label{II.3.1.6-zh}
设 $(f_\alpha)$ 是 $Y$ 上 $\mathcal{S}$ 的一族截面，且
$f_\alpha\in\Gamma(Y,\mathcal{S}_{d_\alpha})$。若由这族截面生成的
$\mathcal{S}$ 的理想层（\egaxref{0.5.1.1-zh}{0, 5.1.1}）从某一阶起
包含所有 $\mathcal{S}_n$，则底空间 $X$ 是各 $X_{f_\alpha}$ 的并。
\end{corollary}

\begin{proof}
事实上，对 $Y$ 的每个仿射开集 $U$，$\varphi^{-1}(U)$ 是各
$X_{f_\alpha}\cap\varphi^{-1}(U)$ 的并
（\egaref{II.2.3.14-zh}{2.3.14}）。
\end{proof}

\begin{corollary}[3.1.7]
\label{II.3.1.7-zh}
设 $\mathcal{A}$ 是一个拟凝聚 $\mathcal{O}_Y$-代数层；令
\[
  \mathcal{S}=\mathcal{A}[T]
  =\mathcal{A}\otimes_{\mathbf{Z}}\mathbf{Z}[T]
\]
，其中 $T$ 是一个不定元（把 $\mathbf{Z}$ 与 $\mathbf{Z}[T]$ 看作
$Y$ 上的常值层）。那么，$X=\operatorname{Proj}(\mathcal{S})$ 典范地
等同于 $\operatorname{Spec}(\mathcal{A})$。特别地，
$\operatorname{Proj}(\mathcal{O}_Y[T])$ 等同于 $Y$。
\end{corollary}

\begin{proof}
对 $Y$ 上处处等于 $T$ 的唯一截面
$f\in\Gamma(Y,\mathcal{S})$ 应用 \egaref{II.3.1.6-zh}{3.1.6}，可见
$X_f=X$。此外，这里 $d=1$，并且
$\mathcal{S}^{(1)}/(f-1)\mathcal{S}^{(1)}
=\mathcal{S}/(f-1)\mathcal{S}$ 典范地同构于 $\mathcal{A}$，故由
\egaref{II.1.2.2-zh}{1.2.2} 得到本推论。
\end{proof}

设 $g\in\Gamma(Y,\mathcal{O}_Y)$；若取
$\mathcal{S}=\mathcal{O}_Y[T]$，则 $g\in\Gamma(Y,\mathcal{S}_0)$；令
\[
  h=gT\in\Gamma(Y,\mathcal{S}_1).
\]
若 $X=\operatorname{Proj}(\mathcal{S})$，则
\egaref{II.3.1.7-zh}{3.1.7} 中定义的典范等同把 $X_h$ 等同于 $Y$ 的
开集 $Y_g$（按 \egaxref{0.5.5.2-zh}{0, 5.5.2} 的意义）。事实上，
可只考虑 $Y=\operatorname{Spec}(A)$ 为仿射的情形；结合
\egaref{II.2.2.5-zh}{2.2.5}，一切便归结为分式环 $A_g$ 典范地等同于
$A[T]/(gT-1)A[T]$（\egaxref{0.1.2.3-zh}{0, 1.2.3}）。
% Locked French authority: ega2-1-fr.tex, 820505 B,
% SHA-256 84EDBE3E83530AF2959B441796337C9DC21EAFCA6A13114A26778760FBF437AC.
% Sequential source scope: lines 4286-4388, subsection II.3.1.8 through 3.1.11 /
% printed pp.51-52.

\begin{proposition}[3.1.8]
\label{II.3.1.8-zh}
设 $\mathcal{S}$ 是一个正次数拟凝聚分次 $\mathcal{O}_Y$-代数层。
\begin{enumerate}
  \item[\rm(i)] 对每个 $d>0$，存在一个从
    $\operatorname{Proj}(\mathcal{S})$ 到
    $\operatorname{Proj}(\mathcal{S}^{(d)})$ 的典范 $Y$-同构。

\oldpage[II]{52}

  \item[\rm(ii)] 设 $\mathcal{S}'$ 是由 $\mathcal{O}_Y$ 和各
    $\mathcal{S}_n$（$n\geq0$）直和而成的分次 $\mathcal{O}_Y$-代数层；
    则 $\operatorname{Proj}(\mathcal{S}')$ 与
    $\operatorname{Proj}(\mathcal{S})$ 典范地 $Y$-同构。
  \item[\rm(iii)] 设 $\mathcal{L}$ 是一个可逆 $\mathcal{O}_Y$-模层
    （\egaxref{0.5.4.1-zh}{0, 5.4.1}），并设
    $\mathcal{S}_{(\mathcal{L})}$ 是由各
    $\mathcal{S}_d\otimes\mathcal{L}^{\otimes d}$（$d\geq0$）直和而成的
    分次 $\mathcal{O}_Y$-代数层；则
    $\operatorname{Proj}(\mathcal{S})$ 与
    $\operatorname{Proj}(\mathcal{S}_{(\mathcal{L})})$ 典范地
    $Y$-同构。
\end{enumerate}
\end{proposition}

\begin{proof}
在三种情形中，都只需在 $Y$ 上局部定义同构；验证它与从开集到更小开集
的限制运算相容是平凡的。因此可以假设 $Y$ 是仿射的；这时 (i) 来自
\egaref{II.2.4.7-zh}{2.4.7, (i)}，(ii) 来自
\egaref{II.2.4.8-zh}{2.4.8}。对于 (iii)，若再假设 $\mathcal{L}$ 同构于
$\mathcal{O}_Y$（这是允许的，因为问题在 $Y$ 上是局部的），则
$\operatorname{Proj}(\mathcal{S})$ 与
$\operatorname{Proj}(\mathcal{S}_{(\mathcal{L})})$ 同构是显然的。为定义
一个\emph{典范}同构，设 $Y=\operatorname{Spec}(A)$，
$\mathcal{S}=\widetilde{S}$，其中 $S$ 是一个分次 $A$-代数；再设 $c$ 是
自由 $A$-模 $L$ 的一个生成元，且 $\mathcal{L}=\widetilde{L}$。于是对
每个 $n>0$，$x_n\mapsto x_n\otimes c^{\otimes n}$ 是从 $S_n$ 到
$S_n\otimes L^{\otimes n}$ 的 $A$-同构，并且这些 $A$-同构定义分次
代数的 $A$-同构
\[
  \varphi_c:S\longrightarrow S_{(L)}
  =\bigoplus_{n\geq0}S_n\otimes L^{\otimes n}.
\]
现设 $f\in S_+$ 是次数为 $d$ 的齐次元素；对每个 $x\in S_{nd}$，有
\[
  \frac{x\otimes c^{\otimes nd}}{(f\otimes c^{\otimes d})^n}
  =\frac{x\otimes(\varepsilon c)^{\otimes nd}}
  {(f\otimes(\varepsilon c)^{\otimes d})^n}
\]
，其中 $\varepsilon\in A$ 是任一可逆元素。这表明，由 $\varphi_c$ 得到的
同构 $S_{(f)}\to(S_{(L)})_{(f\otimes c^d)}$ \emph{不依赖于}所取的
$L$ 的生成元 $c$，并完成证明。
\end{proof}

\begin{env}[3.1.9]
\label{II.3.1.9-zh}
回忆（\egaxref{0.4.1.3-zh}{0, 4.1.3} 与
\egaxref{I.1.3.14-zh}{I, 1.3.14}）：拟凝聚分次
$\mathcal{O}_Y$-代数层 $\mathcal{S}$ \emph{由 $\mathcal{O}_Y$-模层
$\mathcal{S}_1$ 生成}的充分必要条件是，存在由仿射开集组成的 $Y$ 的
一个覆盖 $(U_\alpha)$，使得 $\Gamma(U_\alpha,\mathcal{S})$ 作为
$\Gamma(U_\alpha,\mathcal{S}_0)$ 上的分次代数，由其次数为 $1$ 的齐次
元素全体 $\Gamma(U_\alpha,\mathcal{S}_1)$ 生成。于是，对 $Y$ 的每个
开集 $V$，$\mathcal{S}|V$ 由 $(\mathcal{O}_Y|V)$-模层
$\mathcal{S}_1|V$ 生成。
\end{env}

\begin{proposition}[3.1.10]
\label{II.3.1.10-zh}
假设 $Y$ 有一个有限仿射开覆盖 $(U_i)$，使每个分次代数
$\Gamma(U_i,\mathcal{S})$ 都是 $\Gamma(U_i,\mathcal{O}_Y)$ 上的有限型
代数。那么存在 $d>0$，使 $\mathcal{S}^{(d)}$ 由 $\mathcal{S}_d$ 生成，
且 $\mathcal{S}_d$ 是有限型 $\mathcal{O}_Y$-模层。
\end{proposition}

\begin{proof}
事实上，由 \egaref{II.2.1.6-zh}{2.1.6, (v)}，对每个 $i$，存在一个
整数 $m_i$，使得对每个 $n>0$ 都有
$\Gamma(U_i,\mathcal{S}_{nm_i})=(\Gamma(U_i,\mathcal{S}_{m_i}))^n$。
结合 \egaref{II.2.1.6-zh}{2.1.6, (i)}，只需取 $d$ 为各 $m_i$ 的一个
公倍数。
\end{proof}

\begin{corollary}[3.1.11]
\label{II.3.1.11-zh}
在 \egaref{II.3.1.10-zh}{3.1.10} 的假设下，
$\operatorname{Proj}(\mathcal{S})$ 与某个齐次谱
$\operatorname{Proj}(\mathcal{S}')$ 在 $Y$ 上同构，其中
$\mathcal{S}'$ 是由 $\mathcal{S}'_1$ 生成的分次 $\mathcal{O}_Y$-代数层，
并假设 $\mathcal{S}'_1$ 是有限型 $\mathcal{O}_Y$-模层。
\end{corollary}

\begin{proof}
只需取 $\mathcal{S}'=\mathcal{S}^{(d)}$，其中 $d$ 由
\egaref{II.3.1.10-zh}{3.1.10} 的性质确定，再应用
\egaref{II.3.1.8-zh}{3.1.8, (i)}。
\end{proof}
% Locked French authority: ega2-1-fr.tex, 820505 B,
% SHA-256 84EDBE3E83530AF2959B441796337C9DC21EAFCA6A13114A26778760FBF437AC.
% Sequential source scope: lines 4390-4502, subsection II.3.1.12 through 3.1.14 /
% printed pp.52-54.

\begin{env}[3.1.12]
\label{II.3.1.12-zh}
若 $\mathcal{S}$ 是一个正次数拟凝聚分次 $\mathcal{O}_Y$-代数层，则由
\egaxref{I.5.1.1-zh}{I, 5.1.1}，它的\emph{幂零根}
$\mathcal{N}$ 是拟凝聚 $\mathcal{O}_Y$-模层；称
$\mathcal{N}_+=\mathcal{N}\cap\mathcal{S}_+$ 是
\emph{$\mathcal{S}_+$ 的幂零根}。这是一个拟凝聚分次
$\mathcal{S}_0$-模层，因为立即可归结到 $Y$ 为仿射的情形，这时结论
来自 \egaref{II.2.1.10-zh}{2.1.10}。于是，对每个 $y\in Y$，
$(\mathcal{N}_+)_y$ 是
$(\mathcal{S}_+)_y=(\mathcal{S}_y)_+$ 的幂零根
（\egaxref{I.5.1.1-zh}{I, 5.1.1}）。若
$\mathcal{N}_+=0$，则称分次 $\mathcal{O}_Y$-代数层 $\mathcal{S}$
\emph{本质约化}；这等价于说

\oldpage[II]{53}

对每个 $y\in Y$，$\mathcal{S}_y$ 都是本质约化分次
$\mathcal{O}_y$-代数。对任一分次 $\mathcal{O}_Y$-代数层
$\mathcal{S}$，$\mathcal{S}/\mathcal{N}_+$ 都本质约化。

若对每个 $y\in Y$，$\mathcal{S}_y$ 都是整环，并且还对每个
$y\in Y$ 有 $(\mathcal{S}_y)_+=(\mathcal{S}_+)_y\neq0$，则称
$\mathcal{S}$ 是\emph{整的}。
\end{env}

\begin{proposition}[3.1.13]
\label{II.3.1.13-zh}
设 $\mathcal{S}$ 是一个正次数分次 $\mathcal{O}_Y$-代数层。若
$X=\operatorname{Proj}(\mathcal{S})$，则 $Y$-概形
$X_{\mathrm{red}}$ 典范地同构于
$\operatorname{Proj}(\mathcal{S}/\mathcal{N}_+)$；特别地，若
$\mathcal{S}$ 本质约化，则 $X$ 是约化的。
\end{proposition}

\begin{proof}
$X'=\operatorname{Proj}(\mathcal{S}/\mathcal{N}_+)$ 是约化的，这立即
来自 \egaref{II.2.4.4-zh}{2.4.4, (i)}，因为该性质是局部的。此外，
对每个仿射开集 $U\subset Y$，${\varphi'}^{-1}(U)$ 等于
$(\varphi^{-1}(U))_{\mathrm{red}}$（这里 $\varphi$ 与
$\varphi'$ 分别表示结构态射 $X\to Y$、$X'\to Y$）。立即验证出典范
$U$-态射 ${\varphi'}^{-1}(U)\to\varphi^{-1}(U)$ 与限制运算相容，因而
定义一个闭浸入 $X'\to X$，它是底空间之间的同胚；结论于是来自
\egaxref{I.5.1.2-zh}{I, 5.1.2}。
\end{proof}

\begin{proposition}[3.1.14]
\label{II.3.1.14-zh}
设 $Y$ 是一个整预概形，$\mathcal{S}$ 是一个满足
$\mathcal{S}_0=\mathcal{O}_Y$ 的拟凝聚分次 $\mathcal{O}_Y$-代数层。
\begin{enumerate}
  \item[\rm(i)] 若 $\mathcal{S}$ 是整的
    （\egaref{II.3.1.12-zh}{3.1.12}），则
    $X=\operatorname{Proj}(\mathcal{S})$ 是整的，且结构态射
    $\varphi:X\to Y$ 是支配的。
  \item[\rm(ii)] 再假设 $\mathcal{S}$ 本质约化。反过来，若 $X$ 是整的
    且 $\varphi$ 是支配的，则 $\mathcal{S}$ 是整的。
\end{enumerate}
\end{proposition}

\begin{proof}
\begin{enumerate}
  \item[\rm(i)] 若 $(U_\alpha)$ 是由非空仿射开集组成的 $Y$ 的一个基，
    则只需在把 $Y$ 换成某个 $U_\alpha$、把 $\mathcal{S}$ 换成
    $\mathcal{S}|U_\alpha$ 后证明该命题。事实上，这一方面蕴含各底空间
    $\varphi^{-1}(U_\alpha)$ 是 $X$ 的不可约开集（因而非空），并且对
    每对指标都有
    $\varphi^{-1}(U_\alpha)\cap\varphi^{-1}(U_\beta)\neq\varnothing$
    （因为 $U_\alpha\cap U_\beta$ 包含某个 $U_\gamma$）；因此 $X$ 不可约
    （\egaxref{0.2.1.4-zh}{0, 2.1.4}）。另一方面，由于约化性是局部
    性质，$X$ 将是约化的，故 $X$ 确为整概形，且 $\varphi(X)$ 在 $Y$
    中稠密。

    因此，设 $Y=\operatorname{Spec}(A)$，其中 $A$ 是整环
    （\egaxref{I.5.1.4-zh}{I, 5.1.4}），并设
    $\mathcal{S}=\widetilde{S}$，其中 $S$ 是一个分次 $A$-代数。假设是：
    对每个 $y\in Y$，$\widetilde{S}_y=S_y$ 是一个满足
    $(S_y)_+\neq0$ 的分次整环。只需证明 $S$ 是\emph{整环}；因为这时
    $S_+\neq0$，并可应用 \egaref{II.2.4.4-zh}{2.4.4, (ii)}。事实上，
    设 $f$、$g$ 是 $S$ 中两个 $\neq0$ 的元素，并假设 $fg=0$。于是对
    每个 $y\in Y$，在 $S_y$ 中有 $(f/1)(g/1)=0$，故按假设
    $f/1=0$ 或 $g/1=0$。例如，假设在 $S_y$ 中 $f/1=0$；这意味着存在
    $a\in A$，使 $a\notin\mathfrak{j}_y$ 且 $af=0$。于是对每个
    $z\in Y$，在整环 $S_z$ 中有 $(a/1)(f/1)=0$；由于 $a/1\neq0$
    （因为 $A$ 是整环），所以 $f/1=0$，这蕴含 $f=0$。
  \item[\rm(ii)] 问题在 $Y$ 上是局部的，故仍可假设
    $Y=\operatorname{Spec}(A)$，其中 $A$ 是整环，并且
    $\mathcal{S}=\widetilde{S}$。按照假设，对每个 $y\in Y$，
    $(S_y)_+$ 不含 $\neq0$ 的幂零元素；由假设，对
    $(S_0)_y=A_y$ 也如此。因此，对每个 $y\in Y$，$S_y$ 都是约化环，
    从而先得到 $S$ 本身是约化的
    （\egaxref{I.5.1.1-zh}{I, 5.1.1}）。另一方面，$X$ 为整概形的假设
    蕴含 $S$ 本质整（\egaref{II.2.4.4-zh}{2.4.4, (ii)}）；于是只需
    证明 $S_+$ 在 $A=S_0$ 中的零化子 $\mathfrak{J}$ 等于

\oldpage[II]{54}

    $0$（\egaref{II.2.1.11-zh}{2.1.11}）。否则，对
    $\mathfrak{J}$ 中某个 $h\neq0$ 将有 $(S_h)_+=0$，从而由
    \egaref{II.3.1.1-zh}{3.1.1}，
    $\varphi^{-1}(D(h))=\varnothing$；这与 $\varphi(X)$ 在 $Y$ 中稠密的
    假设矛盾（因为 $h$ 不是幂零元，所以 $D(h)\neq\varnothing$）。
\end{enumerate}
\end{proof}
% Locked French authority: ega2-1-fr.tex, 820505 B,
% SHA-256 84EDBE3E83530AF2959B441796337C9DC21EAFCA6A13114A26778760FBF437AC.
% Sequential source scope: lines 4504-4617, subsection II.3.2 through the
% paragraph after 3.2.4 / printed pp.54-55.

\subsection{与分次 \texorpdfstring{$\mathcal{S}$}{S}-模层相伴的
\texorpdfstring{$\operatorname{Proj}(\mathcal{S})$}{Proj(S)} 上的层}
\label{subsection:II.3.2-zh}

\begin{env}[3.2.1]
\label{II.3.2.1-zh}
设 $Y$ 是一个预概形，$\mathcal{S}$ 是一个正次数拟凝聚分次
$\mathcal{O}_Y$-代数层，$\mathcal{M}$ 是一个拟凝聚分次
$\mathcal{S}$-模层（相对于 $(Y,\mathcal{O}_Y)$ 或相对于环化空间
$(Y,\mathcal{S})$，二者等价，见 \egaxref{I.9.6.1-zh}{I, 9.6.1}）。
沿用 \egaref{II.3.1.1-zh}{3.1.1} 的记号，以
$\widetilde{\mathcal{M}}_U$ 表示拟凝聚 $\mathcal{O}_{X_U}$-模层
$\widetilde{\Gamma(U,\mathcal{M})}$。若 $U'\subset U$，则
$\Gamma(U',\mathcal{M})$ 典范地等同于
$\Gamma(U,\mathcal{M})\otimes_A A'$（\egaxref{I.1.6.4-zh}{I, 1.6.4}）；
因此
$\widetilde{\mathcal{M}}_{U'}
=\rho_{U',U}^*(\widetilde{\mathcal{M}}_U)$
（\egaref{II.2.8.11-zh}{2.8.11}）。
\end{env}

\begin{proposition}[3.2.2]
\label{II.3.2.2-zh}
在 $\operatorname{Proj}(\mathcal{S})=X$ 上存在唯一一个拟凝聚
$\mathcal{O}_X$-模层 $\widetilde{\mathcal{M}}$，使得对 $Y$ 的每个仿射
开集 $U$ 都有
\[
  \eta_U^*(\widetilde{\Gamma(U,\mathcal{M})})
  =\widetilde{\mathcal{M}}|f^{-1}(U)
\]
（其中 $\eta_U$ 表示同构
$f^{-1}(U)\xrightarrow{\sim}\operatorname{Proj}(\Gamma(U,\mathcal{S}))$，
而 $f$ 是结构态射 $X\to Y$）。
\end{proposition}

\begin{proof}
由于 $\rho_{U',U}$ 等同于单射态射
$f^{-1}(U')\to f^{-1}(U)$（\egaref{II.3.1.2.1-zh}{3.1.2.1}），该命题
立即由关系
$\widetilde{\mathcal{M}}_{U'}=\rho_{U',U}^*(\widetilde{\mathcal{M}}_U)$
以及层的粘合原理（\egaxref{0.3.3.1-zh}{0, 3.3.1}）得出。
\end{proof}

称 $\widetilde{\mathcal{M}}$ 是与拟凝聚分次 $\mathcal{S}$-模层
$\mathcal{M}$ 相伴的 $\mathcal{O}_X$-模层。

\begin{proposition}[3.2.3]
\label{II.3.2.3-zh}
设 $\mathcal{M}$ 是一个拟凝聚分次 $\mathcal{S}$-模层，并设
$f\in\Gamma(Y,\mathcal{S}_d)$（$d>0$）。若 $\xi_f$ 是从 $X_f$ 到
$Y$-预概形
$Z_f=\operatorname{Spec}(\mathcal{S}^{(d)}/(f-1)\mathcal{S}^{(d)})$ 的
典范同构（\egaref{II.3.1.4-zh}{3.1.4}），则
$(\xi_f)_*(\widetilde{\mathcal{M}}|X_f)$ 是 $\mathcal{O}_{Z_f}$-模层
$\widetilde{\mathcal{M}^{(d)}/(f-1)\mathcal{M}^{(d)}}$
（\egaref{II.1.4.3-zh}{1.4.3}）。
\end{proposition}

\begin{proof}
问题在 $Y$ 上是局部的；结合图表
\egaref{II.2.8.12.1-zh}{2.8.12.1} 的交换性，立即归结到
\egaref{II.2.2.5-zh}{2.2.5}。
\end{proof}

\begin{proposition}[3.2.4]
\label{II.3.2.4-zh}
$\mathcal{O}_X$-模层 $\widetilde{\mathcal{M}}$ 关于 $\mathcal{M}$ 是一个
从拟凝聚分次 $\mathcal{S}$-模层范畴到拟凝聚 $\mathcal{O}_X$-模层范畴
的正合加性协变函子，并且与归纳极限及直和交换。
\end{proposition}

\begin{proof}
问题在 $Y$ 上是局部的，因而归结到
\egaxref{I.1.3.11-zh}{I, 1.3.11}、
\egaxref{I.1.3.9-zh}{I, 1.3.9} 以及
\egaref{II.2.5.4-zh}{2.5.4}。
\end{proof}

特别地，若 $\mathcal{N}$ 是 $\mathcal{M}$ 的一个拟凝聚分次
$\mathcal{S}$-子模层，则 $\widetilde{\mathcal{N}}$ 典范地等同于
$\widetilde{\mathcal{M}}$ 的一个拟凝聚 $\mathcal{O}_X$-子模层；更特别
地，对 $\mathcal{S}$ 的任一拟凝聚分次理想层 $\mathcal{J}$，
$\widetilde{\mathcal{J}}$ 是 $\mathcal{O}_X$ 的拟凝聚理想层。

若 $\mathcal{M}$ 是一个拟凝聚分次 $\mathcal{S}$-模层，而
$\mathcal{I}$ 是 $\mathcal{O}_Y$ 的一个拟凝聚理想层，则
$\mathcal{I}\mathcal{M}$ 是 $\mathcal{M}$ 的一个拟凝聚分次
$\mathcal{S}$-子模层，并且
\[
\label{II.3.2.4.1-zh}
  \widetilde{\mathcal{I}\mathcal{M}}
  =\mathcal{I}\cdot\widetilde{\mathcal{M}}
\tag{3.2.4.1}
\]
（右边按 \egaxref{0.4.3.5-zh}{0, 4.3.5} 中定义的意义理解）。事实上，
只需在 $Y=\operatorname{Spec}(A)$ 为仿射、
$\mathcal{S}=\widetilde{S}$（其中 $S$ 是一个分次 $A$-代数）、
$\mathcal{M}=\widetilde{M}$ 的情形验证这个公式，

\oldpage[II]{55}

其中 $M$ 是一个分次 $S$-模，并且
$\mathcal{I}=\widetilde{\mathfrak{I}}$，其中 $\mathfrak{I}$ 是 $A$ 的一个
理想。对 $S_+$ 的每个齐次元素 $f$，
\egaref{II.3.2.4.1-zh}{3.2.4.1} 左边在
$D_+(f)=\operatorname{Spec}(S_{(f)})$ 上的限制，与
$(\mathfrak{I}M)_{(f)}=\mathfrak{I}\cdot M_{(f)}$ 相伴；由
\egaxref{I.1.3.13-zh}{I, 1.3.13} 与
\egaxref{I.1.6.9-zh}{I, 1.6.9}，右边的限制亦然。
% Locked French authority: ega2-1-fr.tex, 820505 B,
% SHA-256 84EDBE3E83530AF2959B441796337C9DC21EAFCA6A13114A26778760FBF437AC.
% Sequential source scope: lines 4619-4739, subsection II.3.2.5 through 3.2.8 /
% printed pp.55-56.

\begin{proposition}[3.2.5]
\label{II.3.2.5-zh}
设 $f\in\Gamma(Y,\mathcal{S}_d)$（$d>0$）。在开集 $X_f$ 上，对每个
$n\in\mathbf{Z}$，$(\mathcal{O}_X|X_f)$-模层
$\widetilde{\mathcal{S}(nd)}|X_f$ 典范地同构于
$\mathcal{O}_X|X_f$。特别地，若 $\mathcal{O}_Y$-代数层 $\mathcal{S}$
由 $\mathcal{S}_1$ 生成（\egaref{II.3.1.9-zh}{3.1.9}），则对每个
$n\in\mathbf{Z}$，$\mathcal{O}_X$-模层
$\widetilde{\mathcal{S}(n)}$ 都是可逆的。
\end{proposition}

\begin{proof}
事实上，对 $Y$ 的每个仿射开集 $U$，结合
\egaref{II.3.1.4-zh}{3.1.4}（其中 $\varphi$ 是结构态射 $X\to Y$），
\egaref{II.2.5.7-zh}{2.5.7} 定义了从
$\widetilde{\mathcal{S}(nd)}|(X_f\cap\varphi^{-1}(U))$ 到
$\mathcal{O}_X|(X_f\cap\varphi^{-1}(U))$ 的典范同构。立即验证出，当把
$U$ 限制到仿射开集 $U'\subset U$ 时，这些同构彼此相容；第一个断言
由此得出。为证明第二个断言，只需注意：若 $\mathcal{S}$ 由
$\mathcal{S}_1$ 生成，则存在由仿射开集组成的 $Y$ 的覆盖
$(U_\alpha)$，使 $\Gamma(U_\alpha,\mathcal{S})$ 由
$\Gamma(U_\alpha,\mathcal{S})_1=\Gamma(U_\alpha,\mathcal{S}_1)$ 生成。
由于可逆性是局部性质，这时只需应用
\egaref{II.2.5.9-zh}{2.5.9} 的结果。
\end{proof}

仍对每个 $n\in\mathbf{Z}$ 令
\[
\label{II.3.2.5.1-zh}
  \mathcal{O}_X(n)=\widetilde{\mathcal{S}(n)}
\tag{3.2.5.1}
\]
，并对每个 $\mathcal{O}_X$-模层 $\mathcal{F}$ 令
\[
\label{II.3.2.5.2-zh}
  \mathcal{F}(n)=\mathcal{F}\otimes_{\mathcal{O}_X}\mathcal{O}_X(n).
\tag{3.2.5.2}
\]
由这些定义立即得到，对 $Y$ 的每个开集 $U$ 都有
\[
  \widetilde{(\mathcal{S}|U)(n)}
  =\mathcal{O}_X(n)|f^{-1}(U)
\]
，其中 $f$ 是结构态射 $X\to Y$。

\begin{proposition}[3.2.6]
\label{II.3.2.6-zh}
设 $\mathcal{M}$ 与 $\mathcal{N}$ 是两个拟凝聚分次 $\mathcal{S}$-模层。
存在关于 $\mathcal{M}$ 与 $\mathcal{N}$ 函子性的典范同态
\[
\label{II.3.2.6.1-zh}
  \lambda:\widetilde{\mathcal{M}}
  \otimes_{\mathcal{O}_X}\widetilde{\mathcal{N}}
  \longrightarrow\widetilde{\mathcal{M}\otimes_{\mathcal{S}}\mathcal{N}}
\tag{3.2.6.1}
\]
以及同样关于 $\mathcal{M}$ 与 $\mathcal{N}$ 函子性的典范同态
\[
\label{II.3.2.6.2-zh}
  \mu:\widetilde{\mathcal{H}om_{\mathcal{S}}(\mathcal{M},\mathcal{N})}
  \longrightarrow
  \mathcal{H}om_{\mathcal{O}_X}
  (\widetilde{\mathcal{M}},\widetilde{\mathcal{N}}).
\tag{3.2.6.2}
\]
此外，若 $\mathcal{S}$ 由 $\mathcal{S}_1$ 生成
（\egaref{II.3.1.9-zh}{3.1.9}），则 $\lambda$ 是同构；若再有
$\mathcal{M}$ 有限表示（\egaref{II.3.1.1-zh}{3.1.1}），则 $\mu$ 是
同构。
\end{proposition}

\begin{proof}
当 $Y$ 为仿射时，$\lambda$ 与 $\mu$ 已在
\egaref{II.2.5.11.2-zh}{2.5.11.2} 与
\egaref{II.2.5.12.2-zh}{2.5.12.2} 中定义。结合
\egaref{II.2.8.14-zh}{2.8.14}，这些定义是局部的，故立即推广到这里的
一般情形。
\end{proof}

\begin{corollary}[3.2.7]
\label{II.3.2.7-zh}
若 $\mathcal{S}$ 由 $\mathcal{S}_1$ 生成，则对 $\mathbf{Z}$ 中任意
$m$、$n$，在典范同构意义下有
\[
\label{II.3.2.7.1-zh}
  \mathcal{O}_X(m)\otimes_{\mathcal{O}_X}\mathcal{O}_X(n)
  =\mathcal{O}_X(m+n)
\tag{3.2.7.1}
\]
以及
\[
\label{II.3.2.7.2-zh}
  \mathcal{O}_X(n)=(\mathcal{O}_X(1))^{\otimes n}
\tag{3.2.7.2}
\]
\end{corollary}

\oldpage[II]{56}

\begin{corollary}[3.2.8]
\label{II.3.2.8-zh}
假设 $\mathcal{S}$ 由 $\mathcal{S}_1$ 生成。对每个分次
$\mathcal{S}$-模层 $\mathcal{M}$ 以及每个 $n\in\mathbf{Z}$，在典范
同构意义下有
\[
\label{II.3.2.8.1-zh}
  (\mathcal{M}(n))^{\sim}=\widetilde{\mathcal{M}}(n)
\tag{3.2.8.1}
\]
\end{corollary}

\begin{proof}
这来自 $Y$ 为仿射时的相应性质
（\egaref{II.2.5.14-zh}{2.5.14} 与
\egaref{II.2.5.15-zh}{2.5.15}）以及
\egaref{II.2.8.11-zh}{2.8.11}。
\end{proof}
% Locked French authority: ega2-1-fr.tex, 820505 B,
% SHA-256 84EDBE3E83530AF2959B441796337C9DC21EAFCA6A13114A26778760FBF437AC.
% Sequential source scope: lines 4741-4803, subsection II.3.2.9 through 3.2.10 /
% printed p.56.

\par\medskip\noindent\textit{注记 (3.2.9). ---\ }
\label{II.3.2.9-zh}
\begin{enumerate}
  \item[\rm(i)] 若 $\mathcal{S}=\mathcal{A}[T]$，其中 $\mathcal{A}$ 是
    一个拟凝聚 $\mathcal{O}_Y$-代数层
    （\egaref{II.3.1.7-zh}{3.1.7}），则立即验证出所有可逆
    $\mathcal{O}_X$-模层 $\mathcal{O}_X(n)$ 都典范地同构于
    $\mathcal{O}_X$。

    此外，设 $\mathcal{N}$ 是一个拟凝聚 $\mathcal{A}$-模层，并令
    $\mathcal{M}=\mathcal{N}\otimes_{\mathcal{A}}\mathcal{A}[T]$。由
    \egaref{II.3.2.3-zh}{3.2.3} 与
    \egaref{II.3.1.7-zh}{3.1.7}，在 $X=\operatorname{Proj}(\mathcal{A}[T])$
    与 $X'=\operatorname{Spec}(\mathcal{A})$ 的典范等同中，
    $\mathcal{O}_X$-模层 $\widetilde{\mathcal{M}}$ 等同于与
    $\mathcal{N}$ 相伴的 $\mathcal{O}_{X'}$-模层
    $\widetilde{\mathcal{N}}$（按 \egaref{II.1.4.3-zh}{1.4.3} 的意义）。
  \item[\rm(ii)] 设 $\mathcal{S}$ 是任一分次 $\mathcal{O}_Y$-代数层，
    而 $\mathcal{S}'$ 是满足 $\mathcal{S}'_0=\mathcal{O}_Y$、且对每个
    $n>0$ 都有 $\mathcal{S}'_n=\mathcal{S}_n$ 的分次
    $\mathcal{O}_Y$-代数层。典范同构
    $X=\operatorname{Proj}(\mathcal{S})$ 到
    $X'=\operatorname{Proj}(\mathcal{S}')$
    （\egaref{II.3.1.8-zh}{3.1.8, (ii)}）对每个 $n\in\mathbf{Z}$ 都把
    $\mathcal{O}_X(n)$ 与 $\mathcal{O}_{X'}(n)$ 等同。这来自仿射情形的
    同一命题（\egaref{II.2.5.16-zh}{2.5.16}），以及这些等同在 $Y$ 的
    各仿射开集上与限制运算交换这一事实。同样，令
    $X^{(d)}=\operatorname{Proj}(\mathcal{S}^{(d)})$；从 $X$ 到
    $X^{(d)}$ 的典范同构（\egaref{II.3.1.8-zh}{3.1.8, (i)}）对每个
    $n\in\mathbf{Z}$ 都把 $\mathcal{O}_X(nd)$ 与
    $\mathcal{O}_{X^{(d)}}(n)$ 等同。
\end{enumerate}

\begin{proposition}[3.2.10]
\label{II.3.2.10-zh}
设 $\mathcal{L}$ 是一个可逆 $\mathcal{O}_Y$-模层，而 $g$ 是典范同构
$X_{(\mathcal{L})}=\operatorname{Proj}(\mathcal{S}_{(\mathcal{L})})
\to X=\operatorname{Proj}(\mathcal{S})$
（\egaref{II.3.1.8-zh}{3.1.8, (iii)}）。对每个整数
$n\in\mathbf{Z}$，$g_*(\mathcal{O}_{X_{(\mathcal{L})}}(n))$ 典范地
同构于 $\mathcal{O}_X(n)\otimes_Y\mathcal{L}^{\otimes n}$。
\end{proposition}

\begin{proof}
先假设 $Y$ 为仿射的，其环为 $A$，并且
$\mathcal{L}=\widetilde{L}$，其中 $L$ 是秩一自由 $A$-模。沿用
\egaref{II.3.1.8-zh}{3.1.8, (iii)} 的证明中的记号，对 $f\in S_d$，
定义从 $(S(n))_{(f)}\otimes_A L^{\otimes n}$ 到
$(S_{(L)}(n))_{(f\otimes c^{\otimes d})}$ 的同构：把
$(x/f^k)\otimes c^{\otimes n}$（其中 $x\in S_{kd+n}$）对应到元素
$(x\otimes c^{\otimes(n+kd)})/(f\otimes c^{\otimes d})^k$。立即可见，
这个同构不依赖于为 $L$ 所选的生成元 $c$；此外，对每个
$f\in S_+$ 如此定义的同构与限制运算 $D_+(f)\to D_+(fg)$ 相容。最后，
在一般情形下，返回定义 \egaref{II.3.1.1-zh}{3.1.1} 不难看出：对
$Y$ 的每个仿射开集 $U$ 如此定义的同构，在把 $U$ 换成仿射开集
$U'\subset U$ 时彼此相容。
\end{proof}
% Locked French authority: ega2-1-fr.tex, 820505 B,
% SHA-256 84EDBE3E83530AF2959B441796337C9DC21EAFCA6A13114A26778760FBF437AC.
% Sequential source scope: lines 4805-4938, subsection II.3.3 through 3.3.2 /
% printed pp.56-58.

\subsection{与
\texorpdfstring{$\operatorname{Proj}(\mathcal{S})$}{Proj(S)} 上的层相伴的
分次 \texorpdfstring{$\mathcal{S}$}{S}-模层}
\label{subsection:II.3.3-zh}

\emph{本节始终假设分次 $\mathcal{O}_Y$-代数层 $\mathcal{S}$ 由
$\mathcal{S}_1$ 生成（\egaref{II.3.1.9-zh}{3.1.9}）。}回顾一下，由
\egaref{II.3.1.8-zh}{3.1.8, (i)}，在有限性条件
（\egaref{II.3.1.10-zh}{3.1.10}）下，这一限制并非本质性的。

\begin{env}[3.3.1]
\label{II.3.3.1-zh}
设 $p$ 是结构态射 $X=\operatorname{Proj}(\mathcal{S})\to Y$。对每个
$\mathcal{O}_X$-模层 $\mathcal{F}$，置
\[
\label{II.3.3.1.1-zh}
  \Gamma_*(\mathcal{F})
  =\bigoplus_{n\in\mathbf{Z}}p_*(\mathcal{F}(n))
\tag{3.3.1.1}
\]

\oldpage[II]{57}

并特别置
\[
\label{II.3.3.1.2-zh}
  \Gamma_*(\mathcal{O}_X)
  =\bigoplus_{n\in\mathbf{Z}}p_*(\mathcal{O}_X(n)).
\tag{3.3.1.2}
\]
已知（\egaxref{0.4.2.2-zh}{0, 4.2.2}），对两个 $\mathcal{O}_X$-模层
$\mathcal{F}$ 与 $\mathcal{G}$，存在典范同态
\[
  p_*(\mathcal{F})\otimes_{\mathcal{O}_Y}p_*(\mathcal{G})
  \longrightarrow
  p_*(\mathcal{F}\otimes_{\mathcal{O}_X}\mathcal{G}).
\]
因而由 \egaref{II.3.2.7.1-zh}{3.2.7.1} 可知，
$\Gamma_*(\mathcal{O}_X)$ 带有分次 $\mathcal{O}_Y$-代数层结构，而
\egaref{II.3.2.5.2-zh}{3.2.5.2} 同样在 $\Gamma_*(\mathcal{F})$ 上定义
了一个分次 $\Gamma_*(\mathcal{O}_X)$-模层结构。

由 \egaref{II.3.2.5-zh}{3.2.5} 以及函子 $p_*$ 的左正合性
（\egaxref{0.4.2.1-zh}{0, 4.2.1}），$\Gamma_*(\mathcal{F})$ 关于
$\mathcal{F}$ 是一个从 $\mathcal{O}_X$-模层范畴到分次
$\mathcal{O}_Y$-模层范畴的协变加性左正合函子（后一个范畴中的态射为
次数 $0$ 的同态）。特别地，若 $\mathcal{J}$ 是 $\mathcal{O}_X$ 中的
一个理想层，则 $\Gamma_*(\mathcal{J})$ 等同于
$\Gamma_*(\mathcal{O}_X)$ 的一个分次理想层。
\end{env}

\begin{env}[3.3.2]
\label{II.3.3.2-zh}
设 $\mathcal{M}$ 是一个拟凝聚分次 $\mathcal{S}$-模层。对 $Y$ 的每个
仿射开集 $U$，我们在 \egaref{II.2.6.2-zh}{2.6.2} 中定义了阿贝尔群
同态
\[
  \alpha_{0,U}:\Gamma(U,\mathcal{M}_0)
  \longrightarrow\Gamma(p^{-1}(U),\widetilde{\mathcal{M}}).
\]
这些同态显然与限制运算交换
（\egaref{II.2.8.13.1-zh}{2.8.13.1}），因而定义了阿贝尔群层同态
\[
\label{II.3.3.2.1-zh}
  \alpha_0:\mathcal{M}_0\longrightarrow p_*(\widetilde{\mathcal{M}}).
\tag{3.3.2.1}
\]
这里并未使用 $\mathcal{S}$ 由 $\mathcal{S}_1$ 生成的假设。把这一结果
应用于每个 $\mathcal{M}_n=(\mathcal{M}(n))_0$，并考虑
\egaref{II.3.2.8.1-zh}{3.2.8.1}，便定义了阿贝尔群层同态
\[
\label{II.3.3.2.2-zh}
  \alpha_n:\mathcal{M}_n\longrightarrow
  p_*(\widetilde{\mathcal{M}}(n))
  \qquad\text{对每个 }n\in\mathbf{Z},
\tag{3.3.2.2}
\]
从而得到分次阿贝尔群层的函子性同态（次数为 $0$）
\[
\label{II.3.3.2.3-zh}
  \alpha:\mathcal{M}\longrightarrow\Gamma_*(\widetilde{\mathcal{M}})
\tag{3.3.2.3}
\]
（也记为 $\alpha_{\mathcal{M}}$）。

特别取 $\mathcal{M}=\mathcal{S}$，可验证
$\alpha:\mathcal{S}\to\Gamma_*(\mathcal{O}_X)$ 是分次
$\mathcal{O}_Y$-代数层的同态，而
\egaref{II.3.3.2.3-zh}{3.3.2.3} 是分次模层的一个双同态，它相对于上述
分次代数层同态。

还要指出，每个 $\alpha_n$ 都对应于
（\egaxref{0.4.4.3-zh}{0, 4.4.3}）一个典范 $\mathcal{O}_X$-模层同态
\[
\label{II.3.3.2.4-zh}
  \alpha_n^\sharp:p^*(\mathcal{M}_n)
  \longrightarrow\widetilde{\mathcal{M}}(n).
\tag{3.3.2.4}
\]
容易验证，这一同态正是与分次 $\mathcal{O}_Y$-模层的典范同态（次数为
$0$）
\[
\label{II.3.3.2.5-zh}
  \mathcal{M}_n\otimes_{\mathcal{O}_Y}\mathcal{S}
  \longrightarrow\mathcal{M}(n)
\tag{3.3.2.5}
\]
按函子性对应的同态（\egaref{II.3.2.4-zh}{3.2.4}）。

\oldpage[II]{58}

右端的分次自然来自 $\mathcal{S}$ 的分次。事实上，只需考虑
$Y=\operatorname{Spec}(A)$ 为仿射、$\mathcal{M}=\widetilde{M}$ 且
$\mathcal{S}=\widetilde{S}$ 的情形，其中分次 $A$-代数 $S$ 由 $S_1$
生成，因而当 $f$ 遍历 $S_1$ 时，$D_+(f)$ 构成 $X$ 的一个覆盖。回到
\egaref{II.2.6.2-zh}{2.6.2} 的定义，并考虑
\egaxref{I.1.6.7-zh}{I, 1.6.7}，便可看出同态
\egaref{II.3.3.2.4-zh}{3.3.2.4} 在 $D_+(f)$ 上的限制，对应于
（\egaxref{I.1.3.8-zh}{I, 1.3.8}）$S_{(f)}$-模同态
$M_n\otimes_A S_{(f)}\to(M(n))_{(f)}$\footnote{法文印本此处把余域
误排为 \(\bigl(S(n)\bigr)_{(f)}\)；本译依同态
\egaref{II.3.3.2.5-zh}{3.3.2.5} 的类型采用
\(\bigl(M(n)\bigr)_{(f)}\)。参见源文勘误记录 EGA2-ZH-DEF-0004
及守护者处置 EGA2-CANON-DISP-0004。}；它把 $x\otimes1$
（$x\in M_n$）对应到 $x/1$。这就证明了所述断言。
\end{env}
% Locked French authority: ega2-1-fr.tex, 820505 B,
% SHA-256 84EDBE3E83530AF2959B441796337C9DC21EAFCA6A13114A26778760FBF437AC.
% Sequential source scope: lines 4940-5029, subsection II.3.3.3 through 3.3.5 /
% printed p.58.

\begin{proposition}[3.3.3]
\label{II.3.3.3-zh}
对每个截面 $f\in\Gamma(Y,\mathcal{S}_d)$（$d>0$），$X_f$ 正好是
$\alpha_d(f)$ 不消失的 $X$ 中的点所成的集合；这里将该截面视为
$\mathcal{O}_X(d)$ 的截面（\egaxref{0.5.5.2-zh}{0, 5.5.2}）。
\end{proposition}

\begin{proof}
（$\alpha_d(f)$ 是 $Y$ 上 $p_*(\mathcal{O}_X(d))$ 的一个截面，但按定义，
这样的截面也就是 $X$ 上 $\mathcal{O}_X(d)$ 的一个截面
（\egaxref{0.4.2.1-zh}{0, 4.2.1}）。）由 $X_f$ 的定义
（\egaref{II.3.1.4-zh}{3.1.4}），问题归结到 $Y$ 为仿射的情形，而该情形
已经在 \egaref{II.2.6.3-zh}{2.6.3} 中处理。
\end{proof}

\begin{env}[3.3.4]
\label{II.3.3.4-zh}
从现在起，除本节开头的假设外，我们还假设：对每个拟凝聚
$\mathcal{O}_X$-模层 $\mathcal{F}$，$p_*(\mathcal{F}(n))$ 在 $Y$ 上都是
拟凝聚的，因而
$\Gamma_*(\mathcal{F})=\bigoplus_{n\in\mathbf{Z}}p_*(\mathcal{F}(n))$
也是一个拟凝聚 $\mathcal{O}_Y$-模层
（\egaxref{I.1.4.1-zh}{I, 1.4.1} 与
\egaxref{I.1.3.9-zh}{I, 1.3.9}）。特别地，若 $X$ 在 $Y$ 上为有限型，
则上述情形总会发生（\egaxref{I.9.2.2-zh}{I, 9.2.2}）。由此可知，
$(\Gamma_*(\mathcal{F}))^{\sim}$ 有定义，并且是一个拟凝聚
$\mathcal{O}_X$-模层。对 $Y$ 的每个仿射开集 $U$，有
（\egaxref{I.1.3.9-zh}{I, 1.3.9} 与
\egaref{II.2.5.4-zh}{2.5.4}）
\begin{align*}
  \left(\Gamma\left(U,\bigoplus_{n\in\mathbf{Z}}
  p_*(\mathcal{F}(n))\right)\right)^{\sim}
  &=\bigoplus_{n\in\mathbf{Z}}
  \bigl(\Gamma(U,p_*(\mathcal{F}(n)))\bigr)^{\sim}\\
  &=\bigoplus_{n\in\mathbf{Z}}
  \bigl(\Gamma(p^{-1}(U),\mathcal{F}(n))\bigr)^{\sim}\\
  &=\left(\bigoplus_{n\in\mathbf{Z}}
  \Gamma(p^{-1}(U),\mathcal{F}(n))\right)^{\sim}\\
  &=\bigl(\Gamma_*(\mathcal{F}|p^{-1}(U))\bigr)^{\sim}.
\end{align*}
因此，由 \egaref{II.2.6.4-zh}{2.6.4} 得到典范同态
\[
  \beta_U:\left(\Gamma\left(U,\bigoplus_{n\in\mathbf{Z}}
  p_*(\mathcal{F}(n))\right)\right)^{\sim}
  \longrightarrow\mathcal{F}|p^{-1}(U).
\]

此外，\egaref{II.2.8.13.2-zh}{2.8.13.2} 的交换性表明，这些同态与 $Y$
上的限制运算交换；从而对拟凝聚 $\mathcal{O}_X$-模层得到典范的函子性
同态
\[
\label{II.3.3.4.1-zh}
  \beta:(\Gamma_*(\mathcal{F}))^{\sim}\longrightarrow\mathcal{F}
\tag{3.3.4.1}
\]
（也记为 $\beta_{\mathcal{F}}$）。
\end{env}

\begin{proposition}[3.3.5]
\label{II.3.3.5-zh}
设 $\mathcal{M}$ 是一个拟凝聚分次 $\mathcal{S}$-模层，$\mathcal{F}$ 是
一个拟凝聚 $\mathcal{O}_X$-模层；复合同态
\[
\label{II.3.3.5.1-zh}
  \widetilde{\mathcal{M}}
  \xrightarrow{\widetilde{\alpha}}
  (\Gamma_*(\widetilde{\mathcal{M}}))^{\sim}
  \xrightarrow{\beta}\widetilde{\mathcal{M}}
\tag{3.3.5.1}
\]
\[
\label{II.3.3.5.2-zh}
  \Gamma_*(\mathcal{F})\xrightarrow{\alpha}
  \Gamma_*((\Gamma_*(\mathcal{F}))^{\sim})
  \xrightarrow{\Gamma_*(\beta)}\Gamma_*(\mathcal{F})
\tag{3.3.5.2}
\]
都是恒等同构。
\end{proposition}

\begin{proof}
事实上，问题在 $Y$ 上是局部的，因而归结到
\egaref{II.2.6.5-zh}{2.6.5}。
\end{proof}
% Locked French authority: ega2-1-fr.tex, 820505 B,
% SHA-256 84EDBE3E83530AF2959B441796337C9DC21EAFCA6A13114A26778760FBF437AC.
% Sequential source scope: lines 5031-5126, subsection II.3.4 through 3.4.4 /
% printed p.59.

\subsection{有限性条件}
\label{subsection:II.3.4-zh}

\oldpage[II]{59}

\begin{proposition}[3.4.1]
\label{II.3.4.1-zh}
设 $Y$ 是一个预概形，$\mathcal{S}$ 是一个由 $\mathcal{S}_1$ 生成的拟凝聚
$\mathcal{O}_Y$-代数层（\egaref{II.3.1.9-zh}{3.1.9}），并且还假设
$\mathcal{S}_1$ 为有限型。则
$X=\operatorname{Proj}(\mathcal{S})$ 在 $Y$ 上为有限型。
\end{proposition}

\begin{proof}
事实上，问题在 $Y$ 上是局部的，因而可假设 $Y$ 是以 $A$ 为环的仿射
预概形；这时 $\mathcal{S}=\widetilde{S}$，其中
$S=\Gamma(Y,\mathcal{S})$，且按假设，$S$ 是由
$S_1=\Gamma(Y,\mathcal{S}_1)$ 生成的 $A$-代数，而还可假设 $S_1$ 是
有限型 $A$-模（\egaxref{I.1.3.9-zh}{I, 1.3.9} 与
\egaxref{I.1.3.12-zh}{I, 1.3.12}）。因此 $S$ 是有限型分次 $A$-代数，
问题归结到 \egaref{II.2.7.1-zh}{2.7.1, (ii)}。
\end{proof}

\begin{env}[3.4.2]
\label{II.3.4.2-zh}
设 $\mathcal{S}$ 是一个拟凝聚分次 $\mathcal{O}_Y$-代数层。对拟凝聚分次
$\mathcal{S}$-模层 $\mathcal{M}$，考虑如下有限性条件：
\begin{enumerate}
  \item[(\textbf{TF})] \emph{存在整数 $n$，使 $\mathcal{S}$-模层
    $\bigoplus_{k\geq n}\mathcal{M}_k$ 为有限型。}
  \item[(\textbf{TN})] \emph{存在整数 $n$，使
    $\mathcal{M}_k=0$ 对 $k\geq n$ 成立。}
\end{enumerate}
若 $\mathcal{M}$ 满足 (\textbf{TN})，则
$\widetilde{\mathcal{M}}=0$，因为问题在 $Y$ 上是局部的
（\egaref{II.2.7.2-zh}{2.7.2}）。

设 $\mathcal{M}$、$\mathcal{N}$ 是两个拟凝聚分次 $\mathcal{S}$-模层。
称次数为 $0$ 的同态 $u:\mathcal{M}\to\mathcal{N}$ 是
(\textbf{TN})-\emph{单射的}（相应地，(\textbf{TN})-\emph{满射的}、
(\textbf{TN})-\emph{双射的}），如果存在整数 $n$，使
$u_k:\mathcal{M}_k\to\mathcal{N}_k$ 对 $k\geq n$ 是单射的（相应地，
满射的、双射的）。这时，由 \egaref{II.2.7.2-zh}{2.7.2}，并考虑问题
在 $Y$ 上的局部性及 \egaxref{I.1.3.9-zh}{I, 1.3.9}，同态
$\widetilde{u}:\widetilde{\mathcal{M}}\to\widetilde{\mathcal{N}}$ 是单射的
（相应地，满射的、双射的）。当 $u$ 是 (\textbf{TN})-双射的时，也称
$u$ 是一个 (\textbf{TN})-\emph{同构}。
\end{env}

\begin{proposition}[3.4.3]
\label{II.3.4.3-zh}
设 $Y$ 是一个预概形，$\mathcal{S}$ 是一个由 $\mathcal{S}_1$ 生成的拟凝聚
分次 $\mathcal{O}_Y$-代数层，并假设 $\mathcal{S}_1$ 为有限型。设
$\mathcal{M}$ 是一个拟凝聚分次 $\mathcal{S}$-模层。
\begin{enumerate}
  \item[\rm(i)] 若 $\mathcal{M}$ 满足条件 (\textbf{TF})，则
    $\widetilde{\mathcal{M}}$ 为有限型。
  \item[\rm(ii)] 假设 $\mathcal{M}$ 满足 (\textbf{TF})。则
    $\widetilde{\mathcal{M}}=0$ 当且仅当 $\mathcal{M}$ 满足
    (\textbf{TN})。
\end{enumerate}
\end{proposition}

\begin{proof}
这些问题在 $Y$ 上是局部的，因而归结到如下情形：$Y$ 是以 $A$ 为环的
仿射预概形，$\mathcal{S}=\widetilde{S}$，其中 $S$ 是分次 $A$-代数且
理想 $S_+$ 为有限型；$\mathcal{M}=\widetilde{M}$，其中 $M$ 是分次
$S$-模。这时命题来自 \egaref{II.2.7.3-zh}{2.7.3}。
\end{proof}

\begin{theorem}[3.4.4]
\label{II.3.4.4-zh}
设 $Y$ 是一个预概形，$\mathcal{S}$ 是一个由 $\mathcal{S}_1$ 生成的拟凝聚
分次 $\mathcal{O}_Y$-代数层，并假设 $\mathcal{S}_1$ 为有限型；令
$X=\operatorname{Proj}(\mathcal{S})$。对每个拟凝聚
$\mathcal{O}_X$-模层 $\mathcal{F}$，典范同态 $\beta$
（\egaref{II.3.3.4-zh}{3.3.4}）是同构。
\end{theorem}

\begin{proof}
先注意，由 \egaref{II.3.4.1-zh}{3.4.1}，$\beta$ 有定义。为证明
$\beta$ 是同构，问题归结到如下情形：$Y$ 是以 $A$ 为环的仿射预概形，
$\mathcal{S}=\widetilde{S}$，其中 $S$ 是由 $S_1$ 生成的分次 $A$-代数，
且 $S_1$ 是有限型 $A$-模。这时只需应用
\egaref{II.2.7.5-zh}{2.7.5}。
\end{proof}
% Locked French authority: ega2-1-fr.tex, 820505 B,
% SHA-256 84EDBE3E83530AF2959B441796337C9DC21EAFCA6A13114A26778760FBF437AC.
% Sequential source scope: lines 5128-5191, subsection II.3.4.5 through 3.4.6 /
% printed pp.59-60.

\begin{corollary}[3.4.5]
\label{II.3.4.5-zh}
在 \egaref{II.3.4.4-zh}{3.4.4} 的假设下，每个拟凝聚
$\mathcal{O}_X$-模层 $\mathcal{F}$ 都同构于形如
$\widetilde{\mathcal{M}}$ 的 $\mathcal{O}_X$-模层，其中 $\mathcal{M}$
是一个拟凝聚分次 $\mathcal{S}$-模层。若再有 $\mathcal{F}$ 为有限型，
并假设 $Y$ 是一个拟紧概形，或 $Y$ 的底空间是诺特的，则可取
$\mathcal{M}$ 为有限型。
\end{corollary}

\oldpage[II]{60}

\begin{proof}
第一个断言立即由 \egaref{II.3.4.4-zh}{3.4.4} 得出，只须取
$\mathcal{M}=\Gamma_*(\mathcal{F})$。为证明第二个断言，只须证明
$\mathcal{M}$ 是其有限型分次 $\mathcal{S}$-子模层
$\mathcal{N}_\lambda$ 的归纳极限。事实上，由此可知
$\widetilde{\mathcal{M}}$ 是 $\widetilde{\mathcal{N}}_\lambda$ 的归纳
极限（\egaref{II.3.2.4-zh}{3.2.4}），因而 $\mathcal{F}$ 是
$\beta(\widetilde{\mathcal{N}}_\lambda)$ 的归纳极限；由于 $X$ 是拟紧的
（\egaref{II.3.4.1-zh}{3.4.1} 与
\egaxref{I.6.3.1-zh}{I, 6.3.1}），且 $\mathcal{F}$ 为有限型，
$\mathcal{F}$ 必等于某个 $\beta(\widetilde{\mathcal{N}}_\lambda)$
（\egaxref{0.5.2.3-zh}{0, 5.2.3}）。

为定义以 $\mathcal{M}$ 为归纳极限的 $\mathcal{N}_\lambda$，只须对每个
$n\in\mathbf{Z}$ 考虑拟凝聚 $\mathcal{O}_Y$-模层 $\mathcal{M}_n$。
由关于 $Y$ 的假设（\egaxref{I.9.4.9-zh}{I, 9.4.9}），它是其有限型
$\mathcal{O}_Y$-子模层 $\mathcal{M}_n^{(\mu_n)}$ 的归纳极限。显然，
$\mathcal{P}_{\mu_n}=\mathcal{S}\cdot\mathcal{M}_n^{(\mu_n)}$ 是有限型
分次 $\mathcal{S}$-模层；立即可验证，取 $\mathcal{N}_\lambda$ 为形如
$\mathcal{P}_{\mu_n}$ 的 $\mathcal{S}$-模层的有限和，便得到所需对象。
\end{proof}

\begin{corollary}[3.4.6]
\label{II.3.4.6-zh}
假设 \egaref{II.3.4.4-zh}{3.4.4} 的条件成立，并且 $Y$ 的底空间拟紧；
设 $\mathcal{F}$ 是一个有限型拟凝聚 $\mathcal{O}_X$-模层。则存在
$n_0$，使得对 $n\geq n_0$，典范同态
$\sigma:p^*(p_*(\mathcal{F}(n)))\to\mathcal{F}(n)$
（\egaxref{0.4.4.3-zh}{0, 4.4.3}）是满射。
\end{corollary}

\begin{proof}
事实上，对每个 $y\in Y$，取 $Y$ 中 $y$ 的一个仿射开邻域 $U$。存在整数
$n_0(U)$，使得对 $n\geq n_0(U)$，
$\mathcal{F}(n)|p^{-1}(U)$ 由它在 $p^{-1}(U)$ 上的有限多个截面生成
（\egaref{II.2.7.9-zh}{2.7.9}）。但这些截面是
$p^*(p_*(\mathcal{F}(n)))$ 在 $p^{-1}(U)$ 上的截面的典范像
（\egaxref{0.3.7.1-zh}{0, 3.7.1} 与
\egaxref{0.4.4.3-zh}{0, 4.4.3}），所以
$\mathcal{F}(n)|p^{-1}(U)$ 等于
$p^*(p_*(\mathcal{F}(n)))|p^{-1}(U)$ 的典范像。最后，由于 $Y$ 拟紧，
存在由仿射开集 $U_i$ 组成的 $Y$ 的有限覆盖；取 $n_0$ 为所有
$n_0(U_i)$ 中的最大者，证明即告完成。
\end{proof}
% Locked French authority: ega2-1-fr.tex, 820505 B,
% SHA-256 84EDBE3E83530AF2959B441796337C9DC21EAFCA6A13114A26778760FBF437AC.
% Sequential source scope: lines 5193-5274, subsection II.3.4.7 through 3.4.8 /
% printed pp.60-61.

\par\medskip\noindent\textit{注记 (3.4.7). ---\ }
\label{II.3.4.7-zh}
设 $p=(\psi,\theta):X\to Y$ 是环化空间的一个态射，$\mathcal{F}$ 是一个
$\mathcal{O}_X$-模层。典范同态
$\sigma:p^*(p_*(\mathcal{F}))\to\mathcal{F}$ 为满射这一事实，可以具体
表述如下（\egaxref{0.4.4.1-zh}{0, 4.4.1}）：对每个 $x\in X$，以及
$\mathcal{F}$ 在 $x$ 的某个开邻域 $V$ 上的每个截面 $s$，存在 $Y$ 中
$p(x)$ 的一个开邻域 $U$、$\mathcal{F}$ 在 $p^{-1}(U)$ 上的有限多个截面
$t_i$（$1\leq i\leq m$）、$x$ 的一个邻域
$W\subset V\cap p^{-1}(U)$，以及 $\mathcal{O}_X$ 在 $W$ 上的截面
$a_i$（$1\leq i\leq m$），使得
\[
  s|W=\sum_i a_i\cdot(t_i|W).
\]
当 $Y$ 是一个\emph{仿射概形}且 $p_*(\mathcal{F})$ \emph{拟凝聚}时，
这一条件等价于 $\mathcal{F}$ \emph{由其在 $X$ 上的截面生成}
（\egaxref{0.5.5.1-zh}{0, 5.5.1}）。事实上，若
$Y=\operatorname{Spec}(A)$，可假设 $U=D(f)$，其中 $f\in A$。把
\egaxref{I.1.4.1-zh}{I, 1.4.1} 应用于 $p_*(\mathcal{F})$，便知存在整数
$n>0$ 以及 $\mathcal{F}$ 在 $X$ 上的截面 $s_i$，使 $t_i$ 是
$s_ig^n$ 在 $p^{-1}(U)$ 上的限制，其中 $g=\theta(f)$。由于 $g$ 在
$p^{-1}(U)$ 上可逆，遂有
\[
  s|W=\sum_i b_i\cdot(s_i|W)
\]
其中 $b_i=a_i(g|W)^{-n}$，断言由此得证。因此，当 $Y$ 为仿射时，推论
\egaref{II.3.4.6-zh}{3.4.6} 重新给出
\egaref{II.2.7.9-zh}{2.7.9}。

\oldpage[II]{61}

由此得出：当 $Y$ 是任意预概形时，对一个满足 $p_*(\mathcal{F})$
\emph{拟凝聚}的拟凝聚 $\mathcal{O}_X$-模层 $\mathcal{F}$，以下三个条件
等价：
\begin{enumerate}
  \item[a)] \emph{典范同态
    $\sigma:p^*(p_*(\mathcal{F}))\to\mathcal{F}$ 是满射。}
  \item[b)] \emph{存在一个拟凝聚 $\mathcal{O}_Y$-模层 $\mathcal{G}$
    以及一个满同态 $p^*(\mathcal{G})\to\mathcal{F}$。}
  \item[c)] \emph{对 $Y$ 的每个仿射开集 $U$，
    $\mathcal{F}|p^{-1}(U)$ 由其在 $p^{-1}(U)$ 上的截面生成。}
\end{enumerate}
我们刚证明了 \emph{a)} 与 \emph{c)} 等价。另一方面，由假设
$p_*(\mathcal{F})$ 拟凝聚，\emph{a)} 显然蕴涵 \emph{b)}。反过来，
每个同态 $u:p^*(\mathcal{G})\to\mathcal{F}$ 都分解为
$p^*(\mathcal{G})\to p^*(p_*(\mathcal{F}))
\xrightarrow{\sigma}\mathcal{F}$
（\egaxref{0.3.5.4.4-zh}{0, 3.5.4.4}）；所以，若 $u$ 为满射，则
$\sigma$ 也为满射，这证明了 \emph{b)} 蕴涵 \emph{a)}。

\begin{corollary}[3.4.8]
\label{II.3.4.8-zh}
假设 \egaref{II.3.4.4-zh}{3.4.4} 的条件成立，并进一步假设 $Y$ 是一个
拟紧概形，或 $Y$ 的底空间是诺特的。设 $\mathcal{F}$ 是一个有限型
拟凝聚 $\mathcal{O}_X$-模层；则存在整数 $n_0$，使得对 $n\geq n_0$，
$\mathcal{F}$ 同构于一个形如 $(p^*(\mathcal{G}))(-n)$ 的
$\mathcal{O}_X$-模层的商，其中 $\mathcal{G}$ 是一个有限型拟凝聚
$\mathcal{O}_Y$-模层（依赖于 $n$）。
\end{corollary}

\begin{proof}
由于结构态射 $X\to Y$ 是有限型的分离态射，
$p_*(\mathcal{F}(n))$ 是拟凝聚的
（\egaxref{I.9.2.2-zh}{I, 9.2.2, b)}）；而由关于 $Y$ 的假设
（\egaxref{I.9.4.9-zh}{I, 9.4.9}），它是其有限型拟凝聚
$\mathcal{O}_Y$-子模层的归纳极限。由
\egaref{II.3.4.6-zh}{3.4.6}、
\egaxref{0.4.3.2-zh}{0, 4.3.2} 与
\egaxref{0.5.2.3-zh}{0, 5.2.3} 可知，$\mathcal{F}(n)$ 是形如
$p^*(\mathcal{G})$ 的 $\mathcal{O}_X$-模层的典范像，其中
$\mathcal{G}$ 是 $p_*(\mathcal{F}(n))$ 的一个有限型拟凝聚
$\mathcal{O}_Y$-子模层。于是该推论来自
\egaref{II.3.2.5.2-zh}{3.2.5.2} 与
\egaref{II.3.2.7.1-zh}{3.2.7.1}。
\end{proof}
% Locked French authority: ega2-1-fr.tex, 820505 B,
% SHA-256 84EDBE3E83530AF2959B441796337C9DC21EAFCA6A13114A26778760FBF437AC.
% Sequential source scope: lines 5276-5359, subsection II.3.5 through (3.5.2.1) /
% printed pp.61-62.

\subsection{函子性}
\label{subsection:II.3.5-zh}

\begin{env}[3.5.1]
\label{II.3.5.1-zh}
设 $Y$ 是一个预概形，$\mathcal{S}$、$\mathcal{S}'$ 是两个正次数拟凝聚
分次 $\mathcal{O}_Y$-代数层；置
$X=\operatorname{Proj}(\mathcal{S})$、
$X'=\operatorname{Proj}(\mathcal{S}')$，并设 $p$、$p'$ 分别为 $X$、$X'$
到 $Y$ 的结构态射。设
$\varphi:\mathcal{S}'\to\mathcal{S}$ 是分次代数层的一个
$\mathcal{O}_Y$-同态。对 $Y$ 的每个仿射开集 $U$，置
$S_U=\Gamma(U,\mathcal{S})$、$S'_U=\Gamma(U,\mathcal{S}')$；再置
$A_U=\Gamma(U,\mathcal{O}_Y)$，同态 $\varphi$ 定义了分次 $A_U$-代数的同态
$\varphi_U:S'_U\to S_U$。与之在 $p^{-1}(U)$ 中相应的是一个开集
$G(\varphi_U)$ 以及态射
$\Phi_U:G(\varphi_U)\to p'^{-1}(U)$
（\egaref{II.2.8.1-zh}{2.8.1}）。此外，若 $V\subset U$ 是仿射开集，则图表
\[
\label{II.3.5.1.1-zh}
  \xymatrix{
    S'_U\ar[r]^{\varphi_U}\ar[d]
    & S_U\ar[d]\\
    S'_V\ar[r]_{\varphi_V}
    & S_V
  }
\tag{3.5.1.1}
\]
交换；从定义（\egaref{II.2.8.1-zh}{2.8.1}）立即验证
\[
  G(\varphi_V)=G(\varphi_U)\cap p^{-1}(V)
\]
并且 $\Phi_V$ 是 $\Phi_U$ 在 $G(\varphi_V)$ 上的限制。这样便定义了 $X$
的一个开部 $G(\varphi)$，使得对每个仿射开集 $U\subset Y$ 都有
$G(\varphi)\cap p^{-1}(U)=G(\varphi_U)$，并得到一个 $Y$-态射

\oldpage[II]{62}

$\Phi:G(\varphi)\to X'$；它是仿射的，称为\emph{与 $\varphi$ 相伴}的态射，
并记作 $\operatorname{Proj}(\varphi)$。若对每个 $y\in Y$，都存在 $y$ 的
一个仿射邻域 $U$，使 $\Gamma(U,\mathcal{O}_Y)$-模
$\Gamma(U,\mathcal{S}_+)$ 由
$\varphi(\Gamma(U,\mathcal{S}'_+))$ 生成，则
$G(\varphi_U)=p^{-1}(U)$，从而 $G(\varphi)=X$。
\end{env}

\begin{proposition}[3.5.2]
\label{II.3.5.2-zh}
\begin{enumerate}
  \item[\rm(i)] 若 $\mathcal{M}$ 是拟凝聚分次 $\mathcal{S}$-模层，则存在
    一个函子性的典范同构，把 $\mathcal{O}_{X'}$-模层
    $(\mathcal{M}_{[\varphi]})^{\sim}$ 同构到 $\mathcal{O}_{X'}$-模层
    $\Phi_*(\widetilde{\mathcal{M}}|G(\varphi))$。
  \item[\rm(ii)] 若 $\mathcal{M}'$ 是拟凝聚分次 $\mathcal{S}'$-模层，则存在
    一个函子性的典范同态 $\nu$，从
    $(\mathcal{O}_X|G(\varphi))$-模层
    $\Phi^*(\widetilde{\mathcal{M}'})$ 到
    $(\mathcal{O}_X|G(\varphi))$-模层
    $(\mathcal{M}'\otimes_{\mathcal{S}'}\mathcal{S})^{\sim}|G(\varphi)$。
    若 $\mathcal{S}'$ 由 $\mathcal{S}'_1$ 生成，则 $\nu$ 是同构。
\end{enumerate}
\end{proposition}

\begin{proof}
事实上，当 $Y$ 为仿射时，所考虑的同态已经定义
（\egaref{II.2.8.7-zh}{2.8.7} 与
\egaref{II.2.8.8-zh}{2.8.8}）；在一般情形，只需验证它们与从 $Y$ 的
一个仿射开集限制到更小仿射开集的运算相容，而这立即由
\egaref{II.3.5.1.1-zh}{3.5.1.1} 的交换性得出。
\end{proof}

特别地，对每个 $n\in\mathbf{Z}$，都有典范同态
\[
\label{II.3.5.2.1-zh}
  \Phi^*(\mathcal{O}_{X'}(n))\longrightarrow
  \mathcal{O}_X(n)|G(\varphi).
\tag{3.5.2.1}
\]
% Locked French authority: ega2-1-fr.tex, 820505 B,
% SHA-256 84EDBE3E83530AF2959B441796337C9DC21EAFCA6A13114A26778760FBF437AC.
% Sequential source scope: lines 5361-5419, proposition II.3.5.3 through
% corollary II.3.5.4 / printed p.62.

\begin{proposition}[3.5.3]
\label{II.3.5.3-zh}
设 $Y$、$Y'$ 是两个预概形，$\psi:Y'\to Y$ 是一个态射，$\mathcal{S}$ 是
拟凝聚分次 $\mathcal{O}_Y$-代数层，并置
$\mathcal{S}'=\psi^*(\mathcal{S})$。则 $Y'$-概形
$X'=\operatorname{Proj}(\mathcal{S}')$ 典范地等同于
$\operatorname{Proj}(\mathcal{S})\times_Y Y'$。此外，若 $\mathcal{M}$ 是
拟凝聚分次 $\mathcal{S}$-模层，则 $\mathcal{O}_{X'}$-模层
$(\psi^*(\mathcal{M}))^{\sim}$ 等同于
$\widetilde{\mathcal{M}}\otimes_Y\mathcal{O}_{Y'}$。
\end{proposition}

\begin{proof}
首先注意，$\psi^*(\mathcal{S})$ 与 $\psi^*(\mathcal{M})$ 都是拟凝聚
$\mathcal{O}_{Y'}$-模层，其齐次分量也是如此
（\egaxref{0.5.1.4-zh}{0, 5.1.4}）。设 $U$ 是 $Y$ 的一个仿射开集，
$U'\subset\psi^{-1}(U)$ 是 $Y'$ 的一个仿射开集，$A$、$A'$ 分别是
$U$、$U'$ 的环；于是 $\mathcal{S}|U=\widetilde{S}$，其中 $S$ 是分次
$A$-代数，而 $\mathcal{S}'|U'$ 等同于
$(S\otimes_A A')^{\sim}$（\egaxref{I.1.6.5-zh}{I, 1.6.5}）。第一项断言
遂由 \egaref{II.2.8.10-zh}{2.8.10} 与
\egaxref{I.3.2.6.2-zh}{I, 3.2.6.2} 得出：由前述等同定义的投影
$\operatorname{Proj}(\mathcal{S}'|U')\to
\operatorname{Proj}(\mathcal{S}|U)$ 显然与 $U$、$U'$ 上的限制运算相容，
因而确实定义了态射
$\operatorname{Proj}(\mathcal{S}')\to\operatorname{Proj}(\mathcal{S})$。
现在设
\[
  q:\operatorname{Proj}(\mathcal{S})\to Y,
  \qquad q':\operatorname{Proj}(\mathcal{S}')\to Y'
\]
为结构态射；此时 $q'^{-1}(U')$ 等同于 $q^{-1}(U)\times_U U'$，并且两个
层
$(\psi^*(\mathcal{M}))^{\sim}|q'^{-1}(U')$ 与
$(\widetilde{\mathcal{M}}\otimes_Y\mathcal{O}_{Y'})|q'^{-1}(U')$
都典范地等同于 $(M\otimes_A A')^{\sim}$，其中置
$M=\Gamma(U,\mathcal{M})$；这是
\egaref{II.2.8.10-zh}{2.8.10} 与
\egaxref{I.1.6.5-zh}{I, 1.6.5} 的结果。第二项断言由此得出，因为上述等同
与限制运算的相容性同样立即可见。
\end{proof}

\begin{corollary}[3.5.4]
\label{II.3.5.4-zh}
沿用 \egaref{II.3.5.3-zh}{3.5.3} 的记号，对每个
$n\in\mathbf{Z}$，$\mathcal{O}_{X'}(n)$ 典范地等同于
$\mathcal{O}_X(n)\otimes_Y\mathcal{O}_{Y'}$（其中
$X=\operatorname{Proj}(\mathcal{S})$）。
\end{corollary}

\begin{proof}
事实上，沿用 \egaref{II.3.5.3-zh}{3.5.3} 的记号，显然对每个
$n\in\mathbf{Z}$ 都有
$\psi^*(\mathcal{S}(n))=\mathcal{S}'(n)$。
\end{proof}
% Locked French authority: ega2-1-fr.tex, 820505 B,
% SHA-256 84EDBE3E83530AF2959B441796337C9DC21EAFCA6A13114A26778760FBF437AC.
% Sequential source scope: lines 5421-5489, environment II.3.5.5 /
% printed pp.62-63.

\begin{env}[3.5.5]
\label{II.3.5.5-zh}
沿用前面的记号，以 $\Psi$ 表示典范投影 $X'\to X$，并置
$\mathcal{M}'=\psi^*(\mathcal{M})$；此外，假设 $\mathcal{S}$ 由
$\mathcal{S}_1$ 生成，且 $X$ 在 $Y$ 上为有限型；于是

\oldpage[II]{63}

$\mathcal{S}'$ 由 $\mathcal{S}'_1$ 生成（只要归结到 $Y$、$Y'$ 均为
仿射的情形即可看出），并且 $X'$ 在 $Y'$ 上为有限型
（\egaxref{I.6.3.4-zh}{I, 6.3.4}）。设 $\mathcal{F}$ 是一个
$\mathcal{O}_X$-模层，并置 $\mathcal{F}'=\Psi^*(\mathcal{F})$；由
\egaref{II.3.5.4-zh}{3.5.4} 与
\egaxref{0.4.3.3-zh}{0, 4.3.3}，对每个 $n\in\mathbf{Z}$ 都有
$\mathcal{F}'(n)=\Psi^*(\mathcal{F}(n))$。此外，按如下方式定义典范的
$\psi$-同态
$\theta_n:q_*(\mathcal{F}(n))\to q'_*(\mathcal{F}'(n))$。鉴于图表
\[
  \xymatrix{
    X\ar[d]_{q}
    & X'\ar[l]_{\Psi}\ar[d]^{q'}\\
    Y
    & Y'\ar[l]^{\psi}
  }
\]
交换，需要定义的是同态
$q_*(\mathcal{F}(n))\to
\psi_*(q'_*(\Psi^*(\mathcal{F}(n))))
=q_*(\Psi_*(\Psi^*(\mathcal{F}(n))))$；只需取
$\theta_n=q_*(\rho_n)$，其中 $\rho_n$ 是典范同态
$\mathcal{F}(n)\to\Psi_*(\Psi^*(\mathcal{F}(n)))$
（\egaxref{0.4.4.3-zh}{0, 4.4.3}）。显然，对 $Y$ 的每个仿射开集 $U$
以及 $Y'$ 中满足 $U'\subset\psi^{-1}(U)$ 的每个仿射开集 $U'$，同态
$\theta_n$ 在截面上给出典范同态
（\egaxref{0.3.7.2-zh}{0, 3.7.2}）
$\Gamma(q^{-1}(U),\mathcal{F}(n))\to
\Gamma(q'^{-1}(U'),\mathcal{F}'(n))$。

于是，\egaref{II.2.8.13.2-zh}{2.8.13.2} 的交换性表明，若
$\mathcal{F}$ 拟凝聚，则图表
\[
  \xymatrix{
    \mathcal{F}\ar[r]^{\rho}
    & \mathcal{F}'\\
    (\Gamma_*(\mathcal{F}))^{\sim}
      \ar[u]^{\beta_{\mathcal{F}}}\ar[r]_{\widetilde{\theta}}
    & (\Gamma_*(\mathcal{F}'))^{\sim}
      \ar[u]_{\beta_{\mathcal{F}'}}
  }
\]
交换（上方的水平箭头是典范 $\Psi$-态射
$\mathcal{F}\to\Psi^*(\mathcal{F})$）。

同样，\egaref{II.2.8.13.1-zh}{2.8.13.1} 的交换性表明，图表
\[
  \xymatrix{
    \Gamma_*(\widetilde{\mathcal{M}})\ar[r]^{\theta}
    & \Gamma_*(\widetilde{\mathcal{M}'})\\
    \mathcal{M}\ar[r]_{\rho}\ar[u]^{\alpha_{\mathcal{M}}}
    & \mathcal{M}'\ar[u]_{\alpha_{\mathcal{M}'}}
  }
\]
交换（下方的水平箭头是典范 $\psi$-态射
$\mathcal{M}\to\psi^*(\mathcal{M})$）。
\end{env}
% Locked French authority: ega2-1-fr.tex, 820505 B,
% SHA-256 84EDBE3E83530AF2959B441796337C9DC21EAFCA6A13114A26778760FBF437AC.
% Sequential source scope: lines 5491-5561, environment II.3.5.6 /
% printed pp.63-64.

\begin{env}[3.5.6]
\label{II.3.5.6-zh}
现在考虑两个预概形 $Y$、$Y'$，一个态射 $g:Y'\to Y$，一个拟凝聚分次
$\mathcal{O}_Y$-代数层 $\mathcal{S}$（相应地，一个拟凝聚分次
$\mathcal{O}_{Y'}$-代数层 $\mathcal{S}'$），以及分次代数层的一个
$g$-态射 $u:\mathcal{S}\to\mathcal{S}'$；也就是说，它是分次代数层的
一个 $\mathcal{O}_Y$-同态
$\mathcal{S}\to g_*(\mathcal{S}')$。还知道，这等价于给定分次代数层的
一个 $\mathcal{O}_{Y'}$-同态
$u^\sharp:g^*(\mathcal{S})\to\mathcal{S}'$。由 $u^\sharp$ 典范地得到
$Y'$-态射
$W=\operatorname{Proj}(u^\sharp):G(u^\sharp)\to
\operatorname{Proj}(g^*(\mathcal{S}))$，其中 $G(u^\sharp)$ 是
$X'=\operatorname{Proj}(\mathcal{S}')$ 的一个开部
（\egaref{II.3.5.1-zh}{3.5.1}）。另一方面，置
$X=\operatorname{Proj}(\mathcal{S})$，则
$X''=\operatorname{Proj}(g^*(\mathcal{S}))$ 典范地等同于

\oldpage[II]{64}

$X\times_Y Y'$（\egaref{II.3.5.3-zh}{3.5.3}）。把
$\operatorname{Proj}(u^\sharp)$ 与第一投影
$p:X\times_Y Y'\to X$ 复合，便得到态射
$v:G(u^\sharp)\to X$；我们把它记作 $\operatorname{Proj}(u)$，并且图表
\[
  \xymatrix{
    G(u^\sharp)\ar[r]^{v}\ar[d]
    & X\ar[d]\\
    Y'\ar[r]_{g}
    & Y
  }
\]
交换。

此外，对每个拟凝聚分次 $\mathcal{S}$-模层 $\mathcal{M}$，都有典范的
$v$-态射
\[
\label{II.3.5.6.1-zh}
  \nu:\widetilde{\mathcal{M}}\longrightarrow
  (g^*(\mathcal{M})\otimes_{g^*(\mathcal{S})}\mathcal{S}')^{\sim}
  |G(u^\sharp).
\tag{3.5.6.1}
\]

事实上，$\nu^\sharp$ 由下列同态复合而得：
\[
  v^*(\widetilde{\mathcal{M}})
  =W^*(p^*(\widetilde{\mathcal{M}}))
  \longrightarrow W^*((g^*(\mathcal{M}))^{\sim})
  \longrightarrow
  (g^*(\mathcal{M})\otimes_{g^*(\mathcal{S})}\mathcal{S}')^{\sim}
  |G(u^\sharp)
\]
其中第一条箭头来自 \egaref{II.3.5.3-zh}{3.5.3} 的同构，第二条箭头是
\egaref{II.3.5.2-zh}{3.5.2, (ii)} 的同态；若 $\mathcal{S}$ 由
$\mathcal{S}_1$ 生成，则由 \egaref{II.3.5.2-zh}{3.5.2} 可知
$\nu^\sharp$ 是同构。

作为 \egaref{II.3.5.6.1-zh}{3.5.6.1} 的特例，对每个
$n\in\mathbf{Z}$ 都有典范的 $v$-态射
\[
\label{II.3.5.6.2-zh}
  \nu:\mathcal{O}_X(n)\longrightarrow
  \mathcal{O}_{X'}(n)|G(u^\sharp).
\tag{3.5.6.2}
\]
\end{env}
% Locked French authority: ega2-1-fr.tex, 820505 B,
% SHA-256 84EDBE3E83530AF2959B441796337C9DC21EAFCA6A13114A26778760FBF437AC.
% Sequential source scope: lines 5563-5620, subsection II.3.6 through
% proposition II.3.6.2 / printed pp.64-65.

\subsection{预概形 \texorpdfstring{$\operatorname{Proj}(\mathcal{S})$}{Proj(S)}
的闭子预概形}
\label{subsection:II.3.6-zh}

\begin{env}[3.6.1]
\label{II.3.6.1-zh}
设 $Y$ 是一个预概形，
$\varphi:\mathcal{S}\to\mathcal{S}'$ 是次数为 $0$ 的拟凝聚分次
$\mathcal{O}_Y$-代数层同态。若存在 $n$，使得当 $k\geq n$ 时，
$\varphi_k:\mathcal{S}_k\to\mathcal{S}'_k$ 是满射（相应地，单射、
双射），则称 $\varphi$ 是 (\textbf{TN})-\emph{满射的}（相应地，
(\textbf{TN})-\emph{单射的}、(\textbf{TN})-\emph{双射的}）。在这种情形，
对相应态射
$\Phi:\operatorname{Proj}(\mathcal{S}')\to\operatorname{Proj}(\mathcal{S})$
的研究可归结到 $\varphi$ 为满射（相应地，单射、双射）的情形。仿照
\egaref{II.2.9.1-zh}{2.9.1}（即 $Y$ 为仿射时的特例）并使用
\egaref{II.3.1.8-zh}{3.1.8} 即可证明这一点。还把 (\textbf{TN})-双射的
$\varphi$ 称为一个 (\textbf{TN})-同构。
\end{env}

\begin{proposition}[3.6.2]
\label{II.3.6.2-zh}
设 $Y$ 是一个预概形，$\mathcal{S}$ 是拟凝聚分次 $\mathcal{O}_Y$-代数层，
并置 $X=\operatorname{Proj}(\mathcal{S})$。
\begin{enumerate}
  \item[\rm(i)] 若 $\varphi:\mathcal{S}\to\mathcal{S}'$ 是分次
    $\mathcal{O}_Y$-代数层的一个 (\textbf{TN})-满同态，则相应态射
    $\Phi=\operatorname{Proj}(\varphi)$
    （\egaref{II.3.5.1-zh}{3.5.1}）在整个
    $\operatorname{Proj}(\mathcal{S}')$ 上有定义，并且是
    $\operatorname{Proj}(\mathcal{S}')$ 到 $X$ 的一个闭浸入。若
    $\mathcal{J}$ 是 $\varphi$ 的核，则与 $\Phi$ 相伴的 $X$ 的闭子预概形
    由 $\mathcal{O}_X$ 的拟凝聚理想层
    $\widetilde{\mathcal{J}}$ 定义。
  \item[\rm(ii)] 进一步假设 $\mathcal{S}_0=\mathcal{O}_Y$，
    $\mathcal{S}$ 由 $\mathcal{S}_1$ 生成，并且 $\mathcal{S}_1$ 为有限型。
    设 $X'$ 是 $X=\operatorname{Proj}(\mathcal{S})$ 的一个闭子预概形，由
    $\mathcal{O}_X$ 的拟凝聚理想层 $\mathcal{I}$ 定义。设
    $\mathcal{J}$

\oldpage[II]{65}

    是 $\mathcal{S}$ 的拟凝聚分次理想层，即典范同态
    $\alpha:\mathcal{S}\to\Gamma_*(\mathcal{O}_X)$
    （\egaref{II.3.3.2-zh}{3.3.2}）下 $\Gamma_*(\mathcal{I})$ 的逆像，并置
    $\mathcal{S}'=\mathcal{S}/\mathcal{J}$。则 $X'$ 是与闭浸入
    $\operatorname{Proj}(\mathcal{S}')\to X$ 相伴的闭子预概形
    （\egaxref{I.4.2.1-zh}{I, 4.2.1}），该闭浸入对应于分次
    $\mathcal{O}_Y$-代数层的典范同态
    $\mathcal{S}\to\mathcal{S}'$。
\end{enumerate}
\end{proposition}
% Locked French authority: ega2-1-fr.tex, 820505 B,
% SHA-256 84EDBE3E83530AF2959B441796337C9DC21EAFCA6A13114A26778760FBF437AC.
% Sequential source scope: lines 5622-5675, proof of II.3.6.2 through
% corollary II.3.6.4 / printed p.65.
% EGA2-ZH-DEF-0006 pending: retain source-literal G(varphi)=X at line 5628.
% EGA2-ZH-DEF-0005 pending: retain source-literal reference 3.6.1(i) at line 5674.

\begin{proof}
\begin{enumerate}
  \item[\rm(i)] 可以假设 $\varphi$ 为满射
    （\egaref{II.3.6.1-zh}{3.6.1}）。于是，对 $Y$ 的每个仿射开集 $U$，
    $\Gamma(U,\mathcal{S})\to\Gamma(U,\mathcal{S}')$ 是满射
    （\egaxref{I.1.3.9-zh}{I, 1.3.9}），所以由
    \egaref{II.3.5.1-zh}{3.5.1} 有 $G(\varphi)=X$。问题立即归结到 $Y$
    为仿射时证明该命题，而此时结论来自
    \egaref{II.2.9.2-zh}{2.9.2, (i)}。
  \item[\rm(ii)] 问题归结为证明：由典范单同态
    $\mathcal{J}\to\mathcal{S}$ 得到的同态
    $\widetilde{\mathcal{J}}\to\mathcal{O}_X$，把
    $\widetilde{\mathcal{J}}$ 同构到 $\mathcal{I}$ 上。由于问题在 $Y$
    上是局部的，可以假设 $Y$ 是环 $A$ 的仿射概形；于是
    $\mathcal{S}=\widetilde{S}$，其中 $S$ 是由 $S_1$ 生成的分次
    $A$-代数，而 $S_1$ 在 $A$ 上为有限型。此时只需应用
    \egaref{II.2.9.2-zh}{2.9.2, (ii)}。
\end{enumerate}
\end{proof}

\begin{corollary}[3.6.3]
\label{II.3.6.3-zh}
在 \egaref{II.3.6.2-zh}{3.6.2, (i)} 的条件下，再假设
$\mathcal{S}$ 由 $\mathcal{S}_1$ 生成，则对每个 $n\in\mathbf{Z}$，
$\Phi^*(\mathcal{O}_X(n))$ 典范地等同于
$\mathcal{O}_{X'}(n)$。
\end{corollary}

\begin{proof}
当 $Y$ 为仿射时，已经定义了这样的典范同构
（\egaref{II.2.9.3-zh}{2.9.3}）；在一般情形，只需验证对 $Y$ 的每个
仿射开集 $U$ 如此定义的同构，都与从 $U$ 过渡到仿射开集
$U'\subset U$ 相容，而这是立即的。
\end{proof}

\begin{corollary}[3.6.4]
\label{II.3.6.4-zh}
设 $Y$ 是一个预概形，$\mathcal{S}$ 是由 $\mathcal{S}_1$ 生成的拟凝聚
分次 $\mathcal{O}_Y$-代数层，$\mathcal{M}$ 是一个拟凝聚
$\mathcal{O}_Y$-模层，$u$ 是满的 $\mathcal{O}_Y$-同态
$\mathcal{M}\to\mathcal{S}_1$，而
$\overline{u}:\mathbf{S}_{\mathcal{O}_Y}(\mathcal{M})\to\mathcal{S}$
是延拓 $u$ 的分次 $\mathcal{O}_Y$-代数层同态
（\egaref{II.1.7.4-zh}{1.7.4}）。则与 $\overline{u}$ 对应的态射是
$\operatorname{Proj}(\mathcal{S})$ 到
$\operatorname{Proj}(\mathbf{S}_{\mathcal{O}_Y}(\mathcal{M}))$ 的一个闭浸入。
\end{corollary}

\begin{proof}
事实上，按假设 $\overline{u}$ 为满射，应用
\egaref{II.3.6.1-zh}{3.6.1, (i)} 即可。
\end{proof}
% Locked French authority: ega2-1-fr.tex, 820505 B,
% SHA-256 84EDBE3E83530AF2959B441796337C9DC21EAFCA6A13114A26778760FBF437AC.
% Sequential source scope: lines 5677-5763, subsection II.3.7 through
% environment II.3.7.2 / printed pp.65-66.

\subsection{预概形到齐次谱的态射}
\label{subsection:II.3.7-zh}

\begin{env}[3.7.1]
\label{II.3.7.1-zh}
设 $q:X\to Y$ 是预概形的一个态射，$\mathcal{L}$ 是一个可逆
$\mathcal{O}_X$-模层，$\mathcal{S}$ 是一个正次数拟凝聚分次
$\mathcal{O}_Y$-代数层；于是 $q^*(\mathcal{S})$ 是一个正次数拟凝聚分次
$\mathcal{O}_X$-代数层。考虑拟凝聚分次 $\mathcal{O}_X$-代数层
$\mathcal{S}'=\bigoplus_{n\geq0}\mathcal{L}^{\otimes n}$，并设给定分次代数层的
一个 $\mathcal{O}_X$-同态
\[
  \psi:q^*(\mathcal{S})\longrightarrow
  \mathcal{S}'=\bigoplus_{n\geq0}\mathcal{L}^{\otimes n},
\]
这也等价于给定分次代数层的一个 $q$-态射
\[
  \psi^\flat:\mathcal{S}\longrightarrow q_*(\mathcal{S}').
\]
已知 $\operatorname{Proj}(\mathcal{S}')$ 典范地等同于 $X$
（\egaref{II.3.1.7-zh}{3.1.7} 和
\egaref{II.3.1.8-zh}{3.1.8, (iii)}）；由 $\psi$ 可典范地得到 $X$ 的一个
开集 $G(\psi)$ 和一个 $Y$-态射
\[
\label{II.3.7.1.1-zh}
  r_{\mathcal{L},\psi}:G(\psi)\longrightarrow
  \operatorname{Proj}(\mathcal{S})=P
\tag{3.7.1.1}
\]

\oldpage[II]{66}

称它为\emph{与 $\mathcal{L}$ 和 $\psi$ 相伴的态射}。回顾
\egaref{II.3.5.6-zh}{3.5.6}，这个态射由 $Y$-态射
\[
  \tau=\operatorname{Proj}(\psi):G(\psi)\longrightarrow
  \operatorname{Proj}(q^*(\mathcal{S}))
\]
与第一投影
$\pi:\operatorname{Proj}(q^*(\mathcal{S}))=P\times_YX\to P$ 复合而得。
\end{env}

\begin{env}[3.7.2]
\label{II.3.7.2-zh}
当 $Y=\operatorname{Spec}(A)$ 为仿射时，因而
$\mathcal{S}=\widetilde{S}$，其中 $S$ 是一个正次数分次 $A$-代数；下面具体说明
$r=r_{\mathcal{L},\psi}$。先设 $X=\operatorname{Spec}(B)$ 也为仿射，并且
$\mathcal{L}=\widetilde{L}$，其中 $L$ 是秩为 $1$ 的自由 $B$-模。于是
$q^*(\mathcal{S})=(S\otimes_AB)^{\sim}$
（\egaxref{I.1.6.5-zh}{I, 1.6.5}）。若 $c$ 是 $L$ 的一个生成元，则
$\psi_n:q^*(\mathcal{S}_n)\to\mathcal{L}^{\otimes n}$ 对应于从
$S_n\otimes_AB$ 到 $L^{\otimes n}$ 的同态
$w_n:s\otimes b\mapsto bv_n(s)c^{\otimes n}$，其中
$v_n:S_n\to B$ 是 $A$-模同态，而各 $v_n$ 组成一个代数同态 $S\to B$。
设 $f\in S_d$（$d>0$），并置 $g=v_d(f)$；由
\egaref{II.2.8.10-zh}{2.8.10} 以及 $D_+(f)$ 与
$\operatorname{Spec}(S_{(f)})$ 的等同
（\egaref{II.2.3.6-zh}{2.3.6}），有
$\pi^{-1}(D_+(f))=D_+(f\otimes1)$。另一方面，公式
\egaref{II.2.8.1.1-zh}{2.8.1.1} 表明（计及 $X$ 与
$\operatorname{Proj}(\mathcal{S}')$ 的典范等同）
\[
  \tau^{-1}(D_+(f\otimes1))=D(g),
\]
因而
\[
\label{II.3.7.2.1-zh}
  r^{-1}(D_+(f))=D(g).
\tag{3.7.2.1}
\]

此外，$\tau=\operatorname{Proj}(\psi)$ 限制到 $D(g)$ 后，对应于把
$(s\otimes1)/(f\otimes1)^n$（其中 $s\in S_{nd}$）映到
$v_{nd}(s)/g^n$ 的同态（\egaref{II.2.8.1-zh}{2.8.1}）；而投影 $\pi$
限制到 $D_+(f\otimes1)$ 后，对应于同态
$s/f^n\mapsto(s\otimes1)/(f\otimes1)^n$。由此可知，$r$ 限制到
$D(g)$ 后，对应于 $A$-代数同态 $\omega:S_{(f)}\to B_g$，它对
$s\in S_{nd}$（$n>0$）满足
$\omega(s/f^n)=v_{nd}(s)/g^n$。转到 $X$ 任意的情形（$Y$ 仍为仿射），
并计及 \egaref{II.2.8.1-zh}{2.8.1}，便得到：
\end{env}
% Locked French authority: ega2-1-fr.tex, 820505 B,
% SHA-256 84EDBE3E83530AF2959B441796337C9DC21EAFCA6A13114A26778760FBF437AC.
% Sequential source scope: lines 5765-5837, proposition II.3.7.3 through
% corollary II.3.7.5 and proof / printed pp.66-67.

\begin{proposition}[3.7.3]
\label{II.3.7.3-zh}
若 $Y=\operatorname{Spec}(A)$ 为仿射，且
$\mathcal{S}=\widetilde{S}$，其中 $S$ 是一个分次 $A$-代数，则对每个
$f\in S_d=\Gamma(Y,\mathcal{S}_d)$，都有
\[
\label{II.3.7.3.1-zh}
  r_{\mathcal{L},\psi}^{-1}(D_+(f))=X_{\psi^\flat(f)}
  \qquad\bigl(\text{其中 }\psi^\flat(f)\in
  \Gamma(X,\mathcal{L}^{\otimes d})\bigr)
\tag{3.7.3.1}
\]
并且 $r_{\mathcal{L},\psi}$ 到
$X_{\psi^\flat(f)}\to D_+(f)=\operatorname{Spec}(S_{(f)})$ 的限制，在
\egaxref{I.2.2.4-zh}{I, 2.2.4} 的意义下对应于代数同态
\[
\label{II.3.7.3.2-zh}
  \psi_{(f)}^\flat:S_{(f)}\longrightarrow
  \Gamma(X_{\psi^\flat(f)},\mathcal{O}_X)
\tag{3.7.3.2}
\]
使得对 $s\in S_{nd}=\Gamma(Y,\mathcal{S}_{nd})$，有
\[
\label{II.3.7.3.3-zh}
  \psi_{(f)}^\flat(s/f^n)
  =(\psi^\flat(s)|X_{\psi^\flat(f)})
   (\psi^\flat(f)|X_{\psi^\flat(f)})^{-n}.
\tag{3.7.3.3}
\]
\end{proposition}

若 $G(\psi)=X$，就称 $r_{\mathcal{L},\psi}$ \emph{处处有定义}。显然，
这等价于对每个仿射开集 $U\subset Y$ 都有
$G(\psi)\cap q^{-1}(U)=q^{-1}(U)$；换言之，这个问题\emph{在 $Y$ 上是
局部的}。若 $Y$ 为仿射，则由 \egaref{II.2.8.1-zh}{2.8.1}，
$G(\psi)$ 是所有 $r^{-1}(D_+(f))$ 的并，其中 $f$ 是 $S_+$ 中的齐次元；
因此，由 \egaref{II.3.7.3.1-zh}{3.7.3.1}，各
$X_{\psi^\flat(f)}$ 必须构成 $X$ 的一个\emph{覆盖}。换言之：

\begin{corollary}[3.7.4]
\label{II.3.7.4-zh}
在 \egaref{II.3.7.3-zh}{3.7.3} 的假设下，$r_{\mathcal{L},\psi}$ 处处有定义
的充分必要条件是：对每个 $x\in X$，存在一个整数 $n>0$ 以及
$\mathcal{S}_n$ 在 $Y$ 上的一个截面 $s$，使得置
$t=\psi^\flat(s)\in\Gamma(X,\mathcal{L}^{\otimes n})$ 时，有
$t(x)\neq0$。
\end{corollary}

\oldpage[II]{67}

要注意，若 $\psi$ 是 (\textbf{TN})-满射，这个条件总是成立。

同样，$r_{\mathcal{L},\psi}$ 是否\emph{支配}的问题在 $Y$ 上是局部的，
而且有：

\begin{corollary}[3.7.5]
\label{II.3.7.5-zh}
在 \egaref{II.3.7.3-zh}{3.7.3} 的假设下，$r_{\mathcal{L},\psi}$ 支配的
充分必要条件是：对每个整数 $n>0$，若截面 $s\in S_n$ 使得
$\psi^\flat(s)\in\Gamma(X,\mathcal{L}^{\otimes n})$ 局部幂零，则该截面
本身幂零。
\end{corollary}

\begin{proof}
事实上，所要表达的是：只有当 $D_+(s)$ 为空时，
$r_{\mathcal{L},\psi}^{-1}(D_+(s))$ 才为空。因此，本推论由
\egaref{II.3.7.3.1-zh}{3.7.3.1} 和
\egaref{II.2.3.7-zh}{2.3.7} 得出。
\end{proof}
% Locked French authority: ega2-1-fr.tex, 820505 B,
% SHA-256 84EDBE3E83530AF2959B441796337C9DC21EAFCA6A13114A26778760FBF437AC.
% Sequential source scope: lines 5839-5910, propositions II.3.7.6-3.7.7
% and proofs / printed p.67.

\begin{proposition}[3.7.6]
\label{II.3.7.6-zh}
设 $q:X\to Y$ 是一个态射，$\mathcal{L}$ 是一个可逆
$\mathcal{O}_X$-模层，$\mathcal{S}$、$\mathcal{S}'$ 是两个拟凝聚分次
$\mathcal{O}_Y$-代数层，
$u:\mathcal{S}'\to\mathcal{S}$ 是分次代数层同态，
$\psi:q^*(\mathcal{S})\to\bigoplus_{n\geq0}\mathcal{L}^{\otimes n}$ 是分次
代数层同态，并设 $\psi'=\psi\circ q^*(u)$ 为其复合同态。若
$r_{\mathcal{L},\psi'}$ 处处有定义，则 $r_{\mathcal{L},\psi}$ 也处处有定义；
若 $u$ 是 (\textbf{TN})-满射且 $r_{\mathcal{L},\psi'}$ 是支配的，则
$r_{\mathcal{L},\psi}$ 也是支配的；反之，若 $u$ 是
(\textbf{TN})-单射且 $r_{\mathcal{L},\psi}$ 是支配的，则
$r_{\mathcal{L},\psi'}$ 是支配的。
\end{proposition}

\begin{proof}
事实上，由 \egaref{II.2.8.4-zh}{2.8.4} 有
$G(\psi')\subset G(\psi)$，第一条断言即由此得出。若 $u$ 是
(\textbf{TN})-满射，则
$\operatorname{Proj}(u):\operatorname{Proj}(\mathcal{S})\to
\operatorname{Proj}(\mathcal{S}')$ 处处有定义并且是闭浸入。由于
$r_{\mathcal{L},\psi'}$ 是 $\operatorname{Proj}(u)$ 与
$r_{\mathcal{L},\psi}$ 到 $G(\psi')$ 的限制之复合，便知若
$r_{\mathcal{L},\psi'}$ 是支配的，则 $r_{\mathcal{L},\psi}$ 也是支配的。
最后，若 $u$ 是 (\textbf{TN})-单射，则由
\egaref{II.2.8.3-zh}{2.8.3}，已知 $\operatorname{Proj}(u)$ 是从 $G(u)$
到 $\operatorname{Proj}(\mathcal{S}')$ 的支配态射。由于 $G(\psi')$ 是
$G(u)$ 在 $r_{\mathcal{L},\psi}$ 下的逆像，便看出若
$r_{\mathcal{L},\psi}$ 是支配的，则 $r_{\mathcal{L},\psi'}$ 也是支配的。
\end{proof}

\begin{proposition}[3.7.7]
\label{II.3.7.7-zh}
设 $Y$ 是一个拟紧预概形，$q:X\to Y$ 是一个拟紧态射，$\mathcal{L}$ 是一个
可逆 $\mathcal{O}_X$-模层，而 $\mathcal{S}$ 是拟凝聚分次
$\mathcal{O}_Y$-代数层，它是一个拟凝聚 $\mathcal{O}_Y$-代数层归纳系
$(\mathcal{S}^\lambda)$ 的滤过归纳极限。设
$\varphi_\lambda:\mathcal{S}^\lambda\to\mathcal{S}$ 为典范同态，
$\psi:q^*(\mathcal{S})\to\bigoplus_{n\geq0}\mathcal{L}^{\otimes n}$ 为分次
代数层同态，并置 $\psi_\lambda=\psi\circ q^*(\varphi_\lambda)$。
$r_{\mathcal{L},\psi}$ 处处有定义的充分必要条件是：存在一个 $\lambda$，
使 $r_{\mathcal{L},\psi_\lambda}$ 处处有定义；此时，对
$\mu\geq\lambda$，$r_{\mathcal{L},\psi_\mu}$ 也处处有定义。
\end{proposition}

\begin{proof}
由 \egaref{II.3.7.6-zh}{3.7.6}，该条件是充分的。反之，设
$r_{\mathcal{L},\psi}$ 处处有定义。可以归结到 $Y$ 为仿射的情形；
事实上，若对每个仿射开集 $U\subset Y$，存在 $\lambda(U)$，使
$r_{\mathcal{L},\psi_{\lambda(U)}}$ 到 $q^{-1}(U)$ 的限制处处有定义，则因
$Y$ 拟紧，只需用有限多个仿射开集 $U_i$ 覆盖 $Y$，并由
\egaref{II.3.7.6-zh}{3.7.6} 取一个对所有指标 $i$ 都满足
$\lambda\geq\lambda(U_i)$ 的指标。若 $Y$ 为仿射，假设蕴含：对每个
$x\in X$，存在 $S_n$ 的一个截面 $s^{(x)}$，使得置
$t^{(x)}=\psi^\flat(s^{(x)})$ 时有 $t^{(x)}(x)\neq0$（这里把
$t^{(x)}$ 看成 $\mathcal{L}^{\otimes n}$ 在 $X$ 上的截面）。因此，对 $x$
的某个邻域 $V(x)$ 中的每个 $z$，都有 $t^{(x)}(z)\neq0$。用有限多个
$V(x_i)$ 覆盖 $X$，并设 $s^{(i)}$ 为 $S$ 的相应截面；于是存在一个
$\lambda$，使所有 $s^{(i)}$ 都形如 $\varphi_\lambda(s_\lambda^{(i)})$，
其中对所有 $i$ 都有 $s_\lambda^{(i)}\in S^\lambda$。由
\egaref{II.3.7.4-zh}{3.7.4} 可知 $r_{\mathcal{L},\psi_\lambda}$ 处处有定义。
最后一条断言是 \egaref{II.3.7.6-zh}{3.7.6} 的直接推论。
\end{proof}
% Locked French authority: ega2-1-fr.tex, 820505 B,
% SHA-256 84EDBE3E83530AF2959B441796337C9DC21EAFCA6A13114A26778760FBF437AC.
% Sequential source scope: lines 5912-5931, corollary II.3.7.8 and proof /
% printed pp.67-68.

\begin{corollary}[3.7.8]
\label{II.3.7.8-zh}
在 \egaref{II.3.7.7-zh}{3.7.7} 的假设下，若各
$r_{\mathcal{L},\psi_\lambda}$ 都是支配的，则 $r_{\mathcal{L},\psi}$ 也是
支配的；当各 $\varphi_\lambda$ 都是单射时，逆命题成立。
\end{corollary}

\oldpage[II]{68}

\begin{proof}
第二条断言是 \egaref{II.3.7.6-zh}{3.7.6} 的一个特例。另一方面，为证明
$r_{\mathcal{L},\psi}$ 是支配的，可以只考虑 $Y$ 为仿射的情形。若
$s\in S$ 使 $\psi^\flat(s)$ 局部幂零，由于至少对一个 $\lambda$ 可写成
$s=\varphi_\lambda(s_\lambda)$，便由假设和
\egaref{II.3.7.5-zh}{3.7.5} 得出 $s_\lambda$ 幂零，因而 $s$ 也幂零；于是
\egaref{II.3.7.5-zh}{3.7.5} 的判别准则适用。
\end{proof}
% Locked French authority: ega2-1-fr.tex, 820505 B,
% SHA-256 84EDBE3E83530AF2959B441796337C9DC21EAFCA6A13114A26778760FBF437AC.
% Sequential source scope: lines 5933-6023, remarks II.3.7.9 /
% printed pp.68-69.

\begin{namedtheorem}[3.7.9]{注}
\label{II.3.7.9-zh}
\begin{enumerate}
  \item[(i)] 沿用 \egaref{II.3.7.1-zh}{3.7.1} 的记号，并计及
    \egaref{II.3.2.10-zh}{3.2.10}，对每个 $n\in\mathbf{Z}$ 都有典范同态
    \[
    \label{II.3.7.9.1-zh}
      \theta:r_{\mathcal{L},\psi}^*(\mathcal{O}_P(n))\longrightarrow
      \mathcal{L}^{\otimes n}
    \tag{3.7.9.1}
    \]
    它在 \egaref{II.3.5.6.2-zh}{3.5.6.2} 中以一般方式定义。立即可见，
    在 \egaref{II.3.7.3-zh}{3.7.3} 的条件下，这个同态到
    $X_{\psi^\flat(f)}$ 的限制具体地把 $s/f^k$（$s\in S_{n+kd}$）映到元素
    $(\psi^\flat(s)|X_{\psi^\flat(f)})
    (\psi^\flat(f)|X_{\psi^\flat(f)})^{-k}$。
  \item[(ii)] 设 $\mathcal{F}$ 是一个拟凝聚 $\mathcal{O}_X$-模层，并假设
    $q$ 拟紧且分离，从而对每个 $n\geq0$，
    $q_*(\mathcal{F}\otimes\mathcal{L}^{\otimes n})$ 都是拟凝聚
    $\mathcal{O}_Y$-模层（\egaxref{I.9.2.2-zh}{I, 9.2.2}）。设
    $\mathcal{M}'=\bigoplus_{n\geq0}\mathcal{F}\otimes
    \mathcal{L}^{\otimes n}$，这是一个拟凝聚分次 $\mathcal{S}'$-模层；并考虑
    它的像
    $\mathcal{M}=q_*(\mathcal{M}')=
    \bigoplus_{n\geq0}q_*(\mathcal{F}\otimes\mathcal{L}^{\otimes n})$
    （借助同态 $\psi^\flat$，这是一个拟凝聚 $\mathcal{S}$-模层）。下面说明
    存在一个 $\mathcal{O}_X$-模层的典范同态
    \[
    \label{II.3.7.9.2-zh}
      \xi:r_{\mathcal{L},\psi}^*(\widetilde{\mathcal{M}})
      \longrightarrow\mathcal{F}|G(\psi).
    \tag{3.7.9.2}
    \]

    事实上，在 \egaref{II.3.5.6.1-zh}{3.5.6.1} 中已经定义了典范同态
    \[
    \label{II.3.7.9.3-zh}
      r_{\mathcal{L},\psi}^*(\widetilde{\mathcal{M}})
      \longrightarrow
      (q^*(\mathcal{M})\otimes_{q^*(\mathcal{S})}\mathcal{S}')^{\sim}
      |G(\psi),
    \tag{3.7.9.3}
    \]
    其中把右端看作 $\operatorname{Proj}(\mathcal{S}')$ 上的拟凝聚层。另一方面，
    对每个拟凝聚分次 $\mathcal{S}'$-模层 $\mathcal{M}'$，都有典范同态
    \[
    \label{II.3.7.9.4-zh}
      q^*(q_*(\mathcal{M}'))\otimes_{q^*(\mathcal{S})}\mathcal{S}'
      \longrightarrow\mathcal{M}'
    \tag{3.7.9.4}
    \]
    其定义如下：对 $X$ 的每个开集 $U$、
    $q^*(q_*(\mathcal{M}'_h))$ 在 $U$ 上的每个截面 $t'$，以及
    $\mathcal{S}'_k$ 在 $U$ 上的每个截面 $b'$，把 $t'\otimes b'$ 映到
    $\mathcal{M}'_{h+k}$ 的截面 $b'\sigma(t')$；这里 $\sigma(t')$ 是
    $\mathcal{M}'_h$ 在 $U$ 上的截面，它按
    \egaxref{0.4.4.3-zh}{0, 4.4.3} 典范地对应于 $t'$。由此得到典范同态
    \[
    \label{II.3.7.9.5-zh}
      (q^*(q_*(\mathcal{M}'))\otimes_{q^*(\mathcal{S})}\mathcal{S}')^{\sim}
      |G(\psi)\longrightarrow\widetilde{\mathcal{M}'}|G(\psi)
    \tag{3.7.9.5}
    \]
    最后，由于 $\widetilde{\mathcal{M}'}$ 典范地等同于 $\mathcal{F}$
    （\egaref{II.3.2.9-zh}{3.2.9, (i)}），便确实得到所求的典范同态。

    在 \egaref{II.3.7.3-zh}{3.7.3} 的条件下，这个同态到
    $X_{\psi^\flat(f)}$ 的限制可具体描述如下：给定
    $\mathcal{F}\otimes\mathcal{L}^{\otimes nd}$ 在 $X$ 上的一个截面
    $t_{nd}$，若 $t'_{nd}$ 是把 $t_{nd}$ 看成
    $q_*(\mathcal{F}\otimes\mathcal{L}^{\otimes nd})$ 在 $Y$ 上的截面，
    则元素 $t'_{nd}/f^n$ 对应于 $\mathcal{F}$ 在
    $X_{\psi^\flat(f)}$ 上的截面
    $(t_{nd}|X_{\psi^\flat(f)})
    (\psi^\flat(f)|X_{\psi^\flat(f)})^{-n}$。
\end{enumerate}
\end{namedtheorem}

\oldpage[II]{69}
% Locked French authority: ega2-1-fr.tex, 820505 B,
% SHA-256 84EDBE3E83530AF2959B441796337C9DC21EAFCA6A13114A26778760FBF437AC.
% Sequential source scope: lines 6025-6086, subsection II.3.8 through
% proposition II.3.8.2 and proof / printed p.69.
% EGA2-ZH-DEF-0007 candidate: retain source-literal U_alpha in the proof
% until the guardian disposes the undefined-symbol reading.

\subsection{浸入齐次谱的判别法}
\label{subsection:II.3.8-zh}

\begin{env}[3.8.1]
\label{II.3.8.1-zh}
沿用 \egaref{II.3.7.1-zh}{3.7.1} 的记号，$r_{\mathcal{L},\psi}$ 是否为浸入
（相应地，开浸入、闭浸入）的问题显然\emph{在 $Y$ 上是局部的}。
\end{env}

\begin{proposition}[3.8.2]
\label{II.3.8.2-zh}
在 \egaref{II.3.7.3-zh}{3.7.3} 的假设下，$r_{\mathcal{L},\psi}$ 处处有定义
且为浸入的充分必要条件是：存在一族截面 $s_\alpha\in S_{n_\alpha}$
（$n_\alpha>0$），使得置 $f_\alpha=\psi^\flat(s_\alpha)$ 时满足下列条件：
\begin{enumerate}
  \item[(i)] 各 $X_{f_\alpha}$ 构成 $X$ 的一个覆盖。
  \item[(ii)] 各 $X_{f_\alpha}$ 都是仿射开集。
  \item[(iii)] 对每个 $\alpha$ 以及每个
    $t\in\Gamma(X_{f_\alpha},\mathcal{O}_X)$，存在整数 $m>0$ 和
    $s\in S_{mn_\alpha}$，使得
    \par\noindent\hfil
    $t=(\psi^\flat(s)|X_{f_\alpha})(f_\alpha|X_{f_\alpha})^{-m}$。
    \hfil\par
\end{enumerate}

$r_{\mathcal{L},\psi}$ 处处有定义且为开浸入的充分必要条件是：存在一族
$(s_\alpha)$ 满足 (i)、(ii)、(iii) 以及下列条件：
\begin{enumerate}
  \item[(iv)] 对每个 $n>0$ 和每个满足
    $\psi^\flat(s)|X_{f_\alpha}=0$ 的 $s\in S_{nn_\alpha}$，存在整数 $k>0$，
    使 $s_\alpha^ks=0$。
\end{enumerate}

$r_{\mathcal{L},\psi}$ 处处有定义且为闭浸入的充分必要条件是：存在一族
$(s_\alpha)$ 满足 (i)、(ii)、(iii) 以及下列条件：
\begin{enumerate}
  \item[(v)] 各 $D_+(s_\alpha)$ 构成
    $P=\operatorname{Proj}(S)$ 的一个覆盖。
\end{enumerate}
\end{proposition}

\begin{proof}
$r$ 为浸入（相应地，闭浸入）的充分必要条件是：存在由各
$D_+(s_\alpha)$ 组成的 $r(G(\psi))$（相应地，$P$）的一个覆盖，使得若置
$V_\alpha=r^{-1}(D_+(s_\alpha))$，则 $r$ 到 $U_\alpha$ 的限制是从
$V_\alpha$ 到 $D_+(s_\alpha)$ 的闭浸入
（\egaxref{I.4.2.4-zh}{I, 4.2.4}）。由
\egaref{II.3.7.3.1-zh}{3.7.3.1}，条件 (i) 同时表示 $r$ 处处有定义，
并且各 $D_+(s_\alpha)$ 覆盖 $r(X)$。由于 $D_+(s_\alpha)$ 为仿射，
由 \egaxref{I.4.2.3-zh}{I, 4.2.3}，条件 (ii) 和 (iii) 表示 $r$ 到
$X_{f_\alpha}$ 的限制是到 $D_+(s_\alpha)$ 的闭浸入。最后，由于 (iii) 和
(iv) 表示 $\psi_{(s_\alpha)}^\flat$ 是双射
（采用 \egaref{II.3.7.3.2-zh}{3.7.3.2} 的记号），(ii)、(iii) 和 (iv)
表示对每个 $\alpha$，$r$ 到 $X_{f_\alpha}$ 的限制都是到
$D_+(s_\alpha)$ 上的同构；因此，(i)、(ii)、(iii) 和 (iv) 表示 $r$
是开浸入。
\end{proof}
% Locked French authority: ega2-1-fr.tex, 820505 B,
% SHA-256 84EDBE3E83530AF2959B441796337C9DC21EAFCA6A13114A26778760FBF437AC.
% Sequential source scope: lines 6088-6149, corollary II.3.8.3 and
% proposition II.3.8.4 with proofs / printed pp.69-70.

\begin{corollary}[3.8.3]
\label{II.3.8.3-zh}
在 \egaref{II.3.7.6-zh}{3.7.6} 的假设下，若
$r_{\mathcal{L},\psi'}$ 处处有定义且为浸入，则
$r_{\mathcal{L},\psi}$ 也如此。若还假设 $u$ 是 (\textbf{TN})-满射，并且
$r_{\mathcal{L},\psi'}$ 是开浸入（相应地，闭浸入），则
$r_{\mathcal{L},\psi}$ 也如此。
\end{corollary}

\begin{proof}
事实上，由 \egaref{II.3.8.2-zh}{3.8.2}，存在一族
$s'_\alpha\in S'_{n_\alpha}$，使得置
$f_\alpha=\psi'^\flat(s'_\alpha)$ 时满足条件 (i)、(ii)、(iii)。若置
$s_\alpha=u(s'_\alpha)$，也有 $f_\alpha=\psi^\flat(s_\alpha)$；而若
$t=(\psi'^\flat(s')|X_{f_\alpha})(f_\alpha|X_{f_\alpha})^{-m}$，则置
$s=u(s')$ 后，也有
$t=(\psi^\flat(s)|X_{f_\alpha})(f_\alpha|X_{f_\alpha})^{-m}$。第一条断言
由此得出。第二条断言立即来自 $\operatorname{Proj}(u)$ 是闭浸入这一事实。
\end{proof}

\begin{proposition}[3.8.4]
\label{II.3.8.4-zh}
假设 \egaref{II.3.7.7-zh}{3.7.7} 的条件成立，并且
$q:X\to Y$ 还是有限型态射。则 $r_{\mathcal{L},\psi}$ 处处有定义且为浸入
的充分必要条件是：存在 $\lambda$，使
$r_{\mathcal{L},\psi_\lambda}$ 处处有定义且为浸入；在这种情形，对
$\mu\geq\lambda$，$r_{\mathcal{L},\psi_\mu}$ 都处处有定义且为浸入。
\end{proposition}

\oldpage[II]{70}

\begin{proof}
计及 \egaref{II.3.8.3-zh}{3.8.3}，只需证明：若
$r_{\mathcal{L},\psi}$ 处处有定义且为浸入，则至少对一个 $\lambda$，
$r_{\mathcal{L},\psi_\lambda}$ 也如此。像
\egaref{II.3.7.7-zh}{3.7.7} 中那样论证，并使用 $Y$ 的拟紧性，首先归结到
$Y$ 为仿射的情形。由于 $X$ 拟紧，\egaref{II.3.8.2-zh}{3.8.2} 表明存在
$S$ 中有限一族 $(s_i)$（$s_i\in S_{n_i}$），满足条件 (i)、(ii)、(iii)。
态射 $X_{f_i}\to Y$（其中 $f_i=\psi^\flat(s_i)$）为有限型；事实上，这是
仿射概形之间的态射，因而拟紧（\egaxref{I.6.6.1-zh}{I, 6.6.1}），又因
$q$ 为有限型而局部有限型（\egaxref{I.6.3.2-zh}{I, 6.3.2}），结论来自
\egaxref{I.6.6.3-zh}{I, 6.6.3}。因此，$X_{f_i}$ 的环 $B_i$ 是一个有限型
$A$-代数（\egaxref{I.6.3.3-zh}{I, 6.3.3}）；设 $(t_{ij})$ 是这个代数的
一族生成元。按假设，存在元素 $s'_{ij}\in S_{m_{ij}n_i}$，使得
\[
  t_{ij}=(\psi^\flat(s'_{ij})|X_{f_i})
  (\psi^\flat(s_i)|X_{f_i})^{-m_{ij}}.
\]
按假设还存在一个 $\lambda$ 以及元素
$s_{i\lambda}\in S_{n_i}^\lambda$、
$s'_{ij\lambda}\in S_{m_{ij}n_i}^\lambda$，它们在
$\varphi_\lambda$ 下的像分别为 $s_i$ 和 $s'_{ij}$。显然，
$r_{\mathcal{L},\psi_\lambda}$ 满足
\egaref{II.3.8.2-zh}{3.8.2} 的条件 (i)、(ii)、(iii)。
\end{proof}
% Locked French authority: ega2-1-fr.tex, 820505 B,
% SHA-256 84EDBE3E83530AF2959B441796337C9DC21EAFCA6A13114A26778760FBF437AC.
% Sequential source scope: lines 6151-6178, proposition II.3.8.5 /
% printed p.70.

\begin{proposition}[3.8.5]
\label{II.3.8.5-zh}
设 $Y$ 是一个拟紧概形，或者是底空间为诺特的一个预概形；
$q:X\to Y$ 是一个有限型态射，$\mathcal{L}$ 是一个可逆
$\mathcal{O}_X$-模层，$\mathcal{S}$ 是一个拟凝聚分次
$\mathcal{O}_Y$-代数层，
$\psi:q^*(\mathcal{S})\to\bigoplus_{n\geq0}\mathcal{L}^{\otimes n}$ 是分次
代数层同态。$r_{\mathcal{L},\psi}$ 处处有定义且为浸入的充分必要条件是：
存在一个整数 $n>0$，以及 $\mathcal{S}_n$ 的一个有限型拟凝聚
$\mathcal{O}_Y$-子模层 $\mathcal{E}$，使得：
\begin{enumerate}
  \item[a)] 同态
    $\psi_n\circ q^*(j_n):q^*(\mathcal{E})\to\mathcal{L}^{\otimes n}$
    是满射（其中 $j_n:\mathcal{E}\to\mathcal{S}$ 是典范单射）；
  \item[b)] 以 $\mathcal{S}'$ 表示 $\mathcal{S}$ 中由 $\mathcal{E}$ 生成的
    （分次）$\mathcal{O}_Y$-子代数层，以 $\psi'$ 表示同态
    $\psi\circ q^*(j')$（$j'$ 是从 $\mathcal{S}'$ 到 $\mathcal{S}$ 的
    单射），则 $r_{\mathcal{L},\psi'}$ 处处有定义且为浸入。
\end{enumerate}

当这些条件成立时，$\mathcal{S}_n$ 中每个包含 $\mathcal{E}$ 的拟凝聚
$\mathcal{O}_Y$-子模层 $\mathcal{E}'$ 都具有同一性质；对每个 $k>0$，
$\bigotimes^k\mathcal{E}$ 在 $\mathcal{S}_{kn}$ 中的像也具有同一性质。
\end{proposition}
% Locked French authority: ega2-1-fr.tex, 820505 B,
% SHA-256 84EDBE3E83530AF2959B441796337C9DC21EAFCA6A13114A26778760FBF437AC.
% Sequential source scope: lines 6180-6221, proof of II.3.8.5 through
% the heading of section II.4 / printed pp.70-71.
% EGA2-ZH-DEF-0008 candidate: retain source-literal f|V in the proof
% until the guardian disposes the undefined-symbol reading.

\begin{proof}
该条件的充分性以及最后两条断言，计及
$\operatorname{Proj}(\mathcal{S})$ 与
$\operatorname{Proj}(\mathcal{S}^{(k)})$ 之间的典范同构
（\egaref{II.3.1.8-zh}{3.1.8}），都是
\egaref{II.3.8.3-zh}{3.8.3} 的特例。

设 $(U_i)$ 是由仿射开集组成的 $Y$ 的一个有限开覆盖，并置
$A_i=A(U_i)$。由于 $q^{-1}(U_i)$ 拟紧，$r_{\mathcal{L},\psi}$ 是处处
有定义的浸入这一假设以及 \egaref{II.3.8.2-zh}{3.8.2} 蕴含：存在
$S^{(i)}=\Gamma(U_i,\mathcal{S})$ 中有限一族元素 $(s_{ij})$
（$s_{ij}\in S_{n_{ij}}^{(i)}$），满足条件 (i)、(ii)、(iii)。像在
\egaref{II.3.8.4-zh}{3.8.4} 的证明中那样可见，态射
$X_{f_{ij}}\to U_i$（其中 $f_{ij}=\psi^\flat(s_{ij})$）是有限型的；因此，
$X_{f_{ij}}$ 的环 $B_{ij}$ 是有限型 $A_i$-代数，且有一组形如
$(\psi^\flat(t_{ijk})|X_{f_{ij}})(f_{ij}|X_{f_{ij}})^{-m_{ijk}}$ 的生成元，
其中 $t_{ijk}\in S_{m_{ijk}n_{ij}}^{(i)}$。设 $n$ 是各
$m_{ijk}n_{ij}$ 的一个公倍数；对每个 $(i,j,k)$，以满足
$\rho m_{ijk}n_{ij}=n$ 的幂 $s_{ij}^\rho$ 代替 $s_{ij}$，并以
$s_{ij}^{\rho-m_{ijk}}$ 乘 $t_{ijk}$，便可假设对每个 $i$，所有
$s_{ij}$ 和 $t_{ijk}$ 都属于 $S_n^{(i)}$，并且 $m_{ijk}=1$。设 $E_i$
是这些元素（固定 $i$）生成的 $S_n^{(i)}$ 的 $A_i$-子模。存在
$\mathcal{S}_n$ 的一个有限型凝聚 $\mathcal{O}_Y$-子模层

\oldpage[II]{71}

$\mathcal{E}_i$，使得 $\mathcal{E}_i|U_i=(E_i)^{\sim}$
（\egaxref{I.9.4.7-zh}{I, 9.4.7}）。显然，$\mathcal{S}_n$ 的
$\mathcal{O}_Y$-子模层 $\mathcal{E}$，即各 $\mathcal{E}_i$ 之和，正是
所求对象；这是因为每个截面 $f_{ij}$ 都具有如下性质：对每个
$x\in X_{f_{ij}}$，存在 $x$ 的一个仿射邻域 $V\subset X_{f_{ij}}$，
使 $f|V$ 成为 $\Gamma(V,\mathcal{L}^{\otimes n})$ 的一个基。

\end{proof}

\section{泛射影丛、丰沛层}
\label{section:II.4-zh}
% Locked French authority: ega2-1-fr.tex, 820505 B,
% SHA-256 84EDBE3E83530AF2959B441796337C9DC21EAFCA6A13114A26778760FBF437AC.
% Sequential source scope: lines 6223-6313, subsection II.4.1 through
% environment II.4.1.3 / printed pp.71-72.

\subsection{泛射影丛的定义}
\label{subsection:II.4.1-zh}

\begin{definition}[4.1.1]
\label{II.4.1.1-zh}
设 $Y$ 是一个预概形，$\mathcal{E}$ 是一个拟凝聚 $\mathcal{O}_Y$-模层，
$\mathbf{S}_{\mathcal{O}_Y}(\mathcal{E})$ 是 $\mathcal{E}$ 的
$\mathcal{O}_Y$-对称代数层（\egaref{II.1.7.4-zh}{1.7.4}），它是拟凝聚的
（\egaref{II.1.7.7-zh}{1.7.7}）。把 $Y$-概形
$P=\operatorname{Proj}(\mathbf{S}_{\mathcal{O}_Y}(\mathcal{E}))$ 称为
\emph{由 $\mathcal{E}$ 定义的 $Y$ 上的泛射影丛}，并记作
$\mathbf{P}(\mathcal{E})$。把 $\mathcal{O}_P$-模层 $\mathcal{O}_P(1)$ 称为
\emph{$P$ 上的基本层}。
\end{definition}

当 $Y$ 是环 $A$ 的仿射概形时，$\mathcal{E}=\widetilde{E}$，其中 $E$ 是一个
$A$-模；此时写作 $\mathbf{P}(E)$，而不写作
$\mathbf{P}(\widetilde{E})$。

当取 $\mathcal{E}=\mathcal{O}_Y^n$ 时，写作 $\mathbf{P}_Y^{n-1}$，而不写作
$\mathbf{P}(\mathcal{E})$；若 $Y$ 还是环 $A$ 的仿射概形，也写作
$\mathbf{P}_A^{n-1}$，而不写作 $\mathbf{P}_Y^{n-1}$。由于
$\mathbf{S}_{\mathcal{O}_Y}(\mathcal{O}_Y)$ 典范地等同于
$\mathcal{O}_Y[T]$（\egaref{II.1.7.4-zh}{1.7.4}），
$\mathbf{P}_Y^0$ 典范地等同于 $Y$（\egaref{II.3.1.7-zh}{3.1.7}）；
例 \egaref{II.2.4.3-zh}{2.4.3} 正是 $\mathbf{P}_K^1$。

\begin{env}[4.1.2]
\label{II.4.1.2-zh}
设 $\mathcal{E}$、$\mathcal{F}$ 是两个拟凝聚 $\mathcal{O}_Y$-模层，
$u:\mathcal{E}\to\mathcal{F}$ 是一个 $\mathcal{O}_Y$-同态；它典范地对应于
分次 $\mathcal{O}_Y$-代数层同态
$\mathbf{S}(u):\mathbf{S}_{\mathcal{O}_Y}(\mathcal{E})\to
\mathbf{S}_{\mathcal{O}_Y}(\mathcal{F})$
（\egaref{II.1.7.4-zh}{1.7.4}）。若 $u$ \emph{为满射}，则
$\mathbf{S}(u)$ 也为满射，因而由
\egaref{II.3.6.2-zh}{3.6.2, (i)}，
$\operatorname{Proj}(\mathbf{S}(u))$ 是闭浸入
$\mathbf{P}(\mathcal{F})\to\mathbf{P}(\mathcal{E})$，把它记作
$\mathbf{P}(u)$。所以可以说，$\mathbf{P}(\mathcal{E})$ 是
\emph{$\mathcal{E}$ 的反变函子}，但作为拟凝聚 $\mathcal{O}_Y$-模层的态射
只取满同态。

仍设 $u$ 为满射，并置 $P=\mathbf{P}(\mathcal{E})$、
$Q=\mathbf{P}(\mathcal{F})$ 和 $j=\mathbf{P}(u)$，则在相差一个同构的意义下，
\[
\label{II.4.1.2.1-zh}
  j^*(\mathcal{O}_P(n))=\mathcal{O}_Q(n)
  \qquad\text{对每个 }n\in\mathbf{Z}
\tag{4.1.2.1}
\]
这是 \egaref{II.3.6.3-zh}{3.6.3} 的推论。
\end{env}

\begin{env}[4.1.3]
\label{II.4.1.3-zh}
现在设 $\psi:Y'\to Y$ 是一个态射，并设
$\mathcal{E}'=\psi^*(\mathcal{E})$；于是
$\mathbf{S}_{\mathcal{O}_{Y'}}(\mathcal{E}')=
\psi^*(\mathbf{S}_{\mathcal{O}_Y}(\mathcal{E}))$
（\egaref{II.1.7.5-zh}{1.7.5}）。因而由
\egaref{II.3.5.3-zh}{3.5.3}，在相差一个典范同构的意义下，有
\[
\label{II.4.1.3.1-zh}
  \mathbf{P}(\psi^*(\mathcal{E}))=
  \mathbf{P}(\mathcal{E})\times_YY'
\tag{4.1.3.1}
\]
此外，显然有
\[
  \psi^*((\mathbf{S}_{\mathcal{O}_Y}(\mathcal{E}))(n))
  =(\mathbf{S}_{\mathcal{O}_{Y'}}(\mathcal{E}'))(n)
\]
对每个 $n\in\mathbf{Z}$ 都成立。因此，置 $P=\mathbf{P}(\mathcal{E})$、
$P'=\mathbf{P}(\mathcal{E}')$，由 \egaref{II.3.5.4-zh}{3.5.4}，在相差一个
同构的意义下，有
\[
\label{II.4.1.3.2-zh}
  \mathcal{O}_{P'}(n)=\mathcal{O}_P(n)\otimes_Y\mathcal{O}_{Y'}
  \qquad\text{对每个 }n\in\mathbf{Z}.
\tag{4.1.3.2}
\]
\end{env}

\oldpage[II]{72}
% Locked French authority: ega2-1-fr.tex, 820505 B,
% SHA-256 84EDBE3E83530AF2959B441796337C9DC21EAFCA6A13114A26778760FBF437AC.
% Sequential source scope: lines 6315-6363, proposition II.4.1.4 and proof /
% printed p.72.

\begin{proposition}[4.1.4]
\label{II.4.1.4-zh}
设 $\mathcal{L}$ 是一个可逆 $\mathcal{O}_Y$-模层。对每个拟凝聚
$\mathcal{O}_Y$-模层 $\mathcal{E}$，存在典范 $Y$-同构
$i_{\mathcal{L}}:\mathbf{P}(\mathcal{E})\xrightarrow{\sim}
\mathbf{P}(\mathcal{E}\otimes\mathcal{L})$。此外，若置
$P=\mathbf{P}(\mathcal{E})$、$Q=\mathbf{P}(\mathcal{E}\otimes\mathcal{L})$，
则对每个 $n\in\mathbf{Z}$，
$i_{\mathcal{L}}^*(\mathcal{O}_Q(n))$ 典范同构于
$\mathcal{O}_P(n)\otimes_Y\mathcal{L}^{\otimes n}$。
\end{proposition}

\begin{proof}
先注意，若 $A$ 是一个环，$E$ 是一个 $A$-模，而 $L$ 是一个
\emph{单生成自由} $A$-模，则可典范地定义 $A$-模同态
\[
  \mathbf{S}_n(E\otimes L)\longrightarrow
  \mathbf{S}_n(E)\otimes L^{\otimes n}
\]
它把 $(x_1\otimes y_1)\cdots(x_n\otimes y_n)$ 映到元素
\[
  (x_1x_2\cdots x_n)\otimes(y_1\otimes y_2\otimes\cdots\otimes y_n)
  \qquad(x_i\in E,\ y_i\in L\text{，其中 }1\leq i\leq n)；
\]
归结到 $L=A$ 的情形，立即验证这个同态实际上是同构。由此得到分次
$A$-代数的典范同构
$\mathbf{S}_A(E\otimes L)\xrightarrow{\sim}
\bigoplus_{n\geq0}\mathbf{S}_n(E)\otimes L^{\otimes n}$。回到 (4.1.4) 的
条件，前述考察使我们可以定义分次 $\mathcal{O}_Y$-代数层的典范同构
\[
\label{II.4.1.4.1-zh}
  \mathbf{S}_{\mathcal{O}_Y}(\mathcal{E}\otimes_{\mathcal{O}_Y}\mathcal{L})
  \xrightarrow{\sim}
  \bigoplus_{n\geq0}\mathbf{S}_n(\mathcal{E})
  \otimes_{\mathcal{O}_Y}\mathcal{L}^{\otimes n}
\tag{4.1.4.1}
\]
先把它定义成预层的同构，再计及
\egaref{II.1.7.4-zh}{1.7.4}、
\egaxref{I.1.3.9-zh}{I, 1.3.9} 和
\egaxref{I.1.3.12-zh}{I, 1.3.12}。于是本命题由
\egaref{II.3.1.8-zh}{3.1.8, (iii)} 和
\egaref{II.3.2.10-zh}{3.2.10} 得出。
\end{proof}
% Locked French authority: ega2-1-fr.tex, 820505 B,
% SHA-256 84EDBE3E83530AF2959B441796337C9DC21EAFCA6A13114A26778760FBF437AC.
% Sequential source scope: lines 6365-6399, environment II.4.1.5 through
% proposition II.4.1.6 and proof / printed p.72.

\begin{env}[4.1.5]
\label{II.4.1.5-zh}
在 \egaref{II.4.1.1-zh}{4.1.1} 的假设下，置
$P=\mathbf{P}(\mathcal{E})$，并以 $p$ 表示结构态射 $P\to Y$。按定义有
$\mathcal{E}=(\mathbf{S}_{\mathcal{O}_Y}(\mathcal{E}))_1$，故由
\egaref{II.3.3.2.2-zh}{3.3.2.2} 有典范同态
$\alpha_1:\mathcal{E}\to p_*(\mathcal{O}_P(1))$，进而由
\egaxref{0.4.4.3-zh}{0, 4.4.3} 还有典范同态
\[
\label{II.4.1.5.1-zh}
  \alpha_1^\sharp:p^*(\mathcal{E})\longrightarrow\mathcal{O}_P(1).
\tag{4.1.5.1}
\]
\end{env}

\begin{proposition}[4.1.6]
\label{II.4.1.6-zh}
典范同态 \egaref{II.4.1.5.1-zh}{4.1.5.1} 是满射。
\end{proposition}

\begin{proof}
事实上，在 \egaref{II.3.3.2-zh}{3.3.2} 中已经看到，
$\alpha_1^\sharp$ 函子性地对应于典范同态
\[
  \mathcal{E}\otimes_{\mathcal{O}_Y}
  \mathbf{S}_{\mathcal{O}_Y}(\mathcal{E})
  \longrightarrow
  (\mathbf{S}_{\mathcal{O}_Y}(\mathcal{E}))(1)；
\]
按定义，$\mathcal{E}$ 生成 $\mathbf{S}_{\mathcal{O}_Y}(\mathcal{E})$，
所以这个同态是满射；结论由 \egaref{II.3.2.4-zh}{3.2.4} 得出。
\end{proof}
% Locked French authority: ega2-1-fr.tex, 820505 B,
% SHA-256 84EDBE3E83530AF2959B441796337C9DC21EAFCA6A13114A26778760FBF437AC.
% Sequential source scope: lines 6401-6468, subsection II.4.2 through
% environment II.4.2.1 / printed pp.72-73.

\subsection{预概形到泛射影丛的态射}
\label{subsection:II.4.2-zh}

\begin{env}[4.2.1]
\label{II.4.2.1-zh}
沿用 \egaref{II.4.1.5-zh}{4.1.5} 的记号，设 $X$ 是一个 $Y$-预概形，
$q:X\to Y$ 是结构态射，并设 $r:X\to P$ 是一个 $Y$-\emph{态射}，从而有
交换图
\[
  \xymatrix{
    P \ar[d]_p & X \ar[l]_r \ar[dl]^q
    \\Y
  }
\]

\oldpage[II]{73}

由于函子 $r^*$ 是右正合的（\egaxref{0.4.3.1-zh}{0, 4.3.1}），由
同态 \egaref{II.4.1.5.1-zh}{4.1.5.1}——它由
\egaref{II.4.1.6-zh}{4.1.6} 为满射——得到满同态
\[
  r^*(\alpha_1^\sharp):r^*(p^*(\mathcal{E}))\longrightarrow
  r^*(\mathcal{O}_P(1)).
\]
但是 $r^*(p^*(\mathcal{E}))=q^*(\mathcal{E})$，而
$r^*(\mathcal{O}_P(1))$ 局部同构于
$r^*(\mathcal{O}_P)=\mathcal{O}_X$；换言之，它是 $\mathcal{O}_X$ 上的
一个可逆层 $\mathcal{L}_r$。这样便从 $r$ 定义了典范的满
$\mathcal{O}_X$-同态
\[
\label{II.4.2.1.1-zh}
  \varphi_r:q^*(\mathcal{E})\longrightarrow\mathcal{L}_r.
\tag{4.2.1.1}
\]

当 $Y=\operatorname{Spec}(A)$ 为仿射且
$\mathcal{E}=\widetilde{E}$ 时，还可如下具体描述这个同态。给定
$f\in E$，由 \egaref{II.2.6.3-zh}{2.6.3} 首先有
\[
\label{II.4.2.1.2-zh}
  r^{-1}(D_+(f))=X_{\varphi_r^\flat(f)}.
\tag{4.2.1.2}
\]
另一方面，设 $V$ 是 $X$ 中包含于 $r^{-1}(D_+(f))$ 的一个仿射开集，
$B$ 是它的环，因而为一个 $A$-代数；置 $S=\mathbf{S}_A(E)$。$r$ 到
$V$ 的限制对应于一个 $A$-同态 $\omega:S_{(f)}\to B$，并且
$q^*(\mathcal{E})|V=\widetilde{E\otimes_A B}$、
$\mathcal{L}_r|V=\widetilde{L_r}$，其中
$L_r=(S(1))_{(f)}\otimes_{S_{(f)}}B_{[\omega]}$
（\egaxref{I.1.6.5-zh}{I, 1.6.5}）。$\varphi_r$ 到
$q^*(\mathcal{E})|V$ 的限制对应于 $B$-同态 $u:E\otimes_A B\to L_r$，
它把 $x\otimes1$ 变为
$(x/1)\otimes f=(f/1)\otimes\omega(x/f)$。$\varphi_r$ 到
$\mathcal{O}_X$-代数层同态的典范延拓
\[
  \psi_r:q^*(\mathbf{S}(\mathcal{E}))=
  \mathbf{S}(q^*(\mathcal{E}))\longrightarrow
  \mathbf{S}(\mathcal{L}_r)=\bigoplus_{n\geq0}\mathcal{L}_r^{\otimes n}
\]
因而具有如下性质：$\psi_r$ 到 $q^*(\mathbf{S}_n(\mathcal{E}))|V$ 的限制
对应于同态
$\mathbf{S}_n(E\otimes_A B)=\mathbf{S}_n(E)\otimes_A B\to L_r^{\otimes n}$，
它把 $s\otimes1$ 变为
$(f/1)^{\otimes n}\otimes\omega(s/f^n)$。
\end{env}
% Locked French authority: ega2-1-fr.tex, 820505 B,
% SHA-256 84EDBE3E83530AF2959B441796337C9DC21EAFCA6A13114A26778760FBF437AC.
% Sequential source scope: lines 6470-6537, environment II.4.2.2 through
% proposition II.4.2.3 and proof / printed pp.73-74.

\begin{env}[4.2.2]
\label{II.4.2.2-zh}
反之，设给定一个态射 $q:X\to Y$、一个可逆 $\mathcal{O}_X$-模层
$\mathcal{L}$ 和一个拟凝聚 $\mathcal{O}_Y$-模层 $\mathcal{E}$。同态
$\varphi:q^*(\mathcal{E})\to\mathcal{L}$ 典范地对应于拟凝聚
$\mathcal{O}_X$-代数层同态
\[
  \psi:\mathbf{S}(q^*(\mathcal{E}))=q^*(\mathbf{S}(\mathcal{E}))
  \longrightarrow\bigoplus_{n\geq0}\mathcal{L}^{\otimes n},
\]
因而由 \egaref{II.3.7.1-zh}{3.7.1} 对应于一个 $Y$-态射
$r_{\mathcal{L},\psi}:G(\psi)\to
\operatorname{Proj}(\mathbf{S}(\mathcal{E}))=\mathbf{P}(\mathcal{E})$，把它记作
$r_{\mathcal{L},\varphi}$。若 $\varphi$ 为满射，则 $\psi$ 也为满射，故由
\egaref{II.3.7.4-zh}{3.7.4}，$r_{\mathcal{L},\varphi}$ 处处有定义。此外，
沿用 \egaref{II.4.2.1-zh}{4.2.1} 和 (4.2.2) 的记号：
\end{env}

\begin{proposition}[4.2.3]
\label{II.4.2.3-zh}
给定态射 $q:X\to Y$ 和拟凝聚 $\mathcal{O}_Y$-模层 $\mathcal{E}$，映射
$r\mapsto(\mathcal{L}_r,\varphi_r)$ 与
$(\mathcal{L},\varphi)\mapsto r_{\mathcal{L},\varphi}$ 在如下两个集合之间
建立一一对应：所有 $Y$-态射 $r:X\to\mathbf{P}(\mathcal{E})$ 的集合，以及
所有偶对 $(\mathcal{L},\varphi)$ 的等价类之集合；这里 $\mathcal{L}$ 是一个
可逆 $\mathcal{O}_X$-模层，$\varphi:q^*(\mathcal{E})\to\mathcal{L}$ 是一个
满同态，而两个偶对 $(\mathcal{L},\varphi)$、
$(\mathcal{L}',\varphi')$ 等价，是指存在一个 $\mathcal{O}_X$-同构
$\tau:\mathcal{L}\xrightarrow{\sim}\mathcal{L}'$，使得
$\varphi'=\tau\circ\varphi$。
\end{proposition}

\begin{proof}
先从一个 $Y$-态射 $r:X\to\mathbf{P}(\mathcal{E})$ 出发，构造
$\mathcal{L}_r$ 和 $\varphi_r$（\egaref{II.4.2.1-zh}{4.2.1}），再构造
$r'=r_{\mathcal{L}_r,\varphi_r}$。由
\egaref{II.4.2.1-zh}{4.2.1} 和
\egaref{II.3.7.2-zh}{3.7.2} 立即可知，态射 $r$ 与 $r'$ 相同；这里为定义
\egaref{II.3.7.2-zh}{3.7.2} 中的同态 $v_n$，取 $\mathcal{L}_r$ 的元素
$(f/1)\otimes1$ 作为生成元。反之，从一个偶对
$(\mathcal{L},\varphi)$ 出发并构造

\oldpage[II]{74}

$r''=r_{\mathcal{L},\varphi}$，再构造 $\mathcal{L}_{r''}$ 和
$\varphi_{r''}$。下面证明存在典范同构
$\tau:\mathcal{L}_{r''}\xrightarrow{\sim}\mathcal{L}$，使得
$\varphi=\tau\circ\varphi_{r''}$。为定义它，可以取
$Y=\operatorname{Spec}(A)$ 和 $X=\operatorname{Spec}(B)$ 都为仿射，并沿用
\egaref{II.4.2.1-zh}{4.2.1} 和
\egaref{II.3.7.2-zh}{3.7.2} 的记号：把 $L_{r''}$ 的每个元素
$(x/1)\otimes1$（其中 $x\in E$）映到 $L$ 的元素 $v_1(x)c$。立即验证
$\tau$ 不依赖于 $L$ 所选的生成元 $c$。由于按假设 $v_1$ 为满射，为证明
$\tau$ 是同构，只需证明：若在 $(S(1))_{(f)}$ 中 $x/1=0$，则在 $B_g$
中 $v_1(x)/1=0$。前一关系表示：对某个 $n$，在
$\mathbf{S}_{n+1}(E)$ 中有 $f^nx=0$；由此得到在 $B$ 中
$v_{n+1}(f^nx)=g^nv_1(x)=0$，从而得到结论。最后，显然对两个等价偶对
$(\mathcal{L},\varphi)$、$(\mathcal{L}',\varphi')$，有
$r_{\mathcal{L},\varphi}=r_{\mathcal{L}',\varphi'}$。
\end{proof}
% Locked French authority: ega2-1-fr.tex, 820505 B,
% SHA-256 84EDBE3E83530AF2959B441796337C9DC21EAFCA6A13114A26778760FBF437AC.
% Sequential source scope: lines 6539-6593, theorem II.4.2.4 through
% remark II.4.2.5 / printed p.74.

特别地，对 $X=Y$ 有：

\begin{theorem}[4.2.4]
\label{II.4.2.4-zh}
$\mathbf{P}(\mathcal{E})$ 的 $Y$-截面集合，典范地一一对应于
$\mathcal{E}$ 的拟凝聚 $\mathcal{O}_Y$-子模层 $\mathcal{F}$ 的集合，其中
$\mathcal{E}/\mathcal{F}$ 为可逆模层。
\end{theorem}

要注意，这个性质对应于把“射影空间”经典地定义成一个向量空间的超平面
集合。经典情形是 $Y=\operatorname{Spec}(K)$，其中 $K$ 是一个域，
$\mathcal{E}=\widetilde{E}$，而 $E$ 是有限维 $K$-向量空间；这时，具有
\egaref{II.4.2.4-zh}{4.2.4} 所述性质的各 $\mathcal{F}$ 对应于 $E$ 的
超平面。另一方面，由 \egaxref{I.3.4.5-zh}{I, 3.4.5}，已知
$\mathbf{P}(\mathcal{E})$ 的 $Y$-截面正是 $\mathbf{P}(\mathcal{E})$ 在 $K$
上的有理点。

\begin{remark}[4.2.5]
\label{II.4.2.5-zh}
从 $X$ 到 $P$ 的 $Y$-态射，与其图态射——即 $P\times_YX$ 的 $X$-截面——
之间有典范的一一对应（\egaxref{I.3.3.14-zh}{I, 3.3.14}）；所以反过来，
\egaref{II.4.2.3-zh}{4.2.3} 可由
\egaref{II.4.2.4-zh}{4.2.4} 推出。以
$\operatorname{Hyp}_Y(X,\mathcal{E})$ 表示
$\mathcal{E}\otimes_Y\mathcal{O}_X=q^*(\mathcal{E})$ 的拟凝聚
$\mathcal{O}_X$-子模层 $\mathcal{F}$ 的集合，其中
$q^*(\mathcal{E})/\mathcal{F}$ 是一个可逆 $\mathcal{O}_X$-模层。若
$g:X'\to X$ 是一个 $Y$-态射，则由 $g^*$ 右正合，有
\[
  g^*(q^*(\mathcal{E})/\mathcal{F})=
  g^*(q^*(\mathcal{E}))/g^*(\mathcal{F}),
\]
所以后一层可逆；因此，$\operatorname{Hyp}_Y(X,\mathcal{E})$ 是
$Y$-预概形范畴中的一个反变函子。于是可以把定理
\egaref{II.4.2.4-zh}{4.2.4} 解释为定义了函子
$\operatorname{Hom}_Y(X,\mathbf{P}(\mathcal{E}))$ 与
$\operatorname{Hyp}_Y(X,\mathcal{E})$ 之间的典范同构；它们在
$Y$-预概形范畴中关于变动的 $X$ 都是反变的。这也给出泛射影丛
$P=\mathbf{P}(\mathcal{E})$ 的如下泛性质刻画；它比第 2、3 节的构造更接近
几何直观：对每个态射 $q:X\to Y$ 以及每个可逆 $\mathcal{O}_X$-模层
$\mathcal{L}$，若后者是
$\mathcal{E}\otimes_{\mathcal{O}_Y}\mathcal{O}_X$ 的商，则存在唯一的
$Y$-态射 $r:X\to P$，使得 $\mathcal{L}=r^*(\mathcal{O}_P(1))$。

以后将看到，如何以同样方式定义格拉斯曼概形等对象。
\end{remark}
% Locked French authority: ega2-1-fr.tex, 820505 B,
% SHA-256 84EDBE3E83530AF2959B441796337C9DC21EAFCA6A13114A26778760FBF437AC.
% Sequential source scope: lines 6595-6640, corollary II.4.2.6 and its
% two consequences / printed pp.74-75.

\begin{corollary}[4.2.6]
\label{II.4.2.6-zh}
假设每个可逆 $\mathcal{O}_Y$-模层都是平凡的
（\egaxref{I.2.4.8-zh}{I, 2.4.8}）。设 $V$ 是群
$\operatorname{Hom}_{\mathcal{O}_Y}(\mathcal{E},\mathcal{O}_Y)$，把它看成
环 $A=\Gamma(Y,\mathcal{O}_Y)$ 上的模；并设 $V^*$ 是 $V$ 中所有满同态
组成的子集。则 $\mathbf{P}(\mathcal{E})$ 的 $Y$-截面集合典范地等同于
$V^*/A^*$，其中 $A^*$ 是 $A$ 的单位群。
\end{corollary}

\oldpage[II]{75}

特别地：
\begin{enumerate}
  \item[1\textsuperscript{o}] 当 $Y$ 是一个\emph{局部概形}时，推论
  \egaref{II.4.2.6-zh}{4.2.6} 适用
  （\egaxref{I.2.4.8-zh}{I, 2.4.8}）。设 $Y$ 是任意预概形，$y$ 是 $Y$
  的一个点，$Y'=\operatorname{Spec}(k(y))$。由
  \egaref{II.4.1.3.1-zh}{4.1.3.1}，$\mathbf{P}(\mathcal{E})$ 的纤维
  $p^{-1}(y)$ 等同于 $\mathbf{P}(\mathcal{E}^y)$，其中
  $\mathcal{E}^y=\mathcal{E}_y\otimes_{\mathcal{O}_y}k(y)
  =\mathcal{E}_y/\mathfrak{m}_y\mathcal{E}_y$ 被看成 $k(y)$ 上的向量空间。
  更一般地，若 $K$ 是 $k(y)$ 的一个扩张，则
  $p^{-1}(y)\otimes_{k(y)}K$ 等同于
  $\mathbf{P}(\mathcal{E}^y\otimes_{k(y)}K)$。所以推论
  \egaref{II.4.2.6-zh}{4.2.6} 表明，\emph{$\mathbf{P}(\mathcal{E})$ 取值于
  $k(y)$ 的扩张 $K$ 的几何点}（\egaxref{I.3.4.5-zh}{I, 3.4.5}）所成集合，
  也可称为\emph{$\mathbf{P}(\mathcal{E})$ 在点 $y$ 上、在 $K$ 上有理的
  几何纤维}，等同于与 $K$-向量空间
  $\mathcal{E}^y\otimes_{k(y)}K$ 的\emph{对偶}相伴的\emph{射影空间}。
  \item[2\textsuperscript{o}] 假设 $Y$ 是环 $A$ 的仿射概形，并进一步假设
  每个可逆 $\mathcal{O}_Y$-模层都是平凡的；再取
  $\mathcal{E}=\mathcal{O}_Y^n$。则在
  \egaref{II.4.2.6-zh}{4.2.6} 中，$V$ 等同于 $A^n$
  （\egaxref{I.1.3.8-zh}{I, 1.3.8}），而 $V^*$ 等同于 $A$ 中生成理想
  $A$ 的元素系 $(f_i)_{1\leq i\leq n}$ 的集合。两个这样的元素系定义
  $\mathbf{P}_Y^{n-1}=\mathbf{P}_A^{n-1}$ 的同一个 $Y$-截面，换言之，定义
  $\mathbf{P}_A^{n-1}$ 取值于 $A$ 的同一个点，当且仅当其中一个由另一个
  乘以 $A$ 的一个可逆元得到。
\end{enumerate}
% Locked French authority: ega2-1-fr.tex, 820505 B,
% SHA-256 84EDBE3E83530AF2959B441796337C9DC21EAFCA6A13114A26778760FBF437AC.
% Sequential source scope: lines 6642-6682, terminology discussion and
% remark II.4.2.7 / printed p.75.

这些性质说明，对 $\mathbf{P}(\mathcal{E})$ 使用“泛射影丛”这一术语是恰当的。
应当注意，由此得到的“射影空间”定义事实上与经典定义对偶；这是因为必须能够
对任意拟凝聚 $\mathcal{O}_Y$-模层 $\mathcal{E}$ 定义
$\mathbf{P}(\mathcal{E})$，而不能要求这个模层必为局部自由模层。

\begin{remark}[4.2.7]
\label{II.4.2.7-zh}
将在第~V 章看到：若 $Y$ 局部诺特且连通，并且 $\mathcal{E}$ 局部自由，
置 $P=\mathbf{P}(\mathcal{E})$，则每个可逆 $\mathcal{O}_P$-模层
$\mathcal{L}$ 都同构于形如
$\mathcal{L}'\otimes_{\mathcal{O}_Y}\mathcal{O}_P(m)$ 的
$\mathcal{O}_P$-模层，其中 $\mathcal{L}'$ 是一个在同构意义下唯一确定的
可逆 $\mathcal{O}_Y$-模层，而 $m$ 是一个唯一确定的整数。换言之，
$\mathrm{H}^1(P,\mathcal{O}_P^*)$ 同构于
$\mathbf{Z}\times\mathrm{H}^1(Y,\mathcal{O}_Y^*)$
（\egaxref{0.5.4.7-zh}{0, 5.4.7}）。还将看到
（\egaxref{III.2.1.14-zh}{III, 2.1.14}，并考虑
\egaxref{0.5.4.10-zh}{0, 5.4.10}）：当 $m<0$ 时，
$p_*(\mathcal{L})=0$\footnote{法文印本在本句两个直像式中均误将
\(\mathcal{L}\) 排作 \(\mathcal{L}^{\otimes m}\)；但前句已定义
\(\mathcal{L}\simeq\mathcal{L}'\otimes\mathcal{O}_P(m)\)，故本译在两个式子中均采用
\(p_*(\mathcal{L})\)。参见稳定源文勘误决定
EGA-II-4.2.7-P75-DIRECT-IMAGE-L-001、转介记录 EGA2-ZH-DEF-0009
及守护者去重处置 EGA2DEF86。}；当 $m\geq0$ 时，
$p_*(\mathcal{L})$ 同构于
$\mathcal{L}'\otimes_{\mathcal{O}_Y}
(\mathbf{S}_{\mathcal{O}_Y}(\mathcal{E}))_m$。若 $\mathcal{F}$ 是一个
拟凝聚 $\mathcal{O}_Y$-模层，则每个 $Y$-态射
$\mathbf{P}(\mathcal{E})\to\mathbf{P}(\mathcal{F})$ 都由下列资料确定：
一个可逆 $\mathcal{O}_Y$-模层 $\mathcal{L}'$、一个整数 $m\geq0$，以及
一个 $\mathcal{O}_Y$-同态
$\psi:\mathcal{F}\to\mathcal{L}'\otimes_{\mathcal{O}_Y}
(\mathbf{S}_{\mathcal{O}_Y}(\mathcal{E}))_m$，使得相应的
$\mathcal{O}_{\mathbf{P}(\mathcal{F})}$-模层同态 $\psi^\sharp$ 为满射。
还将看到，若所考虑的 $Y$-态射是同构，则 $m=1$，且 $\mathcal{F}$ 同构于
$\mathcal{E}\otimes_{\mathcal{O}_Y}\mathcal{L}'$
（\egaref{II.4.1.4-zh}{4.1.4} 的逆命题）。由此可以把
$\mathbf{P}(\mathcal{E})$ 的自同构芽层确定为群层
$\mathcal{A}ut(\mathcal{E})$（若 $\mathcal{E}$ 的秩为 $n$，则它局部同构于
$\operatorname{GL}(n,\mathcal{O}_Y)$）对 $\mathcal{O}_Y^*$ 的商。
\end{remark}

% EGA-II-4.2.7-P75-DIRECT-IMAGE-L-001 and EGA2DEF86 authorize only
% the two explicit p_*(L) target readings above; diplomatic French remains unchanged.
% Locked French authority: ega2-1-fr.tex, 820505 B,
% SHA-256 84EDBE3E83530AF2959B441796337C9DC21EAFCA6A13114A26778760FBF437AC.
% Sequential source scope: lines 6684-6713, environments II.4.2.8-4.2.9 /
% printed pp.75-76.

\begin{env}[4.2.8]
\label{II.4.2.8-zh}
沿用 \egaref{II.4.2.1-zh}{4.2.1} 的记号，设 $u:X'\to X$ 是一个态射；
若 $Y$-态射 $r:X\to P$ 对应于同态
$\varphi:q^*(\mathcal{E})\to\mathcal{L}$，则由定义立即可知，$Y$-态射
$r\circ u$ 对应于
$u^*(\varphi):u^*(q^*(\mathcal{E}))\to u^*(\mathcal{L})$。
\end{env}

\begin{env}[4.2.9]
\label{II.4.2.9-zh}
设 $\mathcal{E}$、$\mathcal{F}$ 是两个拟凝聚 $\mathcal{O}_Y$-模层，
$v:\mathcal{E}\to\mathcal{F}$ 是一个满同态，而 $j=\mathbf{P}(v)$ 是相应的
闭浸入
$\mathbf{P}(\mathcal{F})\to\mathbf{P}(\mathcal{E})$
（\egaref{II.4.1.2-zh}{4.1.2}）。若 $Y$-态射
$r:X\to\mathbf{P}(\mathcal{F})$ 对应于同态
$\varphi:q^*(\mathcal{F})\to\mathcal{L}$，则

\oldpage[II]{76}

$Y$-态射 $j\circ r$ 对应于
\[
  q^*(\mathcal{E})\xrightarrow{q^*(v)}q^*(\mathcal{F})
  \xrightarrow{\varphi}\mathcal{L}~;
\]
这仍由 \egaref{II.4.2.1-zh}{4.2.1} 中给出的定义得出。
\end{env}
% Locked French authority: ega2-1-fr.tex, 820505 B,
% SHA-256 84EDBE3E83530AF2959B441796337C9DC21EAFCA6A13114A26778760FBF437AC.
% Sequential source scope: lines 6715-6751, environment II.4.2.10 /
% printed p.76.

\begin{env}[4.2.10]
\label{II.4.2.10-zh}
设 $\psi:Y'\to Y$ 是一个态射，并置
$\mathcal{E}'=\psi^*(\mathcal{E})$。若 $Y$-态射 $r:X\to P$ 对应于同态
$\varphi:q^*(\mathcal{E})\to\mathcal{L}$，则 $Y'$-态射
\[
  r_{(Y')}:X_{(Y')}\longrightarrow P'=\mathbf{P}(\mathcal{E}')
\]
对应于
$\varphi_{(Y')}:q_{(Y')}^*(\mathcal{E}')=
  q^*(\mathcal{E})\otimes_Y\mathcal{O}_{Y'}\to
\mathcal{L}\otimes_Y\mathcal{O}_{Y'}$。事实上，由
\egaref{II.4.1.3.1-zh}{4.1.3.1}，有交换图
\[
  \xymatrix{
    Y' \ar[d]
    & P'=P_{(Y')} \ar[l]_{p_{(Y')}} \ar[d]^{u}
    & X_{(Y')} \ar[l]_{r_{(Y')}} \ar[d]^{v}
    \\Y
    & P \ar[l]_{p}
    & X \ar[l]_{r}
  }
\]
由 \egaref{II.4.1.3.2-zh}{4.1.3.2}，有
\[
  (r_{(Y')})^*(\mathcal{O}_{P'}(1))
  =(r_{(Y')})^*(u^*(\mathcal{O}_P(1)))
  =v^*(r^*(\mathcal{O}_P(1)))
  =v^*(\mathcal{L})=\mathcal{L}\otimes_Y\mathcal{O}_{Y'}~;
\]
另一方面，$u^*(\alpha_1^\sharp)$ 正是典范同态
$\alpha_1^\sharp:(p_{(Y')})^*(\mathcal{E}')\to\mathcal{O}_{P'}(1)$；像
\egaref{II.4.1.6-zh}{4.1.6} 那样，把 $P$ 和 $P'$ 中的典范同态
$\alpha_1^\sharp$ 明写出来即可看出这一点。断言因而成立。
\end{env}
% Locked French authority: ega2-1-fr.tex, 820505 B,
% SHA-256 84EDBE3E83530AF2959B441796337C9DC21EAFCA6A13114A26778760FBF437AC.
% Sequential source scope: lines 6753-6835, subsection II.4.3 through
% environment II.4.3.1 / printed pp.76-77.

\subsection{Segre 态射}
\label{subsection:II.4.3-zh}

\begin{env}[4.3.1]
\label{II.4.3.1-zh}
设 $Y$ 是一个预概形，$\mathcal{E}$、$\mathcal{F}$ 是两个拟凝聚
$\mathcal{O}_Y$-模层；置 $P_1=\mathbf{P}(\mathcal{E})$、
$P_2=\mathbf{P}(\mathcal{F})$，并以 $p_1:P_1\to Y$、$p_2:P_2\to Y$
表示结构态射。设 $Q=P_1\times_YP_2$，并以 $q_1:Q\to P_1$、
$q_2:Q\to P_2$ 表示典范投影；$\mathcal{O}_Q$-模层
\[
  \mathcal{L}=\mathcal{O}_{P_1}(1)\otimes_Y\mathcal{O}_{P_2}(1)
  :=q_1^*(\mathcal{O}_{P_1}(1))\otimes_{\mathcal{O}_Q}
  q_2^*(\mathcal{O}_{P_2}(1))
\]
作为两个可逆 $\mathcal{O}_Q$-模层的张量积，是可逆的
（\egaxref{0.5.4.4-zh}{0, 5.4.4}）。另一方面，若
$r=p_1\circ q_1=p_2\circ q_2$ 是结构态射 $Q\to Y$，则
\[
  r^*(\mathcal{E}\otimes_{\mathcal{O}_Y}\mathcal{F})
  =q_1^*(p_1^*(\mathcal{E}))\otimes_{\mathcal{O}_Q}
  q_2^*(p_2^*(\mathcal{F}))
\]
（\egaxref{0.4.3.3-zh}{0, 4.3.3}）；典范满同态
（\egaref{II.4.1.5.1-zh}{4.1.5.1}）
$p_1^*(\mathcal{E})\to\mathcal{O}_{P_1}(1)$、
$p_2^*(\mathcal{F})\to\mathcal{O}_{P_2}(1)$ 因而由张量积给出典范同态
\[
\label{II.4.3.1.1-zh}
  s:r^*(\mathcal{E}\otimes_{\mathcal{O}_Y}\mathcal{F})\longrightarrow
  \mathcal{L}
\tag{4.3.1.1}
\]
它显然是满射；由此便得到（\egaref{II.4.2.2-zh}{4.2.2}）一个称为
\emph{Segre 态射}的典范态射：
\[
\label{II.4.3.1.2-zh}
  \varsigma:\mathbf{P}(\mathcal{E})\times_Y\mathbf{P}(\mathcal{F})
  \longrightarrow
  \mathbf{P}(\mathcal{E}\otimes_{\mathcal{O}_Y}\mathcal{F}).
\tag{4.3.1.2}
\]
下面在 $Y=\operatorname{Spec}(A)$ 为仿射、
$\mathcal{E}=\widetilde{E}$、$\mathcal{F}=\widetilde{F}$ 的情形把态射
$\varsigma$ 明写出来，其中 $E$、$F$ 是两个 $A$-模，因而
$\mathcal{E}\otimes_{\mathcal{O}_Y}\mathcal{F}=\widetilde{E\otimes_A F}$
（\egaxref{I.1.3.12-zh}{I, 1.3.12}）。置
$R=\mathbf{S}_A(E)$、$S=\mathbf{S}_A(F)$、
$T=\mathbf{S}_A(E\otimes_A F)$；设 $f\in E$、$g\in F$，并考虑 $Q$ 的
仿射开集
\[
  D_+(f)\times_YD_+(g)=\operatorname{Spec}(B).
\]

\oldpage[II]{77}

这里 $B=R_{(f)}\otimes_AS_{(g)}$；$\mathcal{L}$ 到这个仿射开集的限制为
$\widetilde{L}$，其中
\[
  L=(R(1))_{(f)}\otimes_A(S(1))_{(g)},
\]
而元素 $c=(f/1)\otimes(g/1)$ 是把 $L$ 看成自由 $B$-模时的一个生成元
（\egaref{II.2.5.7-zh}{2.5.7}）。同态
\egaref{II.4.3.1.1-zh}{4.3.1.1} 对应于从
$(E\otimes_A F)\otimes_AB$ 到 $L$ 的同态
\[
  (x\otimes y)\otimes b\longmapsto b((x/1)\otimes(y/1)).
\]
沿用 \egaref{II.3.7.2-zh}{3.7.2} 的记号，因而有
$v_1(x\otimes y)=(x/f)\otimes(y/g)$；$\varsigma$ 到
$D_+(f)\times_YD_+(g)$ 的限制，是这个仿射概形到 $D_+(f\otimes g)$
的态射，对应于环同态
$\omega:T_{(f\otimes g)}\to R_{(f)}\otimes_AS_{(g)}$，满足
\[
\label{II.4.3.1.3-zh}
  \omega((x\otimes y)/(f\otimes g))=(x/f)\otimes(y/g)
\tag{4.3.1.3}
\]
其中 $x\in E$、$y\in F$。
\end{env}
% Locked French authority: ega2-1-fr.tex, 820505 B,
% SHA-256 84EDBE3E83530AF2959B441796337C9DC21EAFCA6A13114A26778760FBF437AC.
% Sequential source scope: lines 6837-6894, environment II.4.3.2 with
% lemma II.4.3.2.4 and proof / printed p.77.

\begin{env}[4.3.2]
\label{II.4.3.2-zh}
由 \egaref{II.4.2.3-zh}{4.2.3}，有典范同构
\[
\label{II.4.3.2.1-zh}
  \tau:\varsigma^*(\mathcal{O}_P(1))\xrightarrow{\sim}
  \mathcal{O}_{P_1}(1)\otimes_Y\mathcal{O}_{P_2}(1)
\tag{4.3.2.1}
\]
其中置 $P=\mathbf{P}(\mathcal{E}\otimes_{\mathcal{O}_Y}\mathcal{F})$。
下面证明：对 $x\in\Gamma(Y,\mathcal{E})$ 和
$y\in\Gamma(Y,\mathcal{F})$，有
\[
\label{II.4.3.2.2-zh}
  \tau(\alpha_1(x\otimes y))=\alpha_1(x)\otimes\alpha_1(y).
\tag{4.3.2.2}
\]
事实上，可以归结到 $Y$ 为仿射的情形；这时沿用
\egaref{II.4.3.1-zh}{4.3.1} 和 \egaref{II.2.6.2-zh}{2.6.2} 的记号，
$\alpha_1^{f\otimes g}(x\otimes y)=(x\otimes y)/1$ 属于
$(T(1))_{(f\otimes g)}$，$\alpha_1^f(x)=x/1$ 属于 $(R(1))_{(f)}$，而
$\alpha_1^g(y)=y/1$ 属于 $(S(1))_{(g)}$。由
\egaref{II.4.2.3-zh}{4.2.3} 中给出的 $\tau$ 的定义以及
\egaref{II.4.3.1-zh}{4.3.1} 中对 $v_1$ 的计算，立即得到
\egaref{II.4.3.2.2-zh}{4.3.2.2}。由此推出公式
\[
\label{II.4.3.2.3-zh}
  \varsigma^{-1}(P_{x\otimes y})=(P_1)_x\times_Y(P_2)_y
\tag{4.3.2.3}
\]
这里沿用 \egaref{II.3.1.4-zh}{3.1.4} 的记号。事实上，考虑到
\egaref{II.3.3.3-zh}{3.3.3}，公式
\egaref{II.4.3.2.2-zh}{4.3.2.2} 把所需结论（借助
\egaxref{I.3.2.7-zh}{I, 3.2.7} 和
\egaxref{I.3.2.3-zh}{I, 3.2.3} 回到仿射情形）归结为证明下述引理：

\begin{lemma}[4.3.2.4]
\label{II.4.3.2.4-zh}
设 $B$、$B'$ 是两个 $A$-代数，并设
$Y=\operatorname{Spec}(A)$、$Z=\operatorname{Spec}(B)$、
$Z'=\operatorname{Spec}(B')$；则对任意 $t\in B$、$t'\in B'$，有
$D(t\otimes t')=D(t)\times_YD(t')$。
\end{lemma}

\begin{proof}
事实上，若 $p:Z\times_YZ'\to Z$ 和 $p':Z\times_YZ'\to Z'$ 是典范投影，
则由 \egaxref{I.1.2.2.2-zh}{I, 1.2.2.2} 有
$p^{-1}(D(t))=D(t\otimes1)$ 和
$p'^{-1}(D(t'))=D(1\otimes t')$；再由
\egaxref{I.3.2.7-zh}{I, 3.2.7} 与
\egaxref{I.1.1.9.1-zh}{I, 1.1.9.1}，并注意
$(t\otimes1)(1\otimes t')=t\otimes t'$，便得到结论。
\end{proof}
\end{env}
% Locked French authority: ega2-1-fr.tex, 820505 B,
% SHA-256 84EDBE3E83530AF2959B441796337C9DC21EAFCA6A13114A26778760FBF437AC.
% Sequential source scope: lines 6896-6914, proposition II.4.3.3 and proof /
% printed p.77.

\begin{proposition}[4.3.3]
\label{II.4.3.3-zh}
Segre 态射是闭浸入。
\end{proposition}

\begin{proof}
事实上，这个问题在 $Y$ 上是局部的，因而可以归结到 $Y$ 为仿射的情形。
沿用 \egaref{II.4.3.1-zh}{4.3.1} 和 \egaref{II.4.3.2-zh}{4.3.2}
的记号，由于当 $f$ 遍历 $E$、$g$ 遍历 $F$ 时，各 $f\otimes g$ 生成
$T$，各 $D_+(f\otimes g)$ 构成 $P$ 的拓扑的一组基。另一方面，由
\egaref{II.4.3.2.3-zh}{4.3.2.3} 有
$\varsigma^{-1}(D_+(f\otimes g))=D_+(f)\times_YD_+(g)$。所以只需
（\egaxref{I.4.2.4-zh}{I, 4.2.4}）证明 $\varsigma$ 到
$D_+(f)\times_YD_+(g)$ 的限制是到 $D_+(f\otimes g)$ 的闭浸入。
但沿用同样的记号，公式 \egaref{II.4.3.1.3-zh}{4.3.1.3} 表明
$\omega$ 是满射，证明完毕。
\end{proof}
% Locked French authority: ega2-1-fr.tex, 820505 B,
% SHA-256 84EDBE3E83530AF2959B441796337C9DC21EAFCA6A13114A26778760FBF437AC.
% Sequential source scope: lines 6916-6951, environment II.4.3.4 /
% printed pp.77-78.

\begin{env}[4.3.4]
\label{II.4.3.4-zh}
若把拟凝聚 $\mathcal{O}_Y$-模层的同态限制为满同态，则 Segre 态射关于
$\mathcal{E}$ 和 $\mathcal{F}$ 具有函子性。事实上，需要证明：若

\oldpage[II]{78}

$\mathcal{E}\to\mathcal{E}'$ 是一个满 $\mathcal{O}_Y$-同态，则交换图
\[
  \xymatrix{
    \mathbf{P}(\mathcal{E}')\times_Y\mathbf{P}(\mathcal{F})
      \ar[r]^{j\times1} \ar[d]_{\varsigma}
    & \mathbf{P}(\mathcal{E})\times_Y\mathbf{P}(\mathcal{F})
      \ar[d]^{\varsigma}
    \\\mathbf{P}(\mathcal{E}'\otimes\mathcal{F}) \ar[r]
    & \mathbf{P}(\mathcal{E}\otimes\mathcal{F})
  }
\]
是交换的，其中 $j$ 表示典范闭浸入
$\mathbf{P}(\mathcal{E}')\to\mathbf{P}(\mathcal{E})$。置
$P_1'=\mathbf{P}(\mathcal{E}')$，其余仍沿用
\egaref{II.4.3.1-zh}{4.3.1} 的记号；$j\times1$ 是闭浸入
（\egaxref{I.4.3.1-zh}{I, 4.3.1}），并且在同构意义下有
\[
  (j\times1)^*(\mathcal{O}_{P_1}(1)\otimes\mathcal{O}_{P_2}(1))
  =j^*(\mathcal{O}_{P_1}(1))\otimes\mathcal{O}_{P_2}(1)
  =\mathcal{O}_{P_1'}(1)\otimes\mathcal{O}_{P_2}(1),
\]
这是 \egaref{II.4.1.2.1-zh}{4.1.2.1} 和
\egaxref{I.9.1.5-zh}{I, 9.1.5} 的结果；断言因而由
\egaref{II.4.2.8-zh}{4.2.8} 和 \egaref{II.4.2.9-zh}{4.2.9} 得出。
\end{env}
% Locked French authority: ega2-1-fr.tex, 820505 B,
% SHA-256 84EDBE3E83530AF2959B441796337C9DC21EAFCA6A13114A26778760FBF437AC.
% Sequential source scope: lines 6953-6979, environment II.4.3.5 /
% printed p.78.

\begin{env}[4.3.5]
\label{II.4.3.5-zh}
沿用 \egaref{II.4.3.1-zh}{4.3.1} 的记号，设 $\psi:Y'\to Y$ 是一个态射，
并置 $\mathcal{E}'=\psi^*(\mathcal{E})$、
$\mathcal{F}'=\psi^*(\mathcal{F})$；则 Segre 态射
$\mathbf{P}(\mathcal{E}')\times\mathbf{P}(\mathcal{F}')
\to\mathbf{P}(\mathcal{E}'\otimes\mathcal{F}')$ 等同于
$\varsigma_{(Y')}$。事实上，除继续沿用
\egaref{II.4.3.1-zh}{4.3.1} 的记号外，再置
$P_1'=\mathbf{P}(\mathcal{E}')$、$P_2'=\mathbf{P}(\mathcal{F}')$；由
\egaref{II.4.1.3.1-zh}{4.1.3.1}，已知 $P_i'$ 等同于 $(P_i)_{(Y')}$
（$i=1,2$），所以结构态射 $P_1'\times P_2'\to Y'$ 等同于
$r_{(Y')}$。另一方面，$\mathcal{E}'\otimes\mathcal{F}'$ 等同于
$\psi^*(\mathcal{E}\otimes\mathcal{F})$，故
$\mathbf{P}(\mathcal{E}'\otimes\mathcal{F}')$ 等同于
$(\mathbf{P}(\mathcal{E}\otimes\mathcal{F}))_{(Y')}$
（\egaref{II.4.1.3.1-zh}{4.1.3.1}）。最后，由
\egaref{II.4.1.3.2-zh}{4.1.3.2} 和
\egaxref{I.9.1.11-zh}{I, 9.1.11}，
$\mathcal{O}_{P_1'}(1)\otimes_{Y'}\mathcal{O}_{P_2'}(1)=\mathcal{L}'$
等同于 $\mathcal{L}\otimes_Y\mathcal{O}_{Y'}$。这时典范同态
$(r_{(Y')})^*(\mathcal{E}'\otimes\mathcal{F}')\to\mathcal{L}'$ 等同于
$s_{(Y')}$，断言由 \egaref{II.4.2.10-zh}{4.2.10} 得出。
\end{env}
% Locked French authority: ega2-1-fr.tex, 820505 B,
% SHA-256 84EDBE3E83530AF2959B441796337C9DC21EAFCA6A13114A26778760FBF437AC.
% Sequential source scope: lines 6981-7007, remark II.4.3.6 / printed p.78.

\begin{remark}[4.3.6]
\label{II.4.3.6-zh}
$\mathbf{P}(\mathcal{E})$ 与 $\mathbf{P}(\mathcal{F})$ 的和预概形同样典范地
同构于 $\mathbf{P}(\mathcal{E}\oplus\mathcal{F})$ 的一个闭子预概形。
事实上，满同态
$\mathcal{E}\oplus\mathcal{F}\to\mathcal{E}$、
$\mathcal{E}\oplus\mathcal{F}\to\mathcal{F}$ 对应于闭浸入
$\mathbf{P}(\mathcal{E})\to\mathbf{P}(\mathcal{E}\oplus\mathcal{F})$、
$\mathbf{P}(\mathcal{F})\to\mathbf{P}(\mathcal{E}\oplus\mathcal{F})$；
只需看出，由此得到的 $\mathbf{P}(\mathcal{E}\oplus\mathcal{F})$ 的两个
闭子预概形的底空间没有公共点。这个问题在 $Y$ 上是局部的，因而可以采用
\egaref{II.4.3.1-zh}{4.3.1} 的记号；$\mathbf{S}_n(E)$ 和
$\mathbf{S}_n(F)$ 可等同于 $\mathbf{S}_n(E\oplus F)$ 中两个交仅为 $0$
的子模；若 $\mathfrak{p}$ 是 $\mathbf{S}(E)$ 的一个分次素理想，并且对
每个 $n\geq0$ 都有
$\mathfrak{p}\cap\mathbf{S}_n(E)\neq\mathbf{S}_n(E)$，则在
$\mathbf{S}(E\oplus F)$ 中有一个与它对应的分次素理想，其在
$\mathbf{S}_n(E)$ 上的迹为
$\mathfrak{p}\cap\mathbf{S}_n(E)$，但它包含 $\mathbf{S}_+(F)$；这一点
立即可见。因此，$\operatorname{Proj}(\mathbf{S}(E))$ 和
$\operatorname{Proj}(\mathbf{S}(F))$ 各自的两个点，不可能在
$\operatorname{Proj}(\mathbf{S}(E\oplus F))$ 中有相同的像。
\end{remark}
% Locked French authority: ega2-1-fr.tex, 820505 B,
% SHA-256 84EDBE3E83530AF2959B441796337C9DC21EAFCA6A13114A26778760FBF437AC.
% Sequential source scope: lines 7009-7075, subsection II.4.4 through
% proposition II.4.4.1 and its proof / printed pp.78-79.

\subsection{泛射影丛中的浸入；极丰沛层}
\label{subsection:II.4.4-zh}

\begin{proposition}[4.4.1]
\label{II.4.4.1-zh}
设 $Y$ 是一个拟紧概形，或者是一个底空间为诺特空间的预概形；
$q:X\to Y$ 是一个有限型态射，$\mathcal{L}$ 是一个可逆
$\mathcal{O}_X$-模层。
\begin{enumerate}
  \item[\rm{(i)}] 设 $\mathcal{S}$ 是一个按正次数分次的拟凝聚
  $\mathcal{O}_Y$-代数层，
  $\psi:q^*(\mathcal{S})\to\bigoplus_{n\geq0}\mathcal{L}^{\otimes n}$
  是一个分次代数同态。态射 $r_{\mathcal{L},\psi}$ 处处有定义并且是
  浸入，当且

\oldpage[II]{79}

  仅当存在整数 $n\geq0$ 以及 $\mathcal{S}_n$ 的一个有限型拟凝聚
  $\mathcal{O}_Y$-子模层 $\mathcal{E}$，使得同态
  \[
    \psi'=\psi_n\circ q^*(j):q^*(\mathcal{E})
      \longrightarrow\mathcal{L}^{\otimes n}=\mathcal{L}'
  \]
  为满射，并且态射
  $r_{\mathcal{L}',\psi'}:X\to\mathbf{P}(\mathcal{E})$ 是浸入；这里
  $j$ 是单射 $\mathcal{E}\to\mathcal{S}_n$。
  \item[\rm{(ii)}] 设 $\mathcal{F}$ 是一个拟凝聚
  $\mathcal{O}_Y$-模层，
  $\varphi:q^*(\mathcal{F})\to\mathcal{L}$ 是一个满同态。态射
  $r_{\mathcal{L},\varphi}$ 是浸入
  $X\to\mathbf{P}(\mathcal{F})$，
  当且仅当 $\mathcal{F}$ 存在一个有限型拟凝聚
  $\mathcal{O}_Y$-子模层 $\mathcal{E}$，使同态
  $\varphi'=\varphi\circ q^*(j):q^*(\mathcal{E})\to\mathcal{L}$
  为满射，并且
  $r_{\mathcal{L},\varphi'}:X\to\mathbf{P}(\mathcal{E})$ 是浸入；
  这里 $j:\mathcal{E}\to\mathcal{F}$ 是典范单射。
\end{enumerate}
\end{proposition}

\begin{proof}
\begin{enumerate}
  \item[\rm{(i)}] 由 \egaref{II.3.8.5-zh}{3.8.5}，
  $r_{\mathcal{L},\psi}$ 处处有定义并且是浸入，等价于存在
  $n\geq0$ 和 $\mathcal{E}$，使得若 $\mathcal{S}'$ 是由
  $\mathcal{E}$ 生成的 $\mathcal{S}$ 的子代数层，则同态
  $q^*(\mathcal{E})\to\mathcal{L}^{\otimes n}$ 为满射，而且
  $r_{\mathcal{L}',\psi'}:X\to\operatorname{Proj}(\mathcal{S}')$
  处处有定义并且是浸入。另一方面，有典范满同态
  $\mathbf{S}(\mathcal{E})\to\mathcal{S}'$，与之对应一个闭浸入
  $\operatorname{Proj}(\mathcal{S}')\to\mathbf{P}(\mathcal{E})$
  （\egaref{II.3.6.2-zh}{3.6.2}）；结论随即成立。
  \item[\rm{(ii)}] 因为 $\mathcal{F}$ 是它的有限型拟凝聚子模层
  $\mathcal{E}_\lambda$ 的归纳极限
  （\egaxref{I.9.4.9-zh}{I, 9.4.9}），所以
  $\mathbf{S}(\mathcal{F})$ 是各
  $\mathbf{S}(\mathcal{E}_\lambda)$ 的归纳极限。于是结论由
  \egaref{II.3.8.4-zh}{3.8.4} 得出；只须注意，在
  \egaref{II.3.8.4-zh}{3.8.4} 的证明中，可以把所有 $n_i$ 都取为
  $1$。事实上，假设 $Y$ 仿射；若
  $r=r_{\mathcal{L},\varphi}$ 是浸入，则 $r(X)$ 是
  $\mathbf{P}(\mathcal{F})$ 的拟紧子空间，可以用有限多个形如
  $D_+(f)$（其中 $f\in F$）的 $\mathbf{P}(\mathcal{F})$ 的开集覆盖，
  并且每个 $D_+(f)\cap r(X)$ 都是闭的。
\end{enumerate}
\end{proof}
% Locked French authority: ega2-1-fr.tex, 820505 B,
% SHA-256 84EDBE3E83530AF2959B441796337C9DC21EAFCA6A13114A26778760FBF437AC.
% Sequential source scope: lines 7077-7155, definition II.4.4.2 through
% proposition II.4.4.4 and its proof / printed pp.79-80.

\begin{definition}[4.4.2]
\label{II.4.4.2-zh}
设 $Y$ 是一个预概形，$q:X\to Y$ 是一个态射。若存在一个拟凝聚
$\mathcal{O}_Y$-模层 $\mathcal{E}$ 以及 $X$ 到
$P=\mathbf{P}(\mathcal{E})$ 中的一个 $Y$-浸入 $i$，使得
$\mathcal{L}$ 同构于 $i^*(\mathcal{O}_P(1))$，则称可逆
$\mathcal{O}_X$-模层 $\mathcal{L}$ 对 $q$ 是\emph{极丰沛的}
（或称它对 $Y$ 是\emph{极丰沛的}；若不致混淆，则简称为
\emph{极丰沛的}）。
\end{definition}

等价地（\egaref{II.4.2.3-zh}{4.2.3}），存在一个拟凝聚
$\mathcal{O}_Y$-模层 $\mathcal{E}$ 以及一个满同态
$\varphi:q^*(\mathcal{E})\to\mathcal{L}$，使
$r_{\mathcal{L},\varphi}:X\to\mathbf{P}(\mathcal{E})$ 是浸入。

注意，存在一个对 $Y$ 极丰沛的 $\mathcal{O}_X$-模层，蕴含 $q$ 是分离的
（\egaref{II.3.1.3-zh}{3.1.3} 与
\egaxref{I.5.5.1-zh}{I, 5.5.1, (i) 和 (ii)}）。

\begin{corollary}[4.4.3]
\label{II.4.4.3-zh}
假设存在一个由 $\mathcal{S}_1$ 生成的拟凝聚分次
$\mathcal{O}_Y$-代数层 $\mathcal{S}$，以及一个 $Y$-浸入
$i:X\to P=\operatorname{Proj}(\mathcal{S})$，使得 $\mathcal{L}$
同构于 $i^*(\mathcal{O}_P(1))$。那么 $\mathcal{L}$ 对 $q$ 是
极丰沛的。
\end{corollary}

\begin{proof}
事实上，若 $\mathcal{F}=\mathcal{S}_1$，则典范同态
$\mathbf{S}(\mathcal{F})\to\mathcal{S}$ 为满射。因此，把相应的闭浸入
$\operatorname{Proj}(\mathcal{S})\to\mathbf{P}(\mathcal{F})$
（\egaref{II.3.6.2-zh}{3.6.2}）与浸入 $i$ 复合，便得到浸入
$j:X\to\mathbf{P}(\mathcal{F})=P'$，使得 $\mathcal{L}$ 同构于
$j^*(\mathcal{O}_{P'}(1))$（\egaref{II.3.6.3-zh}{3.6.3}）。
\end{proof}

\begin{proposition}[4.4.4]
\label{II.4.4.4-zh}
设 $q:X\to Y$ 是一个拟紧态射，$\mathcal{L}$ 是一个可逆
$\mathcal{O}_X$-模层。下列性质等价：
\begin{enumerate}
  \item[\rm{a)}] $\mathcal{L}$ 对 $q$ 是极丰沛的。
  \item[\rm{b)}] $q_*(\mathcal{L})$ 是拟凝聚的，典范同态
  $\sigma:q^*(q_*(\mathcal{L}))\to\mathcal{L}$ 为满射，并且态射
  $r_{\mathcal{L},\sigma}:X\to\mathbf{P}(q_*(\mathcal{L}))$ 是浸入。
\end{enumerate}
\end{proposition}

\begin{proof}
因为 $q$ 拟紧，所以若 $q$ 分离，则已知 $q_*(\mathcal{L})$ 是拟凝聚的
（\egaxref{I.9.2.2-zh}{I, 9.2.2}）。

\oldpage[II]{80}

已知（\egaref{II.3.4.7-zh}{3.4.7}），存在满同态
$\varphi:q^*(\mathcal{E})\to\mathcal{L}$（其中 $\mathcal{E}$ 是一个
拟凝聚 $\mathcal{O}_Y$-模层）蕴含 $\sigma$ 为满射。此外，
$\varphi$ 的分解
$q^*(\mathcal{E})\to q^*(q_*(\mathcal{L}))
\xrightarrow{\sigma}\mathcal{L}$
典范地对应于分解
\[
  q^*(\mathbf{S}(\mathcal{E}))
    \longrightarrow q^*(\mathbf{S}(q_*(\mathcal{L})))
    \longrightarrow\bigoplus_{n\geq0}\mathcal{L}^{\otimes n}
\]
故由 \egaref{II.3.8.3-zh}{3.8.3}，假设
$r_{\mathcal{L},\varphi}$ 是浸入，便蕴含
$j=r_{\mathcal{L},\sigma}$ 也是浸入。此外，由
\egaref{II.4.2.4-zh}{4.2.4}，置
$P'=\mathbf{P}(q_*(\mathcal{L}))$ 时，$\mathcal{L}$ 同构于
$j^*(\mathcal{O}_{P'}(1))$。由此可见 a) 与 b) 等价。
\end{proof}
% Locked French authority: ega2-1-fr.tex, 820505 B,
% SHA-256 84EDBE3E83530AF2959B441796337C9DC21EAFCA6A13114A26778760FBF437AC.
% Sequential source scope: lines 7157-7221, corollary II.4.4.5 through
% corollary II.4.4.7 and its proof / printed p.80.

\begin{corollary}[4.4.5]
\label{II.4.4.5-zh}
假设 $q$ 拟紧。$\mathcal{L}$ 对 $Y$ 是极丰沛的，当且仅当 $Y$ 存在一个
开覆盖 $(U_\alpha)$，使得对每个 $\alpha$，
$\mathcal{L}|q^{-1}(U_\alpha)$ 对 $U_\alpha$ 是极丰沛的。
\end{corollary}

\begin{proof}
事实上，\egaref{II.4.4.4-zh}{4.4.4} 的条件 b) 关于 $Y$ 是局部的。
\end{proof}

\begin{proposition}[4.4.6]
\label{II.4.4.6-zh}
设 $Y$ 是一个拟紧概形，或者是一个底空间为诺特空间的预概形；
$q:X\to Y$ 是一个有限型态射，$\mathcal{L}$ 是一个可逆
$\mathcal{O}_X$-模层。此时，\egaref{II.4.4.4-zh}{4.4.4} 的条件 a)、
b) 还分别等价于下列条件：
\begin{enumerate}
  \item[\rm{a')}] 存在一个有限型拟凝聚 $\mathcal{O}_Y$-模层
  $\mathcal{E}$ 以及一个满同态
  $\varphi:q^*(\mathcal{E})\to\mathcal{L}$，使
  $r_{\mathcal{L},\varphi}$ 是浸入。
  \item[\rm{b')}] $q_*(\mathcal{L})$ 存在一个有限型凝聚
  $\mathcal{O}_Y$-子模层 $\mathcal{E}$，它具有 a') 中所述的性质。
\end{enumerate}
\end{proposition}

\begin{proof}
显然，a') 或 b') 蕴含 a)；另一方面，由
\egaref{II.4.4.1-zh}{4.4.1}，a) 蕴含 a')；同样，b) 蕴含 b')。
\end{proof}

\begin{corollary}[4.4.7]
\label{II.4.4.7-zh}
假设 $Y$ 是一个拟紧概形或一个诺特预概形。若 $\mathcal{L}$ 对 $q$
极丰沛，则存在一个拟凝聚分次 $\mathcal{O}_Y$-代数层 $\mathcal{S}$，
使得 $\mathcal{S}_1$ 为有限型并生成 $\mathcal{S}$，并且存在一个稠密
开 $Y$-浸入 $i:X\to P=\operatorname{Proj}(\mathcal{S})$，使
$\mathcal{L}$ 同构于 $i^*(\mathcal{O}_P(1))$。
\end{corollary}

\begin{proof}
事实上，\egaref{II.4.4.6-zh}{4.4.6} 的条件 b') 成立。此时结构态射
$p:\mathbf{P}(\mathcal{E})=P'\to Y$ 是分离且有限型的
（\egaref{II.3.1.3-zh}{3.1.3}），所以当 $Y$ 是拟紧概形（相应地，
诺特预概形）时，$P'$ 是拟紧概形（相应地，诺特预概形）。设 $Z$ 是
$P'$ 的子预概形 $X'$ 的闭包（\egaxref{I.9.5.11-zh}{I, 9.5.11}），
这里这个子预概形与 $X$ 到 $P'$ 中的浸入
$j=r_{\mathcal{L},\varphi}$ 相伴。显然，$j$ 分解为稠密开浸入
$i:X\to Z$，再接典范单射 $Z\to P'$。但是 $Z$ 可与一个预概形
$\operatorname{Proj}(\mathcal{S})$ 等同，其中 $\mathcal{S}$ 是
$\mathbf{S}(\mathcal{E})$ 对一个拟凝聚分次理想层取商所得的分次
$\mathcal{O}_Y$-代数层（\egaref{II.3.6.2-zh}{3.6.2}）；显然
$\mathcal{S}_1$ 为有限型并生成 $\mathcal{S}$。此外，
$\mathcal{O}_Z(1)$ 是 $\mathcal{O}_{P'}(1)$ 沿典范单射的逆像
（\egaref{II.3.6.3-zh}{3.6.3}），所以
$\mathcal{L}=i^*(\mathcal{O}_Z(1))$。
\end{proof}
% Locked French authority: ega2-1-fr.tex, 820505 B,
% SHA-256 84EDBE3E83530AF2959B441796337C9DC21EAFCA6A13114A26778760FBF437AC.
% Sequential source scope: lines 7223-7282, proposition II.4.4.8 through
% corollary II.4.4.9 and its proof / printed pp.80-81.

\begin{proposition}[4.4.8]
\label{II.4.4.8-zh}
设 $q:X\to Y$ 是一个态射，$\mathcal{L}$ 是一个对 $q$ 极丰沛的
$\mathcal{O}_X$-模层，$\mathcal{L}'$ 是一个可逆
$\mathcal{O}_X$-模层，并且存在一个拟凝聚 $\mathcal{O}_Y$-模层
$\mathcal{E}'$ 以及一个满同态
$q^*(\mathcal{E}')\to\mathcal{L}'$。那么
$\mathcal{L}\otimes_{\mathcal{O}_X}\mathcal{L}'$ 对 $q$ 是极丰沛的。
\end{proposition}

\begin{proof}
事实上，假设蕴含存在一个 $Y$-态射
$r':X\to P'=\mathbf{P}(\mathcal{E}')$，使
$\mathcal{L}'=r'^*(\mathcal{O}_{P'}(1))$
（\egaref{II.4.2.1-zh}{4.2.1}）。又由假设，存在一个拟凝聚
$\mathcal{O}_Y$-模层 $\mathcal{E}$ 以及一个

\oldpage[II]{81}

$Y$-浸入 $r:X\to P=\mathbf{P}(\mathcal{E})$，使
$\mathcal{L}=r^*(\mathcal{O}_P(1))$。置
$Q=\mathbf{P}(\mathcal{E}\otimes\mathcal{E}')$，并考虑 Segre 态射
$\varsigma:P\times_YP'\to Q$（\egaref{II.4.3.1-zh}{4.3.1}）。由于
$r$ 是浸入，$(r,r')_Y:X\to P\times_YP'$ 也是浸入
（\egaxref{I.5.3.14-zh}{I, 5.3.14}）；又由于 $\varsigma$ 是浸入
（\egaref{II.4.3.3-zh}{4.3.3}），
$r'':X\xrightarrow{(r,r')_Y}P\times_YP'\xrightarrow{\varsigma}Q$
也是浸入。另一方面，由 \egaref{II.4.3.2.1-zh}{4.3.2.1}，
$\varsigma^*(\mathcal{O}_Q(1))$ 同构于
$\mathcal{O}_P(1)\otimes_Y\mathcal{O}_{P'}(1)$，从而由
\egaxref{I.9.1.4-zh}{I, 9.1.4}，
$r''^*(\mathcal{O}_Q(1))$ 同构于
$\mathcal{L}\otimes\mathcal{L}'$；命题得证。
\end{proof}

\begin{corollary}[4.4.9]
\label{II.4.4.9-zh}
设 $q:X\to Y$ 是一个态射。
\begin{enumerate}
  \item[\rm{(i)}] 设 $\mathcal{L}$ 是一个可逆 $\mathcal{O}_X$-模层，
  $\mathcal{K}$ 是一个可逆 $\mathcal{O}_Y$-模层。$\mathcal{L}$ 对
  $q$ 是极丰沛的，当且仅当
  $\mathcal{L}\otimes q^*(\mathcal{K})$ 也是极丰沛的。
  \item[\rm{(ii)}] 若 $\mathcal{L}$、$\mathcal{L}'$ 是两个对 $q$
  极丰沛的 $\mathcal{O}_X$-模层，则
  $\mathcal{L}\otimes\mathcal{L}'$ 也是极丰沛的；特别地，
  对每个 $n>0$，$\mathcal{L}^{\otimes n}$ 对 $q$ 是极丰沛的。
\end{enumerate}
\end{corollary}

\begin{proof}
(ii) 是 \egaref{II.4.4.8-zh}{4.4.8} 的直接推论，(i) 中条件的必要性也
同样如此。另一方面，若 $\mathcal{L}\otimes q^*(\mathcal{K})$ 是
极丰沛的，则由上述结果，
$(\mathcal{L}\otimes q^*(\mathcal{K}))
\otimes q^*(\mathcal{K}^{-1})$ 也是极丰沛的，而后一个
$\mathcal{O}_X$-模层同构于 $\mathcal{L}$
（\egaxref{0.5.4.3-zh}{0, 5.4.3} 与
\egaxref{0.5.4.5-zh}{5.4.5}）。
\end{proof}
% Locked French authority: ega2-1-fr.tex, 820505 B,
% SHA-256 84EDBE3E83530AF2959B441796337C9DC21EAFCA6A13114A26778760FBF437AC.
% Sequential source scope: lines 7284-7452, proposition II.4.4.10,
% lemma II.4.4.10.1, and the complete proof / printed pp.81-83.

\begin{proposition}[4.4.10]
\label{II.4.4.10-zh}
\begin{enumerate}
  \item[\rm{(i)}] 对任意预概形 $Y$，任意可逆
  $\mathcal{O}_Y$-模层 $\mathcal{L}$ 对恒等态射 $1_Y$ 都是极丰沛的。
  \item[\rm{(i bis)}] 设 $f:X\to Y$ 是一个态射，
  $j:X'\to X$ 是一个浸入。若 $\mathcal{L}$ 是一个对 $f$ 极丰沛的
  $\mathcal{O}_X$-模层，则 $j^*(\mathcal{L})$ 对 $f\circ j$ 是
  极丰沛的。
  \item[\rm{(ii)}] 设 $Z$ 是一个拟紧预概形，$f:X\to Y$ 是一个
  有限型态射，$g:Y\to Z$ 是一个拟紧态射，$\mathcal{L}$ 是一个对
  $f$ 极丰沛的 $\mathcal{O}_X$-模层，$\mathcal{K}$ 是一个对 $g$
  极丰沛的 $\mathcal{O}_Y$-模层。那么存在整数 $n_0>0$，使得对所有
  $n\geq n_0$，$\mathcal{L}\otimes f^*(\mathcal{K}^{\otimes n})$
  对 $g\circ f$ 是极丰沛的。
  \item[\rm{(iii)}] 设 $f:X\to Y$、$g:Y'\to Y$ 是两个态射，并置
  $X'=X_{(Y')}$。若 $\mathcal{L}$ 是一个对 $f$ 极丰沛的
  $\mathcal{O}_X$-模层，则
  $\mathcal{L}'=\mathcal{L}\otimes_Y\mathcal{O}_{Y'}$ 是一个对
  $f_{(Y')}$ 极丰沛的 $\mathcal{O}_{X'}$-模层。
  \item[\rm{(iv)}] 设 $f_i:X_i\to Y_i$（$i=1,2$）是两个
  $S$-态射。若 $\mathcal{L}_i$ 是一个对 $f_i$ 极丰沛的
  $\mathcal{O}_{X_i}$-模层（$i=1,2$），则
  $\mathcal{L}_1\otimes_S\mathcal{L}_2$ 对
  $f_1\times_Sf_2$ 是极丰沛的。
  \item[\rm{(v)}] 设 $f:X\to Y$、$g:Y\to Z$ 是两个态射。若一个
  $\mathcal{O}_X$-模层 $\mathcal{L}$ 对 $g\circ f$ 是极丰沛的，
  则 $\mathcal{L}$ 对 $f$ 是极丰沛的。
  \item[\rm{(vi)}] 设 $f:X\to Y$ 是一个态射，$j$ 是典范单射
  $X_{\mathrm{red}}\to X$。若 $\mathcal{L}$ 是一个对 $f$ 极丰沛的
  $\mathcal{O}_X$-模层，则 $j^*(\mathcal{L})$ 对
  $f_{\mathrm{red}}$ 是极丰沛的。
\end{enumerate}
\end{proposition}

\begin{proof}
性质 (i bis) 立即由定义 \egaref{II.4.4.2-zh}{4.4.2} 得出；仿照
\egaxref{I.5.5.12-zh}{I, 5.5.12} 的论证，还立即可见 (vi) 在形式上由
(i bis) 与 (v) 得出；我们把这个论证留给读者。为证明 (v)，像
\egaxref{I.5.5.12-zh}{I, 5.5.12} 那样考虑分解
$X\xrightarrow{\Gamma_f}X\times_ZY\xrightarrow{p_2}Y$，并注意
$p_2=(g\circ f)\times1_Y$。由假设以及 (i)、(iv)，
$\mathcal{L}\otimes_{\mathcal{O}_Z}\mathcal{O}_Y$ 对 $p_2$ 是
极丰沛的；另一方面，
$\mathcal{L}=\Gamma_f^*(\mathcal{L}\otimes_{\mathcal{O}_Z}
\mathcal{O}_Y)$（\egaxref{I.9.1.4-zh}{I, 9.1.4}），并且
$\Gamma_f$ 是浸入（\egaxref{I.5.3.11-zh}{I, 5.3.11}）；故可应用
(i bis)。

\oldpage[II]{82}

为证明 (i)，在定义 \egaref{II.4.4.2-zh}{4.4.2} 中取
$\mathcal{E}=\mathcal{L}$，并注意此时
$\mathbf{P}(\mathcal{E})$ 可与 $Y$ 等同
（\egaref{II.4.1.4-zh}{4.1.4}）。

下面证明 (iii)。存在一个拟凝聚 $\mathcal{O}_Y$-模层
$\mathcal{E}$ 以及一个 $Y$-浸入
$i:X\to\mathbf{P}(\mathcal{E})=P$，使
$\mathcal{L}=i^*(\mathcal{O}_P(1))$。若置
$\mathcal{E}'=g^*(\mathcal{E})$，则 $\mathcal{E}'$ 是一个拟凝聚
$\mathcal{O}_{Y'}$-模层，并且
$P'=\mathbf{P}(\mathcal{E}')=P_{(Y')}$（\egaref{II.4.1.3.1-zh}{4.1.3.1}）；
$i_{(Y')}$ 是 $X_{(Y')}$ 到 $P'$ 中的浸入
（\egaxref{I.4.3.2-zh}{I, 4.3.2}），而 $\mathcal{L}'$ 同构于
$(i_{(Y')})^*(\mathcal{O}_{P'}(1))$
（\egaref{II.4.2.10-zh}{4.2.10}）。

为证明 (iv)，注意由假设，存在一个 $Y_i$-浸入
$r_i:X_i\to P_i=\mathbf{P}(\mathcal{E}_i)$，其中
$\mathcal{E}_i$ 是一个拟凝聚 $\mathcal{O}_{Y_i}$-模层，并且
$\mathcal{L}_i=r_i^*(\mathcal{O}_{P_i}(1))$（$i=1,2$）。
$r_1\times_Sr_2$ 是 $X_1\times_SX_2$ 到 $P_1\times_SP_2$ 中的一个
$S$-浸入（\egaxref{I.4.3.1-zh}{I, 4.3.1}），并且
$\mathcal{O}_{P_1}(1)\otimes_S\mathcal{O}_{P_2}(1)$ 沿这个浸入的
逆像是 $\mathcal{L}_1\otimes_S\mathcal{L}_2$。另一方面，置
$T=Y_1\times_SY_2$，并分别以 $p_1,p_2$ 表示 $T$ 到 $Y_1,Y_2$ 的
投影。若置 $P_i'=\mathbf{P}(p_i^*(\mathcal{E}_i))$（$i=1,2$），则由
\egaref{II.4.1.3.1-zh}{4.1.3.1} 有
$P_i'=P_i\times_{Y_i}T$，因而
\[
  P_1'\times_TP_2'
  =(P_1\times_{Y_1}T)\times_T(P_2\times_{Y_2}T)
  =P_1\times_{Y_1}(T\times_{Y_2}P_2)
  =P_1\times_{Y_1}(Y_1\times_SP_2)
  =P_1\times_SP_2
\]
精确到典范同构。同样，
$\mathcal{O}_{P_i'}(1)=\mathcal{O}_{P_i}(1)
\otimes_{Y_i}\mathcal{O}_T$
（\egaref{II.4.1.3.2-zh}{4.1.3.2}）；类似的计算——尤其以
\egaxref{I.9.1.9.1-zh}{I, 9.1.9.1} 和
\egaxref{I.9.1.2-zh}{9.1.2} 为依据——表明，在上述等同中，
$\mathcal{O}_{P_1'}(1)\otimes_T\mathcal{O}_{P_2'}(1)$ 与
$\mathcal{O}_{P_1}(1)\otimes_S\mathcal{O}_{P_2}(1)$ 等同。因此可把
$r_1\times_Sr_2$ 看成 $X_1\times_SX_2$ 到
$P_1'\times_TP_2'$ 中的一个 $T$-浸入，而
$\mathcal{O}_{P_1'}(1)\otimes_T\mathcal{O}_{P_2'}(1)$ 沿这个浸入的
逆像是 $\mathcal{L}_1\otimes_S\mathcal{L}_2$。最后像
\egaref{II.4.4.8-zh}{4.4.8} 中那样，利用 Segre 态射结束论证。

只剩证明 (ii)。首先可归结到 $Z$ 是仿射概形的情形，因为一般地，$Z$
有一个由仿射开集组成的有限覆盖 $(U_i)$；若命题对
$\mathcal{K}|g^{-1}(U_i)$、$\mathcal{L}|f^{-1}(g^{-1}(U_i))$ 和某个
整数 $n_i$ 已经证明，则只须取 $n_0$ 为各 $n_i$ 中的最大者，便能由
\egaref{II.4.4.5-zh}{4.4.5} 得到关于 $\mathcal{K}$ 和
$\mathcal{L}$ 的命题。假设蕴含 $f$、$g$ 都是分离态射，故 $X$、$Y$
都是拟紧概形。

由 \egaref{II.4.4.6-zh}{4.4.6}，存在浸入
$r:X\to P=\mathbf{P}(\mathcal{E})$，其中 $\mathcal{E}$ 是一个有限型
拟凝聚 $\mathcal{O}_Y$-模层，并且
$\mathcal{L}=r^*(\mathcal{O}_P(1))$。我们将证明存在一个对复合态射
$P\xrightarrow{h}Y\xrightarrow{g}Z$ 极丰沛的
$\mathcal{O}_P$-模层 $\mathcal{M}$，使得对某个整数 $m$，
$\mathcal{O}_P(1)$ 同构于
$\mathcal{M}\otimes_Y\mathcal{K}^{\otimes(-m)}$。对
$n\geq m+1$，由假设以及把 (iv) 应用于态射 $h:P\to Y$、$1_Y$，
$\mathcal{O}_P(1)\otimes_Y\mathcal{K}^{\otimes n}$ 对 $Z$ 是
极丰沛的；又因 $r$ 是浸入，并且
$\mathcal{L}\otimes f^*(\mathcal{K}^{\otimes n})
=r^*(\mathcal{O}_P(1)\otimes_Y\mathcal{K}^{\otimes n})$，结论便由
(i bis) 得出。为证明关于 $\mathcal{O}_P(1)$ 的断言，我们将用到
下列引理：

\begin{lemma}[4.4.10.1]
\label{II.4.4.10.1-zh}
设 $Z$ 是一个拟紧概形，或者是一个底空间为诺特空间的预概形；
$g:Y\to Z$ 是一个拟紧态射，$\mathcal{K}$ 是一个对 $g$ 极丰沛的
可逆 $\mathcal{O}_Y$-模层，$\mathcal{E}$ 是一个有限型拟凝聚
$\mathcal{O}_Y$-模层。那么存在整数 $m_0$，使得对每个 $m\geq m_0$，
$\mathcal{E}$ 同构于一个形如
$g^*(\mathcal{F})\otimes\mathcal{K}^{\otimes(-m)}$ 的
$\mathcal{O}_Y$-模层的商，其中 $\mathcal{F}$ 是一个依赖于 $m$ 的
有限型拟凝聚 $\mathcal{O}_Z$-模层。
\end{lemma}

这个引理将在 \egaref{II.4.5.10.1-zh}{4.5.10.1} 中证明；读者可以核验，
\egaref{II.4.4.10-zh}{4.4.10} 在第 4.5 节中任何地方都没有用到。

\oldpage[II]{83}

承认这个引理以后，便存在 $P$ 到
\[
  P_1=\mathbf{P}(g^*(\mathcal{F})\otimes\mathcal{K}^{\otimes(-m)})
\]
中的一个闭浸入 $j_1$，使得 $\mathcal{O}_P(1)$ 同构于
$j_1^*(\mathcal{O}_{P_1}(1))$（\egaref{II.4.1.2-zh}{4.1.2}）。
另一方面，存在 $P_1$ 到
$P_2=\mathbf{P}(g^*(\mathcal{F}))$ 的一个同构，它把
$\mathcal{O}_{P_1}(1)$ 与
$\mathcal{O}_{P_2}(1)\otimes_Y\mathcal{K}^{\otimes(-m)}$ 等同
（\egaref{II.4.1.4-zh}{4.1.4}）。因此存在闭浸入
$j_2:P\to P_2$，使得 $\mathcal{O}_P(1)$ 同构于
$j_2^*(\mathcal{O}_{P_2}(1))\otimes_Y
\mathcal{K}^{\otimes(-m)}$。最后，$P_2$ 可与
$P_3\times_ZY$ 等同，其中 $P_3=\mathbf{P}(\mathcal{F})$，并且
$\mathcal{O}_{P_2}(1)$ 可与
$\mathcal{O}_{P_3}(1)\otimes_Z\mathcal{O}_Y$ 等同
（\egaref{II.4.1.3-zh}{4.1.3}）。按定义，
$\mathcal{O}_{P_3}(1)$ 对 $Z$ 极丰沛；$\mathcal{K}$ 也如此，故由
(iv)，$\mathcal{O}_{P_2}(1)\otimes_Y\mathcal{K}$ 对 $Z$ 极丰沛。
再由 (i bis)，
$\mathcal{M}=j_2^*(\mathcal{O}_{P_2}(1)\otimes_Y\mathcal{K})$
也极丰沛，并且 $\mathcal{O}_P(1)$ 同构于
$\mathcal{M}\otimes_Y\mathcal{K}^{\otimes(-m-1)}$；证明完毕。
\end{proof}
% Locked French authority: ega2-1-fr.tex, 820505 B,
% SHA-256 84EDBE3E83530AF2959B441796337C9DC21EAFCA6A13114A26778760FBF437AC.
% Sequential source scope: lines 7454-7474, proposition II.4.4.11
% and its proof / printed p.83.

\begin{proposition}[4.4.11]
\label{II.4.4.11-zh}
设 $f:X\to Y$、$f':X'\to Y$ 是两个态射，$X''$ 是和预概形
$X\sqcup X'$，$f''$ 是态射 $X''\to Y$，它在 $X$（相应地，$X'$）
上与 $f$（相应地，$f'$）一致。设 $\mathcal{L}$（相应地，
$\mathcal{L}'$）是一个可逆 $\mathcal{O}_X$-模层（相应地，可逆
$\mathcal{O}_{X'}$-模层），并设 $\mathcal{L}''$ 是可逆
$\mathcal{O}_{X''}$-模层，它在 $X$ 上与 $\mathcal{L}$ 一致，在
$X'$ 上与 $\mathcal{L}'$ 一致。那么 $\mathcal{L}''$ 对 $f''$ 是
极丰沛的，当且仅当 $\mathcal{L}$ 对 $f$ 极丰沛，并且
$\mathcal{L}'$ 对 $f'$ 极丰沛。
\end{proposition}

\begin{proof}
立即可归结到 $Y$ 为仿射的情形。若 $\mathcal{L}''$ 极丰沛，则由
\egaref{II.4.4.10-zh}{4.4.10, (i bis)}，$\mathcal{L}$ 和
$\mathcal{L}'$ 也极丰沛。反过来，若 $\mathcal{L}$ 和
$\mathcal{L}'$ 极丰沛，则由定义 \egaref{II.4.4.2-zh}{4.4.2} 以及
\egaref{II.4.3.6-zh}{4.3.6}，立即得出 $\mathcal{L}''$ 极丰沛。
\end{proof}
% Locked French authority: ega2-1-fr.tex, 820505 B,
% SHA-256 84EDBE3E83530AF2959B441796337C9DC21EAFCA6A13114A26778760FBF437AC.
% Sequential source scope: lines 7476-7507, subsection II.4.5 through
% environment II.4.5.1 / printed p.83.

\subsection{丰沛层}
\label{subsection:II.4.5-zh}

\begin{env}[4.5.1]
\label{II.4.5.1-zh}
给定一个预概形 $X$ 和一个可逆 $\mathcal{O}_X$-模层
$\mathcal{L}$。对任意 $\mathcal{O}_X$-模层 $\mathcal{F}$，在
$\mathcal{L}$ 不会引起混淆时，记
$\mathcal{F}(n)=\mathcal{F}\otimes_{\mathcal{O}_X}
\mathcal{L}^{\otimes n}$（$n\in\mathbf{Z}$）；另一方面，记
\[
  S=\bigoplus_{n\geq0}\Gamma(X,\mathcal{L}^{\otimes n})
\]
（它是 \egaxref{0.5.4.6-zh}{0, 5.4.6} 中定义的环
$\Gamma_\bullet(\mathcal{L})$ 的分次子环）。把 $X$ 看成一个
$\mathbf{Z}$-预概形，并以 $p$ 表示结构态射
$X\to\operatorname{Spec}(\mathbf{Z})$；则 $\mathcal{O}_X$-分次代数
同态
\[
  p^*(\widetilde{S})\longrightarrow
  \bigoplus_{n\geq0}\mathcal{L}^{\otimes n}
\]
与分次环 $S$ 的自同态之间有一一对应
（\egaxref{I.2.2.5-zh}{I, 2.2.5}）。这样与 $S$ 的恒等自同构对应的同态
\[
  \varepsilon:p^*(\widetilde{S})\longrightarrow
  \bigoplus_{n\geq0}\mathcal{L}^{\otimes n}
\]
称为典范同态。与之对应的态射
$G(\varepsilon)\to\operatorname{Proj}(S)$
（\egaref{II.3.7.1-zh}{3.7.1}）也称为典范态射。
\end{env}
% Locked French authority: ega2-1-fr.tex, 820505 B,
% SHA-256 84EDBE3E83530AF2959B441796337C9DC21EAFCA6A13114A26778760FBF437AC.
% Sequential source scope: lines 7509-7583, theorem II.4.5.2
% and its proof / printed pp.83-84.

\begin{theorem}[4.5.2]
\label{II.4.5.2-zh}
设 $X$ 是一个拟紧概形，或者是一个底空间为诺特空间的预概形；
$\mathcal{L}$ 是一个可逆 $\mathcal{O}_X$-模层，$S$ 是分次环
$\bigoplus_{n\geq0}\Gamma(X,\mathcal{L}^{\otimes n})$。下列条件等价：
\begin{enumerate}
  \item[\rm{a)}] 当 $f$ 遍历 $S_+$ 的齐次元素时，各 $X_f$ 构成
  $X$ 的拓扑的一组基。
  \item[\rm{a')}] 当 $f$ 遍历 $S_+$ 的齐次元素时，其中仿射的
  $X_f$ 构成 $X$ 的一个覆盖。
  \item[\rm{b)}] 典范态射
  $G(\varepsilon)\to\operatorname{Proj}(S)$
  （\egaref{II.4.5.1-zh}{4.5.1}）处处有定义，并且是稠密开浸入。

\oldpage[II]{84}

  \item[\rm{b')}] 典范态射
  $G(\varepsilon)\to\operatorname{Proj}(S)$ 处处有定义，并且把 $X$
  的底空间同胚到 $\operatorname{Proj}(S)$ 的一个子空间上。
  \item[\rm{c)}] 对任意拟凝聚 $\mathcal{O}_X$-模层
  $\mathcal{F}$，以 $\mathcal{F}_n$ 表示由 $\mathcal{F}(n)$ 在
  $X$ 上的截面所生成的 $\mathcal{F}(n)$ 的
  $\mathcal{O}_X$-子模层；则当 $n>0$ 遍历各整数时，
  $\mathcal{F}$ 是各子 $\mathcal{O}_X$-模层
  $\mathcal{F}_n(-n)$ 的和。
  \item[\rm{c')}] 对 $\mathcal{O}_X$ 的任意拟凝聚理想层，性质 c)
  成立。
\end{enumerate}
此外，若 $(f_\alpha)$ 是 $S_+$ 的一族齐次元素，并且各
$X_{f_\alpha}$ 都仿射，则典范态射
$X\to\operatorname{Proj}(S)$ 在 $\bigcup_\alpha X_{f_\alpha}$
上的限制，是 $\bigcup_\alpha X_{f_\alpha}$ 到
$\bigcup_\alpha(\operatorname{Proj}(S))_{f_\alpha}$ 上的同构。
\end{theorem}

\begin{proof}
显然 b) 蕴含 b')；由公式 \egaref{II.3.7.3.1-zh}{3.7.3.1}，b') 蕴含
a)，这里要注意 $\varepsilon^\flat$ 是恒等映射。条件 a) 蕴含 a')：
每个 $x\in X$ 都有一个仿射邻域 $U$，使
$\mathcal{L}|U$ 同构于 $\mathcal{O}_X|U$；若
$f\in\Gamma(X,\mathcal{L}^{\otimes n})$ 满足
$x\in X_f\subset U$，则 $X_f$ 也是其中满足
$(f|U)(x')\neq0$ 的各点 $x'\in U$ 的集合，因而是仿射开集
（\egaxref{I.1.3.6-zh}{I, 1.3.6}）。为证明 a') 蕴含 b)，只须证明
命题的最后一个断言，并再核验：若 $X=\bigcup_\alpha X_{f_\alpha}$，
则 \egaref{II.3.8.2-zh}{3.8.2} 的条件 (iv) 成立。后一点立即由
\egaxref{I.9.3.1-zh}{I, 9.3.1, (i)} 得出。至于
\egaref{II.4.5.2-zh}{4.5.2} 的最后一个断言，因为
$X_{f_\alpha}$ 是
$(\operatorname{Proj}(S))_{f_\alpha}$ 沿态射
$G(\varepsilon)\to\operatorname{Proj}(S)$ 的逆像，只须应用
\egaxref{I.9.3.2-zh}{I, 9.3.2}。因此 a)、a')、b)、b') 等价。

为证明 a') 蕴含 c)，注意：若 $X_f$ 仿射（其中 $f\in S_k$），则
$\mathcal{F}|X_f$ 由它在 $X_f$ 上的截面生成
（\egaxref{I.1.3.9-zh}{I, 1.3.9}）。另一方面，由
\egaxref{I.9.3.1-zh}{I, 9.3.1, (ii)}，这样的截面 $s$ 具有形式
$(t|X_f)\otimes(f|X_f)^{-m}$，其中
$t\in\Gamma(X,\mathcal{F}(km))$。按定义，$t$ 也是
$\mathcal{F}_{km}$ 的一个截面，故 $s$ 确是
$\mathcal{F}_{km}(-km)$ 在 $X_f$ 上的一个截面；这证明了 c)。
显然 c) 蕴含 c')；只剩证明 c') 蕴含 a)。设 $U$ 是 $x\in X$ 的一个
开邻域，并设 $\mathcal{J}$ 是 $\mathcal{O}_X$ 的一个拟凝聚理想层，
它定义 $X$ 的一个以 $X-U$ 为底空间的闭子预概形
（\egaxref{I.5.2.1-zh}{I, 5.2.1}）。假设 c') 蕴含存在一个整数
$n>0$ 以及 $\mathcal{J}(n)$ 在 $X$ 上的一个截面 $f$，使
$f(x)\neq0$。但显然 $f\in S_n$，并且 $x\in X_f\subset U$；这证明了
a)。
\end{proof}
% Locked French authority: ega2-1-fr.tex, 820505 B,
% SHA-256 84EDBE3E83530AF2959B441796337C9DC21EAFCA6A13114A26778760FBF437AC.
% Sequential source scope: lines 7585-7632, definition II.4.5.3,
% corollary II.4.5.4 and its proof / printed pp.84-85.
% EGA2-ZH-DEF-0010 is retained source-literally pending guardian disposition.

当 $X$ 是一个底空间为诺特空间的预概形时，
\egaref{II.4.5.2-zh}{4.5.2} 中的等价条件推出 $X$ 是概形，
因为根据 \egaref{II.4.5.2-zh}{4.5.2, b)}，它同构于概形
$S=\operatorname{Proj}(A)$ 的一个子预概形。

\begin{definition}[4.5.3]
\label{II.4.5.3-zh}
若 $X$ 是一个拟紧概形，并且满足
\egaref{II.4.5.2-zh}{4.5.2} 中的等价条件，则称一个可逆
$\mathcal{O}_X$-模层 $\mathcal{L}$ 是丰沛的。
\end{definition}

由判别法 \egaref{II.4.5.2-zh}{4.5.2, a)} 显然可得：若
$\mathcal{L}$ 是丰沛 $\mathcal{O}_X$-模，则对 $X$ 的任意开子集
$U$，$\mathcal{L}|U$ 是丰沛 $(\mathcal{O}_X|U)$-模。

由 \egaref{II.4.5.2-zh}{4.5.2} 的证明可得，仿射的各
$X_f$ 已构成 $X$ 的拓扑的一组基。此外：

\begin{corollary}[4.5.4]
\label{II.4.5.4-zh}
设 $\mathcal{L}$ 是一个丰沛 $\mathcal{O}_X$-模。对 $X$ 的任意有限
子空间 $Z$ 及 $Z$ 的任意邻域 $U$，存在一个整数 $n$ 和一个
$f\in\Gamma(X,\mathcal{L}^{\otimes n})$，使 $X_f$ 是 $Z$ 的一个包含于
$U$ 的仿射邻域。
\end{corollary}

\begin{proof}
\oldpage[II]{85}

根据 \egaref{II.4.5.2-zh}{4.5.2, b)}，只须证明：对
$\operatorname{Proj}(S)$ 的任意有限子集 $Z'$ 以及 $Z'$ 的任意开
邻域 $U$，存在一个齐次元 $f\in S_+$，使
$Z'\subset(\operatorname{Proj}(S))_f\subset U$
（\egaref{II.2.4.1-zh}{2.4.1}）。由定义，$U$ 在
$\operatorname{Proj}(S)$ 中的补集是形如
$V_+(\mathfrak{I})$ 的闭集 $Y$，其中 $\mathfrak{I}$ 是 $S$ 的一个不包含
$S_+$ 的分次理想（\egaref{II.2.3.2-zh}{2.3.2}）；另一方面，
$Z'$ 的各点按定义是 $S_+$ 的不包含 $\mathfrak{I}$ 的分次素理想
$\mathfrak{p}_i$（\egaref{II.2.3.1-zh}{2.3.1}）。因此存在一个
$f\in\mathfrak{I}$，它不属于任何 $\mathfrak{p}_i$（Bourbaki,
\emph{Alg. comm.}, chap.~II, \S~1, n\textsuperscript{o}~1, prop.~2）；
由于各 $\mathfrak{p}_i$ 都是分次的，\emph{loc. cit.} 中的论证表明，
甚至可以假设 $f$ 是齐次的；这个元便满足所需条件。
\end{proof}
% Locked French authority: ega2-1-fr.tex, 820505 B,
% SHA-256 84EDBE3E83530AF2959B441796337C9DC21EAFCA6A13114A26778760FBF437AC.
% Sequential source scope: lines 7634-7685, proposition II.4.5.5
% and its proof / printed p.85.

\begin{proposition}[4.5.5]
\label{II.4.5.5-zh}
假设 $X$ 是一个拟紧概形，或者是一个底空间为诺特空间的
预概形。\egaref{II.4.5.2-zh}{4.5.2} 中的条件 a) 至 c') 还等价于
下列条件：
\begin{enumerate}
  \item[\rm{d)}] 对任意有限型拟凝聚 $\mathcal{O}_X$-模
  $\mathcal{F}$，存在一个整数 $n_0$，使对任意
  $n\geq n_0$，$\mathcal{F}(n)$ 都由它在 $X$ 上的截面生成。
  \item[\rm{d')}] 对任意有限型拟凝聚 $\mathcal{O}_X$-模
  $\mathcal{F}$，存在两个整数 $n>0$、$k>0$，使 $\mathcal{F}$ 同构于
  $\mathcal{O}_X$-模 $\mathcal{L}^{\otimes(-n)}\otimes\mathcal{O}_X^k$ 的一个商。
  \item[\rm{d'')}] 性质 d') 对 $\mathcal{O}_X$ 的任意有限型拟凝聚
  理想层成立。
\end{enumerate}
\end{proposition}

\begin{proof}
由于 $X$ 是拟紧的，若一个有限型拟凝聚 $\mathcal{O}_X$-模
$\mathcal{F}$ 使得 $\mathcal{F}(n)$（它为有限型）由其在 $X$ 上的截面
生成，则 $\mathcal{F}(n)$ 已由这些截面中的有限个生成
（\egaxref{0.5.2.3-zh}{0, 5.2.3}）。因此 d) 推出 d')，而 d') 显然推出
d'')。由于任意拟凝聚 $\mathcal{O}_X$-模 $\mathcal{G}$ 是它的各有限型
$\mathcal{O}_X$-子模的归纳极限（\egaxref{I.9.4.9-zh}{I, 9.4.9}），
为了核验 \egaref{II.4.5.2-zh}{4.5.2} 中的条件 c')，只须对
$\mathcal{O}_X$ 的有限型拟凝聚理想层核验它，故 d'') 推出 c')。
只剩下证明：若 $\mathcal{L}$ 丰沛，则性质 d) 成立。取 $X$ 的一个
有限覆盖 $X_{f_i}$（$f_i\in S_{n_i}$），可以假设它们都是仿射的；
把各 $f_i$ 替换为适当的幂（这不改变各 $X_{f_i}$），可以假设所有
$n_i$ 都等于同一整数 $m$。由假设，层 $\mathcal{F}|X_{f_i}$ 为有限型，
因而由它在 $X_{f_i}$ 上的有限个截面 $h_{ij}$ 生成
（\egaxref{I.1.3.13-zh}{I, 1.3.13}）；因此存在一个整数 $k_0$，
使对每个数对 $(i,j)$，截面 $h_{ij}\otimes f_i^{\otimes k_0}$ 都能延拓为
$\mathcal{F}(k_0m)$ 在 $X$ 上的一个截面
（\egaxref{I.9.3.1-zh}{I, 9.3.1}）。更进一步，对任意 $k\geq k_0$，
$h_{ij}\otimes f_i^{\otimes k}$ 都能延拓为 $\mathcal{F}(km)$ 在 $X$ 上的
截面，因而对这些 $k$，$\mathcal{F}(km)$ 由它在 $X$ 上的截面生成。
对任意满足 $0<p<m$ 的 $p$，$\mathcal{F}(p)$ 也是有限型的，故存在
一个整数 $k_p$，使对 $k\geq k_p$，
$\mathcal{F}(p)(km)=\mathcal{F}(p+km)$ 由它在 $X$ 上的截面生成。
取 $n_0$ 大于所有 $k_pm$，便得对任意 $n\geq n_0$，
$\mathcal{F}(n)$ 都由它在 $X$ 上的截面生成，因为这样的 $n$ 可写成
$n=km+p$，其中 $k\geq k_p$ 且 $0\leq p<m$。
\end{proof}
% Locked French authority: ega2-1-fr.tex, 820505 B,
% SHA-256 84EDBE3E83530AF2959B441796337C9DC21EAFCA6A13114A26778760FBF437AC.
% Sequential source scope: lines 7687-7759, proposition II.4.5.6,
% corollaries II.4.5.7-II.4.5.8 and remark II.4.5.9 / printed pp.85-86.

\begin{proposition}[4.5.6]
\label{II.4.5.6-zh}
设 $X$ 是一个拟紧概形，$\mathcal{L}$ 是一个可逆 $\mathcal{O}_X$-模层。
\begin{enumerate}
  \item[\rm{(i)}] 设 $n$ 是一个 $>0$ 的整数。$\mathcal{L}$ 丰沛的充要条件是
  $\mathcal{L}^{\otimes n}$ 丰沛。
  \item[\rm{(ii)}] 设 $\mathcal{L}'$ 是一个可逆 $\mathcal{O}_X$-模层，并且对
  任意 $x\in X$，存在一个整数 $n>0$

\oldpage[II]{86}

  和 $\mathcal{L}'^{\otimes n}$ 在 $X$ 上的一个截面 $s'$，使
  $s'(x)\neq0$。那么，若 $\mathcal{L}$ 丰沛，则
  $\mathcal{L}\otimes\mathcal{L}'$ 也丰沛。
\end{enumerate}
\end{proposition}

\begin{proof}
性质 (i) 显然由 \egaref{II.4.5.2-zh}{4.5.2} 的判别条件 a) 得出，
因为 $X_{f^{\otimes n}}=X_f$。另一方面，若 $\mathcal{L}$ 丰沛，则对
任意 $x\in X$ 及 $x$ 的任意邻域 $U$，存在 $m>0$ 和
$f\in\Gamma(X,\mathcal{L}^{\otimes m})$，使
$x\in X_f\subset U$（\egaref{II.4.5.2-zh}{4.5.2, a)}）；若
$f'\in\Gamma(X,\mathcal{L}'^{\otimes n})$ 满足 $f'(x)\neq0$，则对
$s=f^{\otimes n}\otimes f'^{\otimes m}\in
\Gamma(X,(\mathcal{L}\otimes\mathcal{L}')^{\otimes mn})$，有
$s(x)\neq0$，从而 $x\in X_s\subset X_f\subset U$。这证明了
$\mathcal{L}\otimes\mathcal{L}'$ 丰沛
（\egaref{II.4.5.2-zh}{4.5.2, a)}）。
\end{proof}

\begin{corollary}[4.5.7]
\label{II.4.5.7-zh}
两个丰沛 $\mathcal{O}_X$-模层的张量积是丰沛的。
\end{corollary}

\begin{corollary}[4.5.8]
\label{II.4.5.8-zh}
设 $\mathcal{L}$ 是一个丰沛 $\mathcal{O}_X$-模层，$\mathcal{L}'$ 是一个可逆
$\mathcal{O}_X$-模层；则存在一个整数 $n_0>0$，使对 $n\geq n_0$，
$\mathcal{L}^{\otimes n}\otimes\mathcal{L}'$ 丰沛，并且由它在 $X$ 上的截面生成。
\end{corollary}

\begin{proof}
事实上，由 \egaref{II.4.5.5-zh}{4.5.5} 可得，存在一个整数 $m_0$，
使对 $m\geq m_0$，$\mathcal{L}^{\otimes m}\otimes\mathcal{L}'$ 由它在 $X$
上的截面生成；根据 \egaref{II.4.5.6-zh}{4.5.6}，因而可取
$n_0=m_0+1$。
\end{proof}

\begin{remark}[4.5.9]
\label{II.4.5.9-zh}
设 $P=\mathrm{H}^1(X,\mathcal{O}_X^*)$ 是可逆 $\mathcal{O}_X$-模层的类群
（\egaxref{0.5.4.7-zh}{0, 5.4.7}），并设 $P^+$ 是 $P$ 中由丰沛层的类
所构成的子集。假设 $P^+$ 非空。由 \egaref{II.4.5.7-zh}{4.5.7}
和 \egaref{II.4.5.8-zh}{4.5.8} 可得
\[
  P^++P^+\subset P^+
  \qquad\text{且}\qquad
  P^+-P^+=P
\]
换言之，对 $P$ 上一个与其群结构相容的预序结构，
$P^+\cup\{0\}$ 是 $P$ 中的正元素集合；并且根据
\egaref{II.4.5.8-zh}{4.5.8}，这个预序甚至是阿基米德的。因此有时也用
“正层”代替丰沛层，并把这种层的逆称为“负层”（本文不采用这种术语）。
\end{remark}
% Locked French authority: ega2-1-fr.tex, 820505 B,
% SHA-256 84EDBE3E83530AF2959B441796337C9DC21EAFCA6A13114A26778760FBF437AC.
% Sequential source scope: lines 7761-7833, proposition II.4.5.10,
% its proof and proof II.4.5.10.1 of lemma II.4.4.10.1 / printed pp.86-87.

\begin{proposition}[4.5.10]
\label{II.4.5.10-zh}
设 $Y$ 是一个仿射概形，$q:X\to Y$ 是一个分离拟紧态射，
$\mathcal{L}$ 是一个可逆 $\mathcal{O}_X$-模层。
\begin{enumerate}
  \item[\rm{(i)}] 若 $\mathcal{L}$ 对 $q$ 极丰沛，则 $\mathcal{L}$ 丰沛。
  \item[\rm{(ii)}] 再假设态射 $q$ 是有限型的。那么，$\mathcal{L}$ 丰沛的
  充要条件是它具有下列两个等价性质之一：
  \begin{enumerate}
    \item[\rm{e)}] 存在 $n_0>0$，使对任意整数 $n\geq n_0$，
    $\mathcal{L}^{\otimes n}$ 对 $q$ 极丰沛。
    \item[\rm{e')}] 存在 $n>0$，使 $\mathcal{L}^{\otimes n}$ 对 $q$ 极丰沛。
  \end{enumerate}
\end{enumerate}
\end{proposition}

\begin{proof}
第一个断言由极丰沛 $\mathcal{O}_X$-模层的定义
（\egaref{II.4.4.2-zh}{4.4.2}）得出：若 $A$ 是 $Y$ 的环，则存在一个
$A$-模 $E$ 和一个满同态
\[
  \psi:q^*((\mathbf{S}(E))^{\sim})\longrightarrow
  \bigoplus_{n\geq0}\mathcal{L}^{\otimes n}
\]
使 $i=r_{\mathcal{L},\psi}$ 是一个处处有定义的浸入
$X\to P=\mathbf{P}(\widetilde{E})$，并且
$\mathcal{L}=i^*(\mathcal{O}_P(1))$。由于当 $f$ 遍历
$(\mathbf{S}(E))_+$ 中的齐次元时，各 $D_+(f)$ 构成 $P$ 的拓扑的一组基，
并且由 \egaref{II.3.7.3.1-zh}{3.7.3.1} 有
$i^{-1}(D_+(f))=X_{\psi^\flat(f)}$，可见
\egaref{II.4.5.2-zh}{4.5.2} 的条件 a) 成立，故 $\mathcal{L}$ 丰沛。

现在证明：若 $q$ 是有限型的且 $\mathcal{L}$ 丰沛，则它满足条件 e)。
首先，由 \egaref{II.4.5.2-zh}{4.5.2} 的判别条件 b) 和
\egaref{II.4.4.1-zh}{4.4.1, (i)} 可得，存在

\oldpage[II]{87}

一个整数 $k_0>0$，使 $\mathcal{L}^{\otimes k_0}$ 对 $q$ 极丰沛。另一方面，
根据 \egaref{II.4.5.5-zh}{4.5.5}，存在一个整数 $m_0$，使对
$m\geq m_0$，$\mathcal{L}^{\otimes m}$ 由它在 $X$ 上的截面生成。
置 $n_0=k_0+m_0$；若 $n\geq n_0$，则可写成 $n=k_0+m$，其中
$m\geq m_0$，从而
$\mathcal{L}^{\otimes n}
=\mathcal{L}^{\otimes k_0}\otimes\mathcal{L}^{\otimes m}$。
由于 $\mathcal{L}^{\otimes m}$ 由它在 $X$ 上的截面生成，根据判别法
\egaref{II.4.4.8-zh}{4.4.8} 和 \egaref{II.3.4.7-zh}{3.4.7}，
$\mathcal{L}^{\otimes n}$ 对 $q$ 极丰沛。最后，e) 推出 e') 是显然的，而 e')
根据 (i) 和 \egaref{II.4.5.6-zh}{4.5.6, (i)} 推出 $\mathcal{L}$ 丰沛。

\begin{env}[4.5.10.1]
\label{II.4.5.10.1-zh}
\emph{引理 \egaref{II.4.4.10.1-zh}{4.4.10.1} 的证明。}置
$\mathcal{E}(n)=\mathcal{E}\otimes\mathcal{K}^{\otimes n}$；由于 $g$ 是分离的
（\egaref{II.4.4.2-zh}{4.4.2}），\egaref{II.3.4.8-zh}{3.4.8} 的论证
适用，并把问题归结为证明：当 $n$ 充分大时，典范同态
$g^*(g_*(\mathcal{E}(n)))\to\mathcal{E}(n)$ 是满的。此外，由于 $Z$ 拟紧，
\egaref{II.3.4.6-zh}{3.4.6} 的论证把命题归结到 $Z$ 为仿射的情形。
但此时，根据 \egaref{II.4.5.10-zh}{4.5.10, (i)}，$\mathcal{K}$ 丰沛，
而结论由 \egaref{II.4.5.5-zh}{4.5.5, d')} 得出。
\end{env}
\end{proof}
% Locked French authority: ega2-1-fr.tex, 820505 B,
% SHA-256 84EDBE3E83530AF2959B441796337C9DC21EAFCA6A13114A26778760FBF437AC.
% Sequential source scope: lines 7835-7860, corollary II.4.5.11,
% its proof and remark II.4.5.12 / printed p.87.

\begin{corollary}[4.5.11]
\label{II.4.5.11-zh}
设 $Y$ 是一个仿射概形，$q:X\to Y$ 是一个分离的有限型态射，
$\mathcal{L}$ 是一个丰沛 $\mathcal{O}_X$-模层，$\mathcal{L}'$ 是一个可逆
$\mathcal{O}_X$-模层。存在一个整数 $n_0$，使对 $n\geq n_0$，
$\mathcal{L}^{\otimes n}\otimes\mathcal{L}'$ 对 $q$ 极丰沛。
\end{corollary}

\begin{proof}
事实上，存在 $m_0$，使对 $m\geq m_0$，
$\mathcal{L}^{\otimes m}\otimes\mathcal{L}'$ 由它在 $X$ 上的截面生成
（\egaref{II.4.5.8-zh}{4.5.8}）；另一方面，存在 $k_0$，使对
$k\geq k_0$，$\mathcal{L}^{\otimes k}$ 对 $q$ 极丰沛。因此，对
$k\geq k_0$，$\mathcal{L}^{\otimes(k+m_0)}\otimes\mathcal{L}'$ 对 $q$
极丰沛（\egaref{II.4.4.8-zh}{4.4.8} 和
\egaref{II.3.4.7-zh}{3.4.7}）。
\end{proof}

\begin{remark}[4.5.12]
\label{II.4.5.12-zh}
目前尚不知道：若一个 $\mathcal{O}_X$-模层 $\mathcal{L}$ 满足
$\mathcal{L}^{\otimes n}$ 极丰沛（相对于该态射），是否能推出
$\mathcal{L}^{\otimes(n+1)}$ 也具有同一性质。
\end{remark}
% Locked French authority: ega2-1-fr.tex, 820505 B,
% SHA-256 84EDBE3E83530AF2959B441796337C9DC21EAFCA6A13114A26778760FBF437AC.
% Sequential source scope: lines 7862-8017, proposition II.4.5.13,
% lemma II.4.5.13.1 and the complete proof / printed pp.87-89.

\begin{proposition}[4.5.13]
\label{II.4.5.13-zh}
设 $X$ 是一个拟紧预概形，$Z$ 是 $X$ 的一个闭子预概形，它由
$\mathcal{O}_X$ 的一个幂零拟凝聚理想层 $\mathcal{J}$ 定义，$j$ 是典范单射
$Z\to X$。一个可逆 $\mathcal{O}_X$-模层 $\mathcal{L}$ 丰沛的充要条件是
$\mathcal{L}'=j^*(\mathcal{L})$ 是一个丰沛 $\mathcal{O}_Z$-模层。
\end{proposition}

\begin{proof}
这个条件是必要的。事实上，对 $\mathcal{L}^{\otimes n}$ 在 $X$ 上的任意一个
截面 $f$，设 $f'$ 是它的典范像 $f\otimes1$；这是
$\mathcal{L}'^{\otimes n}=\mathcal{L}^{\otimes n}
\otimes_{\mathcal{O}_X}(\mathcal{O}_X/\mathcal{J})$ 在空间 $Z$ 上的一个截面
（其底空间与 $X$ 的相同）。显然 $X_f=Z_{f'}$，因此
\egaref{II.4.5.2-zh}{4.5.2} 的判别法 a) 表明 $\mathcal{L}'$ 丰沛。

为证明这个条件也是充分的，先注意：考虑预概形的有限序列
$X_k=(X,\mathcal{O}_X/\mathcal{J}^{k+1})$，可以归结到 $\mathcal{J}^2=0$ 的情形；
这个序列中的每一个预概形都是下一个预概形的闭子预概形，并由
一个平方为零的理想层定义。此外，$X$ 是一个概形，因为由假设
$X_{\mathrm{red}}$ 是概形（\egaref{II.4.5.3-zh}{4.5.3} 和
\egaxref{I.5.5.1-zh}{I, 5.5.1}）。由 \egaref{II.4.5.2-zh}{4.5.2} 的判别法 a)，
只须证明下列引理。

\begin{lemma}[4.5.13.1]
\label{II.4.5.13.1-zh}
在 \egaref{II.4.5.13-zh}{4.5.13} 的假设下，再假设 $\mathcal{J}$ 的平方为零；
设 $\mathcal{L}$ 是一个可逆 $\mathcal{O}_X$-模层，并设 $g$ 是
$\mathcal{L}'^{\otimes n}$ 在 $Z$ 上的一个截面，使 $Z_g$ 是仿射的。
那么存在一个整数 $m>0$，使 $g^{\otimes m}$ 是
$\mathcal{L}^{\otimes nm}$ 在 $X$ 上的某个截面 $f$ 的典范像。
\end{lemma}

有下列 $\mathcal{O}_X$-模层的正合序列：
\[
  0\longrightarrow\mathcal{J}(n)\longrightarrow
  \mathcal{O}_X(n)=\mathcal{L}^{\otimes n}\longrightarrow
  \mathcal{O}_Z(n)=\mathcal{L}'^{\otimes n}\longrightarrow0
\]

\oldpage[II]{88}

这是因为 $\mathcal{F}(n)$ 是关于 $\mathcal{F}$ 的正合函子；由此得到上同调正合序列
\[
  0\longrightarrow\Gamma(X,\mathcal{J}(n))\longrightarrow
  \Gamma(X,\mathcal{L}^{\otimes n})\longrightarrow
  \Gamma(X,\mathcal{L}'^{\otimes n})\xrightarrow{\partial}
  \mathrm{H}^1(X,\mathcal{J}(n))
\]
它特别把 $g$ 对应到一个元
$\partial g\in\mathrm{H}^1(X,\mathcal{J}(n))$。

注意，由于 $\mathcal{J}^2=0$，可以把 $\mathcal{J}$ 看成一个拟凝聚
$\mathcal{O}_Z$-模层，并且对任意 $k$ 都有
$\mathcal{L}'^{\otimes k}\otimes_{\mathcal{O}_Z}\mathcal{J}(n)
=\mathcal{J}(n+k)$。因此，对任意截面
$s\in\Gamma(X,\mathcal{L}'^{\otimes k})$，与 $s$ 作张量乘法给出一个
$\mathcal{O}_Z$-模层同态 $\mathcal{J}(n)\to\mathcal{J}(n+k)$，从而给出上同调群同态
$\mathrm{H}^i(X,\mathcal{J}(n))\to
\mathrm{H}^i(X,\mathcal{J}(n+k))$。

现在证明：对充分大的 $m>0$，有
\begin{equation}
\label{II.4.5.13.2-zh}
  g^{\otimes m}\otimes\partial g=0
\tag{4.5.13.2}
\end{equation}
事实上，$Z_g$ 是 $Z$ 的一个仿射开子集，故当把 $\mathcal{J}(n)$ 看成
一个 $\mathcal{O}_Z$-模层时，有 $\mathrm{H}^1(Z_g,\mathcal{J}(n))=0$
（\egaxref{I.5.1.9.2-zh}{I, 5.1.9.2}）。特别地，置 $g'=g|Z_g$，并考虑它在映射
$\partial:\Gamma(Z_g,\mathcal{L}'^{\otimes n})\to
\mathrm{H}^1(Z_g,\mathcal{J}(n))$ 下的像，便有 $\partial g'=0$。为具体说明
这个关系，注意在次数 $1$ 上，一个阿贝尔群层的上同调与它的
Čech 上同调相同（G, II, 5.9）。因此，为构造 $\partial g$，要取 $X$
的一个充分细的开覆盖 $(U_\alpha)$；可以假设这个覆盖是有限的，并由
仿射开子集组成。对每个 $\alpha$，取一个截面
$g_\alpha\in\Gamma(U_\alpha,\mathcal{L}^{\otimes n})$，使它在
$\Gamma(U_\alpha,\mathcal{L}'^{\otimes n})$ 中的典范像为 $g|U_\alpha$，再考虑上闭链
$(g_{\alpha|\beta}-g_{\beta|\alpha})$ 的类；这里 $g_{\alpha|\beta}$ 表示
$g_\alpha$ 在 $U_\alpha\cap U_\beta$ 上的限制，而该上闭链取值于
$\mathcal{J}(n)$。此外，可以假设 $\partial g'$ 以同样方式由覆盖
$U_\alpha\cap Z_g$ 和各限制 $g_\alpha|(U_\alpha\cap Z_g)$ 计算
（必要时把 $(U_\alpha)$ 替换为一个更细的覆盖）。于是，关系 $\partial g'=0$
意味着：对每个 $\alpha$，存在一个截面
$h_\alpha\in\Gamma(U_\alpha\cap Z_g,\mathcal{J}(n))$，使
\[
  (g_{\alpha|\beta}-g_{\beta|\alpha})
  |(U_\alpha\cap U_\beta\cap Z_g)
  =h_{\alpha|\beta}-h_{\beta|\alpha}
\]
其中 $h_{\alpha|\beta}$ 表示 $h_\alpha$ 在
$U_\alpha\cap U_\beta\cap Z_g$ 上的限制（G, II, 5.11）。于是，存在一个整数
$m>0$，使对每个 $\alpha$，$g^{\otimes m}\otimes h_\alpha$ 都是某个截面
$t_\alpha\in\Gamma(U_\alpha,\mathcal{J}(n+nm))$ 在
$U_\alpha\cap Z_g$ 上的限制（\egaxref{I.9.3.1-zh}{I, 9.3.1}）。
因此，对每一对指标都有
\[
  g^{\otimes m}\otimes(g_{\alpha|\beta}-g_{\beta|\alpha})
  =t_{\alpha|\beta}-t_{\beta|\alpha},
\]
这就证明了 \egaref{II.4.5.13.2-zh}{4.5.13.2}。

另一方面，注意若 $s\in\Gamma(X,\mathcal{O}_Z(p))$、
$t\in\Gamma(X,\mathcal{O}_Z(q))$，则在群
$\mathrm{H}^1(X,\mathcal{J}(p+q))$ 中有
\begin{equation}
\label{II.4.5.13.3-zh}
  \partial(s\otimes t)=(\partial s)\otimes t+s\otimes(\partial t).
\tag{4.5.13.3}
\end{equation}

事实上，为计算等式两边，仍可取 $X$ 的一个开覆盖 $(U_\alpha)$；
对每个 $\alpha$，分别取截面
$s_\alpha\in\Gamma(U_\alpha,\mathcal{O}_X(p))$ 和
$t_\alpha\in\Gamma(U_\alpha,\mathcal{O}_X(q))$，使它们在
$\Gamma(U_\alpha,\mathcal{O}_Z(p))$ 和
$\Gamma(U_\alpha,\mathcal{O}_Z(q))$ 中的典范像分别为 $s|U_\alpha$ 和
$t|U_\alpha$。于是，\egaref{II.4.5.13.3-zh}{4.5.13.3} 由下列关系得出：
\[
  (s_{\alpha|\beta}\otimes t_{\alpha|\beta})
  -(s_{\beta|\alpha}\otimes t_{\beta|\alpha})
  =(s_{\alpha|\beta}-s_{\beta|\alpha})\otimes t_{\alpha|\beta}
  +s_{\beta|\alpha}\otimes(t_{\alpha|\beta}-t_{\beta|\alpha})
\]
其中采用与上文相同的记号。由此对 $k$ 归纳可得
\begin{equation}
\label{II.4.5.13.4-zh}
  \partial(g^{\otimes k})=(kg^{\otimes(k-1)})\otimes(\partial g)
\tag{4.5.13.4}
\end{equation}

\oldpage[II]{89}

由 \egaref{II.4.5.13.2-zh}{4.5.13.2} 和
\egaref{II.4.5.13.4-zh}{4.5.13.4} 可得
$\partial(g^{\otimes(m+1)})=0$；因此，$g^{\otimes(m+1)}$ 是
$\mathcal{L}^{\otimes n(m+1)}$ 在 $X$ 上的某个截面 $f$ 的典范像，
这就完成了 \egaref{II.4.5.13-zh}{4.5.13} 的证明。
\end{proof}
% Locked French authority: ega2-1-fr.tex, 820505 B,
% SHA-256 84EDBE3E83530AF2959B441796337C9DC21EAFCA6A13114A26778760FBF437AC.
% Sequential source scope: lines 8019-8027, corollary II.4.5.14
% and its concluding justification / printed p.89.

\begin{corollary}[4.5.14]
\label{II.4.5.14-zh}
设 $X$ 是一个诺特概形，$j$ 是典范单射
$X_{\mathrm{red}}\to X$。一个可逆 $\mathcal{O}_X$-模层
$\mathcal{L}$ 丰沛的充要条件是 $j^*(\mathcal{L})$ 是一个丰沛
$\mathcal{O}_{X_{\mathrm{red}}}$-模层。
\end{corollary}

这由 \egaxref{I.6.1.6-zh}{I, 6.1.6} 得出。
% Locked French authority: ega2-1-fr.tex, 820505 B,
% SHA-256 84EDBE3E83530AF2959B441796337C9DC21EAFCA6A13114A26778760FBF437AC.
% Sequential source scope: lines 8029-8058, subsection II.4.6,
% definition II.4.6.1 and proposition II.4.6.2 / printed p.89.

\subsection{相对丰沛层}
\label{II.4.6-zh}

\begin{definition}[4.6.1]
\label{II.4.6.1-zh}
设 $f:X\to Y$ 是一个拟紧态射，$\mathcal{L}$ 是一个可逆
$\mathcal{O}_X$-模层。若存在 $Y$ 的一个由仿射开集组成的覆盖
$(U_\alpha)$，使得置 $X_\alpha=f^{-1}(U_\alpha)$ 时，对每个
$\alpha$，$\mathcal{L}|X_\alpha$ 都是一个丰沛
$\mathcal{O}_{X_\alpha}$-模层，则称 $\mathcal{L}$ \emph{相对于 $f$
丰沛}，或\emph{相对于 $Y$ 丰沛}，或称为\emph{$f$-丰沛}、\emph{$Y$-丰沛}
的（若不会同 \egaref{II.4.5.3-zh}{4.5.3} 中定义的概念混淆，甚至简称为
\emph{丰沛}的）。
\end{definition}

一个 $f$-丰沛 $\mathcal{O}_X$-模层的存在性蕴含 $f$ 必为分离态射
（\egaref{II.4.5.3-zh}{4.5.3} 和
\egaxref{I.5.5.5-zh}{I, 5.5.5}）。

\begin{proposition}[4.6.2]
\label{II.4.6.2-zh}
设 $f:X\to Y$ 是一个拟紧态射，$\mathcal{L}$ 是一个可逆
$\mathcal{O}_X$-模层。若 $\mathcal{L}$ 相对于 $f$ 极丰沛，则它相对于
$f$ 丰沛。
\end{proposition}

这由极丰沛层概念的关于 $Y$ 的局部性
（\egaref{II.4.4.5-zh}{4.4.5}）、定义
\egaref{II.4.6.1-zh}{4.6.1} 和判别法
\egaref{II.4.5.10-zh}{4.5.10, (i)} 得出。
% Locked French authority: ega2-1-fr.tex, 820505 B,
% SHA-256 84EDBE3E83530AF2959B441796337C9DC21EAFCA6A13114A26778760FBF437AC.
% Sequential source scope: lines 8060-8127, proposition II.4.6.3
% and its proof / printed pp.89-90.

\begin{proposition}[4.6.3]
\label{II.4.6.3-zh}
设 $f:X\to Y$ 是一个拟紧态射，$\mathcal{L}$ 是一个可逆
$\mathcal{O}_X$-模层，并设 $\mathcal{S}$ 是分次
$\mathcal{O}_Y$-代数
$\bigoplus_{n\geq0}f_*(\mathcal{L}^{\otimes n})$。下列条件等价：
\begin{enumerate}
  \item[\rm{a)}] $\mathcal{L}$ 是 $f$-丰沛的。
  \item[\rm{b)}] $\mathcal{S}$ 是拟凝聚的，并且典范同态
  $\sigma:f^*(\mathcal{S})\to
  \bigoplus_{n\geq0}\mathcal{L}^{\otimes n}$
  （\egaxref{0.4.4.3-zh}{0, 4.4.3}）使得 $Y$-态射
  $r_{\mathcal{L},\sigma}:G(\sigma)\to\operatorname{Proj}(\mathcal{S})=P$
  处处有定义，并且是稠密开浸入。
  \item[\rm{b')}] 态射 $f$ 是分离的，$Y$-态射
  $r_{\mathcal{L},\sigma}$ 处处有定义，并把 $X$ 的底空间同胚到
  $\operatorname{Proj}(\mathcal{S})$ 的一个子空间上。
\end{enumerate}
此外，当这些条件成立时，对每个 $n\in\mathbf{Z}$，在
\egaref{II.3.7.9.1-zh}{3.7.9.1} 中定义的典范同态
\begin{equation}
\label{II.4.6.3.1-zh}
  r_{\mathcal{L},\sigma}^*(\mathcal{O}_P(n))
  \longrightarrow\mathcal{L}^{\otimes n}
\tag{4.6.3.1}
\end{equation}
是同构。

最后，对任意拟凝聚 $\mathcal{O}_X$-模层 $\mathcal{F}$，若置
$\mathcal{M}=\bigoplus_{n\geq0}f_*(\mathcal{F}\otimes
\mathcal{L}^{\otimes n})$，则在
\egaref{II.3.7.9.2-zh}{3.7.9.2} 中定义的典范同态
\begin{equation}
\label{II.4.6.3.2-zh}
  r_{\mathcal{L},\sigma}^*(\widetilde{\mathcal{M}})
  \longrightarrow\mathcal{F}
\tag{4.6.3.2}
\end{equation}
是同构。
\end{proposition}

已经指出，a) 蕴含 $f$ 是分离的，因而 $\mathcal{S}$ 是拟凝聚的
（\egaxref{I.9.2.2-zh}{I, 9.2.2, a)}）。由于
$r_{\mathcal{L},\sigma}$ 是处处有定义的浸入这一事实关于 $Y$ 是局部的，
为证明 a) 蕴含 b)，可以假设 $Y$ 是仿射的且 $\mathcal{L}$ 丰沛；

\oldpage[II]{90}

此时断言由 \egaref{II.4.5.2-zh}{4.5.2, b)} 得出。显然 b) 蕴含
b')；最后，为证明 b') 蕴含 a)，只须取 $Y$ 的一个由仿射开集
$U_\alpha$ 组成的覆盖，并对每个层
$\mathcal{L}|f^{-1}(U_\alpha)$ 应用判别法
\egaref{II.4.5.2-zh}{4.5.2, b')}。

对于最后两个断言，使用 $\sigma^\flat$ 在这里是恒等映射这一事实，以及同态
\egaref{II.3.7.9.1-zh}{3.7.9.1} 和
\egaref{II.3.7.9.2-zh}{3.7.9.2} 的具体描述；由此立即可得
\egaref{II.4.6.3.1-zh}{4.6.3.1} 是同构。对于
\egaref{II.4.6.3.2-zh}{4.6.3.2}，可以归结到 $Y$ 为仿射、因而
$\mathcal{L}$ 丰沛的情形；显然，同态
\egaref{II.4.6.3.2-zh}{4.6.3.2} 是单射，而判别法
\egaref{II.4.5.2-zh}{4.5.2, c)} 表明它是满射，结论随之得出。
% Locked French authority: ega2-1-fr.tex, 820505 B,
% SHA-256 84EDBE3E83530AF2959B441796337C9DC21EAFCA6A13114A26778760FBF437AC.
% Sequential source scope: lines 8129-8175, corollaries II.4.6.4-II.4.6.7
% and their justifications / printed p.90.

\begin{corollary}[4.6.4]
\label{II.4.6.4-zh}
设 $(U_\alpha)$ 是 $Y$ 的一个开覆盖。$\mathcal{L}$ 相对于 $Y$ 丰沛的充要条件是：
对每个 $\alpha$，$\mathcal{L}|f^{-1}(U_\alpha)$ 相对于
$U_\alpha$ 丰沛。
\end{corollary}

事实上，条件 b) 关于 $Y$ 是局部的。

\begin{corollary}[4.6.5]
\label{II.4.6.5-zh}
设 $\mathcal{K}$ 是一个可逆 $\mathcal{O}_Y$-模层。$\mathcal{L}$ 是
$Y$-丰沛的充要条件是 $\mathcal{L}\otimes f^*(\mathcal{K})$ 也是
如此。
\end{corollary}

这是 \egaref{II.4.6.4-zh}{4.6.4} 的显然推论：取各 $U_\alpha$，使对每个
$\alpha$，$\mathcal{K}|U_\alpha$ 都同构于 $\mathcal{O}_Y|U_\alpha$。

\begin{corollary}[4.6.6]
\label{II.4.6.6-zh}
假设 $Y$ 是仿射的；$\mathcal{L}$ 是 $Y$-丰沛的充要条件是
$\mathcal{L}$ 丰沛。
\end{corollary}

这立即由定义 \egaref{II.4.6.1-zh}{4.6.1} 和判别法
\egaref{II.4.6.3-zh}{4.6.3, b)}、
\egaref{II.4.5.2-zh}{4.5.2, b)} 得出，因为在这里，按定义有
$\operatorname{Proj}(\mathcal{S})
=\operatorname{Proj}(\Gamma(Y,\mathcal{S}))$。

\begin{corollary}[4.6.7]
\label{II.4.6.7-zh}
设 $f:X\to Y$ 是一个拟紧态射。假设存在一个拟凝聚
$\mathcal{O}_Y$-模层 $\mathcal{E}$ 和一个 $Y$-态射
$g:X\to P=\mathbf{P}(\mathcal{E})$，它把 $X$ 的底空间同胚到 $P$
的一个子空间上；则 $\mathcal{L}=g^*(\mathcal{O}_P(1))$ 是
$Y$-丰沛的。
\end{corollary}

事实上，可以假设 $Y$ 是仿射的；此时推论由判别法
\egaref{II.4.5.2-zh}{4.5.2, a)}、公式
\egaref{II.3.7.3.1-zh}{3.7.3.1} 和
\egaref{II.4.2.3-zh}{4.2.3} 得出。
% Locked French authority: ega2-1-fr.tex, 820505 B,
% SHA-256 84EDBE3E83530AF2959B441796337C9DC21EAFCA6A13114A26778760FBF437AC.
% Sequential source scope: lines 8177-8220, proposition II.4.6.8
% and its proof / printed pp.90-91.

\begin{proposition}[4.6.8]
\label{II.4.6.8-zh}
设 $X$ 是一个拟紧概形，或者是一个底空间为诺特空间的预概形，
$f:X\to Y$ 是一个分离拟紧态射。一个可逆 $\mathcal{O}_X$-模层
$\mathcal{L}$ 是 $f$-丰沛的充要条件是下列等价条件之一成立：
\begin{enumerate}
  \item[\rm{c)}] 对任意有限型拟凝聚 $\mathcal{O}_X$-模层
  $\mathcal{F}$，存在一个整数 $n_0>0$，使对任意 $n\geq n_0$，
  典范同态
  $\sigma:f^*(f_*(\mathcal{F}\otimes\mathcal{L}^{\otimes n}))\to
  \mathcal{F}\otimes\mathcal{L}^{\otimes n}$ 是满射。
  \item[\rm{c')}] 对 $\mathcal{O}_X$ 的任意有限型拟凝聚理想层
  $\mathcal{J}$，存在一个整数 $n>0$，使典范同态
  $\sigma:f^*(f_*(\mathcal{J}\otimes\mathcal{L}^{\otimes n}))\to
  \mathcal{J}\otimes\mathcal{L}^{\otimes n}$ 是满射。
\end{enumerate}
\end{proposition}

由于 $X$ 拟紧，$f(X)$ 也拟紧，因而存在 $Y$ 的仿射开集给出的
$f(X)$ 的一个有限覆盖 $(U_i)$。为在 $\mathcal{L}$ 为 $f$-丰沛时证明条件
c)，可以把 $Y$ 替换为各 $U_i$，并把 $X$ 替换为各 $f^{-1}(U_i)$；
因为若对每个 $i$ 都得到一个整数 $n_i$，使得只要 $n\geq n_i$，c) 就对
$U_i$、$f^{-1}(U_i)$ 和 $\mathcal{L}|f^{-1}(U_i)$ 成立，那么取各
$n_i$ 中最大的一个作为 $n_0$，就得到 c) 对 $Y$、$X$ 和
$\mathcal{L}$ 成立。当 $Y$ 是仿射的时，考虑到
\egaref{II.4.6.6-zh}{4.6.6}，条件 c) 由
\egaref{II.4.5.5-zh}{4.5.5, d)} 得出。显然 c) 蕴含 c')。
最后，为证明 c') 蕴含 $\mathcal{L}$ 是 $f$-丰沛的，仍可以归结到
$Y$ 为仿射的情形：事实上，$\mathcal{O}_X|f^{-1}(U_i)$ 的任意有限型
拟凝聚理想层 $\mathcal{J}_i$ 都是 $\mathcal{O}_X$ 的某个有限型凝聚
理想层的限制（\egaxref{I.9.4.7-zh}{I, 9.4.7}），而假设 c') 蕴含

\oldpage[II]{91}

$\mathcal{J}_i\otimes(\mathcal{L}^{\otimes n}|f^{-1}(U_i))$
由它的截面生成（考虑到 \egaxref{I.9.2.2-zh}{I, 9.2.2} 和
\egaref{II.3.4.7-zh}{3.4.7}）；所以只须应用判别法
\egaref{II.4.5.5-zh}{4.5.5, d'')}。
% Locked French authority: ega2-1-fr.tex, 820505 B,
% SHA-256 84EDBE3E83530AF2959B441796337C9DC21EAFCA6A13114A26778760FBF437AC.
% Sequential source scope: lines 8222-8285, propositions II.4.6.9-II.4.6.11,
% corollaries II.4.6.10-II.4.6.12 and their proofs / printed p.91.

\begin{proposition}[4.6.9]
\label{II.4.6.9-zh}
设 $f:X\to Y$ 是一个拟紧态射，$\mathcal{L}$ 是一个可逆
$\mathcal{O}_X$-模层。
\begin{enumerate}
  \item[\rm{(i)}] 设 $n$ 是一个 $>0$ 的整数。$\mathcal{L}$ 是
  $f$-丰沛的充要条件是 $\mathcal{L}^{\otimes n}$ 是 $f$-丰沛的。
  \item[\rm{(ii)}] 设 $\mathcal{L}'$ 是一个可逆
  $\mathcal{O}_X$-模层，并假设存在一个整数 $n>0$，使典范同态
  $\sigma:f^*(f_*(\mathcal{L}'^{\otimes n}))\to
  \mathcal{L}'^{\otimes n}$ 是满射。那么，若 $\mathcal{L}$ 是
  $f$-丰沛的，则 $\mathcal{L}\otimes\mathcal{L}'$ 也是如此。
\end{enumerate}
\end{proposition}

事实上，立即可归结到 $Y$ 为仿射的情形，此时命题是
\egaref{II.4.5.6-zh}{4.5.6} 的直接推论。

\begin{corollary}[4.6.10]
\label{II.4.6.10-zh}
两个 $f$-丰沛 $\mathcal{O}_X$-模层的张量积是 $f$-丰沛的。
\end{corollary}

\begin{proposition}[4.6.11]
\label{II.4.6.11-zh}
设 $Y$ 是一个拟紧预概形，$f:X\to Y$ 是一个有限型态射，
$\mathcal{L}$ 是一个可逆 $\mathcal{O}_X$-模层。$\mathcal{L}$ 是
$f$-丰沛的充要条件是它具有下列两个等价性质之一：
\begin{enumerate}
  \item[\rm{d)}] 存在 $n_0>0$，使对任意整数 $n\geq n_0$，
  $\mathcal{L}^{\otimes n}$ 相对于 $f$ 极丰沛。
  \item[\rm{d')}] 存在 $n>0$，使 $\mathcal{L}^{\otimes n}$ 相对于
  $f$ 极丰沛。
\end{enumerate}
\end{proposition}

若 $\mathcal{L}$ 相对于 $f$ 丰沛，则存在 $Y$ 的一个由仿射开集组成的
有限覆盖 $(U_i)$，使各 $\mathcal{L}|f^{-1}(U_i)$ 都丰沛。由
\egaref{II.4.5.10-zh}{4.5.10} 可得，存在一个整数 $n_0$，使对任意
$n\geq n_0$ 和每个 $i$，$\mathcal{L}^{\otimes n}|f^{-1}(U_i)$
都相对于 $f^{-1}(U_i)\to U_i$ 极丰沛；因此
$\mathcal{L}^{\otimes n}$ 相对于 $f$ 极丰沛
（\egaref{II.4.4.5-zh}{4.4.5}）。反过来，d') 已蕴含
$\mathcal{L}^{\otimes n}$ 是 $f$-丰沛的
（\egaref{II.4.6.2-zh}{4.6.2}），因而 $\mathcal{L}$ 也是如此
（\egaref{II.4.6.9-zh}{4.6.9, (i)}）。

\begin{corollary}[4.6.12]
\label{II.4.6.12-zh}
设 $Y$ 是一个拟紧预概形，$f:X\to Y$ 是一个有限型态射，
$\mathcal{L}$、$\mathcal{L}'$ 是两个可逆 $\mathcal{O}_X$-模层。
若 $\mathcal{L}$ 是 $f$-丰沛的，则存在 $n_0$，使对任意
$n\geq n_0$，$\mathcal{L}^{\otimes n}\otimes\mathcal{L}'$
相对于 $f$ 极丰沛。
\end{corollary}

使用 $Y$ 的一个有限仿射开覆盖和
\egaref{II.4.5.11-zh}{4.5.11}，按照
\egaref{II.4.6.11-zh}{4.6.11} 中的方式论证即可。
% Locked French authority: ega2-1-fr.tex, 820505 B,
% SHA-256 84EDBE3E83530AF2959B441796337C9DC21EAFCA6A13114A26778760FBF437AC.
% Sequential source scope: lines 8287-8337, proposition II.4.6.13
% and the opening reduction of its proof / printed pp.91-92.

\begin{proposition}[4.6.13]
\label{II.4.6.13-zh}
\begin{enumerate}
  \item[\rm{(i)}] 对任意预概形 $Y$，每个可逆
  $\mathcal{O}_Y$-模层 $\mathcal{L}$ 都相对于恒等态射 $1_Y$ 丰沛。
  \item[\rm{(i bis)}] 设 $f:X\to Y$ 是一个拟紧态射，
  $j:X'\to X$ 是一个拟紧态射，并且它把 $X'$ 的底空间同胚到
  $X$ 的一个子空间上。若 $\mathcal{L}$ 是一个相对于 $f$ 丰沛的
  $\mathcal{O}_X$-模层，则 $j^*(\mathcal{L})$ 相对于 $f\circ j$
  丰沛。
  \item[\rm{(ii)}] 设 $Z$ 是一个拟紧预概形，$f:X\to Y$、
  $g:Y\to Z$ 是两个拟紧态射，$\mathcal{L}$ 是一个相对于 $f$
  丰沛的 $\mathcal{O}_X$-模层，$\mathcal{K}$ 是一个相对于 $g$
  丰沛的 $\mathcal{O}_Y$-模层。那么存在一个整数 $n_0>0$，使对每个
  $n\geq n_0$，$\mathcal{L}\otimes f^*(\mathcal{K}^{\otimes n})$
  都相对于 $g\circ f$ 丰沛。
  \item[\rm{(iii)}] 设 $f:X\to Y$ 是一个拟紧态射，
  $g:Y'\to Y$ 是一个态射，并置 $X'=X_{(Y')}$。若 $\mathcal{L}$
  是一个相对于 $f$ 丰沛的 $\mathcal{O}_X$-模层，则
  $\mathcal{L}'=\mathcal{L}\otimes_Y\mathcal{O}_{Y'}$ 是一个
  相对于 $f_{(Y')}$ 丰沛的 $\mathcal{O}_{X'}$-模层。
  \item[\rm{(iv)}] 设 $f_i:X_i\to Y_i$（$i=1,2$）是两个拟紧
  $S$-态射。若 $\mathcal{L}_i$ 是一个相对于 $f_i$ 丰沛的
  $\mathcal{O}_{X_i}$-模层（$i=1,2$），则
  $\mathcal{L}_1\otimes_S\mathcal{L}_2$ 相对于
  $f_1\times_S f_2$ 丰沛。
  \item[\rm{(v)}] 设 $f:X\to Y$、$g:Y\to Z$ 是两个态射，并且
  $g\circ f$ 拟紧。若一个 $\mathcal{O}_X$-模层 $\mathcal{L}$
  相对于 $g\circ f$ 丰沛，并且 $g$ 是分离的，或 $X$ 的底空间
  局部诺特，则 $\mathcal{L}$ 相对于 $f$ 丰沛。

  \oldpage[II]{92}

  \item[\rm{(vi)}] 设 $f:X\to Y$ 是一个拟紧态射，$j$ 是典范单射
  $X_{\mathrm{red}}\to X$。若 $\mathcal{L}$ 是一个相对于 $f$
  丰沛的 $\mathcal{O}_X$-模层，则 $j^*(\mathcal{L})$ 相对于
  $f_{\mathrm{red}}$ 丰沛。
\end{enumerate}
\end{proposition}

先注意，使用与 \egaref{II.4.4.10-zh}{4.4.10} 中相同的论证，并以
\egaref{II.4.6.4-zh}{4.6.4} 代替
\egaref{II.4.4.5-zh}{4.4.5}，可以从 (i)、(i bis) 和 (iv) 推出
(v) 和 (vi)；我们把论证细节留给读者。断言 (i) 显然是
\egaref{II.4.4.10-zh}{4.4.10, (i)} 和
\egaref{II.4.6.2-zh}{4.6.2} 的推论。为证明 (i bis)、(iii) 和
(iv)，将使用下述引理：
% Locked French authority: ega2-1-fr.tex, 820505 B,
% SHA-256 84EDBE3E83530AF2959B441796337C9DC21EAFCA6A13114A26778760FBF437AC.
% Sequential source scope: lines 8339-8450, lemma II.4.6.13.1
% and the remainder of the proof of II.4.6.13 / printed pp.92-93.

\begin{lemma}[4.6.13.1]
\label{II.4.6.13.1-zh}
\begin{enumerate}
  \item[\rm{(i)}] 设 $u:Z\to S$ 是一个态射，$\mathcal{L}$ 是一个
  可逆 $\mathcal{O}_S$-模层，$s$ 是 $\mathcal{L}$ 在 $S$ 上的一个
  截面，$s'$ 是与它典范对应的 $u^*(\mathcal{L})=\mathcal{L}'$
  在 $Z$ 上的截面。那么 $Z_{s'}=u^{-1}(S_s)$。
  \item[\rm{(ii)}] 设 $Z$、$Z'$ 是两个 $S$-预概形，$p$、$p'$ 是
  $T=Z\times_S Z'$ 的两个投影，$\mathcal{L}$（相应地，
  $\mathcal{L}'$）是一个可逆 $\mathcal{O}_Z$-模层（相应地，一个
  可逆 $\mathcal{O}_{Z'}$-模层），$t$（相应地，$t'$）是
  $\mathcal{L}$ 在 $Z$ 上的一个截面（相应地，$\mathcal{L}'$ 在
  $Z'$ 上的一个截面），$s$（相应地，$s'$）是与它对应的
  $p^*(\mathcal{L})$（相应地，$p'^*(\mathcal{L}')$）在
  $Z\times_S Z'$ 上的截面。那么
  $T_{s\otimes s'}=Z_t\times_S Z'_{t'}$。
\end{enumerate}
\end{lemma}

由定义可知，可以归结到所考虑的全部预概形都是仿射的情形。此外，在
(i) 中，可以假设 $\mathcal{L}=\mathcal{O}_S$；此时断言 (i) 立即由
\egaxref{I.1.2.2.2-zh}{I, 1.2.2.2} 得出。同样，在 (ii) 中，可以
限于 $\mathcal{L}=\mathcal{O}_Z$、
$\mathcal{L}'=\mathcal{O}_{Z'}$ 的情形，此时断言归结为引理
\egaref{II.4.3.2.4-zh}{4.3.2.4}。

现在证明 (i bis)。可以假设 $Y$ 是仿射的
（\egaref{II.4.6.4-zh}{4.6.4}），因而 $\mathcal{L}$ 丰沛
（\egaref{II.4.6.6-zh}{4.6.6}）；当 $s$ 遍历各
$\Gamma(X,\mathcal{L}^{\otimes n})$（$n>0$）的并集时，诸 $X_s$
构成 $X$ 的拓扑的一个基（\egaref{II.4.5.2-zh}{4.5.2, a)}），
所以根据假设，诸 $j^{-1}(X_s)$ 构成 $X'$ 的拓扑的一个基；于是由引理
\egaref{II.4.6.13.1-zh}{4.6.13.1, (i)} 和
\egaref{II.4.5.2-zh}{4.5.2, a)} 可知，$j^*(\mathcal{L})$ 丰沛。

接着证明 (iii)。再次可以假设 $Y$ 和 $Y'$ 都是仿射的
（\egaref{II.4.6.4-zh}{4.6.4}），由此投影 $h:X'\to X$ 是仿射的
（\egaref{II.1.5.5-zh}{1.5.5}）。由于 $\mathcal{L}$ 丰沛
（\egaref{II.4.6.6-zh}{4.6.6}），当 $s$ 遍历
$\mathcal{L}^{\otimes n}$（$n>0$）在 $X$ 上满足 $X_s$ 仿射的诸截面时，
诸 $X_s$ 覆盖 $X$（\egaref{II.4.5.2-zh}{4.5.2, a')}）；因此诸
$h^{-1}(X_s)$ 是仿射的（\egaref{II.1.2.5-zh}{1.2.5}），并覆盖
$X'$；于是再由引理 \egaref{II.4.6.13.1-zh}{4.6.13.1, (i)} 和
\egaref{II.4.5.2-zh}{4.5.2, a')} 可知，$\mathcal{L}'$ 丰沛，而态射
$f_{(Y')}$ 是拟紧的（\egaxref{I.6.6.4-zh}{I, 6.6.4, (iii)}）。

为证明 (iv)，先注意 $f_1\times_S f_2$ 是拟紧的
（\egaxref{I.6.6.4-zh}{I, 6.6.4, (iv)}）。此外，可以假设
$S$、$Y_1$ 和 $Y_2$ 都是仿射的
（\egaref{II.4.6.4-zh}{4.6.4} 和
\egaxref{I.3.2.7-zh}{I, 3.2.7}），因而 $\mathcal{L}_i$ 丰沛
（$i=1,2$）（\egaref{II.4.6.6-zh}{4.6.6}）。当 $s_i$ 遍历
$\mathcal{L}_i^{\otimes n_i}$ 的满足 $(X_i)_{s_i}$ 仿射的诸截面时，
诸开集 $(X_1)_{s_1}\times_S(X_2)_{s_2}$ 覆盖
$X_1\times_S X_2$（\egaref{II.4.5.2-zh}{4.5.2, a')}）。而且，以
适当的幂替换 $s_1$ 和 $s_2$ 并不改变 $(X_i)_{s_i}$，所以总可以假设
$n_1=n_2$。于是由 \egaref{II.4.6.13.1-zh}{4.6.13.1, (ii)} 和
\egaref{II.4.5.2-zh}{4.5.2, a')} 可知，
$\mathcal{L}_1\otimes_S\mathcal{L}_2$ 丰沛；由于
$Y_1\times_S Y_2$ 是仿射的
（\egaref{II.4.6.6-zh}{4.6.6}），断言随之得出。

还须证明 (ii)。采用与 \egaref{II.4.4.10-zh}{4.4.10} 中相同的论证，
但在这里使用 \egaref{II.4.6.4-zh}{4.6.4}，可以归结到 $Z$ 为仿射的
情形。此时 $\mathcal{K}$ 丰沛，而 $Y$ 拟紧，所以存在有限多个截面
$s_i\in\Gamma(Y,\mathcal{K}^{\otimes k_i})$，使诸 $Y_{s_i}$

\oldpage[II]{93}

都是仿射的并覆盖 $Y$
（\egaref{II.4.5.2-zh}{4.5.2, a')}）；用适当的幂替换各 $s_i$，
还可以假设所有 $k_i$ 都等于同一个整数 $k$。设 $s'_i$ 是与 $s_i$
典范对应的 $f^*(\mathcal{K}^{\otimes k})$ 在 $X$ 上的截面，从而诸
% SOURCE-DEFECT CANDIDATE, kept source-literal: French line 8420 prints
% “4.6.1.13, (i)”; the immediately preceding lemma is 4.6.13.1.
$X_{s'_i}=f^{-1}(Y_{s_i})$（4.6.1.13，(i)）覆盖 $X$。由于
$\mathcal{L}|X_{s'_i}$ 丰沛
（\egaref{II.4.6.4-zh}{4.6.4} 和
\egaref{II.4.6.6-zh}{4.6.6}），对每个 $i$，存在有限多个截面
$t_{ij}\in\Gamma(X,\mathcal{L}^{\otimes n_{ij}})$，使诸
$X_{t_{ij}}$ 都是仿射的、包含在 $X_{s'_i}$ 中并覆盖
$X_{s'_i}$（\egaref{II.4.5.2-zh}{4.5.2, a')}）；还可以假设所有
$n_{ij}$ 都等于同一个数 $n$。现在，$X$ 是分离且拟紧的，所以存在一个
整数 $m>0$，并且对每个 $(i,j)$ 存在一个截面
\[
  u_{ij}\in\Gamma(X,\mathcal{L}^{\otimes n}\otimes
  f^*(\mathcal{K}^{\otimes mk}))
\]
使 $t_{ij}\otimes s_i'^{\otimes m}$ 是 $u_{ij}$ 在 $X_{s'_i}$
上的限制（\egaxref{I.9.3.1-zh}{I, 9.3.1}）；此外，
$X_{u_{ij}}=X_{t_{ij}}$，所以诸 $X_{u_{ij}}$ 都是仿射的并覆盖
$X$。还可以假设 $m$ 具有形式 $nr$；若置 $n_0=rk$，便由
\egaref{II.4.5.2-zh}{4.5.2, a')} 看出
$\mathcal{L}\otimes_{\mathcal{O}_X}f^*(\mathcal{K}^{\otimes n_0})$
丰沛。此外，存在 $h_0>0$，使只要 $h\geq h_0$，
$\mathcal{K}^{\otimes h}$ 就由它在 $Y$ 上的截面生成
（\egaref{II.4.5.5-zh}{4.5.5}）；因而更有，只要 $h\geq h_0$，
$f^*(\mathcal{K}^{\otimes h})$ 就由它在 $X$ 上的截面生成，这是逆像的
定义所致（\egaxref{0.3.7.1-zh}{0, 3.7.1} 和
\egaref{II.4.4.1-zh}{4.4.1}）。由此根据
\egaref{II.4.5.6-zh}{4.5.6} 可知，只要 $h\geq h_0$，
$\mathcal{L}\otimes f^*(\mathcal{K}^{\otimes n_0+h})$ 就丰沛，
这就完成了证明。
% Locked French authority: ega2-1-fr.tex, 820505 B,
% SHA-256 84EDBE3E83530AF2959B441796337C9DC21EAFCA6A13114A26778760FBF437AC.
% Sequential source scope: lines 8452-8499, remark II.4.6.14,
% proposition II.4.6.15 and corollary II.4.6.16 / printed p.93.

\begin{remark}[4.6.14]
\label{II.4.6.14-zh}
在 (ii) 的条件下，不要以为 $\mathcal{L}\otimes f^*(\mathcal{K})$
相对于 $g\circ f$ 丰沛；事实上，由于
$\mathcal{L}\otimes f^*(\mathcal{K}^{-1})$ 相对于 $f$ 也是丰沛的
（\egaref{II.4.6.5-zh}{4.6.5}），便会由此推出 $\mathcal{L}$ 相对于
$g\circ f$ 丰沛；特别取 $g$ 为恒等态射，便会得到每个可逆
$\mathcal{O}_X$-模层相对于 $f$ 都是丰沛的，而一般并非如此
（见 \egaref{II.5.1.6-zh}{5.1.6}、
\egaref{II.5.3.4-zh}{5.3.4, (i)} 和
\egaref{II.5.3.1-zh}{5.3.1}）。
\end{remark}

\begin{proposition}[4.6.15]
\label{II.4.6.15-zh}
设 $f:X\to Y$ 是一个拟紧态射，$\mathcal{J}$ 是
$\mathcal{O}_X$ 的一个局部幂零拟凝聚理想层，$Z$ 是由
$\mathcal{J}$ 定义的 $X$ 的闭子预概形，$j:Z\to X$ 是典范单射。
一个可逆 $\mathcal{O}_X$-模层 $\mathcal{L}$ 相对于 $f$ 丰沛的充要条件是
$j^*(\mathcal{L})$ 相对于 $f\circ j$ 丰沛。
\end{proposition}

事实上，由于问题在 $Y$ 上是局部的
（\egaref{II.4.6.4-zh}{4.6.4}），可以假设 $Y$ 是仿射的；此时 $X$ 拟紧，
因而可以假设 $\mathcal{J}$ 幂零。考虑到
\egaref{II.4.6.6-zh}{4.6.6}，本命题便正是
\egaref{II.4.5.13-zh}{4.5.13}。

\begin{corollary}[4.6.16]
\label{II.4.6.16-zh}
设 $X$ 是一个局部诺特预概形，$f:X\to Y$ 是一个拟紧态射。
一个可逆 $\mathcal{O}_X$-模层 $\mathcal{L}$ 相对于 $f$ 丰沛的充要条件是，
它沿典范单射 $X_{\mathrm{red}}\to X$ 的逆像 $\mathcal{L}'$
相对于 $f_{\mathrm{red}}$ 丰沛。
\end{corollary}

已经看到该条件是必要的
（\egaref{II.4.6.13-zh}{4.6.13, (vi)}）；反过来，若该条件成立，
为了证明 $\mathcal{L}$ 相对于 $f$ 丰沛，可以归结到 $Y$ 为仿射的情形
（\egaref{II.4.6.4-zh}{4.6.4}）；于是 $Y_{\mathrm{red}}$ 也仿射，
故 $\mathcal{L}'$ 丰沛（\egaref{II.4.6.6-zh}{4.6.6}），而由
\egaref{II.4.5.13-zh}{4.5.13}，$\mathcal{L}$ 也丰沛，因为此时
$X$ 是诺特的，并且 $X_{\mathrm{red}}$ 是 $X$ 的一个闭子预概形，
由一个幂零拟凝聚理想层定义
（\egaxref{I.6.1.6-zh}{I, 6.1.6}）。
% Locked French authority: ega2-1-fr.tex, 820505 B,
% SHA-256 84EDBE3E83530AF2959B441796337C9DC21EAFCA6A13114A26778760FBF437AC.
% Sequential source scope: lines 8501-8518, proposition II.4.6.17
% and its proof / printed pp.93-94.

\begin{proposition}[4.6.17]
\label{II.4.6.17-zh}
采用 \egaref{II.4.4.11-zh}{4.4.11} 的记号和假设。
$\mathcal{L}''$ 相对于 $f''$ 丰沛的充要条件是 $\mathcal{L}$ 相对于
$f$ 丰沛，并且 $\mathcal{L}'$ 相对于 $f'$ 丰沛。
\end{proposition}

\oldpage[II]{94}

条件的必要性由 \egaref{II.4.6.13-zh}{4.6.13, (i bis)} 得出。
为证明条件充分，可以归结到 $Y$ 为仿射的情形；这时，把
\egaref{II.4.5.2-zh}{4.5.2, a)} 的判别法应用于
$\mathcal{L}$、$\mathcal{L}'$ 和 $\mathcal{L}''$，并注意到
$\mathcal{L}$ 在 $X$ 上的一个截面可以（以 $0$）延拓成
$\mathcal{L}''$ 在 $X''$ 上的一个截面，便得 $\mathcal{L}''$ 丰沛。
% Locked French authority: ega2-1-fr.tex, 820505 B,
% SHA-256 84EDBE3E83530AF2959B441796337C9DC21EAFCA6A13114A26778760FBF437AC.
% Sequential source scope: lines 8520-8538, proposition II.4.6.18
% and its proof / printed p.94.

\begin{proposition}[4.6.18]
\label{II.4.6.18-zh}
设 $Y$ 是一个拟紧预概形，$\mathcal{S}$ 是一个有限型拟凝聚分次
$\mathcal{O}_Y$-代数层，$X=\operatorname{Proj}(\mathcal{S})$，
$f:X\to Y$ 是结构态射。那么 $f$ 是有限型的，并且存在一个整数
$d>0$，使 $\mathcal{O}_X(d)$ 可逆且为 $f$-丰沛的。
\end{proposition}

由 \egaref{II.3.1.10-zh}{3.1.10}，存在一个整数 $d>0$，使
$\mathcal{S}^{(d)}$ 由 $\mathcal{S}_d$ 生成。已知在 $X$ 与
$X^{(d)}=\operatorname{Proj}(\mathcal{S}^{(d)})$ 之间的典范同构中，
$\mathcal{O}_X(d)$ 与 $\mathcal{O}_{X^{(d)}}(1)$ 等同
（\egaref{II.3.2.9-zh}{3.2.9, (ii)}）。因此，问题归结为
$\mathcal{S}$ 由 $\mathcal{S}_1$ 生成的情形；这时，本命题由
\egaref{II.4.4.3-zh}{4.4.3} 和
\egaref{II.4.6.2-zh}{4.6.2} 得出，同时要注意 $f$ 是有限型态射
（\egaref{II.3.4.1-zh}{3.4.1}）。
% Locked French authority: ega2-1-fr.tex, 820505 B,
% SHA-256 84EDBE3E83530AF2959B441796337C9DC21EAFCA6A13114A26778760FBF437AC.
% Sequential source scope: lines 8540-8570, Section II.5 opening,
% subsection II.5.1 opening and Definition II.5.1.1 / printed p.94.

\section{拟仿射态射；拟射影态射。紧合态射；射影态射}
\label{section:II.5-zh}

\subsection{拟仿射态射。}
\label{subsection:II.5.1-zh}

\begin{definition}[5.1.1]
\label{II.5.1.1-zh}
称同构于某个仿射概形在其一个拟紧开集上所诱导的子概形的概形为
\emph{拟仿射概形}。称态射 $f:X\to Y$ 为\emph{拟仿射的}，或者称
$X$（借助 $f$ 看作 $Y$-预概形）为一个\emph{$Y$-拟仿射概形}，如果
$Y$ 存在一个由仿射开集组成的覆盖 $(U_\alpha)$，使各
$f^{-1}(U_\alpha)$ 都是拟仿射概形。
\end{definition}

显然，拟仿射态射是\emph{分离的}
（\egaxref{I.5.5.5-zh}{I, 5.5.5} 和
\egaxref{I.5.5.8-zh}{5.5.8}）且是\emph{拟紧的}
（\egaxref{I.6.6.1-zh}{I, 6.6.1}）；每个仿射态射显然都是拟仿射的。

回忆一下，对于任意预概形 $X$，若置
$A=\Gamma(X,\mathcal{O}_X)$，恒等同态
$A\to A=\Gamma(X,\mathcal{O}_X)$ 定义一个称为\emph{典范态射}的态射
$X\to\operatorname{Spec}(A)$（\egaxref{I.2.2.4-zh}{I, 2.2.4}）；
此外，它正是 \egaref{II.4.5.1-zh}{4.5.1} 在
$\mathcal{L}=\mathcal{O}_X$ 这一特殊情形中所定义的典范态射，因为
$\operatorname{Proj}(A[T])$ 与 $\operatorname{Spec}(A)$ 典范等同
（\egaref{II.3.1.7-zh}{3.1.7}）。
% Locked French authority: ega2-1-fr.tex, 820505 B,
% SHA-256 84EDBE3E83530AF2959B441796337C9DC21EAFCA6A13114A26778760FBF437AC.
% Sequential source scope: lines 8572-8614, Proposition II.5.1.2,
% its proof and the following observation / printed pp.94-95.

\begin{proposition}[5.1.2]
\label{II.5.1.2-zh}
设 $X$ 是一个拟紧概形，或一个底空间为诺特空间的预概形，$A$ 是环
$\Gamma(X,\mathcal{O}_X)$。下列条件等价：
\begin{enumerate}
  \item[\rm{a)}] $X$ 是拟仿射概形。
  \item[\rm{b)}] 典范态射 $u:X\to\operatorname{Spec}(A)$ 是开浸入。
  \item[\rm{b')}] 典范态射 $u:X\to\operatorname{Spec}(A)$ 把 $X$
  同胚到 $\operatorname{Spec}(A)$ 的底空间的一个子空间上。
  \item[\rm{c)}] $\mathcal{O}_X$-模层 $\mathcal{O}_X$ 相对于 $u$
  极丰沛（\egaref{II.4.4.2-zh}{4.4.2}）。
  \item[\rm{c')}] $\mathcal{O}_X$-模层 $\mathcal{O}_X$ 丰沛
  （\egaref{II.4.5.1-zh}{4.5.1}）。
  \item[\rm{d)}] 当 $f$ 遍历 $A$ 时，诸 $X_f$ 构成 $X$ 的拓扑的一个基。
  \item[\rm{d')}] 当 $f$ 遍历 $A$ 时，诸仿射的 $X_f$ 构成 $X$ 的一个覆盖。

\oldpage[II]{95}

  \item[\rm{e)}] 每个拟凝聚 $\mathcal{O}_X$-模层都由其在 $X$ 上的截面生成。
  \item[\rm{e')}] $\mathcal{O}_X$ 的每个有限型拟凝聚理想层都由其在
  $X$ 上的截面生成。
\end{enumerate}
\end{proposition}

显然，b) 蕴含 a)，而由把 \egaref{II.4.4.4-zh}{4.4.4} 的判别法 b)
应用于恒等态射（同时考虑到命题陈述前的说明），a) 蕴含 c)；另一方面，
c) 蕴含 c')（\egaref{II.4.5.10-zh}{4.5.10, (i)}），并且由
\egaref{II.4.5.2-zh}{4.5.2, 判别法 b) 和 b')}，c')、b) 和 b')
等价。最后，由 \egaref{II.4.5.2-zh}{4.5.2, 判别法 a)、a')、c)}
和 \egaref{II.4.5.5-zh}{4.5.5, 判别法 d'')}，c') 分别等价于
判别法 d)、d')、e)、e')。

还应注意，在上述记号下，诸仿射的 $X_f$ 构成 $X$ 的拓扑的一个
\emph{基}，并且典范态射 $u$ 是\emph{支配的}
（\egaref{II.4.5.2-zh}{4.5.2}）。
% Locked French authority: ega2-1-fr.tex, 820505 B,
% SHA-256 84EDBE3E83530AF2959B441796337C9DC21EAFCA6A13114A26778760FBF437AC.
% Sequential source scope: lines 8616-8651, Corollaries II.5.1.3-II.5.1.5
% and their proofs / printed p.95.

\begin{corollary}[5.1.3]
\label{II.5.1.3-zh}
设 $X$ 是一个拟紧预概形。若存在从 $X$ 到一个仿射概形 $Y$ 的态射
$v:X\to Y$，它把 $X$ 同胚到 $Y$ 的一个开子空间上，则 $X$ 是
拟仿射的。
\end{corollary}

事实上，存在 $\mathcal{O}_Y$ 在 $Y$ 上的一族截面 $(g_\alpha)$，使诸
$D(g_\alpha)$ 构成 $v(X)$ 的拓扑的一个基；若
$v=(\psi,\theta)$，并置 $f_\alpha=\theta(g_\alpha)$，则
$X_{f_\alpha}=\psi^{-1}(D(g_\alpha))$
（\egaxref{I.2.2.4.1-zh}{I, 2.2.4.1}），所以诸 $X_{f_\alpha}$
构成 $X$ 的拓扑的一个基，从而满足
\egaref{II.5.1.2-zh}{5.1.2} 的判别法 d)。

\begin{corollary}[5.1.4]
\label{II.5.1.4-zh}
若 $X$ 是拟仿射概形，则每个可逆 $\mathcal{O}_X$-模层都极丰沛
（相对于典范态射），因而更是丰沛的。
\end{corollary}

事实上，这样的模层 $\mathcal{L}$ 由其在 $X$ 上的截面生成
（\egaref{II.5.1.2-zh}{5.1.2, e)}），所以
$\mathcal{L}\otimes\mathcal{O}_X=\mathcal{L}$ 极丰沛
（\egaref{II.4.4.8-zh}{4.4.8}）。

\begin{corollary}[5.1.5]
\label{II.5.1.5-zh}
设 $X$ 是一个拟紧预概形。若存在一个可逆 $\mathcal{O}_X$-模层
$\mathcal{L}$，使 $\mathcal{L}$ 和 $\mathcal{L}^{-1}$ 都丰沛，
则 $X$ 是拟仿射概形。
\end{corollary}

事实上，此时 $\mathcal{O}_X=\mathcal{L}\otimes\mathcal{L}^{-1}$
丰沛（\egaref{II.4.5.7-zh}{4.5.7}）。
% Locked French authority: ega2-1-fr.tex, 820505 B,
% SHA-256 84EDBE3E83530AF2959B441796337C9DC21EAFCA6A13114A26778760FBF437AC.
% Sequential source scope: lines 8653-8692, Proposition II.5.1.6
% and its proof / printed pp.95-96.

\begin{proposition}[5.1.6]
\label{II.5.1.6-zh}
设 $f:X\to Y$ 是一个拟紧态射。下列条件等价：
\begin{enumerate}
  \item[\rm{a)}] $f$ 是拟仿射的。
  \item[\rm{b)}] $\mathcal{O}_Y$-代数层
  $f_*(\mathcal{O}_X)=\mathcal{A}(X)$ 是拟凝聚的，并且与恒等同态
  $\mathcal{A}(X)\to\mathcal{A}(X)$ 对应的典范态射
  $X\to\operatorname{Spec}(\mathcal{A}(X))$
  （\egaref{II.1.2.7-zh}{1.2.7}）是开浸入。
  \item[\rm{b')}] $\mathcal{O}_Y$-代数层 $\mathcal{A}(X)$ 是
  拟凝聚的，并且典范态射
  $X\to\operatorname{Spec}(\mathcal{A}(X))$ 把 $X$ 同胚到
  $\operatorname{Spec}(\mathcal{A}(X))$ 的一个子空间上。
  \item[\rm{c)}] $\mathcal{O}_X$-模层 $\mathcal{O}_X$ 对 $f$ 极丰沛。
  \item[\rm{c')}] $\mathcal{O}_X$-模层 $\mathcal{O}_X$ 对 $f$ 丰沛。
  \item[\rm{d)}] 态射 $f$ 是分离的，并且对每个拟凝聚
  $\mathcal{O}_X$-模层 $\mathcal{F}$，典范同态
  $\sigma:f^*(f_*(\mathcal{F}))\to\mathcal{F}$
  （\egaxref{0.4.4.3-zh}{0, 4.4.3}）都是满射。
\end{enumerate}
此外，当 $f$ 是拟仿射的时，每个可逆 $\mathcal{O}_X$-模层
$\mathcal{L}$ 都相对于 $f$ 极丰沛。
\end{proposition}

a) 与 c') 的等价性，由 $f$-丰沛性在 $Y$ 上的局部性
（\egaref{II.4.6.4-zh}{4.6.4}）、定义
\egaref{II.5.1.1-zh}{5.1.1} 和判别法
\egaref{II.5.1.2-zh}{5.1.2, c')} 得出。其余性质在 $Y$ 上都是局部的

\oldpage[II]{96}

，因而立即由 \egaref{II.5.1.2-zh}{5.1.2} 和
\egaref{II.5.1.4-zh}{5.1.4} 得出；这里还要注意，当 $f$ 分离时，
$f_*(\mathcal{F})$ 是拟凝聚的
（\egaxref{I.9.2.2-zh}{I, 9.2.2, a)}）。
% Locked French authority: ega2-1-fr.tex, 820505 B,
% SHA-256 84EDBE3E83530AF2959B441796337C9DC21EAFCA6A13114A26778760FBF437AC.
% Sequential source scope: lines 8694-8730, Corollaries II.5.1.7-II.5.1.9
% and their proofs / printed p.96.

\begin{corollary}[5.1.7]
\label{II.5.1.7-zh}
设 $f:X\to Y$ 是一个拟仿射态射。对 $Y$ 的每个开集 $U$，$f$ 的限制
$f^{-1}(U)\to U$ 都是拟仿射的。
\end{corollary}

\begin{corollary}[5.1.8]
\label{II.5.1.8-zh}
设 $Y$ 是一个仿射概形，$f:X\to Y$ 是一个拟紧态射。$f$ 是拟仿射的
充要条件是 $X$ 是拟仿射概形。
\end{corollary}

这是 \egaref{II.5.1.6-zh}{5.1.6} 和
\egaref{II.4.6.6-zh}{4.6.6} 的直接推论。

\begin{corollary}[5.1.9]
\label{II.5.1.9-zh}
设 $Y$ 是一个拟紧概形，或一个底空间为诺特空间的预概形，
$f:X\to Y$ 是一个有限型态射。若 $f$ 是拟仿射的，则
$\mathcal{A}(X)=f_*(\mathcal{O}_X)$ 存在一个有限型拟凝聚
$\mathcal{O}_Y$-子代数层 $\mathcal{B}$
（\egaxref{I.9.6.2-zh}{I, 9.6.2}），使与典范单射
$\mathcal{B}\to\mathcal{A}(X)$ 对应的态射
$X\to\operatorname{Spec}(\mathcal{B})$
（\egaref{II.1.2.7-zh}{1.2.7}）是浸入。此外，$\mathcal{A}(X)$ 的
每个包含 $\mathcal{B}$ 的有限型拟凝聚 $\mathcal{O}_Y$-子代数层
$\mathcal{B}'$ 都具有同样的性质。
\end{corollary}

事实上，$\mathcal{A}(X)$ 是它的各有限型拟凝聚
$\mathcal{O}_Y$-子代数层的归纳极限
（\egaxref{I.9.6.5-zh}{I, 9.6.5}）；考虑到
$\operatorname{Spec}(\mathcal{A}(X))$ 与
$\operatorname{Proj}(\mathcal{A}(X)[T])$ 的等同
（\egaref{II.3.1.7-zh}{3.1.7}），结论便是
\egaref{II.3.8.4-zh}{3.8.4} 的一个特殊情形。
% Locked French authority: ega2-1-fr.tex, 820505 B,
% SHA-256 84EDBE3E83530AF2959B441796337C9DC21EAFCA6A13114A26778760FBF437AC.
% Sequential source scope: lines 8732-8760, Proposition II.5.1.10
% and its proof / printed p.96.

\begin{proposition}[5.1.10]
\label{II.5.1.10-zh}
\begin{enumerate}
  \item[\rm{(i)}] 若一个拟紧态射 $X\to Y$ 把 $X$ 的底空间同胚到
  $Y$ 的底空间的一个子空间上（特别地，若它是闭浸入），则它是拟仿射的。
  \item[\rm{(ii)}] 两个拟仿射态射的复合是拟仿射的。
  \item[\rm{(iii)}] 若 $f:X\to Y$ 是一个拟仿射 $S$-态射，则对
  底预概形的每个扩张 $S'\to S$，
  $f_{(S')}:X_{(S')}\to Y_{(S')}$ 都是拟仿射态射。
  \item[\rm{(iv)}] 若 $f:X\to Y$ 和 $g:X'\to Y'$ 是两个拟仿射
  $S$-态射，则 $f\times_S g$ 是拟仿射的。
  \item[\rm{(v)}] 若 $f:X\to Y$、$g:Y\to Z$ 是两个态射，使
  $g\circ f$ 拟仿射，并且 $g$ 是分离的，或 $X$ 的底空间局部诺特，
  则 $f$ 是拟仿射的。
  \item[\rm{(vi)}] 若 $f$ 是拟仿射态射，则
  $f_{\mathrm{red}}$ 也是拟仿射态射。
\end{enumerate}
\end{proposition}

考虑判别法 \egaref{II.5.1.6-zh}{5.1.6, c')}，(i)、(iii)、(iv)、
(v) 和 (vi) 分别立即由
\egaref{II.4.6.13-zh}{4.6.13, (i bis)、(iii)、(iv)、(v) 和 (vi)}
得出。为证明 (ii)，可以归结到 $Z$ 为仿射的情形；此时，把
\egaref{II.4.6.13-zh}{4.6.13, (ii)} 应用于
$\mathcal{L}=\mathcal{O}_X$ 和 $\mathcal{K}=\mathcal{O}_Y$，
便直接得到该断言。
% Locked French authority: ega2-1-fr.tex, 820505 B,
% SHA-256 84EDBE3E83530AF2959B441796337C9DC21EAFCA6A13114A26778760FBF437AC.
% Sequential source scope: lines 8762-8787, Remark II.5.1.11,
% Proposition II.5.1.12 and its proof / printed pp.96-97.

\begin{remark}[5.1.11]
\label{II.5.1.11-zh}
设 $f:X\to Y$、$g:Y\to Z$ 是两个态射，并且 $X\times_ZY$ 局部诺特。
那么图浸入 $\Gamma_f:X\to X\times_ZY$ 是拟仿射的，因为它是拟紧的
（\egaxref{I.6.3.5-zh}{I, 6.3.5}）；而
\egaxref{I.5.5.12-zh}{I, 5.5.12} 表明，在 (v) 中，即使删去
$g$ 分离这一假设，结论仍然成立。
\end{remark}

\begin{proposition}[5.1.12]
\label{II.5.1.12-zh}
设 $f:X\to Y$ 是一个拟紧态射，$g:X'\to X$ 是一个拟仿射态射。
若 $\mathcal{L}$ 是一个对 $f$ 丰沛的 $\mathcal{O}_X$-模层，则
$g^*(\mathcal{L})$ 是一个对 $f\circ g$ 丰沛的
$\mathcal{O}_{X'}$-模层。
\end{proposition}

事实上，由于 $\mathcal{O}_{X'}$ 对 $g$ 极丰沛，并且问题在 $Y$ 上
是局部的（\egaref{II.4.6.4-zh}{4.6.4}），由
\egaref{II.4.6.13-zh}{4.6.13, (ii)} 可知，存在一个整数 $n$
（这里假设 $Y$ 仿射），使 $g^*(\mathcal{L})^{\otimes n}$ 对
$f\circ g$ 丰沛；因此 $g^*(\mathcal{L})$ 对 $f\circ g$ 丰沛
（\egaref{II.4.6.9-zh}{4.6.9}）。

\oldpage[II]{97}
% Locked French authority: ega2-1-fr.tex, 820505 B,
% SHA-256 84EDBE3E83530AF2959B441796337C9DC21EAFCA6A13114A26778760FBF437AC.
% Sequential source scope: lines 8789-8824, subsection II.5.2 opening
% and Theorem II.5.2.1 / printed p.97.

\subsection{塞尔判别法。}
\label{subsection:II.5.2-zh}

\begin{theorem}[5.2.1]
\label{II.5.2.1-zh}
\emph{（塞尔判别法。）}设 $X$ 是一个拟紧概形，或一个底空间为
诺特空间的预概形。下列条件等价：
\begin{enumerate}
  \item[\rm{a)}] $X$ 是仿射概形。
  \item[\rm{b)}] 存在一族元素
  $f_\alpha\in A=\Gamma(X,\mathcal{O}_X)$，使各 $X_{f_\alpha}$
  都是仿射的，并且诸 $f_\alpha$ 在 $A$ 中生成的理想等于 $A$。
  \item[\rm{c)}] 在拟凝聚 $\mathcal{O}_X$-模层范畴中，函子
  $\Gamma(X,\mathcal{F})$ 关于 $\mathcal{F}$ 是正合的；换言之，若
  \[
    0\to\mathcal{F}'\to\mathcal{F}\to\mathcal{F}''\to0 \tag{*}
  \]
  是拟凝聚 $\mathcal{O}_X$-模层的正合列，则序列
  \[
    0\to\Gamma(X,\mathcal{F}')\to\Gamma(X,\mathcal{F})
      \to\Gamma(X,\mathcal{F}'')\to0
  \]
  是正合的。
  \item[\rm{c')}] 对每个拟凝聚 $\mathcal{O}_X$-模层的正合列 (*)，
  只要其中的 $\mathcal{F}$ 同构于有限乘积 $\mathcal{O}_X^n$ 的一个
  $\mathcal{O}_X$-子模层，性质 c) 就成立。
  \item[\rm{d)}] 对每个拟凝聚 $\mathcal{O}_X$-模层
  $\mathcal{F}$，都有 $\mathrm{H}^1(X,\mathcal{F})=0$。
  \item[\rm{d')}] 对 $\mathcal{O}_X$ 的每个拟凝聚理想层
  $\mathcal{J}$，都有 $\mathrm{H}^1(X,\mathcal{J})=0$。
\end{enumerate}
\end{theorem}
% Locked French authority: ega2-1-fr.tex, 820505 B,
% SHA-256 84EDBE3E83530AF2959B441796337C9DC21EAFCA6A13114A26778760FBF437AC.
% Sequential source scope: lines 8826-8890, proof of Theorem II.5.2.1
% / printed pp.97-98.

显然 a) 蕴含 b)；另一方面，b) 蕴含诸 $X_{f_\alpha}$ 覆盖 $X$，
因为依假设截面 $1$ 是诸 $f_\alpha$ 的线性组合，从而诸
$D(f_\alpha)$ 覆盖 $\operatorname{Spec}(A)$。于是
\egaref{II.4.5.2-zh}{4.5.2} 的最后一个断言表明
$X\to\operatorname{Spec}(A)$ 是同构。

已知 a) 蕴含 c)（\egaxref{I.1.3.11-zh}{I, 1.3.11}），而 c) 显然
蕴含 c')。下面证明 c') 蕴含 b)。首先，c') 蕴含如下事实：对每个闭点
$x\in X$ 以及 $x$ 的每个开邻域 $U$，存在 $f\in A$，使
% EGA2-CANON-DISP-0010 / EGA2DEF94: guardian-confirmed reversible reader
% correction; French diplomatic source remains unchanged.
$x\in X_f\subset U$\footnote{法文印本此处写作
$x\in X_f\subset X-U$。但所取开邻域含有该点，而下文构造的截面在其
补集上为零、在该点的芽等于单位，故相应主开集必须包含于该邻域。本译依
源文勘误记录 EGA2-ZH-DEF-0013 及守护者处置
EGA2-CANON-DISP-0010 作可逆校正；法文外交式来源保持不变。}。事实上，令 $\mathcal{J}$（相应地，
$\mathcal{J}'$）为 $\mathcal{O}_X$ 的拟凝聚理想层，它定义 $X$ 的
约化闭子预概形，而该子预概形的底空间是 $X-U$（相应地，
$(X-U)\cup\{x\}$）（\egaxref{I.5.2.1-zh}{I, 5.2.1}）；显然有
$\mathcal{J}'\subset\mathcal{J}$，并且
$\mathcal{J}''=\mathcal{J}/\mathcal{J}'$ 是一个支集为 $\{x\}$ 的
拟凝聚 $\mathcal{O}_X$-模层，且 $\mathcal{J}''_x=k(x)$。将假设 c')
应用于正合列 $0\to\mathcal{J}'\to\mathcal{J}\to\mathcal{J}''\to0$，
可知 $\Gamma(X,\mathcal{J})\to\Gamma(X,\mathcal{J}'')$ 是满射。
$\mathcal{J}''$ 中在 $x$ 处的芽为 $1_x$ 的截面，因而是某个截面
$f\in\Gamma(X,\mathcal{J})\subset\Gamma(X,\mathcal{O}_X)$ 的像；
依定义有 $f(x)=1_x$，并且对 $X-U$ 中的点有 $f(y)=0$，这就证明了
上述断言。此外，若 $U$ 是仿射的，则 $X_f$ 也是仿射的
（\egaxref{I.1.3.6-zh}{I, 1.3.6}），所以诸仿射的 $X_f$
（$f\in A$）之并是一个包含 $X$ 的所有闭点的开集 $Z$；由于 $X$
是拟紧的柯尔莫哥洛夫空间，必有 $Z=X$
（\egaxref{0.2.1.3-zh}{0, 2.1.3}）。又因 $X$ 拟紧，存在有限多个元素
$f_i\in A$（$1\leq i\leq n$），使诸 $X_{f_i}$ 都是仿射的并覆盖
$X$。现在考虑由截面 $f_i$ 定义的同态
$\mathcal{O}_X^n\to\mathcal{O}_X$
（\egaxref{0.5.1.1-zh}{0, 5.1.1}）；对每个 $x\in X$，诸
$(f_i)_x$ 中至少有一个可逆，所以该同态是满射，于是得到正合列
$0\to\mathcal{R}\to\mathcal{O}_X^n\to\mathcal{O}_X\to0$，其中
$\mathcal{R}$ 是 $\mathcal{O}_X^n$ 的拟凝聚 $\mathcal{O}_X$-子模层。
于是

\oldpage[II]{98}

由 c') 可知，相应的同态
$\Gamma(X,\mathcal{O}_X^n)\to\Gamma(X,\mathcal{O}_X)$ 是满射，
这就证明了 b)。

最后，a) 蕴含 d)（\egaxref{I.5.1.9.2-zh}{I, 5.1.9.2}），而 d)
显然蕴含 d')。只须证明 d') 蕴含 c')。设 $\mathcal{F}'$ 是
$\mathcal{O}_X^n$ 的拟凝聚 $\mathcal{O}_X$-子模层。滤过
$0\subset\mathcal{O}_X\subset\mathcal{O}_X^2\subset\cdots\subset
\mathcal{O}_X^n$ 在 $\mathcal{F}'$ 上诱导出由
$\mathcal{F}'_k=\mathcal{O}_X^k\cap\mathcal{F}'$
（$0\leq k\leq n$）组成的滤过；这些都是拟凝聚
$\mathcal{O}_X$-模层（\egaxref{I.4.1.1-zh}{I, 4.1.1}），而
$\mathcal{F}'_{k+1}/\mathcal{F}'_k$ 同构于
$\mathcal{O}_X^{k+1}/\mathcal{O}_X^k=\mathcal{O}_X$ 的一个拟凝聚
$\mathcal{O}_X$-子模层，也就是说，同构于 $\mathcal{O}_X$ 的一个
拟凝聚理想层。因此，假设 d') 蕴含
$\mathrm{H}^1(X,\mathcal{F}'_{k+1}/\mathcal{F}'_k)=0$；于是由
上同调正合列
$\mathrm{H}^1(X,\mathcal{F}'_k)\to\mathrm{H}^1(X,\mathcal{F}'_{k+1})
\to\mathrm{H}^1(X,\mathcal{F}'_{k+1}/\mathcal{F}'_k)=0$，可以对 $k$
作归纳，证明对每个 $k$ 都有
$\mathrm{H}^1(X,\mathcal{F}'_k)=0$。证毕。
% Locked French authority: ega2-1-fr.tex, 820505 B,
% SHA-256 84EDBE3E83530AF2959B441796337C9DC21EAFCA6A13114A26778760FBF437AC.
% Sequential source scope: lines 8892-8950, Remark II.5.2.1.1 and
% Corollary II.5.2.2 with proof / printed p.98.

\begin{remark}[5.2.1.1]
\label{II.5.2.1.1-zh}
当 $X$ 是诺特预概形时，在 c') 和 d') 的陈述中，可以把“拟凝聚”
换成“凝聚”。事实上，在证明 c') 蕴含 b) 时，$\mathcal{J}$ 和
$\mathcal{J}'$ 此时都是凝聚理想层；另一方面，凝聚模层的每个拟凝聚
子模层都是凝聚的（\egaxref{I.6.1.1-zh}{I, 6.1.1}），结论由此得出。
\end{remark}

\begin{corollary}[5.2.2]
\label{II.5.2.2-zh}
设 $f:X\to Y$ 是分离拟紧态射。下列条件等价：
\begin{enumerate}
  \item[\rm{a)}] $f$ 是仿射态射。
  \item[\rm{b)}] 在拟凝聚 $\mathcal{O}_X$-模层范畴中，函子 $f_*$
  是正合的。
  \item[\rm{c)}] 对每个拟凝聚 $\mathcal{O}_X$-模层
  $\mathcal{F}$，都有 $\mathrm{R}^1f_*(\mathcal{F})=0$。
  \item[\rm{c')}] 对 $\mathcal{O}_X$ 的每个拟凝聚理想层
  $\mathcal{J}$，都有 $\mathrm{R}^1f_*(\mathcal{J})=0$。
\end{enumerate}
\end{corollary}

根据函子 $\mathrm{R}^1f_*$ 的定义（T, 3.7.3），这些条件全都是
$Y$ 上的局部条件，故可设 $Y$ 是仿射的。如果 $f$ 是仿射的，那么
$X$ 也是仿射的，而性质 b) 正是
（\egaxref{I.1.6.4-zh}{I, 1.6.4}）。反过来，来证明 b) 蕴含 a)：
对每个拟凝聚 $\mathcal{O}_X$-模层 $\mathcal{F}$，
$f_*(\mathcal{F})$ 是拟凝聚 $\mathcal{O}_Y$-模层
（\egaxref{I.9.2.2-zh}{I, 9.2.2, a)}）。
% EGA2-CANON-DISP-0011 / EGA2DEF94: guardian-confirmed reversible reader
% correction; French diplomatic source remains unchanged.
按假设，函子 $f_*$ 关于 $\mathcal{F}$ 是正合的\footnote{法文印本此处把
对象 $f_*(\mathcal{F})$ 称为函子；但条件 b) 所断言正合的是直像函子本身，
而给定模层只是其变元。本译依源文勘误记录
EGA2-ZH-DEF-0014 及守护者处置 EGA2-CANON-DISP-0011 作可逆校正；
法文外交式来源保持不变。}；又因为
$Y$ 是仿射的，函子 $\Gamma(Y,\mathcal{G})$ 关于 $\mathcal{G}$
（在拟凝聚 $\mathcal{O}_Y$-模层范畴中）是正合的
（\egaxref{I.1.3.11-zh}{I, 1.3.11}）。因此
$\Gamma(Y,f_*(\mathcal{F}))=\Gamma(X,\mathcal{F})$ 关于
$\mathcal{F}$ 是正合的；根据
（\egaref{II.5.2.1-zh}{5.2.1, c)}），这就证明了所断言的结论。

如果 $f$ 是仿射的，那么对 $Y$ 的每个仿射开集 $U$，
$f^{-1}(U)$ 都是仿射的（\egaref{II.1.3.2-zh}{1.3.2}），因而
$\mathrm{H}^1(f^{-1}(U),\mathcal{F})=0$
（\egaref{II.5.2.1-zh}{5.2.1, d)}）；按定义，这蕴含
$\mathrm{R}^1f_*(\mathcal{F})=0$。最后，设条件 c') 成立；
勒雷谱序列（G, II, 4.17.1 和 I, 4.5.1）的低次项正合列特别给出
正合列
\[
  0\to\mathrm{H}^1(Y,f_*(\mathcal{J}))\to
  \mathrm{H}^1(X,\mathcal{J})\to
  \mathrm{H}^0(Y,\mathrm{R}^1f_*(\mathcal{J})).
\]
由于 $Y$ 是仿射的，且 $f_*(\mathcal{J})$ 是拟凝聚的
（\egaxref{I.9.2.2-zh}{I, 9.2.2, a)}），有
$\mathrm{H}^1(Y,f_*(\mathcal{J}))=0$
（\egaref{II.5.2.1-zh}{5.2.1}）。故假设 c') 蕴含
$\mathrm{H}^1(X,\mathcal{J})=0$；再由
（\egaref{II.5.2.1-zh}{5.2.1}）可知 $X$ 是仿射概形。
% Locked French authority: ega2-1-fr.tex, 820505 B,
% SHA-256 84EDBE3E83530AF2959B441796337C9DC21EAFCA6A13114A26778760FBF437AC.
% Sequential source scope: lines 8951-8976, Corollary II.5.2.3 and
% Remark II.5.2.4 / printed pp.98-99.

\begin{corollary}[5.2.3]
\label{II.5.2.3-zh}
若 $f:X\to Y$ 是仿射态射，则对每个拟凝聚 $\mathcal{O}_X$-模层
$\mathcal{F}$，典范同态
$\mathrm{H}^1(Y,f_*(\mathcal{F}))\to\mathrm{H}^1(X,\mathcal{F})$
是双射。
\end{corollary}

\oldpage[II]{99}

事实上，勒雷谱序列的低次项给出正合列
\[
  0\to\mathrm{H}^1(Y,f_*(\mathcal{F}))\to
  \mathrm{H}^1(X,\mathcal{F})\to
  \mathrm{H}^0(Y,\mathrm{R}^1f_*(\mathcal{F}))
\]
而结论由（\egaref{II.5.2.2-zh}{5.2.2}）得出。

\begin{remark}[5.2.4]
\label{II.5.2.4-zh}
在第三章第 1 节中，我们将证明：若 $X$ 是仿射的，则对每个
$i>0$ 以及每个拟凝聚 $\mathcal{O}_X$-模层 $\mathcal{F}$，都有
$\mathrm{H}^i(X,\mathcal{F})=0$。
\end{remark}
% Locked French authority: ega2-1-fr.tex, 820505 B,
% SHA-256 84EDBE3E83530AF2959B441796337C9DC21EAFCA6A13114A26778760FBF437AC.
% Sequential source scope: lines 8978-9000, subsection II.5.3 opening and
% Definition II.5.3.1 / printed p.99.

\subsection{拟射影态射。}
\label{subsection:II.5.3-zh}

\begin{definition}[5.3.1]
\label{II.5.3.1-zh}
如果态射 $f:X\to Y$ 是有限型的，并且存在一个相对于 $f$ 丰沛的
可逆 $\mathcal{O}_X$-模层，就称 $f$ 是\emph{拟射影态射}；也称
$X$（经由 $f$ 视为 $Y$-预概形）\emph{在 $Y$ 上拟射影}，或称 $X$
是\emph{拟射影 $Y$-概形}。
\end{definition}

应当注意，这一概念不具有在 $Y$ 上的局部性：永田 [26] 和广中的
反例表明，即使 $X$ 和 $Y$ 是代数闭域上的非奇异代数概形，$Y$ 的
每个点也都可以有一个仿射邻域 $U$，使得 $f^{-1}(U)$ 在 $U$ 上
拟射影，而 $f$ 本身却不拟射影。

还应注意，拟射影态射必定是分离的
（\egaref{II.4.6.1-zh}{4.6.1}）。当 $Y$ 拟紧时，说 $f$ 拟射影，
等价于说 $f$ 是有限型的，并且存在一个相对于 $f$ 极丰沛的
$\mathcal{O}_X$-模层（\egaref{II.4.6.2-zh}{4.6.2} 和
\egaref{II.4.6.11-zh}{4.6.11}）。此外：

% Locked French authority: ega2-1-fr.tex, 820505 B,
% SHA-256 84EDBE3E83530AF2959B441796337C9DC21EAFCA6A13114A26778760FBF437AC.
% Sequential source scope: lines 9002-9027, Proposition II.5.3.2 and
% its consequences / printed p.99.

\begin{proposition}[5.3.2]
\label{II.5.3.2-zh}
设 $Y$ 是拟紧概形，或底空间为诺特空间的预概形，并设 $X$ 是
$Y$-预概形。下列条件等价：
\begin{enumerate}
  \item[\rm{a)}] $X$ 是拟射影 $Y$-概形。
  \item[\rm{b)}] $X$ 在 $Y$ 上是有限型的，并且存在有限型拟凝聚
  $\mathcal{O}_Y$-模层 $\mathcal{E}$，使 $X$ 在 $Y$ 上同构于
  $\mathbf{P}(\mathcal{E})$ 的一个子预概形。
  \item[\rm{c)}] $X$ 在 $Y$ 上是有限型的，并且存在拟凝聚分次
  $\mathcal{O}_Y$-代数 $\mathcal{S}$，使 $\mathcal{S}_1$ 为有限型
  并生成 $\mathcal{S}$，且 $X$ 在 $Y$ 上同构于
  $\operatorname{Proj}(\mathcal{S})$ 的某个处处稠密开集上所诱导的
  子预概形。
\end{enumerate}
\end{proposition}

这立即由上述说明以及（\egaref{II.4.4.3-zh}{4.4.3}）、
（\egaref{II.4.4.6-zh}{4.4.6}）和
（\egaref{II.4.4.7-zh}{4.4.7}）得出。

应当注意，当 $Y$ 是诺特预概形时，在
（\egaref{II.5.3.2-zh}{5.3.2}）的条件 b) 和 c) 中，可以删去
$X$ 在 $Y$ 上有限型这一假设，因为此时该条件自动成立
（\egaxref{I.6.3.5-zh}{I, 6.3.5}）。

% Locked French authority: ega2-1-fr.tex, 820505 B,
% SHA-256 84EDBE3E83530AF2959B441796337C9DC21EAFCA6A13114A26778760FBF437AC.
% Sequential source scope: lines 9029-9044, Corollary II.5.3.3 and proof
% / printed p.99.

\begin{corollary}[5.3.3]
\label{II.5.3.3-zh}
设 $Y$ 是拟紧概形，并且存在一个丰沛的 $\mathcal{O}_Y$-模层
$\mathcal{L}$（\egaref{II.4.5.3-zh}{4.5.3}）。一个 $Y$-概形
$X$ 拟射影的充要条件是：它在 $Y$ 上为有限型，并同构于形如
$\mathbf{P}_Y^r$ 的射影丛的一个 $Y$-子概形。
\end{corollary}

事实上，若 $\mathcal{E}$ 是有限型拟凝聚 $\mathcal{O}_Y$-模层，
则 $\mathcal{E}$ 同构于某个 $\mathcal{O}_Y$-模层
$\mathcal{L}^{\otimes(-n)}\otimes_{\mathcal{O}_Y}\mathcal{O}_Y^k$
的商（\egaref{II.4.5.5-zh}{4.5.5}），故
$\mathbf{P}(\mathcal{E})$ 同构于 $\mathbf{P}_Y^{k-1}$ 的一个闭子概形
（\egaref{II.4.1.2-zh}{4.1.2} 和
\egaref{II.4.1.4-zh}{4.1.4}）。

% Locked French authority: ega2-1-fr.tex, 820505 B,
% SHA-256 84EDBE3E83530AF2959B441796337C9DC21EAFCA6A13114A26778760FBF437AC.
% Sequential source scope: lines 9046-9089, Proposition II.5.3.4 through
% Corollary II.5.3.6 / printed pp.99-100.

\begin{proposition}[5.3.4]
\label{II.5.3.4-zh}
\begin{enumerate}
  \item[\rm{(i)}] 有限型拟仿射态射（特别地，拟紧浸入或有限型仿射态射）
  是拟射影的。
  \item[\rm{(ii)}] 若 $f:X\to Y$ 和 $g:Y\to Z$ 都是拟射影的，并且
  $Z$ 拟紧，则 $g\circ f$ 是拟射影的。

\oldpage[II]{100}

  \item[\rm{(iii)}] 若 $f:X\to Y$ 是拟射影 $S$-态射，则对底预概形的
  每个扩张 $S'\to S$，$f_{(S')}:X_{(S')}\to Y_{(S')}$ 都是拟射影的。
  \item[\rm{(iv)}] 若 $f:X\to Y$ 和 $g:X'\to Y'$ 是两个拟射影
  $S$-态射，则 $f\times_Sg$ 是拟射影的。
  \item[\rm{(v)}] 若 $f:X\to Y$、$g:Y\to Z$ 是两个态射，使
  $g\circ f$ 拟射影，并且 $g$ 是分离的或 $X$ 局部诺特，则
  $f$ 是拟射影的。
  \item[\rm{(vi)}] 若 $f$ 是拟射影态射，则 $f_{\mathrm{red}}$
  也是拟射影态射。
\end{enumerate}
\end{proposition}

(i) 由（\egaref{II.5.1.6-zh}{5.1.6}）和
（\egaref{II.5.1.10-zh}{5.1.10, (i)}）得出。其余断言都是定义
（\egaref{II.5.3.1-zh}{5.3.1}）、有限型态射的性质
（\egaxref{I.6.3.4-zh}{I, 6.3.4}）以及
（\egaref{II.4.6.13-zh}{4.6.13}）的直接推论。

\begin{remark}[5.3.5]
\label{II.5.3.5-zh}
应当注意，可能有 $f_{\mathrm{red}}$ 拟射影而 $f$ 不拟射影的情形；
即使假设 $Y$ 是有限维 $\mathbf{C}$-代数的谱，并且 $f$ 是紧合态射，
这种情形仍可能发生。
\end{remark}

\begin{corollary}[5.3.6]
\label{II.5.3.6-zh}
若 $X$ 和 $X'$ 是两个拟射影 $Y$-概形，则 $X\amalg X'$ 是拟射影
$Y$-概形。
\end{corollary}

这由（\egaref{II.4.6.18-zh}{4.6.18}）得出。
% Locked French authority: ega2-1-fr.tex, 820505 B,
% SHA-256 84EDBE3E83530AF2959B441796337C9DC21EAFCA6A13114A26778760FBF437AC.
% Sequential source scope: lines 9091-9120, subsection II.5.4 opening and
% Definition II.5.4.1 / printed p.100.

\subsection{紧合态射和普遍闭态射。}
\label{subsection:II.5.4-zh}
\label{II.5.4-zh}

\begin{definition}[5.4.1]
\label{II.5.4.1-zh}
若预概形态射 $f:X\to Y$ 满足下列两个条件，就称它是\emph{紧合态射}：
\begin{enumerate}
  \item[\rm{a)}] $f$ 是分离的有限型态射。
  \item[\rm{b)}] 对每个预概形 $Y'$ 以及每个态射 $Y'\to Y$，投影
  $f_{(Y')}:X\times_YY'\to Y'$ 都是闭态射
  （\egaxref{I.2.2.6-zh}{I, 2.2.6}）。
\end{enumerate}
在这种情况下，也称 $X$（视为结构态射为 $f$ 的 $Y$-预概形）
\emph{在 $Y$ 上紧合}。
\end{definition}

显然，条件 a) 和 b) 都是 $Y$ 上的局部条件。为了验证 $X\times_YY'$
的闭子集 $Z$ 在投影 $q:X\times_YY'\to Y'$ 下的像在 $Y'$ 中闭，
只须验证：对 $Y'$ 的每个仿射开集 $U'$，$q(Z)\cap U'$ 在 $U'$ 中闭。
由于 $q(Z)\cap U'=q(Z\cap q^{-1}(U'))$，且 $q^{-1}(U')$ 等同于
$X\times_YU'$（\egaxref{I.4.4.1-zh}{I, 4.4.1}），可见在验证定义
（\egaref{II.5.4.1-zh}{5.4.1}）的条件 b) 时，可以限于 $Y'$ 是
仿射概形的情形。后面将看到（\egaref{II.5.6.3-zh}{5.6.3}）：
若 $Y$ 局部诺特，甚至只须在 $Y'$ 在 $Y$ 上有限型时验证 b)。

显然，每个紧合态射都是\emph{闭态射}。

% Locked French authority: ega2-1-fr.tex, 820505 B,
% SHA-256 84EDBE3E83530AF2959B441796337C9DC21EAFCA6A13114A26778760FBF437AC.
% Sequential source scope: lines 9122-9160, Proposition II.5.4.2 and proof
% / printed pp.100-101.

\begin{proposition}[5.4.2]
\label{II.5.4.2-zh}
\begin{enumerate}
  \item[\rm{(i)}] 闭浸入是紧合态射。
  \item[\rm{(ii)}] 两个紧合态射的复合是紧合的。
  \item[\rm{(iii)}] 若 $X$、$Y$ 是 $S$-预概形，而
  $f:X\to Y$ 是紧合 $S$-态射，则对底预概形的每个扩张
  $S'\to S$，$f_{(S')}:X_{(S')}\to Y_{(S')}$ 都是紧合的。
  \item[\rm{(iv)}] 若 $f:X\to Y$ 和 $g:X'\to Y'$ 是两个紧合
  $S$-态射，则 $S$-态射
  $f\times_Sg:X\times_SX'\to Y\times_SY'$ 是紧合的。
\end{enumerate}
\end{proposition}

\oldpage[II]{101}

只须证明 (i)、(ii) 和 (iii)（\egaxref{I.3.5.1-zh}{I, 3.5.1}）。
在这三种情形的每一种中，定义（\egaref{II.5.4.1-zh}{5.4.1}）的
条件 a) 都由先前的结果
（\egaxref{I.5.5.1-zh}{I, 5.5.1} 和
\egaxref{I.6.3.4-zh}{6.3.4}）推出；只剩验证条件 b)。在情形 (i) 中，
这是显然的，因为若 $X\to Y$ 是闭浸入，则
$X\times_YY'\to Y\times_YY'=Y'$ 也是闭浸入
（\egaxref{I.4.3.2-zh}{I, 4.3.2} 和
\egaref{II.3.3.3-zh}{3.3.3}）。为了证明 (ii)，考虑两个紧合态射
$X\to Y$、$Y\to Z$，以及一个态射 $Z'\to Z$。可以写成
$X\times_ZZ'=X\times_Y(Y\times_ZZ')$
（\egaxref{I.3.3.9.1-zh}{I, 3.3.9.1}），故投影
$X\times_ZZ'\to Z'$ 分解为
$X\times_Y(Y\times_ZZ')\to Y\times_ZZ'\to Z'$。考虑到开头的说明，
(ii) 便由两个闭态射的复合是闭态射这一事实得出。最后，对每个态射
$S'\to S$，$X_{(S')}$ 等同于 $X\times_YY_{(S')}$
（\egaxref{I.3.3.11-zh}{I, 3.3.11}）；对每个态射
$Z\to Y_{(S')}$，可以写成
\[
  X_{(S')}\times_{Y_{(S')}}Z
  =(X\times_YY_{(S')})\times_{Y_{(S')}}Z=X\times_YZ~;
\]
依假设，$X\times_YZ\to Z$ 是闭态射，这就证明了 (iii)。

% Locked French authority: ega2-1-fr.tex, 820505 B,
% SHA-256 84EDBE3E83530AF2959B441796337C9DC21EAFCA6A13114A26778760FBF437AC.
% Sequential source scope: lines 9162-9190, Corollary II.5.4.3 and proof
% / printed p.101.

\begin{corollary}[5.4.3]
\label{II.5.4.3-zh}
设 $f:X\to Y$、$g:Y\to Z$ 是两个态射，并且 $g\circ f$ 紧合。
\begin{enumerate}
  \item[\rm{(i)}] 若 $g$ 分离，则 $f$ 紧合。
  \item[\rm{(ii)}] 若 $g$ 分离且为有限型，而 $f$ 是满射，则 $g$ 紧合。
\end{enumerate}
\end{corollary}

(i) 由（\egaref{II.5.4.2-zh}{5.4.2}）经由一般方法
（\egaxref{I.5.5.12-zh}{I, 5.5.12}）得出。为了证明 (ii)，只须验证
定义（\egaref{II.5.4.1-zh}{5.4.1}）的条件 b)。对每个态射
$Z'\to Z$，交换图
\[
  \xymatrix{
    X\times_ZZ'\ar[r]^{f\times1_{Z'}}\ar[dr]_p &
    Y\times_ZZ'\ar[d]^{p'}\\
    & Z'
  }
\]
（其中 $p$ 和 $p'$ 是投影）是交换的
（\egaxref{I.3.2.1-zh}{I, 3.2.1}）；此外，因为 $f$ 是满射，
$f\times1_{Z'}$ 也是满射（\egaxref{I.3.5.2-zh}{I, 3.5.2}），而
依假设 $p$ 是闭态射。于是，$Y\times_ZZ'$ 的每个闭子集 $F$ 都是
$X\times_ZZ'$ 的某个闭子集 $E$ 在 $f\times1_{Z'}$ 下的像，故依假设
$p'(F)=p(E)$ 在 $Z'$ 中闭，推论得证。

% Locked French authority: ega2-1-fr.tex, 820505 B,
% SHA-256 84EDBE3E83530AF2959B441796337C9DC21EAFCA6A13114A26778760FBF437AC.
% Sequential source scope: lines 9192-9234, Corollaries II.5.4.4-5 and
% their proofs / printed pp.101-102.

\begin{corollary}[5.4.4]
\label{II.5.4.4-zh}
若 $X$ 是在 $Y$ 上紧合的预概形，而 $\mathcal{S}$ 是拟凝聚
$\mathcal{O}_Y$-代数，则每个 $Y$-态射
$f:X\to\operatorname{Proj}(\mathcal{S})$ 都是紧合的（从而更是闭的）。
\end{corollary}

事实上，结构态射 $p:\operatorname{Proj}(\mathcal{S})\to Y$ 是分离的，
而依假设 $p\circ f$ 紧合。

\begin{corollary}[5.4.5]
\label{II.5.4.5-zh}
设 $f:X\to Y$ 是分离的有限型态射。设
$(X_i)_{1\leq i\leq n}$（相应地，$(Y_i)_{1\leq i\leq n}$）是
$X$（相应地，$Y$）的一族有限多个闭子预概形，$j_i$（相应地，$h_i$）
是典范嵌入 $X_i\to X$（相应地，$Y_i\to Y$）。假设 $X$ 的底空间
是诸 $X_i$ 之并，并且对每个 $i$，存在态射 $f_i:X_i\to Y_i$，使图
\[
  \xymatrix{
    X_i\ar[r]^{f_i}\ar[d]_{j_i} & Y_i\ar[d]^{h_i}\\
    X\ar[r]_f & Y
  }
\]
交换。那么，$f$ 紧合的充要条件是每个 $f_i$ 都紧合。
\end{corollary}

\oldpage[II]{102}

若 $f$ 紧合，则 $f\circ j_i$ 也紧合，因为 $j_i$ 是闭浸入
（\egaref{II.5.4.2-zh}{5.4.2}）；而 $h_i$ 是闭浸入，因而是分离态射，
所以由（\egaref{II.5.4.3-zh}{5.4.3}）可知 $f_i$ 紧合。反过来，
设每个 $f_i$ 都紧合，并考虑由诸 $X_i$ 之和所得的预概形 $Z$；令
$u$ 为态射 $Z\to X$，其在每个 $X_i$ 上就是 $j_i$。$f\circ u$
在每个 $X_i$ 上的限制等于 $f\circ j_i=h_i\circ f_i$；因为 $h_i$
和 $f_i$ 都紧合，所以该限制紧合
（\egaref{II.5.4.2-zh}{5.4.2}）。于是由定义
（\egaref{II.5.4.1-zh}{5.4.1}）立即可知 $f\circ u$ 紧合。又因依假设
$u$ 是满射，由（\egaref{II.5.4.3-zh}{5.4.3}）可知 $f$ 紧合。

% Locked French authority: ega2-1-fr.tex, 820505 B,
% SHA-256 84EDBE3E83530AF2959B441796337C9DC21EAFCA6A13114A26778760FBF437AC.
% Sequential source scope: lines 9236-9263, Corollary II.5.4.6 and
% reduction II.5.4.7 / printed p.102.

\begin{corollary}[5.4.6]
\label{II.5.4.6-zh}
设 $f:X\to Y$ 是分离的有限型态射；$f$ 紧合的充要条件是
$f_{\mathrm{red}}:X_{\mathrm{red}}\to Y_{\mathrm{red}}$ 紧合。
\end{corollary}

这是（\egaref{II.5.4.5-zh}{5.4.5}）在 $n=1$、
$X_1=X_{\mathrm{red}}$ 和 $Y_1=Y_{\mathrm{red}}$ 时的特例
（\egaxref{I.5.1.5-zh}{I, 5.1.5}）。

\begin{env}[5.4.7]
\label{II.5.4.7-zh}
若 $X$ 和 $Y$ 是诺特预概形，而 $f:X\to Y$ 是分离的有限型态射，
则为了验证 $f$ 紧合，可以归结为验证整预概形之间的支配态射紧合。
事实上，设 $X_i$（$1\leq i\leq n$）是 $X$ 的有限多个不可约分支；
对每个 $i$，考虑 $X$ 的唯一约化闭子预概形，其底空间为 $X_i$，
仍将它记作 $X_i$（\egaxref{I.5.2.1-zh}{I, 5.2.1}）。另一方面，
令 $Y_i$ 为 $Y$ 的唯一约化闭子预概形，其底空间为
$\overline{f(X_i)}$。若 $g_i$（相应地，$h_i$）是嵌入态射
$X_i\to X$（相应地，$Y_i\to Y$），则有
$f\circ g_i=h_i\circ f_i$，其中 $f_i$ 是支配态射 $X_i\to Y_i$
（\egaxref{I.5.2.2-zh}{I, 5.2.2}）。这便满足
（\egaref{II.5.4.5-zh}{5.4.5}）的应用条件；因此 $f$ 紧合的充要条件
是所有 $f_i$ 都紧合。
\end{env}

% Locked French authority: ega2-1-fr.tex, 820505 B,
% SHA-256 84EDBE3E83530AF2959B441796337C9DC21EAFCA6A13114A26778760FBF437AC.
% Sequential source scope: lines 9265-9292, Corollary II.5.4.8 and proof
% / printed p.102.

\begin{corollary}[5.4.8]
\label{II.5.4.8-zh}
设 $X$、$Y$ 是在 $S$ 上分离且有限型的 $S$-预概形，并设
$f:X\to Y$ 是 $S$-态射。$f$ 紧合的充要条件是：对每个
$S$-预概形 $S'$，态射
$f\times_S1_{S'}:X\times_SS'\to Y\times_SS'$ 都是闭态射。
\end{corollary}

先注意，若 $g:X\to S$ 和 $h:Y\to S$ 是结构态射，则按定义
$g=h\circ f$，故 $f$ 分离且为有限型
（\egaxref{I.5.5.1-zh}{I, 5.5.1} 和
\egaxref{I.6.3.4-zh}{6.3.4}）。若 $f$ 紧合，则
$f\times_S1_{S'}$ 也紧合（\egaref{II.5.4.2-zh}{5.4.2}）；从而
$f\times_S1_{S'}$ 更是闭态射。反过来，设陈述中的条件成立，并设
$Y'$ 是 $Y$-预概形；
$Y'$ 也可视为 $S$-预概形。由于 $Y\to S$ 分离，$X\times_YY'$
等同于 $X\times_SY'$ 的一个闭子预概形
（\egaxref{I.5.4.2-zh}{I, 5.4.2}）。在交换图
\[
  \xymatrix{
    X\times_YY'\ar[r]^{f\times1_{Y'}}\ar[d] &
    Y\times_YY'=Y'\ar[d]\\
    X\times_SY'\ar[r]_{f\times_S1_{Y'}} & Y\times_SY'
  }
\]
中，竖直箭头都是闭浸入。因此立即可知：若 $f\times_S1_{Y'}$
是闭态射，则 $f\times_Y1_{Y'}$ 也是闭态射。
% Locked French authority: ega2-1-fr.tex, 820505 B,
% SHA-256 84EDBE3E83530AF2959B441796337C9DC21EAFCA6A13114A26778760FBF437AC.
% Sequential source scope: lines 9294-9309, Remark II.5.4.9
% / printed pp.102-103.

\begin{remark}[5.4.9]
\label{II.5.4.9-zh}
若态射 $f:X\to Y$ 满足定义（\egaref{II.5.4.1-zh}{5.4.1}）的条件
b)，就称它是\emph{普遍闭态射}。读者会注意到：

\oldpage[II]{103}

在命题（\egaref{II.5.4.2-zh}{5.4.2}）至
（\egaref{II.5.4.8-zh}{5.4.8}）中，可以处处把“紧合”换成“普遍闭”，
而各结果仍然成立；并且在
（\egaref{II.5.4.3-zh}{5.4.3}）、
（\egaref{II.5.4.5-zh}{5.4.5}）、
（\egaref{II.5.4.6-zh}{5.4.6}）和
（\egaref{II.5.4.8-zh}{5.4.8}）的假设中，可以省去有限性条件。
\end{remark}

% Locked French authority: ega2-1-fr.tex, 820505 B,
% SHA-256 84EDBE3E83530AF2959B441796337C9DC21EAFCA6A13114A26778760FBF437AC.
% Sequential source scope: lines 9311-9354, construction II.5.4.10 and
% its two consequences / printed p.103.

\begin{env}[5.4.10]
\label{II.5.4.10-zh}
设 $f:X\to Y$ 是有限型态射。若 $f$ 在 $X$ 的某个以 $Z$ 为底空间的
闭子预概形（\egaxref{I.5.2.1-zh}{I, 5.2.1}）上的限制紧合，就称 $X$
的闭子集 $Z$ \emph{在 $Y$ 上紧合}（或称\emph{$Y$-紧合}，或称
\emph{相对于 $f$ 紧合}）。由于这个限制此时是分离的，由
（\egaref{II.5.4.6-zh}{5.4.6}）和
（\egaxref{I.5.5.1-zh}{I, 5.5.1, (vi)}）可知，上述性质与选择哪一个
以 $Z$ 为底空间的 $X$ 的闭子预概形无关。若 $g:X'\to X$ 是紧合态射，
则 $g^{-1}(Z)$ 是 $X'$ 的紧合闭子集：事实上，若 $T$ 是 $X$ 的一个
以 $Z$ 为底空间的闭子预概形，只须注意到 $g$ 在 $X'$ 的闭子预概形
$g^{-1}(T)$ 上的限制是紧合态射 $g^{-1}(T)\to T$
（\egaref{II.5.4.2-zh}{5.4.2, (iii)}），然后应用
（\egaref{II.5.4.2-zh}{5.4.2, (ii)}）。另一方面，若 $X''$ 是有限型
$Y$-概形，而 $u:X\to X''$ 是 $Y$-态射，则 $u(Z)$ 是 $X''$ 的紧合
闭子集。事实上，取 $T$ 为 $X$ 的以 $Z$ 为底空间的约化闭子预概形；
因为 $f$ 在 $T$ 上的限制紧合，由
（\egaref{II.5.4.3-zh}{5.4.3, (i)}）可知 $u$ 在 $T$ 上的限制也紧合，
故 $u(Z)$ 在 $X''$ 中闭。设 $T''$ 是 $X''$ 的一个以 $u(Z)$ 为底空间
的闭子预概形（\egaxref{I.5.2.1-zh}{I, 5.2.1}），使得 $u|T$ 分解为
\[
  T\xrightarrow{v}T''\xrightarrow{j}X'',
\]
其中 $j$ 是典范嵌入（\egaxref{I.5.2.2-zh}{I, 5.2.2}），而 $v$
因而是紧合满射（\egaref{II.5.4.5-zh}{5.4.5}）。若 $g$ 是结构态射
$X''\to Y$ 在 $T''$ 上的限制，则 $g$ 分离且为有限型，并有
$f|T=g\circ v$；故由（\egaref{II.5.4.3-zh}{5.4.3, (ii)}）可知
$g$ 紧合，这就证明了所断言的结论。
\end{env}

特别地，由这些说明可知：若 $Z$ 是 $X$ 的 $Y$-紧合闭子集，则：
\begin{enumerate}
  \item[\rm{1\textsuperscript{o}}] 对 $X$ 的每个闭子预概形 $X'$，
  $Z\cap X'$ 都是 $X'$ 的 $Y$-紧合闭子集。
  \item[\rm{2\textsuperscript{o}}] 若 $X$ 是有限型 $Y$-概形 $X''$
  的一个子预概形，则 $Z$ 也是 $X''$ 的 $Y$-紧合闭子集（特别地，
  在 $X''$ 中闭）。
\end{enumerate}
% Locked French authority: ega2-1-fr.tex, 820505 B,
% SHA-256 84EDBE3E83530AF2959B441796337C9DC21EAFCA6A13114A26778760FBF437AC.
% Sequential source scope: lines 9356-9381, subsection II.5.5 opening and
% Proposition II.5.5.1 / printed p.103.

\subsection{射影态射。}
\label{subsection:II.5.5-zh}

\begin{proposition}[5.5.1]
\label{II.5.5.1-zh}
设 $X$ 是 $Y$-预概形。下列条件等价：
\begin{enumerate}
  \item[\rm{a)}] $X$ 在 $Y$ 上同构于某个射影丛
  $\mathbf{P}(\mathcal{E})$ 的一个闭子预概形，其中 $\mathcal{E}$
  是有限型拟凝聚 $\mathcal{O}_Y$-模层。
  \item[\rm{b)}] 存在拟凝聚分次 $\mathcal{O}_Y$-代数
  $\mathcal{S}$，使 $\mathcal{S}_1$ 为有限型并生成 $\mathcal{S}$，
  且 $X$ 在 $Y$ 上同构于 $\operatorname{Proj}(\mathcal{S})$。
\end{enumerate}
\end{proposition}

由（\egaref{II.3.6.2-zh}{3.6.2, (ii)}），条件 a) 蕴含 b)：若
$\mathcal{J}$ 是 $\mathbf{S}(\mathcal{E})$ 的拟凝聚分次理想层，则
拟凝聚分次 $\mathcal{O}_Y$-代数
$\mathcal{S}=\mathbf{S}(\mathcal{E})/\mathcal{J}$ 由
$\mathcal{S}_1$ 生成，而后者是 $\mathcal{E}$ 的典范像，故为有限型
$\mathcal{O}_Y$-模层。将（\egaref{II.3.6.2-zh}{3.6.2}）应用于
$\mathcal{M}\to\mathcal{S}_1$ 为恒等映射的情形，可知条件 b) 蕴含 a)。

% Locked French authority: ega2-1-fr.tex, 820505 B,
% SHA-256 84EDBE3E83530AF2959B441796337C9DC21EAFCA6A13114A26778760FBF437AC.
% Sequential source scope: lines 9383-9442, Definition II.5.5.2 and
% Theorem II.5.5.3 with proof / printed pp.104.

\oldpage[II]{104}

\begin{definition}[5.5.2]
\label{II.5.5.2-zh}
若 $Y$-预概形 $X$ 满足（\egaref{II.5.5.1-zh}{5.5.1}）中等价的
条件 a) 和 b)，就称它\emph{在 $Y$ 上射影}，或称它是
\emph{射影 $Y$-概形}。若态射 $f:X\to Y$ 使 $X$ 成为射影
$Y$-概形，就称该态射是\emph{射影态射}。
\end{definition}

显然，若 $f:X\to Y$ 是射影态射，则存在一个\emph{相对于 $f$
极丰沛的} $\mathcal{O}_X$-模层
（\egaref{II.4.4.2-zh}{4.4.2}）。

\begin{theorem}[5.5.3]
\label{II.5.5.3-zh}
\begin{enumerate}
  \item[\rm{(i)}] 每个射影态射都是拟射影且紧合的。
  \item[\rm{(ii)}] 反过来，设 $Y$ 是拟紧概形，或底空间为诺特空间的
  预概形；则每个既拟射影又紧合的态射 $f:X\to Y$ 都是射影态射。
\end{enumerate}
\end{theorem}

(i) 显然，若 $f:X\to Y$ 射影，则它为有限型且拟射影（因而特别
分离）。另一方面，由（\egaref{II.5.5.1-zh}{5.5.1, b)}）和
（\egaref{II.3.5.3-zh}{3.5.3}）立即可知：若 $f$ 射影，则对每个态射
$Y'\to Y$，$f\times_Y1_{Y'}:X\times_YY'\to Y'$ 也射影。因此，要证明
$f$ 普遍闭，只须证明射影态射 $f$ 是\emph{闭态射}。这个问题在 $Y$
上是局部的，故可设 $Y=\operatorname{Spec}(A)$，从而由
（\egaref{II.5.5.1-zh}{5.5.1}）可设
$X=\operatorname{Proj}(S)$，其中 $S$ 是由 $S_1$ 中有限多个元素生成的
分次 $A$-代数。对每个 $y\in Y$，纤维 $f^{-1}(y)$ 等同于
$\operatorname{Proj}(S)\times_Y\operatorname{Spec}(k(y))$
（\egaxref{I.3.6.1-zh}{I, 3.6.1}），从而等同于
$\operatorname{Proj}(S\otimes_Ak(y))$
（\egaref{II.2.8.10-zh}{2.8.10}）。因此，$f^{-1}(y)$ 为空的充要条件
是 $S\otimes_Ak(y)$ 满足条件 (TN)
（\egaref{II.2.7.4-zh}{2.7.4}），换言之，对充分大的 $n$，有
$S_n\otimes_Ak(y)=0$。但由于 $(S_n)_y$ 是有限型
$\mathcal{O}_y$-模，按中山引理，上述条件等价于对充分大的 $n$ 有
$(S_n)_y=0$。若 $\mathfrak a_n$ 是 $A$-模 $S_n$ 在 $A$ 中的零化子，
则上述条件又等价于：对充分大的 $n$，
$\mathfrak a_n\not\subset\mathfrak j_y$
（\egaxref{0.1.7.4-zh}{0, 1.7.4}）。又因依假设
$S_nS_1=S_{n+1}$，有 $\mathfrak a_n\subset\mathfrak a_{n+1}$；若
$\mathfrak a$ 是诸 $\mathfrak a_n$ 之并，便得到
$f(X)=V(\mathfrak a)$，这证明了 $f(X)$ 在 $Y$ 中闭。现在，若
$X'$ 是 $X$ 的任意闭子集，则存在一个以 $X'$ 为底空间的 $X$ 的
闭子预概形（\egaxref{I.5.2.1-zh}{I, 5.2.1}）；并且由
（\egaref{II.5.5.1-zh}{5.5.1, a)}）显然可知，态射
$X'\to X\xrightarrow{f}Y$ 射影，故 $f(X')$ 在 $Y$ 中闭。

(ii) 关于 $Y$ 的假设和 $f$ 拟射影这一事实蕴含：存在有限型拟凝聚
$\mathcal{O}_Y$-模层 $\mathcal{E}$ 以及 $Y$-浸入
$j:X\to\mathbf{P}(\mathcal{E})$
（\egaref{II.5.3.2-zh}{5.3.2}）。但因为 $f$ 紧合，由
（\egaref{II.5.4.4-zh}{5.4.4}）可知 $j$ 是闭浸入，因此 $f$ 射影。
% Locked French authority: ega2-1-fr.tex, 820505 B,
% SHA-256 84EDBE3E83530AF2959B441796337C9DC21EAFCA6A13114A26778760FBF437AC.
% Sequential source scope: lines 9444-9485, Remarks II.5.5.4
% / printed pp.104-105.

\par\medskip\noindent\textit{说明（5.5.4）。---\ }
\label{II.5.5.4-zh}
\begin{enumerate}
  \item[\rm{(i)}] 设态射 $f:X\to Y$ 满足：
  1\textsuperscript{o} $f$ 紧合；2\textsuperscript{o} 存在一个相对于
  $f$ 极丰沛的 $\mathcal{O}_X$-模层 $\mathcal{L}$；
  3\textsuperscript{o} 拟凝聚 $\mathcal{O}_Y$-模层
  $\mathcal{E}=f_*(\mathcal{L})$ 为有限型。则 $f$ 是射影态射：
  事实上，由（\egaref{II.4.4.4-zh}{4.4.4}）可得 $Y$-浸入
  $r:X\to\mathbf{P}(\mathcal{E})$；又因 $f$ 紧合，由
  （\egaref{II.5.4.4-zh}{5.4.4}）可知 $r$ 是闭浸入。第三章第 3 节
  将证明：当 $Y$ 局部诺特时，上述条件 3\textsuperscript{o} 是另外
  两个条件的推论。因此，在这种情况下，条件 1\textsuperscript{o}
  和 2\textsuperscript{o} 刻画射影态射；若 $Y$ 拟紧，则可将条件
  2\textsuperscript{o} 换成存在一个相对于 $f$ 丰沛的
  $\mathcal{O}_X$-模层这一假设
  （\egaref{II.4.6.11-zh}{4.6.11}）。
  \item[\rm{(ii)}] 设 $Y$ 是拟紧概形，并存在一个丰沛的
  $\mathcal{O}_Y$-模层。$Y$-概形 $X$ 射影的充要条件是：它在 $Y$
  上同构于形如 $\mathbf{P}_Y^r$ 的射影丛的一个闭 $Y$-子概形。
  该条件显然充分。

\oldpage[II]{105}

  反过来，若 $X$ 在 $Y$ 上射影，则它拟射影，故存在一个从 $X$ 到
  某个 $\mathbf{P}_Y^r$ 的 $Y$-浸入 $j$
  （\egaref{II.5.3.3-zh}{5.3.3}）；由
  （\egaref{II.5.4.4-zh}{5.4.4}）和
  （\egaref{II.5.5.3-zh}{5.5.3}）可知该浸入是闭的。
  \item[\rm{(iii)}] （\egaref{II.5.5.3-zh}{5.5.3}）中的论证表明：
  对每个预概形 $Y$ 和每个整数 $r\geq0$，结构态射
  $\mathbf{P}_Y^r\to Y$ 都是满射。事实上，若置
  $\mathcal{S}=\mathbf{S}_{\mathcal{O}_Y}(\mathcal{O}_Y^{r+1})$，
  则显然有 $\mathcal{S}_y=\mathbf{S}_{k(y)}((k(y))^{r+1})$
  （\egaref{II.1.7.3-zh}{1.7.3}），故对每个 $y\in Y$ 和每个
  $n\geq0$，都有 $(\mathcal{S}_n)_y\neq0$。
  \item[\rm{(iv)}] 由永田的例子 \cite{II-26} 可知，存在不是
  拟射影态射的紧合态射。
\end{enumerate}

% Locked French authority: ega2-1-fr.tex, 820505 B,
% SHA-256 84EDBE3E83530AF2959B441796337C9DC21EAFCA6A13114A26778760FBF437AC.
% Sequential source scope: lines 9487-9543, Proposition II.5.5.5 and proof
% / printed p.105.

\begin{proposition}[5.5.5]
\label{II.5.5.5-zh}
\begin{enumerate}
  \item[\rm{(i)}] 闭浸入是射影态射。
  \item[\rm{(ii)}] 若 $f:X\to Y$ 和 $g:Y\to Z$ 是射影态射，并且
  $Z$ 是拟紧概形，或底空间为诺特空间的预概形，则 $g\circ f$
  是射影态射。
  \item[\rm{(iii)}] 若 $f:X\to Y$ 是射影 $S$-态射，则对底预概形的
  每个扩张 $S'\to S$，$f_{(S')}:X_{(S')}\to Y_{(S')}$ 都是射影态射。
  \item[\rm{(iv)}] 若 $f:X\to Y$ 和 $g:X'\to Y'$ 是射影
  $S$-态射，则 $f\times_Sg$ 也是射影态射。
  \item[\rm{(v)}] 若 $g\circ f$ 是射影态射，并且 $g$ 分离，则
  $f$ 是射影态射。
  \item[\rm{(vi)}] 若 $f$ 是射影态射，则 $f_{\mathrm{red}}$
  也是射影态射。
\end{enumerate}
\end{proposition}

(i) 立即由（\egaref{II.3.1.7-zh}{3.1.7}）得出。由于 (ii) 中对
$Z$ 加了限制，这里必须分别证明 (iii) 和 (iv)
（参见（\egaxref{I.3.5.1-zh}{I, 3.5.1}））。为了证明 (iii)，可归结到
$S=Y$ 的情形（\egaxref{I.3.3.11-zh}{I, 3.3.11}），此时该断言立即由
（\egaref{II.5.5.1-zh}{5.5.1, b)}）和
（\egaref{II.3.5.3-zh}{3.5.3}）得出。为了证明 (iv)，立即可归结到
$X=\mathbf{P}(\mathcal{E})$、$X'=\mathbf{P}(\mathcal{E}')$ 的情形，
其中 $\mathcal{E}$（相应地，$\mathcal{E}'$）是有限型拟凝聚
$\mathcal{O}_Y$-模层（相应地，$\mathcal{O}_{Y'}$-模层）。令 $p$、
$p'$ 分别为 $T=Y\times_SY'$ 到 $Y$ 和 $Y'$ 的典范投影。根据
（\egaref{II.4.1.3.1-zh}{4.1.3.1}），有
$\mathbf{P}(p^*(\mathcal{E}))=\mathbf{P}(\mathcal{E})\times_YT$、
$\mathbf{P}(p'^*(\mathcal{E}'))=\mathbf{P}(\mathcal{E}')\times_{Y'}T$，
从而
\begin{align*}
\mathbf{P}(p^*(\mathcal{E}))\times_T\mathbf{P}(p'^*(\mathcal{E}'))
&=(\mathbf{P}(\mathcal{E})\times_YT)\times_T
(T\times_{Y'}\mathbf{P}(\mathcal{E}'))\\
&=\mathbf{P}(\mathcal{E})\times_Y
(T\times_{Y'}\mathbf{P}(\mathcal{E}'))
=\mathbf{P}(\mathcal{E})\times_S\mathbf{P}(\mathcal{E}')
\end{align*}
这里把 $T$ 换成 $Y\times_SY'$，并使用了
（\egaxref{I.3.3.9.1-zh}{I, 3.3.9.1}）。而
$p^*(\mathcal{E})$ 和 $p'^*(\mathcal{E}')$ 都在 $T$ 上为有限型
（\egaxref{0.5.2.4-zh}{0, 5.2.4}），故
$p^*(\mathcal{E})\otimes_{\mathcal{O}_T}p'^*(\mathcal{E}')$
也为有限型；又因
$\mathbf{P}(p^*(\mathcal{E}))\times_T\mathbf{P}(p'^*(\mathcal{E}'))$
等同于
$\mathbf{P}(p^*(\mathcal{E})\otimes_{\mathcal{O}_T}p'^*(\mathcal{E}'))$
的一个闭子预概形（\egaref{II.4.3.3-zh}{4.3.3}），这就完成了
(iv) 的证明。为了得到 (v) 和 (vi)，可以应用
（\egaxref{I.5.5.13-zh}{I, 5.5.13}），因为由
（\egaref{II.5.5.1-zh}{5.5.1, a)}）可知，射影 $Y$-概形的每个
闭子预概形都是射影 $Y$-概形。

只剩证明 (ii)；鉴于关于 $Z$ 的假设，它由
（\egaref{II.5.5.3-zh}{5.5.3}）、
（\egaref{II.5.3.4-zh}{5.3.4, (ii)}）和
（\egaref{II.5.4.2-zh}{5.4.2, (ii)}）得出。

% Locked French authority: ega2-1-fr.tex, 820505 B,
% SHA-256 84EDBE3E83530AF2959B441796337C9DC21EAFCA6A13114A26778760FBF437AC.
% Sequential source scope: lines 9545-9592, Propositions II.5.5.6-7 and
% Corollary II.5.5.8 / printed pp.105-106.

\begin{proposition}[5.5.6]
\label{II.5.5.6-zh}
若 $X$ 和 $X'$ 是两个射影 $Y$-概形，则 $X\amalg X'$ 是射影
$Y$-概形。
\end{proposition}

这是（\egaref{II.5.5.2-zh}{5.5.2}）和
（\egaref{II.4.3.6-zh}{4.3.6}）的显然推论。

\begin{proposition}[5.5.7]
\label{II.5.5.7-zh}
设 $X$ 是射影 $Y$-概形，$\mathcal{L}$ 是相对于 $Y$ 丰沛的
$\mathcal{O}_X$-模层；对 $\mathcal{L}$ 在 $X$ 上的每个截面 $f$，
$X_f$ 都在 $Y$ 上仿射。
\end{proposition}

\oldpage[II]{106}

由于问题在 $Y$ 上是局部的，可以假设 $Y=\operatorname{Spec}(A)$；
此外，$X_{f^{\otimes n}}=X_f$，故以某个适当的
$\mathcal{L}^{\otimes n}$ 代替 $\mathcal{L}$，可以假设
$\mathcal{L}$ 对结构态射 $q:X\to Y$ 极丰沛
（\egaref{II.4.6.11-zh}{4.6.11}）。于是典范同态
$\sigma:q^*(q_*(\mathcal{L}))\to\mathcal{L}$ 是满射，而相应态射
\[
r=r_{\mathcal{L},\sigma}:X\to P=\mathbf{P}(q_*(\mathcal{L}))
\]
是一个浸入，满足 $\mathcal{L}=r^*(\mathcal{O}_P(1))$
（\egaref{II.4.4.4-zh}{4.4.4}）；此外，由于 $X$ 在 $Y$ 上紧合，
浸入 $r$ 是闭的（\egaref{II.5.4.4-zh}{5.4.4}）。另一方面，
按定义有 $f\in\Gamma(Y,q_*(\mathcal{L}))$，而 $\sigma^\flat$ 是
$q_*(\mathcal{L})$ 的恒等映射；于是由公式
（\egaref{II.3.7.3.1-zh}{3.7.3.1}）可知
$X_f=r^{-1}(D_+(f))$。因此，$X_f$ 是仿射概形 $D_+(f)$ 的闭子预概形，
从而确为仿射概形。

特别取 $Y=X$，便得到下述推论（还须计及
（\egaref{II.4.6.13-zh}{4.6.13, (i)}））；此外，它的直接证明也是
立即的：

\begin{corollary}[5.5.8]
\label{II.5.5.8-zh}
设 $X$ 是预概形，$\mathcal{L}$ 是可逆 $\mathcal{O}_X$-模层。
对 $\mathcal{L}$ 在 $X$ 上的每个截面 $f$，$X_f$ 都在 $X$ 上仿射
（因而当 $X$ 是仿射概形时，它是仿射概形）。
\end{corollary}
% Locked French authority: ega2-1-fr.tex, 820505 B,
% SHA-256 84EDBE3E83530AF2959B441796337C9DC21EAFCA6A13114A26778760FBF437AC.
% Sequential source scope: lines 9594-9638, subsection II.5.6 opening,
% Theorem II.5.6.1 and proof step A / printed pp.106-107.

\subsection{周氏引理。}
\label{subsection:II.5.6-zh}

\begin{theorem}[5.6.1]
\label{II.5.6.1-zh}
（周氏引理。）设 $S$ 是预概形，$X$ 是有限型 $S$-概形。假设下列
条件之一成立：
\begin{enumerate}
  \item[\rm{a)}] $S$ 是诺特的。
  \item[\rm{b)}] $S$ 是拟紧概形，而 $X$ 只有有限多个不可约分支。
\end{enumerate}
在这些条件下：
\begin{enumerate}
  \item[\rm{(i)}] 存在一个在 $S$ 上拟射影的 $S$-概形 $X'$，以及一个
  射影满 $S$-态射 $f:X'\to X$。
  \item[\rm{(ii)}] 可以选取 $X'$ 和 $f$，使存在开集 $U\subset X$，
  对于它，$U'=f^{-1}(U)$ 在 $X'$ 中稠密，并且 $f$ 在 $U'$ 上的限制
  是同构 $U'\simeq U$。
  \item[\rm{(iii)}] 若 $X$ 是约化的（相应地，不可约的，整的），
  则可以设 $X'$ 也是约化的（相应地，不可约的，整的）。
\end{enumerate}
\end{theorem}

证明分为若干步骤。

A) 首先可以归结到 $X$ 不可约的情形。事实上，在假设 a) 下，$X$
是诺特的；所以在两种假设下，$X$ 的不可约分支 $X_i$ 都只有有限多个。
若已经对 $X$ 的以各 $X_i$ 为底空间的约化闭子预概形证明了定理，
而 $X'_i$ 和 $f_i:X'_i\to X_i$ 是对应于 $X_i$ 的预概形和态射，则由
诸 $X'_i$ 之和所得的预概形 $X'$，以及其在每个 $X'_i$ 上的限制为
$j_i\circ f_i$ 的态射 $f:X'\to X$（$j_i$ 是典范嵌入 $X_i\to X$），
满足要求。事实上，若所有 $X'_i$ 都约化，则 $X'$ 显然约化；另一方面，
取 $U$ 为诸集合 $U_i\cap(\bigcup_{j\neq i}X_j)^{\complement}$ 之并，
便满足 (ii)。最后，因为各 $X'_i$ 在 $S$ 上拟射影，$X'$ 也具有此性质

\oldpage[II]{107}

（\egaref{II.5.3.6-zh}{5.3.6}）；同样，各态射 $X'_i\to X$ 由
（\egaref{II.5.5.5-zh}{5.5.5, (i) 和 (ii)}）可知是射影的，故
$f$ 是射影态射（\egaref{II.5.5.6-zh}{5.5.6}），并且按定义显然为满射。
% Locked French authority: ega2-1-fr.tex, 820505 B,
% SHA-256 84EDBE3E83530AF2959B441796337C9DC21EAFCA6A13114A26778760FBF437AC.
% Sequential source scope: lines 9640-9698, proof step B and constructions
% (5.6.1.1)-(5.6.1.4) / printed p.107.

B) 现在设 $X$ 不可约。因为结构态射 $r:X\to S$ 为有限型，所以
存在由仿射开集组成的 $S$ 的有限覆盖 $(S_i)$；而对每个 $i$，存在由
仿射开集组成的 $r^{-1}(S_i)$ 的有限覆盖 $(T_{ij})$。态射
$T_{ij}\to S_i$ 是有限型仿射态射，故为拟射影态射
（\egaref{II.5.3.4-zh}{5.3.4, (i)}）。在假设 a)、b) 的每一种下，
浸入 $S_i\to S$ 都拟紧，因而是拟射影态射
（\egaref{II.5.3.4-zh}{5.3.4, (ii)}）；所以 $r$ 在 $T_{ij}$ 上的
限制是拟射影态射（\egaref{II.5.3.4-zh}{5.3.4, (iii)}）。把诸
$T_{ij}$ 记为 $U_k$，其中 $(1\leq k\leq n)$。对每个指标 $k$，存在开浸入
$\varphi_k:U_k\to P_k$，其中 $P_k$ 在 $S$ 上射影
（\egaref{II.5.3.2-zh}{5.3.2} 和
\egaref{II.5.5.2-zh}{5.5.2}）。令 $U=\bigcap_kU_k$；因为 $X$
不可约且各 $U_k$ 非空，$U$ 非空，因而在 $X$ 中处处稠密。诸
$\varphi_k$ 在 $U$ 上的限制定义了态射
\[
\varphi:U\to P=P_1\times_SP_2\times_S\cdots\times_SP_n
\]
使诸图
\[
\tag{5.6.1.1}
\label{II.5.6.1.1-zh}
\xymatrix{
U\ar[r]^\varphi\ar[d]_{j_k} & P\ar[d]^{p_k}\\
U_k\ar[r]_{\varphi_k} & P_k
}
\]
交换，其中 $j_k$ 是典范嵌入 $U\to U_k$，而 $p_k$ 是典范投影
$P\to P_k$。若 $j$ 是典范嵌入 $U\to X$，则态射
$\psi=(j,\varphi)_S:U\to X\times_SP$ 是浸入
（\egaxref{I.5.3.14-zh}{I, 5.3.14}）。在假设 a) 下，$X\times_SP$
局部诺特（\egaref{II.3.4.1-zh}{3.4.1} 和
\egaxref{I.6.3.7-zh}{I, 6.3.7 和 6.3.8}）；在假设 b) 下，
$X\times_SP$ 是拟紧概形
（\egaxref{I.5.5.1-zh}{I, 5.5.1 和 6.6.4}）。在两种情况下，与
$\psi$ 相联系的子预概形 $Z$（其底空间为 $\psi(U)$）在
$X\times_SP$ 中的闭包 $X'$ 都存在，并且 $\psi$ 分解为
\[
\tag{5.6.1.2}
\label{II.5.6.1.2-zh}
\psi:U\xrightarrow{\psi'}X'\xrightarrow{h}X\times_SP
\]
其中 $\psi'$ 是开浸入，而 $h$ 是闭浸入
（\egaxref{I.9.5.10-zh}{I, 9.5.10}）。设
$q_1:X\times_SP\to X$ 和 $q_2:X\times_SP\to P$ 是典范投影；定义
\[
\tag{5.6.1.3}
\label{II.5.6.1.3-zh}
f:X'\xrightarrow{h}X\times_SP\xrightarrow{q_1}X
\]
以及
\[
\tag{5.6.1.4}
\label{II.5.6.1.4-zh}
g:X'\xrightarrow{h}X\times_SP\xrightarrow{q_2}P.
\]
下面将看到，$X'$ 和 $f$ 满足要求。
% Locked French authority: ega2-1-fr.tex, 820505 B,
% SHA-256 84EDBE3E83530AF2959B441796337C9DC21EAFCA6A13114A26778760FBF437AC.
% Sequential source scope: lines 9700-9725, proof step C
% / printed pp.107-108.

C) 首先证明 $f$ 是射影满射，并且 $f$ 在
$U'=f^{-1}(U)$ 上的限制是从 $U'$ 到 $U$ 的同构。因为各 $P_k$ 在
$S$ 上射影，$P$ 也具有此性质
（\egaref{II.5.5.5-zh}{5.5.5, (iv)}）；故 $X\times_SP$ 在 $X$ 上
射影（\egaref{II.5.5.5-zh}{5.5.5, (iii)}），而作为
$X\times_SP$ 的闭子预概形，$X'$ 也具有此性质。另一方面，
$f\circ\psi'=q_1\circ(h\circ\psi')=q_1\circ\psi=j$，故
$f(X')$ 包含 $X$ 的处处稠密开集 $U$；但 $f$ 是闭态射
（\egaref{II.5.5.3-zh}{5.5.3}），所以 $f(X')=X$。再注意，
$q_1^{-1}(U)=U\times_SP$ 是 $X\times_SP$ 的一个开集上所诱导的；
而按定义，预概形 $U'=h^{-1}(U\times_SP)$ 是 $X'$ 在开集 $U'$
上所诱导的。因此，它是子预概形 $Z$ 相对于 $U\times_SP$ 的闭包

\oldpage[II]{108}

（\egaxref{I.9.5.8-zh}{I, 9.5.8}）。而浸入 $\psi$ 分解为
\[
U\xrightarrow{\Gamma_\varphi}U\times_SP\xrightarrow{j\times1}X\times_SP,
\]
又因 $P$ 在 $S$ 上分离，图态射 $\Gamma_\varphi$ 是闭浸入
（\egaxref{I.5.4.3-zh}{I, 5.4.3}），故 $Z$ 是 $U\times_SP$ 的
闭子预概形，从而 $U'=Z$。又因 $\psi$ 是浸入，$f$ 在 $U'$ 上的
限制是到 $U$ 的同构，其逆为 $\psi'$；最后，按 $X'$ 的定义，
$U'$ 在 $X'$ 中稠密。
% Locked French authority: ega2-1-fr.tex, 820505 B,
% SHA-256 84EDBE3E83530AF2959B441796337C9DC21EAFCA6A13114A26778760FBF437AC.
% Sequential source scope: lines 9727-9755, proof step D opening
% and diagram (5.6.1.5) / printed p.108.

D) 现在证明 $g$ 是浸入；因为 $P$ 在 $S$ 上射影，这将证明 $X'$ 在
$S$ 上拟射影。置
\begin{align*}
V_k&=\varphi_k(U_k) &&\text{（$P_k$ 的开集）},\\
W_k&=p_k^{-1}(V_k) &&\text{（$P$ 的开集）},\\
U'_k&=f^{-1}(U_k),\qquad U''_k=g^{-1}(W_k)
&&\text{（$X'$ 的开集）}.
\end{align*}
显然，诸 $U'_k$ 构成 $X'$ 的一个开覆盖；我们先证明诸 $U''_k$
也如此，方法是证明 $U'_k\subset U''_k$。为此，只需证明图
\[
\tag{5.6.1.5}
\label{II.5.6.1.5-zh}
\xymatrix{
U'_k\ar[r]^{g|U'_k}\ar[d]_{f|U'_k} & P\ar[d]^{p_k}\\
U_k\ar[r]_{\varphi_k} & P_k
}
\]
交换。预概形 $U'_k=h^{-1}(U_k\times_SP)$ 是 $X'$ 在开集 $U'_k$
上所诱导的，因而是 $Z=U'\subset U'_k$ 相对于 $U'_k$ 的闭包
（\egaxref{I.9.5.8-zh}{I, 9.5.8}）。因为 $P_k$ 是 $S$-概形，为证明
图（\egaref{II.5.6.1.5-zh}{5.6.1.5}）交换，只需证明将此图与典范
嵌入 $U'\to U'_k$ 复合后所得的图交换；或者，由于 $U'$ 与 $U$
同构，等价地，只需与 $\psi$ 复合（\egaxref{I.9.5.6-zh}{I, 9.5.6}）。
但按定义，此时得到的正是图（\egaref{II.5.6.1.1-zh}{5.6.1.1}），
故结论成立。
% Locked French authority: ega2-1-fr.tex, 820505 B,
% SHA-256 84EDBE3E83530AF2959B441796337C9DC21EAFCA6A13114A26778760FBF437AC.
% Sequential source scope: lines 9757-9789, proof steps D-E
% / printed pp.108-109.

因此，诸 $W_k$ 构成 $g(X')$ 的一个开覆盖。为证明 $g$ 是浸入，
只需证明每个限制 $g|U''_k$ 都是到 $W_k$ 中的浸入
（\egaxref{I.4.2.4-zh}{I, 4.2.4}）。为此，考虑态射
$u_k:W_k\xrightarrow{p_k}V_k\xrightarrow{\varphi_k^{-1}}U_k\to X$；
因为 $X$ 在 $S$ 上分离，图态射
$\Gamma_{u_k}:W_k\to X\times_SW_k$ 是闭浸入
（\egaxref{I.5.4.3-zh}{I, 5.4.3}），所以图
$T_k=\Gamma_{u_k}(W_k)$ 是 $X\times_SW_k$ 的闭子预概形。若证明
这个子预概形包含 $U'$，则由（\egaxref{I.9.5.8-zh}{I, 9.5.8}），
它也包含 $X'$ 在 $X'$ 的开集 $X''_k$ 上所诱导的子预概形。由于
$q_2$ 在 $T_k$ 上的限制是到 $W_k$ 的同构，$g$ 在 $X''_k$ 上的限制
便是到 $W_k$ 中的浸入，结论即得证。令 $v_k$ 为典范嵌入
$U'\to X\times_SW_k$；我们必须证明存在态射 $w_k:U'\to W_k$，使得
$v_k=\Gamma_{u_k}\circ w_k$。按乘积的定义，只需证明
$q_1\circ v_k=u_k\circ q_2\circ v_k$
（\egaxref{I.3.2.1-zh}{I, 3.2.1}）；或者，在右边与同构

\oldpage[II]{109}

$\psi':U\to U'$ 复合后，只需证明
$q_1\circ\psi=u_k\circ q_2\circ\psi$。但由于 $q_1\circ\psi=j$、
$q_2\circ\psi=\varphi$，结合 $u_k$ 的定义，该断言由图
（\egaref{II.5.6.1.1-zh}{5.6.1.1}）的交换性得出。

E) 显然，在上述构造中，由于 $U$、从而 $U'$ 不可约，$X'$ 也不可约，
因而态射 $f$ 是双有理的（\egaxref{I.2.2.9-zh}{I, 2.2.9}）。此外，
若 $X$ 约化，则 $U'$ 也约化，故 $X'$ 约化
（\egaxref{I.9.5.9-zh}{I, 9.5.9}）。证明完毕。
% Locked French authority: ega2-1-fr.tex, 820505 B,
% SHA-256 84EDBE3E83530AF2959B441796337C9DC21EAFCA6A13114A26778760FBF437AC.
% Sequential source scope: lines 9791-9813, Corollary II.5.6.2
% / printed p.109.

\begin{corollary}[5.6.2]
\label{II.5.6.2-zh}
假设（\egaref{II.5.6.1-zh}{5.6.1}）的假设 a)、b) 之一成立。
$X$ 在 $S$ 上紧合的充要条件是：存在一个在 $S$ 上射影的概形 $X'$，
以及一个满 $S$-态射 $f:X'\to X$（此时由
（\egaref{II.5.5.5-zh}{5.5.5, (v)}）可知它是射影的）。在这种情形下，
还可以选取 $f$，使得 $X$ 中存在稠密开集 $U$，而 $f$ 在
$f^{-1}(U)$ 上的限制是同构 $f^{-1}(U)\simeq U$，并且 $f^{-1}(U)$
在 $X'$ 中稠密。此外，若 $X$ 不可约（相应地，约化），则可以假设
$X'$ 也如此；当 $X$ 和 $X'$ 都不可约时，$f$ 是双有理态射。
\end{corollary}

由（\egaref{II.5.5.3-zh}{5.5.3}）和
（\egaref{II.5.4.3-zh}{5.4.3, (ii)}），该条件是充分的。它也是必要的：
沿用（\egaref{II.5.6.1-zh}{5.6.1}）的记号，若 $X$ 在 $S$ 上紧合，
$X'$ 也在 $S$ 上紧合；因为它在 $X$ 上射影，故在 $X$ 上紧合
（\egaref{II.5.5.3-zh}{5.5.3}），结论由
（\egaref{II.5.4.2-zh}{5.4.2, (ii)}）得出。此外，由于 $X'$ 在
$S$ 上拟射影，由（\egaref{II.5.5.3-zh}{5.5.3}）可知它在 $S$ 上射影。
% Locked French authority: ega2-1-fr.tex, 820505 B,
% SHA-256 84EDBE3E83530AF2959B441796337C9DC21EAFCA6A13114A26778760FBF437AC.
% Sequential source scope: lines 9815-9846, Corollary II.5.6.3
% and first proof part / printed p.109.

\begin{corollary}[5.6.3]
\label{II.5.6.3-zh}
设 $S$ 是局部诺特预概形，$X$ 是在 $S$ 上有限型的 $S$-概形，
$f_0:X\to S$ 是其结构态射。$X$ 在 $S$ 上紧合的充要条件是：
对每个有限型态射 $S'\to S$，
$(f_0)_{(S')}:X_{(S')}\to S'$ 都是闭态射。甚至只需对每个形如
$S'=S\otimes_{\mathbf Z}\mathbf Z[T_1,\ldots,T_n]$（$T_i$ 为不定元）的
$S$-预概形验证此条件。
\end{corollary}

此条件显然是必要的；下面证明它是充分的。由于问题在 $S$ 和 $S'$
上都是局部的（\egaref{II.5.4.1-zh}{5.4.1}），可以假设 $S$ 和 $S'$
都是仿射诺特的。由周氏引理，存在射影 $S$-概形 $P$、浸入
$j:X'\to P$ 和射影满态射 $f:X'\to X$，使图
\[
\xymatrix{
X\ar[d]_{f_0} & X'\ar[l]_f\ar[d]^j\\
S & P\ar[l]_r
}
\]
交换。因为 $P$ 在 $S$ 上有限型，第一个假设蕴含投影
$q_2:X\times_SP\to P$ 是闭态射。但浸入 $j$ 是 $q_2$ 与从
$X'\times_SP$ 到 $X\times_SP$ 的态射 $f\times1$ 的复合；而 $f$
是射影的，故是紧合的（\egaref{II.5.5.3-zh}{5.5.3}），所以
$f\times1$ 是闭态射。由此可知 $j$ 是闭浸入，因而是紧合的
（\egaref{II.5.4.2-zh}{5.4.2, (ii)}）。此外，结构态射 $r:P\to S$
是射影的，故是紧合的（\egaref{II.5.5.3-zh}{5.5.3}），所以
$f_0\circ f=r\circ j$ 是紧合的
（\egaref{II.5.4.2-zh}{5.4.2, (ii)}）；最后，由于 $f$ 是满射，
由（\egaref{II.5.4.3-zh}{5.4.3}）可知 $f_0$ 紧合。
% Locked French authority: ega2-1-fr.tex, 820505 B,
% SHA-256 84EDBE3E83530AF2959B441796337C9DC21EAFCA6A13114A26778760FBF437AC.
% Sequential source scope: lines 9848-9867, Corollary II.5.6.3
% final proof part / printed pp.109-110.

为只用第二个假设证明该命题，只需证明它蕴含第一个假设。事实上，
若 $S'$ 是在 $S=\operatorname{Spec}(A)$ 上有限型的仿射概形，

\oldpage[II]{110}

则 $S'=\operatorname{Spec}(A[c_1,\ldots,c_n])$
（\egaxref{I.6.3.3-zh}{I, 6.3.3}），所以 $S'$ 同构于
$S''=\operatorname{Spec}(A[T_1,\ldots,T_n])$（$T_i$ 为不定元）的一个
闭子预概形。在交换图
\[
\xymatrix{
X\times_SS'\ar[r]^{1_X\times j}\ar[d]_{(f_0)_{(S')}} &
X\times_SS''\ar[d]^{(f_0)_{(S'')}}\\
S'\ar[r]_j & S''
}
\]
中，$j$ 和 $1_X\times j$ 都是闭浸入
（\egaxref{I.4.3.1-zh}{I, 4.3.1}），而 $(f_0)_{(S'')}$ 依假设是闭态射；
因此 $(f_0)_{(S')}$ 也是闭态射。
% Locked French authority: ega2-1-fr.tex, 820505 B,
% SHA-256 84EDBE3E83530AF2959B441796337C9DC21EAFCA6A13114A26778760FBF437AC.
% Sequential source scope: lines 9869-9910, Section II.6 opening and
% Definition II.6.1.1 through environment II.6.1.2 / printed pp.110-111.

\section{整态射与有限态射}
\label{section:II.6-zh}

\subsection{在另一预概形上整的预概形。}
\label{subsection:II.6.1-zh}

\begin{definition}[6.1.1]
\label{II.6.1.1-zh}
设 $X$ 是 $S$-预概形，$f:X\to S$ 是其结构态射。若存在由仿射开集
组成的 $S$ 的覆盖 $(S_\alpha)$，使得对每个 $\alpha$，所诱导的预概形
$f^{-1}(S_\alpha)$ 是仿射概形，且其环 $B_\alpha$ 是
$S_\alpha$ 的环 $A_\alpha$ 上的整代数，则称 $X$ \emph{在 $S$ 上整}，
或称 $f$ 为\emph{整态射}。若 $X$ 在 $S$ 上整且在 $S$ 上有限型，
则称 $X$ \emph{在 $S$ 上有限}，或称 $f$ 为\emph{有限态射}。
\end{definition}

当 $S$ 是环 $A$ 的仿射概形时，也称“在 $A$ 上整（相应地，有限）”，
而不称“在 $S$ 上整（相应地，有限）”。

\begin{env}[6.1.2]
\label{II.6.1.2-zh}
显然，若 $X$ 在 $S$ 上整，则它在 $S$ 上仿射。对于在 $S$ 上仿射的
预概形 $X$，它在 $S$ 上整（相应地，有限）的充要条件是：与之相联的
拟凝聚 $\mathcal{O}_S$-代数 $\mathcal{A}(X)$ 具有如下性质：存在由仿射
开集组成的 $S$ 的覆盖 $(S_\alpha)$，使得对每个 $\alpha$，
$\Gamma(S_\alpha,\mathcal{A}(X))$ 是
$\Gamma(S_\alpha,\mathcal{O}_S)$ 上的整代数（相应地，整且有限型的
代数）。具有此性质的拟凝聚 $\mathcal{O}_S$-代数称为在
$\mathcal{O}_S$ 上\emph{整}（相应地，\emph{有限}）。因此，给定一个在
$S$ 上整（相应地，有限）的预概形，等价于
（\egaref{II.1.3.1-zh}{1.3.1}）给定一个在 $\mathcal{O}_S$ 上整
（相应地，有限）的拟凝聚 $\mathcal{O}_S$-代数。注意，拟凝聚
$\mathcal{O}_S$-代数 $\mathcal{B}$ 有限，当且仅当它是有限型
$\mathcal{O}_S$-模（\egaxref{I.1.3.9-zh}{I, 1.3.9}）；这也等价于说
$\mathcal{B}$ 是整且有限型的 $\mathcal{O}_S$-代数，因为环 $A$ 上
整且有限型的代数是有限型 $A$-模。
\end{env}
% Locked French authority: ega2-1-fr.tex, 820505 B,
% SHA-256 84EDBE3E83530AF2959B441796337C9DC21EAFCA6A13114A26778760FBF437AC.
% Sequential source scope: lines 9912-9955, Proposition II.6.1.3 through
% Lemma II.6.1.4.1 and proof / printed pp.110-111.

\begin{proposition}[6.1.3]
\label{II.6.1.3-zh}
设 $S$ 是局部诺特预概形。对于在 $S$ 上仿射的预概形 $X$，它在
$S$ 上有限的充要条件是 $\mathcal{O}_S$-代数 $\mathcal{A}(X)$ 是凝聚的。
\end{proposition}

结合前一评注，这归结为注意到：若 $S$ 局部诺特，则有限型拟凝聚
$\mathcal{O}_S$-模恰好就是凝聚 $\mathcal{O}_S$-模
（\egaxref{I.1.5.1-zh}{I, 1.5.1}）。

\begin{proposition}[6.1.4]
\label{II.6.1.4-zh}
设 $X$ 是在 $S$ 上整（相应地，有限）的预概形，$f:X\to S$ 是结构态射。
则对每个环为 $A$ 的仿射开集 $U\subset S$，$f^{-1}(U)$ 都是仿射概形，
而其环 $B$ 是 $A$ 上的整代数（相应地，有限代数）。
\end{proposition}

\oldpage[II]{111}

先证明下列引理：

\begin{lemma}[6.1.4.1]
\label{II.6.1.4.1-zh}
设 $A$ 是环，$M$ 是 $A$-模，$(g_i)_{1\leq i\leq m}$ 是 $A$ 中的有限
元素系，且诸 $D(g_i)$（$1\leq i\leq m$）覆盖
$\operatorname{Spec}(A)$。若对每个 $i$，$M_{g_i}$ 是有限型
$A_{g_i}$-模，则 $M$ 是有限型 $A$-模。
\end{lemma}

事实上，可以假设 $M_{g_i}$ 有有限生成元系 $(m_{ij}/g_i^n)$，其中
$m_{ij}\in M$，并且 $n$ 对所有指标 $i$ 都相同。证明诸 $m_{ij}$ 生成
$M$。令 $M'$ 为它们生成的 $M$ 的 $A$-子模，并取 $m$ 为 $M$ 的元素。
依假设，对每个 $i$，存在 $a_{ij}\in A$ 和与 $i$ 无关的整数 $p$，使在
$M_{g_i}$ 中有 $m/1=(\sum_j a_{ij}m_{ij})/g_i^p$；这蕴含存在整数
$r\geq p$，使对每个 $i$ 都有 $g_i^r m\in M'$。由于
$D(g_i^r)=D(g_i)$ 覆盖 $\operatorname{Spec}(A)$，诸 $g_i^r$ 生成的
$A$ 的理想等于 $A$；换言之，存在 $a_i\in A$，使
$\sum_i a_i g_i^r=1$。因此 $m=(\sum_i a_i g_i^r)m\in M'$，引理得证。
% Locked French authority: ega2-1-fr.tex, 820505 B,
% SHA-256 84EDBE3E83530AF2959B441796337C9DC21EAFCA6A13114A26778760FBF437AC.
% Sequential source scope: lines 9957-9988, proof of Proposition II.6.1.4
% / printed p.111.

现在，已知（\egaref{II.1.3.2-zh}{1.3.2}）$f^{-1}(U)$ 是仿射的。
若 $\varphi$ 是与 $f$ 相应的同态 $A\to B$，则存在由开集 $D(g_i)$
（$g_i\in A$）组成的 $U$ 的有限覆盖，使得若 $h_i=\varphi(g_i)$，
则 $B_{h_i}$ 是 $A_{g_i}$ 上的整代数（相应地，整有限代数）。事实上，
存在由仿射开集 $V_\alpha\subset U$ 组成的 $U$ 的覆盖，使得若
$A_\alpha=A(V_\alpha)$，则 $B_\alpha=A(f^{-1}(V_\alpha))$ 是
$A_\alpha$ 上的整代数（相应地，有限代数）。每个 $x\in U$ 属于某个
$V_\alpha$，故存在 $g\in A$，使 $x\in D(g)\subset V_\alpha$；若
$g_\alpha$ 是 $g$ 在 $A_\alpha$ 中的像，则
$A(D(g))=A_g=(A_\alpha)_{g_\alpha}$。令 $h=\varphi(g)$，并令
$h_\alpha$ 为 $g_\alpha$ 在 $B_\alpha$ 中的像；则
\[
A(D(h))=B_h=(B_\alpha)_{h_\alpha}
\]
而由于 $B_\alpha$ 在 $A_\alpha$ 上整（相应地，有限），
$(B_\alpha)_{h_\alpha}$ 在 $(A_\alpha)_{g_\alpha}$ 上整
（相应地，有限）。现在只需利用 $U$ 拟紧，便得到所需的覆盖。

先假设诸 $B_{h_i}$ 在 $A_{g_i}$ 上整且有限。作为 $A_{g_i}$-模，
$B_{h_i}$ 也写成 $B_{g_i}$，故引理
（\egaref{II.6.1.4.1-zh}{6.1.4.1}）表明，此时 $B$ 是有限型 $A$-模。

现在只假设每个 $B_{h_i}$ 在 $A_{g_i}$ 上整。取 $b\in B$，并令 $C$
为 $b$ 生成的 $B$ 的 $A$-子代数。对每个 $i$，$C_{h_i}$ 是
$B_{h_i}$ 中由 $b/1$ 生成的 $A_{g_i}$-代数；由假设，每个 $C_{h_i}$
都是有限型 $A_{g_i}$-模，故由
（\egaref{II.6.1.4.1-zh}{6.1.4.1}）知 $C$ 是有限型 $A$-模；这证明
$B$ 在 $A$ 上整。
% Locked French authority: ega2-1-fr.tex, 820505 B,
% SHA-256 84EDBE3E83530AF2959B441796337C9DC21EAFCA6A13114A26778760FBF437AC.
% Sequential source scope: lines 9990-10033, Proposition II.6.1.5 and proof
% / printed pp.111-112.

\begin{proposition}[6.1.5]
\label{II.6.1.5-zh}
\begin{enumerate}
\item[(i)] 闭浸入是有限的（因而更是整的）。
\item[(ii)] 两个有限态射（相应地，整态射）的复合是有限的
（相应地，整的）。
\item[(iii)] 若 $f:X\to Y$ 是有限（相应地，整）的 $S$-态射，则对基的
每个扩张 $S'\to S$，$f_{(S')}:X_{(S')}\to Y_{(S')}$ 都是有限的
（相应地，整的）。
\item[(iv)] 若 $f:X\to Y$ 和 $g:X'\to Y'$ 是两个有限（相应地，整）的
$S$-态射，则
$f\times_S g:X\times_S X'\to Y\times_S Y'$ 也如此。
\item[(v)] 若 $f:X\to Y$ 和 $g:Y\to Z$ 是两个态射，使得 $g\circ f$
有限（相应地，整），且 $g$ 分离，则 $f$ 有限（相应地，整）。
\item[(vi)] 若 $f:X\to Y$ 是有限（相应地，整）的态射，则
$f_{\mathrm{red}}$ 也如此。
\end{enumerate}
\end{proposition}

\oldpage[II]{112}

由（\egaxref{I.5.5.12-zh}{I, 5.5.12}），只需证明 (i)、(ii) 和 (iii)。
为证明闭浸入 $X\to S$ 是有限的，可以归结到
$S=\operatorname{Spec}(A)$ 的情形；此时只需注意，商环
$A/\mathfrak{J}$ 是单生成 $A$-模。为证明两个有限（相应地，整）的态射
$X\to Y$、$Y\to Z$ 的复合是有限的（相应地，整的），仍可假设 $Z$
（从而 $X$、$Y$（\egaref{II.1.3.4-zh}{1.3.4}））是仿射的；这时断言
等价于说：若 $B$ 是有限（相应地，整）的 $A$-代数，而 $C$ 是有限
（相应地，整）的 $B$-代数，则 $C$ 是有限（相应地，整）的 $A$-代数；
这是显然的。最后，为证明 (iii)，可以归结到 $S=Y$ 的情形，因为
$X_{(S')}$ 与 $X\times_Y Y_{(S')}$ 同一
（\egaxref{I.3.3.11-zh}{I, 3.3.11}）；还可假设
$S=\operatorname{Spec}(A)$、$S'=\operatorname{Spec}(A')$。于是 $X$
是环为 $B$ 的仿射概形（\egaref{II.1.3.4-zh}{1.3.4}），而
$X_{(S')}$ 是环为 $A'\otimes_A B$ 的仿射概形；只需注意，若 $B$ 是
有限（相应地，整）的 $A$-代数，则 $A'\otimes_A B$ 是有限
（相应地，整）的 $A'$-代数。

还要注意，若 $X$ 和 $Y$ 是两个在 $S$ 上有限（相应地，整）的
$S$-预概形，则它们的\emph{和} $X\amalg Y$ 是在 $S$ 上有限
（相应地，整）的预概形；因为这归结为看到：若 $B$ 和 $C$ 是两个
$A$-代数，且在 $A$ 上有限（相应地，整），则 $B\times C$ 也如此。
% Locked French authority: ega2-1-fr.tex, 820505 B,
% SHA-256 84EDBE3E83530AF2959B441796337C9DC21EAFCA6A13114A26778760FBF437AC.
% Sequential source scope: lines 10035-10056, Corollaries II.6.1.6-II.6.1.7
% / printed p.112.

\begin{corollary}[6.1.6]
\label{II.6.1.6-zh}
若 $X$ 是在 $S$ 上整（相应地，有限）的预概形，则对每个开集
$U\subset S$，$f^{-1}(U)$ 都在 $U$ 上整（相应地，有限）。
\end{corollary}

这是（\egaref{II.6.1.5-zh}{6.1.5, (iii)}）的特例。

\begin{corollary}[6.1.7]
\label{II.6.1.7-zh}
设 $f:X\to Y$ 是有限态射。则对每个 $y\in Y$，纤维 $f^{-1}(y)$ 是
$k(y)$ 上的有限代数概形，因而其底空间是离散且有限的。
\end{corollary}

事实上，作为 $k(y)$-预概形，$f^{-1}(y)$ 与
$X\times_Y\operatorname{Spec}(k(y))$ 同一
（\egaxref{I.3.6.1-zh}{I, 3.6.1}），后者在
$\operatorname{Spec}(k(y))$ 上有限
（\egaref{II.6.1.5-zh}{6.1.5, (iii)}）；因此它是仿射概形，其环是
$k(y)$ 上有限秩的代数（\egaref{II.6.1.4-zh}{6.1.4}）。命题于是由
（\egaxref{I.6.4.4-zh}{I, 6.4.4}）得出。
% Locked French authority: ega2-1-fr.tex, 820505 B,
% SHA-256 84EDBE3E83530AF2959B441796337C9DC21EAFCA6A13114A26778760FBF437AC.
% Sequential source scope: lines 10058-10115, Corollary II.6.1.8 through
% Proposition II.6.1.10 and proof / printed pp.112-113.

\begin{corollary}[6.1.8]
\label{II.6.1.8-zh}
设 $X$、$S$ 是两个整预概形，$f:X\to S$ 是优势态射。若 $f$ 是整态射
（相应地，有限态射），则 $X$ 上的有理函数域 $R(X)$ 是 $S$ 上的
有理函数域 $R(S)$ 的代数扩张（相应地，有限次代数扩张）。
\end{corollary}

事实上，设 $s$ 是 $S$ 的泛点；$k(s)$-预概形 $f^{-1}(s)$ 在
$\operatorname{Spec}(k(s))$ 上整（相应地，有限）
（\egaref{II.6.1.5-zh}{6.1.5, (iii)}），并且按假设包含 $X$ 的泛点
$x$；$x$ 在 $f^{-1}(s)$ 中的局部环等于 $k(x)$
（\egaxref{I.3.6.5-zh}{I, 3.6.5}），而它是一个在 $k(s)$ 上整
（相应地，有限）的代数的局部环（\egaref{II.6.1.4-zh}{6.1.4}），
故得此推论。

\begin{remark}[6.1.9]
\label{II.6.1.9-zh}
假设 $g$ 是\emph{分离的}，对于
\egaref{II.6.1.5-zh}{6.1.5, (v)} 的成立至关重要：事实上，如果
$Y$ 在 $Z$ 上不分离，则恒等态射 $1_Y$ 是复合态射
$Y\xrightarrow{\Delta_Y}Y\times_ZY\xrightarrow{p_1}Y$，但 $\Delta_Y$
不是整态射，这由 \egaref{II.6.1.10-zh}{6.1.10} 可知：
\end{remark}

\begin{proposition}[6.1.10]
\label{II.6.1.10-zh}
每个整态射都是普遍闭的。
\end{proposition}

设 $f:X\to Y$ 是整态射；根据 \egaref{II.6.1.5-zh}{6.1.5, (iii)}，
只须证明 $f$ 是闭的。设 $Z$ 是 $X$ 的一个闭子集；存在 $X$ 的一个

\oldpage[II]{113}

以 $Z$ 为底空间的子预概形（\egaxref{I.5.2.1-zh}{I, 5.2.1}），
因此由 \egaref{II.6.1.5-zh}{6.1.5, (i) 和 (ii)} 可知，可归结为证明
$f(X)$ 在 $Y$ 中是闭的。根据
\egaref{II.6.1.5-zh}{6.1.5, (vii)}，可假设 $X$ 和 $Y$ 都约化；
此外，若 $T$ 是 $Y$ 的约化闭子预概形，其底空间为
$\overline{f(X)}$（\egaxref{I.5.2.1-zh}{I, 5.2.1}），则已知 $f$
可分解为 $X\xrightarrow{g}T\xrightarrow{j}Y$，其中 $j$ 是单射态射
（\egaxref{I.5.2.2-zh}{I, 5.2.2}）；由于 $j$ 是分离的
（\egaxref{I.5.5.1-zh}{I, 5.5.1, (i)}），由
\egaref{II.6.1.5-zh}{6.1.5, (v)} 可知 $g$ 是整态射。因此可假设
$f(X)$ 在 $Y$ 中稠密。最后，由于问题在 $Y$ 上是局部的，可归结到
$Y=\operatorname{Spec}(A)$ 的情形。此时 $X=\operatorname{Spec}(B)$，
其中 $B$ 是在 $A$ 上整的 $A$-代数
（\egaref{II.6.1.4-zh}{6.1.4}）；此外 $A$ 是约化的
（\egaxref{I.5.1.4-zh}{I, 5.1.4}），而 $f(X)$ 在 $Y$ 中稠密这一假设
蕴含与 $f$ 对应的同态 $\varphi:A\to B$ 是单射
（\egaxref{I.1.2.7-zh}{I, 1.2.7}）。在这些条件下，说 $f(X)=Y$，
就是指 $A$ 的每个素理想都是 $B$ 中某个素理想在 $A$ 中的收缩；
这正是 Cohen--Seidenberg 第一定理
（\egaxcite[t.~I, p.~257, th.~3]{I-13}）。
% Locked French authority: ega2-1-fr.tex, 820505 B,
% SHA-256 84EDBE3E83530AF2959B441796337C9DC21EAFCA6A13114A26778760FBF437AC.
% Sequential source scope: lines 10117-10153, Corollary II.6.1.11 and
% Proposition II.6.1.12 through the lead-in to Lemma II.6.1.12.1 / p.113.

\begin{corollary}[6.1.11]
\label{II.6.1.11-zh}
每个有限态射 $f:X\to Y$ 都是射影态射。
\end{corollary}

事实上，由于 $f$ 是仿射的，$\mathcal{O}_X$ 是相对于 $f$ 很丰沛的
$\mathcal{O}_X$-模（\egaref{II.5.1.2-zh}{5.1.2}）；此外
$f_*(\mathcal{O}_X)$ 是有限型拟凝聚 $\mathcal{O}_Y$-模
（\egaref{II.6.1.2-zh}{6.1.2}）；最后，$f$ 是分离、有限型且普遍闭的
（\egaref{II.6.1.10-zh}{6.1.10}），因而满足判据
\egaref{II.5.5.4-zh}{5.5.4, (ii)} 的适用条件。

\begin{proposition}[6.1.12]
\label{II.6.1.12-zh}
设 $f:X'\to X$ 是有限态射，并令
$\mathcal{B}=f_*(\mathcal{O}_{X'})$（它是拟凝聚 $\mathcal{O}_X$-代数，
并且是有限型 $\mathcal{O}_X$-模）。设 $\mathcal{F}'$ 是拟凝聚
$\mathcal{O}_{X'}$-模；则 $\mathcal{F}'$ 是秩 $r$ 的局部自由模，
当且仅当 $f_*(\mathcal{F}')$ 是秩 $r$ 的局部自由 $\mathcal{B}$-模。
\end{proposition}

显然，如果 $f_*(\mathcal{F}')|U$ 同构于 $\mathcal{B}^r|U$
（其中 $U$ 是 $X$ 的开集），则 $\mathcal{F}'|f^{-1}(U)$ 同构于
$\mathcal{O}_{X'}^r|f^{-1}(U)$（\egaref{II.1.4.2-zh}{1.4.2}）。
反过来，假设 $\mathcal{F}'$ 是秩 $r$ 的局部自由模，并证明
$f_*(\mathcal{F}')$ 作为 $\mathcal{B}$-模局部同构于
$\mathcal{B}^r$。设 $x$ 是 $X$ 的一个点；当 $U$ 遍历 $x$ 的一个
仿射邻域基本系时，$f^{-1}(U)$ 遍历有限集合 $f^{-1}(x)$ 的一个
仿射邻域基本系（\egaref{II.1.2.5-zh}{1.2.5}），这是因为 $f$ 是闭的
（\egaref{II.6.1.10-zh}{6.1.10}）。于是该命题由下述引理推出：
% Locked French authority: ega2-1-fr.tex, 820505 B,
% SHA-256 84EDBE3E83530AF2959B441796337C9DC21EAFCA6A13114A26778760FBF437AC.
% Sequential source scope: lines 10155-10187, Lemma II.6.1.12.1 and proof
% / printed pp.113-114.

\begin{lemma}[6.1.12.1]
\label{II.6.1.12.1-zh}
设 $Y$ 是预概形，$\mathcal{E}$ 是秩 $r$ 的局部自由
$\mathcal{O}_Y$-模，$Z$ 是 $Y$ 的一个有限子集，且包含在仿射开集
$V$ 中。则存在 $Z$ 的一个邻域 $U\subset V$，使得
$\mathcal{E}|U$ 同构于 $\mathcal{O}_Y^r|U$。
\end{lemma}

显然可假设 $Y$ 是仿射的；对于每个 $z_i\in Z$，闭包
$\overline{\{z_i\}}$ 中至少存在一个闭点 $z'_i$
（\egaxref{0.2.1.3-zh}{0, 2.1.3}）；若 $Z'$ 是这些 $z'_i$ 构成的集合，
则 $Z'$ 的每个邻域也是 $Z$ 的邻域，因而可假设 $Z$ 在 $Y$ 中闭且离散。
考虑 $Y$ 的以 $Z$ 为底空间的约化闭子预概形
（\egaxref{I.5.2.1-zh}{I, 5.2.1}），并令 $j:Z\to Y$ 为典范单射；
$j^*(\mathcal{E})=\mathcal{E}\otimes_Y\mathcal{O}_Z$ 在离散概形 $Z$
上是秩 $r$ 的局部自由模，因而同构于 $\mathcal{O}_Z^r$；换言之，
存在 $\mathcal{E}\otimes_Y\mathcal{O}_Z$ 在 $Z$ 上的 $r$ 个截面
$s_i$（$1\leq i\leq r$），使得由这些截面定义的同态
$\mathcal{O}_Z^r\to\mathcal{E}\otimes_Y\mathcal{O}_Z$ 是双射。
但是 $Y=\operatorname{Spec}(A)$ 是仿射的，$Z$ 由 $A$ 的理想
$\mathfrak{J}$ 定义，并且 $\mathcal{E}=\widetilde{M}$，其中 $M$ 是
$A$-模；这些 $s_i$ 是

\oldpage[II]{114}

$M\otimes_A(A/\mathfrak{J})$ 的元素，因而是 $r$ 个元素
$t_i\in M=\Gamma(Y,\mathcal{E})$ 的像。对于每个 $z_j\in Z$，于是存在
$z_j$ 的一个邻域 $V_j$，使得 $t_i$ 在 $V_j$ 上的限制定义一个同构
$\mathcal{O}_Y^r|V_j\to\mathcal{E}|V_j$
（\egaxref{0.5.5.4-zh}{0, 5.5.4}）；因此，由这些 $V_j$ 的并所构成的
邻域 $U$ 即为所求。
% Locked French authority: ega2-1-fr.tex, 820505 B,
% SHA-256 84EDBE3E83530AF2959B441796337C9DC21EAFCA6A13114A26778760FBF437AC.
% Sequential source scope: lines 10189-10220, Proposition II.6.1.13 and
% reduction to Corollary II.6.1.14 / printed p.114.

\begin{proposition}[6.1.13]
\label{II.6.1.13-zh}
设 $g:X'\to X$ 是预概形的整态射，$Y$ 是正规且局部整的预概形，
$f$ 是从 $Y$ 到 $X'$ 的有理映射，并且 $g\circ f$ 是处处有定义的
有理映射（\egaxref{I.7.2.1-zh}{I, 7.2.1}）；则 $f$ 处处有定义。
\end{proposition}

如果 $f_1$ 和 $f_2$ 是类 $f$ 中的两个态射（从 $Y$ 的稠密开集到
$X'$），则显然 $g\circ f_1$ 和 $g\circ f_2$ 是等价态射，这说明可用
记号 $g\circ f$ 表示它们的等价类。还要回顾，当进一步假设 $Y$
局部诺特时，$Y$ 正规这一假设已经蕴含 $Y$ 局部整
（\egaxref{I.6.1.13-zh}{I, 6.1.13}）。

为证明 \egaref{II.6.1.13-zh}{6.1.13}，首先注意到该问题在 $Y$ 上
是局部的，因而可假设在 $g\circ f$ 的类中存在态射 $h:Y\to X$。
考虑逆像 $Y'=X'_{(h)}=X'_{(Y)}$，并注意到态射
$g'=g_{(Y)}:Y'\to Y$ 是整态射
（\egaref{II.6.1.5-zh}{6.1.5, (iii)}）。利用从 $Y$ 到 $X'$ 的
有理映射与 $Y'$ 的 $Y$-有理截面之间的对应
（\egaxref{I.7.1.2-zh}{I, 7.1.2}），可知问题归结为证明
\egaref{II.6.1.13-zh}{6.1.13} 中 $X=Y$ 的特殊情形，即得到下面的

\begin{corollary}[6.1.14]
\label{II.6.1.14-zh}
设 $X$ 是正规且局部整的预概形，$g:X'\to X$ 是整态射，$f$ 是
$X'$ 的一个 $X$-有理截面。则 $f$ 处处有定义。
\end{corollary}
% Locked French authority: ega2-1-fr.tex, 820505 B,
% SHA-256 84EDBE3E83530AF2959B441796337C9DC21EAFCA6A13114A26778760FBF437AC.
% Sequential source scope: lines 10222-10255, proof of Corollary II.6.1.14,
% Corollary II.6.1.15 and conclusion of Proposition II.6.1.13 / p.114.

由于问题在 $X$ 上是局部的，可假设 $X$ 是整的；于是 $f$ 可等同于
从 $X$ 的一个开集 $U$ 到 $X'$ 的态射
（\egaxref{I.7.2.2-zh}{I, 7.2.2}），该态射是 $g^{-1}(U)$ 的
$U$-截面。由于 $g$ 是分离的，$f$ 是从 $U$ 到 $g^{-1}(U)$ 的闭浸入
（\egaxref{I.5.4.6-zh}{I, 5.4.6}）；设 $Z$ 为 $g^{-1}(U)$ 的与 $f$
对应的闭子预概形（\egaxref{I.4.2.1-zh}{I, 4.2.1}），它同构于 $U$，
因而是整的；设 $X_1$ 为 $X'$ 的约化子预概形，其底空间是 $Z$ 在
$X'$ 中的闭包 $\overline{Z}$
（\egaxref{I.5.2.1-zh}{I, 5.2.1}）；于是 $Z$ 是 $X_1$ 某个开集上的
诱导子预概形（\egaxref{I.5.2.3-zh}{I, 5.2.3}），并且由于它不可约，
$X_1$ 也不可约，后者因而是整的。于是 $f$ 可看作 $X_1$ 的
$X$-有理截面；由于把 $g$ 限制在 $X_1$ 上所得态射是整态射
（\egaref{II.6.1.5-zh}{6.1.5, (i) 和 (ii)}），最终归结为证明
\egaref{II.6.1.14-zh}{6.1.14} 中 $X'=X_1$ 的特殊情形，换言之，
归结为下面的

\begin{corollary}[6.1.15]
\label{II.6.1.15-zh}
设 $X$ 是整正规预概形，$X'$ 是整预概形，$g:X'\to X$ 是整态射。
如果存在 $X'$ 的一个 $X$-有理截面 $f$，则 $g$ 是同构。
\end{corollary}

由于问题在 $X$ 上是局部的，可假设 $X$ 是仿射的，其环为整环 $A$；
于是 $X'$ 是仿射的，其环 $A'$ 在 $A$ 上整
（\egaref{II.6.1.4-zh}{6.1.4}），并且是整环；此外，
\egaref{II.6.1.14-zh}{6.1.14} 中的论证表明，$X$ 中存在一个稠密开集，
同构于 $X'$ 中的一个稠密开集，因而 $A$ 和 $A'$ 具有相同的分式域。
另一方面，根据 \egaxref{I.8.2.1.1-zh}{I, 8.2.1.1} 以及各
$\mathcal{O}_x$ 都整闭这一假设，环 $A$ 是整闭的，故 $A'=A$；
这就完成了 \egaref{II.6.1.13-zh}{6.1.13} 的证明。
% Locked French authority: ega2-1-fr.tex, 820505 B,
% SHA-256 84EDBE3E83530AF2959B441796337C9DC21EAFCA6A13114A26778760FBF437AC.
% Sequential source scope: lines 10257-10289, subsection II.6.2 opening and
% Proposition II.6.2.1 with proof / printed pp.114-115.

\subsection{拟有限态射。}
\label{subsection:II.6.2-zh}

\begin{proposition}[6.2.1]
\label{II.6.2.1-zh}
设 $f:X\to Y$ 是局部有限型态射，$x$ 是 $X$ 的一个点。下列条件等价：

\oldpage[II]{115}

\begin{enumerate}
\item[a)] 点 $x$ 在其纤维 $f^{-1}(f(x))$ 中是孤立的。
\item[b)] 环 $\mathcal{O}_x$ 是拟有限 $\mathcal{O}_{f(x)}$-模
（\egaxref{0.7.4.1-zh}{0, 7.4.1}）。
\end{enumerate}
\end{proposition}

该问题显然在 $X$ 和 $Y$ 上都是局部的，故可假设
$X=\operatorname{Spec}(A)$ 和 $Y=\operatorname{Spec}(B)$ 都仿射，
其中 $A$ 是有限型 $B$-代数（\egaxref{I.6.3.3-zh}{I, 6.3.3}）。
此外，可以用 $X\times_Y\operatorname{Spec}(\mathcal{O}_{f(x)})$
代替 $X$，而不改变纤维 $f^{-1}(f(x))$ 或局部环 $\mathcal{O}_x$
（\egaxref{I.3.6.5-zh}{I, 3.6.5}）；因此可假设 $B$ 是局部环
（等于 $\mathcal{O}_{f(x)}$）。若 $\mathfrak{n}$ 是 $B$ 的极大理想，
则 $f^{-1}(f(x))$ 是环为 $A/\mathfrak{n}A$ 的仿射概形，并且在
$k(f(x))=B/\mathfrak{n}$ 上有限型
（\egaxref{I.6.4.11-zh}{I, 6.4.11}）。在此情形下，若 a) 成立，
还可假设 $f^{-1}(f(x))$ 约化为点 $x$；于是
$A/\mathfrak{n}A$ 在 $B/\mathfrak{n}$ 上秩有限
（\egaxref{I.6.4.4-zh}{I, 6.4.4}），换言之，$A$ 是拟有限 $B$-模。
反过来，若情况如此，则 $f^{-1}(f(x))$ 是阿廷仿射概形，因而离散
（\egaxref{I.6.4.4-zh}{I, 6.4.4}）；所以 $x$ 在其纤维中孤立，
这就表明 b) 蕴含 a)。
% Locked French authority: ega2-1-fr.tex, 820505 B,
% SHA-256 84EDBE3E83530AF2959B441796337C9DC21EAFCA6A13114A26778760FBF437AC.
% Sequential source scope: lines 10291-10318, Corollary II.6.2.2 and
% Definition II.6.2.3 / printed p.115.

\begin{corollary}[6.2.2]
\label{II.6.2.2-zh}
设 $f:X\to Y$ 是有限型态射。下列条件等价：
\begin{enumerate}
\item[a)] 每个点 $x\in X$ 在其纤维 $f^{-1}(f(x))$ 中是孤立的
（换言之，子空间 $f^{-1}(f(x))$ 是离散的）。
\item[b)] 对每个 $x\in X$，预概形 $f^{-1}(f(x))$ 是有限
$k(f(x))$-预概形。
\item[c)] 对每个 $x\in X$，环 $\mathcal{O}_x$ 是拟有限
$\mathcal{O}_{f(x)}$-模。
\end{enumerate}
\end{corollary}

a) 与 c) 的等价性由 \egaref{II.6.2.1-zh}{6.2.1} 得出。另一方面，
由于 $f^{-1}(f(x))$ 是代数 $k(f(x))$-预概形
（\egaxref{I.6.4.11-zh}{I, 6.4.11}），a) 与 b) 的等价性由
\egaxref{I.6.4.4-zh}{I, 6.4.4} 得出。

\begin{definition}[6.2.3]
\label{II.6.2.3-zh}
当 $f:X\to Y$ 是满足 \egaref{II.6.2.2-zh}{6.2.2} 中等价条件的
有限型态射时，称 $f$ 为\emph{拟有限态射}，或称 $X$ 在 $Y$ 上拟有限。
\end{definition}

显然，每个有限态射都是拟有限的（\egaref{II.6.1.8-zh}{6.1.8}）。
% Locked French authority: ega2-1-fr.tex, 820505 B,
% SHA-256 84EDBE3E83530AF2959B441796337C9DC21EAFCA6A13114A26778760FBF437AC.
% Sequential source scope: lines 10320-10341, Proposition II.6.2.4
% / printed p.115.

\begin{proposition}[6.2.4]
\label{II.6.2.4-zh}
\begin{enumerate}
\item[(i)] 浸入 $X\to Y$，若为闭浸入，或 $X$ 是诺特的，
则它是拟有限态射。
\item[(ii)] 若 $f:X\to Y$ 和 $g:Y\to Z$ 是拟有限态射，则
$g\circ f$ 是拟有限态射。
\item[(iii)] 若 $X$ 和 $Y$ 是 $S$-预概形，$f:X\to Y$ 是拟有限
$S$-态射，则对于基的每个扩张 $g:S'\to S$，
$f_{(S')}:X_{(S')}\to Y_{(S')}$ 都是拟有限态射。
\item[(iv)] 若 $f:X\to Y$ 和 $g:X'\to Y'$ 是两个拟有限 $S$-态射，则
\[
f\times_S g:X\times_S Y\to X'\times_S Y'
\]
是拟有限态射。
\item[(v)] 设 $f:X\to Y$ 和 $g:Y\to Z$ 是两个态射，使得 $g\circ f$
是拟有限态射；若此外 $g$ 分离，或 $X$ 诺特，或 $X\times_ZY$
局部诺特，则 $f$ 是拟有限态射。
\item[(vi)] 若 $f$ 是拟有限态射，则 $f_{\mathrm{red}}$ 也如此。
\end{enumerate}
\end{proposition}
% Locked French authority: ega2-1-fr.tex, 820505 B,
% SHA-256 84EDBE3E83530AF2959B441796337C9DC21EAFCA6A13114A26778760FBF437AC.
% Sequential source scope: lines 10343-10370, proof of Proposition II.6.2.4
% / printed pp.115-116.

若 $f:X\to Y$ 是浸入，则每个纤维都至多约化为一个点，因而断言 (i)
由 \egaxref{I.6.3.4-zh}{I, 6.3.4 (i) 和 6.3.5} 得出。为证明 (ii)，
先注意 $h=g\circ f$ 是有限型态射
（\egaxref{I.6.3.4-zh}{I, 6.3.4 (ii)}）；此外，若 $z=h(x)$ 且
$y=f(x)$，则 $y$ 在 $g^{-1}(z)$ 中孤立，故存在 $Y$ 中 $y$ 的开邻域
$V$，它不与 $g^{-1}(z)$ 中异于 $y$ 的任何点相交；所以
$f^{-1}(V)$ 是 $x$ 的开邻域，并且不与任何

\oldpage[II]{116}

$f^{-1}(y')$ 相交，其中 $y'\neq y$ 且该点属于 $g^{-1}(z)$；由于 $x$
在 $f^{-1}(y)$ 中孤立，它也在
$h^{-1}(z)=f^{-1}(g^{-1}(z))$ 中孤立。为证明 (iii)，可以归结到
$Y=S$ 的情形（\egaxref{I.3.3.11-zh}{I, 3.3.11}）；仍先注意
$f'=f_{(S')}$ 是有限型态射
（\egaxref{I.6.3.4-zh}{I, 6.3.4, (iii)}）；另一方面，若
$x'\in X'=X_{(S')}$，并置 $y'=f'(x')$ 和 $y=g(y')$，则
$f'^{-1}(y')$ 与 $f^{-1}(y)\otimes_{k(y)}k(y')$ 同一
（\egaxref{I.3.6.5-zh}{I, 3.6.5}）；依假设，$f^{-1}(y)$ 在 $k(y)$
上秩有限，故 $f'^{-1}(y')$ 在 $k(y')$ 上秩有限，从而离散。
断言 (iv)、(v)、(vi) 可由 (i)、(ii)、(iii) 借助一般方法
（\egaxref{I.5.5.12-zh}{I, 5.5.12}）推出，但 (v) 中的假设不是
“$g$ 分离”这一假设时除外；为处理这种情形，先注意到，若 $x$ 在
$f^{-1}(g^{-1}(g(f(x))))$ 中孤立，则更是在 $f^{-1}(f(x))$ 中孤立；
而 $f$ 是有限型态射这一事实于是由
\egaxref{I.6.3.6-zh}{I, 6.3.6} 得出。
% Locked French authority: ega2-1-fr.tex, 820505 B,
% SHA-256 84EDBE3E83530AF2959B441796337C9DC21EAFCA6A13114A26778760FBF437AC.
% Sequential source scope: lines 10372-10400, Proposition II.6.2.5 and proof
% / printed p.116.

\begin{proposition}[6.2.5]
\label{II.6.2.5-zh}
设 $A$ 是完备诺特局部环，$X$ 是局部有限型 $A$-概形，$x$ 是 $X$
的一个点，位于 $Y=\operatorname{Spec}(A)$ 的闭点 $y$ 之上，并假设
$x$ 在其纤维 $f^{-1}(y)$ 中孤立（其中 $f$ 是结构态射 $X\to Y$）。
则 $\mathcal{O}_x$ 是有限型 $A$-模，并且 $X$ 在 $Y$ 上同构于
$X'=\operatorname{Spec}(\mathcal{O}_x)$（它是有限 $Y$-概形）与某个
$A$-概形 $X''$ 的和（I, 3.1）。
\end{proposition}

由 \egaref{II.6.2.1-zh}{6.2.1} 可知，$\mathcal{O}_x$ 是拟有限
$A$-模。由于 $\mathcal{O}_x$ 是诺特的
（\egaxref{I.6.3.7-zh}{I, 6.3.7}），且同态
$A\to\mathcal{O}_x$ 是局部同态，$A$ 完备这一假设蕴含
$\mathcal{O}_x$ 是有限型 $A$-模
（\egaxref{0.7.4.3-zh}{0, 7.4.3}）。令
$X'=\operatorname{Spec}(\mathcal{O}_x)$ 为 $X$ 在点 $x$ 处的局部概形
（\egaxref{I.2.4.1-zh}{I, 2.4.1}），并令 $g:X'\to X$ 为典范态射。
由于复合 $X'\xrightarrow{g}X\xrightarrow{f}Y$ 是有限态射
（\egaref{II.6.1.1-zh}{6.1.1}），且 $f$ 分离，故 $g$ 有限
（\egaref{II.6.1.5-zh}{6.1.5, (v)}），所以 $g(X')$ 在 $X$ 中闭
（\egaref{II.6.1.10-zh}{6.1.10}）。另一方面，由于 $g$ 是有限型态射，
依 $g$ 的定义和 \egaxref{I.6.5.4-zh}{I, 6.5.4}，$g$ 在 $X'$ 的闭点
$x'$ 处局部是同构；但由于 $X'$ 是 $x'$ 的唯一开邻域，这意味着 $g$
是开浸入，故 $g(X')$ 在 $X$ 中也开；证明完毕。
% Locked French authority: ega2-1-fr.tex, 820505 B,
% SHA-256 84EDBE3E83530AF2959B441796337C9DC21EAFCA6A13114A26778760FBF437AC.
% Sequential source scope: lines 10402-10419, Corollary II.6.2.6 and
% Remark II.6.2.7 / printed p.116.

\begin{corollary}[6.2.6]
\label{II.6.2.6-zh}
设 $A$ 是完备诺特局部环，$Y=\operatorname{Spec}(A)$，
$f:X\to Y$ 是分离拟有限态射。则 $X$ 在 $Y$ 上同构于一个和
$X'\amalg X''$，其中 $X'$ 是有限 $Y$-概形，而 $X''$ 是拟有限
$Y$-概形，并且若 $y$ 是 $Y$ 的闭点，则
$X''\cap f^{-1}(y)=\varnothing$。
\end{corollary}

事实上，依假设，纤维 $f^{-1}(y)$ 有限且离散，故此推论由
\egaref{II.6.2.5-zh}{6.2.5} 对该纤维的点数作归纳而得。

\begin{remark}[6.2.7]
\label{II.6.2.7-zh}
我们将在第五章看到：若 $Y$ 局部诺特，则分离拟有限态射
$X\to Y$ 必为拟仿射态射。
\end{remark}
% Locked French authority: ega2-1-fr.tex, 820505 B,
% SHA-256 84EDBE3E83530AF2959B441796337C9DC21EAFCA6A13114A26778760FBF437AC.
% Sequential source scope: lines 10421-10479, subsection II.6.3 opening,
% Proposition II.6.3.1 and Corollary II.6.3.2 / printed pp.116-117.

\subsection{预概形的整闭包。}
\label{subsection:II.6.3-zh}

\begin{proposition}[6.3.1]
\label{II.6.3.1-zh}
设 $(X,\mathcal{A})$ 是环化空间，$\mathcal{B}$ 是一个交换
$\mathcal{A}$-代数，$f$ 是 $\mathcal{B}$ 在 $X$ 上的一个截面。
下列性质等价：
\begin{enumerate}
\item[a)] 由所有 $f^n$（$n\geq 0$）生成的 $\mathcal{B}$ 的
$\mathcal{A}$-子模（\egaxref{0.5.1.1-zh}{0, 5.1.1}）是有限型的。
\item[b)] 存在 $\mathcal{B}$ 的一个 $\mathcal{A}$-子代数
$\mathcal{C}$，它作为 $\mathcal{A}$-模是有限型的，并且
$f\in\Gamma(X,\mathcal{C})$。
\item[c)] 对每个 $x\in X$，$f_x$ 在茎 $\mathcal{A}_x$ 上是整的。
\end{enumerate}
\end{proposition}

\oldpage[II]{117}

由于由所有 $f^n$ 生成的 $\mathcal{B}$ 的 $\mathcal{A}$-子模本身
就是一个 $\mathcal{A}$-代数，故 a) 蕴含 b) 是显然的。另一方面，
b) 蕴含对每个 $x\in X$，$\mathcal{A}_x$-模 $\mathcal{C}_x$
是有限型的；由此可知，代数 $\mathcal{C}_x$ 的每个元素，特别是
$f_x$，都在 $\mathcal{A}_x$ 上是整的。最后，若对某个 $x\in X$
有形如
\[
f_x^n+(a_1)_xf_x^{n-1}+\ldots+(a_n)_x=0
\]
的关系，其中 $a_i$（$1\leq i\leq n$）是 $\mathcal{A}$ 在 $x$
的某个开邻域 $U$ 上的截面，则截面
$f^n|U+a_1\cdot f^{n-1}|U+\ldots+a_n$ 在 $x$ 的某个邻域
$V\subset U$ 上为零。于是立即得到，所有 $f^k|V$（$k\geq0$）
都是 $f^j|V$（$0\leq j\leq n-1$）以
$\Gamma(V,\mathcal{A})$ 中元素为系数的线性组合；因此 c) 蕴含 a)。

当 \egaref{II.6.3.1-zh}{6.3.1} 的等价条件成立时，称截面 $f$
\emph{在 $\mathcal{A}$ 上是整的}。

\begin{corollary}[6.3.2]
\label{II.6.3.2-zh}
在 \egaref{II.6.3.1-zh}{6.3.1} 的假设下，存在 $\mathcal{B}$ 的
一个唯一的 $\mathcal{A}$-子模 $\mathcal{A}'$，使得对每个
$x\in X$，$\mathcal{A}'_x$ 都是在 $\mathcal{A}_x$ 上整的芽
$f_x\in\mathcal{B}_x$ 的集合。对每个开集 $U\subset X$，
$\mathcal{A}'$ 在 $U$ 上的截面恰为那些在
$\mathcal{A}|U$ 上整的截面 $f\in\Gamma(U,\mathcal{B})$。
\end{corollary}

$\mathcal{A}'$ 的存在性是直接的：只需令
$\Gamma(U,\mathcal{A}')$ 为所有满足下述条件的
$f\in\Gamma(U,\mathcal{B})$ 的集合：对每个 $x\in U$，$f_x$
都在 $\mathcal{A}_x$ 上整。第二个断言立即由
\egaref{II.6.3.1-zh}{6.3.1} 得出。

显然，$\mathcal{A}'$ 是 $\mathcal{B}$ 的一个
$\mathcal{A}$-子代数；称它为\emph{$\mathcal{A}$ 在
$\mathcal{B}$ 中的整闭包}。
% Locked French authority: ega2-1-fr.tex, 820505 B,
% SHA-256 84EDBE3E83530AF2959B441796337C9DC21EAFCA6A13114A26778760FBF437AC.
% Sequential source scope: lines 10481-10512, environment II.6.3.3
% / printed p.117.

\begin{env}[6.3.3]
\label{II.6.3.3-zh}
设 $(X,\mathcal{A})$、$(Y,\mathcal{B})$ 是两个环化空间，
\[
g=(\psi,\theta):X\to Y
\]
是一个态射。设 $\mathcal{C}$（相应地，$\mathcal{D}$）是一个
$\mathcal{A}$-代数（相应地，一个 $\mathcal{B}$-代数），并设
\[
u:\mathcal{D}\to\mathcal{C}
\]
是一个 $g$-态射（\egaxref{0.4.4.1-zh}{0, 4.4.1}）。若
$\mathcal{A}'$（相应地，$\mathcal{B}'$）是 $\mathcal{A}$
（相应地，$\mathcal{B}$）在 $\mathcal{C}$（相应地，
$\mathcal{D}$）中的整闭包，则 $u$ 在 $\mathcal{B}'$ 上的
\emph{限制}是一个 $g$-态射
\[
u':\mathcal{B}'\to\mathcal{A}'.
\]

事实上，若 $j$ 是典范单射 $\mathcal{B}'\to\mathcal{D}$，则只需证明
\[
v=u^\sharp\circ g^*(j):g^*(\mathcal{B}')\to\mathcal{C}
\]
把 $g^*(\mathcal{B}')$ 映入 $\mathcal{A}'$。而
$(g^*(\mathcal{B}'))_x=\mathcal{B}'_{\psi(x)}\otimes_{\mathcal{B}_{\psi(x)}}
\mathcal{A}_x$ 的每个元素依 $\mathcal{B}'$ 的定义都在
$\mathcal{A}_x$ 上整，故其在 $v_x$ 下的像也是如此，这就证明了该断言。
\end{env}
% Locked French authority: ega2-1-fr.tex, 820505 B,
% SHA-256 84EDBE3E83530AF2959B441796337C9DC21EAFCA6A13114A26778760FBF437AC.
% Sequential source scope: lines 10513-10543, Proposition II.6.3.4 and
% its consequences / printed pp.117-118.

\begin{proposition}[6.3.4]
\label{II.6.3.4-zh}
设 $X$ 是预概形，$\mathcal{A}$ 是拟凝聚 $\mathcal{O}_X$-代数。
则 $\mathcal{O}_X$ 在 $\mathcal{A}$ 中的整闭包
$\mathcal{O}'_X$ 是拟凝聚 $\mathcal{O}_X$-代数，并且对 $X$
的任意仿射开集 $U$，$\Gamma(U,\mathcal{O}'_X)$ 是
$\Gamma(U,\mathcal{O}_X)$ 在 $\Gamma(U,\mathcal{A})$ 中的整闭包。
\end{proposition}

只须考虑 $X=\operatorname{Spec}(B)$ 为仿射且
$\mathcal{A}=\widetilde{A}$ 的情形，其中 $A$ 是

\oldpage[II]{118}

$B$-代数；设 $B'$ 是 $B$ 在 $A$ 中的整闭包。问题归结为证明：
对每个 $x\in X$，$A_x$ 中在 $B_x$ 上整的元素必属于 $B'_x$。
这是因为，对任一交换环 $C$，在一个 $C$-代数中取整闭包与取分式环
（相对于 $C$ 的一个乘法子集）这两种运算可交换
（\egaxcite[t.~I, p.~261 et 257]{I-13}）。

于是，$X$-概形 $X'=\operatorname{Spec}(\mathcal{O}'_X)$ 称为
\emph{$X$ 相对于 $\mathcal{A}$ 的整闭包}（或\emph{相对于
$\operatorname{Spec}(\mathcal{A})$ 的整闭包}）；显然，$X'$ 在
$X$ 上是\emph{整的}（\egaref{II.6.1.2-zh}{6.1.2}）。

由 \egaref{II.6.3.4-zh}{6.3.4} 立即推出：若
$f:X'\to X$ 是结构态射，则对 $X$ 的任意开集 $U$，
$f^{-1}(U)$ 是 $X$ 在 $U$ 上诱导的预概形相对于
$\mathcal{A}|U$ 的\emph{整闭包}。
% Locked French authority: ega2-1-fr.tex, 820505 B,
% SHA-256 84EDBE3E83530AF2959B441796337C9DC21EAFCA6A13114A26778760FBF437AC.
% Sequential source scope: lines 10545-10567, environment II.6.3.5
% / printed p.118.

\begin{env}[6.3.5]
\label{II.6.3.5-zh}
设 $X$、$Y$ 是两个预概形，$f:X\to Y$ 是态射，
$\mathcal{A}$（相应地，$\mathcal{B}$）是拟凝聚
$\mathcal{O}_X$-代数（相应地，$\mathcal{O}_Y$-代数），并设
$u:\mathcal{B}\to\mathcal{A}$ 是一个 $f$-态射。我们已在
\egaref{II.6.3.3-zh}{6.3.3} 中看到，由此得到一个 $f$-态射
$u':\mathcal{O}'_Y\to\mathcal{O}'_X$，其中
$\mathcal{O}'_X$（相应地，$\mathcal{O}'_Y$）是
$\mathcal{O}_X$（相应地，$\mathcal{O}_Y$）在
$\mathcal{A}$（相应地，$\mathcal{B}$）中的整闭包。因此，若
$X'$（相应地，$Y'$）是 $X$（相应地，$Y$）相对于
$\mathcal{A}$（相应地，$\mathcal{B}$）的整闭包，则由 $u$
典范地得到一个态射 $f'=\operatorname{Spec}(u'):X'\to Y'$
（\egaref{II.1.5.6-zh}{1.5.6}），使下图交换：
\[
\label{II.6.3.5.1-zh}
\xymatrix{
X'\ar[r]^{f'}\ar[d]&Y'\ar[d]\\
X\ar[r]_{f}&Y
}
\tag{6.3.5.1}
\]
\end{env}
% Locked French authority: ega2-1-fr.tex, 820505 B,
% SHA-256 84EDBE3E83530AF2959B441796337C9DC21EAFCA6A13114A26778760FBF437AC.
% Sequential source scope: lines 10569-10604, environment II.6.3.6
% / printed p.118.

\begin{env}[6.3.6]
\label{II.6.3.6-zh}
假设 $X$ 只有\emph{有限}多个不可约分支 $X_i$
（$1\leq i\leq r$），其泛点为 $\xi_i$；特别考虑 $X$
相对于一个拟凝聚 $\mathcal{R}(X)$-代数 $\mathcal{A}$ 的整闭包
（无论将其视为 $\mathcal{O}_X$-代数还是
$\mathcal{R}(X)$-代数，两者等价）。已知
（\egaxref{I.7.3.5-zh}{I, 7.3.5}），$\mathcal{A}$ 是 $r$ 个
拟凝聚 $\mathcal{O}_X$-代数 $\mathcal{A}_i$ 的直积；
$\mathcal{A}_i$ 的支撑包含于 $X_i$，而 $\mathcal{A}_i$ 在
$X_i$ 上诱导的层是一个单纯层，其纤维 $A_i$ 是
$\mathcal{O}_{\xi_i}$-代数。于是显然
（\egaref{II.6.3.4-zh}{6.3.4}），$\mathcal{O}_X$ 在
$\mathcal{A}$ 中的整闭包 $\mathcal{O}'_X$ 是
$\mathcal{O}_X$ 在各 $\mathcal{A}_i$ 中的整闭包
$\mathcal{O}^{(i)}_X$ 的\emph{直积}；从而，$X$ 相对于
$\mathcal{A}$ 的整闭包
$X'=\operatorname{Spec}(\mathcal{O}'_X)$ 是各
$\operatorname{Spec}(\mathcal{O}^{(i)}_X)=X'_i$
（$1\leq i\leq r$）的 $X$-概形\emph{不交并}。

再假设 $\mathcal{O}_X$-代数 $\mathcal{A}$ 是\emph{约化的}，
或者等价地，每个代数 $A_i$ 都是约化的，因而可视为
\emph{域} $k(\xi_i)$ 上的代数（该域等于以 $X_i$ 为底空间的
$X$ 的约化子预概形的有理函数域）；则由
\egaref{II.1.3.8-zh}{1.3.8}，每个 $X'_i$ 都是约化
$X$-概形，而且 $X'$ 也是 $X_{\mathrm{red}}$ 的整闭包。
进一步假设每个代数 $A_i$ 都是\emph{有限多个域 $K_{ij}$
的直积}（$1\leq j\leq s_i$）；若 $\mathcal{K}_{ij}$ 是
$\mathcal{A}_i$ 中对应于 $K_{ij}$ 的子代数，则显然
$\mathcal{O}^{(i)}_X$ 是 $\mathcal{O}_X$ 在各
$\mathcal{K}_{ij}$ 中的整闭包 $\mathcal{O}^{(ij)}_X$
的\emph{直积}。因此，$X'_i$ 是各 $X$-概形
$X'_{ij}=\operatorname{Spec}(\mathcal{O}^{(ij)}_X)$
（$1\leq j\leq s_i$）的\emph{不交并}。此外，在这些假设和记号下：
\end{env}
% Locked French authority: ega2-1-fr.tex, 820505 B,
% SHA-256 84EDBE3E83530AF2959B441796337C9DC21EAFCA6A13114A26778760FBF437AC.
% Sequential source scope: lines 10606-10633, Proposition II.6.3.7 and proof
% / printed p.119.

\oldpage[II]{119}

\begin{proposition}[6.3.7]
\label{II.6.3.7-zh}
每个 $X'_{ij}$ 都是整且正规的 $X$-预概形，且其有理函数域
$R(X'_{ij})$ 可典范地等同于 $k(\xi_i)$ 在 $K_{ij}$ 中的代数闭包
$K'_{ij}$。
\end{proposition}

由上述结果，可设 $X$ 是整的，因而 $r=1$ 且 $s_1=1$，从而唯一的代数
$A_1$ 是一个域 $K$；令 $\xi$ 为 $X$ 的泛点，并令 $f:X'\to X$ 为结构
态射。对 $X$ 的任一非空仿射开集 $U$，$f^{-1}(U)$ 可等同于整环
$B_U=\Gamma(U,\mathcal{O}_X)$ 在域 $K$ 中的整闭包 $B'_U$ 的谱
（\egaref{II.6.3.4-zh}{6.3.4}）；由于环 $B'_U$ 是整闭整环，其谱上
各点的局部环也都是整闭整环，故按定义，$f^{-1}(U)$ 是一个
\emph{整且正规的}预概形（\egaxref{0.4.1.4-zh}{0, 4.1.4} 与
\egaxref{I.5.1.4-zh}{I, 5.1.4}）。此外，由于 $(0)$ 是 $B_U$ 的素理想
$(0)$ 上方 $B'_U$ 中唯一的素理想
（\egaxcite[t.~I, p.~259]{I-13}），$f^{-1}(\xi)$ 仅由一个点 $\xi'$ 组成，
而 $k(\xi')$ 是 $B'_U$ 的分式域 $K'$，也就是 $k(\xi)$ 在 $K$ 中的
代数闭包。最后，$X'$ 是不可约的：当 $U$ 遍历 $X$ 的所有非空仿射
开集时，各 $f^{-1}(U)$ 构成 $X'$ 的一个由不可约开集组成的开覆盖；
而其中任意两个开集的交 $f^{-1}(U\cap V)$ 都含有 $\xi'$，因而非空，
故由 \egaxref{0.2.1.4-zh}{0, 2.1.4} 即得结论。
% Locked French authority: ega2-1-fr.tex, 820505 B,
% SHA-256 84EDBE3E83530AF2959B441796337C9DC21EAFCA6A13114A26778760FBF437AC.
% Sequential source scope: lines 10635-10650, Corollary II.6.3.8 and
% normalization terminology / printed p.119.

\begin{corollary}[6.3.8]
\label{II.6.3.8-zh}
设 $X$ 是约化预概形，只有有限个不可约分支 $X_i$
（$1\leq i\leq r$），并令 $\xi_i$ 为 $X_i$ 的泛点。$X$ 相对于
$\mathcal{R}(X)$ 的整闭包 $X'$ 是 $r$ 个整且正规的 $X$-预概形
$X'_i$ 的不交并。若 $f:X'\to X$ 是结构态射，则 $f^{-1}(\xi_i)$ 仅由
$X'_i$ 的泛点 $\xi'_i$ 组成，并且
$k(\xi'_i)=k(\xi_i)$；换言之，$f$ 是双有理态射。
\end{corollary}

在此情形下，称 $X'$ 为约化预概形 $X$ 的\emph{正规化}；注意，由于
$f$ 是双有理整态射，它是\emph{满射}
（\egaref{II.6.1.10-zh}{6.1.10}）。此时，$X'=X$ 当且仅当 $X$ 正规。
当 $X$ 是整预概形时，由 \egaref{II.6.3.8-zh}{6.3.8} 可知，其正规化
$X'$ 也是整的。
% Locked French authority: ega2-1-fr.tex, 820505 B,
% SHA-256 84EDBE3E83530AF2959B441796337C9DC21EAFCA6A13114A26778760FBF437AC.
% Sequential source scope: lines 10652-10694, environment II.6.3.9
% / printed pp.119-120.

\begin{env}[6.3.9]
\label{II.6.3.9-zh}
设 $X$、$Y$ 是两个整预概形，$f:X\to Y$ 是支配态射，并令
$L=R(X)$、$K=R(Y)$ 分别为 $X$ 与 $Y$ 的有理函数域；$f$ 典范地对应
于一个单射 $K\to L$，而当把 $K$（相应地，$L$）等同于简单层
$\mathcal{R}(Y)$（相应地，$\mathcal{R}(X)$）时，此单射是一个
$f$-态射。设 $K_1$（相应地，$L_1$）是 $K$（相应地，$L$）的一个
扩张，并设给定一个单同态 $K_1\to L_1$，使图
\[
\xymatrix{
K_1\ar[r]&L_1\\
K\ar[u]\ar[r]&L\ar[u]
}
\]
交换；若把 $K_1$（相应地，$L_1$）看作 $Y$（相应地，$X$）上的
简单层，从而看作一个 $\mathcal{R}(Y)$-代数（相应地，一个
$\mathcal{R}(X)$-代数），这就是说 $K_1\to L_1$ 是一个 $f$-态射。
在此设定下，若 $X'$（相应地，$Y'$）是 $X$（相应地，$Y$）相对于
$L_1$（相应地，$K_1$）的整闭包，则 $X'$（相应地，$Y'$）是整且
正规的预概形（\egaref{II.6.3.6-zh}{6.3.6}），其有理函数域可典范地
等同于 $L$（相应地，$K$）在 $L_1$（相应地，$K_1$）中的代数闭包

\oldpage[II]{120}

$L'$（相应地，$K'$）；并且存在一个典范态射（必为支配态射）
$f':X'\to Y'$，使图 \egaref{II.6.3.5.1-zh}{6.3.5.1} 交换。

最重要的情形是取 $L_1=L$，因而 $K_1$ 是 $K$ 的一个包含于 $L$ 中的
扩张，并假设 $X$ 是整且\emph{正规}的，从而 $X'=X$。上述结果于是
表明：当 $X$ 正规，且 $Y'$ 是 $Y$ 相对于某个域
$K_1\subset L=R(X)$ 的整闭包时，任一支配态射 $f:X\to Y$ 都可分解为
\[
f:X\xrightarrow{f'}Y'\to Y
\]
其中 $f'$ 是支配态射；此外，一旦给定单同态 $K_1\to L$，$f'$ 必然
唯一，这一点可通过归约到 $X$ 与 $Y$ 都是仿射的情形看出。因此，
给定 $Y$、$L$ 以及一个 $K$-单同态 $K_1\to L$ 时，$Y$ 相对于 $K_1$
的整闭包 $Y'$ 是一个\emph{泛问题}的解。
\end{env}
% Locked French authority: ega2-1-fr.tex, 820505 B,
% SHA-256 84EDBE3E83530AF2959B441796337C9DC21EAFCA6A13114A26778760FBF437AC.
% Sequential source scope: lines 10696-10716, Remark II.6.3.10
% / printed p.120.

\begin{remark}[6.3.10]
\label{II.6.3.10-zh}
仍取 \egaref{II.6.3.6-zh}{6.3.6} 的假设，并进一步假设每个代数
$A_i$ 在 $k(\xi_i)$ 上都具有\emph{有限秩}（这蕴含 $A_i$ 是有限个域
的直积）；在某些情形下，可以断言结构态射 $X'\to X$ 不仅是整态射，
而且还是\emph{有限态射}。这里只讨论 $X$ \emph{约化}的情形；由于
该问题在 $X$ 上是局部的，还可假设 $X$ 是环为 $C$ 的仿射预概形，
且 $C$ 只有有限个极小素理想 $\mathfrak{p}_i$
（$1\leq i\leq r$），使得 $C_i=C/\mathfrak{p}_i$ 都是整环；此时，
若每个 $C_i$ 在其分式域的任一有限次扩张中的整闭包都是有限型
$C$-模，则 $X'$ 在 $X$ 上有限
（\egaref{II.6.3.4-zh}{6.3.4}）。已知当 $C$ 是\emph{域上的有限型
代数}（\egaxcite[t.~I, p.~267, th.~9]{I-13}）、\emph{$\mathbf{Z}$ 上的
有限型代数}（\egaxcite[t.~I, p.~93, th.~3]{I-9}），或\emph{完备诺特
局部环上的有限型代数}（[25], p.~298）时，此条件总是成立。因此，
当 $X$ 是域上、$\mathbf{Z}$ 上或完备诺特局部环上的\emph{有限型
概形}时，$X'\to X$ 是有限态射。
\end{remark}
\subsection{\texorpdfstring{$\mathcal{O}_X$}{O\_X}-模自同态的行列式。}
\label{subsection:II.6.4-zh}

\begin{env}[6.4.1]
\label{II.6.4.1-zh}
设 $A$ 是（交换）环，$E$ 是秩为 $n$ 的自由 $A$-模，$u$ 是
$E$ 的一个自同态。回顾如下：为了定义 $u$ 的\emph{特征多项式}，
考虑秩为 $n$ 的自由 $A[T]$-模 $E\otimes_AA[T]$（$T$ 为不定元）
的自同态 $u\otimes1$，并置
\[
\label{II.6.4.1.1-zh}
P(u,T)=\det(T\cdot I-(u\otimes1))
\tag{6.4.1.1}
\]
（其中 $I$ 是 $E\otimes_AA[T]$ 的恒等自同构）。有
\[
\label{II.6.4.1.2-zh}
P(u,T)=T^n-\sigma_1(u)T^{n-1}+\ldots+(-1)^n\sigma_n(u)
\tag{6.4.1.2}
\]
其中 $\sigma_i(u)$ 是 $A$ 的一个元素，并且是由 $u$ 相对于
$E$ 的任意一个基的矩阵元素构成的、整系数的 $i$ 次齐次多项式；
称 $\sigma_i(u)$ 为 $u$ 的\emph{初等对称函数}。特别地，有
$\sigma_1(u)=\operatorname{Tr}u$ 和 $\sigma_n(u)=\det u$。回顾
凯莱--哈密顿定理，有
\[
\label{II.6.4.1.3-zh}
P(u,u)=u^n-\sigma_1(u)u^{n-1}+\ldots+(-1)^n\sigma_n(u)=0
\tag{6.4.1.3}
\]

\oldpage[II]{121}

这也可写成
\[
\label{II.6.4.1.4-zh}
(\det u)\cdot I_E=uQ(u)
\tag{6.4.1.4}
\]
（其中 $I_E$ 是 $E$ 的恒等自同构），这里
\[
\label{II.6.4.1.5-zh}
Q(u)=(-1)^{n+1}(u^{n-1}-\sigma_1(u)u^{n-2}+\ldots+(-1)^{n-1}\sigma_{n-1}(u)).
\tag{6.4.1.5}
\]

设 $\varphi:A\to B$ 是环同态，它使 $B$ 成为一个 $A$-代数；
考虑自由且秩为 $n$ 的 $B$-模 $E_{(B)}=E\otimes_AB$，以及 $u$
在 $E_{(B)}$ 上的自同态扩张 $u\otimes1$。立即可得，对每个指标
$i$，都有 $\sigma_i(u\otimes1)=\varphi(\sigma_i(u))$。
\end{env}
\begin{env}[6.4.2]
\label{II.6.4.2-zh}
现在假设 $A$ 是一个\emph{整环}，其分式域为 $K$，并且 $E$
是一个\emph{有限型} $A$-模（但不一定自由）。设 $n$ 是 $E$ 的
\emph{秩}，即自由 $K$-模 $E\otimes_AK$ 的秩。$E$ 的每个自同态
$u$ 都典范地对应于 $E\otimes_AK$ 的自同态 $u\otimes1$。沿用术语，
仍称多项式 $P(u\otimes1,T)$ 为 $u$ 的\emph{特征多项式}，并将其记为
$P(u,T)$；其系数 $\sigma_i(u\otimes1)$（它们属于 $K$）仍记为
$\sigma_i(u)$，并称为 $u$ 的\emph{初等对称函数}。特别地，依定义
$\det u=\det(u\otimes1)$。在这些记号下，从
\egaref{II.6.4.1.3-zh}{6.4.1.3} 到
\egaref{II.6.4.1.5-zh}{6.4.1.5} 的公式都有意义且仍然成立；这里把
$u^j$ 理解为同态 $E\to E\otimes_AK$，它由典范同态
$x\mapsto x\otimes1$ 后接 $E\otimes_AK$ 的自同态
$u^j\otimes1=(u\otimes1)^j$ 复合而成。

若 $F$ 是 $E$ 的挠子模，并且 $E_0=E/F$，则 $u(F)\subset F$，
故取商后，$u$ 给出 $E_0$ 的一个自同态 $u_0$。此外，
$E\otimes_AK$ 可与 $E_0\otimes_AK$ 等同，而 $u\otimes1$ 可与
$u_0\otimes1$ 等同，因此对 $1\leq i\leq n$，有
$\sigma_i(u)=\sigma_i(u_0)$。

当 $E$ \emph{无挠}时，$E$ 典范地等同于 $E\otimes_AK$ 的一个
$A$-子模，并且关系 $u\otimes1=0$ 等价于 $u=0$。当 $E$ 是自由
$A$-模时，由以上所述，在 \egaref{II.6.4.1-zh}{6.4.1} 和
\egaref{II.6.4.2-zh}{6.4.2} 中给出的 $\sigma_i(u)$ 的两个定义
彼此相同，这就说明了所采用记号的合理性。

还应注意，当 $E$ 是一个\emph{挠模}时，换言之，当
$E_0=\{0\}$ 时，$E_0$ 的外代数约化为 $K$，而 $E_0$ 的唯一
自同态 $u_0$ 的行列式\emph{等于 $1$}。
\end{env}
\begin{proposition}[6.4.3]
\label{II.6.4.3-zh}
设 $A$ 是整环，$E$ 是有限型 $A$-模，$u$ 是 $E$ 的一个自同态。
则 $u$ 的初等对称函数 $\sigma_i(u)$（特别是 $\det u$）都是
$K$ 中在 $A$ 上整的元素。
\end{proposition}

设 $A'$ 是 $A$ 的整闭包。由于 $A'[T]$ 是整闭环（[24], p.~99），
它就是 $A[T]$ 在其分式域 $K(T)$ 中的整闭包。用
$T\cdot I-u\otimes1$ 代替 $u$，并用 $A[T]$ 代替 $A$，可见问题归结为
证明 $\det u$ 在 $A$ 上整。若 $n$ 是 $E$ 的秩，则
$\det u=\det(\bigwedge^n u)$ 且
$(\bigwedge^n u)\otimes1=\bigwedge^n(u\otimes1)$，因此可假设
$n=1$。但在此情形下，映射 $u\mapsto\det u$ 是从 $A$-模
$\operatorname{Hom}_A(E,E)$ 到 $K$ 的一个同态。由于 $E$ 是有限型的，
$\operatorname{Hom}_A(E,E)$ 同构于有限型 $A$-模 $E^n$ 的一个子模
（若 $E$ 有一个由 $n$ 个元素组成的生成系）；所以诸元素 $\det u$
属于 $K$ 的某个有限型 $A$-子模，从而都在 $A$ 上整。

\oldpage[II]{122}
\begin{corollary}[6.4.4]
\label{II.6.4.4-zh}
在 \egaref{II.6.4.3-zh}{6.4.3} 的假设下，若再假设 $A$ 正规，
则各 $\sigma_i(u)$（特别是 $\operatorname{Tr}u$ 与 $\det u$）
都属于 $A$。
\end{corollary}

\begin{proposition}[6.4.5]
\label{II.6.4.5-zh}
设 $A$ 是整环，$E$ 是秩为 $n$ 的有限型 $A$-模，$u$ 是 $E$
的一个自同态，并且各 $\sigma_i(u)$ 都属于 $A$。若 $u$ 是 $E$
的自同构，则 $\det u$ 必须在 $A$ 中可逆；当 $E$ 无挠时，
此条件也是充分的。
\end{proposition}

该条件是\emph{充分的}，因为由假设以及公式
\egaref{II.6.4.1.4-zh}{6.4.1.4} 和
\egaref{II.6.4.1.5-zh}{6.4.1.5}（由于 $E$ 无挠，这些公式在
$E$ 中成立，而不只是在 $E\otimes_AK$ 中成立），可知
$(\det u)^{-1}Q(u)$ 是 $u$ 的逆。

该条件是\emph{必要的}，因为若 $u$ 可逆，则由
\egaref{II.6.4.3-zh}{6.4.3}，$\det(u^{-1})$ 属于 $A$ 在其分式域
$K$ 中的整闭包 $A'$，并且显然是在 $A'$ 中 $\det u$ 的逆元。
于是我们的断言由下述引理推出：

\begin{lemma}[6.4.5.1]
\label{II.6.4.5.1-zh}
设 $A$ 是环 $A'$ 的一个子环，并且 $A'$ 在 $A$ 上整。若元素
$x\in A$ 在 $A'$ 中可逆，则它在 $A$ 中也可逆。
\end{lemma}

否则，$x$ 将属于 $A$ 的某个极大理想 $\mathfrak{m}$；由
Cohen--Seidenberg 第一定理
（\egaxcite[t.~I, p.~257, th.~3]{I-13}），存在 $A'$ 的极大理想
$\mathfrak{m}'$，使得 $\mathfrak{m}=\mathfrak{m}'\cap A$；
于是 $x\in\mathfrak{m}'$，这就产生矛盾。
\begin{corollary}[6.4.6]
\label{II.6.4.6-zh}
设 $A$ 是整闭整环，$E$ 是有限型无挠 $A$-模，$u$ 是 $E$
的一个自同态。则 $u$ 是 $E$ 的自同构，当且仅当 $\det u$
在 $A$ 中可逆。
\end{corollary}

这由 \egaref{II.6.4.4-zh}{6.4.4} 和
\egaref{II.6.4.5-zh}{6.4.5} 得出。

\begin{remark}[6.4.7]
\label{II.6.4.7-zh}
以后将需要前述结果的一种推广。考虑一个\emph{约化}诺特环 $A$；
设 $\mathfrak{p}_\alpha$（$1\leq\alpha\leq r$）是它的极小素理想，
$K_\alpha$ 是整环 $A/\mathfrak{p}_\alpha$ 的分式域，$K$ 是 $A$
的全分式环，即各域 $K_\alpha$ 的\emph{直积}。设 $E$ 是有限型
$A$-模，并\emph{假设} $E\otimes_AK$ 是秩为 $n$ 的自由 $K$-模
（在这里，这不再由其他假设推出）；等价地说，所有
$K_\alpha$-向量空间 $E\otimes_AK_\alpha=E_\alpha$ 都具有
\emph{相同的维数 $n$}。若 $u$ 是 $E$ 的一个自同态，仍令
$P(u,T)=P(u\otimes1,T)$ 且
$\sigma_j(u)=\sigma_j(u\otimes1)$，特别地，
$\det u=\det(u\otimes1)$；因此各 $\sigma_j(u)$ 都是 $K$ 的元素。
显然，$E\otimes_AK$ 是各 $E_\alpha$ 的直和，且后者在
$u\otimes1$ 下稳定；$u\otimes1$ 在 $E_\alpha$ 上的限制正是
$u$ 在这个 $K_\alpha$-向量空间上的延拓 $u_\alpha$；由此可知，
$\sigma_j(u)$ 是 $K$ 中这样的元素：它在各 $K_\alpha$ 中的分量
为 $\sigma_j(u_\alpha)$。由于 $A$ 在 $K$ 中的整闭包是 $A$
在各 $K_\alpha$ 中的整闭包的\emph{直积}，所以各
$\sigma_j(u)$ 在 $A$ 上都是\emph{整的}。

\begin{lemma}[6.4.7.1]
\label{II.6.4.7.1-zh}
当 $u$ 遍历 $\operatorname{Hom}_A(E,E)$ 时，由所有元素
$\sigma_j(u)$（$1\leq j\leq n$）生成的 $K$ 的 $A$-子代数，
作为 $A$-模是有限型的。
\end{lemma}

只须证明，由各 $P(u,T)$ 生成的 $K[T]$ 的子-$A[T]$-代数作为
$A[T]$-模是有限型的。事实上，若 $F_i(T)$
（$1\leq i\leq m$）构成这个 $A[T]$-模的一组生成元，则
各 $F_i(T)$ 的系数都在 $A$ 上整，因而它们生成一个作为
$A$-模为有限型的 $A$-代数
（\egaxcite[t.~I, p.~255, th.~1]{I-13}）。因此，可以

\oldpage[II]{123}

用 $A[T]$（它是诺特环）代替 $A$，并用
$E\otimes_AA[T]=E'$ 代替 $E$；此时
$E'\otimes_{A[T]}K[T]=E\otimes_AK[T]$ 是秩为 $n$ 的自由
$K[T]$-模。回到原来的记号，只须证明：当 $u$ 遍历
$\operatorname{Hom}_A(E,E)$ 时，由各元素 $\det u$ 生成的
$A$-模是有限型的。\emph{更强地}（因为有限型 $A$-模的每个
子模都是有限型的），只须证明：当 $v$ 遍历 $\bigwedge^nE$
的自同态集合时，由各 $\det v$ 生成的 $A$-模是有限型的；
换言之，问题又归结为 $n=1$ 的情形。但在此情形下，该命题由
$\operatorname{Hom}_A(E,E)$ 是有限型 $A$-模以及
$v\mapsto\det v$ 是从这个 $A$-模到 $K$ 的同态这两个事实得出。

设 $F$ 是典范同态 $E\to E\otimes_AK$ 的核，并令 $E_0=E/F$；
如上可见，$E\otimes_AK$ 可与 $E_0\otimes_AK$ 等同，
$u(F)\subset F$；若 $u_0$ 是由 $u$ 取商后在 $E_0$ 上诱导的
自同态，则 $u\otimes1$ 可与 $u_0\otimes1$ 等同，并且对每个 $j$
都有 $\sigma_j(u)=\sigma_j(u_0)$。\emph{若 $F=0$}，则公式
\egaref{II.6.4.1.3-zh}{6.4.1.3} 和
\egaref{II.6.4.1.5-zh}{6.4.1.5} 都有意义，并且当把 $E$
等同于 $E\otimes_AK$ 的一个子模、把各 $u^j$ 视为同态
$E\to E\otimes_AK$ 时仍然成立；因此，命题
\egaref{II.6.4.5-zh}{6.4.5} 连同其证明也推广到这种情形。
\end{remark}
\begin{env}[6.4.8]
\label{II.6.4.8-zh}
设 $(X,\mathcal{A})$ 是环化空间，$\mathcal{E}$ 是一个\emph{局部自由}
（有限秩）$\mathcal{A}$-模，$u$ 是 $\mathcal{E}$ 的一个自同态。依假设，
$X$ 的拓扑有一个基 $\mathfrak{B}$，使得对每个
$V\in\mathfrak{B}$，$\mathcal{E}|V$ 都同构于
$\mathcal{A}^n|V$（其中 $n$ 可随 $V$ 而变）。令 $u$ 为
$\mathcal{E}$ 的一个自同态；于是对每个 $V\in\mathfrak{B}$，$u_V$
都是 $\Gamma(V,\mathcal{A})$-模 $\Gamma(V,\mathcal{E})$ 的一个
自同态，而该模依假设是自由的；因此行列式 $\det u_V$ 有定义，并且
属于 $\Gamma(V,\mathcal{A})$。此外，若 $e_1,\ldots,e_n$ 构成
$\Gamma(V,\mathcal{E})$ 的一个基，则它们在任一开集 $W\subset V$
上的限制构成 $\Gamma(W,\mathcal{E})$ 在
$\Gamma(W,\mathcal{A})$ 上的一个基，故 $\det u_W$ 是
$\det u_V$ 在 $W$ 上的限制。因此，$\mathcal{A}$ 在 $X$ 上恰有一个
截面；我们将它记为 $\det u$，并称为 $u$ 的\emph{行列式}，且
$\det u$ 在每个 $V\in\mathfrak{B}$ 上的限制都是 $\det u_V$。显然，对每个
$x\in X$，都有 $(\det u)_x=\det u_x$；对于 $\mathcal{E}$ 的两个
自同态 $u$、$v$，有
\[
\label{II.6.4.8.1-zh}
\det(u\circ v)=(\det u)(\det v)
\tag{6.4.8.1}
\]
以及
\[
\label{II.6.4.8.2-zh}
\det(1_{\mathcal{E}})=1_{\mathcal{A}}
\tag{6.4.8.2}
\]
并且，若 $\mathcal{E}$ 的秩为常数（当 $X$ 连通时确是如此
（\egaxref{0.5.4.1-zh}{0, 5.4.1}））且等于 $n$，则
\[
\label{II.6.4.8.3-zh}
\det(s\cdot u)=s^n\det u
\tag{6.4.8.3}
\]
对每个 $s\in\Gamma(X,\mathcal{A})$ 都成立（注意，当 $n\geq1$ 时，
$\det(0)=0_{\mathcal{A}}$；但当 $n=0$ 时，
$\det(0)=1_{\mathcal{A}}$）。此外，$u$ 是 $\mathcal{E}$ 的一个
\emph{自同构}，当且仅当 $\det u$ 在
$\Gamma(X,\mathcal{A})$ 中\emph{可逆}。

若 $\mathcal{E}$ 的秩为常数，则可用同样的方法定义初等对称函数
$\sigma_i(u)$，它们是 $\Gamma(X,\mathcal{A})$ 的元素；关系
\egaref{II.6.4.1.3-zh}{6.4.1.3} 至
\egaref{II.6.4.1.5-zh}{6.4.1.5} 仍然成立。

\oldpage[II]{124}

这样便定义了乘法幺半群的一个同态
$\det:\operatorname{Hom}_{\mathcal{A}}(\mathcal{E},\mathcal{E})
\to\Gamma(X,\mathcal{A})$；注意按定义有
$\operatorname{Hom}_{\mathcal{A}}(\mathcal{E},\mathcal{E})=
\Gamma(X,\mathcal{H}om_{\mathcal{A}}(\mathcal{E},\mathcal{E}))$，
由此可见，在此定义中可将 $X$ 换成 $X$ 的任一开集 $U$，从而立即
定义出乘法幺半群层的一个同态
$\det:\mathcal{H}om_{\mathcal{A}}(\mathcal{E},\mathcal{E})\to\mathcal{A}$。
当 $\mathcal{E}$ 的秩为常数时，同样可定义集合层的同态
$\sigma_i:\mathcal{H}om_{\mathcal{A}}(\mathcal{E},\mathcal{E})
\to\mathcal{A}$；当 $i=1$ 时，同态
$\sigma_1=\operatorname{Tr}$ 是一个 $\mathcal{A}$-模同态。

设 $(Y,\mathcal{B})$ 是另一个环化空间，并设
$f:(X,\mathcal{A})\to(Y,\mathcal{B})$ 是环化空间的一个态射；若
$\mathcal{F}$ 是局部自由 $\mathcal{B}$-模，则
$f^*(\mathcal{F})$ 是局部自由 $\mathcal{A}$-模（若
$\mathcal{F}$ 的秩为常数，则前者与后者具有相同的秩）
（\egaxref{0.5.4.5-zh}{0, 5.4.5}）。对 $\mathcal{F}$ 的任一
自同态 $v$，$f^*(v)$ 是 $f^*(\mathcal{F})$ 的一个自同态；由定义
立即可知，$\det f^*(v)$ 是 $\mathcal{A}=f^*(\mathcal{B})$ 在 $X$
上的截面，它典范地对应于
$\det v\in\Gamma(Y,\mathcal{B})$。换言之，同态
$f^*(\det):f^*(\mathcal{H}om_{\mathcal{B}}(\mathcal{F},\mathcal{F}))
\to f^*(\mathcal{B})=\mathcal{A}$ 是复合
\[
f^*(\mathcal{H}om_{\mathcal{B}}(\mathcal{F},\mathcal{F}))
\xrightarrow{\gamma^\sharp}
\mathcal{H}om_{\mathcal{A}}(f^*(\mathcal{F}),f^*(\mathcal{F}))
\xrightarrow{\det}\mathcal{A}
\]
（\egaxref{0.4.4.6-zh}{0, 4.4.6}）。对于各 $\sigma_i$，也有类似的
结果。
\end{env}
\begin{env}[6.4.9]
\label{II.6.4.9-zh}
现设 $X$ 是\emph{局部整预概形}，因而 $X$ 上的有理函数层
$\mathcal{R}(X)$ 是一个局部简单的域层
（\egaxref{I.7.3.4-zh}{I, 7.3.4}），并且作为
$\mathcal{O}_X$-模是拟凝聚的。若 $\mathcal{E}$ 是\emph{有限型}
拟凝聚 $\mathcal{O}_X$-模，则
$\mathcal{E}'=\mathcal{E}\otimes_{\mathcal{O}_X}\mathcal{R}(X)$
是局部自由 $\mathcal{R}(X)$-模
（\egaxref{I.7.3.6-zh}{I, 7.3.6}）；对 $\mathcal{E}$ 的任一
自同态 $u$，$u\otimes1_{\mathcal{R}(X)}$ 因而是
$\mathcal{E}'$ 的一个自同态，而 $\det(u\otimes1)$ 是
$\mathcal{R}(X)$ 在 $X$ 上的一个截面；仍称它为 $u$ 的
\emph{行列式}，并仍记为 $\det u$。由
\egaref{II.6.4.3-zh}{6.4.3} 可知，$\det u$ 是
$\mathcal{O}_X$ 在 $\mathcal{R}(X)$ 中的\emph{整闭包}
（\egaref{II.6.3.2-zh}{6.3.2}）的一个截面；若进一步假设 $X$
\emph{正规}，则 $\det u$ 是 $X$ 上的一个
\emph{$\mathcal{O}_X$ 的截面}；若再假设 $\mathcal{E}$
\emph{无挠}，则由 \egaref{II.6.4.6-zh}{6.4.6}，$u$ 是
$\mathcal{E}$ 的自同构，当且仅当 $\det u$ \emph{可逆}。公式
\egaref{II.6.4.8.1-zh}{6.4.8.1} 至
\egaref{II.6.4.8.3-zh}{6.4.8.3} 仍然成立；把同态
$u\mapsto\det u$ 应用于
$\mathcal{H}om_{\mathcal{O}_X}(\mathcal{E},\mathcal{E})$ 的截面模，
可得到层的同态
$\det:\mathcal{H}om_{\mathcal{O}_X}(\mathcal{E},\mathcal{E})
\to\mathcal{R}(X)$；当 $X$ 正规时，该同态取值于
$\mathcal{O}_X$。当 $\mathcal{E}'$ 的\emph{秩为常数}时，对于其余
初等对称函数 $\sigma_j(u)$ 也有类似的定义和结果；若进一步假设
$X$ \emph{正规}，则各 $\sigma_j(u)$ 都是 $X$ 上的
\emph{$\mathcal{O}_X$ 的截面}。

最后，设 $X$、$Y$ 是两个整预概形，并设 $f:X\to Y$ 是一个
\emph{支配}态射。已知存在典范同态
$f^*(\mathcal{R}(Y))\to\mathcal{R}(X)$
（\egaxref{I.7.3.8-zh}{I, 7.3.8}\footnote{见\emph{勘误与补遗}，
清单~1。}）；由此，对任一有限型拟凝聚 $\mathcal{O}_Y$-模
$\mathcal{F}$，得到典范同态
$\theta:f^*(\mathcal{F}\otimes_{\mathcal{O}_Y}\mathcal{R}(Y))
\to f^*(\mathcal{F})\otimes_{\mathcal{O}_X}\mathcal{R}(X)$。若 $v$ 是
$\mathcal{F}$ 的一个自同态，

\oldpage[II]{125}

则 $f^*(v\otimes1_{\mathcal{R}(Y)})$ 是
$f^*(\mathcal{F}\otimes_{\mathcal{O}_Y}\mathcal{R}(Y))$ 的一个
自同态，并且有交换图
\[
\xymatrix{
f^*(\mathcal{F}\otimes_{\mathcal{O}_Y}\mathcal{R}(Y))
  \ar[r]^{f^*(v\otimes1)}\ar[d]_{\theta}
& f^*(\mathcal{F}\otimes_{\mathcal{O}_Y}\mathcal{R}(Y))\ar[d]^{\theta}\\
f^*(\mathcal{F})\otimes_{\mathcal{O}_X}\mathcal{R}(X)
  \ar[r]_{f^*(v)\otimes1}
& f^*(\mathcal{F})\otimes_{\mathcal{O}_X}\mathcal{R}(X).
}
\]
由此容易得出，$\det f^*(v)$ 是 $\mathcal{R}(Y)$ 的截面
$\det v$ 在同态 $f^*(\mathcal{R}(Y))\to\mathcal{R}(X)$ 下的典范像；
事实上，可以立即归约到
$X=\operatorname{Spec}(A)$、$Y=\operatorname{Spec}(B)$ 都是仿射的
情形，其中 $A$、$B$ 是整环，其分式域分别为 $K$、$L$，同态
$B\to A$ 为单射，因而延拓为单同态 $L\to K$；若
$\mathcal{F}=\widetilde{M}$，其中 $M$ 是有限型 $B$-模，则
$M\otimes_B L$ 在 $L$ 上的秩与
$(M\otimes_B A)\otimes_A K$ 在 $K$ 上的秩\emph{相等}，并且对于
$M$ 的任一自同态 $u$，$\det((u\otimes1)\otimes1)$ 是
$\det(u\otimes1)$ 在 $K$ 中的像，结论由此得证。
\end{env}
\begin{env}[6.4.10]
\label{II.6.4.10-zh}
最后，设 $X$ 是\emph{局部诺特约化预概形}，因而 $X$ 上的有理函数层
$\mathcal{R}(X)$ 仍是拟凝聚 $\mathcal{O}_X$-模
（\egaxref{I.7.3.4-zh}{I, 7.3.4}）；另一方面，设 $\mathcal{E}$ 是
一个\emph{凝聚} $\mathcal{O}_X$-模，使得
$\mathcal{E}'=\mathcal{E}\otimes_{\mathcal{O}_X}\mathcal{R}(X)$
\emph{局部自由}（且秩有限）。由
\egaref{II.6.4.7-zh}{6.4.7}，若 $\mathcal{E}'$ 的秩为常数，则对
$\mathcal{E}$ 的任一自同态 $u$，可以定义各 $\sigma_j(u)$；它们是
$\mathcal{R}(X)$ 在 $X$ 上的截面。当不再假设 $\mathcal{E}'$ 的秩
为常数时，仍可定义同态
$\det:\mathcal{H}om_{\mathcal{O}_X}(\mathcal{E},\mathcal{E})
\to\mathcal{R}(X)$。
\end{env}
\subsection{可逆层的范数。}
\label{subsection:II.6.5-zh}

\begin{env}[6.5.1]
\label{II.6.5.1-zh}
设 $(X,\mathcal{A})$ 是一个环化空间，$\mathcal{B}$ 是一个
（交换）$\mathcal{A}$-代数。$\mathcal{A}$-模 $\mathcal{B}$ 典范地
等同于 $\mathcal{H}om_{\mathcal{A}}(\mathcal{B},\mathcal{B})$ 的一个
$\mathcal{A}$-子模；$\mathcal{B}$ 在 $X$ 的开集 $U$ 上的一个截面
$f$，等同于由该截面给出的乘法。若 $(X,\mathcal{A})$ 与
$\mathcal{B}$ 满足 \egaref{II.6.4.8-zh}{6.4.8}、
\egaref{II.6.4.9-zh}{6.4.9} 或 \egaref{II.6.4.10-zh}{6.4.10}
中列出的某个条件，则可以定义 $\det(f)$（并且在某些情形下定义
$\sigma_j(f)$）；它们是 $U$ 上 $\mathcal{A}$ 或
$\mathcal{R}(X)$ 的截面。我们把它们分别称为 $f$ 的\emph{范数}
和 $f$ 的\emph{初等对称函数}，并将前者记为
$N_{\mathcal{B}/\mathcal{A}}(f)$。以下假设满足下列条件之一：
\begin{enumerate}
\item[(I)] \emph{$\mathcal{B}$ 是局部自由的（有限秩）
$\mathcal{A}$-模。}
\item[(II)] \emph{$(X,\mathcal{A})$ 是局部诺特约化预概形，
$\mathcal{B}$ 是一个凝聚 $\mathcal{A}$-模，使得
$\mathcal{B}\otimes_{\mathcal{A}}\mathcal{R}(X)$ 是局部自由的
$\mathcal{R}(X)$-模，并且对每个开集 $U\subset X$ 上的截面
$f\in\Gamma(U,\mathcal{B})$，$N_{\mathcal{B}/\mathcal{A}}(f)$
都是 $U$ 上 $\mathcal{A}$ 的截面。}
\end{enumerate}

\oldpage[II]{126}

应注意，当局部诺特预概形 $X$ 是\emph{正规}时，这最后一个条件
自动满足（\egaref{II.6.4.9-zh}{6.4.9}）。

另一方面，假设
$\mathcal{B}\otimes_{\mathcal{A}}\mathcal{R}(X)$ 局部自由，还可
如下表述。记 $X_\alpha$ 为 $X$ 的约化闭子预概形，其底空间是
$X$ 的不可约分支（\egaxref{I.5.2.1-zh}{I, 5.2.1}）；于是它们
是局部诺特整预概形。每个 $x\in X$ 属于有限多个子空间
$X_\alpha$；另一方面，
$\mathcal{B}\otimes_{\mathcal{A}}\mathcal{R}(X_\alpha)$ 是秩为
常数 $k_\alpha$ 的局部自由 $\mathcal{R}(X_\alpha)$-模
（\egaxref{I.7.3.6-zh}{I, 7.3.6}）。说
$\mathcal{B}\otimes_{\mathcal{A}}\mathcal{R}(X)$ 是局部自由的
$\mathcal{R}(X)$-模，意味着对每个 $x\in X$，\emph{满足
$x\in X_\alpha$ 的那些指标所对应的秩 $k_\alpha$ 都相等}。
这个问题是局部的，可归结为 $X=\operatorname{Spec}(C)$ 的情形，
其中 $C$ 是约化诺特环，$\mathcal{B}=\widetilde{D}$，而 $D$ 是
一个 $C$-代数并且作为 $C$-模是有限型的。若
$\mathfrak{p}_i$（$1\leq i\leq m$）是 $C$ 的极小素理想，则
$C$ 的全分式环 $L$ 是整环 $C/\mathfrak{p}_i$ 的分式域
$K_i$ 的直积，而 $D\otimes_C L$ 是各 $D\otimes_C K_i$ 的直和，
由此得到结论。

显然，在假设 (I) 或 (II) 下，我们由此定义了乘法幺半群层的一个
同态 $N_{\mathcal{B}/\mathcal{A}}:\mathcal{B}\to\mathcal{A}$，
若不会引起混淆，也记为 $N$，并称为\emph{范数同态}。因此，对于
同一开集 $U$ 上 $\mathcal{B}$ 的两个截面 $f$、$g$，有
\[
\label{II.6.5.1.1-zh}
N_{\mathcal{B}/\mathcal{A}}(fg)
=N_{\mathcal{B}/\mathcal{A}}(f)N_{\mathcal{B}/\mathcal{A}}(g)
\tag{6.5.1.1}
\]
这里等式右边是 $U$ 上 $\mathcal{A}$ 的相应截面；并且
\[
\label{II.6.5.1.2-zh}
N_{\mathcal{B}/\mathcal{A}}(1_{\mathcal{B}})=1_{\mathcal{A}};
\tag{6.5.1.2}
\]
最后，对 $U$ 上 $\mathcal{A}$ 的任一截面 $s$，若
$\mathcal{B}$ 的秩为常数且等于 $n$，则
\[
\label{II.6.5.1.3-zh}
N_{\mathcal{B}/\mathcal{A}}(s\cdot1_{\mathcal{B}})=s^n
\tag{6.5.1.3}
\]
（当 $s=0_{\mathcal{A}}$ 时，该公式给出：若 $n\geq1$，则
$N(0_{\mathcal{B}})=0_{\mathcal{A}}$；若 $n=0$，则
$N(0_{\mathcal{B}})=N(1_{\mathcal{B}})=1_{\mathcal{A}}$）。

在假设 (I) 下，$f\in\Gamma(U,\mathcal{B})$ 可逆，当且仅当
$N(f)\in\Gamma(U,\mathcal{A})$ 可逆。在假设 (II) 下，这个条件是
必要的；当进一步假设
$\mathcal{B}\to\mathcal{B}\otimes_{\mathcal{A}}\mathcal{R}(X)$
是\emph{单射}，并且满足下面更强的条件时，它也是充分的（由
\egaref{II.6.4.7-zh}{6.4.7}）：
\begin{enumerate}
\item[(II \emph{bis})] \emph{$(X,\mathcal{A})$ 是局部诺特约化预概形，
$\mathcal{B}$ 是凝聚 $\mathcal{A}$-模，使得
$\mathcal{B}\otimes_{\mathcal{A}}\mathcal{R}(X)$ 是局部自由的
$\mathcal{R}(X)$-模；并且对任意开集上的截面
$f\in\Gamma(U,\mathcal{B})$，只要
$\mathcal{B}\otimes_{\mathcal{A}}\mathcal{R}(X)|U$ 在
$\mathcal{R}(X)|U$ 上具有常数秩 $n$，则 $\sigma_j(f)$
（$1\leq j\leq n$）都是 $U$ 上 $\mathcal{A}$ 的截面。}
\end{enumerate}
（还应注意，当 $X$ 是\emph{正规}时，这个条件也满足。）
\end{env}

\begin{env}[6.5.2]
\label{II.6.5.2-zh}
假设 \egaref{II.6.5.1-zh}{6.5.1} 的条件 (I) 或 (II) 之一成立，
并设 $\mathcal{L}'$ 是一个\emph{可逆} $\mathcal{B}$-模。我们将按如下
方式，典范地（除唯一同构外）赋予它一个可逆 $\mathcal{A}$-模
$\mathcal{L}$。记 $\mathcal{A}^*$（相应地，$\mathcal{B}^*$）为
$\mathcal{A}$（相应地，$\mathcal{B}$）的子层，使得对每个开集
$U\subset X$，$\Gamma(U,\mathcal{A}^*)$（相应地，
$\Gamma(U,\mathcal{B}^*)$）是 $\Gamma(U,\mathcal{A})$
\oldpage[II]{127}
（相应地，$\Gamma(U,\mathcal{B})$）中的可逆元素的集合；它们是
乘法群层，而限制在 $\mathcal{B}^*$ 上的
$N_{\mathcal{B}/\mathcal{A}}$ 是群层的一个同态
$\mathcal{B}^*\to\mathcal{A}^*$
（\egaref{II.6.5.1-zh}{6.5.1}）。设 $\mathfrak{L}$ 为所有二元组
$(U_\lambda,\eta_\lambda)$ 所成的集合，它们满足如下性质：
$U_\lambda$ 是 $X$ 的开集，且 $\eta_\lambda$ 是
$(\mathcal{B}|U_\lambda)$-模的同构
$\mathcal{L}'|U_\lambda\simeq\mathcal{B}|U_\lambda$。
由假设，诸 $U_\lambda$ 构成 $X$ 的一个覆盖；对任意两个指标
$\lambda$、$\mu$，置
$\omega_{\lambda\mu}=(\eta_\lambda|U_\lambda\cap U_\mu)
\circ(\eta_\mu|U_\lambda\cap U_\mu)^{-1}$，
它是 $\mathcal{B}|U_\lambda\cap U_\mu$ 的自同构，并典范地等同于
$U_\lambda\cap U_\mu$ 上 $\mathcal{B}^*$ 的一个截面，而
$(\omega_{\lambda\mu})$ 是以 $\mathcal{B}^*$ 为值、相对于覆盖
$\mathfrak{U}=(U_\lambda)$ 的一个 $1$-上循环
（\egaxref{0.5.4.7-zh}{0, 5.4.7}）。由于
$N_{\mathcal{B}/\mathcal{A}}:\mathcal{B}^*\to\mathcal{A}^*$ 是一个
同态，$(N_{\mathcal{B}/\mathcal{A}}\omega_{\lambda\mu})$ 便是以
$\mathcal{A}^*$ 为值、相对于 $\mathfrak{U}$ 的一个 $1$-上循环；
因此它对应于一个可逆 $\mathcal{A}$-模 $\mathcal{L}$（除唯一同构外）
（\egaxref{0.5.4.7-zh}{0, 5.4.7}）。我们把它记为
$N_{\mathcal{B}/\mathcal{A}}(\mathcal{L}')$，并称为可逆
$\mathcal{B}$-模 $\mathcal{L}'$ 的\emph{范数}。

设 $\mathfrak{M}$ 是 $\mathfrak{L}$ 的一个子集，使得相应的
$U_\lambda$ 仍构成 $X$ 的一个覆盖，并设 $\mathfrak{B}$ 为这个
覆盖。把上循环 $(\omega_{\lambda\mu})$ 限制到 $\mathfrak{B}$，
得到一个以 $\mathcal{A}^*$ 为值、相对于 $\mathfrak{B}$ 的
$1$-上循环 $(N_{\mathcal{B}/\mathcal{A}}\omega_{\lambda\mu})$；它是
$\mathfrak{U}$ 上的 $1$-上循环 $(N_{\mathcal{B}/\mathcal{A}}\omega_{\lambda\mu})$
的限制。显然，由 $\mathfrak{B}$ 的这个 $1$-上循环定义的可逆 $\mathcal{A}$-模
与 $N_{\mathcal{B}/\mathcal{A}}(\mathcal{L}')$ 之间有一个典范同构；
这使我们能够用 $\mathfrak{U}$ 的任意子覆盖来定义这个可逆 $\mathcal{A}$-模。
由此立即可见，若 $\mathcal{L}'_1$、$\mathcal{L}'_2$ 是两个可逆
$\mathcal{B}$-模，则根据
\egaref{II.6.5.1.1-zh}{6.5.1.1} 和 \egaref{II.6.5.1.2-zh}{6.5.1.2}，有
\[
\label{II.6.5.2.1-zh}
N(\mathcal{L}'_1\otimes_{\mathcal{B}}\mathcal{L}'_2)
=N(\mathcal{L}'_1)\otimes_{\mathcal{A}}N(\mathcal{L}'_2)
\tag{6.5.2.1}
\]
以及
\[
\label{II.6.5.2.2-zh}
N_{\mathcal{B}/\mathcal{A}}(\mathcal{B})=\mathcal{A}
\tag{6.5.2.2}
\]
并且
\[
\label{II.6.5.2.3-zh}
N((\mathcal{L}')^{-1})=(N(\mathcal{L}'))^{-1}
\tag{6.5.2.3}
\]
这些等式均是在典范同构意义下成立的。此外，由
\egaref{II.6.5.1.3-zh}{6.5.1.3} 可知，若 $\mathcal{L}$ 是一个
可逆 $\mathcal{A}$-模，并且在情形 (I) 中 $\mathcal{B}$ 在
$\mathcal{A}$ 上的秩为常数 $n$（相应地，在情形 (II) 中
$\mathcal{B}\otimes_{\mathcal{A}}\mathcal{R}(X)$ 在
$\mathcal{R}(X)$ 上的秩为常数 $n$），则在典范同构意义下
\[
\label{II.6.5.2.4-zh}
N_{\mathcal{B}/\mathcal{A}}(\mathcal{L}\otimes_{\mathcal{A}}\mathcal{B})
=\mathcal{L}^{\otimes n}.
\tag{6.5.2.4}
\]
\end{env}
\begin{env}[6.5.3]
\label{II.6.5.3-zh}
现在证明，$N_{\mathcal{B}/\mathcal{A}}(\mathcal{L}')$ 是可逆
$\mathcal{B}$-模范畴中的一个\emph{协变函子}。事实上，设
$h':\mathcal{L}'_1\to\mathcal{L}'_2$ 是可逆 $\mathcal{B}$-模的一个
同态，并设 $\mathfrak{B}=(U_\lambda)$ 是 $X$ 的一个开覆盖，使得对每个
$\lambda$，存在 $(\mathcal{B}|U_\lambda)$-模同构
$\eta^{(1)}_\lambda:\mathcal{L}'_1|U_\lambda\simeq\mathcal{B}|U_\lambda$
和
$\eta^{(2)}_\lambda:\mathcal{L}'_2|U_\lambda\simeq\mathcal{B}|U_\lambda$；
于是，对每个 $\lambda$，存在 $\mathcal{B}|U_\lambda$ 的一个自同态
$h'_\lambda$，满足
$h'_\lambda\circ\eta^{(1)}_\lambda
=\eta^{(2)}_\lambda\circ(h'|U_\lambda)$，
并且显然可以将 $h'_\lambda$ 视为 $\mathcal{B}$ 在
$U_\lambda$ 上的一个截面（\egaxref{0.5.1.1-zh}{0, 5.1.1}）。因此，
对任意一对指标 $(\lambda,\mu)$，在
$U_\lambda\cap U_\mu$ 上，
$(\eta^{(2)}_\lambda)^{-1}\circ h'_\lambda\circ\eta^{(1)}_\lambda$
与 $(\eta^{(2)}_\mu)^{-1}\circ h'_\mu\circ\eta^{(1)}_\mu$ 的限制相同。
由此，对于取值于 $\mathcal{B}^*$、分别对应于
$\mathcal{L}'_1$ 和 $\mathcal{L}'_2$ 的 $1$-上循环
$(\omega^{(1)}_{\lambda\mu})$ 和
$(\omega^{(2)}_{\lambda\mu})$，有关系
\[
\omega^{(2)}_{\lambda\mu}h'_\mu
=h'_\lambda\omega^{(1)}_{\lambda\mu}
\]

\oldpage[II]{128}

令 $h_\lambda=N(h'_\lambda)$，于是有相应的关系
\[
N(\omega^{(2)}_{\lambda\mu})h_\mu
=h_\lambda N(\omega^{(1)}_{\lambda\mu})
\]
从而各 $h_\lambda$ 定义了一个同态
$N(\mathcal{L}'_1)\to N(\mathcal{L}'_2)$，记为
$N_{\mathcal{B}/\mathcal{A}}(h')$ 或 $N(h')$。在假设 (I) 下，
$h'$ 是\emph{同构}的充要条件（因为这是一个局部问题）是
$N_{\mathcal{B}/\mathcal{A}}(h')$ 为同构。在假设 (II) 下，这个条件
仍然是必要的；若假设 (II \emph{bis}) 成立且
$\mathcal{B}\to\mathcal{B}\otimes_{\mathcal{A}}\mathcal{R}(X)$
是单射，则它也是充分的。

特别取 $\mathcal{L}'_1=\mathcal{B}$；此时
$\mathcal{B}\to\mathcal{L}'$ 的同态根据
\egaxref{0.5.1.1-zh}{0, 5.1.1} 与 $\mathcal{L}'$ 在 $X$ 上的截面相对应，
因而得到典范映射
\[
N_{\mathcal{B}/\mathcal{A}}:\Gamma(X,\mathcal{L}')
\to\Gamma(X,N_{\mathcal{B}/\mathcal{A}}(\mathcal{L}'))。
\]
再由 \egaref{II.6.5.1.1-zh}{6.5.1.1} 可知，若
$f'_1\in\Gamma(X,\mathcal{L}'_1)$、$f'_2\in\Gamma(X,\mathcal{L}'_2)$，
则
\[
\label{II.6.5.3.1-zh}
N(f'_1\otimes f'_2)=N(f'_1)\otimes N(f'_2).
\tag{6.5.3.1}
\]
对任意可逆 $\mathcal{A}$-模 $\mathcal{L}$ 及其截面
$f\in\Gamma(X,\mathcal{L})$，考虑到
\egaref{II.6.5.2.4-zh}{6.5.2.4}，有
\[
\label{II.6.5.3.2-zh}
N_{\mathcal{B}/\mathcal{A}}(f\otimes1_{\mathcal{B}})
=f^{\otimes n}
\tag{6.5.3.2}
\]
其中，在假设 (I) 下 $\mathcal{B}$ 的秩恒为 $n$（相应地，在假设
(II) 下 $\mathcal{B}\otimes_{\mathcal{A}}\mathcal{R}(X)$ 的秩恒为
$n$）。最后，与 $\mathcal{L}'$ 在 $X$ 上的截面 $f'$ 相对应的同态
$\mathcal{B}\to\mathcal{L}'$ 是同构，当且仅当对每个 $x\in X$，
$f'_x$ 是 $\mathcal{L}'_x$ 的一个基；在假设 (I) 下，这个条件等价于
对每个 $x$，$(N(f'))_x$ 是 $(N(\mathcal{L}'))_x$ 的一个基；
在假设 (II) 下，该条件仍然是必要的，而当 $\mathcal{B}$ 满足
(II \emph{bis}) 且
$\mathcal{B}\to\mathcal{B}\otimes_{\mathcal{A}}\mathcal{R}(X)$
为单射时，它也是充分的。
\end{env}

\begin{env}[6.5.4]
\label{II.6.5.4-zh}
设 $(X,\mathcal{A})$、$(X',\mathcal{A}')$ 是两个赋环空间，
$f:X'\to X$ 是一个态射，$\mathcal{B}$ 是一个
$\mathcal{A}$-代数，$\mathcal{B}'=f^*(\mathcal{B})$ 是其作为
$\mathcal{A}'$-代数的逆像。假设下列条件之一成立：
\begin{enumerate}
\item[$1^\circ$] $\mathcal{B}$ 满足
\egaref{II.6.5.1-zh}{6.5.1} 的假设 (I)。
\item[$2^\circ$] $(X,\mathcal{A})$ 和 $\mathcal{B}$ 满足
\egaref{II.6.5.1-zh}{6.5.1} 的假设 (II)；$(X',\mathcal{A}')$ 是
局部诺特约化预概形；若以 $X_\alpha$ 和 $X'_\beta$ 分别表示
$X$ 和 $X'$ 的约化闭子预概形，其底空间分别是这些空间的不可约
分支，则 $f$ 在每个 $X'_\beta$ 上的限制是从 $X'_\beta$ 到某个
$X_\alpha$ 的\emph{支配}态射。
\end{enumerate}
在这些条件下，$\mathcal{B}'$ 分别满足
\egaref{II.6.5.1-zh}{6.5.1} 的假设 (I) 或假设 (II)；第一种断言
是直接的。为证明第二种，只须说明：对每个 $x'\in X'$，对于满足
$x'\in X'_\beta$ 的那些 $\beta$，各
$\mathcal{B}'\otimes_{\mathcal{O}_{X'}}\mathcal{R}(X'_\beta)$
的秩彼此\emph{相同}。事实上，若 $f$ 在 $X'_\beta$ 上的限制是到
$X_\alpha$ 中的优势态射，则
$\mathcal{B}'\otimes_{\mathcal{O}_{X'}}\mathcal{R}(X'_\beta)$ 的秩
等于 $\mathcal{B}\otimes_{\mathcal{O}_X}\mathcal{R}(X_\alpha)$ 的秩
（如在 \egaref{II.6.4.9-zh}{6.4.9} 中那样立即归结到仿射情形），
故根据假设 (II) 及
\egaref{II.6.5.1-zh}{6.5.1} 得到所述断言。

\oldpage[II]{129}

由此，再根据 \egaref{II.6.4.8-zh}{6.4.8}、
\egaref{II.6.4.9-zh}{6.4.9} 和
\egaref{II.6.4.10-zh}{6.4.10}，若 $s$ 是 $\mathcal{B}$ 在
$U\subset X$ 上的一个截面，$s'$ 是 $\mathcal{B}'$ 在
$f^{-1}(U)$ 上对应的截面，则
$N_{\mathcal{B}'/\mathcal{A}'}(s')$ 是 $\mathcal{A}'$ 在
$f^{-1}(U)$ 上的截面，并且它对应于 $\mathcal{A}$ 在 $U$ 上的截面
$N_{\mathcal{B}/\mathcal{A}}(s)$。

若 $\mathcal{M}$ 是可逆 $\mathcal{B}$-模，则由上述结果，若
$\mathcal{M}'=f^*(\mathcal{M})$（它是可逆 $\mathcal{B}'$-模），则有
$N_{\mathcal{B}'/\mathcal{A}'}(\mathcal{M}')
=f^*(N_{\mathcal{B}/\mathcal{A}}(\mathcal{M}))$
相差一个典范同构。
\end{env}
\begin{env}[6.5.5]
\label{II.6.5.5-zh}
现设 $(X,\mathcal{A})$ 是一个\emph{预概形}。如所知，给定一个拟凝聚的
$\mathcal{A}$-代数 $\mathcal{B}$，它作为 $\mathcal{A}$-模是\emph{有限型}的，
等价于给定一个\emph{有限}态射 $g:X'\to X$，使得
$g_*(\mathcal{O}_{X'})=\mathcal{B}$，且只在 $X$-同构意义下确定
（\egaref{II.6.1.2-zh}{6.1.2} 和 \egaxref{II.1.3.1-zh}{1.3.1}）。此外，
给定一个拟凝聚的 $\mathcal{O}_{X'}$-模层 $\mathcal{F}'$，等价于给定一个
拟凝聚的 $\mathcal{B}$-模层 $\mathcal{F}$，使得
$g_*(\mathcal{F}')=\mathcal{F}$（\egaref{II.1.4.3-zh}{1.4.3}）；而
$\mathcal{F}'$ 可逆，当且仅当 $\mathcal{F}$ 可逆
（\egaref{II.6.1.12-zh}{6.1.12}）。因此，若要用有限态射 $g$ 的语言
表述前述结果，就必须假设 $g_*(\mathcal{O}_{X'})$ 是一个
（有限型）局部自由的 $\mathcal{O}_X$-模，或者 $(X,\mathcal{O}_X)$ 与
$g_*(\mathcal{O}_{X'})$ 满足假设 (II)。于是对每个可逆的
$\mathcal{O}_{X'}$-模 $\mathcal{L}'$，置
\[
\label{II.6.5.5.1-zh}
N_{X'/X}(\mathcal{L}')
=N_{g_*(\mathcal{O}_{X'})/\mathcal{O}_X}(g_*(\mathcal{L}'))
\tag{6.5.5.1}
\]
并称其为 $\mathcal{L}'$（相对于 $g$）的\emph{范数}。同样，若
$h':\mathcal{L}'_1\to\mathcal{L}'_2$ 是可逆
$\mathcal{O}_{X'}$-模之间的同态，置
\[
\label{II.6.5.5.2-zh}
N_{X'/X}(h')
=N_{g_*(\mathcal{O}_{X'})/\mathcal{O}_X}(g_*(h')):
N_{X'/X}(\mathcal{L}'_1)\to N_{X'/X}(\mathcal{L}'_2).
\tag{6.5.5.2}
\]
特别地，取 $\mathcal{L}'_1=\mathcal{O}_{X'}$，得到典范映射
\[
\label{II.6.5.5.3-zh}
N_{X'/X}:\Gamma(X',\mathcal{L}')
\to\Gamma(X,N_{X'/X}(\mathcal{L}')).
\tag{6.5.5.3}
\]
我们把这些翻译中的大部分留给读者，只明确说明下列几项：
\end{env}

\begin{proposition}[6.5.6]
\label{II.6.5.6-zh}
设 $g:X'\to X$ 是有限态射，并假设 $g_*(\mathcal{O}_{X'})$ 是局部自由的
$\mathcal{O}_X$-模，或者 $(X,\mathcal{O}_X)$ 与
$g_*(\mathcal{O}_{X'})$ 满足 (II \emph{bis})（特别地，当 $X$ 局部诺特且
正规时即如此）。对于可逆 $\mathcal{O}_{X'}$-模之间的同态
$h':\mathcal{L}'_1\to\mathcal{L}'_2$，它是同构，当且仅当在第一种假设下
$N_{X'/X}(h')$ 是同构；在第二种假设下，这个条件仍是必要的，并且当
同态
$g_*(\mathcal{O}_{X'})\to
g_*(\mathcal{O}_{X'})\otimes_{\mathcal{O}_X}\mathcal{R}(X)$
为单射时也是充分的。
\end{proposition}

这里使用了如下事实：$g_*(h')$ 是同构，当且仅当 $h'$ 是同构
（\egaref{II.1.4.2-zh}{1.4.2}）。

\begin{corollary}[6.5.7]
\label{II.6.5.7-zh}
设 $g:X'\to X$ 是有限态射，并假设 $g_*(\mathcal{O}_{X'})$ 是局部自由的
$\mathcal{O}_X$-模，或者 $(X,\mathcal{O}_X)$ 与
$g_*(\mathcal{O}_{X'})$ 满足 (II \emph{bis})，且
$g_*(\mathcal{O}_{X'})\to
g_*(\mathcal{O}_{X'})\otimes_{\mathcal{O}_X}\mathcal{R}(X)$
为单射。设 $\mathcal{L}'$ 是一个可逆的 $\mathcal{O}_{X'}$-模，$f'$ 是
$\mathcal{L}'$ 在 $X'$ 上的一个截面，且
$f=N_{X'/X}(f')$ 是相应的 $\mathcal{L}=N_{X'/X}(\mathcal{L}')$
在 $X$ 上的截面（\egaref{II.6.5.5.3-zh}{6.5.5.3}）。

\oldpage[II]{130}

于是有 $g(X'-X'_{f'})=X-X_f$，并且 $X_f$ 是 $X$ 中满足
$g^{-1}(U)\subset X'_{f'}$ 的最大开集 $U$。
\end{corollary}

事实上，$g(X'-X'_{f'})$ 在 $X$ 中是闭的
（\egaref{II.6.1.10-zh}{6.1.10}）；所以只须证明最后一个断言。关系
$U\subset X_f$ 等价于由 $f|U$ 定义的同态
$\mathcal{O}_X|U\to\mathcal{L}|U$ 是同构。根据
\egaref{II.6.5.6-zh}{6.5.6}，这又等价于由
$f'|g^{-1}(U)$ 定义的同态
$\mathcal{O}_{X'}|g^{-1}(U)\to\mathcal{L}'|g^{-1}(U)$ 是同构，
也就是 $g^{-1}(U)\subset X'_{f'}$。

\begin{proposition}[6.5.8]
\label{II.6.5.8-zh}
设 $g:X'\to X$ 是有限态射，$f:Y\to X$ 是一个态射；设
$Y'=X'_{(Y)}$、$g'=g_{(Y)}$、$f'=f_{(X')}$，从而有交换图
\[
\xymatrix{
X'\ar[d]_g & Y'\ar[l]_{f'}\ar[d]^{g'}\\
X & Y\ar[l]^f
}
\]
假设 $g_*(\mathcal{O}_{X'})$ 局部自由，或者 $(X,\mathcal{O}_X)$ 与
$g_*(\mathcal{O}_{X'})$ 满足 (II)，$Y$ 是局部诺特约化预概形，并且
$f$ 在 $Y$ 的每个不可约分支上的限制，是到 $X$ 的某个不可约分支的
支配态射。则对每个可逆的 $\mathcal{O}_{X'}$-模 $\mathcal{L}'$，有
\[
N_{Y'/Y}({f'}^*(\mathcal{L}'))=f^*(N_{X'/X}(\mathcal{L}'))
\]
相差一个典范同构。
\end{proposition}

注意，根据 \egaref{II.1.5.2-zh}{1.5.2}，有
$f^*(g_*(\mathcal{L}'))=g'_*({f'}^*(\mathcal{L}'))$；特别地，
$g'_*(\mathcal{O}_{Y'})=f^*(g_*(\mathcal{O}_{X'}))$。因此，若
$g_*(\mathcal{O}_{X'})$ 局部自由，则 $g'_*(\mathcal{O}_{Y'})$ 也局部自由。
结论随即由定义和 \egaref{II.6.5.4-zh}{6.5.4} 得出。

\begin{remark}[6.5.9]
\label{II.6.5.9-zh}
我们将在后文推广上述范数的概念，并将其与除子的直像概念联系起来。
\end{remark}
\subsection{应用：丰沛判据。}
\label{subsection:II.6.6-zh}

\begin{proposition}[6.6.1]
\label{II.6.6.1-zh}
设 $Y$ 是一个预概形，$f:X\to Y$ 是拟紧态射，$g:X'\to X$ 是有限满射态射。
假设下列两种情形之一成立：$g_*(\mathcal{O}_{X'})$ 是局部自由的
$\mathcal{O}_X$-模；或者 $(X,\mathcal{O}_X)$ 与
$g_*(\mathcal{O}_{X'})$ 满足 (II \emph{bis})。那么，对任意对 $f\circ g$
丰沛的可逆 $\mathcal{O}_{X'}$-模 $\mathcal{L}'$，令
$N_{X'/X}(\mathcal{L}')=\mathcal{L}$，则其对 $f$ 丰沛。
\end{proposition}

可以假设 $Y$ 是仿射的（\egaref{II.4.6.4-zh}{4.6.4}）；于是依据
（\egaref{II.4.6.6-zh}{4.6.6}），上述命题等价于下述推论。

\begin{corollary}[6.6.2]
\label{II.6.6.2-zh}
设 $X$ 是拟紧预概形，$g:X'\to X$ 是有限满射态射，并且
$g_*(\mathcal{O}_{X'})$ 是局部自由的 $\mathcal{O}_X$-模，或者
$(X,\mathcal{O}_X)$ 与 $g_*(\mathcal{O}_{X'})$ 满足 (II \emph{bis})。
那么，对任意丰沛的 $\mathcal{O}_{X'}$-模 $\mathcal{L}'$，
$\mathcal{L}=N_{X'/X}(\mathcal{L}')$ 是丰沛的。
\end{corollary}

在第二种假设下，还可以附加假设典范同态
$g_*(\mathcal{O}_{X'})\to
g_*(\mathcal{O}_{X'})\otimes_{\mathcal{O}_X}\mathcal{R}(X)$
是\emph{单射}。事实上，若非如此，令 $\mathcal{T}$ 为该同态的核；

\oldpage[II]{131}

它是 $g_*(\mathcal{O}_{X'})=\mathcal{B}$ 的一个凝聚理想层
（\egaxref{I.6.1.1-zh}{I, 6.1.1}），并置
$X''=\operatorname{Spec}(\mathcal{B}/\mathcal{T})$。于是有交换图
\[
\xymatrix{
X''\ar[rr]^h\ar[dr]_{g'} && X'\ar[dl]^g\\
& X
}
\]
其中 $h$ 是闭浸入（\egaref{II.1.4.10-zh}{1.4.10}）。此外，我们知道
$\mathcal{T}$ 的支集是 $X$ 中的一个稀疏闭集
（\egaxref{0.5.2.2-zh}{0, 5.2.2}、\egaxref{I.7.4.6-zh}{I, 7.4.6}），
因此对于 $X$ 的每个不可约分支的泛点 $x$，存在 $x$ 的一个仿射开邻域
$U$，使得
$\mathcal{B}|U=(\mathcal{B}/\mathcal{T})|U$。由于按假设 $g$ 是满射，
可得 $x\in g'(X'')$；故 $g'$ 是支配态射，而它又是有限态射，所以是\emph{满射}
（\egaref{II.6.1.10-zh}{6.1.10}）。按定义有
\[
g'_*(\mathcal{O}_{X''})\otimes_{\mathcal{O}_X}\mathcal{R}(X)
=(\mathcal{B}/\mathcal{T})\otimes_{\mathcal{O}_X}\mathcal{R}(X)
=g_*(\mathcal{O}_{X'})\otimes_{\mathcal{O}_X}\mathcal{R}(X),
\]
所以 $(X,\mathcal{O}_X)$ 与 $g'_*(\mathcal{O}_{X''})$ 满足
(II \emph{bis})，并且
\[
g'_*(\mathcal{O}_{X''})\longrightarrow
g'_*(\mathcal{O}_{X''})\otimes_{\mathcal{O}_X}\mathcal{R}(X)
\]
是单射。最后，$h^*(\mathcal{L}')=\mathcal{L}''$ 是一个丰沛的
$\mathcal{O}_{X''}$-模（\egaref{II.4.6.13-zh}{4.6.13, (i \emph{bis})}），并且
$N_{X''/X}(\mathcal{L}'')=N_{X'/X}(\mathcal{L}')$。事实上，为定义这两个
$\mathcal{O}_X$-可逆模，可以取 $X$ 的同一个仿射开覆盖 $(U_\lambda)$，使得
$g_*(\mathcal{L}')$ 和 $g'_*(\mathcal{L}'')$ 在 $U_\lambda$ 上的限制分别同构于
$\mathcal{B}|U_\lambda$ 和 $(\mathcal{B}/\mathcal{T})|U_\lambda$。显然，每个同构
\[
\eta_\lambda:g_*(\mathcal{L}')|U_\lambda\longrightarrow
\mathcal{B}|U_\lambda
\]
典范地对应于一个同构
\[
\eta'_\lambda:g'_*(\mathcal{L}'')|U_\lambda\longrightarrow
(\mathcal{B}/\mathcal{T})|U_\lambda,
\]
从而，若 $(\omega_{\lambda\mu})$ 与 $(\omega'_{\lambda\mu})$ 是分别对应于同构系
$(\eta_\lambda)$ 与 $(\eta'_\lambda)$ 的 $1$-上循环
（\egaref{II.6.5.2-zh}{6.5.2}），则 $\omega'_{\lambda\mu}$ 是
$\omega_{\lambda\mu}\in\Gamma(U_\lambda\cap U_\mu,\mathcal{B})$ 在
$\Gamma(U_\lambda\cap U_\mu,\mathcal{B}/\mathcal{T})$ 中的典范像。由
$\mathcal{T}$ 的定义可得
\[
N_{\mathcal{B}/\mathcal{A}}\omega_{\lambda\mu}
=N_{(\mathcal{B}/\mathcal{T})/\mathcal{A}}\omega'_{\lambda\mu}
\]
（其中 $\mathcal{A}=\mathcal{O}_X$），这就得到所述等式。

因此，在假设 (II \emph{bis}) 时，我们可以假设同态
$g_*(\mathcal{O}_{X'})\to
g_*(\mathcal{O}_{X'})\otimes_{\mathcal{O}_X}\mathcal{R}(X)$
是单射。只需证明：当 $f$ 遍历 $\mathcal{L}^{\otimes n}$（$n>0$）的
$X$-上截面时，各 $X_f$ 构成 $X$ 的拓扑基（\egaref{II.4.5.2-zh}{4.5.2}）。
为此，取 $x\in X$ 及 $x$ 的任意邻域 $U$；由于 $g^{-1}(x)$ 是有限集
（\egaref{II.6.1.7-zh}{6.1.7}），且 $\mathcal{L}'$ 丰沛，存在整数
$n>0$ 及 $X'$ 上 ${\mathcal{L}'}^{\otimes n}$ 的一个截面 $f'$，使得
$X'_{f'}$ 是包含 $g^{-1}(x)$ 且包含于 $g^{-1}(U)$ 的邻域
（\egaref{II.4.5.4-zh}{4.5.4}）。因为
\[
\mathcal{L}^{\otimes n}=N_{X'/X}({\mathcal{L}'}^{\otimes n}),
\]
只需取 $f=N_{X'/X}(f')$；事实上，此时
$X-X_f=g(X'-X'_{f'})$（\egaref{II.6.5.7-zh}{6.5.7}），所以
$x\in X_f\subset U$。
\begin{corollary}[6.6.3]
\label{II.6.6.3-zh}
在 \egaref{II.6.6.1-zh}{6.6.1} 的假设下，一个可逆
$\mathcal{O}_X$-模层 $\mathcal{L}$ 对 $f$ 丰沛，当且仅当
$\mathcal{L}'=g^*(\mathcal{L})$ 对 $f\circ g$ 丰沛。
\end{corollary}

必要性成立，因为 $g$ 是仿射态射
（\egaref{II.5.1.12-zh}{5.1.12}）。为证明充分性，可设 $Y$ 仿射
（\egaref{II.4.6.4-zh}{4.6.4}），从而 $X$ 和 $X'$ 拟紧且
$\mathcal{L}'$ 丰沛
（\egaref{II.4.6.6-zh}{4.6.6}），于是只须证明 $\mathcal{L}$ 丰沛。

然而，点集

\oldpage[II]{132}

$x\in X$ 在第一种（相应地，第二种）假设下满足：
$g_*(\mathcal{O}_{X'})$（相应地，
$g_*(\mathcal{O}_{X'})\otimes_{\mathcal{O}_X}\mathcal{R}(X)$）
在其邻域上具有给定秩 $n$；该点集在 $X$ 中既开又闭。因此，
$X$ 是有限多个这类开集的不交并，故可设 X 等于其中一个开集
（\egaref{II.4.6.17-zh}{4.6.17}）。但此时
$N_{X'/X}(\mathcal{L}')=\mathcal{L}^{\otimes n}$，故由
\egaref{II.6.6.2-zh}{6.6.2}，$\mathcal{L}^{\otimes n}$ 丰沛，
于是 $\mathcal{L}$ 也丰沛（\egaref{II.4.5.6-zh}{4.5.6}）。

\begin{corollary}[6.6.4]
\label{II.6.6.4-zh}
假设 \egaref{II.6.6.1-zh}{6.6.1} 的条件成立，并进一步假设
$f:X\to Y$ 是有限型态射。则 $f$ 是拟射影态射，当且仅当
$f\circ g$ 是拟射影态射。若进一步假设 $Y$ 是拟紧概形，或是底空间
为诺特空间的预概形，则 $f$ 是射影态射，当且仅当
$f\circ g$ 是射影态射。
\end{corollary}

该假设蕴含 $f\circ g$ 是有限型态射。根据拟射影态射的定义
（\egaref{II.5.3.1-zh}{5.3.1}），第一个断言由
\egaref{II.6.6.1-zh}{6.6.1} 和
\egaref{II.6.6.3-zh}{6.6.3} 得出。结合这一结果以及
\egaref{II.5.5.3-zh}{5.5.3, (ii)}，还须证明：当假设 $f$ 是
拟射影态射时，$f$ 是紧合态射，当且仅当 $f\circ g$ 是紧合态射。
然而，此时 $f$ 是分离态射（\egaref{II.5.3.1-zh}{5.3.1}）且为
有限型；由于 $g$ 是满射，所述断言由
\egaref{II.5.4.2-zh}{5.4.2, (ii)} 和
\egaref{II.5.4.3-zh}{5.4.3, (ii)} 得出。

特别地：

\begin{corollary}[6.6.5]
\label{II.6.6.5-zh}
设 $X$ 是域 $K$ 上的有限型预概形，$K'$ 是 $K$ 的有限次扩域。
则 $X$ 在 $K$ 上为射影概形（相应地，拟射影概形），当且仅当
$X'=X\otimes_K K'$ 在 $K'$ 上为射影概形（相应地，拟射影概形）。
\end{corollary}

必要性确实成立
（\egaref{II.5.3.4-zh}{5.3.4, (iii)} 和
\egaref{II.5.5.5-zh}{5.5.5, (iii)}）。反过来，设上述条件成立，
并令 $g:X'\to X$ 为典范投影。显然，$g$ 是有限态射
（\egaref{II.6.1.5-zh}{6.1.5, (iii)}）且为满射
（\egaxref{I.3.5.2-zh}{I, 3.5.2, (ii)}）。此外，
$g_*(\mathcal{O}_{X'})$ 是局部自由的 $\mathcal{O}_X$-模层，因为它
同构于 $\mathcal{O}_X\otimes_K K'$
（\egaref{II.1.5.2-zh}{1.5.2}）。于是由假设以及
\egaref{II.6.1.11-zh}{6.1.11} 和
\egaref{II.5.5.5-zh}{5.5.5, (ii)}，$X'$ 在\emph{$K$}上为射影概形
（相应地，拟射影概形）；再由
\egaref{II.6.6.4-zh}{6.6.4} 得出 $X$ 在 $K$ 上为射影概形
（相应地，拟射影概形）。

在第五章中，我们将证明，当 $K'$ 是 $K$ 的\emph{任意}扩域时，
\egaref{II.6.6.5-zh}{6.6.5} 的结论仍然成立。

本编号的末尾用于证明判据
（\egaref{II.6.6.11-zh}{6.6.11}）；这是
\egaref{II.6.6.1-zh}{6.6.1} 的一个相当技术性的加强，初次阅读时
可以略去。

\begin{lemma}[6.6.6]
\label{II.6.6.6-zh}
设 $X$ 是约化诺特预概形，$\mathcal{E}$ 是相干
$\mathcal{O}_X$-模层，并且
$\mathcal{E}\otimes_{\mathcal{O}_X}\mathcal{R}(X)$ 是秩恒为 $n$
的局部自由 $\mathcal{R}(X)$-模层。则存在约化诺特预概形 $Z$ 及
有限双有理态射 $h:Z\to X$，具有如下性质：集合层的同态
$\sigma_i:\mathcal{H}om_{\mathcal{O}_X}(\mathcal{E},\mathcal{E})
\to\mathcal{R}(X)$（$1\leq i\leq n$）
（参见 \egaref{II.6.4.10-zh}{6.4.10}）将
$\mathcal{H}om_{\mathcal{O}_X}(\mathcal{E},\mathcal{E})$ 映入相干
$\mathcal{O}_X$-代数 $h_*(\mathcal{O}_Z)$。
\end{lemma}

事实上，考虑 $X$ 的一个仿射开集 $U$，其环
$A(U)=A$；设 $E=\Gamma(U,\mathcal{E})$，并令 $C_U$ 为
$R(U)$ 的子代数，由 $u$ 遍历 $\operatorname{Hom}_A(E,E)$ 时的
各 $\sigma_i(u)$ 生成；我们已在
\egaref{II.6.4.7.1-zh}{6.4.7.1} 中看到，这个 $A$-代数具有
\emph{有限秩}。此外，显然，当从仿射开集 $U$ 限制到其仿射开子集
$U'\subset U$ 时，代数 $C_U$ 的构造与限制运算相交换。于是由此定义了
$\mathcal{R}(X)$ 的一个\emph{有限} $\mathcal{O}_X$-子代数层
$\mathcal{C}$，满足对 $X$ 的每个仿射开集 $U$，
$\Gamma(U,\mathcal{C})=C_U$。于是

\oldpage[II]{133}

取 $Z=\operatorname{Spec}(\mathcal{C})$，并令 $h$ 为结构态射，
故它是有限态射（\egaref{II.6.1.2-zh}{6.1.2}）；由于
$\mathcal{C}$ 是约化的，$Z$ 是约化诺特预概形
（\egaref{II.1.3.8-zh}{1.3.8}）。最后，按定义，$C_U$ 的全分式环
是 $R(U)$；又由于 $C_U$ 包含于 $A(U)$ 在 $R(U)$ 中的整闭包内，
$A(U)$ 的极小素理想与 $C_U$ 的极小素理想之间存在一一对应
（\egaxcite[t.~I, p.~259]{I-13}），这就证明 $h$ 是双有理态射并完成证明。

\begin{corollary}[6.6.7]
\label{II.6.6.7-zh}
在 \egaref{II.6.6.6-zh}{6.6.6} 的假设下，设 $W$ 是 $X$ 的一个开集，
并且对每个 $x\in W$，或者 $X$ 在点 $x$ 处正规，或者
$\mathcal{E}_x$ 是自由 $\mathcal{O}_x$-模。则可取 $h$，使得
$h$ 在 $h^{-1}(W)$ 上的限制是从 $h^{-1}(W)$ 到 $W$ 的同构。
\end{corollary}

事实上，两个假设中的任一个都蕴含：若 $U\subset W$ 是仿射开集，
则按 \egaref{II.6.6.6-zh}{6.6.6} 的记号，对每个 $x\in U$ 有
$(\sigma_i(u))_x\in A_x$
（\egaref{II.6.4.3-zh}{6.4.3}），故 $\sigma_i(u)\in A$；
结论于是由 \egaref{II.6.6.6-zh}{6.6.6} 中给出的 $h$ 的定义推出。

\begin{env}[6.6.8]
\label{II.6.6.8-zh}
设 $X$ 是约化诺特预概形，$g:X'\to X$ 是有限满射态射，从而
$\mathcal{B}=g_*(\mathcal{O}_{X'})$ 是相干
$\mathcal{O}_X$-代数；再假设
$\mathcal{B}\otimes_{\mathcal{O}_X}\mathcal{R}(X)$ 是秩恒为 $n$
的局部自由 $\mathcal{R}(X)$-模层。于是可以应用引理
\egaref{II.6.6.6-zh}{6.6.6}，取 $\mathcal{E}=\mathcal{B}$；按照
\egaref{II.6.6.6-zh}{6.6.6} 的记号，得到集合层乘法幺半群的同态
$\sigma_n:\mathcal{H}om_{\mathcal{O}_X}(\mathcal{B},\mathcal{B})
\to h_*(\mathcal{O}_Z)$，并将此同态与典范同态
$\mathcal{B}\to
\mathcal{H}om_{\mathcal{O}_X}(\mathcal{B},\mathcal{B})$
复合；后一个同态来自
\egaref{II.6.5.1-zh}{6.5.1}，因而得到集合层乘法幺半群的同态：
\[
\label{II.6.6.8.1-zh}
N':\mathcal{B}=g_*(\mathcal{O}_{X'})
\longrightarrow
h_*(\mathcal{O}_Z)=\mathcal{C}.
\tag{6.6.8.1}
\]

由此，对任意可逆 $\mathcal{O}_{X'}$-模层 $\mathcal{L}'$，
$g_*(\mathcal{L}')$ 是一个可逆 $\mathcal{B}$-模
（\egaref{II.6.1.12-zh}{6.1.12}），而由
\egaref{II.6.5.2-zh}{6.5.2} 的方法，可以函子性地给
$\mathcal{L}'$ 关联一个可逆 $\mathcal{C}$-模，记为
$N'(g_*(\mathcal{L}'))$。
\end{env}
\begin{lemma}[6.6.9]
\label{II.6.6.9-zh}
设 $X$ 是约化诺特预概形，$g:X'\to X$ 是有限满射态射，并且
$g_*(\mathcal{O}_{X'})\otimes_{\mathcal{O}_X}\mathcal{R}(X)$ 是秩恒为常数的
局部自由 $\mathcal{R}(X)$-模层。则存在约化诺特预概形 $Z$ 及有限双有理态射
$h:Z\to X$，具有如下性质：对任意丰沛的
$\mathcal{O}_{X'}$-模层 $\mathcal{L}'$，满足
$h_*(\mathcal{M})=N'(g_*(\mathcal{L}'))$ 的可逆 $\mathcal{O}_Z$-模层
$\mathcal{M}$ 是丰沛的（记号见 \egaref{II.6.6.8-zh}{6.6.8}）。
\end{lemma}

首先假设同态
$\mathcal{B}\to
\mathcal{B}\otimes_{\mathcal{O}_X}\mathcal{R}(X)$ 是\emph{单射}。
按 \egaref{II.6.6.6-zh}{6.6.6} 定义 $Z$ 和 $h$（取
$\mathcal{E}=g_*(\mathcal{O}_{X'})$）。取 $z\in Z$；需证明存在整数 $m>0$ 及
$Z$ 上 $\mathcal{M}^{\otimes m}$ 的一个截面 $t$，使得 $Z_t$ 是包含 $z$ 的仿射开集
（\egaref{II.4.5.2-zh}{4.5.2}）。令 $x=h(z)$，取含 $x$ 的 $X$ 的仿射开集
$U$；于是 $h^{-1}(U)$ 是 $z$ 的仿射开邻域，只需找到 $t$ 使
$z\in Z_t\subset h^{-1}(U)$，因为此时 $Z_t$ 必为仿射
（\egaref{II.5.5.8-zh}{5.5.8}）。按假设，存在整数 $n>0$ 及一个截面 $s'$，它是
${\mathcal{L}'}^{\otimes n}$ 在 $X'$ 上的截面，使得
\[
\label{II.6.6.9.1-zh}
g^{-1}(x)\subset X'_{s'}\subset g^{-1}(U)
\tag{6.6.9.1}
\]
根据 \egaref{II.4.5.4-zh}{4.5.4}。按定义，$s'$ 也是
$g_*({\mathcal{L}'}^{\otimes n})$ 在 $X$ 上的截面；如同
\egaref{II.6.5.2-zh}{6.5.2} 中那样，它对应于
$s=N'(s')$，即 $N'(g_*({\mathcal{L}'}^{\otimes n}))$ 在

\oldpage[II]{134}

$X$ 上的截面。我们将证明，若 $t$ 是把 $s$ 视为
$\mathcal{M}^{\otimes n}$ 在 $Z$ 上的截面，则 $t$ 满足要求。置
\[
\label{II.6.6.9.2-zh}
V=X-g(X'-X'_{s'})
\tag{6.6.9.2}
\]
根据 \egaref{II.6.6.9.1-zh}{6.6.9.1} 及
\egaref{II.6.1.10-zh}{6.1.10}，这是含 $x$ 且包含于 $U$ 的 $X$ 的开集。
我们将证明
\[
\label{II.6.6.9.3-zh}
h^{-1}(V)\subset Z_t\subset h^{-1}(U),
\tag{6.6.9.3}
\]
从而完成证明。等价地，令 $T$ 为满足 $s_y$ 可逆的 $y\in X$ 的集合，
则该集合包含 $V$ 且包含于 $U$。首先取包含于 $V$ 的仿射开集 $W$；
$g^{-1}(W)$ 是 $X'$ 中的仿射开集，并且由
\egaref{II.6.6.9.2-zh}{6.6.9.2}，对每个 $y'\in g^{-1}(W)$，$s'_{y'}$
都可逆；由 $X$ 和 $\mathcal{B}$ 的假设，可应用
\egaref{II.6.4.7-zh}{6.4.7} 的结果，看到若 $y=g(y')$，则 $s_y$ 可逆，
换言之 $V\subset T$。另一方面，由 \egaref{II.6.4.7-zh}{6.4.7} 也得到，
反之，若 $s_y$ 可逆，则 $s'_{y'}$ 亦可逆；由
\egaref{II.6.6.9.1-zh}{6.6.9.1}，这蕴含 $y'\in g^{-1}(U)$，从而
$y\in U$，故在此情形 $T\subset U$。

同 \egaref{II.6.6.2-zh}{6.6.2} 中的论证一样转到一般情形：将 $X'$ 换成
$X''$，使得
\[
g'_*(\mathcal{O}_{X''})\longrightarrow
g'_*(\mathcal{O}_{X''})\otimes_{\mathcal{O}_X}\mathcal{R}(X)
\]
为单射，并将 $\mathcal{L}'$ 换成一个丰沛的
$\mathcal{O}_{X''}$-模层 $\mathcal{L}''$，满足
$N'(g_*(\mathcal{L}'))=N'(g'_*(\mathcal{L}''))$。因此，引理
\egaref{II.6.6.9-zh}{6.6.9} 在所有情形下均得证（取适当的 $h$）。

\begin{corollary}[6.6.10]
\label{II.6.6.10-zh}
假设 \egaref{II.6.6.9-zh}{6.6.9} 的条件成立；对于任意使 $g^*(\mathcal{L})$
丰沛的 $\mathcal{O}_X$-可逆模层 $\mathcal{L}$，$h^*(\mathcal{L})$ 是丰沛的。
\end{corollary}

事实上，令 $\mathcal{L}'=g^*(\mathcal{L})$，则
$g_*(\mathcal{L}')=\mathcal{L}\otimes_{\mathcal{O}_X}\mathcal{B}$
（\egaxref{0.5.4.10-zh}{0, 5.4.10}），从而
\[
N'(g_*(\mathcal{L}'))=
(\mathcal{L}\otimes_{\mathcal{O}_X}\mathcal{C})^{\otimes n}
\]
（与 \egaref{II.6.5.2.4-zh}{6.5.2.4} 的论证相同）。因此
$\mathcal{M}=(h^*(\mathcal{L}))^{\otimes n}$，且 $\mathcal{M}$ 丰沛，故
$h^*(\mathcal{L})$ 亦丰沛（\egaref{II.4.5.6-zh}{4.5.6}）。

\begin{proposition}[6.6.11]
\label{II.6.6.11-zh}
设 $Y$ 是仿射概形，$X$ 是约化诺特预概形，$f:X\to Y$ 是拟紧态射，
$g:X'\to X$ 是有限满射态射。设 $W$ 是 $X$ 的开子集，并且对每个 $x\in W$，
或者 $X$ 在点 $x$ 处正规，或者存在 $x$ 的开邻域 $T\subset W$，
使 $(g_*(\mathcal{O}_{X'}))|T$ 成为局部自由的
$(\mathcal{O}_X|T)$-模。则存在约化的 $Y$-预概形 $Z$ 及有限双有理的
$Y$-态射 $h:Z\to X$，使得 $h$ 在 $h^{-1}(W)$ 上的限制是同构
$h^{-1}(W)\xrightarrow{\sim}W$，并具有如下性质：对任意
$\mathcal{O}_X$-可逆模层 $\mathcal{L}$，若 $g^*(\mathcal{L})$ 对
$f\circ g$ 丰沛，则 $h^*(\mathcal{L})$ 对 $f\circ h$ 丰沛。
\end{proposition}

由于 $Y$ 仿射，$g^*(\mathcal{L})$ 丰沛；于是只需证明，对于适当选择的 $h$，
$h^*(\mathcal{L})$ 丰沛（\egaref{II.4.6.6-zh}{4.6.6}）。我们将证明，
可以将 $g$ 换成有限满射态射 $g':X''\to X$，使得
${g'}^*(\mathcal{L})$ 丰沛，并且
$g'_*(\mathcal{O}_{X''})\otimes_{\mathcal{O}_X}\mathcal{R}(X)$ 是秩恒为常数的
\emph{局部自由} $\mathcal{R}(X)$-模层；这样就归结为
\egaref{II.6.6.10-zh}{6.6.10} 的条件，命题便得证。

为此，令 $\mathcal{B}=g_*(\mathcal{O}_{X'})$；记 $X_i$
（$1\leq i\leq n$）为 $X$ 的约化闭预概形，其底空间是 $X$ 的不可约分支
（\egaxref{I.5.2.1-zh}{I, 5.2.1}）；按假设它们是整的。令 $X'_i$ 为 $X'$ 的闭预概形

\oldpage[II]{135}

$g^{-1}(X_i)$，令 $g_i:X'_i\to X_i$ 为 $g$ 在 $X'_i$ 上的限制；它是有限的
（\egaref{II.6.1.5-zh}{6.1.5, (iii)}）且满射。令 $k_i$ 为
$\mathcal{O}_{X_i}$-代数 $\mathcal{B}_i=(g_i)_*(\mathcal{O}_{X'_i})$ 的\emph{秩}。
由于 $\mathcal{B}_i\otimes_{\mathcal{O}_{X_i}}\mathcal{R}(X_i)$ 是常值预层
（\egaxref{I.7.3.5-zh}{I, 7.3.5}），对于 $X$ 的任一只遇到\emph{唯一}不可约分支
$X_i$ 的开集 $U$，$k_i$ 也是 $(\mathcal{O}_X|U)$-代数 $\mathcal{B}|U$ 的秩。
若 $T$ 是 $X$ 中的开子集，使 $\mathcal{B}|T$ 同构于 $\mathcal{O}_X^m|T$，
则由前述备注，对所有满足 $T\cap X_i\neq\varnothing$ 的指标 $i$，数 $k_i$ 都等于 $m$。
令 $U$ 为 $X$ 中由 $\mathcal{B}$ 在其邻域内为局部自由 $\mathcal{O}_X$-模的点组成的开集，
并令 $U_j$（$1\leq j\leq s$）为它的连通分支；这些分支在 $X$ 中是开集，且数目有限
（因为 $U$ 诺特）。令 $V_j$ 为 $X'$ 的闭预概形，即在开集 $g^{-1}(U_j)$ 上诱导的预概形的\emph{闭包}
（\egaxref{I.9.5.11-zh}{I, 9.5.11}）。由上面的结论，对于满足 $X_i\cap U_j\neq\varnothing$ 的所有指标 $i$，
秩 $k_i$ 都等于同一个整数 $m_j$；还应注意，同一个 $X_i$ 不可能遇到两个指标不同的 $U_j$。
令 $i_\lambda$ 为满足 $X_i\cap U=\varnothing$ 的指标 $i$。令 $k$ 为所有 $k_i$ 的乘积，
置 $n_i=k/k_i$，并按如下方式定义预概形 $X''$。
对于每个 $j$（$1\leq j\leq s$），取 $k/m_j$ 个同构于 $V_j$ 的预概形；对于每个 $\lambda$，
取 $k/k_{i_\lambda}$ 个同构于 $X'_{i_\lambda}$ 的预概形；$X''$ 是所有这些预概形的\emph{不交并}。
定义态射 $g'':X''\to X'$，其在 $X''$ 的每个和项上的限制均为典范嵌入；显然，
$g''$ 是有限支配态射，因而是满射（因为有限态射是闭的
（\egaref{II.6.1.10-zh}{6.1.10}））。置 $g'=g\circ g''$；这是有限满射态射
$X''\to X$；且 ${g'}^*(\mathcal{L})={g''}^*(g^*(\mathcal{L}))$，故
${g'}^*(\mathcal{L})$ 是丰沛的 $\mathcal{O}_{X''}$-模层
（\egaref{II.5.1.12-zh}{5.1.12}）。显然，对于这个新的预概形 $X''$，
为 $X'$ 定义的那些 $k_i$ 所对应的秩都\emph{等于 $k$}；根据
（\egaxref{I.7.3.3-zh}{I, 7.3.3}），立即得到：对于 $X$ 的任意仿射开集 $T$，
\[
\bigl(g'_*(\mathcal{O}_{X''})
\otimes_{\mathcal{O}_X}\mathcal{R}(X)\bigr)|T
\]
是一个 $(\mathcal{R}(X)|T)$-模，同构于
$(\mathcal{R}(X)|T)^k$。

\begin{corollary}[6.6.12]
\label{II.6.6.12-zh}
若在 \egaref{II.6.6.11-zh}{6.6.11} 的表述中 $W=X$，则一个
$\mathcal{O}_X$-可逆模层 $\mathcal{L}$ 对 $f$ 丰沛，当且仅当
$g^*(\mathcal{L})$ 对 $f\circ g$ 丰沛。
\end{corollary}

\begin{remark}[6.6.13]
\label{II.6.6.13-zh}
在第三章中我们将看到：若 $Y$ 是诺特的，$f$ 为有限型，并且 $f$ 在 $X$ 的约化闭预概形上的限制是\emph{紧合态射}，
而该闭预概形的底空间为 $X-W$，则 \egaref{II.6.6.12-zh}{6.6.12} 的结论仍成立。
但是，在第五章中我们将给出域 $K$ 上代数概形 $X$ 的例子（结构态射
$X\to\operatorname{Spec}(K)$ 不是紧合态射），其正规化 $X'$ 是拟仿射的，但并非拟仿射
（因此 $\mathcal{O}_X$ 不丰沛，尽管 $\mathcal{O}_{X'}$ 丰沛
（\egaref{II.5.1.12-zh}{5.1.12}），且态射 $X'\to X$ 是有限满射
（\egaref{II.6.3.10-zh}{6.3.10}））。下一号将看到，将“拟仿射”替换为“仿射”时，
这种情形不可能发生。
\end{remark}
% Locked French authority: ega2-1-fr.tex, 820505 B,
% SHA-256 84EDBE3E83530AF2959B441796337C9DC21EAFCA6A13114A26778760FBF437AC.
% Sequential source scope: lines 12070-12155, subsection II.6.7 opening and
% Theorem II.6.7.1 through the reduction to integral X and Y / printed p.136.

\subsection{Chevalley 定理。}
\label{subsection:II.6.7-zh}

我们将借助 Serre 判别法（\egaref{II.5.2.1-zh}{5.2.1}）证明下面的定理；
在代数概形的情形，C.~Chevalley 已用其他方法证明了它。

\oldpage[II]{136}

\begin{theorem}[6.7.1]
\label{II.6.7.1-zh}
设 $X$ 是仿射概形，$Y$ 是诺特预概形，$f:X\to Y$ 是有限满射态射。
则 $Y$ 是仿射概形。
\end{theorem}

显然，$f_{\mathrm{red}}:X_{\mathrm{red}}\to Y_{\mathrm{red}}$ 是有限的
（\egaref{II.6.1.5-zh}{6.1.5, (vi)}）。由于 $X_{\mathrm{red}}$ 是
仿射概形，又由于 $Y$ 是诺特的，说 $Y$ 仿射或说 $Y_{\mathrm{red}}$
仿射是等价的（\egaxref{I.6.1.7-zh}{I, 6.1.7}），故可假设 $Y$
是\emph{约化的}。对 $Y$ 的每个闭子集 $Y'$，存在唯一一个底空间为
$Y'$ 的 $Y$ 的约化子预概形（\egaxref{I.5.1.2-zh}{I, 5.1.2}）。
它的逆像 $f^{-1}(Y')$ 典范同构于 $X\times_Y Y'$
（\egaxref{I.4.4.1-zh}{I, 4.4.1}），并且作为 $X$ 的闭子预概形是
仿射的；$f$ 在 $f^{-1}(Y')$ 上的限制可与
$f\times_Y 1_{Y'}$ 等同，并且是有限满射态射
（\egaref{II.6.1.5-zh}{6.1.5, (iii)}）。因此，根据诺特归纳原理
（\egaxref{0.2.2.2-zh}{0, 2.2.2}），并考虑到
\egaxref{I.6.1.7-zh}{I, 6.1.7}，只需在下述假设下证明该定理：
对于\emph{每个闭子集 $Y'\neq Y$}，底空间为 $Y'$ 的 $Y$ 的每个
闭子预概形都是仿射的。由此可知，\emph{对于每个支集（闭集）
$Z$ 不同于 $Y$ 的凝聚 $\mathcal{O}_Y$-模 $\mathcal{F}$，都有
$\mathrm{H}^1(Y,\mathcal{F})=0$}。事实上，存在一个底空间为 $Z$
的 $Y$ 的闭子预概形 $Y'$，使得若 $j:Y'\to Y$ 是典范单射，则
$\mathcal{F}=j_*(j^*(\mathcal{F}))$
（\egaxref{I.9.3.5-zh}{I, 9.3.5}）；于是由
\egaref{II.5.2.3-zh}{5.2.3} 以及
\egaxref{I.5.1.9.2-zh}{I, 5.1.9.2}，有
$\mathrm{H}^1(Y,\mathcal{F})
=\mathrm{H}^1(Y',j^*(\mathcal{F}))=0$。

先假设 $Y$ 可约，并设 $Y'$ 是 $Y$ 的一个不可约分支，
$Y''=Y-Y'$；仍以 $Y'$ 表示 $Y$ 的约化闭子预概形，其底空间为
$Y'$，并以 $j$ 表示典范单射 $Y'\to Y$。设 $\mathcal{F}$ 是一个
凝聚 $\mathcal{O}_Y$-模，并考虑典范同态
\[
\rho:\mathcal{F}\longrightarrow
\mathcal{F}'=j_*(j^*(\mathcal{F}))
\]
（\egaxref{0.4.4.3-zh}{0, 4.4.3}）。由
\egaxref{0.5.3.10-zh}{0, 5.3.10} 和
\egaxref{0.5.3.12-zh}{0, 5.3.12}，$\mathcal{F}'$ 是凝聚
$\mathcal{O}_Y$-模，因为若以 $\mathcal{J}$ 表示定义子预概形
$Y'$ 的 $\mathcal{O}_Y$ 的理想层，则
$j_*(\mathcal{O}_{Y'})=\mathcal{O}_Y/\mathcal{J}$。因此，
$\mathcal{G}=\operatorname{Ker}\rho$ 和
$\mathcal{K}=\operatorname{Im}\rho$ 也都是凝聚
$\mathcal{O}_Y$-模（\egaxref{0.5.3.4-zh}{0, 5.3.4}）。而按定义，
在 $Y'$ 的泛点 $y$ 处，$\mathcal{F}'$ 的纤维
$\mathcal{F}'_y$ 等于 $\mathcal{F}_y$；这是因为 $y$ 是 $Y'$
的内点，并且由于 $Y$ 约化，有 $\mathcal{J}_y=0$。由此可知，
$y$ 不属于 $\mathcal{G}$ 的支集（闭集）；另一方面，
$\mathcal{F}'$ 的支集（因而更有理由，$\mathcal{K}$ 的支集）包含于
$Y'$ 中。换言之，$\mathcal{G}$ 与 $\mathcal{K}$ 的支集都
\emph{不同于 $Y$}。于是
$\mathrm{H}^1(Y,\mathcal{G})=\mathrm{H}^1(Y,\mathcal{K})=0$，
而把上同调长正合列应用于正合列
$0\to\mathcal{G}\to\mathcal{F}\to\mathcal{K}\to0$，便得到
$\mathrm{H}^1(Y,\mathcal{F})=0$。最后由 Serre 判别法
（\egaref{II.5.2.1-zh}{5.2.1}）得出结论。

因此可假设 $Y$ 不可约，从而是\emph{整的}。还可假设 $X$
也是\emph{整的}：事实上，记 $X_i$ 为 $X$ 的约化闭子预概形，
其底空间是 $X$ 的不可约分支（\egaxref{I.5.2.1-zh}{I, 5.2.1}），
并记 $f_i$ 为 $f$ 在 $X_i$ 上的限制。至少有一个 $f_i$ 是支配态射；
又因它是有限态射（\egaref{II.6.1.5-zh}{6.1.5}），故为满射
（\egaref{II.6.1.10-zh}{6.1.10}）。另一方面，$X_i$ 是仿射概形，
所以在定理的表述中可以用 $X_i$ 代替 $X$。
% Locked French authority: ega2-1-fr.tex, 820505 B,
% SHA-256 84EDBE3E83530AF2959B441796337C9DC21EAFCA6A13114A26778760FBF437AC.
% Sequential source scope: lines 12156-12217, Lemma II.6.7.1.1 and proof /
% printed p.137. Printed mathematical and cross-reference anomalies are retained.

\begin{lemma}[6.7.1.1]
\label{II.6.7.1.1-zh}
设 $X$、$Y$ 是两个整诺特预概形，$x$（相应地，$y$）是 $X$
（相应地，$Y$）的泛点，$f:X\to Y$ 是有限满射态射。设
$\mathcal{L}$ 是一个可逆 $\mathcal{O}_X$-模层，并且存在 $y$ 的一个
仿射开邻域 $U$ 和一个截面 $g\in\Gamma(X,\mathcal{L})$，使得

\oldpage[II]{137}

$x\in X_g\subset f^{-1}(U)$。则存在两个整数 $m>0$、$n>0$，一个同态
$u:\mathcal{O}_Y^m\to f_*(\mathcal{L}^{\otimes n})$，以及 $y$ 的一个
开邻域 $V$，使得限制 $u|V$ 是从 $\mathcal{O}_Y^m|V$ 到
$f_*(\mathcal{L}^{\otimes n})|V$ 的同构。
\end{lemma}

设 $C$ 是 $U$ 的（整）环，$k=\mathcal{O}_y$ 是它的分式域。由于 $f$
是有限态射，$U'=f^{-1}(U)$ 是仿射的
（\egaref{II.1.3.2-zh}{1.3.2}）；设 $D$ 是它的（整）环，其分式域为
$K=\mathcal{O}_x$。按假设，$D$ 是有限型 $C$-模
（\egaref{II.6.1.4-zh}{6.1.4}），因而 $K$ 是 $k$ 的有限秩扩张。纤维
\[
f^{-1}(y)=X\times_Y\operatorname{Spec}(k(y))
=X\times_Y\operatorname{Spec}(\mathcal{O}_y)
\]
可与 $\operatorname{Spec}(K)$ 等同
（\egaxref{I.3.6.5-zh}{I, 3.6.6}）；取 $D$ 中的元素 $s_i$
（$1\leq i\leq m$），使其构成 $K$ 在 $k$ 上的一组基。存在 $n>0$，
使得 $X_g$ 上 $\mathcal{L}^{\otimes n}$ 的截面
$(s_i|X_g)g^{\otimes n}$ 可延拓为 $X$ 上 $\mathcal{L}^{\otimes n}$
的截面 $b_i$（$1\leq i\leq m$）
（\egaxref{I.9.3.1-zh}{I, 9.3.1}）。按定义，各 $b_i$ 也都是
$f_*(\mathcal{L}^{\otimes n})$ 在 $Y$ 上的截面，因而定义了同态
$u:\mathcal{O}_Y^m\to f_*(\mathcal{L}^{\otimes n})$
（\egaxref{0.5.1.1-zh}{0, 5.1.1}）；我们将证明 $u$ 满足要求。
有 $\mathcal{L}^{\otimes n}|U'=\widetilde{M}$，其中 $M$ 是有限型
$D$-模；若 $\varphi$ 是与 $f$ 的限制态射 $f^{-1}(U)\to U$ 对应的
单射 $C\to D$，则 $M_{[\varphi]}$ 是有限型 $C$-模。因此
\[
f_*(\mathcal{L}^{\otimes n})|U=(M_{[\varphi]})^\sim
\]
（\egaxref{I.1.6.3-zh}{I, 1.6.3}）是相干的；由于 $U$ 是 $Y$ 的任意
仿射开集，$f_*(\mathcal{L}^{\otimes n})$ 是相干的。此外，
$u|U=\widetilde{\theta}$，其中 $\theta$ 是一个 $C$-模同态
$C^m\to M_{[\varphi]}$，且 $u_y=\theta_y$ 是同态
\[
\theta\otimes1:K^m=C^m\otimes K\longrightarrow
M_{[\varphi]}\otimes K,
\]
而后一个同态按定义是一个\emph{同构}，因为各 $(b_i)_x$ 构成
$M_{[\varphi]}\otimes K$ 在 $k$ 上的一组基，并且按假设，$g_x$ 是
$(\mathcal{L}^{\otimes n})_x$ 的一个生成元。于是
$\mathcal{O}_Y$-模层 $\operatorname{Ker}u$ 和
$\operatorname{Coker}u$ 的支集都不包含 $y$；由于这些
$\mathcal{O}_Y$-模层是相干的
（\egaxref{0.5.3.3-zh}{0, 5.3.3}），其支集是闭的
（\egaxref{0.5.2.2-zh}{0, 5.2.2}），引理得证。

在此情形下，取 $\mathcal{L}=\mathcal{O}_X$，引理
（\egaref{II.6.7.1.1-zh}{6.7.1.1}）的假设便得到满足，因为 $X$ 是仿射的
（\egaxref{I.1.1.10-zh}{I, 1.1.10}）；我们置
$\mathcal{A}=\mathcal{O}_Y$、$\mathcal{B}=f_*(\mathcal{O}_X)$。由 Serre 判别法
（\egaref{II.5.2.1.1-zh}{5.2.1.1}），只需证明：对于任意凝聚
$\mathcal{O}_Y$-模层 $\mathcal{F}$，都有
$\mathrm{H}^1(Y,\mathcal{F})=0$；甚至只需在
$\mathcal{F}\subset\mathcal{O}_Y$ 时证明这一点，而由于 $Y$ 是整的，
这蕴含 $\mathcal{F}$ 无挠；事实上，我们将证明，对于\emph{任意}凝聚且
\emph{无挠}的 $\mathcal{O}_Y$-模层 $\mathcal{F}$，都有
$\mathrm{H}^1(Y,\mathcal{F})=0$。同态
$u:\mathcal{A}^m\to\mathcal{B}$ 定义了同态
\[
v:\mathcal{G}
=\mathcal{H}om_{\mathcal{A}}(\mathcal{B},\mathcal{F})
\longrightarrow
\mathcal{H}om_{\mathcal{A}}(\mathcal{A}^m,\mathcal{F})
=\mathcal{F}^m.
\]
先证明 $v$ 是\emph{单射}：按假设，
$\mathcal{T}=\operatorname{Coker}u$ 的支集与 $V$ 不相交，因而
是一个挠 $\mathcal{O}_Y$-模层
（\egaxref{I.7.4.6-zh}{I, 7.4.6}）；正合列
\[
\mathcal{A}^m\longrightarrow\mathcal{B}\longrightarrow\mathcal{T}
\longrightarrow0
\]
由函子 $\mathcal{H}om_{\mathcal{A}}$ 的左正合性给出正合列
\[
0\longrightarrow
\mathcal{H}om_{\mathcal{A}}(\mathcal{T},\mathcal{F})
\longrightarrow\mathcal{G}\xrightarrow{v}\mathcal{F}^m.
\]
但由于 $\mathcal{F}$ 无挠，有
$\mathcal{H}om_{\mathcal{A}}(\mathcal{T},\mathcal{F})=0$
（\egaxref{0.5.2.6-zh}{0, 5.2.6}），故断言成立。因此有正合列
\[
0\longrightarrow\mathcal{G}\longrightarrow\mathcal{F}^m
\longrightarrow\operatorname{Coker}v\longrightarrow0
\]

\oldpage[II]{138}

其中 $\mathcal{G}$ 和 $\operatorname{Coker}v$ 都是凝聚
$\mathcal{O}_Y$-模层
（\egaxref{0.5.3.4-zh}{0, 5.3.4} 和
\egaxref{0.5.3.5-zh}{5.3.5}）；由上同调正合列，只需证明
$\mathrm{H}^1(Y,\mathcal{G})
=\mathrm{H}^1(Y,\operatorname{Coker}v)=0$，便可推出
$\mathrm{H}^1(Y,\mathcal{F}^m)
=(\mathrm{H}^1(Y,\mathcal{F}))^m=0$，进而推出
$\mathrm{H}^1(Y,\mathcal{F})=0$。然而，限制 $v|V$ 是同构，
所以 $\operatorname{Coker}v$ 的支集不等于 $Y$；由开头的假设，
$\mathrm{H}^1(Y,\operatorname{Coker}v)=0$。另一方面，
$\mathcal{G}$ 是凝聚的 $\mathcal{B}$-模层
（\egaxref{I.9.6.4-zh}{I, 9.6.4}）；由于 $X$ 在 $Y$ 上是仿射的，
存在拟凝聚的 $\mathcal{O}_X$-模层 $\mathcal{K}$，使得
$\mathcal{G}$ 同构于 $f_*(\mathcal{K})$
（\egaref{II.1.4.3-zh}{1.4.3}）；由于 $X$ 是仿射的，有
$\mathrm{H}^1(X,\mathcal{K})=0$
（\egaxref{I.5.1.9.2-zh}{I, 5.1.9.2}），故由
（\egaref{II.5.2.3-zh}{5.2.3}）也有
$\mathrm{H}^1(Y,\mathcal{G})=0$，这就完成了定理
（\egaref{II.6.7.1-zh}{6.7.1}）的证明。

\begin{corollary}[6.7.2]
\label{II.6.7.2-zh}
设 $X$ 为诺特预概形，$(X_i)_{1\leq i\leq n}$ 是由闭集组成的空间
$X$ 的一个有限覆盖。则 $X$ 为仿射的，当且仅当对每个 $i$，存在
$X$ 的一个闭子预概形，其底空间为 $X_i$，且该闭子预概形为仿射的。
\end{corollary}

事实上，若条件成立，令 $X'$ 为各 $X_i$ 的\emph{不交并}概形；
显然 $X'$ 是仿射的，并可定义满射态射
$f:X'\to X$，使 $f$ 在 $X_i$ 上的限制为典范嵌入。由
（\egaref{II.6.7.1-zh}{6.7.1}），问题归结为验证 $f$ 是\emph{有限}态射，
而这已在（\egaref{II.6.1.5-zh}{6.1.5}）中证明。

\begin{corollary}[6.7.3]
\label{II.6.7.3-zh}
一个诺特预概形 $X$ 为仿射的，当且仅当以 $X$ 的各不可约分支为底空间的
约化闭子预概形均为仿射的。
\end{corollary}

\section{赋值判别法}
\label{section:II.7-zh}

本节给出态射的分离性与紧合性的赋值判别法，也就是说，这些判别法涉及形如
$\operatorname{Spec}(A)$ 的可变辅助概形，其中 $A$ 是赋值环。在适当的“诺特型”
假设下，这些判别法可以只考虑 $A$ 为\emph{离散}赋值环的情形。无疑，
这将是我们以后唯一需要应用的情形；而在一般情形下引入任意赋值环，
只是为了与经典理论建立联系。
\subsection{赋值环回顾。}
\label{subsection:II.7.1-zh}

\begin{env}[7.1.1]
\label{II.7.1.1-zh}
在刻画赋值环的各种等价性质中，我们将使用如下性质：若环 $A$ 是一个不是域的整环，
并且在包含于 $A$ 的分式域 $K$ 中且不同于 $K$ 的诸局部环所成的集合里，$A$ 关于支配关系是
\emph{极大的}（\egaxref{I.8.1.1-zh}{I, 8.1.1}），则称之为\emph{赋值环}。
回顾一下，赋值环是\emph{整闭的}。若 $A$ 是赋值环，则对于 $A$ 的每个非零素理想
$\mathfrak{p}\neq0$，$A_{\mathfrak{p}}$ 也是赋值环。
\end{env}

\begin{env}[7.1.2]
\label{II.7.1.2-zh}
设 $K$ 是一个域，$A$ 是 $K$ 的一个不是域的局部子环；

\oldpage[II]{139}

则存在一个以 $K$ 为分式域并支配 $A$ 的赋值环
（\egaxcite[第 1-07 页，引理 2]{I-1}）。

另一方面，设 $B$ 是赋值环，$k$ 是它的剩余域，$K$ 是它的分式域，$L$ 是 $k$ 的一个扩张。
则存在一个支配 $B$ 的\emph{完备}赋值环 $C$，其剩余域为 $L$。事实上，$L$ 是某个纯超越扩张
$L'=k(T_\mu)_{\mu\in M}$ 的代数扩张；已知可以把由 $B$ 所对应的 $K$ 的赋值延拓为
$K'=K(T_\mu)_{\mu\in M}$ 的一个赋值，使 $L'$ 成为该赋值的剩余域
（\cite[第 98 页]{II-24}）。用该延拓赋值的赋值环之完备化替换 $B$，即可归结到 $B$ 完备且
$L$ 是 $k$ 的代数闭包的情形。若 $\overline{K}$ 是 $K$ 的一个代数闭包，则可把定义 $B$ 的赋值
延拓到 $\overline{K}$；相应的剩余域是 $k$ 的一个代数闭包，这一点可由把 $k[T]$ 中首一多项式
的系数提升到 $\overline{K}$ 中看出。因此最终归结到 $L=k$ 的情形，此时取 $B$ 的完备化作为
$C$ 即可。
\end{env}

\begin{env}[7.1.3]
\label{II.7.1.3-zh}
设 $K$ 是一个域，$A$ 是 $K$ 的一个子环；$A$ 在 $K$ 中的整闭包 $A'$ 等于所有包含 $A$
且以 $K$ 为分式域的赋值环之交（\egaxcite[第 51 页，定理 2]{I-11}）。命题
\egaref{II.7.1.2-zh}{7.1.2} 可等价地表述为如下几何形式：
\end{env}
\begin{proposition}[7.1.4]
\label{II.7.1.4-zh}
设 $Y$ 是预概形，$p:X\to Y$ 是态射，$x$ 是 $X$ 的一点，$y=p(x)$，而 $y'\neq y$
是 $y$ 的一个特化（\egaxref{0.2.1.2-zh}{0, 2.1.2}）。则存在一个局部概形 $Y'$，它是
某个赋值环的谱，并存在一个分离态射 $f:Y'\to Y$，使得若以 $a$ 表示 $Y'$ 的唯一闭点，
以 $b$ 表示 $Y'$ 的泛点，则 $f(a)=y'$ 且 $f(b)=y$。此外，还可假设下列两个附加性质之一成立：
\begin{enumerate}
\item[\rm{(i)}] $Y'$ 是某个\emph{完备}赋值环的谱，该赋值环的剩余域是代数闭域，并且存在一个
$k(y)$-同态 $k(x)\to k(b)$。
\item[\rm{(ii)}] 存在一个 $k(y)$-同构
$k(x)\xrightarrow{\sim}k(b)$。
\end{enumerate}
\end{proposition}

令 $Y_1$ 为 $Y$ 的约化闭预概形，其底空间为 $\overline{\{y\}}$
（\egaxref{I.5.2.1-zh}{I, 5.2.1}），并令 $X_1$ 为逆像 $p^{-1}(Y_1)$ 所定义的闭预概形；
按假设 $y'\in\overline{\{y\}}$，而 $k(x)$ 在 $X$ 与 $X_1$ 中相同，故可假设 $Y$ 为\emph{整的}，
其泛点为 $y$。于是 $\mathcal{O}_{y'}$ 是一个不是域的整局部环，其分式域为
$\mathcal{O}_y=k(y)$，且 $k(x)$ 是 $k(y)$ 的扩张。为实现条件 $f(a)=y'$、$f(b)=y$ 以及
附加条件 (i)（相应地 (ii)），取 $Y'=\operatorname{Spec}(A')$，其中 $A'$ 是一个支配
$\mathcal{O}_{y'}$ 的赋值环：在 (i) 中，要求它完备且其剩余域为包含 $k(x)$ 的代数闭扩张；
在 (ii) 中，则取支配 $\mathcal{O}_{y'}$ 的赋值环 $A'$，并要求 $k(x)$ 是其分式域。$A'$ 的存在性由
\egaref{II.7.1.2-zh}{7.1.2} 保证。
\begin{env}[7.1.5]
\label{II.7.1.5-zh}
回顾一下，若局部环 $A$ 存在一个不同于极大理想 $\mathfrak{m}$ 的素理想，并且 $A$ 的每个
不同于 $\mathfrak{m}$ 的素理想都是\emph{极小}素理想，则称该局部环的\emph{维数为 $1$}；
当 $A$ 是\emph{整环}时，这等价于说 $\mathfrak{m}$ 与 $(0)$ 是仅有的素理想，并且
$\mathfrak{m}\neq(0)$。换言之，$Y=\operatorname{Spec}(A)$ 恰由两个

\oldpage[II]{140}

点 $a$、$b$ 组成：$a$ 是唯一的\emph{闭点}，$\mathfrak{j}_a=\mathfrak{m}$，且
$k(a)=k$ 是\emph{剩余域} $k=A/\mathfrak{m}$；$b$ 是 $Y$ 的\emph{泛点}，
$\mathfrak{j}_b=(0)$，集合 $\{b\}$ 是 $Y$ 中除 $\varnothing$ 与 $Y$ 以外唯一的开集
（因而处处\emph{稠密}），而 $k(b)=K$ 是 $A$ 的\emph{分式域}。
\end{env}

\begin{env}[7.1.6]
\label{II.7.1.6-zh}
对于一个诺特且维数为 $1$ 的局部环 $A$，已知下列条件等价
（\egaxcite[第 2-08 与 17-01 页]{I-1}）：
\begin{enumerate}
\item[\rm{a)}] $A$ 正规；
\item[\rm{b)}] $A$ 正则；
\item[\rm{c)}] $A$ 是赋值环。
\end{enumerate}
此外，此时 $A$ 是\emph{离散赋值环}。命题 \egaref{II.7.1.2-zh}{7.1.2} 与
\egaref{II.7.1.3-zh}{7.1.3} 对离散赋值环有如下类似结果：
\end{env}
% Locked French authority: ega2-1-fr.tex, 820505 B,
% SHA-256 84EDBE3E83530AF2959B441796337C9DC21EAFCA6A13114A26778760FBF437AC.
% Sequential source scope: lines 12446-12481, Proposition II.7.1.7 and proof.

\begin{proposition}[7.1.7]
\label{II.7.1.7-zh}
设 $A$ 是一个不是域的诺特整局部环，$K$ 是它的分式域，$L$ 是 $K$
的一个有限生成扩张。则存在一个以 $L$ 为分式域并支配 $A$ 的离散赋值环。
\end{proposition}

先假设 $L=K$。设 $\mathfrak{m}$ 是 $A$ 的极大理想，
$(x_1,\ldots,x_n)$ 是 $\mathfrak{m}$ 的一组非零生成元，并设 $B$ 是 $K$
的子环 $A[x_2/x_1,\ldots,x_n/x_1]$；该环是诺特的。立即可见，$B$ 的理想
$\mathfrak{m}B$ 等于主理想 $x_1B$。若 $\mathfrak{p}$ 是 $x_1B$ 的一个极小
素理想，则已知 $\mathfrak{p}$ 的高度为 $1$
（\egaxcite[t.~I, p.~277]{I-13}）；换言之，$B_{\mathfrak{p}}$ 是一个
\emph{维数为 $1$}的诺特局部环。显然，
$\mathfrak{p}B_{\mathfrak{p}}\cap A$ 是 $A$ 的一个理想，它包含
$\mathfrak{m}$ 而不包含 $1$，故等于 $\mathfrak{m}$；因此
$B_{\mathfrak{p}}$ \emph{支配} $A$
（\egaxref{I.8.1.1-zh}{I, 8.1.1}）。由 Krull--Akizuki 定理
（\cite[p.~293]{II-25}），$B_{\mathfrak{p}}$ 的整闭包 $C$ 是诺特环
（尽管 $C$ 不一定是有限型 $B_{\mathfrak{p}}$-模）。若 $\mathfrak{n}$
是 $C$ 的一个极大理想，则 $C_{\mathfrak{n}}$ 是维数为 $1$ 的正规诺特局部环
（\cite[p.~295]{II-25}），因而是一个离散赋值环；它支配
$B_{\mathfrak{p}}$，从而当然也支配 $A$。

现在设 $L$ 是 $K$ 的一个有限生成扩张。根据以上所述，可归结到 $A$ 已经是
离散赋值环的情形。设 $w$ 是与 $A$ 相伴的 $K$ 的一个赋值；存在 $L$ 的一个
离散赋值 $w'$ \emph{延拓} $w$：事实上，对 $L$ 的生成元数目作归纳，可归结到
$L=K(\alpha)$ 的情形，而此时该命题是经典结果
（\cite[p.~106]{II-24}）。
% Locked French authority: ega2-1-fr.tex, 820505 B,
% SHA-256 84EDBE3E83530AF2959B441796337C9DC21EAFCA6A13114A26778760FBF437AC.
% Sequential source scope: lines 12482-12506, Corollary II.7.1.8 and proof.

\begin{corollary}[7.1.8]
\label{II.7.1.8-zh}
设 $A$ 是一个诺特整环，$K$ 是它的分式域，$L$ 是 $K$ 的一个有限型
扩张。则 $A$ 在 $L$ 中的整闭包，等于所有以 $L$ 为分式域且包含
$A$ 的离散赋值环之交。
\end{corollary}

事实上，任一这样的离散赋值环都是正规环，因而\emph{更不待言}，它包含
$L$ 中每个在 $A$ 上整的元素。因此只需证明：若 $x\in L$ 在 $A$ 上
不是整的，则存在一个以 $L$ 为分式域、包含 $A$ 而不包含 $x$ 的离散
赋值环 $C$。关于 $x$ 的假设实际上意味着
$x\notin B=A[1/x]$；换言之，$1/x$ 在诺特环 $B$ 中不是可逆元。因此，
$B$ 有一个包含 $1/x$ 的素理想 $\mathfrak{p}$。局部整环
$B_{\mathfrak{p}}$ 是诺特的并且包含在 $L$ 中，而后者是
$B_{\mathfrak{p}}$ 的分式域的有限型扩张（该分式域包含 $K$）。由
（\egaref{II.7.1.7-zh}{7.1.7}），遂有一个离散赋值环 $C$，它支配
$B_{\mathfrak{p}}$ 并以 $L$ 为分式域；由于
$1/x\in\mathfrak{p}B_{\mathfrak{p}}$ 属于 $C$ 的极大理想，故
$x\notin C$，从而证明完毕。

（\egaref{II.7.1.7-zh}{7.1.7}）的几何形式如下：
\oldpage[II]{141}
\begin{proposition}[7.1.9]
\label{II.7.1.9-zh}
设 $Y$ 是局部诺特预概形，$p:X\to Y$ 是局部有限型态射，$x$ 是 $X$ 的一点，
$y=p(x)$，$y'\neq y$ 是 $y$ 的一个特化。则存在一个局部概形 $Y'$，它是某个
离散赋值环的谱，并存在一个分离态射 $f:Y'\to Y$ 以及一个从 $Y'$ 到 $X$ 的
$Y$-有理映射 $g$，使得若以 $a$ 表示 $Y'$ 的闭点，以 $b$ 表示它的泛点，则有
$f(a)=y'$、$f(b)=y$、$g(b)=x$，并且在交换图
\[
\tag{7.1.9.1}
\label{II.7.1.9.1-zh}
\begin{tikzcd}[column sep=large,row sep=large]
& k(x) \arrow[dl,"\gamma"'] \\
k(b) & k(y) \arrow[u,"\pi"] \arrow[l,"\varphi"']
\end{tikzcd}
\]
中（其中 $\pi$、$\varphi$、$\gamma$ 分别是对应于 $p$、$f$、$g$ 的同态），
$\gamma$ 是双射。
\end{proposition}

如同（\egaref{II.7.1.4-zh}{7.1.4}）中那样，考虑到
（\egaxref{I.6.3.4-zh}{I, 6.3.4, (iv)}），可以归结到 $Y$ 是以 $y$ 为泛点的
整预概形的情形；又由于问题对 $X$ 和 $Y$ 都是局部的，可以假设 $p$ 是有限型态射。
于是便处于（\egaref{II.7.1.4-zh}{7.1.4}）的情形，并且还具有附加性质：
$k(x)$ 是 $k(y)$ 的\emph{有限型}扩张
（\egaxref{I.6.4.11-zh}{I, 6.4.11}），且 $\mathcal{O}_{y'}$ 是诺特的；
因此可以应用（\egaref{II.7.1.7-zh}{7.1.7}），取
$Y'=\operatorname{Spec}(A')$，其中 $A'$ 是支配 $\mathcal{O}_{y'}$ 的离散赋值环，
且其分式域为 $k(x)$。这样就定义了交换图
（\egaref{II.7.1.9.1-zh}{7.1.9.1}），其中 $\gamma$ 是双射，而 $\pi$ 和
$\varphi$ 分别对应于态射 $p$ 和 $f$。此外，由于 $X$ 和 $Y$ 都是局部诺特的
（\egaxref{I.6.6.2-zh}{I, 6.6.2}），并且 $Y'$ 是整的，存在唯一一个从 $Y'$ 到
$X$ 的 $Y$-有理映射 $g$ 与同构 $\gamma$ 对应
（\egaxref{I.7.1.15-zh}{I, 7.1.15}），这就完成了证明。

\subsection{分离性的赋值判别法。}
\label{subsection:II.7.2-zh}

\begin{proposition}[7.2.1]
\label{II.7.2.1-zh}
设 $X$、$Y$ 是两个预概形，$f:X\to Y$ 是一个拟紧态射。下列两个条件等价：
\begin{enumerate}
\item[\rm a)] 态射 $f$ 是闭的。
\item[\rm b)] 对于每个 $x\in X$ 以及 $y=f(x)$ 的每个不同于 $y$ 的特化 $y'$，
存在 $x$ 的一个特化 $x'$，使得 $f(x')=x$。
\end{enumerate}
\end{proposition}

条件 b) 表明 $f(\overline{\{x\}})=\overline{\{y\}}$，因而它是 a) 的推论。为了证明 b) 蕴含
a)，考虑底空间 $X$ 的一个闭子集 $X'$；令 $Y'=\overline{f(X')}$，并证明
$Y'=f(X')$。分别考虑 $X$ 与 $Y$ 中以 $X'$ 与 $Y'$ 为底空间的约化闭预概形
（\egaxref{I.5.2.1-zh}{I, 5.2.1}）；于是存在一个态射 $f':X'\to Y'$，使图表
\[
\begin{tikzcd}[column sep=large,row sep=large]
X' \arrow[r,"f'"] \arrow[d] & Y' \arrow[d] \\
X \arrow[r,"f"'] & Y
\end{tikzcd}
\]
交换（\egaxref{I.5.2.2-zh}{I, 5.2.2}），并且由于 $f$ 是拟紧的，$f'$ 也是拟紧的。因此，
问题归结为证明：若 $f$ 是一个拟紧且\emph{占优}的态射，则

\oldpage[II]{142}

条件 b) 蕴含 $f(X)=Y$。事实上，设 $y'$ 是 $Y$ 的一个点，$y$ 是 $Y$ 的一个包含 $y'$ 的
不可约分支的泛点；依条件 b)，只须证明 $f^{-1}(y)$ 非空。但已知这一性质是 $f$ 占优且拟紧的
推论（\egaxref{I.6.6.5-zh}{I, 6.6.5}）。

\begin{corollary}[7.2.2]
\label{II.7.2.2-zh}
设 $f:X\to Y$ 是一个拟紧浸入态射。底空间 $X$ 在 $Y$ 中为闭集，当且仅当对于它的每个点，
它都包含该点在 $Y$ 中的每个特化。
\end{corollary}
% Locked French authority: ega2-1-fr.tex, 820505 B,
% SHA-256 84EDBE3E83530AF2959B441796337C9DC21EAFCA6A13114A26778760FBF437AC.
% Sequential source scope: lines 12592-12634, Proposition II.7.2.3 and proof.

\begin{proposition}[7.2.3]
\label{II.7.2.3-zh}
设 $Y$ 是预概形（相应地，局部诺特预概形），$f:X\to Y$ 是态射
（相应地，局部有限型态射）。下列条件等价：
\begin{enumerate}
\item[\rm a)] $f$ 是分离的。
\item[\rm b)] 对角态射 $X\to X\times_YX$ 是拟紧的，并且对每个形如
$Y'=\operatorname{Spec}(A)$ 的 $Y$-预概形，其中 $A$ 是赋值环
（相应地，离散赋值环），从 $Y'$ 到 $X$ 的两个 $Y$-态射若在 $Y'$
的泛点处相同，则它们相等。
\item[\rm c)] 对角态射 $X\to X\times_YX$ 是拟紧的，并且对每个形如
$Y'=\operatorname{Spec}(A)$ 的 $Y$-预概形，其中 $A$ 是赋值环
（相应地，离散赋值环），$X'=X_{(Y')}$ 的两个 $Y'$-截面若在 $Y'$
的泛点处相同，则它们相等。
\end{enumerate}
\end{proposition}

b) 与 c) 的等价性来自从 $Y'$ 到 $X$ 的 $Y$-态射与 $X'$ 的
$Y'$-截面之间的一一对应（\egaxref{I.3.3.14-zh}{I, 3.3.14}）。
若 $X$ 在 $Y$ 上分离，则由于 $Y'$ 是整的，根据
\egaxref{I.7.2.2.1-zh}{I, 7.2.2.1}，条件 b) 成立。还需证明 b)
蕴含对角态射 $\Delta:X\to X\times_YX$ 是闭态射；这等价于证明它满足
\egaref{II.7.2.2-zh}{7.2.2} 的判别法。设 $z$ 是对角集
$\Delta(X)$ 的一个点，而 $z'\neq z$ 是 $z$ 在 $X\times_YX$ 中的一个
特化。于是根据 \egaref{II.7.1.4-zh}{7.1.4}，存在一个赋值环 $A$
以及一个从 $Y'=\operatorname{Spec}(A)$ 到 $X\times_YX$ 的态射 $f$，
它把 $Y'$ 的闭点 $a$ 映到 $z'$，把 $Y'$ 的泛点 $b$ 映到 $z$；该态射
使 $Y'$ 成为一个 $(X\times_YX)$-预概形，因而当然也成为一个
$Y$-预概形。将 $f$ 与 $X\times_YX$ 的两个投影复合，得到从 $Y'$ 到
$X$ 的两个 $Y$-态射 $g_1$、$g_2$；依假设，它们在点 $b$ 处相同，
故二者都等于同一个态射 $g$。这意味着
（\egaxref{I.5.3.1-zh}{I, 5.3.1}）$f$ 分解为 $f=\Delta\circ g$，
因而 $z'\in\Delta(X)$。当假设 $Y$ 局部诺特且 $f$ 局部有限型时，
$X\times_YX$ 是局部诺特的（\egaxref{I.6.6.7-zh}{I, 6.6.7}）；
于是根据 \egaref{II.7.1.9-zh}{7.1.9}，可重复以上论证，并在其中假设
$A$ 是离散赋值环。

\begin{remark}[7.2.4]
\label{II.7.2.4-zh}
\begin{enumerate}
\item[\rm (i)] 当 $Y$ 局部诺特且 $f$ 局部有限型时，态射 $\Delta$ 拟紧这一假设
总是成立，因为此时 $X\times_YX$ 局部诺特
（\egaxref{I.6.6.4-zh}{I, 6.6.4, (i)}）。在一般情形下，该假设还表示：
对于 $X$ 的任意仿射开集覆盖 $(U_\alpha)$，各集合
$U_\alpha\cap U_\beta$ 都是\emph{拟紧的}。
\item[\rm (ii)] 为使 $f$ 成为分离态射，只须在赋值环 $A$ \emph{完备}且其剩余域
\emph{代数闭}的情形下验证条件 b) 或条件 c)；这确实由
（\egaref{II.7.2.3-zh}{7.2.3}）和
（\egaref{II.7.1.4-zh}{7.1.4}）的证明得出。
\end{enumerate}
\end{remark}

\oldpage[II]{143}
% Locked French authority: ega2-1-fr.tex, 820505 B,
% SHA-256 84EDBE3E83530AF2959B441796337C9DC21EAFCA6A13114A26778760FBF437AC.
% Sequential source scope: lines 12654-12712, subsection II.7.3 opening
% through Environment II.7.3.2.

\subsection{紧合性的赋值判别法。}
\label{subsection:II.7.3-zh}

\begin{proposition}[7.3.1]
\label{II.7.3.1-zh}
设 $A$ 是赋值环，$Y=\operatorname{Spec}(A)$，$b$ 是 $Y$ 的泛点，
$X$ 是整概形，$f:X\to Y$ 是闭态射，并且 $f^{-1}(b)$ 约化为一个点
$x$，相应的同态 $k(b)\to k(x)$ 是双射。则 $f$ 是同构。
\end{proposition}

由于 $f$ 是闭的且为支配态射，有 $f(X)=Y$。只需
（\egaxref{I.4.2.2-zh}{I, 4.2.2}）证明：对 $Y$ 中每个
$y'\neq b$，存在\emph{唯一一个}点 $x'$，使得 $f(x')=y'$ 且相应的
同态 $\mathcal{O}_{y'}\to\mathcal{O}_{x'}$ 是双射，因为此时 $f$
将是同胚。事实上，若 $f(x')=y'$，则 $\mathcal{O}_{x'}$ 是包含于
$K=k(x)=k(b)$ 中并支配 $\mathcal{O}_{y'}$ 的一个局部环；后者是局部环
$A_{y'}$，因而是以 $K$ 为分式域的赋值环
（\egaref{II.7.1.1-zh}{7.1.1}）。另一方面，
$\mathcal{O}_{x'}\neq K$，因为 $x'$ 不是 $X$ 的泛点
（\egaxref{0.2.1.3-zh}{0, 2.1.3}）；由此得到
$\mathcal{O}_{x'}=\mathcal{O}_{y'}$。由于 $X$ 是整概形，关系
$\mathcal{O}_{x'}=\mathcal{O}_{x''}$ 蕴含 $x'=x''$
（\egaxref{I.8.2.2-zh}{I, 8.2.2}），证明完毕。

\begin{env}[7.3.2]
\label{II.7.3.2-zh}
设 $A$ 是赋值环，$Y=\operatorname{Spec}(A)$，$b$ 是 $Y$ 的泛点，
从而 $\mathcal{O}_b=k(b)$ 等于 $A$ 的分式域 $K$；再设 $f:X\to Y$
是态射。已知（\egaxref{I.7.1.4-zh}{I, 7.1.4}），$X$ 的
\emph{$Y$-有理截面}与点 $b$ 处的 $Y$-截面的\emph{芽}（这些截面
定义在 $b$ 的邻域中）一一对应，因而有典范映射
\[
\tag{7.3.2.1}
\label{II.7.3.2.1-zh}
\Gamma_{\mathrm{rat}}(X/Y)\longrightarrow
\Gamma(f^{-1}(b)/\operatorname{Spec}(K))
\]
而依定义（\egaxref{I.3.4.5-zh}{I, 3.4.5}），
$\Gamma(f^{-1}(b)/\operatorname{Spec}(K))$ 的元素可与
\emph{$f^{-1}(b)=X\otimes_AK$ 中在 $K$ 上有理的点}等同。当 $f$
\emph{分离}时，由 \egaxref{I.5.4.7-zh}{I, 5.4.7}，映射
\egaref{II.7.3.2.1-zh}{7.3.2.1} 是\emph{单射}，因为 $Y$ 是整概形。

把 \egaref{II.7.3.2.1-zh}{7.3.2.1} 与典范映射
$\Gamma(X/Y)\to\Gamma_{\mathrm{rat}}(X/Y)$
（\egaxref{I.7.1.2-zh}{I, 7.1.2}）复合，得到典范映射
\[
\tag{7.3.2.2}
\label{II.7.3.2.2-zh}
\Gamma(X/Y)\longrightarrow \Gamma(f^{-1}(b)/\operatorname{Spec}(K)).
\]
当 $f$ \emph{分离}时，该映射也仍为\emph{单射}
（\egaxref{I.5.4.7-zh}{I, 5.4.7}）。
\end{env}

% Locked French authority: ega2-1-fr.tex, 820505 B,
% SHA-256 84EDBE3E83530AF2959B441796337C9DC21EAFCA6A13114A26778760FBF437AC.
% Sequential source scope: lines 12713-12776, II.7.3.3 through II.7.3.5.

\begin{proposition}[7.3.3]
\label{II.7.3.3-zh}
设 $A$ 是一个分式域为 $K$ 的赋值环，$Y=\operatorname{Spec}(A)$，$b$
是 $Y$ 的泛点，$f:X\to Y$ 是一个\emph{分离}且\emph{闭}的态射。则典范映射
（\egaref{II.7.3.2.2-zh}{7.3.2.2}）是双射（这等价于说它是\emph{满射}，
并且蕴含：相对于 $Y$ 而言，$X$ 的有理截面均\emph{处处有定义}）。
\end{proposition}

于是，设 $x$ 是 $f^{-1}(b)$ 中一个\emph{在 $K$ 上有理}的点。由于 $f$
是分离态射，态射 $f^{-1}(b)\to\operatorname{Spec}(K)$（它与 $f$ 对应）
也是分离态射
（\egaxref{I.5.5.1-zh}{I, 5.5.1, (iv)}）；又因为 $f^{-1}(b)$ 的每个截面
都是闭浸入（\egaxref{I.5.4.6-zh}{I, 5.4.6}），故 $\{x\}$
\emph{在 $f^{-1}(b)$ 中}是闭的。考虑既约闭子预概形 $X'$；它是 $X$ 的
子预概形，其底空间为闭包 $\overline{\{x\}}$，即 $\{x\}$ 在 $X$ 中的闭包。显然，$f$ 到
$X'$ 的限制满足（\egaref{II.7.3.1-zh}{7.3.1}）的假设，因而是从 $X'$
到 $Y$ 的一个\emph{同构}；其逆同构便是所求的 $Y$-截面（取值于 $X$）。

\begin{env}[7.3.4]
\label{II.7.3.4-zh}
为了陈述以下两个结果，我们将使用一种会在第四章中得到说明和讨论的术语：若 $F$ 是
\oldpage[II]{144}
预概形 $Y$ 的一个子集，则所谓 $F$ 在 $Y$ 中的\emph{余维数}，记作
$\operatorname{codim}_YF$，是诸整数
$\dim(\mathcal{O}_z)$ 的下确界，其中 $z$ 遍历 $F$。
\end{env}

\begin{corollary}[7.3.5]
\label{II.7.3.5-zh}
设 $Y$ 是一个局部诺特既约预概形，$N$ 是由点 $y\in Y$ 组成的集合，
在这些点处 $Y$ 不正则（\egaxref{0.4.1.4-zh}{0, 4.1.4}）；假设
$\operatorname{codim}_YN\geq2$。设 $f:X\to Y$ 是一个有限型、分离且闭
的态射，$g$ 是相对于 $Y$ 的、取值于 $X$ 的一个有理截面；若 $Y'$ 是
$Y$ 中 $g$ 无定义的那些点的集合（该集合是\emph{闭的}
（\egaxref{I.7.2.1-zh}{I, 7.2.1}）），则
$\operatorname{codim}_YY'\geq2$。
\end{corollary}

只需证明 $g$ 在每个点 $z\in Y$ 都有定义，只要
$\dim\mathcal{O}_z\leq1$。若 $\dim\mathcal{O}_z=0$，则 $z$ 是 $Y$ 的某个不可约分支的泛点
（\egaxref{I.1.1.14-zh}{I, 1.1.14}），故属于 $Y$ 的每个处处稠密开集，
特别地属于 $g$ 的定义域。于是设 $\dim\mathcal{O}_z=1$；依假设，
$\mathcal{O}_z$ 是一个正则诺特局部环，故由
（\egaref{II.7.1.6-zh}{7.1.6}），它是一个\emph{离散赋值环}。设
$Z=\operatorname{Spec}(\mathcal{O}_z)$；由于 $U=Y-Y'$ 处处稠密，
$U\cap Z$ 非空（\egaxref{I.2.4.2-zh}{I, 2.4.2}）；设 $g'$ 是从 $Z$ 到
$X$ 的有理映射，它由 $g$ 诱导
（\egaxref{I.7.2.8-zh}{I, 7.2.8}）；只须证明 $g'$ 是一个\emph{态射}
（\egaxref{I.7.2.9-zh}{I, 7.2.9}）。但可以把 $g'$ 看成一个有理 $Z$-截面，
其所在的 $Z$-预概形是 $f^{-1}(Z)=X\times_YZ$；显然，态射
$f^{-1}(Z)\to Z$（它与 $f$ 对应）是闭态射，而由
（\egaxref{I.5.5.1-zh}{I, 5.5.1, (i)}）可知它是分离态射；于是由
（\egaref{II.7.3.3-zh}{7.3.3}）得出 $g'$ 处处有定义；由于 $Z$ 既约，
并且 $X$ 在 $Y$ 上分离，故 $g'$ 是一个态射
（\egaxref{I.7.2.2-zh}{I, 7.2.2}）。
% Locked French authority: ega2-1-fr.tex, 820505 B,
% SHA-256 84EDBE3E83530AF2959B441796337C9DC21EAFCA6A13114A26778760FBF437AC.
% Sequential source scope: lines 12777-12804, Corollary II.7.3.6,
% Remark II.7.3.7, and the transition to the next result.

\begin{corollary}[7.3.6]
\label{II.7.3.6-zh}
设 $S$ 是局部诺特预概形，$X$ 和 $Y$ 是两个 $S$-预概形；假设
$Y$ 是约化的，并且使 $Y$ 非正则的点 $y\in Y$ 所成的集合 $N$
满足
$\operatorname{codim}_YN\geq2$；最后，假设结构态射 $X\to S$ 是
紧合的。设 $f$ 是从 $Y$ 到 $X$ 的一个 $S$-有理映射，并设 $Y'$
为 $Y$ 中 $f$ 无定义的点所成的集合；则
$\operatorname{codim}_YY'\geq2$。
\end{corollary}

已知（\egaxref{I.7.1.2-zh}{I, 7.1.2}），可以把从 $Y$ 到 $X$ 的
$S$-有理映射与 $X\times_SY$ 的 $Y$-有理截面等同；由于结构态射
$X\times_SY\to Y$ 是闭的（\egaref{II.5.4.1-zh}{5.4.1}），可以应用
\egaref{II.7.3.5-zh}{7.3.5}，从而得到本推论。

\begin{remark}[7.3.7]
\label{II.7.3.7-zh}
特别地，当 $Y$ 是\emph{正规}的（\egaxref{0.4.1.4-zh}{0, 4.1.4}）时，
由 \egaref{II.7.1.6-zh}{7.1.6}，\egaref{II.7.3.5-zh}{7.3.5} 和
\egaref{II.7.3.6-zh}{7.3.6} 中对 $Y$ 所作的假设均成立。
\end{remark}

我们可以用 \egaref{II.7.3.3-zh}{7.3.3} 的一个逆命题来刻画普遍闭
（相应地，紧合）态射：
% Locked French authority: ega2-1-fr.tex, 820505 B,
% SHA-256 84EDBE3E83530AF2959B441796337C9DC21EAFCA6A13114A26778760FBF437AC.
% Sequential source scope: lines 12805-12867, Theorem II.7.3.8 and proof.

\begin{theorem}[7.3.8]
\label{II.7.3.8-zh}
设 $Y$ 是预概形（相应地，局部诺特预概形），$f:X\to Y$ 是拟紧分离态射
（相应地，有限型态射）。下列条件等价：
\begin{enumerate}
\item[\rm a)] $f$ 是普遍闭态射（相应地，紧合态射）。
\item[\rm b)] 对每个形如 $Y'=\operatorname{Spec}(A)$ 的 $Y$-概形，
其中 $A$ 是分式域为 $K$ 的赋值环（相应地，离散赋值环），与典范单射
$A\to K$ 对应的典范映射
\[
\operatorname{Hom}_Y(Y',X)\longrightarrow
\operatorname{Hom}_Y(\operatorname{Spec}(K),X)
\]
是满射（相应地，双射）。
\oldpage[II]{145}
\item[\rm c)] 对每个形如 $Y'=\operatorname{Spec}(A)$ 的 $Y$-概形，
其中 $A$ 是赋值环（相应地，离散赋值环），相对于 $Y'$-预概形
$X_{(Y')}$ 的典范映射 \egaref{II.7.3.2.2-zh}{7.3.2.2} 是满射
（相应地，双射）。
\end{enumerate}
\end{theorem}

b) 与 c) 的等价性立即由 \egaxref{I.3.3.14-zh}{I, 3.3.14} 得出；
a) 蕴含 c)，因为在两种假设中的任一种下，a) 都蕴含
$f_{(Y')}$ 是分离的（\egaxref{I.5.5.1-zh}{I, 5.5.1, (iv)}）且为
闭态射，故只需应用 \egaref{II.7.3.3-zh}{7.3.3}。还需证明 b)
蕴含 a)。先考虑 $Y$ 任意而 $f$ 分离且拟紧的情形。若 $f$ 满足条件
b)，则对任意 $Y$-预概形 $Y''$，
$f_{(Y'')}:X_{(Y'')}\to Y''$ 也满足该条件；这是因为 b) 与 c) 等价，
并且对每个态射 $Y'\to Y''$ 都有
$X_{(Y'')}\times_{Y''}Y'=X\times_YY'$
（\egaxref{I.3.3.9.1-zh}{I, 3.3.9.1}）。此外，当 $f$ 分离且拟紧时，
$f_{(Y'')}$ 也分离且拟紧
（\egaxref{I.5.5.1-zh}{I, 5.5.1, (iv)} 和
\egaxref{I.6.6.4-zh}{6.6.4, (iii)}），所以问题归结为证明 b) 蕴含
$f$ 是\emph{闭态射}。为此，只需验证
\egaref{II.7.2.1-zh}{7.2.1} 的条件 b)。设 $x\in X$，并设
$y'$ 是 $y=f(x)$ 的一个不同于 $y$ 的特化。由
\egaref{II.7.1.4-zh}{7.1.4}，存在一个概形 $Y'$，它是某个赋值环的谱，
以及一个分离态射 $g:Y'\to Y$，使得若 $a$ 表示 $Y'$ 的闭点，$b$
表示其泛点，则 $g(a)=y'$、$g(b)=y$，并且存在一个 $k(y)$-同态
$k(x)\to k(b)$。后者典范地对应于一个 $Y$-态射
$\operatorname{Spec}(k(b))\to X$
（\egaxref{I.2.4.6-zh}{I, 2.4.6}），故由 b)，存在一个 $Y$-态射
$h:Y'\to X$ 与前述态射相对应。于是 $h(b)=x$；若置 $h(a)=x'$，
则 $x'$ 是 $x$ 的一个特化，并且
$f(x')=f(h(a))=g(a)=y'$。

现在设 $Y$ 局部诺特且 $f$ 为有限型。首先，由
\egaref{II.7.2.3-zh}{7.2.3}，假设 b) 蕴含 $f$ 是\emph{分离的}，
因为对角态射 $X\to X\times_YX$ 是拟紧的
（\egaref{II.7.2.4-zh}{7.2.4}）。此外，考虑到
\egaref{II.5.6.3-zh}{5.6.3}，为了验证 $f$ 紧合，只需证明对每个
\emph{有限型} $Y$-预概形 $Y''$，
$f_{(Y'')}:X_{(Y'')}\to Y''$ 都是\emph{闭态射}。此时 $Y''$
局部诺特，故可重复第一种情形中的论证，其中取 $Y'$ 为离散赋值环的谱，
并应用 \egaref{II.7.1.9-zh}{7.1.9} 来代替
\egaref{II.7.1.4-zh}{7.1.4}。

% Locked French authority: ega2-1-fr.tex, 820505 B,
% SHA-256 84EDBE3E83530AF2959B441796337C9DC21EAFCA6A13114A26778760FBF437AC.
% Sequential source scope: lines 12868-12912, Remark II.7.3.9.

\begin{remark}[7.3.9]
\label{II.7.3.9-zh}
\begin{enumerate}
\item[\rm (i)] 当 $Y$ 是任意预概形且 $f$ 是分离态射时，为使 $f$
普遍闭，只需条件 b) 或条件 c) 对那些\emph{完备}且剩余域\emph{代数闭}
的赋值环 $A$ 成立；这事实上由上述证明和
\egaref{II.7.1.4-zh}{7.1.4} 得出。
\item[\rm (ii)] 由 \egaref{II.7.3.8-zh}{7.3.8} 的判别法 c)，可以得到
射影态射 $X\to Y$ 是闭态射（\egaref{II.5.5.3-zh}{5.5.3}）这一事实
的一个更接近经典方法的新证明。事实上，可以假设 $Y$ 是仿射的，
从而把 $X$ 等同于某个射影丛 $\mathbf{P}_Y^n$ 的闭子预概形
（\egaref{II.5.3.3-zh}{5.3.3}）；为证明 $X\to Y$ 是闭态射，只需验证
结构态射 $\mathbf{P}_Y^n\to Y$ 也是如此，而
\egaref{II.7.3.8-zh}{7.3.8} 的判别法 c) 与
\egaref{II.4.1.3.1-zh}{4.1.3.1} 共同表明，问题归结为证明下述事实：
\emph{若 $Y$ 是赋值环 $A$ 的谱，其分式域为 $K$，则
$\mathbf{P}_Y^n$ 的每个 $K$-值点，都来自 $\mathbf{P}_Y^n$ 的一个
$A$-值点（通过限制到 $Y$ 的泛点）}。然而，每个可逆
$\mathcal{O}_Y$-模都是平凡的（\egaxref{I.2.4.8-zh}{I, 2.4.8}）；
因此，由 \egaref{II.4.2.6-zh}{4.2.6}，$\mathbf{P}_Y^n$ 的一个
$K$-值点可以等同于 $K$ 中元素组 $(\zeta c_0,\zeta c_1,\ldots,\zeta
c_n)$ 所在的等价类，其中 $\zeta\neq0$，而 $c_i$ 是 $K$ 中不全为零
的元素。现在，用 $A$ 中一个具有
\oldpage[II]{146}
适当赋值的元素乘所有 $c_i$，便可以假设所有 $c_i$ 都属于 $A$，并且
其中至少一个可逆。但此时，由 \egaref{II.4.2.6-zh}{4.2.6}，
系 $(c_0,\ldots,c_n)$ 也定义了 $\mathbf{P}_Y^n$ 的一个 $A$-值点，
这就证明了我们的断言。
\item[\rm (iii)] 当把 $Y$-预概形 $X$ 的资料看作等价于在
$Y$-预概形 $Y'$ 上的函子
\[
X(Y')=\operatorname{Hom}_Y(Y',X)
\]
的资料时，判别法 \egaref{II.7.2.3-zh}{7.2.3} 和
\egaref{II.7.3.8-zh}{7.3.8} 尤为方便；例如，这些判别法将使我们能够
证明，在某些条件下，“Picard 概形”是紧合的。
\end{enumerate}
\end{remark}

% Locked French authority: ega2-1-fr.tex, 820505 B,
% SHA-256 84EDBE3E83530AF2959B441796337C9DC21EAFCA6A13114A26778760FBF437AC.
% Sequential source scope: lines 12913-12988, Corollary II.7.3.10 and proof.

\begin{corollary}[7.3.10]
\label{II.7.3.10-zh}
设 $Y$ 是整概形（相应地，局部诺特整概形），$X$ 是整概形，
$f:X\to Y$ 是支配态射。
\begin{enumerate}
\item[\rm (i)] 若 $f$ 拟紧且普遍闭，则任何以 $X$ 上的有理函数域
$R(X)$ 为分式域并支配 $Y$ 的某个局部环的赋值环，也支配 $X$ 的某个
局部环。
\item[\rm (ii)] 反过来，假设 $f$ 为有限型，并且每个以 $R(X)$ 为
分式域的赋值环（相应地，每个离散赋值环）都满足 (i) 中所述的性质。
则 $f$ 是紧合态射。
\end{enumerate}
\end{corollary}

首先注意，在所有情形下，这些假设都蕴含 $f$ 是分离的
（\egaxref{I.5.5.9-zh}{I, 5.5.9}）。
\begin{enumerate}
\item[\rm (i)] 设 $K=R(Y)$、$L=R(X)$，$y$ 是 $Y$ 的一个点，$A$
是一个以 $L$ 为分式域并支配 $\mathcal{O}_y$ 的赋值环。于是单射
$\mathcal{O}_y\to A$ 定义了一个从 $Y'=\operatorname{Spec}(A)$ 到
$Y$ 的态射 $h$（\egaxref{I.2.4.4-zh}{I, 2.4.4}），使得
$h(a)=y$，其中 $a$ 表示 $Y'$ 的闭点。此外，若 $\eta$ 是 $Y$ 的
泛点（它也是 $\operatorname{Spec}(\mathcal{O}_y)$ 的泛点），并以
$b$ 表示 $Y'$ 的泛点，则 $h(b)=\eta$（因为依假设 $K\subset L$）。
若 $\xi$ 是 $X$ 的泛点，则依假设 $k(\xi)=k(b)=L$，从而有一个
$Y$-态射 $g:\operatorname{Spec}(L)\to X$，使得 $g(b)=\xi$；由
\egaref{II.7.3.8-zh}{7.3.8}，$g$ 来自一个 $Y$-态射
$g':Y'\to X$。若 $x=g'(a)$，则显然 $A$ 支配 $\mathcal{O}_x$。
\item[\rm (ii)] 由于问题在 $Y$ 上是局部的，总可以假设 $Y$ 仿射
（相应地，仿射且诺特）。因为 $f$ 是有限型的，在两种情形下都可以应用
Chow 引理（\egaref{II.5.6.1-zh}{5.6.1}）。因此，存在射影态射
$p:P\to Y$、浸入态射 $j:X'\to P$，以及射影、满射且双有理的态射
$g:X'\to X$（其中 $X'$ 是整的），使得图
\[
\begin{tikzcd}[column sep=large,row sep=large]
P \arrow[d,"p"'] & X' \arrow[l,"j"'] \arrow[d,"g"] \\
Y & X \arrow[l,"f"]
\end{tikzcd}
\]
交换。只需证明 $j$ 是\emph{闭}浸入；事实上，此时
$f\circ g=p\circ j$ 是射影态射，因而是紧合态射；又因 $g$ 是满射，
$f$ 也将是紧合态射（\egaref{II.5.4.3-zh}{5.4.3}）。设 $Z$ 是 $P$
的约化闭子预概形，其底空间为 $\overline{j(X')}$
（\egaxref{I.5.2.1-zh}{I, 5.2.1}）。由于 $X'$ 是整的，$j$ 分解为
$i\circ h$，其中 $i:Z\to P$ 是典范单射，$h:X'\to Z$ 是支配开浸入
（\egaxref{I.5.2.3-zh}{I, 5.2.3}），并且 $Z$ 是整的；

\oldpage[II]{147}
此外，$Z$ 在 $Y$ 上是射影的。由此可见，可以归结到 $P$ \emph{整}
且 $j$ \emph{支配并且双有理}的情形，而问题完全归结为证明 $j$
\emph{满射}。设 $z\in P$；$\mathcal{O}_z$ 是一个整局部环
（相应地，诺特整局部环），其分式域为
\[
L=R(P)=R(X')=R(X).
\]
可以只考虑 $z$ 不是 $P$ 的泛点的情形。因此，由
\egaref{II.7.1.2-zh}{7.1.2} 和 \egaref{II.7.1.7-zh}{7.1.7}，
存在一个以 $L$ 为分式域并支配 $\mathcal{O}_z$ 的赋值环
（相应地，离散赋值环）$A$。更有理由，置 $y=p(z)$，则 $A$ 支配
$\mathcal{O}_y$，故由假设，存在 $x\in X$，使得 $A$ 支配
$\mathcal{O}_x$。由于 $g$ 是紧合态射，证明的第一部分表明，$A$
还支配某个 $\mathcal{O}_{x'}$，其中 $x'\in X'$。于是
$\mathcal{O}_z$ 与 $\mathcal{O}_{j(x')}=\mathcal{O}_{x'}$ 同源
（\egaxref{I.8.1.4-zh}{I, 8.1.4}）；由于 $P$ 是概形，这就蕴含
$z=j(x')$（\egaxref{I.8.2.2-zh}{I, 8.2.2}），证明完毕。
\end{enumerate}

% Locked French authority: ega2-1-fr.tex, 820505 B,
% SHA-256 84EDBE3E83530AF2959B441796337C9DC21EAFCA6A13114A26778760FBF437AC.
% Sequential source scope: lines 12989-13040, Corollary II.7.3.11,
% Remark II.7.3.12, and terminal oldpage 148.

\begin{corollary}[7.3.11]
\label{II.7.3.11-zh}
设 $X$、$Y$ 是两个整概形，$f:X\to Y$ 是支配、准紧且普遍闭的态射。
再假设 $Y$ 是以（整）环 $B$ 为坐标环的仿射概形。则
$\Gamma(X,\mathcal{O}_X)$ 典范同构于 $B$ 在 $R(X)$ 中的整闭包的
一个子环。
\end{corollary}

事实上（\egaxref{I.8.2.1.1-zh}{I, 8.2.1.1}），
$\Gamma(X,\mathcal{O}_X)$ 可等同于所有 $\mathcal{O}_x$（其中
$x\in X$）的交；由 \egaref{II.7.3.10-zh}{7.3.10}、
\egaref{II.7.1.2-zh}{7.1.2} 和 \egaref{II.7.1.3-zh}{7.1.3}，
$\Gamma(X,\mathcal{O}_X)$ 因而包含于所有包含 $B$ 且以 $R(X)$ 为
分式域的赋值环的交中；此时结论由 \egaref{II.7.1.3-zh}{7.1.3}
得出。

\begin{remark}[7.3.12]
\label{II.7.3.12-zh}
在 \egaref{II.7.3.11-zh}{7.3.11} 的假设下，并且当假设 $R(X)$ 是
$R(Y)$ 的有限型扩张时，在许多情形中可以断定
$\Gamma(X,\mathcal{O}_X)$ 是环 $B=\Gamma(Y,\mathcal{O}_Y)$ 上的
一个\emph{有限型模}。例如，当 $B$ 是一个\emph{域上的有限型代数}
时就是如此，因为已知 $B$ 在其分式域的一个有限型扩张中的整闭包是
有限型 $B$-模（[13], t.~I, p.~267, th.~9）；此时结论由
\egaref{II.7.3.11-zh}{7.3.11} 和 $B$ 是诺特环这一事实得出。

特别地，\emph{域 $K$ 上的一个紧合仿射概形 $X$ 是有限的}。事实上，
由 \egaref{II.1.6.4-zh}{1.6.4}、\egaref{II.5.4.6-zh}{5.4.6} 和
\egaxref{I.6.4.4-zh}{I, 6.4.4, c)}，可以归结到 $X$ \emph{约化}的
情形。此外，只需证明 $X$ 的每个闭子预概形在 $K$ 上有限，其中该
闭子预概形的底空间是 $X$ 的一个不可约分支（这样的分支只有有限个）；
因而，考虑到 \egaref{II.5.4.5-zh}{5.4.5}，最终归结到 $X$
\emph{为整概形}的情形。但此时，结论由上述评注得出。

在第三章中，我们将用其他方法再次得到后一命题，并把它作为更一般结果
的推论；这些结果表明，若 $f:X\to Y$ 是紧合态射且 $Y$ 局部诺特，则对
每个凝聚 $\mathcal{O}_X$-模 $\mathcal{F}$，$f_*(\mathcal{F})$ 都是
凝聚的（\egaxref{III.4.4.2-zh}{III, 4.4.2}）。

最后请注意，判别法 \egaref{II.7.3.10-zh}{7.3.10} 在经典代数几何中
被用作紧合态射的\emph{定义}。我们只是出于这一原因才提及它，因为
在我们所知的一切应用中，判别法 \egaref{II.7.3.8-zh}{7.3.8} 似乎
更便于使用。
\end{remark}

\oldpage[II]{148}
% Locked French authority: ega2-1-fr.tex, 820505 B,
% SHA-256 84EDBE3E83530AF2959B441796337C9DC21EAFCA6A13114A26778760FBF437AC.
% Sequential source scope: lines 13041-13131, II.7.4 opening through II.7.4.3.

\subsection{代数曲线与一维函数域}
\label{subsection:II.7.4-zh}

本节旨在说明如何用概形语言表述代数曲线的经典概念（例如 C.~Chevalley
的著作 [23] 中所讲述的概念）。在本节中，\emph{$k$ 始终表示一个域，所考虑
的一切概形都是有限型 $k$-概形，所有态射都是 $k$-态射}。

\begin{proposition}[7.4.1]
\label{II.7.4.1-zh}
设 $X$ 是 $k$ 上的有限型预概形（因而是诺特的）；设 $x_i$
$(1\leq i\leq n)$ 是不可约分支 $X_i$（即 $X$ 的各不可约分支）的泛点，并令
$K_i=k(x_i)$ $(1\leq i\leq n)$。下列条件等价：
\begin{enumerate}
\item[\rm a)] 每个 $K_i$ 都是 $k$ 的超越次数等于 1 的扩张。
\item[\rm b)] 对于每个闭点 $x$（视为 $X$ 的点），局部环 $\mathcal{O}_x$ 的维数
为 1（\egaref{II.7.1.5-zh}{7.1.5}）。
\item[\rm c)] $X$ 的不同于各 $X_i$ 的不可约闭子集，恰为 $X$ 的闭点。
\end{enumerate}
\end{proposition}

由于 $X$ 拟紧，每个不可约闭子集 $F$（它是 $X$ 的子集）都包含一个闭点
（\egaxref{0.2.1.3-zh}{0, 2.1.3}）。由
（\egaxref{I.2.4.2-zh}{I, 2.4.2}），$\mathcal{O}_x$ 的素理想与 $X$
中包含 $x$ 的不可约闭子集之间存在双射对应
（\egaxref{I.1.1.14-zh}{I, 1.1.14}）；b) 与 c) 的等价性立即由此得出。
另一方面，若 $\mathfrak{p}_\alpha$ $(1\leq\alpha\leq r)$ 是诺特局部环
$\mathcal{O}_x$ 的极小素理想，则局部环
$\mathcal{O}_x/\mathfrak{p}_\alpha$ 是整环，其分式域是某些 $K_i$，所对应
的指标满足 $x\in X_i$。此外，已知（[1], p.~4-06, th.~2），有限型整局部
$k$-代数的维数，等于其分式域在 $k$ 上的超越次数。最后，
$\mathcal{O}_x$ 的维数是各 $\mathcal{O}_x/\mathfrak{p}_\alpha$ 的维数的
上确界；而条件 a) 蕴含这些维数等于 1，故 a) 蕴含 b)；反过来，若
$\mathcal{O}_x$ 的维数为 1，则任一 $\mathfrak{p}_\alpha$ 都不可能等于
$\mathcal{O}_x$ 的极大理想，否则 $\mathcal{O}_x$ 的维数便为 0；因此每个
$\mathcal{O}_x/\mathfrak{p}_\alpha$ 的维数都是 1，这就证明了 b) 蕴含 a)。

应注意，在（\egaref{II.7.4.1-zh}{7.4.1}）的条件下，集合 $X$
\emph{或为空集，或为无限集}，这立即由
（\egaxref{I.6.4.4-zh}{I, 6.4.4}）得出。

\begin{definition}[7.4.2]
\label{II.7.4.2-zh}
所谓 $k$ 上的代数曲线，是指 $k$ 上满足
（\egaref{II.7.4.1-zh}{7.4.1}）各条件的非空代数概形。
\end{definition}

用将在第四章中引入的维数语言来说，这表示：$k$ 上的代数曲线是一个非空
代数 $k$-概形，并且\emph{它的所有不可约分支的维数均为 1}。

应注意，若 $X$ 是 $k$ 上的一条代数曲线，则既约闭子预概形
$X_i$ $(1\leq i\leq n)$（它们是 $X$ 的子预概形），其底空间为 $X$ 的不可约分支，也都是 $k$
上的代数曲线。

\begin{corollary}[7.4.3]
\label{II.7.4.3-zh}
设 $X$ 是一条不可约代数曲线。$X$ 唯一的非闭点是它的泛点。$X$
中不同于 $X$ 的闭子集，恰为闭点的有限集合；它们也恰好是 $X$ 中
不处处稠密的子集。
\end{corollary}

若点 $x\in X$ 不是闭点，则它在 $X$ 中的闭包是 $X$ 的不可约闭子集；
因此该闭包必为整个 $X$（\egaref{II.7.4.1-zh}{7.4.1}），从而 $x$
是 $X$ 的泛点。设 $F$ 是 $X$ 中不同于 $X$ 的闭子集；它不可能包含

\oldpage[II]{149}

$X$ 的泛点，故它的所有点都是闭点（既在 $X$ 中是闭的，\emph{更不待言}，
在 $F$ 中也是闭的）；考虑 $X$ 的以 $F$ 为底空间的既约闭子预概形
（\egaxref{I.5.2.1-zh}{I, 5.2.1}），遂由
（\egaxref{I.6.2.2-zh}{I, 6.2.2}）得出 $F$ 是有限且离散的。因而，$X$
的任一无限子集在 $X$ 中的闭包必等于 $X$。

当 $X$ 是任意代数曲线时，把（\egaref{II.7.4.3-zh}{7.4.3}）应用于 $X$
的各不可约分支，即知 $X$ 的非闭点恰为这些分支的泛点。
% Locked French authority: ega2-1-fr.tex, 820505 B,
% SHA-256 84EDBE3E83530AF2959B441796337C9DC21EAFCA6A13114A26778760FBF437AC.
% Sequential source scope: lines 13132-13192, Corollaries II.7.4.4--II.7.4.6.

\begin{corollary}[7.4.4]
\label{II.7.4.4-zh}
设 $X$ 和 $Y$ 是 $k$ 上的两条不可约代数曲线，$f:X\to Y$ 是
$k$-态射。则 $f$ 为支配态射，当且仅当对每个 $y\in Y$，
$f^{-1}(y)$ 都是有限的。
\end{corollary}

事实上，若 $f$ 不是支配态射，则由 \egaref{II.7.4.3-zh}{7.4.3}，
$f(X)$ 必为 $Y$ 的一个\emph{有限}子集；因此，若对 $Y$ 的每个点
$f^{-1}(y)$ 都是有限的，便会有 $X$ 有限，这是不可能的，因为这与
\egaref{II.7.4.1-zh}{7.4.1} 矛盾。反之，若 $f$ 是支配态射，则对每个
不同于 $Y$ 的泛点 $\eta$ 的 $y\in Y$，由于 $\{y\}$ 在 $Y$ 中闭
（\egaref{II.7.4.3-zh}{7.4.3}），$f^{-1}(y)$ 在 $X$ 中闭；另一方面，
依假设，$f^{-1}(y)$ 不含 $X$ 的泛点 $\xi$，故由
\egaref{II.7.4.3-zh}{7.4.3}，它是有限的。最后，为说明当 $f$ 是
支配态射时 $f^{-1}(\eta)$ 是有限的，注意纤维 $f^{-1}(\eta)$ 是
$k(\eta)$ 上的有限型不可约概形，其泛点为 $\xi$
（\egaxref{I.6.3.9-zh}{I, 6.3.9} 和
\egaxref{I.6.4.11-zh}{6.4.11}）。由于 $k(\xi)$ 和 $k(\eta)$ 都是
$k$ 的有限型扩张，并且超越次数同为 1，$k(\xi)$ 必为 $k(\eta)$ 的
有限次扩张；因此 $\xi$ 在 $f^{-1}(\eta)$ 中闭
（\egaxref{I.6.4.2-zh}{I, 6.4.2}），从而 $f^{-1}(\eta)$ 仅由点
$\xi$ 组成。

我们将在第三章中看到：若诺特预概形之间的\emph{紧合}态射
$f:X\to Y$ 满足对每个 $y\in Y$，$f^{-1}(y)$ 都是有限的，则该态射
必为\emph{有限}态射；因而由 \egaref{II.7.4.4-zh}{7.4.4} 可知，
从一条不可约代数曲线到一条代数曲线的支配紧合态射是\emph{有限}态射。

\begin{corollary}[7.4.5]
\label{II.7.4.5-zh}
设 $X$ 是 $k$ 上的代数曲线。则 $X$ 正则，当且仅当 $X$ 正规，
也当且仅当其闭点的局部环都是离散赋值环。
\end{corollary}

这立即由 \egaref{II.7.4.1-zh}{7.4.1} 的条件 b) 和
\egaref{II.7.1.6-zh}{7.1.6} 得出。

\begin{corollary}[7.4.6]
\label{II.7.4.6-zh}
设 $X$ 是约化代数曲线，$\mathcal{A}$ 是约化的凝聚
$\mathcal{R}(X)$-代数；则 $X$ 关于 $\mathcal{A}$ 的整闭包 $X'$
（\egaref{II.6.3.4-zh}{6.3.4}）是一条正规代数曲线，并且典范态射
$X'\to X$ 是有限态射。
\end{corollary}

$X'\to X$ 为有限态射这一事实由 \egaref{II.6.3.10-zh}{6.3.10}
得出；因此 $X'$ 是 $k$-代数概形。此外，若 $x_i$
$(1\leq i\leq n)$ 是 $X$ 的各不可约分支的泛点，$x'_j$
$(1\leq j\leq m)$ 是 $X'$ 的各不可约分支的泛点，则每个
$k(x'_j)$ 都是某个 $k(x_i)$ 的有限代数扩张
（\egaref{II.6.3.6-zh}{6.3.6}），因而在 $k$ 上的超越次数为 1。
所以 $X'$ 确实是 $k$ 上的代数曲线；此外，已知 $X'$ 是有限个整正规
概形的不交并（\egaref{II.6.3.6-zh}{6.3.6} 和
\egaref{II.6.3.7-zh}{6.3.7}）。
% Locked French authority: ega2-1-fr.tex, 820505 B,
% SHA-256 84EDBE3E83530AF2959B441796337C9DC21EAFCA6A13114A26778760FBF437AC.
% Sequential source scope: lines 13193-13250, II.7.4.7--II.7.4.10.

\begin{env}[7.4.7]
\label{II.7.4.7-zh}
若 $k$ 上的代数曲线 $X$ 在 $k$ 上是\emph{紧合的}，就称它是
\emph{完备的}。
\end{env}

\begin{corollary}[7.4.8]
\label{II.7.4.8-zh}
$k$ 上的约化代数曲线 $X$ 完备，当且仅当它的正规化 $X'$ 完备。
\end{corollary}

\oldpage[II]{150}
事实上，此时典范态射 $f:X'\to X$ 是有限态射
（\egaref{II.7.4.6-zh}{7.4.6}），因而是紧合态射
（\egaref{II.6.1.11-zh}{6.1.11}）和满射
（\egaref{II.6.3.8-zh}{6.3.8}）。若 $g:X\to\operatorname{Spec}(k)$
是结构态射，则由 \egaref{II.5.4.2-zh}{5.4.2, (iii)} 和
\egaref{II.5.4.3-zh}{5.4.3, (ii)}，$g$ 与 $g\circ f$ 同时为紧合
态射，因为依假设 $g$ 是分离的。

\begin{proposition}[7.4.9]
\label{II.7.4.9-zh}
设 $X$ 是 $k$ 上的正规代数曲线，$Y$ 是在 $k$ 上紧合的代数
$k$-概形。则每个从 $X$ 到 $Y$ 的 $k$-有理映射都处处有定义；换言之，
它是一个态射。
\end{proposition}

事实上，由 \egaref{II.7.3.7-zh}{7.3.7}，在这样的映射无定义的点
$x\in X$ 处，$\mathcal{O}_x$ 的维数必须为 $\geq2$；因此这类点的
集合为空。最后一个断言来自 \egaxref{I.7.2.3-zh}{I, 7.2.3}。

\begin{corollary}[7.4.10]
\label{II.7.4.10-zh}
$k$ 上的正规代数曲线在 $k$ 上拟射影。
\end{corollary}

由于 $X$ 是有限条整正规代数曲线的不交并
（\egaref{II.6.3.8-zh}{6.3.8}），可归结到 $X$ 为整概形的情形
（\egaref{II.5.3.6-zh}{5.3.6}）。由于 $X$ 拟紧，它可由有限多个仿射
开集 $U_i$ $(1\leq i\leq n)$ 覆盖；又因每个这样的开集在 $k$ 上都是
有限型的，存在整数 $n_i$ 和一个 $k$-浸入
$f_i:U_i\to\mathbf{P}_k^{n_i}$
（\egaref{II.5.3.3-zh}{5.3.3} 和
\egaref{II.5.3.4-zh}{5.3.4, (i)}）。由于 $U_i$ 在 $X$ 中稠密，
由 \egaref{II.7.4.9-zh}{7.4.9}，$f_i$ 延拓为一个 $k$-态射
$g_i:X\to\mathbf{P}_k^{n_i}$，从而得到从 $X$ 到各
$\mathbf{P}_k^{n_i}$ 在 $k$ 上的积 $P$ 的一个 $k$-态射
$g=(g_1,\ldots,g_n)_k$。此外，对每个指标 $i$，由于 $g_i$ 在 $U_i$
上的限制是浸入，$g$ 在 $U_i$ 上的限制也是浸入
（\egaxref{I.5.3.14-zh}{I, 5.3.14}）。诸 $U_i$ 覆盖 $X$，且 $g$
是分离的（\egaxref{I.5.5.1-zh}{I, 5.5.1, (v)}），所以 $g$ 是从
$X$ 到 $P$ 的浸入（\egaxref{I.8.2.8-zh}{I, 8.2.8}）。由于 Segre
态射（\egaref{II.4.3.3-zh}{4.3.3}）给出从 $P$ 到某个
$\mathbf{P}_k^N$ 的浸入，这就证明了 $X$ 是拟射影的。
% Locked French authority: ega2-1-fr.tex, 820505 B,
% SHA-256 84EDBE3E83530AF2959B441796337C9DC21EAFCA6A13114A26778760FBF437AC.
% Sequential source scope: lines 13251-13293, Corollary II.7.4.11,
% oldpage 151, and Remark II.7.4.12.

\begin{corollary}[7.4.11]
\label{II.7.4.11-zh}
正规代数曲线 $X$ 同构于某条正规完备代数曲线 $\widehat{X}$ 在一个
稠密开集上所诱导的概形；该曲线在唯一同构意义下确定。
\end{corollary}

若 $X_1$、$X_2$ 是两条正规完备曲线，则由
\egaref{II.7.4.9-zh}{7.4.9} 立即可知，从 $X_1$ 中的稠密开集 $U_1$
到 $X_2$ 中的稠密开集 $U_2$ 的任何同构，都唯一延拓为从 $X_1$ 到
$X_2$ 的同构；这就证明了唯一性的断言。为证明 $\widehat{X}$ 的存在性，
只需注意，可以把 $X$ 看成射影空间 $\mathbf{P}_k^n$ 的一个子概形
（\egaref{II.7.4.10-zh}{7.4.10}）。设 $\overline{X}$ 是 $X$ 在
$\mathbf{P}_k^n$ 中的闭包（\egaxref{I.9.5.11-zh}{I, 9.5.11}）；
由于 $X$ 是 $\overline{X}$ 在 $\overline{X}$ 的一个稠密开集上所诱导
的概形（\egaxref{I.9.5.10-zh}{I, 9.5.10}），$X$ 的各不可约分支的
泛点 $x_i$ 也是 $\overline{X}$ 的各不可约分支的泛点，并且对这两个
概形，$k(x_i)$ 都相同。因此，由 \egaref{II.7.4.1-zh}{7.4.1}，
$\overline{X}$ 是 $k$ 上的代数曲线；它是约化的
（\egaxref{I.9.5.9-zh}{I, 9.5.9}），并且在 $k$ 上是射影的
（\egaref{II.5.5.1-zh}{5.5.1}），从而是完备的
（\egaref{II.5.5.3-zh}{5.5.3}）。此时取 $\widehat{X}$ 为
$\overline{X}$ 的\emph{正规化}；它仍是完备的
（\egaref{II.7.4.8-zh}{7.4.8}）。此外，若
$h:\widehat{X}\to\overline{X}$ 是典范态射，则因为 $X$ 正规，$h$
在 $h^{-1}(X)$ 上的限制是到 $X$ 的同构
（\egaref{II.6.3.4-zh}{6.3.4}）；又因为 $h^{-1}(X)$ 包含
$\widehat{X}$ 的各不可约分支的泛点
（\egaref{II.6.3.8-zh}{6.3.8}），它在 $\widehat{X}$ 中稠密，证明
至此完成。

\oldpage[II]{151}

\begin{remark}[7.4.12]
\label{II.7.4.12-zh}
我们将在第五章中证明，即使不假设曲线正规（甚至不假设它约化），
\egaref{II.7.4.10-zh}{7.4.10} 的结论仍然成立；我们还将证明，一条
代数曲线（无论约化与否）是仿射的，当且仅当其赋予约化结构的各不可约
分支都不是完备的。
\end{remark}
% Locked French authority: ega2-1-fr.tex, 820505 B,
% SHA-256 84EDBE3E83530AF2959B441796337C9DC21EAFCA6A13114A26778760FBF437AC.
% Sequential source scope: lines 13294-13326, Corollary II.7.4.13 and
% Environment II.7.4.14.

\begin{corollary}[7.4.13]
\label{II.7.4.13-zh}
设 $X$ 是函数域 $R(X)=K$ 的正规不可约曲线，$Y$ 是函数域
$L=R(Y)$ 的整完备曲线。支配 $k$-态射 $X\to Y$ 与 $k$-单同态
$L\to K$ 之间存在典范一一对应。
\end{corollary}

由 \egaref{II.7.4.9-zh}{7.4.9}，从 $X$ 到 $Y$ 的 $k$-有理映射
可与 $k$-态射 $u:X\to Y$ 等同。支配态射 $u$ 的特征是 $u(x)=y$
（其中 $x$ 和 $y$ 分别表示 $X$ 和 $Y$ 的泛点），故本推论由这些说明
以及 \egaxref{I.7.1.13-zh}{I, 7.1.13} 得出。

\begin{env}[7.4.14]
\label{II.7.4.14-zh}
当取 $Y$ 为\emph{射影直线}
$\mathbf{P}_k^1=\operatorname{Proj}(k[T_0,T])$ 时，可以进一步明确
\egaref{II.7.4.13-zh}{7.4.13} 的结果，其中 $T_0$ 和 $T$ 是两个
不定元。它确实是整概形（\egaref{II.2.4.4-zh}{2.4.4}），而 $Y$
在开集 $D_+(T_0)$ 上所诱导的概形同构于
$\operatorname{Spec}(k[T])$（\egaref{II.2.3.6-zh}{2.3.6}）；因此，
$Y$ 的泛点是 $k[T]$ 的理想 $(0)$，其有理函数域是 $k(T)$，这表明
$Y$ 是 $k$ 上的完备代数曲线。此外，$S=k[T_0,T]$ 中包含 $T_0$
且不同于 $S_+$ 的唯一分次素理想是主理想 $(T_0)$，故
$Y=\mathbf{P}_k^1$ 中 $D_+(T_0)$ 的补集约化为\emph{一个闭点}，
称为“无穷远点”，并记为 $\infty$（关于向量丛与射影丛之间关系的一般
研究，见 \egaref{subsection:II.8.4-zh}{8.4}）。采用这些记号：
\end{env}

% Locked French authority: ega2-1-fr.tex, 820505 B,
% SHA-256 84EDBE3E83530AF2959B441796337C9DC21EAFCA6A13114A26778760FBF437AC.
% Sequential source scope: lines 13327-13369, Corollary II.7.4.15,
% Lemma II.7.4.15.1, and oldpage 152.

\begin{corollary}[7.4.15]
\label{II.7.4.15-zh}
设 $X$ 是函数域为 $R(X)=K$ 的正规不可约曲线。存在从 $K$ 到所有
从 $X$ 到 $\mathbf{P}_k^1$ 且不同于取值为 $\infty$ 的常值态射 $u$
所成集合的典范双射。态射 $u$ 为支配态射，当且仅当 $K$ 中与它对应
的元素在 $k$ 上超越。
\end{corollary}

此陈述立即由 \egaref{II.7.4.9-zh}{7.4.9} 和下述引理得出。

\begin{lemma}[7.4.15.1]
\label{II.7.4.15.1-zh}
设 $X$ 是 $k$ 上的整预概形，$K=R(X)$ 是其有理函数域。存在从集合
$K$ 到所有从 $X$ 到 $\mathbf{P}_k^1$ 且不同于取值为 $\infty$ 的
常值态射的有理映射 $u$ 所成集合的典范双射。这样的有理映射为支配
映射，当且仅当 $K$ 中与它对应的元素在 $k$ 上超越。
\end{lemma}

首先，从 $X$ 到 $\mathbf{P}_k^1$ 的有理映射与 $\mathbf{P}_k^1$ 中
取值于 $k$ 的扩张 $K$ 的点一一对应
（\egaxref{I.7.1.12-zh}{I, 7.1.12}）。若这样的点落在
$\mathbf{P}_k^1$ 的泛点上（\egaxref{I.3.4.5-zh}{I, 3.4.5}），则
相应的有理映射显然是支配的。否则，由于 $\mathbf{P}_k^1$ 中每个
不同于泛点的点都是闭点（\egaref{II.7.4.3-zh}{7.4.3}），$u$ 的
定义域 $U$ 在类 $u$ 中唯一的态射 $U\to\mathbf{P}_k^1$
（\egaxref{I.7.2.2-zh}{I, 7.2.2}）下的像仅由 $\mathbf{P}_k^1$ 的
一个闭点 $y$ 组成；因此该态射（它未必在 $X$ 上处处有定义）不是
支配态射。滥用语言，此时称有理映射 $u$ “为常值映射，取值为 $y$”。
剩下只需在 $\mathbf{P}_k^1$ 中所落之点
（\egaxref{I.3.4.5-zh}{I, 3.4.5}）不同于 $\infty$ 的 $K$-值点与
$K$ 的元素之间建立一一对应，并验证这种点落在 $\mathbf{P}_k^1$ 的
泛点上，当且仅当它对应于一个

\oldpage[II]{152}
在 $k$ 上超越的元素。由 \egaref{II.4.2.6-zh}{4.2.6，例
1\textsuperscript{o}}，这一验证是立即的。
% Locked French authority: ega2-1-fr.tex, 820505 B,
% SHA-256 84EDBE3E83530AF2959B441796337C9DC21EAFCA6A13114A26778760FBF437AC.
% Sequential source scope: lines 13370-13390, II.7.4.16-II.7.4.17.

\begin{corollary}[7.4.16]
\label{II.7.4.16-zh}
设 $X$、$Y$ 是 $k$ 上两条正规、完备且不可约的代数曲线；设
$K=R(X)$、$L=R(Y)$ 是它们的有理函数域。所有 $k$-同构
$X\xrightarrow{\sim}Y$ 组成的集合与所有 $k$-同构
$L\xrightarrow{\sim}K$ 组成的集合之间，存在一个典范双射对应。
\end{corollary}

这是（\egaref{II.7.4.13-zh}{7.4.13}）的一个显然推论。

\begin{env}[7.4.17]
\label{II.7.4.17-zh}
推论（\egaref{II.7.4.16-zh}{7.4.16}）表明，$k$ 上一条正规、完备且
不可约的代数曲线，\emph{在唯一同构意义下由其有理函数域 $K$ 所确定}；
依定义，$K$ 是 $k$ 的有限型扩张，超越次数为 $1$（经典上称之为一个
\emph{一元代数函数域}）。此外：
\end{env}
% Locked French authority: ega2-1-fr.tex, 820505 B,
% SHA-256 84EDBE3E83530AF2959B441796337C9DC21EAFCA6A13114A26778760FBF437AC.
% Sequential source scope: lines 13391-13433, II.7.4.18-II.7.4.19.

\begin{proposition}[7.4.18]
\label{II.7.4.18-zh}
任给扩张 $K$，它是 $k$ 的有限型扩张且超越次数为 $1$，则存在一条
正规、完备且不可约的代数曲线 $X$，使得 $R(X)=K$（它在唯一同构
意义下确定）。$X$ 的局部环所成集合（\egaxref{I.8.2.1-zh}{I, 8.2.1}）
可与由 $K$ 以及所有包含 $k$ 且以 $K$ 为分式域的赋值环组成的集合
等同。
\end{proposition}

事实上，$K$ 是纯超越扩张 $k(T)$ 的有限次扩张；后者是 $k$ 的扩张，
如前所见，它可与射影直线 $Y=\mathbf{P}_k^1$ 的有理函数域等同。设 $X$ 是
$Y$ 关于 $K$ 的\emph{整闭包}（\egaref{II.6.3.4-zh}{6.3.4}）；$X$ 是
函数域为 $K$ 的正规代数曲线（\egaref{II.6.3.7-zh}{6.3.7}），并且它
是完备的，因为态射 $X\to Y$ 是有限态射
（\egaref{II.7.4.6-zh}{7.4.6}）。$X$ 的局部环 $\mathcal{O}_x$ 如下：
当 $x$ 是泛点时，它是域 $K$；当 $x$ 不同于泛点时，$\mathcal{O}_x$
是一个包含 $k$ 且以 $K$ 为分式域的离散赋值环
（\egaref{II.7.4.5-zh}{7.4.5}）。反之，设 $A$ 是这样的环；由于态射
$X\to\operatorname{Spec}(k)$ 是固有的，且 $A$ 支配 $k$，故它支配
$X$ 的某个局部环 $\mathcal{O}_x$（\egaref{II.7.3.10-zh}{7.3.10}）；
后者是一个以 $K$ 为分式域的赋值环，因而必等于 $A$。

\begin{namedtheorem}[7.4.19]{注}
\label{II.7.4.19-zh}
由（\egaref{II.7.4.16-zh}{7.4.16}）和
（\egaref{II.7.4.18-zh}{7.4.18}）可知，\emph{给定 $k$ 上一条正规、
完备且不可约的代数曲线，本质上等价于给定 $k$ 的一个有限型且超越
次数为 $1$ 的扩张 $K$}。应当注意，若 $k'$ 是基域 $k$ 的一个扩张，
则 $X\otimes_k k'$ 仍是 $k'$ 上的一条完备代数曲线
（\egaref{II.5.4.2-zh}{5.4.2，（iii）}），但一般而言它既非约化的，
也非不可约的。然而，若 $K$ 是 $k$ 的可分扩张，并且 $k$ 在 $K$ 中
是代数闭的，则它既约化又不可约（在我们不采用的经典术语中，这表述
为 $K$ 是 $k$ 的一个“正则扩张”）。但即便在这种情形下，
$X\otimes_k k'$ 仍可能不是正规的。读者可在第四章中找到这些问题的
详细讨论。
\end{namedtheorem}
% Locked French authority: ega2-1-fr.tex, 820505 B,
% SHA-256 84EDBE3E83530AF2959B441796337C9DC21EAFCA6A13114A26778760FBF437AC.
% Sequential source scope: lines 13434-13486, Chapter II.8 and subsection
% II.8.1 headings through Environment II.8.1.1.

\section{概形的暴涨；投影锥；泛射影闭包}
\label{section:II.8-zh}

\subsection{预概形的暴涨}
\label{subsection:II.8.1-zh}

\begin{env}[8.1.1]
\label{II.8.1.1-zh}
设 $Y$ 是预概形，并且对每个整数 $n\geq0$，设 $\mathcal{I}_n$ 是
$\mathcal{O}_Y$ 的拟凝聚理想；假设满足下列条件：
\begin{equation}
\label{II.8.1.1.1-zh}
  \mathcal{I}_0=\mathcal{O}_Y,\qquad
  \mathcal{I}_n\subset\mathcal{I}_m\quad\text{当 }m\leq n
\tag{8.1.1.1}
\end{equation}
\begin{equation}
\label{II.8.1.1.2-zh}
  \mathcal{I}_m\mathcal{I}_n\subset\mathcal{I}_{m+n}
  \quad\text{对任意 }m,n.
\tag{8.1.1.2}
\end{equation}

\oldpage[II]{153}
注意，这些假设蕴含
\begin{equation}
\label{II.8.1.1.3-zh}
  \mathcal{I}_1^n\subset\mathcal{I}_n.
\tag{8.1.1.3}
\end{equation}

置
\begin{equation}
\label{II.8.1.1.4-zh}
  \mathcal{S}=\bigoplus_{n\geq0}\mathcal{I}_n.
\tag{8.1.1.4}
\end{equation}

由 \egaref{II.8.1.1.1-zh}{8.1.1.1} 和
\egaref{II.8.1.1.2-zh}{8.1.1.2} 可知，$\mathcal{S}$ 是拟凝聚分次
$\mathcal{O}_Y$-代数，因而定义了 $Y$-概形
$X=\operatorname{Proj}(\mathcal{S})$。若 $\mathcal{J}$ 是
$\mathcal{O}_Y$ 的\emph{可逆}理想，则
$\mathcal{I}_n\otimes_{\mathcal{O}_Y}\mathcal{J}^{\otimes n}$ 典范
等同于 $\mathcal{I}_n\mathcal{J}^n$。因此，若把 $\mathcal{I}_n$
替换为 $\mathcal{I}_n\mathcal{J}^n$，从而把 $\mathcal{S}$ 替换为
拟凝聚 $\mathcal{O}_Y$-代数 $\mathcal{S}_{(\mathcal{J})}$，则
$X_{(\mathcal{J})}=\operatorname{Proj}(\mathcal{S}_{(\mathcal{J})})$
典范同构于 $X$（\egaref{II.3.1.8-zh}{3.1.8}）。
\end{env}

% Locked French authority: ega2-1-fr.tex, 820505 B,
% SHA-256 84EDBE3E83530AF2959B441796337C9DC21EAFCA6A13114A26778760FBF437AC.
% Sequential source scope: lines 13487-13544, II.8.1.2-II.8.1.3.

\begin{env}[8.1.2]
\label{II.8.1.2-zh}
假设 $Y$ \emph{局部整}，因而有理函数层 $\mathcal{R}(Y)$ 是一个拟凝聚
$\mathcal{O}_Y$-代数层（\egaxref{I.7.3.7-zh}{I, 7.3.7}）。我们称一个
$\mathcal{O}_Y$-子模层 $\mathcal{I}$（它包含于 $\mathcal{R}(Y)$）为
$\mathcal{R}(Y)$ 的\emph{分式理想层}，如果它是\emph{有限型}的
（\egaxref{0.5.2.1-zh}{0, 5.2.1}）。假设对每个 $n\geq0$，给定一个拟凝聚
分式理想层 $\mathcal{I}_n$，它视为 $\mathcal{R}(Y)$ 的子模层，使得
$\mathcal{I}_0=\mathcal{O}_Y$，并且条件
（\egaref{II.8.1.1.2-zh}{8.1.1.2}）成立（但第二个条件
（\egaref{II.8.1.1.1-zh}{8.1.1.1}）不一定成立）；于是仍可由公式
（\egaref{II.8.1.1.4-zh}{8.1.1.4}）定义一个拟凝聚分次
$\mathcal{O}_Y$-代数层，以及相应的 $Y$-概形
$X=\operatorname{Proj}(\mathcal{S})$；仍有从 $X$ 到 $X_{(\mathcal{J})}$ 的
典范同构，其中 $\mathcal{J}$ 是 $\mathcal{R}(Y)$ 的任意可逆分式理想层。
\end{env}

\begin{definition}[8.1.3]
\label{II.8.1.3-zh}
设 $Y$ 是一个预概形（相应地，一个局部整预概形），而 $\mathcal{I}$ 是
$\mathcal{O}_Y$ 的一个拟凝聚理想层（相应地，$\mathcal{R}(Y)$ 的一个拟凝聚
分式理想层）。称 $Y$-概形
$X=\operatorname{Proj}(\bigoplus_{n\geq0}\mathcal{I}^n)$ 是\emph{对理想层
$\mathcal{I}$ 作暴涨}所得，或称为\emph{$Y$ 关于 $\mathcal{I}$ 的暴涨预概形}。
当 $\mathcal{I}$ 是 $\mathcal{O}_Y$ 的拟凝聚理想层，而 $Y'$ 是 $Y$ 的
由 $\mathcal{I}$ 定义的闭子预概形时，也称 $X$ 是一个 $Y$-概形，它由对
$Y'$ 作暴涨而得到。
\end{definition}

依定义，$\mathcal{S}=\bigoplus_{n\geq0}\mathcal{I}^n$ 由
$\mathcal{S}_1=\mathcal{I}$ 生成；若 $\mathcal{I}$ 是\emph{有限型}
$\mathcal{O}_Y$-模层，则 $X$ 在 $Y$ 上是\emph{射影的}
（\egaref{II.5.5.2-zh}{5.5.2}）。不对 $\mathcal{I}$ 作任何假设时，
$\mathcal{O}_X$-模层 $\mathcal{O}_X(1)$ 是\emph{可逆的}
（\egaref{II.3.2.5-zh}{3.2.5}），并且由
（\egaref{II.4.4.3-zh}{4.4.3}），它对于结构态射 $X\to Y$ 是
\emph{极丰沛的}。

注意，若 $f:X\to Y$ 是结构态射，则 $f$ 在
$f^{-1}(Y-Y')$ 上的限制，是到 $Y-Y'$ 的一个\emph{同构}；这里
$\mathcal{I}$ 是 $\mathcal{O}_Y$ 的理想层，且 $Y'$ 是它所定义的闭子预概形：
事实上，由于问题在 $Y$ 上是局部的，只需假设
$\mathcal{I}=\mathcal{O}_Y$，而我们的断言便由
（\egaref{II.3.1.7-zh}{3.1.7}）得出。

当把 $\mathcal{I}$ 替换为 $\mathcal{I}^d$（$d>0$）时，暴涨 $Y$-概形
$X$ 被一个与之典范同构的 $Y$-概形 $X'$ 替换
（\egaref{II.8.1.1-zh}{8.1.1}）；同样，对于每个可逆理想层（相应地，可逆
分式理想层）$\mathcal{J}$，暴涨预概形 $X_{(\mathcal{J})}$（相对于理想层
$\mathcal{I}\mathcal{J}$）典范同构于 $X$
（\egaref{II.8.1.1-zh}{8.1.1}）。

特别地，当 $\mathcal{I}$ 是理想层（相应地，分式理想层）且\emph{可逆}时，
所得的 $Y$-概形（即对 $\mathcal{I}$ 作暴涨所得者）\emph{同构}于 $Y$
（\egaref{II.3.1.7-zh}{3.1.7}）。
% Locked French authority: ega2-1-fr.tex, 820505 B,
% SHA-256 84EDBE3E83530AF2959B441796337C9DC21EAFCA6A13114A26778760FBF437AC.
% Sequential source scope: lines 13545-13589, Proposition II.8.1.4 and proof.

\begin{proposition}[8.1.4]
\label{II.8.1.4-zh}
设 $Y$ 是整预概形。
\begin{enumerate}
\item[\rm(i)] 对于任意序列 $(\mathcal{I}_n)$，其各项均为
$\mathcal{R}(Y)$ 的拟凝聚分式理想层，并满足
（\egaref{II.8.1.1.2-zh}{8.1.1.2}），

\oldpage[II]{154}
若 $\mathcal{I}_0=\mathcal{O}_Y$，则 $Y$-概形
$X=\operatorname{Proj}(\bigoplus_{n\geq0}\mathcal{I}_n)$ 是整的，并且
结构态射 $f:X\to Y$ 是支配态射。
\item[\rm(ii)] 设 $\mathcal{I}$ 是 $\mathcal{R}(Y)$ 的一个拟凝聚分式
理想层，并设 $X$ 是 $Y$-概形，它由对 $Y$ 相对于 $\mathcal{I}$ 作暴涨
而得。若 $\mathcal{I}\neq0$，则结构态射 $f:X\to Y$ 是双有理且
满射的。
\end{enumerate}
\end{proposition}

\begin{enumerate}
\item[\rm(i)] 这是因为
$\mathcal{S}=\bigoplus_{n\geq0}\mathcal{I}_n$ 是一个\emph{整的}
$\mathcal{O}_Y$-代数层（\egaref{II.3.1.12-zh}{3.1.12} 和
\egaref{II.3.1.14-zh}{3.1.14}）；事实上，对每个 $y\in Y$，
$\mathcal{O}_y$ 都是整环（\egaxref{I.5.1.4-zh}{I, 5.1.4}）。
\item[\rm(ii)] 由（i），$X$ 是整的；再设 $x$ 和 $y$ 分别是 $X$ 和
$Y$ 的泛点，则 $f(x)=y$，而我们必须证明 $k(x)$ 的秩为 $1$，这是
相对于 $k(y)$ 而言的。然而，$x$ 也是纤维 $f^{-1}(y)$ 的泛点；若
$\psi$ 是典范态射
$Z\to Y$，其中 $Z=\operatorname{Spec}(k(y))$，则预概形 $f^{-1}(y)$
可与 $\operatorname{Proj}(\mathcal{S}')$ 等同，其中
$\mathcal{S}'=\psi^*(\mathcal{S})$（\egaref{II.3.5.3-zh}{3.5.3}）。
但显然有 $\mathcal{S}'=\bigoplus_{n\geq0}(\mathcal{I}_y)^n$；又因
$\mathcal{I}$ 是 $\mathcal{R}(Y)$ 的非零拟凝聚分式理想层，故
$\mathcal{I}_y\neq0$（\egaxref{I.7.3.6-zh}{I, 7.3.6}），从而
$\mathcal{I}_y=k(y)$；因此 $\operatorname{Proj}(\mathcal{S}')$ 可与
$\operatorname{Spec}(k(y))$ 等同（\egaref{II.3.1.7-zh}{3.1.7}），结论
随之得出。
\end{enumerate}

我们还要证明（\egaref{II.8.1.4-zh}{8.1.4}）的一个\emph{逆命题}；
这一证明将在（\egaxref{III.2.3.8-zh}{III, 2.3.8}）中给出。
% Locked French authority: ega2-1-fr.tex, 820505 B,
% SHA-256 84EDBE3E83530AF2959B441796337C9DC21EAFCA6A13114A26778760FBF437AC.
% Sequential source scope: lines 13590-13672, II.8.1.5.

\begin{env}[8.1.5]
\label{II.8.1.5-zh}
回到（\egaref{II.8.1.1-zh}{8.1.1}）的情形和记号。依定义，单射同态
$\mathcal{I}_{n+1}\to\mathcal{I}_n$
（\egaref{II.8.1.1.1-zh}{8.1.1.1}）对每个 $k\in\mathbf{Z}$ 定义一个次数为
$0$ 的分次 $\mathcal{S}$-模层单射同态
\begin{equation}
\label{II.8.1.5.1-zh}
  u_k:\mathcal{S}_+(k+1)\to\mathcal{S}(k)~;
\tag{8.1.5.1}
\end{equation}
由于 $\mathcal{S}_+(k+1)$ 与 $\mathcal{S}(k+1)$ 典范地 (TN)-同构，
$u_k$ 典范地对应于一个 $\mathcal{O}_X$-模层单射同态
（\egaref{II.3.4.2-zh}{3.4.2}）：
\begin{equation}
\label{II.8.1.5.2-zh}
  \widetilde{u}_k:\mathcal{O}_X(k+1)\to\mathcal{O}_X(k).
\tag{8.1.5.2}
\end{equation}

另一方面，回顾（\egaref{II.3.2.6-zh}{3.2.6}），我们已经定义了典范同态
\begin{equation}
\label{II.8.1.5.3-zh}
  \lambda:\mathcal{O}_X(h)\otimes_{\mathcal{O}_X}\mathcal{O}_X(k)
  \to\mathcal{O}_X(h+k)
\tag{8.1.5.3}
\end{equation}
并且由于图表
\[
\xymatrix{
  \mathcal{S}(h)\otimes_{\mathcal{S}}\mathcal{S}(k)
    \otimes_{\mathcal{S}}\mathcal{S}(l)\ar[r]\ar[d]
  & \mathcal{S}(h+k)\otimes_{\mathcal{S}}\mathcal{S}(l)\ar[d]\\
  \mathcal{S}(h)\otimes_{\mathcal{S}}\mathcal{S}(k+l)\ar[r]
  & \mathcal{S}(h+k+l)
}
\]
交换，由 $\lambda$ 的函子性（\egaref{II.3.2.6-zh}{3.2.6}）可知，同态
（\egaref{II.8.1.5.3-zh}{8.1.5.3}）在
\begin{equation}
\label{II.8.1.5.4-zh}
  \mathcal{S}_X=\bigoplus_{n\in\mathbf{Z}}\mathcal{O}_X(n)
\tag{8.1.5.4}
\end{equation}

\oldpage[II]{155}
上定义了一个拟凝聚分次 $\mathcal{O}_X$-代数层结构。此外，图表
\[
\xymatrix{
  \mathcal{S}(h)\otimes_{\mathcal{S}}\mathcal{S}_+(k+1)\ar[r]\ar[d]_{1\otimes u_k}
  & \mathcal{S}_+(h+k+1)\ar[d]^{u_{k+h}}\\
  \mathcal{S}(h)\otimes_{\mathcal{S}}\mathcal{S}(k)\ar[r]
  & \mathcal{S}(h+k)
}
\]
交换；于是 $\lambda$ 的函子性表明下图交换：
\begin{equation}
\label{II.8.1.5.5-zh}
\xymatrix{
  \mathcal{O}_X(h)\otimes_{\mathcal{O}_X}\mathcal{O}_X(k+1)
    \ar[r]^{\lambda}\ar[d]_{1\otimes\widetilde{u}_k}
  & \mathcal{O}_X(h+k+1)\ar[d]^{\widetilde{u}_{k+h}}\\
  \mathcal{O}_X(h)\otimes_{\mathcal{O}_X}\mathcal{O}_X(k)
    \ar[r]^{\lambda}
  & \mathcal{O}_X(h+k)
}
\tag{8.1.5.5}
\end{equation}
其中水平箭头是典范同态。因此可以说，各 $\widetilde{u}_k$ 定义了一个
\emph{单射同态}（次数为 $0$），即一个\emph{分次 $\mathcal{S}_X$-模层同态}
\begin{equation}
\label{II.8.1.5.6-zh}
  \widetilde{u}:\mathcal{S}_X(1)\to\mathcal{S}_X.
\tag{8.1.5.6}
\end{equation}
\end{env}

% Locked French authority: ega2-1-fr.tex, 820505 B,
% SHA-256 84EDBE3E83530AF2959B441796337C9DC21EAFCA6A13114A26778760FBF437AC.
% Sequential source scope: lines 13673-13708, Environment II.8.1.6
% through oldpage 156.

\begin{env}[8.1.6]
\label{II.8.1.6-zh}
沿用 \egaref{II.8.1.5-zh}{8.1.5} 的记号。现在注意，对 $n\geq0$，
复合同态
$\widetilde{v}_n=\widetilde{u}_{n-1}\widetilde{u}_{n-2}\cdots\widetilde{u}_0$
是一个\emph{单射}同态 $\mathcal{O}_X(n)\to\mathcal{O}_X$；我们用
$\mathcal{I}_{n,X}$ 表示它的像，因而这是 $\mathcal{O}_X$ 的一个
拟凝聚理想，并且\emph{同构}于 $\mathcal{O}_X(n)$。此外，图表
\[
\xymatrix{
  \mathcal{O}_X(m)\otimes_{\mathcal{O}_X}\mathcal{O}_X(n)
    \ar[r]^{\lambda}\ar[d]_{\widetilde{v}_m\otimes\widetilde{v}_n}
  & \mathcal{O}_X(m+n)\ar[d]^{\widetilde{v}_{m+n}}\\
  \mathcal{O}_X\ar[r]_{\mathrm{id}}
  & \mathcal{O}_X
}
\]
对 $m\geq0$、$n\geq0$ 是交换的。由此得到下列包含关系：
\begin{equation}
\label{II.8.1.6.1-zh}
  \mathcal{I}_{0,X}=\mathcal{O}_X,\qquad
  \mathcal{I}_{n,X}\subset\mathcal{I}_{m,X}
  \quad\text{当 }0\leq m\leq n
\tag{8.1.6.1}
\end{equation}
\begin{equation}
\label{II.8.1.6.2-zh}
  \mathcal{I}_{m,X}\mathcal{I}_{n,X}\subset\mathcal{I}_{m+n,X}
  \quad\text{当 }m\geq0,\ n\geq0.
\tag{8.1.6.2}
\end{equation}
\end{env}

\oldpage[II]{156}

\begin{proposition}[8.1.7]
\label{II.8.1.7-zh}
设 $Y$ 是预概形，$\mathcal{I}$ 是 $\mathcal{O}_Y$ 的一个拟凝聚理想，
$X=\operatorname{Proj}(\bigoplus_{n\geq0}\mathcal{I}^n)$ 是对 $\mathcal{I}$
作暴涨所得的 $Y$-概形。于是，对于每个 $n>0$，有典范同构
\begin{equation}
\label{II.8.1.7.1-zh}
  \mathcal{O}_X(n)\xrightarrow{\sim}\mathcal{I}^n\mathcal{O}_X
  =\mathcal{I}_{n,X}
\tag{8.1.7.1}
\end{equation}
（参见（\egaxref{0.4.3.5-zh}{0, 4.3.5}）），因而当 $n>0$ 时，
$\mathcal{I}^n\mathcal{O}_X$ 是一个极丰沛可逆 $\mathcal{O}_X$-模。
\end{proposition}

最后一个断言是立即的，因为 $\mathcal{O}_X(1)$ 是可逆的
（\egaref{II.3.2.5-zh}{3.2.5}），并且按定义它相对于 $Y$ 是极丰沛的
（\egaref{II.4.4.3-zh}{4.4.3} 与 \egaref{II.4.4.9-zh}{4.4.9}）。另一方面，
按定义，$v_n$ 的像恰为 $\mathcal{I}^n\mathcal{S}$，故
（\egaref{II.8.1.7.1-zh}{8.1.7.1}）由函子 $\widetilde{\mathcal{M}}$ 的正合性
（\egaref{II.3.2.4-zh}{3.2.4}）以及公式
（\egaref{II.3.2.4.1-zh}{3.2.4.1}）得出。
% Locked French authority: ega2-1-fr.tex, 820505 B,
% SHA-256 84EDBE3E83530AF2959B441796337C9DC21EAFCA6A13114A26778760FBF437AC.
% Sequential source scope: lines 13735-13775, Corollary II.8.1.8 and
% Environment II.8.1.9.

\begin{corollary}[8.1.8]
\label{II.8.1.8-zh}
在 \egaref{II.8.1.7-zh}{8.1.7} 的假设下，若 $f:X\to Y$ 是结构态射，
$Y'$ 是由 $\mathcal{I}$ 定义的 $Y$ 的闭子预概形，则 $X$ 的闭子预概形
$X'=f^{-1}(Y')$ 由 $\mathcal{I}\mathcal{O}_X$ 定义（它典范同构于
$\mathcal{O}_X(1)$），从而有典范正合列
\begin{equation}
\label{II.8.1.8.1-zh}
  0\to\mathcal{O}_X(1)\to\mathcal{O}_X\to\mathcal{O}_{X'}\to0.
\tag{8.1.8.1}
\end{equation}
\end{corollary}

这由 \egaref{II.8.1.7.1-zh}{8.1.7.1} 和
\egaxref{I.4.4.5-zh}{I, 4.4.5} 得出。

\begin{env}[8.1.9]
\label{II.8.1.9-zh}
在 \egaref{II.8.1.7-zh}{8.1.7} 的假设下，可以进一步明确
$\mathcal{I}_{n,X}$ 的结构。事实上，注意同态
\[
  \widetilde{u}_{-1}:\mathcal{O}_X\to\mathcal{O}_X(-1)
\]
典范对应于 $X$ 上 $\mathcal{O}_X(-1)$ 的一个截面 $s$；我们称之为
（相对于 $\mathcal{I}$ 的）\emph{典范截面}
（\egaxref{0.5.1.1-zh}{0, 5.1.1}）。另一方面，图表
\egaref{II.8.1.5.5-zh}{8.1.5.5} 中的水平箭头都是同构
（\egaref{II.3.2.7-zh}{3.2.7}）；在该图表中分别把 $h$、$k$ 换成
$k$、$-1$，便得到
$\widetilde{u}_k=1_k\otimes\widetilde{u}_{-1}$（$1_h$ 表示
$\mathcal{O}_X(h)$ 的恒等映射）。换言之，同态 $\widetilde{u}_k$
正是\emph{与典范截面 $s$ 作张量乘法}（对每个
$k\in\mathbf{Z}$）。同态 $\widetilde{u}$
（\egaref{II.8.1.5.6-zh}{8.1.5.6}）也可以用同样的方式解释。

因此，对每个 $n\geq0$，同态
$\widetilde{v}_n:\mathcal{O}_X(n)\to\mathcal{O}_X$ 正是与
$s^{\otimes n}$ 作张量乘法；由此得到：
\end{env}
% Locked French authority: ega2-1-fr.tex, 820505 B,
% SHA-256 84EDBE3E83530AF2959B441796337C9DC21EAFCA6A13114A26778760FBF437AC.
% Sequential source scope: lines 13776-13806, II.8.1.10-II.8.1.11 partial.

\begin{corollary}[8.1.10]
\label{II.8.1.10-zh}
采用（\egaref{II.8.1.8-zh}{8.1.8}）的记号，$X'$ 的底空间是所有
$x\in X$ 中满足 $s(x)=0$ 的点所成集合，其中 $s$ 表示
$\mathcal{O}_X(-1)$ 的典范截面。
\end{corollary}

事实上，若 $c_x$ 是纤维 $(\mathcal{O}_X(1))_x$ 在一点 $x$ 处的一个
生成元，则 $s_x\otimes c_x$ 可典范地等同于 $\mathcal{I}_{1,X}$ 在点
$x$ 处的纤维的一个生成元，因而它是可逆元，当且仅当
$s_x\notin\mathfrak{m}_x(\mathcal{O}_X(-1))_x$，换言之，当且仅当
$s(x)\neq0$。

\begin{proposition}[8.1.11]
\label{II.8.1.11-zh}
设 $Y$ 是整预概形，$\mathcal{I}$ 是 $\mathcal{R}(Y)$ 的一个拟凝聚
分式理想层，而 $X$ 是 $Y$-概形，它由对 $\mathcal{I}$ 作暴涨而得。
于是 $\mathcal{I}\mathcal{O}_X$ 是可逆 $\mathcal{O}_X$-模，并且相对于
$Y$ 是极丰沛的。
\end{proposition}

由于问题关于 $Y$ 是局部的（\egaref{II.4.4.5-zh}{4.4.5}），只需考虑
$Y=\operatorname{Spec}(A)$ 的情形，其中 $A$ 是分式域为 $K$ 的整环，
且 $\mathcal{I}=\widetilde{\mathfrak{I}}$，这里 $\mathfrak{I}$ 是 $K$ 的
一个分式理想；于是存在元素 $a\neq0$，它属于 $A$，并且满足
$a\mathfrak{I}\subset A$。置 $S=\bigoplus_{n\geq0}\mathfrak{I}^n$；映射
$x\mapsto ax$ 是一个 $A$-同构，它把 $\mathfrak{I}^{n+1}=(S(1))_n$ 映到
$a\mathfrak{I}^{n+1}=a\mathfrak{I}S_n\subset\mathfrak{I}^n=S_n$ 上，
% Locked French authority: ega2-1-fr.tex, 820505 B,
% SHA-256 84EDBE3E83530AF2959B441796337C9DC21EAFCA6A13114A26778760FBF437AC.
% Sequential source scope: lines 13807-13819, terminal prose of II.8.1.

\oldpage[II]{157}
因而定义了一个次数为 $0$ 的分次 $S$-模 (TN)-同构
$S_+(1)\to a\mathfrak{I}S$。另一方面，$x\mapsto a^{-1}x$ 是一个次数为
$0$ 的分次 $S$-模同构
$a\mathfrak{I}S\xrightarrow{\sim}\mathfrak{I}S$。于是通过复合
（\egaref{II.3.2.4-zh}{3.2.4}），得到一个 $\mathcal{O}_X$-模层同构
$\mathcal{O}_X(1)\xrightarrow{\sim}\mathcal{I}\mathcal{O}_X$；又因为 $S$
由 $S_1=\mathfrak{I}$ 生成，所以 $\mathcal{O}_X(1)$ 是可逆的
（\egaref{II.3.2.5-zh}{3.2.5}），并且是极丰沛的
（\egaref{II.4.4.3-zh}{4.4.3} 和 \egaref{II.4.4.9-zh}{4.4.9}），我们的
断言由此得证。

% Locked French authority: ega2-1-fr.tex, 820505 B,
% SHA-256 84EDBE3E83530AF2959B441796337C9DC21EAFCA6A13114A26778760FBF437AC.
% Sequential source scope: lines 13820-13882, subsection II.8.2 heading
% through Environment II.8.2.1.

\subsection{分次环中局部化的预备结果}
\label{subsection:II.8.2-zh}

\begin{env}[8.2.1]
\label{II.8.2.1-zh}
设 $S$ 是一个分次环，此处暂不假设其次数为正。置
\begin{equation}
\label{II.8.2.1.1-zh}
  S^{\geq}=\bigoplus_{n\geq0}S_n,
  \qquad
  S^{\leq}=\bigoplus_{n\leq0}S_n,
\tag{8.2.1.1}
\end{equation}
它们都是 $S$ 的分次子环，其次数分别全为非负和非正。若 $f$ 是 $S$
中次数为 $d$（可正可负）的齐次元，则分式环 $S_f=S'$ 仍带有分次环
结构：对 $n\in\mathbf{Z}$，取 $S'_n$ 为所有 $x/f^k$ 构成的集合，其中
$x\in S_{n+kd}$ 且 $k\geq0$；仍置 $S_{(f)}=S'_0$，并分别以
$S_f^{\geq}$ 和 $S_f^{\leq}$ 表示 $S'^{\geq}$ 和 $S'^{\leq}$。若 $d>0$，则
\begin{equation}
\label{II.8.2.1.2-zh}
  (S^{\geq})_f=S_f
\tag{8.2.1.2}
\end{equation}
成立。事实上，若 $x\in S_{n+kd}$ 且 $n+kd<0$，则可写成
$x/f^k=xf^h/f^{h+k}$，而当 $h$ 充分大且 $>0$ 时，
$n+(h+k)d>0$。因此，由定义可得
\begin{equation}
\label{II.8.2.1.3-zh}
  (S^{\geq})_{(f)}=(S_f^{\geq})_0=S_{(f)}.
\tag{8.2.1.3}
\end{equation}

若 $M$ 是一个分次 $S$-模，同样置
\begin{equation}
\label{II.8.2.1.4-zh}
  M^{\geq}=\bigoplus_{n\geq0}M_n,
  \qquad
  M^{\leq}=\bigoplus_{n\leq0}M_n
\tag{8.2.1.4}
\end{equation}
它们分别是分次 $S^{\geq}$-模和分次 $S^{\leq}$-模，其交为 $S_0$-模
$M_0$。若 $f\in S_d$，则仍把 $M_f$ 定义为分次 $S_f$-模：其次数为
$n$ 的元素取为 $z/f^k$，其中 $z\in M_{n+kd}$ 且 $k\geq0$；以
$M_{(f)}$ 表示 $M_f$ 中所有次数为 $0$ 的元素，它是一个
$S_{(f)}$-模；并分别写 $M_f^{\geq}$ 和 $M_f^{\leq}$ 来代替
$(M_f)^{\geq}$ 和 $(M_f)^{\leq}$。若 $d>0$，与上面相同可见
\begin{equation}
\label{II.8.2.1.5-zh}
  (M^{\geq})_f=M_f
\tag{8.2.1.5}
\end{equation}
且
\begin{equation}
\label{II.8.2.1.6-zh}
  (M^{\geq})_{(f)}=(M_f^{\geq})_0=M_{(f)}.
\tag{8.2.1.6}
\end{equation}
\end{env}
% Locked French authority: ega2-1-fr.tex, 820505 B,
% SHA-256 84EDBE3E83530AF2959B441796337C9DC21EAFCA6A13114A26778760FBF437AC.
% Sequential source scope: lines 13884-13925, II.8.2.2-II.8.2.3.

\begin{env}[8.2.2]
\label{II.8.2.2-zh}
设 $z$ 是一个不定元，我们称它为\emph{齐次化变量}。若 $S$ 是一个分次环
（次数可正可负），则多项式代数\footnote{这里不会与用记号
$\widehat{S}$ 表示一个环的分离完备化的用法相混淆。}
\begin{equation}
\label{II.8.2.2.1-zh}
  \widehat{S}=S[z]
\tag{8.2.2.1}
\end{equation}

\oldpage[II]{158}
它是一个分次 $S$-代数：对于 $fz^n$（$n\geq0$），其中 $f$ 是齐次元，规定
\begin{equation}
\label{II.8.2.2.2-zh}
  \deg(fz^n)=n+\deg f.
\tag{8.2.2.2}
\end{equation}
\end{env}

\begin{lemma}[8.2.3]
\label{II.8.2.3-zh}
\begin{enumerate}
\item[\rm(i)] 有如下（非分次）环的典范同构
\begin{equation}
\label{II.8.2.3.1-zh}
  \widehat{S}_{(z)}\xrightarrow{\sim}
  \widehat{S}/(z-1)\widehat{S}\xrightarrow{\sim}S.
\tag{8.2.3.1}
\end{equation}
\item[\rm(ii)] 有如下（非分次）环的典范同构
\begin{equation}
\label{II.8.2.3.2-zh}
  \widehat{S}_{(f)}\xrightarrow{\sim}S_f^{\leq}
\tag{8.2.3.2}
\end{equation}
对每个 $f\in S_d$ 都成立，其中 $d>0$。
\end{enumerate}
\end{lemma}
% Locked French authority: ega2-1-fr.tex, 820505 B,
% SHA-256 84EDBE3E83530AF2959B441796337C9DC21EAFCA6A13114A26778760FBF437AC.
% Sequential source scope: lines 13927-13972, proof text following
% II.8.2.3 through Lemma II.8.2.5.

\egaref{II.8.2.3.1-zh}{8.2.3.1} 中的第一个同构是在
\egaref{II.2.2.5-zh}{2.2.5} 中定义的，第二个同构是平凡的；这样定义
的同构 $\widehat{S}_{(z)}\xrightarrow{\sim}S$ 把
$xz^n/z^{n+k}$（其中 $\deg(x)=k$（$k\geq-n$））对应到元素 $x$。
同态 \egaref{II.8.2.3.2-zh}{8.2.3.2} 把
$xz^n/f^k$（其中 $\deg(x)=kd-n$）对应到 $S_f^{\leq}$ 中次数为
$-n$ 的元素 $x/f^k$；这里所得确实是同构，这仍是显然的。

\begin{env}[8.2.4]
\label{II.8.2.4-zh}
设 $M$ 是分次 $S$-模。显然，$S$-模
\begin{equation}
\label{II.8.2.4.1-zh}
  \widehat{M}=M\otimes_S\widehat{S}=M\otimes_S S[z]
\tag{8.2.4.1}
\end{equation}
是各 $S$-模 $M\otimes_S Sz^n$ 的直和，因而也是各阿贝尔群
$M_k\otimes_S Sz^n$（$k\in\mathbf{Z}$，$n\geq0$）的直和；在
$\widehat{M}$ 上赋予分次 $\widehat{S}$-模结构，规定
\begin{equation}
\label{II.8.2.4.2-zh}
  \deg(x\otimes z^n)=n+\deg x
\tag{8.2.4.2}
\end{equation}
对 $M$ 中每个齐次 $x$ 成立。我们把证明
\egaref{II.8.2.3-zh}{8.2.3} 的类似结果留给读者：
\end{env}

\begin{lemma}[8.2.5]
\label{II.8.2.5-zh}
\begin{enumerate}
\item[\rm(i)] 存在（非分次）模的典范双同构
\begin{equation}
\label{II.8.2.5.1-zh}
  \widehat{M}_{(z)}\xrightarrow{\sim}M.
\tag{8.2.5.1}
\end{equation}
\item[\rm(ii)] 对每个 $f\in S_d$（$d>0$），存在（非分次）模的双同构
\begin{equation}
\label{II.8.2.5.2-zh}
  \widehat{M}_{(f)}\xrightarrow{\sim}M_f^{\leq}.
\tag{8.2.5.2}
\end{equation}
\end{enumerate}
\end{lemma}
% Locked French authority: ega2-1-fr.tex, 820505 B,
% SHA-256 84EDBE3E83530AF2959B441796337C9DC21EAFCA6A13114A26778760FBF437AC.
% Sequential source scope: lines 13974-14021, II.8.2.6-II.8.2.7.

\begin{env}[8.2.6]
\label{II.8.2.6-zh}
设 $S$ 是一个次数为\emph{正}的分次环，并考察 $S$ 中如下递减分次
理想列
\begin{equation}
\label{II.8.2.6.1-zh}
  S_{[n]}=\bigoplus_{m\geq n}S_m\qquad(n\geq0)
\tag{8.2.6.1}
\end{equation}
（特别地，有 $S_{[0]}=S$，$S_{[1]}=S_+$）。由于显然有
$S_{[m]}S_{[n]}\subset S_{[m+n]}$，故可取
\begin{equation}
\label{II.8.2.6.2-zh}
  S^\natural=\bigoplus_{n\geq0}S_n^\natural
  \quad\text{且}\quad S_n^\natural=S_{[n]}.
\tag{8.2.6.2}
\end{equation}
来定义一个\emph{分次环} $S^\natural$。因此，$S_0^\natural$ 是视为
\emph{非分次环}的环 $S$，从而 $S^\natural$ 是一个
$S_0^\natural$-代数。对于任意齐次元 $f\in S_d$（$d>0$），我们以
$f^\natural$ 表示元素 $f$，这里把它视为属于 $S_{[d]}=S_d^\natural$。采用
这些记号：
\end{env}

\oldpage[II]{159}

\begin{lemma}[8.2.7]
\label{II.8.2.7-zh}
设 $S$ 是一个次数为正的分次环，$f$ 是 $S_d$ 的一个齐次元
（$d>0$）。有如下典范环同构
\begin{equation}
\label{II.8.2.7.1-zh}
  S_f\xrightarrow{\sim}\bigoplus_{n\in\mathbf{Z}}S(n)_{(f)}
\tag{8.2.7.1}
\end{equation}
\begin{equation}
\label{II.8.2.7.2-zh}
  (S_f^{\geq})_{f/1}\xrightarrow{\sim}S_f
\tag{8.2.7.2}
\end{equation}
\begin{equation}
\label{II.8.2.7.3-zh}
  S^\natural_{(f^\natural)}\xrightarrow{\sim}S_f^{\geq}
\tag{8.2.7.3}
\end{equation}
其中前两个是分次环同构。
\end{lemma}
% Locked French authority: ega2-1-fr.tex, 820505 B,
% SHA-256 84EDBE3E83530AF2959B441796337C9DC21EAFCA6A13114A26778760FBF437AC.
% Sequential source scope: lines 14023-14077, proof of II.8.2.7 through II.8.2.9.

由定义立即有 $(S_f)_n=(S(n)_f)_0$，因而得到同构
（\egaref{II.8.2.7.1-zh}{8.2.7.1}），它无非就是恒等映射。另一方面，
由于 $f/1$ 在 $S_f$ 中可逆，故有典范同构
$S_f\xrightarrow{\sim}(S_f^{\geq})_{f/1}=(S_f)_{f/1}$；这是把
（\egaref{II.8.2.1.2-zh}{8.2.1.2}）应用于 $S_f$ 所得；其逆同构依定义
就是同构（\egaref{II.8.2.7.2-zh}{8.2.7.2}）。最后，若
$x=\sum_{m\geq n}y_m$ 是 $S_{[n]}$ 的一个元素，且 $n=kd$，则把元素
$x/(f^\natural)^k$ 对应到元素 $\sum_m y_m/f^k$，后者属于
$S_f^{\geq}$；立即验证可知，这定义了一个同构
（\egaref{II.8.2.7.3-zh}{8.2.7.3}）。

\begin{env}[8.2.8]
\label{II.8.2.8-zh}
若 $M$ 是一个分次 $S$-模，则类似地，对每个 $n\in\mathbf{Z}$，置
\begin{equation}
\label{II.8.2.8.1-zh}
  M_{[n]}=\bigoplus_{m\geq n}M_m
\tag{8.2.8.1}
\end{equation}
而由于 $S_{[m]}M_{[n]}\subset M_{[m+n]}$（$m\geq0$），可定义一个分次
$S^\natural$-模 $M^\natural$，取
\begin{equation}
\label{II.8.2.8.2-zh}
  M^\natural=\bigoplus_{n\in\mathbf{Z}}M_n^\natural,
  \quad\text{且}\quad M_n^\natural=M_{[n]}.
\tag{8.2.8.2}
\end{equation}
\end{env}

我们仍把如下结果的证明留给读者：

\begin{lemma}[8.2.9]
\label{II.8.2.9-zh}
采用（\egaref{II.8.2.7-zh}{8.2.7}）与
（\egaref{II.8.2.8-zh}{8.2.8}）的记号，有如下模的典范双同构
\begin{equation}
\label{II.8.2.9.1-zh}
  M_f\xrightarrow{\sim}\bigoplus_{n\in\mathbf{Z}}M(n)_{(f)}
\tag{8.2.9.1}
\end{equation}
\begin{equation}
\label{II.8.2.9.2-zh}
  (M_f^{\geq})_{f/1}\xrightarrow{\sim}M_f
\tag{8.2.9.2}
\end{equation}
\begin{equation}
\label{II.8.2.9.3-zh}
  M^\natural_{(f^\natural)}\xrightarrow{\sim}M_f^{\geq}
\tag{8.2.9.3}
\end{equation}
其中前两个是分次模的双同构。
\end{lemma}
% Locked French authority: ega2-1-fr.tex, 820505 B,
% SHA-256 84EDBE3E83530AF2959B441796337C9DC21EAFCA6A13114A26778760FBF437AC.
% Sequential source scope: lines 14079-14125, Lemma II.8.2.10 and
% proof item (i) through its conclusion.

\begin{lemma}[8.2.10]
\label{II.8.2.10-zh}
设 $S$ 是正次数分次环。
\begin{enumerate}
\item[\rm(i)] $S^\natural$ 是有限型（相应地，诺特）
$S_0^\natural$-代数，当且仅当 $S$ 是有限型（相应地，诺特）
$S_0$-代数。
\item[\rm(ii)] 对 $n\geq n_0$ 有
$S_{n+1}^\natural=S_1^\natural S_n^\natural$，当且仅当对
$n\geq n_0$ 有 $S_{n+1}=S_1S_n$。
\item[\rm(iii)] 对 $n\geq n_0$ 有 $S_n^\natural=S_1^{\natural n}$，
当且仅当对 $n\geq n_0$ 有 $S_n=S_1^n$。
\item[\rm(iv)] 若 $(f_\alpha)$ 是 $S_+$ 的一族齐次元素，使 $S_+$
等于由各 $f_\alpha$ 生成的 $S_+$-理想在 $S_+$ 中的根，则
$S_+^\natural$ 等于由各 $f_\alpha^\natural$ 生成的
$S_+^\natural$-理想在 $S_+^\natural$ 中的根。
\end{enumerate}
\end{lemma}

\begin{enumerate}
\item[\rm(i)] 若 $S^\natural$ 是有限型 $S_0^\natural$-代数，则由
\egaref{II.2.1.6-zh}{2.1.6}，(i)，$S_+=S_1^\natural$ 是
$S=S_0^\natural$ 上的有限型模，故 $S$ 是有限型 $S_0$-代数
（\egaref{II.2.1.4-zh}{2.1.4}）；若 $S^\natural$ 是诺特环，则
$S_0^\natural=S$ 也同样是诺特环
（\egaref{II.2.1.5-zh}{2.1.5}）。反之，若 $S$ 是一个

\oldpage[II]{160}
有限型 $S_0$-代数，则已知（\egaref{II.2.1.6-zh}{2.1.6}，(ii)），
存在 $h>0$ 和 $m_0>0$，使得对 $n\geq m_0$ 有
$S_{n+h}=S_hS_n$；显然可以假设 $m_0\geq h$。此外，各 $S_m$ 都是
有限型 $S_0$-模（\egaref{II.2.1.6-zh}{2.1.6}，(i)）。在此情形下，
若 $n\geq m_0+h$，则
$S_n^\natural=S_hS_{n-h}^\natural
=S_h^\natural S_{n-h}^\natural$；若 $m<m_0+h$，置
$E=S_{m_0}+\cdots+S_{m_0+h-1}$，则
\[
  S_m^\natural=S_m+\cdots+S_{m_0+h-1}+S_hE+S_h^2E+\cdots.
\]
对 $1\leq m\leq m_0$，令 $G_m$ 为各 $S_0$-模 $S_i$
（$m\leq i\leq m_0+h-1$）的有限生成元系之并，并把它看作
$S_{[m]}$ 的子集。对 $m_0+1\leq m\leq m_0+h-1$，同样令 $G_m$
为各 $S_0$-模 $S_i$（$m\leq i\leq m_0+h-1$）以及 $S_hE$ 的
有限生成元系之并，并把它看作 $S_{[m]}$ 的子集。显然，对
$1\leq m\leq m_0+h-1$ 有 $S_m^\natural=S_0^\natural G_m$；因此，
对 $1\leq m\leq m_0+h-1$ 的各 $G_m$ 取并所得的 $G$，是
$S_0^\natural$-代数 $S^\natural$ 的一个生成元系。由此可知，若
$S=S_0^\natural$ 是诺特环，则 $S^\natural$ 也是诺特环。
% Locked French authority: ega2-1-fr.tex, 820505 B,
% SHA-256 84EDBE3E83530AF2959B441796337C9DC21EAFCA6A13114A26778760FBF437AC.
% Sequential source scope: lines 14127-14162, proof of II.8.2.10(ii)-(iv).

\item[\rm(ii)] 显然，若有 $S_{n+1}=S_1S_n$，其中 $n\geq n_0$，则
$S_{n+1}^\natural=S_1S_n^\natural$，并且\emph{更有}
$S_{n+1}^\natural=S_1^\natural S_n^\natural$ 对 $n\geq n_0$ 成立。
反之，后一关系可写成
\[
  S_{n+1}+S_{n+2}+\cdots=(S_1+S_2+\cdots)(S_n+S_{n+1}+\cdots)
\]
比较两边次数为 $n+1$ 的项（在 $S$ 中），便得 $S_{n+1}=S_1S_n$。

\item[\rm(iii)] 若有 $S_n=S_1^n$，其中 $n\geq n_0$，则
$S_n^\natural=S_1^n+S_1^{n+1}+\cdots$；由于 $S_1^\natural$ 包含
$S_1+S_1^2+\cdots$，故 $S_n^\natural\subset S_1^{\natural n}$，从而对
$S_n^\natural=S_1^{\natural n}$ 对 $n\geq n_0$ 成立。反之，
$S_1^{\natural n}=(S_1+S_2+\cdots)^n$ 中次数为 $n$ 的项（在 $S$ 中），
只有 $S_1^n$ 中的项；因此，关系 $S_n^\natural=S_1^{\natural n}$ 蕴含
$S_n=S_1^n$。

\item[\rm(iv)] 只需证明：若把元素 $g\in S_{k+h}$ 视为
$S_k^\natural$ 的元素（$k>0$，$h\geq0$），则存在整数 $n>0$，使得在
$S_{kn}^\natural$ 中，$g^n$ 是诸 $f_\alpha^\natural$ 的线性组合，
其系数属于 $S^\natural$。依假设，存在整数 $m_0$，使得当
$m\geq m_0$ 时，在\emph{$S$ 中}有
$g^m=\sum_\alpha c_{\alpha m}f_\alpha$，而出现在此公式中的指标
$\alpha$ \emph{与 $m$ 无关}；此外，显然可假设 $c_{\alpha m}$ 都是
齐次的，并且
\[
  \deg(c_{\alpha m})=m(k+h)-\deg f_\alpha
\]
在 $S$ 中成立。于是把 $m_0$ 取得充分大，使得不等式
$km_0>\deg f_\alpha$ 对所有 $f_\alpha$ 都成立，而它们出现在
$g^{m_0}$ 中；对每个
$\alpha$，令 $c'_{\alpha m}$ 表示元素 $c_{\alpha m}$，并把它视为次数为
$km-\deg(f_\alpha)$ 的元素（在 $S^\natural$ 中）；于是，在
$S^\natural$ 中有 $g^m=\sum_\alpha c'_{\alpha m}f_\alpha^\natural$，
证明完毕。
\end{enumerate}
% Locked French authority: ega2-1-fr.tex, 820505 B,
% SHA-256 84EDBE3E83530AF2959B441796337C9DC21EAFCA6A13114A26778760FBF437AC.
% Sequential source scope: lines 14164-14225, Environment II.8.2.11
% through the proof of Proposition II.8.2.12.

\begin{env}[8.2.11]
\label{II.8.2.11-zh}
考虑分次 $S_0$-代数
\begin{equation}
\label{II.8.2.11.1-zh}
  S^\natural\otimes_S S_0=S^\natural/S_+S^\natural
  =\bigoplus_{n\geq0}S_{[n]}/S_+S_{[n]}.
\tag{8.2.11.1}
\end{equation}
由于 $S_n$ 是 $S_{[n]}/S_+S_{[n]}$ 的商 $S_0$-模，故有分次
$S_0$-代数的典范同态
\begin{equation}
\label{II.8.2.11.2-zh}
  S^\natural\otimes_S S_0\to S
\tag{8.2.11.2}
\end{equation}
它显然是\emph{满射的}，因而由 \egaref{II.2.9.2-zh}{2.9.2} 对应于
典范\emph{闭浸入}
\begin{equation}
\label{II.8.2.11.3-zh}
  \operatorname{Proj}(S)\to
  \operatorname{Proj}(S^\natural\otimes_S S_0).
\tag{8.2.11.3}
\end{equation}
\end{env}

\oldpage[II]{161}

\begin{proposition}[8.2.12]
\label{II.8.2.12-zh}
典范态射 \egaref{II.8.2.11.3-zh}{8.2.11.3} 是双射。要使同态
\egaref{II.8.2.11.2-zh}{8.2.11.2} 是 (TN)-双射的，必要且充分的是
存在 $n_0$，使得对 $n\geq n_0$ 有 $S_{n+1}=S_1S_n$。若后一条件
成立，则 \egaref{II.8.2.11.3-zh}{8.2.11.3} 是同构；当 $S$ 是诺特环
时，其逆命题也成立。
\end{proposition}

为证明第一个断言，只需（\egaref{II.2.8.3-zh}{2.8.3}）证明同态
\egaref{II.8.2.11.2-zh}{8.2.11.2} 的核 $\mathfrak{I}$ 由
\emph{幂零元素}组成。事实上，若 $f\in S_{[n]}$ 的模
$S_+S_{[n]}$ 类属于该核，则这意味着 $f\in S_{[n+1]}$；此时把元素
$f^{n+1}$ 看作属于 $S_{[n(n+1)]}$，它属于
$S_+S_{[n(n+1)]}$，因为它可写成 $f\cdot f^n$；因此
$f^{n+1}$ 的模 $S_+S_{[n(n+1)]}$ 类为零，这证明了我们的断言。
由于对 $n\geq n_0$ 有 $S_{n+1}=S_1S_n$ 这一假设，等价于对
$n\geq n_0$ 有
$S_{n+1}^{\natural}=S_1^{\natural}S_n^{\natural}$
（\egaref{II.8.2.10-zh}{8.2.10}，(ii)），故依定义，该假设等价于
\egaref{II.8.2.11.2-zh}{8.2.11.2} 是 (TN)-单射的，因而是
(TN)-双射的；于是由 \egaref{II.2.9.1-zh}{2.9.1}，
\egaref{II.8.2.11.3-zh}{8.2.11.3} 是同构。反之，若
\egaref{II.8.2.11.3-zh}{8.2.11.3} 是同构，则
$\operatorname{Proj}(S^{\natural}\otimes_S S_0)$ 上的层
$\widetilde{\mathfrak{I}}$ 为零
（\egaref{II.2.9.2-zh}{2.9.2}，(i)）。由于
$S^{\natural}\otimes_S S_0$ 作为 $S^{\natural}$ 的商是诺特环
（\egaref{II.8.2.10-zh}{8.2.10}，(i)），由
\egaref{II.2.7.3-zh}{2.7.3} 可知 $\mathfrak{I}$ 满足条件 (TN)，
故对 $n\geq n_0$ 有
$S_{n+1}^{\natural}=S_1^{\natural}S_n^{\natural}$；再由
\egaref{II.8.2.10-zh}{8.2.10}，(ii)，证明完成。
% Locked French authority: ega2-1-fr.tex, 820505 B,
% SHA-256 84EDBE3E83530AF2959B441796337C9DC21EAFCA6A13114A26778760FBF437AC.
% Sequential source scope: lines 14227-14267, II.8.2.13-II.8.2.14.1.

\begin{env}[8.2.13]
\label{II.8.2.13-zh}
现在考虑典范单射 $(S_+)^n\to S_{[n]}$；它们定义了一个次数为 $0$ 的
分次环单射同态
\begin{equation}
\label{II.8.2.13.1-zh}
  \bigoplus_{n\geq0}(S_+)^n\to S^{\natural}.
\tag{8.2.13.1}
\end{equation}
\end{env}

\begin{proposition}[8.2.14]
\label{II.8.2.14-zh}
同态（\egaref{II.8.2.13.1-zh}{8.2.13.1}）为 (TN)-同构的充要条件是：
存在 $n_0$，使得 $S_n=S_1^n$ 对每个 $n\geq n_0$ 都成立。在这种情形，
与（\egaref{II.8.2.13.1-zh}{8.2.13.1}）对应的态射处处有定义，并且是同构
\[
  \operatorname{Proj}(S^{\natural})\xrightarrow{\sim}
  \operatorname{Proj}\left(\bigoplus_{n\geq0}(S_+)^n\right)~;
\]
当 $S$ 是诺特环时，逆命题也成立。
\end{proposition}

由（\egaref{II.8.2.10-zh}{8.2.10}, (iii)）与
（\egaref{II.2.9.1-zh}{2.9.1}），前两个断言是显然的。第三个断言由
（\egaref{II.8.2.10-zh}{8.2.10}, (i) 和 (iii)）和下一引理推出：

\begin{lemma}[8.2.14.1]
\label{II.8.2.14.1-zh}
设 $T$ 是一个次数为正的分次环，并且是一个有限型 $T_0$-代数。若与单射同态
$\bigoplus_{n\geq0}T_1^n\to T$ 对应的态射处处有定义，并且是同构
\[
  \operatorname{Proj}(T)\to
  \operatorname{Proj}\left(\bigoplus_{n\geq0}T_1^n\right),
\]
则存在 $n_0$，使得 $T_n=T_1^n$ 对 $n\geq n_0$ 成立。
\end{lemma}
% Locked French authority: ega2-1-fr.tex, 820505 B,
% SHA-256 84EDBE3E83530AF2959B441796337C9DC21EAFCA6A13114A26778760FBF437AC.
% Sequential source scope: lines 14269-14295, proof of Lemma II.8.2.14.1.

事实上，设 $g_i$（$1\leq i\leq r$）是 $T_0$-模 $T_1$ 的生成元。
该假设首先蕴含各 $D_+(g_i)$ 覆盖 $\operatorname{Proj}(T)$
（\egaref{II.2.8.1-zh}{2.8.1}）。设 $(h_j)_{1\leq j\leq s}$ 是
$T_+$ 的一族齐次元素，$\deg(h_j)=n_j$，它们与各 $g_i$ 一起组成
理想 $T_+$ 的一个生成元系，或者等价地
（\egaref{II.2.1.3-zh}{2.1.3}），组成 $T$ 作为 $T_0$-代数的一个
生成元系。若置 $T'=\bigoplus_{n\geq0}T_1^n$，则依假设，环
$T_{(g_i)}$ 的元素 $h_j/g_i^{n_j}$ 必属于子环 $T'_{(g_i)}$；因此
存在整数 $k$，使得对每个 $j$ 都有
$T_1^k h_j\subset T_1^{k+n_j}$。对 $r$ 作归纳可得，对每个
$r\geq1$，都有 $T_1^k h_j^r\subset T'$；再由各 $h_j$ 的定义，
便有 $T_1^kT\subset T'$。另一方面，对每个 $j$，存在整数 $m_j$，
使 $h_j^{m_j}$ 属于由各 $g_i$ 生成的 $T$ 的理想
（\egaref{II.2.3.14-zh}{2.3.14}），故 $h_j^{m_j}\in T_1T$，并且

\oldpage[II]{162}
$h_j^{m_jk}\in T_1^kT\subset T'$。因此存在整数 $m_0\geq k$，使得
对 $m\geq m_0$ 有 $h_j^m\in T_1^{mn_j}$。此时，若 $q$ 是各整数
$n_j$ 中的最大者，则数 $n_0=qsm_0+k$ 满足要求。事实上，当
$n\geq n_0$ 时，$T_n$ 的一个元素是一些属于 $T_1^\alpha u$ 的
单项式之和，其中 $u$ 是各 $h_j$ 的幂的乘积；若 $\alpha\geq k$，
由上述结果可知 $T_1^\alpha u\subset T_1^n$。在相反情形中，各
$h_j$ 的指数中有一个是 $\geq m_0$，因而
$u\in T_1^\beta v$，其中 $\beta\geq k$，而 $v$ 仍是各 $h_j$ 的幂
的乘积；于是归结到前一种情形，并最终看到，在所有情形中都有
$T_1^\alpha u\subset T_1^n$。
% Locked French authority: ega2-1-fr.tex, 820505 B,
% SHA-256 84EDBE3E83530AF2959B441796337C9DC21EAFCA6A13114A26778760FBF437AC.
% Sequential source scope: lines 14297-14308, Remark II.8.2.15.

\begin{remark}[8.2.15]
\label{II.8.2.15-zh}
条件 $S_n=S_1^n$ 对 $n\geq n_0$ 成立，显然蕴含
$S_{n+1}=S_1S_n$ 对 $n\geq n_0$ 成立；但其逆命题不成立，即使假设
$S$ 是诺特环也是如此。例如，设 $K$ 是一个域，
$A=K[x]$，$B=K[y]/y^2K[y]$，其中 $x$ 和 $y$ 是两个不定元，规定
$x$ 的次数为 $1$，$y$ 的次数为 $2$；再设 $S=A\otimes_KB$，于是
$S$ 是 $K$ 上的分次代数，它有由元素 $1$、$x^n$（$n\geq1$）以及
$x^ny$（$n\geq0$）组成的一组基。显然，$S_{n+1}=S_1S_n$ 对
$n\geq2$ 成立，但 $S_1^n=Kx^n$，而 $S_n=Kx^n+Kx^{n-2}y$ 对
$n\geq2$ 成立。
\end{remark}

% Locked French authority: ega2-1-fr.tex, 820505 B,
% SHA-256 84EDBE3E83530AF2959B441796337C9DC21EAFCA6A13114A26778760FBF437AC.
% Sequential source scope: lines 14310-14352, subsection II.8.3 and II.8.3.1.

\subsection{投影锥}
\label{subsection:II.8.3-zh}

\begin{env}[8.3.1]
\label{II.8.3.1-zh}
设 $Y$ 为一个预概形；在本编号中，只考虑 \emph{$Y$-预概形} 和
\emph{$Y$-态射}。设 $\mathcal{S}$ 为一个按 \emph{正次数} 分次的拟凝聚
$\mathcal{O}_Y$-代数；\emph{此外，假设 $\mathcal{S}_0=\mathcal{O}_Y$}。
按照（\egaref{II.8.2.2-zh}{8.2.2}）中引入的记号，置
\begin{equation}
\label{II.8.3.1.1-zh}
  \widehat{\mathcal{S}}=\mathcal{S}[z]
  =\mathcal{S}\otimes_{\mathcal{O}_Y}\mathcal{O}_Y[z]
\tag{8.3.1.1}
\end{equation}
并将其视为一个按正次数分次的 $\mathcal{O}_Y$-代数，其中次数由公式
（\egaref{II.8.2.2.2-zh}{8.2.2.2}）定义。因此，对任一仿射开集 $U$
（视为 $Y$ 的开集），有
\[
  \Gamma(U,\widehat{\mathcal{S}})=(\Gamma(U,\mathcal{S}))[z].
\]
下文置
\begin{equation}
\label{II.8.3.1.2-zh}
  X=\operatorname{Proj}(\mathcal{S}),\qquad
  C=\operatorname{Spec}(\mathcal{S}),\qquad
  \widehat{C}=\operatorname{Proj}(\widehat{\mathcal{S}})
\tag{8.3.1.2}
\end{equation}
（在 $C$ 的定义中，把 $\mathcal{S}$ 看成非分次的 $\mathcal{O}_Y$-代数），
并称 $C$（相应地称 $\widehat{C}$）为由 $\mathcal{S}$ 定义的
\emph{仿射锥}（相应地为 \emph{射影锥}）；也用“锥”代替“仿射锥”。滥用语言，
仍称 $C$（相应地称 $\widehat{C}$）为 \emph{$X$ 的仿射投影锥}
（相应地为 \emph{$X$ 的射影投影锥}），这里默认预概形 $X$ 以
$\operatorname{Proj}(\mathcal{S})$ 的形式给出；最后，称 $\widehat{C}$ 为
$C$ 的 \emph{泛射影闭包}（这里理解为 $\mathcal{S}$ 的数据是 $C$ 的结构的
一部分）。
\end{env}
% Locked French authority: ega2-1-fr.tex, 820505 B,
% SHA-256 84EDBE3E83530AF2959B441796337C9DC21EAFCA6A13114A26778760FBF437AC.
% Sequential source scope: lines 14354-14404, Proposition II.8.3.2 and
% construction of i, epsilon, and j.

\begin{proposition}[8.3.2]
\label{II.8.3.2-zh}
存在典范的 $Y$-态射
\begin{equation}
\label{II.8.3.2.1-zh}
  Y\xrightarrow{\varepsilon}C\xrightarrow{i}\widehat{C}
\tag{8.3.2.1}
\end{equation}
\begin{equation}
\label{II.8.3.2.2-zh}
  X\xrightarrow{j}\widehat{C}
\tag{8.3.2.2}
\end{equation}

\oldpage[II]{163}
使得 $\varepsilon$ 和 $j$ 都是闭浸入，而 $i$ 是仿射态射，并且是支配
开浸入，满足
\begin{equation}
\label{II.8.3.2.3-zh}
  i(C)=\widehat{C}-j(X)~;
\tag{8.3.2.3}
\end{equation}
此外，$\widehat{C}$ 是 $\widehat{C}$ 中包含 $i(C)$ 的最小闭子预概形。
\end{proposition}

为定义 $i$，考察 $\widehat{C}$ 的开集
\begin{equation}
\label{II.8.3.2.4-zh}
  \widehat{C}_z=
  \operatorname{Spec}(\widehat{\mathcal{S}}/(z-1)\widehat{\mathcal{S}})
\tag{8.3.2.4}
\end{equation}
（\egaref{II.3.1.4-zh}{3.1.4}），其中 $z$ 可典范地等同于
$\widehat{\mathcal{S}}$ 在 $Y$ 上的一个截面。于是，同构
$i:C\xrightarrow{\sim}\widehat{C}_z$ 对应于典范同构
（\egaref{II.8.2.3.1-zh}{8.2.3.1}）
\[
  \widehat{\mathcal{S}}/(z-1)\widehat{\mathcal{S}}
  \xrightarrow{\sim}\mathcal{S}.
\]
态射 $\varepsilon$ 对应于增广同态
$\mathcal{S}\to\mathcal{S}_0=\mathcal{O}_Y$，其核为
$\mathcal{S}_+$（\egaref{II.1.2.7-zh}{1.2.7}）；由于该同态是满射，
$\varepsilon$ 是闭浸入（\egaref{II.1.4.10-zh}{1.4.10}）。最后，$j$
同样对应于（\egaref{II.3.5.1-zh}{3.5.1}）次数为 $0$ 的满同态
$\widehat{\mathcal{S}}\to\mathcal{S}$；它在 $\mathcal{S}$ 上限制为恒等，
并在 $z\widehat{\mathcal{S}}$ 上为零，而后者正是它的核；$j$ 处处有
定义，并且由（\egaref{II.3.6.2-zh}{3.6.2}）可知它是闭浸入。
% Locked French authority: ega2-1-fr.tex, 820505 B,
% SHA-256 84EDBE3E83530AF2959B441796337C9DC21EAFCA6A13114A26778760FBF437AC.
% Sequential source scope: lines 14406-14449, continuation of the proof of
% Proposition II.8.3.2 through Environment II.8.3.2.6.

为证明 \egaref{II.8.3.2-zh}{8.3.2} 的其他断言，显然可以归结到
$Y=\operatorname{Spec}(A)$ 为仿射概形且 $\mathcal{S}=\widetilde{S}$
的情形，其中 $S$ 是分次 $A$-代数，从而
$\widehat{\mathcal{S}}=(\widehat{S})^{\sim}$。此时，$S_+$ 的齐次元素
$f$ 可等同于 $Y$ 上 $\widehat{\mathcal{S}}$ 的截面，而
\egaref{II.2.3.3-zh}{2.3.3} 中记为 $D_+(f)$ 的 $\widehat{C}$ 的开集，
也写作 $\widehat{C}_f$（\egaref{II.3.1.4-zh}{3.1.4}）；同样，
\egaxref{I.1.1.1-zh}{I, 1.1.1} 中记为 $D(f)$ 的 $C$ 的开集，也写作
$C_f$（\egaxref{0.5.5.2-zh}{0, 5.5.2}）。在此情形下，由
\egaref{II.2.3.14-zh}{2.3.14} 和 $\widehat{S}$ 的定义可知，开集
$\widehat{C}_z=i(C)$ 以及各 $\widehat{C}_f$（$f$ 是 $S_+$ 中的
齐次元素）构成 $\widehat{C}$ 的一个\emph{覆盖}。此外，在这些记号下，
有
\begin{equation}
\label{II.8.3.2.5-zh}
  i^{-1}(\widehat{C}_f)=C_f~;
\tag{8.3.2.5}
\end{equation}
事实上，
$\widehat{C}_f\cap i(C)=\widehat{C}_f\cap\widehat{C}_z
=\widehat{C}_{fz}=\operatorname{Spec}(\widehat{S}_{(fz)})$。现在，若
$d=\deg(f)$，则 $\widehat{S}_{(fz)}$ 典范同构于
$(\widehat{S}_{(z)})_{f/z^d}$（\egaref{II.2.2.2-zh}{2.2.2}）；由同构
\egaref{II.8.2.3.1-zh}{8.2.3.1} 的定义可知，
$(\widehat{S}_{(z)})_{f/z^d}$ 在相应分式环同构下的像恰为 $S_f$。
由于 $C_f=\operatorname{Spec}(S_f)$，这就证明了
\egaref{II.8.3.2.5-zh}{8.3.2.5}，同时也表明态射 $i$ 是仿射的。
此外，把 $i$ 在 $C_f$ 上的限制看作到 $\widehat{C}_f$ 的态射时，
它对应于典范同态 $\widehat{S}_{(f)}\to\widehat{S}_{(fz)}$
（\egaxref{I.1.7.3-zh}{I, 1.7.3}）；由上述结果和
\egaref{II.8.2.3.2-zh}{8.2.3.2}，可以表述如下结果：

\begin{env}[8.3.2.6]
\label{II.8.3.2.6-zh}
若 $Y=\operatorname{Spec}(A)$ 是仿射概形，且
$\mathcal{S}=\widetilde{S}$，则对 $S_+$ 中每个齐次 $f$，
$\widehat{C}_f$ 典范等同于 $\operatorname{Spec}(S_f^{\leq})$；此时，
$i$ 的限制态射 $C_f\to\widehat{C}_f$ 对应于典范单射
$S_f^{\leq}\to S_f$。
\end{env}
% Locked French authority: ega2-1-fr.tex, 820505 B,
% SHA-256 84EDBE3E83530AF2959B441796337C9DC21EAFCA6A13114A26778760FBF437AC.
% Sequential source scope: lines 14451-14485, completion of the proof of
% Proposition II.8.3.2, including Lemma II.8.3.2.7.

现在注意（当 $Y$ 仿射时），$\widehat{C}=\operatorname{Proj}(\widehat{S})$
中 $\widehat{C}_z$ 的补集

\oldpage[II]{164}
依定义是 $\widehat{S}$ 中包含 $z$ 的分次素理想所成的集合；由 $j$
的定义，这正是 $j(X)$，从而证明了
\egaref{II.8.3.2.3-zh}{8.3.2.3}。

最后，为证明 \egaref{II.8.3.2-zh}{8.3.2} 的最后一个断言，仍可假设
$Y$ 仿射。沿用前面的记号，注意在环 $\widehat{S}$ 中，$z$ 不是
$0$ 的因子；由于 $\overline{i(C)}=\widehat{C}$，只需证明下述引理。

\begin{lemma}[8.3.2.7]
\label{II.8.3.2.7-zh}
设 $T$ 是正次数分次环，$Z=\operatorname{Proj}(T)$，$g$ 是 $T$ 中
次数为 $d>0$ 的齐次元素。若 $g$ 在 $T$ 中不是 $0$ 的因子，则 $Z$
是 $Z$ 中包含 $Z_g=D_+(g)$ 的最小闭子预概形。
\end{lemma}

由 \egaxref{I.4.1.9-zh}{I, 4.1.9}，问题在 $Z$ 上是局部的；因此，
对每个齐次元素 $h\in T_e$（$e>0$），只需证明 $Z_h$ 是 $Z_h$ 中
包含 $Z_{gh}$ 的最小闭子预概形。由定义和
\egaxref{I.4.3.2-zh}{I, 4.3.2} 可知，该条件等价于典范同态
$T_{(h)}\to T_{(gh)}$ 是\emph{单射的}。现在，这个同态可等同于典范
同态 $T_{(h)}\to(T_{(h)})_{g^e/h^d}$
（\egaref{II.2.2.3-zh}{2.2.3}）。但由于 $g^e$ 在 $T$ 中不是
$0$ 的因子，$g^e/h^d$ 在 $T_h$ 中也不是零因子
（\emph{更不用说}在 $T_{(h)}$ 中）；因为若
$(g^e/h^d)(t/h^m)=0$，其中 $t\in T$ 且 $m>0$，则存在 $n>0$，使得
$h^ng^et=0$，从而 $h^nt=0$，因而 $t/h^m=0$ 在 $T_h$ 中成立。
这就完成了证明（\egaxref{0.1.2.2-zh}{0, 1.2.2}）。
% Locked French authority: ega2-1-fr.tex, 820505 B,
% SHA-256 84EDBE3E83530AF2959B441796337C9DC21EAFCA6A13114A26778760FBF437AC.
% Sequential source scope: lines 14487-14508, Environment II.8.3.3.

\begin{env}[8.3.3]
\label{II.8.3.3-zh}
常把仿射锥 $C$ 与射影锥 $\widehat{C}$ 在开集 $i(C)$ 上诱导的子预概形
等同起来，所借助的是开浸入 $i$。$C$ 的、与闭浸入 $\varepsilon$ 相应
的闭子预概形称为\emph{$C$ 的顶点预概形}；又说，$\varepsilon$ 是一个
$Y$-截面，其取值于 $C$，并称其为\emph{顶点截面}或\emph{$C$ 的零截面}；
可把 $Y$ 与 $C$ 的顶点预概形等同起来，所借助的是 $\varepsilon$。此外，
$i\circ\varepsilon$ 是一个 $Y$-截面，其取值于 $\widehat{C}$，因而也是闭浸入
（\egaxref{I.5.4.6-zh}{I, 5.4.6}）；它对应于次数为 $0$ 的典范满同态
$\widehat{\mathcal{S}}=\mathcal{S}[z]\to\mathcal{O}_Y[z]$
（参见（\egaref{II.3.1.7-zh}{3.1.7}）），其核为
$\mathcal{S}_+[z]=\mathcal{S}_+\widehat{\mathcal{S}}$。与该闭浸入相应
的 $\widehat{C}$ 的子预概形仍称为\emph{$\widehat{C}$ 的顶点预概形}，
并称 $i\circ\varepsilon$ 为\emph{$\widehat{C}$ 的顶点截面}；借助
$i\circ\varepsilon$，可把它与 $Y$ 等同起来。最后，$\widehat{C}$ 的、
与 $j$ 相应的闭子预概形称为\emph{$\widehat{C}$ 的无穷远轨迹}，
并可与 $X$ 等同起来，所借助的是 $j$。
\end{env}
% Locked French authority: ega2-1-fr.tex, 820505 B,
% SHA-256 84EDBE3E83530AF2959B441796337C9DC21EAFCA6A13114A26778760FBF437AC.
% Sequential source scope: lines 14510-14558, Environment II.8.3.4.

\begin{env}[8.3.4]
\label{II.8.3.4-zh}
分别由 $C$（相应地由 $\widehat{C}$）在下列 \emph{开集} 上诱导的子预概形
\begin{equation}
\label{II.8.3.4.1-zh}
  E=C-\varepsilon(Y),\qquad
  \widehat{E}=\widehat{C}-i(\varepsilon(Y))
\tag{8.3.4.1}
\end{equation}
分别（滥用语言）称为 \emph{去顶点仿射锥} 和 \emph{去顶点射影锥}，它们由
$\mathcal{S}$ 定义；应当注意，尽管使用了这一术语，$E$ 在 $Y$ 上
\emph{不一定是仿射的}，而 $\widehat{E}$ 在 $Y$ 上也不一定是射影的
（参见（\egaref{II.8.4.3-zh}{8.4.3}））。将 $C$ 等同于 $i(C)$ 时，
在底空间上有
\begin{equation}
\label{II.8.3.4.2-zh}
  C\cup\widehat{E}=\widehat{C},\qquad C\cap\widehat{E}=E
\tag{8.3.4.2}
\end{equation}
因此可以把 $\widehat{C}$ 看作对开子预概形 $C$ 与 $\widehat{E}$ 作
\emph{粘合} 而得到；此外，由（\egaref{II.8.3.2.3-zh}{8.3.2.3}），
\begin{equation}
\label{II.8.3.4.3-zh}
  E=\widehat{E}-j(X).
\tag{8.3.4.3}
\end{equation}

\oldpage[II]{165}
当 $Y=\operatorname{Spec}(A)$ 为仿射的时，采用
（\egaref{II.8.3.2-zh}{8.3.2}）的记号，有
\begin{equation}
\label{II.8.3.4.4-zh}
  E=\bigcup C_f,\qquad
  \widehat{E}=\bigcup\widehat{C}_f,\qquad
  C_f=C\cap\widehat{C}_f
\tag{8.3.4.4}
\end{equation}
其中 $f$ 遍历 $S_+$ 的齐次元素集合（也可以只遍历这个集合的一个子集 $M$，
它生成 $S_+$ 的一个理想，该理想在 $S_+$ 中的根理想为 $S_+$ 本身；换言之，
使得诸 $X_f$（其中 $f\in M$）覆盖 $X$
（\egaref{II.2.3.14-zh}{2.3.14}））。于是，$C$ 与 $\widehat{C}_f$ 沿
$C_f$ 的粘合由单射态射 $C_f\to C$、$C_f\to\widehat{C}_f$ 所确定；
正如（\egaref{II.8.3.2.6-zh}{8.3.2.6}）所见，它们分别对应于典范同态
$S\to S_f$、$S_f^{\leq}\to S_f$。
\end{env}
% Locked French authority: ega2-1-fr.tex, 820505 B,
% SHA-256 84EDBE3E83530AF2959B441796337C9DC21EAFCA6A13114A26778760FBF437AC.
% Sequential source scope: lines 14560-14610, Proposition II.8.3.5 and
% its proof, including Environment II.8.3.5.3.

\begin{proposition}[8.3.5]
\label{II.8.3.5-zh}
沿用 \egaref{II.8.3.1-zh}{8.3.1} 和 \egaref{II.8.3.4-zh}{8.3.4}
的记号，与典范单射
$\varphi:\mathcal{S}\to\widehat{\mathcal{S}}=\mathcal{S}[z]$ 相伴的态射
（\egaref{II.3.5.1-zh}{3.5.1}）是仿射满射态射（称为\emph{典范收缩}）
\begin{equation}
\label{II.8.3.5.1-zh}
  p:\widehat{E}\to X
\tag{8.3.5.1}
\end{equation}
并且有
\begin{equation}
\label{II.8.3.5.2-zh}
  p\circ j=1_X.
\tag{8.3.5.2}
\end{equation}
\end{proposition}

为证明该命题，可以归结到 $Y$ 仿射的情形。考虑到
\egaref{II.8.3.4.4-zh}{8.3.4.4} 给出的 $\widehat{E}$ 的表达式，
$p$ 的定义域 $G(\varphi)$ 等于 $\widehat{E}$ 这一事实，将由下列断言
中的第一个得出：

\begin{env}[8.3.5.3]
\label{II.8.3.5.3-zh}
若 $Y=\operatorname{Spec}(A)$ 是仿射概形，并且
$\mathcal{S}=\widetilde{S}$，则对每个齐次 $f\in S_+$，有
\begin{equation}
\label{II.8.3.5.4-zh}
  p^{-1}(X_f)=\widehat{C}_f
\tag{8.3.5.4}
\end{equation}
而 $p$ 在 $\widehat{C}_f=\operatorname{Spec}(S_f^{\leq})$ 上的限制，
看作从 $\widehat{C}_f$ 到 $X_f$ 的态射时，对应于典范单射
$S_{(f)}\to S_f^{\leq}$。若还有 $f\in S_1$，则 $\widehat{C}_f$
同构于 $X_f\otimes_{\mathbf{Z}}\mathbf{Z}[T]$（$T$ 是不定元）。
\end{env}

事实上，公式 \egaref{II.8.3.5.4-zh}{8.3.5.4} 只是
\egaref{II.2.8.1.1-zh}{2.8.1.1} 的一个特例，而第二个断言正是当 $Y$
仿射时 $\operatorname{Proj}(\varphi)$ 的定义
（\egaref{II.2.8.1-zh}{2.8.1}）。此时，公式
\egaref{II.8.3.5.2-zh}{8.3.5.2} 以及 $p$ 为满射这一事实，来自典范
同态的复合 $\mathcal{S}\to\widehat{\mathcal{S}}\to\mathcal{S}$ 是
$\mathcal{S}$ 上的恒等映射。最后，\egaref{II.8.3.5.3-zh}{8.3.5.3}
的最后一个断言来自如下事实：当 $f\in S_1$ 时，$S_f^{\leq}$ 同构于
$S_{(f)}[T]$（\egaref{II.2.2.1-zh}{2.2.1}）。
% Locked French authority: ega2-1-fr.tex, 820505 B,
% SHA-256 84EDBE3E83530AF2959B441796337C9DC21EAFCA6A13114A26778760FBF437AC.
% Sequential source scope: lines 14612-14640, Corollary II.8.3.6 and proof.

\begin{corollary}[8.3.6]
\label{II.8.3.6-zh}
把 $p$ 限制到 $E$ 所得的限制态射
\begin{equation}
\label{II.8.3.6.1-zh}
  \pi:E\to X
\tag{8.3.6.1}
\end{equation}
是满射的仿射态射。若 $Y$ 是仿射的，且 $f$ 是 $S_+$ 中的齐次元，则有
\begin{equation}
\label{II.8.3.6.2-zh}
  \pi^{-1}(X_f)=C_f
\tag{8.3.6.2}
\end{equation}
并且 $\pi$ 在 $C_f$ 上的限制对应于典范单射 $S_{(f)}\to S_f$。此外，
若 $f\in S_1$，则 $C_f$ 同构于
$X_f\otimes_{\mathbf{Z}}\mathbf{Z}[T,T^{-1}]$（$T$ 为不定元）。
\end{corollary}

公式（\egaref{II.8.3.6.2-zh}{8.3.6.2}）立即由
（\egaref{II.8.3.5.3-zh}{8.3.5.3}）和
（\egaref{II.8.3.2.5-zh}{8.3.2.5}）得出，并证明了 $\pi$ 的满性；
另一方面，已经看到，浸入 $i$ 限制到 $C_f$ 后，对应于

\oldpage[II]{166}
单射 $S_f^{\leq}\to S_f$（\egaref{II.8.3.2-zh}{8.3.2}）。最后，末一
断言由以下事实得出：对 $f\in S_1$，$S_f$ 同构于
$S_{(f)}[T,T^{-1}]$（\egaref{II.2.2.1-zh}{2.2.1}）。
% Locked French authority: ega2-1-fr.tex, 820505 B,
% SHA-256 84EDBE3E83530AF2959B441796337C9DC21EAFCA6A13114A26778760FBF437AC.
% Sequential source scope: lines 14642-14662, Remark II.8.3.7.

\begin{remark}[8.3.7]
\label{II.8.3.7-zh}
当 $Y$ 仿射时，由浸入 $\varepsilon$ 的定义
（\egaref{II.8.3.2-zh}{8.3.2}），$E$ 的底空间的点是 $S$ 中不包含
$S_+$ 的素理想 $\mathfrak{p}$（未必分次）。对于这样的理想
$\mathfrak{p}$，各 $\mathfrak{p}\cap S_n$ 显然满足
\egaref{II.2.1.9-zh}{2.1.9} 的条件，因此存在唯一的 $S$ 的
\emph{分次}素理想 $\mathfrak{q}$，使得对每个 $n$ 都有
$\mathfrak{q}\cap S_n=\mathfrak{p}\cap S_n$；此时，底空间之间的映射
$\pi:E\to X$ 可以借助下列关系来解释：
\begin{equation}
\label{II.8.3.7.1-zh}
  \pi(\mathfrak{p})=\mathfrak{q}.
\tag{8.3.7.1}
\end{equation}
事实上，为验证该关系，只需考虑 $S_+$ 中满足
$\mathfrak{p}\in D(f)$ 的一个齐次 $f$，并注意到 $\mathfrak{q}_{(f)}$
是 $\mathfrak{p}_f$ 在单射 $S_{(f)}\to S_f$ 下的逆像。
\end{remark}
% Locked French authority: ega2-1-fr.tex, 820505 B,
% SHA-256 84EDBE3E83530AF2959B441796337C9DC21EAFCA6A13114A26778760FBF437AC.
% Sequential source scope: lines 14664-14675, Corollary II.8.3.8 and proof.

\begin{corollary}[8.3.8]
\label{II.8.3.8-zh}
若 $\mathcal{S}$ 由 $\mathcal{S}_1$ 生成，则态射 $p$ 和 $\pi$ 都是
有限型的；对每个 $x\in X$，纤维 $p^{-1}(x)$ 同构于
$\operatorname{Spec}(k(x)[T])$，而纤维 $\pi^{-1}(x)$ 同构于
$\operatorname{Spec}(k(x)[T,T^{-1}])$。
\end{corollary}

这立即由 \egaref{II.8.3.5-zh}{8.3.5} 和
\egaref{II.8.3.6-zh}{8.3.6} 得出；只需注意，当 $Y$ 仿射且 $S$ 由
$S_1$ 生成时，对 $f\in S_1$ 的各 $X_f$ 构成 $X$ 的一个覆盖
（\egaref{II.2.3.14-zh}{2.3.14}）。
% Locked French authority: ega2-1-fr.tex, 820505 B,
% SHA-256 84EDBE3E83530AF2959B441796337C9DC21EAFCA6A13114A26778760FBF437AC.
% Sequential source scope: lines 14677-14714, opening of Remark II.8.3.9.
% This bounded slice intentionally leaves the remark environment open.

\begin{remark}[8.3.9]
\label{II.8.3.9-zh}
与分次 $\mathcal{O}_Y$-代数 $\mathcal{O}_Y[T]$（$T$ 为不定元）对应的
去顶点仿射锥可等同于
$G_m=\operatorname{Spec}(\mathcal{O}_Y[T,T^{-1}])$，因为它就是 $C_T$，
如（\egaref{II.8.3.2-zh}{8.3.2}）所见（更一般的结果见
（\egaref{II.8.4.4-zh}{8.4.4}））。这个预概形典范地带有
“\emph{$Y$-交换群概形}”的结构。这一概念稍后将予以详述，此处可简要概括
如下。所谓 $Y$-群概形，是一个 $Y$-概形 $G$，配备两个 $Y$-态射
$p:G\times_YG\to G$ 和 $s:G\to G$；它们满足与群中合成律和取逆映射的
公理形式相同的条件：图
\[
  \xymatrix{
    G\times G\times G \ar[r]^{p\times1}\ar[d]_{1\times p}
      & G\times G \ar[d]^{p} \\
    G\times G \ar[r]_{p} & G
  }
\]
应当交换（“结合律”），并且还应满足一个条件，它对应于在群中下列两个映射
\[
  (x,y)\mapsto(x,x^{-1},y)\mapsto(x,x^{-1}y)\mapsto x(x^{-1}y)
\]
和
\[
  (x,y)\mapsto(x,x^{-1},y)\mapsto(x,yx^{-1})\mapsto(yx^{-1})x
\]
都化为 $(x,y)\mapsto y$ 这一事实；例如，与第一个复合映射对应的态射序列为
\[
  G\times G\xrightarrow{(1,s)\times1}G\times G\times G
  \xrightarrow{1\times p}G\times G\xrightarrow{p}G
\]
第二个序列类似地由读者写出。

% Locked French authority: ega2-1-fr.tex, 820505 B,
% SHA-256 84EDBE3E83530AF2959B441796337C9DC21EAFCA6A13114A26778760FBF437AC.
% Sequential source scope: lines 14715-14741, continuation of Remark II.8.3.9.

\oldpage[II]{167}
显然（\egaxref{I.3.4.3-zh}{I, 3.4.3}），给定一个 $Y$-群概形结构，
其底层是一个 $Y$-概形 $G$，等价于对\emph{任意} $Y$-预概形 $Z$，
在集合 $\operatorname{Hom}_Y(Z,G)$ 上给定一个\emph{群}结构；这些结构
须满足：对于任意 $Y$-态射 $Z\to Z'$，相应的映射
$\operatorname{Hom}_Y(Z',G)\to\operatorname{Hom}_Y(Z,G)$ 是群同态。在
这里所考察的 $G_m$ 这一特殊情形中，$\operatorname{Hom}_Y(Z,G)$ 可与
所有 $Z$-截面所成集合等同，这些截面属于 $Z\times_YG_m$
（\egaxref{I.3.3.14-zh}{I, 3.3.14}），因而可与所有 $Z$-截面所成集合
等同，这些截面属于 $\operatorname{Spec}(\mathcal{O}_Z[T,T^{-1}])$；
最后，与（\egaxref{I.3.3.15-zh}{I, 3.3.15}）中相同的论证表明，
此集合可典范地与环 $\Gamma(Z,\mathcal{O}_Z)$ 的\emph{可逆}元所成集合
等同，而该集合上的群结构正是由环 $\Gamma(Z,\mathcal{O}_Z)$ 中的乘法
给出的结构。读者可以验证，上面所考察的态射 $p$ 和 $s$ 是按如下方式
得到的：由（\egaref{II.1.2.7-zh}{1.2.7}）和
（\egaref{II.1.4.6-zh}{1.4.6}），它们对应于 $\mathcal{O}_Y$-代数同态
\begin{align*}
  \pi &: \mathcal{O}_Y[T,T^{-1}]
  \to\mathcal{O}_Y[T,T^{-1},T',T'^{-1}],\\
  \sigma &: \mathcal{O}_Y[T,T^{-1}]\to\mathcal{O}_Y[T,T^{-1}]
\end{align*}
并完全由 $\pi(T)=TT'$ 和 $\sigma(T)=T^{-1}$ 这两个数据所确定。
% Locked French authority: ega2-1-fr.tex, 820505 B,
% SHA-256 84EDBE3E83530AF2959B441796337C9DC21EAFCA6A13114A26778760FBF437AC.
% Sequential source scope: lines 14743-14765, continuation of open Remark II.8.3.9.

在此情形下，可以把 $G_m$ 看作任意\emph{仿射锥}
$C=\operatorname{Spec}(\mathcal{S})$ 的一个“\emph{泛算子域}”，其中
$\mathcal{S}$ 是正次数分次的拟凝聚 $\mathcal{O}_Y$-代数。这意味着，
可以典范地定义一个 $Y$-态射 $G_m\times_YC\to C$，它具有带有算子群
的集合的外作用律所具有的形式性质。或者，像上面对群概形所作的那样，
对每个 $Y$-预概形 $Z$，在 $\operatorname{Hom}_Y(Z,C)$ 上给出一个外作用律，
以群 $\operatorname{Hom}_Y(Z,G_m)$ 为算子集，满足带有算子群的集合的
通常公理，并且满足与 $Y$-态射 $Z\to Z'$ 的相容性条件。在当前情形，
态射 $G_m\times_YC\to C$ 由给定如下 $\mathcal{O}_Y$-代数同态来定义：
\[
  \mathcal{S}\to
  \mathcal{S}\otimes_{\mathcal{O}_Y}\mathcal{O}_Y[T,T^{-1}]
  =\mathcal{S}[T,T^{-1}]
\]
它把每个截面 $s_n\in\Gamma(U,\mathcal{S}_n)$（$U$ 是 $Y$ 的开集）
对应到截面
$s_nT^n\in\Gamma(U,\mathcal{S}\otimes_{\mathcal{O}_Y}
\mathcal{O}_Y[T,T^{-1}])$。
% Locked French authority: ega2-1-fr.tex, 820505 B,
% SHA-256 84EDBE3E83530AF2959B441796337C9DC21EAFCA6A13114A26778760FBF437AC.
% Sequential source scope: lines 14767-14791, conclusion of Remark II.8.3.9.

反过来，设给定一个先验\emph{未分次}的拟凝聚
$\mathcal{O}_Y$-代数 $\mathcal{S}$，并且在
$C=\operatorname{Spec}(\mathcal{S})$ 上给定一个“集合\emph{带有算子群}的
$Y$-概形”结构，其算子域为群 $Y$-概形 $G_m$；那么可由此在
$\mathcal{S}$ 上典范地导出一个 $\mathcal{O}_Y$-代数\emph{分次}。事实上，
给定一个 $Y$-态射 $G_m\times_YC\to C$，等价于给定一个
$\mathcal{O}_Y$-代数同态
$\psi:\mathcal{S}\to\mathcal{S}[T,T^{-1}]$；因此可写成
$\psi=\sum_{n\in\mathbf{Z}}\psi_nT^n$，其中
$\psi_n:\mathcal{S}\to\mathcal{S}$ 是 $\mathcal{O}_Y$-模同态
（对每个截面 $s\in\Gamma(U,\mathcal{S})$，其中 $U$ 是 $Y$ 的开集，
除有限多个指标外均有 $\psi_n(s)=0$）。于是可以验证，带有算子群的
集合的公理蕴含：$\psi_n(\mathcal{S})=\mathcal{S}_n$ 定义了
$\mathcal{S}$ 上的一个 $\mathcal{O}_Y$-代数分次（次数可正可负），而
$\psi_n$ 正是相应的投影算子。由此，对任意仿射 $Y$-概形，都有一种
“\emph{仿射锥}”结构的概念，它以“几何”方式定义，而无须预先给定分次。

\oldpage[II]{168}
这里不再进一步发展这一观点；我们把与上述说明相对应的定义和结果的
精确表述留给读者。
\end{remark}
% Locked French authority: ega2-1-fr.tex, 820505 B,
% SHA-256 84EDBE3E83530AF2959B441796337C9DC21EAFCA6A13114A26778760FBF437AC.
% Sequential source scope: lines 14793-14815, subsection II.8.4 heading,
% Environment II.8.4.1 and Proposition II.8.4.2.

\subsection{向量丛的泛射影闭包}
\label{subsection:II.8.4-zh}

\begin{env}[8.4.1]
\label{II.8.4.1-zh}
设 $Y$ 是预概形，$\mathcal{E}$ 是一个拟凝聚 $\mathcal{O}_Y$-模。若取
$\mathcal{S}$ 为分次 $\mathcal{O}_Y$-代数
$\mathbf{S}_{\mathcal{O}_Y}(\mathcal{E})$，则定义
\egaref{II.8.3.1.1-zh}{8.3.1.1} 表明，$\widehat{\mathcal{S}}$ 等同于
$\mathbf{S}_{\mathcal{O}_Y}(\mathcal{E}\oplus\mathcal{O}_Y)$。由
$\mathcal{S}$ 定义的仿射锥 $\operatorname{Spec}(\mathcal{S})$ 按定义是
$\mathbf{V}(\mathcal{E})$，而 $\operatorname{Proj}(\mathcal{S})$ 按定义是
$\mathbf{P}(\mathcal{E})$，故有：
\end{env}

\begin{proposition}[8.4.2]
\label{II.8.4.2-zh}
$Y$ 上向量丛 $\mathbf{V}(\mathcal{E})$ 的泛射影闭包典范同构于
$\mathbf{P}(\mathcal{E}\oplus\mathcal{O}_Y)$，而后者的无穷远轨迹典范同构于
$\mathbf{P}(\mathcal{E})$。
\end{proposition}
% Locked French authority: ega2-1-fr.tex, 820505 B,
% SHA-256 84EDBE3E83530AF2959B441796337C9DC21EAFCA6A13114A26778760FBF437AC.
% Sequential source scope: lines 14817-14837, Remark II.8.4.3.

\begin{remark}[8.4.3]
\label{II.8.4.3-zh}
例如，取 $\mathcal{E}=\mathcal{O}_Y^r$，其中 $r\geq2$；于是，由
$\mathcal{S}$ 定义的去顶点仿射锥 $E$ 与去顶点射影锥 $\widehat{E}$，当
$Y\neq\varnothing$ 时，在 $Y$ 上既不是仿射的，也不是射影的。后一个断言
是立即的，因为 $\widehat{C}=\mathbf{P}(\mathcal{O}_Y^{r+1})$ 在 $Y$ 上是
射影的，而 $E$ 与 $\widehat{E}$ 的底空间都是 $\widehat{C}$ 中的非闭开集，
所以典范浸入 $E\to\widehat{E}$ 与 $\widehat{E}\to\widehat{C}$ 都不是
射影的（\egaref{II.5.5.3-zh}{5.5.3}），再由
（\egaref{II.5.5.5-zh}{5.5.5, (v)}）即得结论。另一方面，设
$Y=\operatorname{Spec}(A)$ 是仿射的，并例如取 $r=2$；此时
$C=\operatorname{Spec}(A[T_1,T_2])$，而 $E$ 是 $C$ 在开集
$D(T_1)\cup D(T_2)$ 上诱导的预概形；我们已经看到，该开集不是仿射的
（\egaxref{I.5.5.11-zh}{I, 5.5.11}）；\emph{更何况} $\widehat{E}$ 不可能
是仿射的，因为 $E$ 正是 $\widehat{E}$ 上截面 $z$ 不消失的开集
（\egaref{II.8.3.2-zh}{8.3.2}）。

不过：
\end{remark}
% Locked French authority: ega2-1-fr.tex, 820505 B,
% SHA-256 84EDBE3E83530AF2959B441796337C9DC21EAFCA6A13114A26778760FBF437AC.
% Sequential source scope: lines 14839-14861, Proposition II.8.4.4 statement.

\begin{proposition}[8.4.4]
\label{II.8.4.4-zh}
若 $\mathcal{L}$ 是可逆 $\mathcal{O}_Y$-模，则对于与
$C=\mathbf{V}(\mathcal{L})$ 对应的去顶点锥 $E$ 和 $\widehat{E}$，有典范同构
\begin{equation}
\label{II.8.4.4.1-zh}
  \operatorname{Spec}\left(
    \bigoplus_{n\in\mathbf{Z}}\mathcal{L}^{\otimes n}
  \right)\xrightarrow{\sim}E
\tag{8.4.4.1}
\end{equation}
以及
\begin{equation}
\label{II.8.4.4.2-zh}
  \mathbf{V}(\mathcal{L}^{-1})\xrightarrow{\sim}\widehat{E}.
\tag{8.4.4.2}
\end{equation}

此外，$\mathbf{V}(\mathcal{L})$ 的泛射影闭包与
$\mathbf{V}(\mathcal{L}^{-1})$ 的泛射影闭包之间存在一个典范同构；该同构把
前者的零截面（相应地，无穷远轨迹）变为后者的无穷远轨迹（相应地，零截面）。
\end{proposition}
% Locked French authority: ega2-1-fr.tex, 820505 B,
% SHA-256 84EDBE3E83530AF2959B441796337C9DC21EAFCA6A13114A26778760FBF437AC.
% Sequential source scope: lines 14863-14895, opening proof of Proposition II.8.4.4.

\begin{proof}
这里有
$\mathcal{S}=\bigoplus_{n\geq0}\mathcal{L}^{\otimes n}$；典范单射
\[
  \mathcal{S}\to\bigoplus_{n\in\mathbf{Z}}\mathcal{L}^{\otimes n}
\]
定义了一个典范支配态射
\begin{equation}
\label{II.8.4.4.3-zh}
  \operatorname{Spec}\left(
    \bigoplus_{n\in\mathbf{Z}}\mathcal{L}^{\otimes n}
  \right)
  \to\mathbf{V}(\mathcal{L})
  =\operatorname{Spec}\left(
    \bigoplus_{n\geq0}\mathcal{L}^{\otimes n}
  \right)
\tag{8.4.4.3}
\end{equation}
只须证明该态射把概形
$\operatorname{Spec}(\bigoplus_{n\in\mathbf{Z}}\mathcal{L}^{\otimes n})$
同构地映到 $E$ 上。问题关于 $Y$ 是局部的，故可设
$Y=\operatorname{Spec}(A)$ 是仿射的，

\oldpage[II]{169}
并且 $\mathcal{L}=\mathcal{O}_Y$；于是
$\mathcal{S}=(A[T])^{\sim}$，且
$\bigoplus_{n\in\mathbf{Z}}\mathcal{L}^{\otimes n}
=(A[T,T^{-1}])^{\sim}$。而 $A[T,T^{-1}]$ 是 $A[T]$ 的分式环
$A[T]_T$，所以 \egaref{II.8.4.4.3-zh}{8.4.4.3} 把第一项等同于
$C=\mathbf{V}(\mathcal{L})$ 在开集 $D(T)$ 上诱导的预概形。此开集在
$C$ 中的补集 $V(T)$，是由理想 $TA[T]$ 定义的 $C$ 的闭子预概形的
底空间，也就是 $C$ 的零截面；因此 $E=D(T)$。
% Locked French authority: ega2-1-fr.tex, 820505 B,
% SHA-256 84EDBE3E83530AF2959B441796337C9DC21EAFCA6A13114A26778760FBF437AC.
% Sequential source scope: lines 14897-14923, conclusion of the proof of Proposition II.8.4.4.

另一方面，同构 \egaref{II.8.4.4.2-zh}{8.4.4.2} 将由最后一个断言推出：
$\mathbf{V}(\mathcal{L}^{-1})$ 是其泛射影闭包中无穷远轨迹的补集，而
$\widehat{E}$ 是 $C=\mathbf{V}(\mathcal{L})$ 的泛射影闭包中零截面的补集。
这些泛射影闭包分别是
$\mathbf{P}(\mathcal{L}^{-1}\oplus\mathcal{O}_Y)$ 和
$\mathbf{P}(\mathcal{L}\oplus\mathcal{O}_Y)$；但可写成
$\mathcal{L}\oplus\mathcal{O}_Y
=\mathcal{L}\otimes(\mathcal{L}^{-1}\oplus\mathcal{O}_Y)$。所求典范同构的
存在性于是由 \egaref{II.4.1.4-zh}{4.1.4} 得出，问题归结为说明该同构
交换零截面与无穷远轨迹。为此可局部化到
$Y=\operatorname{Spec}(A)$ 为仿射的情形，并取
$L=Ac$、$L^{-1}=Ac'$，使典范同构 $L\otimes L^{-1}\to A$ 把
$c\otimes c'$ 映到 $A$ 的元素 $1$。此时，$\mathbf{S}(L\oplus A)$ 是
$A[z]$ 与 $\bigoplus_{n\geq0}Ac^{\otimes n}$ 的张量积，
$\mathbf{S}(L^{-1}\oplus A)$ 是 $A[z]$ 与
$\bigoplus_{n\geq0}Ac'^{\otimes n}$ 的张量积，而
\egaref{II.4.1.4-zh}{4.1.4} 中定义的同构把
$z^h\otimes c'^{\otimes(n-h)}$ 映到
$z^{n-h}\otimes c^{\otimes h}$。在
$\mathbf{P}(\mathcal{L}^{-1}\oplus\mathcal{O}_Y)$ 中，无穷远轨迹是截面
$z$ 的零点集，而零截面是截面 $c'$ 的零点集；对
$\mathbf{P}(\mathcal{L}\oplus\mathcal{O}_Y)$ 也有类似定义，故结论立即由
上述具体说明得出。
\end{proof}
% Locked French authority: ega2-1-fr.tex, 820505 B,
% SHA-256 84EDBE3E83530AF2959B441796337C9DC21EAFCA6A13114A26778760FBF437AC.
% Sequential source scope: lines 14925-14980, subsection II.8.5 and
% environment II.8.5.1 / printed pp.169-170.

\subsection{函子性}
\label{subsection:II.8.5-zh}

\begin{env}[8.5.1]
\label{II.8.5.1-zh}
设 $Y$、$Y'$ 是两个预概形，$q:Y'\to Y$ 是一个态射，
$\mathcal{S}$（相应地，$\mathcal{S}'$）是一个按\emph{正次数}分次的拟凝聚
$\mathcal{O}_Y$-代数层（相应地，$\mathcal{O}_{Y'}$-代数层）。考虑分次
代数层的一个 $q$-态射
\begin{equation}
\label{II.8.5.1.1-zh}
  \varphi:\mathcal{S}\to\mathcal{S}'.
\tag{8.5.1.1}
\end{equation}
由（\egaref{II.1.5.6-zh}{1.5.6}）可知，典范地有与之相应的态射
\[
  \Phi=\operatorname{Spec}(\varphi):
  \operatorname{Spec}(\mathcal{S}')\to\operatorname{Spec}(\mathcal{S})
\]
使图
\begin{equation}
\label{II.8.5.1.2-zh}
  \xymatrix{
    C'\ar[r]^{\Phi}\ar[d] & C\ar[d]\\
    Y'\ar[r]_{q} & Y
  }
\tag{8.5.1.2}
\end{equation}
交换，其中置 $C=\operatorname{Spec}(\mathcal{S})$、
$C'=\operatorname{Spec}(\mathcal{S}')$。
\emph{此外，假设
$\mathcal{S}_0=\mathcal{O}_Y$、$\mathcal{S}'_0=\mathcal{O}_{Y'}$}；设
$\varepsilon:Y\to C$ 和 $\varepsilon':Y'\to C'$ 是典范浸入
（\egaref{II.8.3.2-zh}{8.3.2}）；于是有交换图
\begin{equation}
\label{II.8.5.1.3-zh}
  \xymatrix{
    Y'\ar[r]^{q}\ar[d]_{\varepsilon'} & Y\ar[d]^{\varepsilon}\\
    C'\ar[r]_{\Phi} & C
  }
\tag{8.5.1.3}
\end{equation}

\oldpage[II]{170}

它对应于图
\[
  \xymatrix{
    \mathcal{S}\ar[r]^{\varphi}\ar[d] & \mathcal{S}'\ar[d]\\
    \mathcal{O}_Y\ar[r] & \mathcal{O}_{Y'}
  }
\]
其中竖直箭头是增广同态；该图的交换性由 $\varphi$ 是一个\emph{分次}
代数同态这一假设得出。
\end{env}
% Locked French authority: ega2-1-fr.tex, 820505 B,
% SHA-256 84EDBE3E83530AF2959B441796337C9DC21EAFCA6A13114A26778760FBF437AC.
% Sequential source scope: lines 14982-15005, Proposition II.8.5.2 and proof.

\begin{proposition}[8.5.2]
\label{II.8.5.2-zh}
若 $E$（相应地 $E'$）是由 $\mathcal{S}$（相应地 $\mathcal{S}'$）定义的
去顶点仿射锥，则 $\Phi^{-1}(E)\subset E'$；若进一步
$\operatorname{Proj}(\varphi):G(\varphi)\to\operatorname{Proj}(\mathcal{S})$
处处有定义（换言之，若
$G(\varphi)=\operatorname{Proj}(\mathcal{S}')$），则
$\Phi^{-1}(E)=E'$，且反之亦然。
\end{proposition}

\begin{proof}
第一个断言由图表 \egaref{II.8.5.1.3-zh}{8.5.1.3} 的交换性得出。为证明
第二个断言，可归结到 $Y=\operatorname{Spec}(A)$ 和
$Y'=\operatorname{Spec}(A')$ 均为仿射，且
$\mathcal{S}=\widetilde{S}$、$\mathcal{S}'=\widetilde{S'}$ 的情形。对
$S_+$ 中每个齐次元 $f$，置 $f'=\varphi(f)$，则
$\Phi^{-1}(C_f)=C'_{f'}$（\egaxref{I.2.2.4.1-zh}{I, 2.2.4.1}）。条件
$G(\varphi)=\operatorname{Proj}(S')$ 意味着：在 $S'_+$ 中，由所有
$f'=\varphi(f)$ 生成的理想的根就是 $S'_+$ 本身
（\egaref{II.2.8.1-zh}{2.8.1} 和
\egaref{II.2.3.14-zh}{2.3.14}）；这又等价于各 $C'_{f'}$ 覆盖 $E'$
（\egaref{II.8.3.4.4-zh}{8.3.4.4}）。
\end{proof}
% Locked French authority: ega2-1-fr.tex, 820505 B,
% SHA-256 84EDBE3E83530AF2959B441796337C9DC21EAFCA6A13114A26778760FBF437AC.
% Sequential source scope: lines 15007-15062, Environment II.8.5.3.

\begin{env}[8.5.3]
\label{II.8.5.3-zh}
$q$-态射 $\varphi$ 典范地延拓为分次代数的一个 $q$-态射
\begin{equation}
\label{II.8.5.3.1-zh}
  \widehat{\varphi}:\widehat{\mathcal{S}}
  \to\widehat{\mathcal{S}'}
\tag{8.5.3.1}
\end{equation}
其中规定 $\widehat{\varphi}(z)=z$。于是由此得到一个态射
\[
  \widehat{\Phi}=\operatorname{Proj}(\widehat{\varphi}):
  G(\widehat{\varphi})\to\widehat{C}
  =\operatorname{Proj}(\widehat{\mathcal{S}})
\]
使图表
\[
  \xymatrix{
    G(\widehat{\varphi})\ar[r]^{\widehat{\Phi}}\ar[d] &
    \widehat{C}\ar[d]\\
    Y'\ar[r]_{q} & Y
  }
\]
交换（\egaref{II.3.5.6-zh}{3.5.6}）。由定义立即可知，若以
$i:C\to\widehat{C}$ 和 $i':C'\to\widehat{C'}$ 表示典范开浸入
（\egaref{II.8.3.2-zh}{8.3.2}），则
$i'(C')\subset G(\widehat{\varphi})$，且图表
\begin{equation}
\label{II.8.5.3.2-zh}
  \xymatrix{
    C'\ar[r]^{\Phi}\ar[d]_{i'} & C\ar[d]^{i}\\
    G(\widehat{\varphi})\ar[r]_{\widehat{\Phi}} & \widehat{C}
  }
\tag{8.5.3.2}
\end{equation}
交换。最后，置
$X=\operatorname{Proj}(\mathcal{S})$、$X'=\operatorname{Proj}(\mathcal{S}')$，
并设 $j:X\to\widehat{C}$、$j':X'\to\widehat{C'}$ 为典范闭浸入
（\egaref{II.8.3.2-zh}{8.3.2}）。仍由这些浸入的定义可得
$j'(G(\varphi))\subset G(\widehat{\varphi})$，并且图表

\oldpage[II]{171}
\begin{equation}
\label{II.8.5.3.3-zh}
  \xymatrix{
    G(\varphi)\ar[r]^{\operatorname{Proj}(\varphi)}\ar[d]_{j'} &
    X\ar[d]^{j}\\
    G(\widehat{\varphi})\ar[r]_{\widehat{\Phi}} & \widehat{C}
  }
\tag{8.5.3.3}
\end{equation}
交换。
\end{env}
% Locked French authority: ega2-1-fr.tex, 820505 B,
% SHA-256 84EDBE3E83530AF2959B441796337C9DC21EAFCA6A13114A26778760FBF437AC.
% Sequential source scope: lines 15064-15104, Proposition II.8.5.4 and proof.

\begin{proposition}[8.5.4]
\label{II.8.5.4-zh}
若 $\widehat{E}$（相应地 $\widehat{E'}$）是由 $\mathcal{S}$（相应地
$\mathcal{S}'$）定义的去顶点射影锥，则
$\widehat{\Phi}^{-1}(\widehat{E})\subset\widehat{E'}$；此外，若
$p:\widehat{E}\to X$ 和 $p':\widehat{E'}\to X'$ 是典范收缩，则
$p'(\widehat{\Phi}^{-1}(\widehat{E}))\subset G(\varphi)$，且图表
\begin{equation}
\label{II.8.5.4.1-zh}
  \xymatrix{
    \widehat{\Phi}^{-1}(\widehat{E})
      \ar[r]^{\widehat{\Phi}}\ar[d]_{p'} &
    \widehat{E}\ar[d]^{p}\\
    G(\varphi)\ar[r]_{\operatorname{Proj}(\varphi)} & X
  }
\tag{8.5.4.1}
\end{equation}
交换。若 $\operatorname{Proj}(\varphi)$ 处处有定义，则
$\widehat{\Phi}$ 也处处有定义，并且
$\widehat{\Phi}^{-1}(\widehat{E})=\widehat{E'}$。
\end{proposition}

\begin{proof}
第一个断言由图表 \egaref{II.8.5.1.3-zh}{8.5.1.3} 和
\egaref{II.8.5.3.2-zh}{8.5.3.2} 的交换性得出；其后的两个断言由
典范收缩的定义（\egaref{II.8.3.5-zh}{8.3.5}）以及
$\widehat{\varphi}$ 的定义得出。另一方面，为说明当
$\operatorname{Proj}(\varphi)$ 处处有定义时 $\widehat{\Phi}$ 也处处有
定义，可归结到 $Y=\operatorname{Spec}(A)$ 和
$Y'=\operatorname{Spec}(A')$ 均为仿射，且
$\mathcal{S}=\widetilde{S}$、$\mathcal{S}'=\widetilde{S'}$ 的情形。假设是：
当 $f$ 遍历 $S_+$ 的齐次元时，在 $S'_+$ 中由各
$\varphi(f)$ 生成的理想在 $S'_+$ 中的根是 $S'_+$ 本身。由此立即得出，
由 $z$ 和各 $\varphi(f)$ 生成的理想在
$(S'[z])_+$ 中的根是 $(S'[z])_+$ 本身，故得所述断言。这同时证明
$\widehat{E'}$ 是各 $\widehat{C'}_{\varphi(f)}$ 的并，因而等于
$\widehat{\Phi}^{-1}(\widehat{E})$。
\end{proof}
% Locked French authority: ega2-1-fr.tex, 820505 B,
% SHA-256 84EDBE3E83530AF2959B441796337C9DC21EAFCA6A13114A26778760FBF437AC.
% Sequential source scope: lines 15106-15119, Corollary II.8.5.5 and proof.

\begin{corollary}[8.5.5]
\label{II.8.5.5-zh}
当 $\operatorname{Proj}(\varphi)$ 处处有定义时，$\widehat{C}$ 的无穷远轨迹
（相应地顶点预概形）底空间在 $\widehat{\Phi}$ 下的原像，是
$\widehat{C'}$ 的无穷远轨迹（相应地顶点预概形）的底空间。
\end{corollary}

\begin{proof}
结合关系 \egaref{II.8.3.4.1-zh}{8.3.4.1} 和
\egaref{II.8.3.4.2-zh}{8.3.4.2}，这立即由
\egaref{II.8.5.4-zh}{8.5.4} 和 \egaref{II.8.5.2-zh}{8.5.2} 得出。
\end{proof}
% Locked French authority: ega2-1-fr.tex, 820505 B,
% SHA-256 84EDBE3E83530AF2959B441796337C9DC21EAFCA6A13114A26778760FBF437AC.
% Sequential source scope: lines 15121-15174, subsection II.8.6 and
% environment II.8.6.1 / printed pp.171-172.

\subsection{去顶点锥的一个典范同构}
\label{subsection:II.8.6-zh}

\begin{env}[8.6.1]
\label{II.8.6.1-zh}
设 $Y$ 是一个预概形，$\mathcal{S}$ 是一个按正次数分次的拟凝聚
$\mathcal{O}_Y$-代数层，\emph{且
$\mathcal{S}_0=\mathcal{O}_Y$}，又设 $X$ 是 $Y$-概形
$\operatorname{Proj}(\mathcal{S})$。我们将把
\egaref{subsection:II.8.5-zh}{8.5} 的结果应用于如下情形：取 $Y'=X$，
$q:X\to Y$ 为结构态射；令
\begin{equation}
\label{II.8.6.1.1-zh}
  \mathcal{S}_X=\bigoplus_{n\in\mathbf{Z}}\mathcal{O}_X(n)
\tag{8.6.1.1}
\end{equation}

\oldpage[II]{172}
这是一个分次拟凝聚 $\mathcal{O}_X$-代数层，其乘法由典范同态
\egaref{II.3.2.6.1-zh}{3.2.6.1}
\[
  \mathcal{O}_X(m)\otimes_{\mathcal{O}_X}\mathcal{O}_X(n)
  \to\mathcal{O}_X(m+n)
\]
定义；其结合性由交换图 \egaref{II.2.5.11.4-zh}{2.5.11.4} 保证。
取 $\mathcal{S}'$ 为 $\mathcal{S}_X$ 的按正次数分次的拟凝聚
$\mathcal{O}_X$-子代数层
$\mathcal{S}_X^{\geq}=\bigoplus_{n\geq0}\mathcal{O}_X(n)$。

最后，考虑典范 $q$-态射
\begin{equation}
\label{II.8.6.1.2-zh}
  \alpha:\mathcal{S}\to\mathcal{S}_X^{\geq}
\tag{8.6.1.2}
\end{equation}
它在 \egaref{II.3.3.2.3-zh}{3.3.2.3} 中定义为同态
$\mathcal{S}\to q_*(\mathcal{S}_X)$，但显然把 $\mathcal{S}$ 映入
$q_*(\mathcal{S}_X^{\geq})$。我们置
\begin{equation}
\label{II.8.6.1.3-zh}
  C_X=\operatorname{Spec}(\mathcal{S}_X^{\geq}),\qquad
  \widehat{C}_X=\operatorname{Proj}(\mathcal{S}_X^{\geq}[z]),\qquad
  X'=\operatorname{Proj}(\mathcal{S}_X^{\geq})
\tag{8.6.1.3}
\end{equation}
并以 $E_X$ 和 $\widehat{E}_X$ 分别表示相应的去顶点仿射锥与去顶点射影锥；
以 $\varepsilon_X:X\to C_X$、$i_X:C_X\to\widehat{C}_X$、
$j_X:X'\to\widehat{C}_X$、$p_X:\widehat{E}_X\to X'$、
$\pi_X:E_X\to X'$ 表示在
\egaref{subsection:II.8.3-zh}{8.3} 中定义的典范态射。
\end{env}
% Locked French authority: ega2-1-fr.tex, 820505 B,
% SHA-256 84EDBE3E83530AF2959B441796337C9DC21EAFCA6A13114A26778760FBF437AC.
% Sequential source scope: lines 15176-15219, Proposition II.8.6.2 and first proof portion.

\begin{proposition}[8.6.2]
\label{II.8.6.2-zh}
结构态射 $u:X'\to X$ 是同构，且态射
$\operatorname{Proj}(\alpha)$ 处处有定义并与 $u$ 相同。态射
$\operatorname{Proj}(\widehat{\alpha}):\widehat{C}_X\to\widehat{C}$
处处有定义，并且它在 $\widehat{E}_X$ 和 $E_X$ 上的限制分别是到
$\widehat{E}$ 和 $E$ 的同构。最后，当把 $X'$ 与 $X$ 借助 $u$
等同起来时，态射 $p_X$ 和 $\pi_X$ 分别与 $X$-预概形
$\widehat{E}_X$ 和 $E_X$ 的结构态射相等同。
\end{proposition}

\begin{proof}
显然可以只考虑 $Y=\operatorname{Spec}(A)$ 为仿射且
$\mathcal{S}=\widetilde{S}$ 的情形；此时 $X$ 是各仿射开集 $X_f$ 的并，
其中 $f$ 遍历 $S_+$ 的齐次元集合，而 $X_f$ 的环为 $S_{(f)}$。此外，
由 \egaref{II.8.2.7.1-zh}{8.2.7.1} 可得
\begin{equation}
\label{II.8.6.2.1-zh}
  \Gamma(X_f,\mathcal{S}_X^{\geq})=S_f^{\geq}.
\tag{8.6.2.1}
\end{equation}

因此 $u^{-1}(X_f)=\operatorname{Proj}(S_f^{\geq})$。但若
$f\in S_d$（$d>0$），则 $\operatorname{Proj}(S_f^{\geq})$ 典范同构于
$\operatorname{Proj}((S_f^{\geq})^{(d)})$
（\egaref{II.2.4.7-zh}{2.4.7}）；另一方面，
$(S_f^{\geq})^{(d)}=(S^{(d)})_f^{\geq}$ 与 $S_{(f)}[T]$ 相等同
（\egaref{II.2.2.1-zh}{2.2.1}），所用映射为 $T\to f/1$。由此
（\egaref{II.3.1.7-zh}{3.1.7}）可知，结构态射
$u^{-1}(X_f)\to X_f$ 是同构，这就得到第一个断言。为证明第二个断言，
注意，限制 $u^{-1}(X_f)\cap G(\alpha)\to X=\operatorname{Proj}(S)$，
也就是 $\operatorname{Proj}(\alpha)$ 的限制，对应于典范映射
$x\to x/1$，这是从 $S$ 到 $S_f^{\geq}$ 的映射
（\egaref{II.2.6.2-zh}{2.6.2}）。由此首先得到
$G(\alpha)=X'$；再考虑到
$u^{-1}(X_f)=(u^{-1}(X_f))_{f/1}$，由
\egaref{II.2.8.1.1-zh}{2.8.1.1} 可知，
$u^{-1}(X_f)$ 在 $\operatorname{Proj}(\alpha)$ 下的像包含于 $X_f$，
而 $\operatorname{Proj}(\alpha)$ 在 $u^{-1}(X_f)$ 上的限制，视为到
$X_f=\operatorname{Spec}(S_{(f)})$ 的态射，正与 $u$ 的限制相同。
% Locked French authority: ega2-1-fr.tex, 820505 B,
% SHA-256 84EDBE3E83530AF2959B441796337C9DC21EAFCA6A13114A26778760FBF437AC.
% Sequential source scope: lines 15220-15239, conclusion of proof of Proposition II.8.6.2.

最后，将公式 \egaref{II.8.3.5.4-zh}{8.3.5.4} 应用于 $p_X$ 而非 $p$，
可见
$p_X^{-1}(u^{-1}(X_f))=
\operatorname{Spec}((S_f^{\geq})_{f/1}^{\leq})$；依
\egaref{II.8.5.4.1-zh}{8.5.4.1}，这个开集是态射
$\operatorname{Proj}(\widehat{\alpha})$ 对
$p^{-1}(X_f)=\operatorname{Spec}(S_f^{\leq})$
（\egaref{II.8.3.5.3-zh}{8.3.5.3}）的原像。考虑到
\egaref{II.8.2.3.2-zh}{8.2.3.2}，
$\operatorname{Proj}(\widehat{\alpha})$ 在
$p_X^{-1}(u^{-1}(X_f))$ 上的限制，对应于
\egaref{II.8.2.7.2-zh}{8.2.7.2} 中同构的逆同构在
$S_f^{\leq}$ 上的限制，由此得到第三个断言；最后一个断言由定义显然成立。

\oldpage[II]{173}
还应注意，由交换图 \egaref{II.8.5.3.2-zh}{8.5.3.2} 可知，
\emph{在 $C_X$ 上限制 $\operatorname{Proj}(\widehat{\alpha})$ 所得的态射恰为
$\operatorname{Spec}(\alpha)$}。
\end{proof}
% Locked French authority: ega2-1-fr.tex, 820505 B,
% SHA-256 84EDBE3E83530AF2959B441796337C9DC21EAFCA6A13114A26778760FBF437AC.
% Sequential source scope: lines 15241-15272, Corollary II.8.6.3 and proof.

\begin{corollary}[8.6.3]
\label{II.8.6.3-zh}
视为 $X$-概形时，$\widehat{E}_X$ 典范同构于
$\operatorname{Spec}(\mathcal{S}_X^{\leq})$，而 $E_X$ 典范同构于
$\operatorname{Spec}(\mathcal{S}_X)$。
\end{corollary}

\begin{proof}
已知态射 $p_X$ 和 $\pi_X$ 都是仿射的
（\egaref{II.8.3.5-zh}{8.3.5} 和 \egaref{II.8.3.6-zh}{8.3.6}），因此
根据 \egaref{II.1.3.1-zh}{1.3.1}，只须在
$Y=\operatorname{Spec}(A)$ 为仿射且 $\mathcal{S}=\widetilde{S}$ 时验证
推论。第一个断言由典范同构
\egaref{II.8.2.7.2-zh}{8.2.7.2}
$(S_f^{\geq})_{f/1}^{\leq}\xrightarrow{\sim}S_f^{\leq}$ 的存在性以及
这些同构与从 $f$ 过渡到 $fg$（$f$ 和 $g$ 是 $S_+$ 中的齐次元）相容
而得。类似地，把公式 \egaref{II.8.3.6.2-zh}{8.3.6.2} 应用于
$\pi_X$ 而非 $\pi$，可知对 $S_+$ 中的齐次元 $f$，
$\pi_X^{-1}(u^{-1}(X_f))=
\operatorname{Spec}((S_f^{\geq})_{f/1})$；第二个断言由典范同构
\egaref{II.8.2.7.2-zh}{8.2.7.2}
$(S_f^{\geq})_{f/1}\xrightarrow{\sim}S_f$ 的存在性推出。

因此可以说，把 $\widehat{C}_X$ 视为 $X$-概形时，它由 $X$ 上两个仿射
$X$-概形
$C_X=\operatorname{Spec}(\mathcal{S}_X^{\geq})$ 和
$\widehat{E}_X=\operatorname{Spec}(\mathcal{S}_X^{\leq})$ 粘合而成，二者的
交是开集 $E_X=\operatorname{Spec}(\mathcal{S}_X)$。
\end{proof}
% Locked French authority: ega2-1-fr.tex, 820505 B,
% SHA-256 84EDBE3E83530AF2959B441796337C9DC21EAFCA6A13114A26778760FBF437AC.
% Sequential source scope: lines 15274-15302, Corollary II.8.6.4 and proof.

\begin{corollary}[8.6.4]
\label{II.8.6.4-zh}
假设 $\mathcal{O}_X(1)$ 是可逆 $\mathcal{O}_X$-模，并且
$\mathcal{S}_X$ 同构于
$\bigoplus_{n\in\mathbf{Z}}(\mathcal{O}_X(1))^{\otimes n}$（特别地，当
$\mathcal{S}$ 由 $\mathcal{S}_1$ 生成时，这一条件成立
（\egaref{II.3.2.5-zh}{3.2.5} 和
\egaref{II.3.2.7-zh}{3.2.7}））。于是，去顶点射影锥
$\widehat{E}$ 等同于 $X$ 上秩为一的向量丛
$\mathbf{V}(\mathcal{O}_X(-1))$，而去顶点仿射锥 $E$ 等同于该向量丛
在零截面补集上诱导的子预概形。在这一等同下，典范收缩
$\widehat{E}\to X$ 等同于 $X$-概形
$\mathbf{V}(\mathcal{O}_X(-1))$ 的结构态射。最后，存在一个典范
$Y$-态射 $\mathbf{V}(\mathcal{O}_X(1))\to C$；它在
$\mathbf{V}(\mathcal{O}_X(1))$ 的零截面补集上的限制，是该补集到
去顶点仿射锥 $E$ 上的同构。
\end{corollary}

\begin{proof}
事实上，若置 $\mathcal{L}=\mathcal{O}_X(1)$，则
$\mathcal{S}_X^{\geq}$ 等于
$\mathbf{S}_{\mathcal{O}_X}(\mathcal{L})$；因此由
\egaref{II.8.6.3-zh}{8.6.3}，$\widehat{E}_X$ 典范等同于
$\mathbf{V}(\mathcal{L}^{-1})$，而 $C_X$ 等同于
$\mathbf{V}(\mathcal{L})$。态射 $\mathbf{V}(\mathcal{L})\to C$ 是
$\operatorname{Proj}(\widehat{\alpha})$ 的限制，所以推论的各断言都是
\egaref{II.8.6.2-zh}{8.6.2} 的特殊情形。
\end{proof}
% Locked French authority: ega2-1-fr.tex, 820505 B,
% SHA-256 84EDBE3E83530AF2959B441796337C9DC21EAFCA6A13114A26778760FBF437AC.
% Sequential source scope: lines 15304-15310, concluding paragraph of II.8.6.

还应注意，$C$ 的顶点预概形底空间在态射
$\mathbf{V}(\mathcal{O}_X(1))\to C$ 下的原像，是
$\mathbf{V}(\mathcal{O}_X(1))$ 的零截面底空间
（\egaref{II.8.5.5-zh}{8.5.5}）；但一般而言，$C$ 和
$\mathbf{V}(\mathcal{O}_X(1))$ 中相应的子预概形并不同构。下面将研究
这一问题。
% Locked French authority: ega2-1-fr.tex, 820505 B,
% SHA-256 84EDBE3E83530AF2959B441796337C9DC21EAFCA6A13114A26778760FBF437AC.
% Sequential source scope: lines 15312-15349, subsection II.8.7 and Environment II.8.7.1.

\subsection{投影锥的暴涨}
\label{subsection:II.8.7-zh}

\begin{env}[8.7.1]
\label{II.8.7.1-zh}
在 \egaref{II.8.6.1-zh}{8.6.1} 的条件下，置
$r=\operatorname{Proj}(\widehat{\alpha})$，则有交换图表
\begin{equation}
\label{II.8.7.1.1-zh}
  \xymatrix{
    X\ar[r]^{i_X\circ\varepsilon_X}\ar[d]_{q} &
    \widehat{C}_X\ar[d]^{r}\\
    Y\ar[r]_{i\circ\varepsilon} & \widehat{C}
  }
\tag{8.7.1.1}
\end{equation}

\oldpage[II]{174}
这是 \egaref{II.8.5.1.3-zh}{8.5.1.3} 和
\egaref{II.8.5.3.2-zh}{8.5.3.2} 的结果；此外，由
\egaref{II.8.6.2-zh}{8.6.2}，$r$ 在零截面补集
$\widehat{C}_X-i_X(\varepsilon_X(X))$ 上的限制，是到零截面补集
$\widehat{C}-i(\varepsilon(Y))$ 上的一个\emph{同构}。为简便起见，若
假设 $Y$ 为仿射，$\mathcal{S}$ 是有限型的并由 $\mathcal{S}_1$ 生成，则
$X$ 在 $Y$ 上射影，且 $\widehat{C}_X$ 在 $X$ 上射影
（\egaref{II.5.5.1-zh}{5.5.1}），故 $\widehat{C}_X$ 在 $Y$ 上射影
（\egaref{II.5.5.5-zh}{5.5.5, (ii)}），从而更在 $\widehat{C}$ 上射影
（\egaref{II.5.5.5-zh}{5.5.5, (v)}）。因此有一个射影 $Y$-态射
$r:\widehat{C}_X\to\widehat{C}$（其在 $C_X$ 上的限制是射影 $Y$-态射
$C_X\to C$），它\emph{把 $X$ 收缩到 $Y$ 中}，并且限制到
\emph{$X$ 与 $Y$ 的补集}时诱导一个\emph{同构}。因此 $C_X$ 与 $C$ 之间
存在一种关系，类似于暴涨预概形与其原预概形之间的关系
（\egaref{II.8.1.3-zh}{8.1.3}）。我们将实际证明，$C_X$ 可等同于某个
分次 $\mathcal{O}_C$-代数的齐次谱。
\end{env}
% Locked French authority: ega2-1-fr.tex, 820505 B,
% SHA-256 84EDBE3E83530AF2959B441796337C9DC21EAFCA6A13114A26778760FBF437AC.
% Sequential source scope: lines 15351-15409, Environment II.8.7.2.

\begin{env}[8.7.2]
\label{II.8.7.2-zh}
沿用 \egaref{II.8.6.1-zh}{8.6.1} 的记号。对每个 $n\geqslant 0$，考虑
分次 $\mathcal{O}_Y$-代数 $\mathcal{S}$ 的拟凝聚理想
\begin{equation}
\label{II.8.7.2.1-zh}
  \mathcal{S}_{[n]}=\bigoplus_{m\geqslant n}\mathcal{S}_m
\tag{8.7.2.1}
\end{equation}
显然有
\begin{equation}
\label{II.8.7.2.2-zh}
  \mathcal{S}_{[0]}=\mathcal{S},\qquad
  \mathcal{S}_{[n]}\subset\mathcal{S}_{[m]}
  \qquad\text{对 }m\leqslant n
\tag{8.7.2.2}
\end{equation}
以及
\begin{equation}
\label{II.8.7.2.3-zh}
  \mathcal{S}_{[n]}\mathcal{S}_{[m]}\subset\mathcal{S}_{[m+n]}.
\tag{8.7.2.3}
\end{equation}

考虑与 $\mathcal{S}_{[n]}$ 相伴的 $\mathcal{O}_C$-模；它是
$\mathcal{O}_C=\widetilde{\mathcal{S}}$ 的拟凝聚理想
（\egaref{II.1.4.4-zh}{1.4.4}）
\begin{equation}
\label{II.8.7.2.4-zh}
  \mathcal{I}_n=(\mathcal{S}_{[n]})^{\sim}.
\tag{8.7.2.4}
\end{equation}

于是，利用 \egaref{II.1.4.4-zh}{1.4.4} 和
\egaref{II.1.4.8.1-zh}{1.4.8.1}，从
\egaref{II.8.7.2.2-zh}{8.7.2.2} 和
\egaref{II.8.7.2.3-zh}{8.7.2.3} 推出类似公式
\begin{equation}
\label{II.8.7.2.5-zh}
  \mathcal{I}_0=\mathcal{O}_C,\qquad
  \mathcal{I}_n\subset\mathcal{I}_m
  \qquad\text{对 }m\leqslant n
\tag{8.7.2.5}
\end{equation}
以及
\begin{equation}
\label{II.8.7.2.6-zh}
  \mathcal{I}_n\mathcal{I}_m\subset\mathcal{I}_{n+m}.
\tag{8.7.2.6}
\end{equation}

因此满足 \egaref{II.8.1.1-zh}{8.1.1} 的条件，这引出拟凝聚分次
$\mathcal{O}_C$-代数
\begin{equation}
\label{II.8.7.2.7-zh}
  \mathcal{S}^{\natural}
  =\bigoplus_{n\geqslant 0}\mathcal{I}_n
  =\left(\bigoplus_{n\geqslant 0}\mathcal{S}_{[n]}\right)^{\sim}.
\tag{8.7.2.7}
\end{equation}
\end{env}
% Locked French authority: ega2-1-fr.tex, 820505 B,
% SHA-256 84EDBE3E83530AF2959B441796337C9DC21EAFCA6A13114A26778760FBF437AC.
% Sequential source scope: lines 15411-15440, Proposition II.8.7.3 and
% opening of its proof / printed pp.174-175.

\begin{proposition}[8.7.3]
\label{II.8.7.3-zh}
存在一个典范 $C$-同构
\begin{equation}
\label{II.8.7.3.1-zh}
  h:C_X\xrightarrow{\sim}\operatorname{Proj}(\mathcal{S}^{\natural}).
\tag{8.7.3.1}
\end{equation}
\end{proposition}

\begin{proof}
先假设 $Y=\operatorname{Spec}(A)$ 为仿射的，于是
$\mathcal{S}=\widetilde{S}$，其中 $S$ 是一个按正次数分次的 $A$-代数，且
$C=\operatorname{Spec}(S)$。定义
\egaref{II.8.2.6.2-zh}{8.2.7.4} 表明，在
\egaref{II.8.2.6-zh}{8.2.6} 的记号下，有
$\mathcal{S}^{\natural}=(S^{\natural})^{\sim}$。为定义
\egaref{II.8.7.3.1-zh}{8.7.3.1}，考虑一个齐次元
$f\in S_d$（$d>0$）以及相应的元素 $f^{\natural}\in S^{\natural}$
（\egaref{II.8.2.6-zh}{8.2.6}）；于是 $S$-同构
\egaref{II.8.2.7.3-zh}{8.2.7.3} 定义了一个 $C$-同构
\begin{equation}
\label{II.8.7.3.2-zh}
  \operatorname{Spec}(S_f^{\geqslant})
  \xrightarrow{\sim}
  \operatorname{Spec}(S_{(f^{\natural})}^{\natural}).
\tag{8.7.3.2}
\end{equation}

\oldpage[II]{175}
% Locked French authority: ega2-1-fr.tex, 820505 B,
% SHA-256 84EDBE3E83530AF2959B441796337C9DC21EAFCA6A13114A26778760FBF437AC.
% Sequential source scope: lines 15441-15494, continuation and closure of the proof of II.8.7.3.

但是沿用 \egaref{II.8.6.2-zh}{8.6.2} 的记号，若
$v:C_X\to X$ 是结构态射，则由
\egaref{II.8.6.2.1-zh}{8.6.2.1} 可知
$v^{-1}(X_f)=\operatorname{Spec}(S_f^{\geqslant})$。另一方面有
$\operatorname{Spec}(S_{(f^{\natural})}^{\natural})=D_+(f^{\natural})$，
所以 \egaref{II.8.7.3.2-zh}{8.7.3.2} 定义了一个同构
$v^{-1}(X_f)\to D_+(f^{\natural})$。此外，若 $g\in S_e$（$e>0$），则图表
\[
  \xymatrix{
    v^{-1}(X_{fg})\ar[r]^{\sim}\ar[d] &
    D_+(f^{\natural}g^{\natural})\ar[d]\\
    v^{-1}(X_f)\ar[r]^{\sim} & D_+(f^{\natural})
  }
\]
交换；这也由同构 \egaref{II.8.2.7.3-zh}{8.2.7.3} 的定义立即得出。
最后，按定义 $S_+$ 由各齐次 $f$ 生成，故由
\egaref{II.8.2.10-zh}{8.2.10, (iv)} 和
\egaref{II.2.3.14-zh}{2.3.14} 可知，各 $D_+(f^{\natural})$ 构成
$\operatorname{Proj}(S^{\natural})$ 的一个覆盖，而各 $v^{-1}(X_f)$ 构成
$C_X$ 的一个覆盖，因为各 $X_f$ 构成 $X$ 的一个覆盖；这就在此情形下
完成了同构 \egaref{II.8.7.3.1-zh}{8.7.3.1} 的定义。

为在一般情形下证明 \egaref{II.8.7.3-zh}{8.7.3}，只需验证如下事实：若
$U$、$U'$ 是 $Y$ 的两个仿射开集，满足 $U'\subset U$，其环分别为
$A$ 和 $A'$，并置
$\mathcal{S}|U=\widetilde{S}$、$\mathcal{S}|U'=\widetilde{S'}$，则图表
\begin{equation}
\label{II.8.7.3.3-zh}
  \xymatrix{
    C_{U'}\ar[r]\ar[d] & \operatorname{Proj}(S'^{\natural})\ar[d]\\
    C_U\ar[r] & \operatorname{Proj}(S^{\natural})
  }
\tag{8.7.3.3}
\end{equation}
交换。而 $S'$ 典范等同于 $S\otimes_A A'$，因而
$S'^{\natural}$ 典范等同于
\[
  S^{\natural}\otimes_S S'=S^{\natural}\otimes_A A'~;
\]
所以
$\operatorname{Proj}(S'^{\natural})=\operatorname{Proj}(S^{\natural})
\times_U U'$（\egaref{II.2.8.10-zh}{2.8.10}）。同样，若
$X=\operatorname{Proj}(S)$、$X'=\operatorname{Proj}(S')$，则有
$X'=X\times_U U'$ 以及
$\mathcal{S}_{X'}=\mathcal{S}_X\otimes_{\mathcal{O}_U}\mathcal{O}_{U'}$
（\egaref{II.3.5.4-zh}{3.5.4}），换言之，
$\mathcal{S}_{X'}=j^*(\mathcal{S}_X)$，其中 $j$ 是投影 $X'\to X$。
因此（\egaref{II.1.5.2-zh}{1.5.2}）有
$C_{U'}=C_U\times_X X'=C_U\times_U U'$，从而图表
\egaref{II.8.7.3.3-zh}{8.7.3.3} 的交换性是立即的。
\end{proof}
% Locked French authority: ega2-1-fr.tex, 820505 B,
% SHA-256 84EDBE3E83530AF2959B441796337C9DC21EAFCA6A13114A26778760FBF437AC.
% Sequential source scope: lines 15496-15545, opening of Remark II.8.7.4(i).

\begin{remark}[8.7.4]
\label{II.8.7.4-zh}
\begin{enumerate}
  \item[\rm{(i)}] \egaref{II.8.7.3-zh}{8.7.3} 论证的末尾可立即推广如下。
  设 $g:Y'\to Y$ 是一个态射，
  $\mathcal{S}'=g^*(\mathcal{S})$，$X'=\operatorname{Proj}(\mathcal{S}')$；
  则有交换图表
  \begin{equation}
  \label{II.8.7.4.1-zh}
    \xymatrix{
      C_{X'}\ar[r]\ar[d] &
      \operatorname{Proj}(\mathcal{S}'^{\natural})\ar[d]\\
      C_X\ar[r] & \operatorname{Proj}(\mathcal{S}^{\natural})
    }
  \tag{8.7.4.1}
  \end{equation}

  另一方面，设 $\varphi:\mathcal{S}''\to\mathcal{S}$ 是分次
  $\mathcal{O}_Y$-代数层的一个同态；若置
  $X''=\operatorname{Proj}(\mathcal{S}'')$，并假设
  $u=\operatorname{Proj}(\varphi):X\to X''$ 处处有定义。此外，还有

  \oldpage[II]{176}
  一个 $Y$-态射 $v:C\to C''$（其中
  $C''=\operatorname{Spec}(\mathcal{S}'')$），使得
  $\mathcal{A}(v)=\varphi$；由于 $\varphi$ 是分次代数的同态，故由
  $\varphi$ 得到分次代数的一个 $v$-态射
  $\psi:\mathcal{S}''^{\natural}\to\mathcal{S}^{\natural}$
  （\egaref{II.1.4.1-zh}{1.4.1}）。此外，由
  \egaref{II.8.2.10-zh}{8.2.10, (iv)} 以及关于 $\varphi$ 的假设，
  容易得出 $\operatorname{Proj}(\psi)$ 处处有定义。最后，考虑到
  \egaref{II.3.5.6.1-zh}{3.5.6.1}，有一个典范的 $u$-态射
  $\mathcal{S}_{X''}\to\mathcal{S}_X$，因而由
  \egaref{II.1.5.6-zh}{1.5.6} 得到一个态射 $w:C_{X''}\to C_X$。
  在此情形下，图表
  \begin{equation}
  \label{II.8.7.4.2-zh}
    \xymatrix{
      C_{X''}\ar[r]^{\sim}\ar[d]_{w} &
      \operatorname{Proj}(\mathcal{S}''^{\natural})
        \ar[d]^{\operatorname{Proj}(\psi)}\\
      C_X\ar[r]^{\sim} & \operatorname{Proj}(\mathcal{S}^{\natural})
    }
  \tag{8.7.4.2}
  \end{equation}
  交换；将其归结到 $Y$ 为仿射的情形即可立即验证。
% Locked French authority: ega2-1-fr.tex, 820505 B,
% SHA-256 84EDBE3E83530AF2959B441796337C9DC21EAFCA6A13114A26778760FBF437AC.
% Sequential source scope: lines 15546-15586, continuation and closure of Remark II.8.7.4.

  \item[\rm{(ii)}] 注意，由
  \egaref{II.8.7.2.5-zh}{8.7.2.5} 和
  \egaref{II.8.7.2.6-zh}{8.7.2.6}，对每个 $m>0$ 有
  $\mathcal{I}_1^m\subset\mathcal{I}_m\subset\mathcal{I}_1$。
  按定义，$\mathcal{I}_1=(\mathcal{S}_+)^{\sim}$，因而
  $\mathcal{I}_1$ 在 $C$ 中定义闭子预概形
  $\varepsilon(Y)$（\egaref{II.1.4.10-zh}{1.4.10} 和
  \egaref{II.8.3.2-zh}{8.3.2}）；由此可知，对每个 $m>0$，
  \emph{$\mathcal{O}_C/\mathcal{I}_m$ 的支集包含在顶点预概形
  $\varepsilon(Y)$ 的底空间中}；在去顶点仿射锥 $E$ 的逆像上，结构态射
  $\operatorname{Proj}(\mathcal{S}^{\natural})\to C$ 因而化为一个
  \emph{同构}（这也由 \egaref{II.8.7.3-zh}{8.7.3} 和
  \egaref{II.8.7.1-zh}{8.7.1} 得出）。此外，把 $C$ 典范地等同于
  $\widehat{C}$ 的一个开集（\egaref{II.8.3.3-zh}{8.3.3}），可显然把
  $\mathcal{O}_C$ 的理想 $\mathcal{I}_m$ 延拓为
  $\mathcal{O}_{\widehat{C}}$ 的理想 $\mathcal{J}_m$，其条件是在
  $\widehat{C}$ 的开集 $\widehat{E}$ 上与
  $\mathcal{O}_{\widehat{C}}$ 相等。若置
  $\mathcal{T}=\bigoplus_{n\geqslant 0}\mathcal{J}_m$，这是一个拟凝聚分次
  $\mathcal{O}_{\widehat{C}}$-代数，则可把同构
  \egaref{II.8.7.3.1-zh}{8.7.3.1} 延拓为一个 $\widehat{C}$-同构
  \begin{equation}
  \label{II.8.7.4.3-zh}
    \widehat{C}_X\xrightarrow{\sim}\operatorname{Proj}(\mathcal{T}).
  \tag{8.7.4.3}
  \end{equation}

  事实上，在 $\widehat{E}$ 上，由上述结果，
  $\operatorname{Proj}(\mathcal{T})$ 典范地等同于
  $\widehat{E}$；于是，通过要求同构
  \egaref{II.8.7.4.3-zh}{8.7.4.3} 在 $\widehat{E}$ 上与典范同构
  $\widehat{E}_X\to\widehat{E}$（\egaref{II.8.6.2-zh}{8.6.2}）相合，
  来定义该同构；显然，这时该同构与
  \egaref{II.8.7.3.1-zh}{8.7.3.1} 在 $E$ 上相合。
\end{enumerate}
\end{remark}
% Locked French authority: ega2-1-fr.tex, 820505 B,
% SHA-256 84EDBE3E83530AF2959B441796337C9DC21EAFCA6A13114A26778760FBF437AC.
% Sequential source scope: lines 15588-15625, Corollary II.8.7.5 and proof.

\begin{corollary}[8.7.5]
\label{II.8.7.5-zh}
假设存在 $n_0>0$，使得
\begin{equation}
\label{II.8.7.5.1-zh}
  \mathcal{S}_{n+1}=\mathcal{S}_1\mathcal{S}_n
  \qquad\text{对 }n\geqslant n_0.
\tag{8.7.5.1}
\end{equation}
则 $C_X$ 的顶点子预概形（同构于 $X$）是 $C$ 的顶点子预概形
（同构于 $Y$）在典范态射 $r:C_X\to C$ 下的逆像。反之，若此性质成立，
并且还假设 $Y$ 是诺特的、$\mathcal{S}$ 为有限型，则存在 $n_0>0$，
使得 \egaref{II.8.7.5.1-zh}{8.7.5.1} 成立。
\end{corollary}

\begin{proof}
第一项断言在 $Y$ 上是局部的，故为证明它，可假设
$Y=\operatorname{Spec}(A)$ 是仿射的，而 $\mathcal{S}=\widetilde{S}$，
其中 $S$ 是正次数分次 $A$-代数。断言于是由
\egaref{II.8.2.12-zh}{8.2.12} 得出，因为
$\operatorname{Proj}(S^{\natural}\otimes_S S_0)=
C_X\times_C\varepsilon(Y)$（依照等同
\egaref{II.8.7.3.1-zh}{8.7.3.1}），也就是说，此预概形是
$\varepsilon(Y)$ 在 $C_X$ 中的逆像
（\egaref{I.4.4.1-zh}{I, 4.4.1}）。当 $Y$ 是仿射诺特的且 $S$ 为有限型时，
逆命题也由 \egaref{II.8.2.12-zh}{8.2.12} 得出。

\oldpage[II]{177}
若 $Y$ 是诺特的（不必仿射）且 $\mathcal{S}$ 为有限型，则 $Y$ 有一个由
诺特仿射开集 $U_i$ 构成的有限覆盖；由上述结果，对每个 $i$，存在整数
$n_i$，使得当 $n\geqslant n_i$ 时，
$\mathcal{S}_{n+1}|U_i=(\mathcal{S}_1|U_i)(\mathcal{S}_n|U_i)$；于是这些
$n_i$ 中的最大者满足 \egaref{II.8.7.5.1-zh}{8.7.5.1}。
\end{proof}
% Locked French authority: ega2-1-fr.tex, 820505 B,
% SHA-256 84EDBE3E83530AF2959B441796337C9DC21EAFCA6A13114A26778760FBF437AC.
% Sequential source scope: lines 15627-15656, Environment II.8.7.6.

\begin{env}[8.7.6]
\label{II.8.7.6-zh}
现在考虑由在仿射锥 $C$ 中\emph{暴涨顶点子预概形}
$\varepsilon(Y)$ 而得到的 $C$-预概形 $Z$；按定义
（\egaref{II.8.1.3-zh}{8.1.3}），它就是预概形
$\operatorname{Proj}(\bigoplus_{n\geqslant 0}\mathcal{S}_+^n)$；典范单射
\begin{equation}
\label{II.8.7.6.1-zh}
  \iota:\bigoplus_{n\geqslant 0}\mathcal{S}_+^n
  \longrightarrow\mathcal{S}^{\natural}
\tag{8.7.6.1}
\end{equation}
（借助等同 \egaref{II.8.7.3-zh}{8.7.3}）定义一个典范的支配
$C$-态射
\begin{equation}
\label{II.8.7.6.2-zh}
  G(\iota)\longrightarrow Z
\tag{8.7.6.2}
\end{equation}
其中 $G(\iota)$ 是 $C_X$ 的一个开集（\egaref{II.3.5.1-zh}{3.5.1}）。
注意，可能有 $G(\iota)\ne C_X$；例子是
$Y=\operatorname{Spec}(K)$，其中 $K$ 是一个域，
$\mathcal{S}=\widetilde{S}$ 且 $S=K[\mathbf{y}]$，这里
$\mathbf{y}$ 是一个\emph{次数为 2}的不定元。若 $R_n$ 表示集合
$(S_+)^n$，把它视为 $S_{[n]}=S_n^{\natural}$ 的子集，则
$S_+^{\natural}$ 不是 $S^{\natural}$ 中由诸 $R_n$ 的并所生成的理想的根
（参见 \egaref{II.2.3.14-zh}{2.3.14}）。
\end{env}
% Locked French authority: ega2-1-fr.tex, 820505 B,
% SHA-256 84EDBE3E83530AF2959B441796337C9DC21EAFCA6A13114A26778760FBF437AC.
% Sequential source scope: lines 15658-15698, Corollary II.8.7.7, proof, and Remark II.8.7.8.

\begin{corollary}[8.7.7]
\label{II.8.7.7-zh}
假设存在 $n_0>0$，使得
\begin{equation}
\label{II.8.7.7.1-zh}
  \mathcal{S}_n=\mathcal{S}_1^n
  \qquad\text{对 }n\geqslant n_0.
\tag{8.7.7.1}
\end{equation}
则典范态射 \egaref{II.8.7.6.2-zh}{8.7.6.2} 处处有定义，并且是同构
$C_X\xrightarrow{\sim}Z$。反之，若此性质成立，并且还假设 $Y$ 是诺特的、
$\mathcal{S}$ 为有限型，则存在 $n_0$，使得
\egaref{II.8.7.7.1-zh}{8.7.7.1} 成立。
\end{corollary}

\begin{proof}
事实上，第一项断言在 $Y$ 上是局部的，因而由
\egaref{II.8.2.14-zh}{8.2.14} 得出；如
\egaref{II.8.7.5-zh}{8.7.5} 中那样论证，逆命题也是如此。
\end{proof}

\begin{remark}[8.7.8]
\label{II.8.7.8-zh}
由于条件 \egaref{II.8.7.7.1-zh}{8.7.7.1} 蕴含
\egaref{II.8.7.5.1-zh}{8.7.5.1}，可见当它成立时，不仅 $C_X$ 可等同于
由暴涨仿射锥 $C$ 的顶点（等同于 $Y$）所得到的预概形，而且 $C_X$ 的
顶点（等同于 $X$）可等同于 $C$ 的顶点 $Y$ 的逆像闭子预概形。此外，
假设 \egaref{II.8.7.7.1-zh}{8.7.7.1} 蕴含：在
$X=\operatorname{Proj}(\mathcal{S})$ 上，$\mathcal{O}_X$-模
$\mathcal{O}_X(n)$ 是可逆的（\egaref{II.3.2.5-zh}{3.2.5} 和
\egaref{II.3.2.9-zh}{3.2.9}），且
$\mathcal{O}_X(n)=\mathcal{L}^{\otimes n}$，其中
$\mathcal{L}=\mathcal{O}_X(1)$（\egaref{II.3.2.7-zh}{3.2.7} 和
\egaref{II.3.2.9-zh}{3.2.9}）；按定义
（\egaref{II.8.6.1.1-zh}{8.6.1.1}），$C_X$ 因而是 $X$ 上的
\emph{向量丛} $\mathbf{V}(\mathcal{L})$，而其顶点是该向量丛的
\emph{零截面}。
\end{remark}
% Locked French authority: ega2-1-fr.tex, 820505 B,
% SHA-256 84EDBE3E83530AF2959B441796337C9DC21EAFCA6A13114A26778760FBF437AC.
% Sequential source scope: lines 15700-15745, subsection II.8.8 and
% Environment II.8.8.1 / printed pp.177-178.

\subsection{丰沛层与收缩}
\label{subsection:II.8.8-zh}

\begin{env}[8.8.1]
\label{II.8.8.1-zh}
设 $Y$ 是一个预概形，$f:X\to Y$ 是一个\emph{分离且拟紧}的态射，
$\mathcal{L}$ 是一个\emph{相对于 $f$ 丰沛}的可逆
$\mathcal{O}_X$-模。考虑按正次数分次的 $\mathcal{O}_Y$-代数层
\begin{equation}
\label{II.8.8.1.1-zh}
  \mathcal{S}=\mathcal{O}_Y\oplus
  \bigoplus_{n\geqslant 1}f_*(\mathcal{L}^{\otimes n})
\tag{8.8.1.1}
\end{equation}

\oldpage[II]{178}
它是拟凝聚的（\egaref{I.9.2.2-zh}{I, 9.2.2, a}）。有分次
$\mathcal{O}_X$-代数层的典范同态
\begin{equation}
\label{II.8.8.1.2-zh}
  \tau:f^*(\mathcal{S})\longrightarrow
  \bigoplus_{n\geqslant 0}\mathcal{L}^{\otimes n}
\tag{8.8.1.2}
\end{equation}
在次数 $\geqslant 1$ 上，它与典范同态
$\sigma:f^*(f_*(\mathcal{L}^{\otimes n}))\to\mathcal{L}^{\otimes n}$
（\egaxref{0.4.4.3-zh}{0, 4.4.3}）一致，而在次数 $0$ 上是恒等映射。
$\mathcal{L}$ 为 $f$-丰沛的假设于是蕴含
（\egaref{II.4.6.3-zh}{4.6.3} 和 \egaref{II.3.6.1-zh}{3.6.1}）：
相应的 $Y$-态射
\begin{equation}
\label{II.8.8.1.3-zh}
  r=r_{\mathcal{L},\tau}:X\longrightarrow
  P=\operatorname{Proj}(\mathcal{S})
\tag{8.8.1.3}
\end{equation}
处处有定义，并且是一个\emph{支配开浸入}；此外有
\begin{equation}
\label{II.8.8.1.4-zh}
  r^*(\mathcal{O}_P(n))=\mathcal{L}^{\otimes n}
  \qquad\text{对每个 }n\in\mathbf{Z}.
\tag{8.8.1.4}
\end{equation}
\end{env}
% Locked French authority: ega2-1-fr.tex, 820505 B,
% SHA-256 84EDBE3E83530AF2959B441796337C9DC21EAFCA6A13114A26778760FBF437AC.
% Sequential source scope: lines 15747-15779, Proposition II.8.8.2.

\begin{proposition}[8.8.2]
\label{II.8.8.2-zh}
设 $C=\operatorname{Spec}(\mathcal{S})$ 是\emph{由
$\mathcal{S}$ 定义的仿射锥}；若 $\mathcal{L}$ 是 $f$-丰沛的，则存在一个典范
$Y$-态射
\begin{equation}
\label{II.8.8.2.1-zh}
  g:V=\mathbf{V}(\mathcal{L})\longrightarrow C
\tag{8.8.2.1}
\end{equation}
使得图表
\begin{equation}
\label{II.8.8.2.2-zh}
  \xymatrix{
    X\ar[r]^{j}\ar[d]_{f} &
    \mathbf{V}(\mathcal{L})\ar[r]^{\pi}\ar[d]^{g} & X\ar[d]^{f}\\
    Y\ar[r]_{\varepsilon} & C\ar[r]_{\psi} & Y
  }
\tag{8.8.2.2}
\end{equation}
交换；其中 $\psi$ 和 $\pi$ 是结构态射，$j$ 和 $\varepsilon$ 是分别把
$X$ 和 $Y$ 映到 $\mathbf{V}(\mathcal{L})$ 的零截面以及 $C$ 的顶点预概形
上的典范浸入。此外，$g$ 在 $\mathbf{V}(\mathcal{L})-j(X)$ 上的限制是
开浸入
\begin{equation}
\label{II.8.8.2.3-zh}
  \mathbf{V}(\mathcal{L})-j(X)\longrightarrow
  E=C-\varepsilon(Y)
\tag{8.8.2.3}
\end{equation}
到去顶点仿射锥 $E$ 中，该锥与 $\mathcal{S}$ 相应。
\end{proposition}
% Locked French authority: ega2-1-fr.tex, 820505 B,
% SHA-256 84EDBE3E83530AF2959B441796337C9DC21EAFCA6A13114A26778760FBF437AC.
% Sequential source scope: lines 15781-15842, proof of Proposition II.8.8.2.

\begin{proof}
沿用 \egaref{II.8.8.1-zh}{8.8.1} 的记号，令
$\mathcal{S}_P^{\geqslant}=\bigoplus_{n\geqslant 0}\mathcal{O}_P(n)$，并令
$C_P=\operatorname{Spec}(\mathcal{S}_P^{\geqslant})$。已知
（\egaref{II.8.6.2-zh}{8.6.2}），存在一个典范态射
$h=\operatorname{Spec}(\alpha):C_P\to C$，使图表
\begin{equation}
\label{II.8.8.2.4-zh}
  \xymatrix{
    C_P\ar[r]\ar[d]_{h} & P\ar[d]^{p}\\
    C\ar[r]_{\psi} & Y
  }
\tag{8.8.2.4}
\end{equation}
交换；此外，若 $\varepsilon_P:P\to C_P$ 是典范浸入，则图表
\begin{equation}
\label{II.8.8.2.5-zh}
  \xymatrix{
    P\ar[r]^{\varepsilon_P}\ar[d]_{p} & C_P\ar[d]^{h}\\
    Y\ar[r]_{\varepsilon} & C
  }
\tag{8.8.2.5}
\end{equation}
交换（\egaref{II.8.7.1.1-zh}{8.7.1.1}）；最后，$h$ 在去顶点仿射锥
$E_P$ 上的限制是一个\emph{同构}
$E_P\xrightarrow{\sim}E$（\egaref{II.8.6.2-zh}{8.6.2}）。另一方面，
由 \egaref{II.8.8.1.4-zh}{8.8.1.4} 得
\[
  r^*(\mathcal{S}_P^{\geqslant})=
  \mathbf{S}_{\mathcal{O}_X}(\mathcal{L})
\]

\oldpage[II]{179}
因而有一个典范 $P$-态射
$q:\mathbf{V}(\mathcal{L})\to C_P$；交换图
\begin{equation}
\label{II.8.8.2.6-zh}
  \xymatrix{
    \mathbf{V}(\mathcal{L})\ar[r]^{\pi}\ar[d]_{q} & X\ar[d]^{r}\\
    C_P\ar[r] & P
  }
\tag{8.8.2.6}
\end{equation}
将 $\mathbf{V}(\mathcal{L})$ 等同于纤维积 $C_P\times_P X$
（\egaref{II.1.5.2-zh}{1.5.2}）。由于 $r$ 是开浸入，故 $q$ 也是开浸入
（\egaref{I.4.3.2-zh}{I, 4.3.2}）。此外，由
\egaref{II.8.5.2-zh}{8.5.2}，$q$ 在
$\mathbf{V}(\mathcal{L})-j(X)$ 上的限制将该预概形映入 $E_P$，且图表
\begin{equation}
\label{II.8.8.2.7-zh}
  \xymatrix{
    X\ar[r]^{j}\ar[d]_{r} & \mathbf{V}(\mathcal{L})\ar[d]^{q}\\
    P\ar[r]_{\varepsilon_P} & C_P
  }
\tag{8.8.2.7}
\end{equation}
交换（这是 \egaref{II.8.5.1.3-zh}{8.5.1.3} 的特殊情形）。取 $g$ 为复合态射
$h\circ q$，则 \egaref{II.8.8.2-zh}{8.8.2} 的各断言立即由这些事实得出。
\end{proof}
% Locked French authority: ega2-1-fr.tex, 820505 B,
% SHA-256 84EDBE3E83530AF2959B441796337C9DC21EAFCA6A13114A26778760FBF437AC.
% Sequential source scope: lines 15844-15877, Remark II.8.8.3.

\begin{remark}[8.8.3]
\label{II.8.8.3-zh}
进一步假设 $Y$ 是一个\emph{诺特}预概形，$f$ 是一个\emph{固有}态射。
由于这时 $r$ 是\emph{固有的}（\egaref{II.5.4.4-zh}{5.4.4}），因而是闭的，
而且它是支配开浸入，所以 $r$ 必为一个\emph{同构}
$X\xrightarrow{\sim}P$。此外，我们将在第三章
（\egaref{III.2.3.5.1-zh}{III, 2.3.5.1}）看到，这时 $\mathcal{S}$ 必为一个
\emph{有限型} $\mathcal{O}_Y$-代数。于是由此得出，
$\mathcal{S}^{\natural}$ 是一个\emph{有限型}
$\mathcal{S}_0^{\natural}$-代数（\egaref{II.8.2.10-zh}{8.2.10, (i)} 和
\egaref{II.8.7.2.7-zh}{8.7.2.7}）；由于 $C_P$ 在 $C$ 上同构于
$\operatorname{Proj}(\mathcal{S}^{\natural})$
（\egaref{II.8.7.3-zh}{8.7.3}），可见态射 $h:C_P\to C$ 是
\emph{射影的}；由于态射 $r$ 是同构，
$q:\mathbf{V}(\mathcal{L})\to C_P$ 也是同构，由此可知态射
$g:\mathbf{V}(\mathcal{L})\to C$ 是\emph{射影的}。此外，由于 $h$ 在
$E_P$ 上的限制是到 $E$ 上的一个同构，而 $q$ 是同构，所以 $g$ 的限制
\egaref{II.8.8.2.3-zh}{8.8.2.3} 是一个\emph{同构}
$\mathbf{V}(\mathcal{L})-j(X)\xrightarrow{\sim}E$。

若进一步假设 $\mathcal{L}$ 对于 $f$ 是\emph{极丰沛的}，我们也将在第三章
（\egaref{III.2.3.5.1-zh}{III, 2.3.5.1}）看到，存在整数 $n_0>0$，使得
$\mathcal{S}_n=\mathcal{S}_1^n$ 对 $n\geqslant n_0$ 成立。于是由
\egaref{II.8.7.7-zh}{8.7.7} 可知，$\mathbf{V}(\mathcal{L})$ 可等同于
由\emph{在仿射锥 $C$ 中暴涨顶点子预概形}（等同于 $Y$）而得到的预概形
$Z$，并且 $\mathbf{V}(\mathcal{L})$ 的\emph{零截面}（等同于 $Y$）是
$C$ 的顶点子预概形 $Y$ 的\emph{逆像}。事实上，上述结果的一部分无须
诺特假设也能建立：
\end{remark}
% Locked French authority: ega2-1-fr.tex, 820505 B,
% SHA-256 84EDBE3E83530AF2959B441796337C9DC21EAFCA6A13114A26778760FBF437AC.
% Sequential source scope: lines 15879-15917, Corollary II.8.8.4 and opening proof.

\begin{corollary}[8.8.4]
\label{II.8.8.4-zh}
设 $Y$ 是一个预概形（相应地，一个拟紧概形），$f:X\to Y$ 是固有态射，
$\mathcal{L}$ 是对于 $f$ 丰沛的可逆 $\mathcal{O}_X$-模。则态射
\egaref{II.8.8.2.1-zh}{8.8.2.1} 是固有的（相应地，是射影的），而其限制
\egaref{II.8.8.2.3-zh}{8.8.2.3} 是一个同构。
\end{corollary}

\begin{proof}
为证明 $g$ 是固有的，可化归到 $Y$ 为仿射的情形，因而只须考虑 $Y$ 是
拟紧概形的情形。与 \egaref{II.8.8.3-zh}{8.8.3} 中相同的论证首先表明，
$r$ 是一个\emph{同构} $X\xrightarrow{\sim}P$；于是 $q$ 也是同构，并且由于
$h$ 在 $E_P$ 上的限制是同构 $E_P\xrightarrow{\sim}E$，已经可见
\egaref{II.8.8.2.3-zh}{8.8.2.3} 是同构。还须证明 $g$ 是\emph{射影的}。

由于按假设 $f$ 是有限型的，可把 \egaref{II.3.8.5-zh}{3.8.5} 应用于
\egaref{II.8.8.1.2-zh}{8.8.1.2} 中的同态

\oldpage[II]{180}
$\tau$：于是存在整数 $d>0$ 和 $\mathcal{S}_d$ 的一个有限型拟凝聚
$\mathcal{O}_Y$-子模 $\mathcal{E}$，使得若 $\mathcal{S}'$ 是由
$\mathcal{E}$ 生成的 $\mathcal{S}$ 的 $\mathcal{O}_Y$-子代数，并且
$\tau'=\tau\circ f^*(\varphi)$（$\varphi$ 是典范单射
$\mathcal{S}'\to\mathcal{S}$），则 $r'=r_{\mathcal{L},\tau'}$ 是一个浸入
\[
  X\longrightarrow P'=\operatorname{Proj}(\mathcal{S}').
\]
此外，由于 $\varphi$ 是单射，$r'$ 也是一个\emph{支配浸入}
（\egaref{II.3.7.6-zh}{3.7.6}）；与对 $r$ 所作的相同论证于是表明，$r'$ 是
一个\emph{满闭浸入}；由于 $r'$ 分解为
$X\xrightarrow{r}\operatorname{Proj}(\mathcal{S})
\xrightarrow{\Phi}\operatorname{Proj}(\mathcal{S}')$，其中
$\Phi=\operatorname{Proj}(\varphi)$，由此得出 $\Phi$ 也是一个
% Locked French authority: ega2-1-fr.tex, 820505 B,
% SHA-256 84EDBE3E83530AF2959B441796337C9DC21EAFCA6A13114A26778760FBF437AC.
% Sequential source scope: lines 15918-15957, continuation of the proof of Corollary II.8.8.4.

\emph{满闭浸入}。但这蕴含 $\Phi$ 是一个\emph{同构}；事实上，可化归到
$Y=\operatorname{Spec}(A)$ 为仿射、$\mathcal{S}=\widetilde{S}$、
$\mathcal{S}'=\widetilde{S'}$ 的情形，其中 $S$ 是一个分次 $A$-代数，
$S'$ 是 $S$ 的一个分次子代数。对每个齐次元 $t\in S'$，
$S'_{(t)}$ 是 $S_{(t)}$ 的一个子环；回到
$\operatorname{Proj}(\varphi)$ 的定义（\egaref{II.2.8.1-zh}{2.8.1}），
可见问题化为证明：若 $B'$ 是环 $B$ 的一个子环，而与典范单射
$B'\to B$ 相应的态射 $\operatorname{Spec}(B)\to\operatorname{Spec}(B')$
是闭浸入，则此态射必为一个\emph{同构}；但这由
\egaref{I.4.2.3-zh}{I, 4.2.3} 得出。

此外，有
$\Phi^*(\mathcal{O}_{P'}(n))=\mathcal{O}_P(n)$
（\egaref{II.3.5.2-zh}{3.5.2, (ii)} 和
\egaref{II.3.5.4-zh}{3.5.4}），故
$r'^*(\mathcal{O}_{P'}(n))$ 同构于 $\mathcal{L}^{\otimes n}$
（\egaref{II.4.6.3-zh}{4.6.3}）。置
$\mathcal{S}''=\mathcal{S}'^{(d)}$，于是
（\egaref{II.3.1.8-zh}{3.1.8, (i)}）$X$ 典范地等同于
$P''=\operatorname{Proj}(\mathcal{S}'')$，而
$\mathcal{L}''=\mathcal{L}^{\otimes d}$ 等同于 $\mathcal{O}_{P''}(1)$
（\egaref{II.3.2.9-zh}{3.2.9, (ii)}）。

在此情形，若 $C''=\operatorname{Spec}(\mathcal{S}'')$，则
$\mathcal{S}_{P''}^{\geqslant}=
\bigoplus_{n\geqslant 0}\mathcal{O}_{P''}(n)$ 等同于
$\bigoplus_{n\geqslant 0}\mathcal{L}''^{\otimes n}$，因而
$C_{P''}=\operatorname{Spec}(\mathcal{S}_{P''}^{\geqslant})$ 等同于
$\mathbf{V}(\mathcal{L}'')$；另一方面，由
\egaref{II.8.7.3-zh}{8.7.3}，$C_{P''}$ 在 $C''$ 上同构于
$\operatorname{Proj}(\mathcal{S}''^{\natural})$；按 $\mathcal{S}''$ 的定义，
$\mathcal{S}''^{\natural}$ 由 $\mathcal{S}_1''^{\natural}$ 生成，并且
$\mathcal{S}_1''^{\natural}$ 在
$\mathcal{S}_0''^{\natural}=\mathcal{S}''$ 上是有限型的
（\egaref{II.8.2.10-zh}{8.2.10, (i) 和 (iii)}），故
$\operatorname{Proj}(\mathcal{S}''^{\natural})$ 在 $C''$ 上是
\emph{射影的}（\egaref{II.5.5.1-zh}{5.5.1}）。于是考虑图表
% Locked French authority: ega2-1-fr.tex, 820505 B,
% SHA-256 84EDBE3E83530AF2959B441796337C9DC21EAFCA6A13114A26778760FBF437AC.
% Sequential source scope: lines 15958-15997, continuation and closure of the proof of Corollary II.8.8.4.

\begin{equation}
\label{II.8.8.4.1-zh}
  \xymatrix{
    \mathbf{V}(\mathcal{L})\ar[r]^{g}\ar[d]_{u} &
    \operatorname{Spec}(\mathcal{S})=C\ar[d]^{v}\\
    \mathbf{V}(\mathcal{L}'')\ar[r]_{g''} &
    \operatorname{Spec}(\mathcal{S}'')=C''
  }
\tag{8.8.4.1}
\end{equation}
其中，$g$ 和 $g''$ 由 \egaref{II.1.5.6-zh}{1.5.6} 对应于典范
$f$-态射
\[
  \mathcal{S}\longrightarrow
  \bigoplus_{n\geqslant 0}\mathcal{L}^{\otimes n}
  \quad\text{且}\quad
  \mathcal{S}''\longrightarrow
  \bigoplus_{n\geqslant 0}\mathcal{L}''^{\otimes n}
\]
（\egaref{II.3.3.2.3-zh}{3.3.2.3}）（参见下文
\egaref{II.8.8.5-zh}{8.8.5}）；$v$ 对应于嵌入态射
$\mathcal{S}''\to\mathcal{S}$，而 $u$ 对应于嵌入态射
$\bigoplus_{n\geqslant 0}\mathcal{L}^{\otimes nd}\to
\bigoplus_{n\geqslant 0}\mathcal{L}^{\otimes n}$；由
\egaref{II.3.3.2-zh}{3.3.2} 立即可见此图表交换。我们刚刚看到 $g''$ 是
射影态射；另一方面，$u$ 是一个\emph{有限}态射。事实上，问题在 $X$ 上
是局部的，故可假设 $X$ 是环 $A$ 的仿射概形且
$\mathcal{L}=\mathcal{O}_X$；这时一切化为注意到，环 $A[T]$ 是其子环
$A[T^d]$ 上的有限型模（$T$ 为不定元）。由于 $Y$ 是拟紧概形，而
$C''$ 在 $Y$ 上仿射，$C''$ 也是拟紧概形，

\oldpage[II]{181}
因此 $g''\circ u$ 是射影态射（\egaref{II.5.5.5-zh}{5.5.5, (ii)}）；
由 \egaref{II.8.8.4.1-zh}{8.8.4.1} 的交换性，$v\circ g$ 也是射影态射；
又因为 $v$ 是仿射的，因而是分离的，最终可知 $g$ 是射影的
（\egaref{II.5.5.5-zh}{5.5.5, (v)}）。
\end{proof}
% Locked French authority: ega2-1-fr.tex, 820505 B,
% SHA-256 84EDBE3E83530AF2959B441796337C9DC21EAFCA6A13114A26778760FBF437AC.
% Sequential source scope: lines 15999-16062, Environment II.8.8.5.

\begin{env}[8.8.5]
\label{II.8.8.5-zh}
回到 \egaref{II.8.8.1-zh}{8.8.1} 的情形。我们将看到，可以给态射
$g:\mathbf{V}(\mathcal{L})\to C$ 另一个定义；这个定义对任意可逆
$\mathcal{O}_X$-模 $\mathcal{L}$（不必丰沛）都适用。为此，考虑与
\egaref{II.8.8.1.2-zh}{8.8.1.2} 的态射 $\tau$ 相应的 $f$-态射
\begin{equation}
\label{II.8.8.5.1-zh}
  \tau^\flat:\mathcal{S}\longrightarrow
  \bigoplus_{n\geqslant 0}\mathcal{L}^{\otimes n}
\tag{8.8.5.1}
\end{equation}
由此（\egaref{II.1.5.6-zh}{1.5.6}）得出一个态射 $g':V\to C$，使得若
$\pi:V\to X$ 和 $\psi:C\to Y$ 是结构态射，则图表
\begin{equation}
\label{II.8.8.5.2-zh}
  \xymatrix{
    X\ar[d]_{f} & V\ar[l]^{\pi}\ar[d]^{g'}\\
    Y & C\ar[l]_{\psi}
  }
  \qquad
  \xymatrix{
    X\ar[r]^{j}\ar[d]_{f} & V\ar[d]^{g'}\\
    Y\ar[r]_{\varepsilon} & C
  }
\tag{8.8.5.2}
\end{equation}
交换（\egaref{II.8.5.1.2-zh}{8.5.1.2} 和
\egaref{II.8.5.1.3-zh}{8.5.1.3}）。我们将证明，（若假设
$\mathcal{L}$ 对于 $f$ 丰沛）态射 $g$ 与 $g'$ 相同。

事实上，问题在 $Y$ 上是局部的，故可假设
$Y=\operatorname{Spec}(A)$ 为仿射的，并且（由于
\egaref{II.8.8.1.3-zh}{8.8.1.3}）把 $X$ 等同于
$P=\operatorname{Proj}(S)$ 的一个开集，其中
\[
  S=A\oplus\bigoplus_{n\geqslant 1}
  \Gamma(X,\mathcal{L}^{\otimes n})~;
\]
于是由 \egaref{II.8.8.1.4-zh}{8.8.1.4} 推出，对每个
$n\in\mathbf{Z}$，有
$\Gamma(X,\mathcal{O}_P(n))=\Gamma(X,\mathcal{L}^{\otimes n})$。
考虑 $h=\operatorname{Spec}(\alpha)$ 的定义，其中 $\alpha$ 是典范
$p$-态射
$\widetilde{S}\to\mathcal{S}_P^{\geqslant}$
（\egaref{II.8.6.1.2-zh}{8.6.1.2}）；所须验证的是，
$\alpha^\sharp:p^*(\widetilde{S})\to\mathcal{S}_P^{\geqslant}$ 在
$X$ 上的限制等于 $\tau$。结合
\egaxref{0.4.4.3-zh}{0, 4.4.3}，问题化为验证：把典范同态
$\alpha_n:S_n\to\Gamma(P,\mathcal{O}_P(n))$ 与限制同态
\[
  \Gamma(P,\mathcal{O}_P(n))\longrightarrow
  \Gamma(X,\mathcal{O}_P(n))=
  \Gamma(X,\mathcal{L}^{\otimes n})
\]
复合后，对每个 $n>0$ 都得到恒等映射；而这立即由代数 $S$ 的定义与
$\alpha_n$ 的定义（\egaref{II.2.6.2-zh}{2.6.2}）得出。
\end{env}
% Locked French authority: ega2-1-fr.tex, 820505 B,
% SHA-256 84EDBE3E83530AF2959B441796337C9DC21EAFCA6A13114A26778760FBF437AC.
% Sequential source scope: lines 16064-16105, Proposition II.8.8.6 and opening of its proof.

\begin{proposition}[8.8.6]
\label{II.8.8.6-zh}
假设（沿用 \egaref{II.8.8.5-zh}{8.8.5} 的记号）当置
$f=(f_0,\lambda)$ 时，同态
$\lambda:\mathcal{O}_Y\to f_*(\mathcal{O}_X)$ 是双射；则：
\begin{enumerate}
  \item[\rm{(i)}] 若置 $g=(g_0,\mu)$，则
  $\mu:\mathcal{O}_C\to g_*(\mathcal{O}_V)$ 是同构。
  \item[\rm{(ii)}] 若 $X$ 是整的（相应地，局部整且正规的），则
  $C$ 是整的（相应地，正规的）。
\end{enumerate}
\end{proposition}

\begin{proof}
事实上，此时 $f$-态射 $\tau^\flat$ 是一个\emph{同构}
\[
  \tau^\flat:\mathcal{S}=\psi_*(\mathcal{O}_C)
  \longrightarrow f_*(\pi_*(\mathcal{O}_V))
  =\psi_*(g_*(\mathcal{O}_V))
\]
而 $Y$-态射 $g$ 可看作满足如下条件的态射：同态 $\mathcal{A}(g)$
（\egaref{II.1.1.2-zh}{1.1.2}）等于 $\tau^\flat$。为验证
$\mu$ 是 $\mathcal{O}_C$-模的同构，只需
（\egaref{II.1.4.2-zh}{1.4.2}）验证
\[
  \mathcal{A}(\mu):\psi_*(\mathcal{O}_C)
  \longrightarrow\psi_*(g_*(\mathcal{O}_V))
\]
是同构。但由定义（\egaref{II.1.1.2-zh}{1.1.2}），有
$\mathcal{A}(\mu)=\mathcal{A}(g)$，故得 (i) 的结论。

为证明 (ii)，可归结到 $Y$ 为仿射的情形，因而
$\mathcal{S}=\widetilde{S}$，其中

\oldpage[II]{182}
$S=\bigoplus_{n\geqslant 0}\Gamma(X,\mathcal{L}^{\otimes n})$；
$X$ 是整的这一假设蕴含环 $S$ 是整环
（\egaref{I.7.4.4-zh}{I, 7.4.4}），因而 $C$ 也是整的
（\egaref{I.5.1.4-zh}{I, 5.1.4}）。为证明 $C$ 是正规的，我们将使用
下面的引理：
% Locked French authority: ega2-1-fr.tex, 820505 B,
% SHA-256 84EDBE3E83530AF2959B441796337C9DC21EAFCA6A13114A26778760FBF437AC.
% Sequential source scope: lines 16107-16137, Lemma II.8.8.6.1 and closure of the proof of II.8.8.6.

\begin{lemma}[8.8.6.1]
\label{II.8.8.6.1-zh}
设 $Z$ 是一个正规整预概形。则环 $\Gamma(Z,\mathcal{O}_Z)$ 是整闭整环。
\end{lemma}

\begin{proof}
事实上，由 \egaref{I.8.2.1.1-zh}{I, 8.2.1.1}，
$\Gamma(Z,\mathcal{O}_Z)$ 是预概形 Z 的有理函数域 $R(Z)$ 中诸整闭环
$\mathcal{O}_z$（$z\in Z$）的交。
\end{proof}

在此情形，先证明 $V$ 是\emph{局部整且正规的}；为此可只考虑
$X=\operatorname{Spec}(A)$ 为仿射的情形，其中环 $A$ 是整闭整环
（\egaref{II.6.3.8-zh}{6.3.8}），且 $\mathcal{L}=\mathcal{O}_X$。
这时 $V=\operatorname{Spec}(A[T])$，而 $A[T]$ 是整闭整环
（[24]，第 99 页），这就建立了我们的断言。另一方面，对 $C$ 的每个
仿射开集 $U$，$g^{-1}(U)$ 是拟紧的，因为态射 $g$ 是拟紧的；由于 $V$
是局部整的，$g^{-1}(U)$ 的连通分支是 $g^{-1}(U)$ 中的整开预概形，因而
只有有限多个；并且由于 $V$ 是正规的，这些预概形也是正规的
（\egaref{II.6.3.8-zh}{6.3.8}）。于是，依照 (i)，
$\Gamma(U,\mathcal{O}_C)$ 等于
$\Gamma(g^{-1}(U),\mathcal{O}_V)$，它是有限多个整闭整环的直积
（\egaref{II.8.8.6.1-zh}{8.8.6.1}），这表明 $C$ 是正规的
（\egaref{II.6.3.4-zh}{6.3.4}）。
\end{proof}
% Locked French authority: ega2-1-fr.tex, 820505 B,
% SHA-256 84EDBE3E83530AF2959B441796337C9DC21EAFCA6A13114A26778760FBF437AC.
% Sequential source scope: lines 16139-16201, subsection II.8.9 through the explanatory paragraph after Corollary II.8.9.2.

\subsection{格劳尔特丰沛性判据：陈述}
\label{subsection:II.8.9-zh}

我们将证明，\egaref{II.8.8.2-zh}{8.8.2} 中所证明的性质\emph{刻画}了
对于 $f$ 丰沛的 $\mathcal{O}_X$-模；确切地说，要证明下述判据：

\begin{theorem}[8.9.1]
\label{II.8.9.1-zh}
\emph{（格劳尔特判据）。}设 $Y$ 是一个预概形，$p:X\to Y$ 是分离拟紧态射，
$\mathcal{L}$ 是可逆 $\mathcal{O}_X$-模。$\mathcal{L}$ 对于 $p$ 丰沛，
当且仅当存在一个 $Y$-预概形 $C$、$C$ 的一个 $Y$-截面
$\varepsilon:Y\to C$，以及一个 $Y$-态射
$q:\mathbf{V}(\mathcal{L})\to C$，使得下列性质成立：
\begin{enumerate}
  \item[\rm{(i)}] 图表
  \begin{equation}
  \label{II.8.9.1.1-zh}
    \xymatrix{
      X\ar[r]^{j}\ar[d]_{p} & \mathbf{V}(\mathcal{L})\ar[d]^{q}\\
      Y\ar[r]_{\varepsilon} & C
    }
  \tag{8.9.1.1}
  \end{equation}
  交换，其中 $j$ 是向量丛 $\mathbf{V}(\mathcal{L})$ 的零截面。
  \item[\rm{(ii)}] $q$ 在 $\mathbf{V}(\mathcal{L})-j(X)$ 上的限制是一个
  拟紧开浸入
  \[
    \mathbf{V}(\mathcal{L})-j(X)\longrightarrow C
  \]
  且其像不与 $\varepsilon(Y)$ 相交。
\end{enumerate}
\end{theorem}

注意，若 $C$ 在 $Y$ 上是\emph{分离的}，则在条件 (ii) 中可去掉开浸入
拟紧这一假设；事实上，为看出后一性质是其他条件的推论，可只考虑 $Y$
为仿射的情形，这时断言由 \egaref{I.5.5.1-zh}{I, 5.5.1, (i)} 和
\egaref{I.5.5.10-zh}{5.5.10} 得出。当假设 $X$ 是诺特的时，也可去掉

\oldpage[II]{183}
同一假设，因为这时 $V$ 也是诺特的，而断言由
\egaref{I.6.3.5-zh}{I, 6.3.5} 得出。

\begin{corollary}[8.9.2]
\label{II.8.9.2-zh}
若态射 $p:X\to Y$ 是固有的，则在 \egaref{II.8.9.1-zh}{8.9.1} 的陈述中，
可假设 $q$ 是固有的，并把“开浸入”替换为“同构”。
\end{corollary}

形象地说（当 $p:X\to Y$ 是固有的时），$\mathcal{L}$ 对于 $p$ 丰沛，
当且仅当可以把向量丛 $\mathbf{V}(\mathcal{L})$ 的零截面“收缩”为基预概形
$Y$。一个重要的特殊情形是 $Y$ 为某个域的谱；此时“收缩”操作就是把
$\mathbf{V}(\mathcal{L})$ 的零截面\emph{收缩为一个点}。
% Locked French authority: ega2-1-fr.tex, 820505 B,
% SHA-256 84EDBE3E83530AF2959B441796337C9DC21EAFCA6A13114A26778760FBF437AC.
% Sequential source scope: lines 16203-16271, Environment II.8.9.3.

\begin{env}[8.9.3]
\label{II.8.9.3-zh}
\egaref{II.8.9.1-zh}{8.9.1} 中各条件的必要性以及推论
\egaref{II.8.9.2-zh}{8.9.2}，立即由
\egaref{II.8.8.2-zh}{8.8.2} 和 \egaref{II.8.8.4-zh}{8.8.4} 得出。

为证明 \egaref{II.8.9.1-zh}{8.9.1} 中各条件的充分性，我们将考虑一个
稍微更一般的情形。为此，沿用 \egaref{II.8.8.2-zh}{8.8.2} 的记号，置
\[
  \mathcal{S}'=\bigoplus_{n\geqslant 0}\mathcal{L}^{\otimes n}
\]
以及
\[
  V=\mathbf{V}(\mathcal{L})=\operatorname{Spec}(\mathcal{S}').
\]
闭子预概形 $j(X)$ 是 $\mathbf{V}(\mathcal{L})$ 的零截面，它由
$\mathcal{O}_V$ 的拟凝聚理想层
$\mathcal{J}=(\mathcal{S}'_+)^{\sim}$ 定义
（\egaref{II.1.4.10-zh}{1.4.10}）。这个 $\mathcal{O}_V$-模是
\emph{可逆的}，因为问题在 $X$ 上是局部的，而这化为注意到：在多项式环
$A[T]$ 中，理想 $TA[T]$ 是秩一自由 $A[T]$-模。此外，仍因问题在
$X$ 上是局部的，立即有
\[
  \mathcal{L}=j^*(\mathcal{J})
\]
以及
\[
  j_*(\mathcal{L})=\mathcal{J}/\mathcal{J}^2.
\]
另一方面，若
\[
  \pi:\mathbf{V}(\mathcal{L})\longrightarrow X
\]
是结构态射，则有
$\pi_*(\mathcal{J})=\mathcal{S}'_+$ 和
$\pi_*(\mathcal{J}/\mathcal{J}^2)=\mathcal{L}$；因此有两个典范同态
$\mathcal{L}\to\pi_*(\mathcal{J})\to\mathcal{L}$：第一个是典范单射
$\mathcal{L}\to\mathcal{S}'_+$，第二个是 $\mathcal{S}'_+$ 到
$\mathcal{S}'_1=\mathcal{L}$ 上的典范投影，而它们的复合是恒等映射。
此外，还可把
\[
  \pi_*(\mathcal{J})=\mathcal{S}'_+
  =\bigoplus_{n\geqslant 1}\mathcal{L}^{\otimes n}
\]
典范地嵌入\emph{乘积}
\[
  \prod_{n\geqslant 1}\mathcal{L}^{\otimes n}
  =\varprojlim_n\pi_*(\mathcal{J}/\mathcal{J}^{n+1})
\]
中（因为
$\pi_*(\mathcal{J}/\mathcal{J}^{n+1})=
\mathcal{L}\oplus\mathcal{L}^{\otimes2}\oplus\cdots\oplus
\mathcal{L}^{\otimes n}$），因而又有两个典范同态
\begin{equation}
\label{II.8.9.3.1-zh}
  \mathcal{L}\longrightarrow
  \varprojlim_n\pi_*(\mathcal{J}/\mathcal{J}^{n+1})
  \longrightarrow\mathcal{L}
\tag{8.9.3.1}
\end{equation}
它们的复合是恒等映射。

在此情形下，我们将证明的 \egaref{II.8.9.1-zh}{8.9.1} 的推广如下：
\end{env}
% Locked French authority: ega2-1-fr.tex, 820505 B,
% SHA-256 84EDBE3E83530AF2959B441796337C9DC21EAFCA6A13114A26778760FBF437AC.
% Sequential source scope: lines 16273-16318, Proposition II.8.9.4 / printed pp.183-184.

\begin{proposition}[8.9.4]
\label{II.8.9.4-zh}
设 $Y$ 是一个预概形，$V$ 是一个 $Y$-预概形，$X$ 是 $V$ 的闭子预概形，
它由理想层 $\mathcal{J}$ 定义；$\mathcal{O}_V$ 的这个理想层是可逆
$\mathcal{O}_V$-模。若 $j:X\to V$ 是

\oldpage[II]{184}
典范单射，则置
$\mathcal{L}=j^*(\mathcal{J})=
\mathcal{J}\otimes_{\mathcal{O}_V}\mathcal{O}_X$，从而
$j_*(\mathcal{L})=\mathcal{J}/\mathcal{J}^2$。假设结构态射
$p:X\to Y$ 是分离且拟紧的，并且满足下列条件：
\begin{enumerate}
  \item[\rm{(i)}] 存在一个有限型 $Y$-态射 $\pi:V\to X$，使得
  $\pi\circ j=1_X$，因而
  $\pi_*(\mathcal{J}/\mathcal{J}^2)=\mathcal{L}$。
  \item[\rm{(ii)}] 存在一个 $\mathcal{O}_X$-模同态
  \[
    \varphi:\mathcal{L}\longrightarrow
    \varprojlim_n\pi_*(\mathcal{J}/\mathcal{J}^{n+1})
  \]
  使得复合映射
  \[
    \mathcal{L}\xrightarrow{\varphi}
    \varprojlim_n\pi_*(\mathcal{J}/\mathcal{J}^{n+1})
    \xrightarrow{\alpha}\pi_*(\mathcal{J}/\mathcal{J}^2)=\mathcal{L}
  \]
  为恒等映射（其中 $\alpha$ 是典范同态）。
  \item[\rm{(iii)}] 存在一个 $Y$-预概形 $C$、一个 $Y$-截面
  $\varepsilon$（其目标为 $C$），以及一个 $Y$-态射 $q:V\to C$，使图表
  \begin{equation}
  \label{II.8.9.4.1-zh}
    \xymatrix{
      X\ar[r]^{j}\ar[d]_{p} & V\ar[d]^{q}\\
      Y\ar[r]_{\varepsilon} & C
    }
  \tag{8.9.4.1}
  \end{equation}
  交换。
  \item[\rm{(iv)}] $q$ 在 $W=V-j(X)$ 上的限制是到 $C$ 中的拟紧开浸入，
  且其像不与 $\varepsilon(Y)$ 相交。
\end{enumerate}
在这些条件下，$\mathcal{L}$ 对于 $p$ 丰沛。
\end{proposition}
% Locked French authority: ega2-1-fr.tex, 820505 B,
% SHA-256 84EDBE3E83530AF2959B441796337C9DC21EAFCA6A13114A26778760FBF437AC.
% Sequential source scope: lines 16320-16364, subsection II.8.10 and Lemma II.8.10.1.

\subsection{格劳尔特丰沛性判据：证明}
\label{subsection:II.8.10-zh}

\begin{lemma}[8.10.1]
\label{II.8.10.1-zh}
设 $\pi:V\to X$ 是一个态射，$j:X\to V$ 是一个 $X$-截面，其取值于
$V$，并且是闭浸入；设 $\mathcal{J}$ 是 $\mathcal{O}_V$ 的拟凝聚理想，
它定义 $V$ 中与 $j$ 相伴的闭子预概形。
\begin{enumerate}
  \item[\rm{(i)}] 对每个 $n\geqslant 0$，
  $\pi_*(\mathcal{O}_V/\mathcal{J}^{n+1})$ 和
  $\pi_*(\mathcal{J}/\mathcal{J}^{n+1})$ 都是拟凝聚
  $\mathcal{O}_X$-模，并且有
  $\pi_*(\mathcal{O}_V/\mathcal{J})=\mathcal{O}_X$、
  $\pi_*(\mathcal{J}/\mathcal{J}^2)=j^*(\mathcal{J})$。
  \item[\rm{(ii)}] 若 $X=\{\xi\}=\operatorname{Spec}(k)$，其中 $k$ 是域，
  则 $\varprojlim_n\pi_*(\mathcal{O}_V/\mathcal{J}^{n+1})$ 同构于局部环
  $\mathcal{O}_{j(\xi)}$ 关于 $\mathfrak{m}_{j(\xi)}$-预进拓扑的分离完备化。
  \item[\rm{(iii)}] 假设 $\mathcal{J}$ 是可逆
  $\mathcal{O}_V$-模（这蕴含
  \[
    \mathcal{L}=j^*(\mathcal{J})=
    \pi_*(\mathcal{J}/\mathcal{J}^2)
  \]
  是可逆 $\mathcal{O}_X$-模），并且存在同态
  \[
    \varphi:\mathcal{L}\longrightarrow
    \varprojlim_n\pi_*(\mathcal{J}/\mathcal{J}^{n+1})
  \]
  使得复合
  \[
    \mathcal{L}\xrightarrow{\varphi}
    \varprojlim_n\pi_*(\mathcal{J}/\mathcal{J}^{n+1})
    \xrightarrow{\alpha}\pi_*(\mathcal{J}/\mathcal{J}^2)
  \]
  （$\alpha$ 是典范同态）为恒等映射。若置
  $\mathcal{S}=\bigoplus_{n\geqslant0}\mathcal{L}^{\otimes n}$，则从
  $\varphi$ 典范地得到如下 $\mathcal{O}_X$-代数同构：完备化
  $\widehat{\mathcal{S}}$，即 $\mathcal{S}$ 相对于其典范滤过的完备化
  （该完备化同构于直积
  $\prod_{n\geqslant0}\mathcal{L}^{\otimes n}$）同构到
  $\varprojlim_n\pi_*(\mathcal{O}_V/\mathcal{J}^{n+1})$。
\end{enumerate}
\end{lemma}
% Locked French authority: ega2-1-fr.tex, 820505 B,
% SHA-256 84EDBE3E83530AF2959B441796337C9DC21EAFCA6A13114A26778760FBF437AC.
% Sequential source scope: lines 16366-16408, opening proof of Lemma II.8.10.1 through part (i).

\begin{proof}
首先注意，$\mathcal{O}_V$-模
$\mathcal{O}_V/\mathcal{J}^{n+1}$ 的支集是 $j(X)$，而
$\mathcal{J}/\mathcal{J}^{n+1}$ 的支集包含在 $j(X)$ 中。在 (ii) 的情形，
$j(X)$ 是 $V$ 的闭点 $j(\xi)$，

\oldpage[II]{185}
而按定义，$\pi_*(\mathcal{O}_V/\mathcal{J}^{n+1})$ 是
$\mathcal{O}_V/\mathcal{J}^{n+1}$ 在点 $j(\xi)$ 处的纤维；也就是说，置
$C=\mathcal{O}_{j(\xi)}$，并以 $\mathfrak{m}$ 表示 $C$ 的极大理想，
它就是 $C$-模 $C/\mathfrak{m}^{n+1}$；于是断言 (ii) 是显然的。

为证明 (i)，注意问题在 $X$ 上是局部的；故可限于 $X$ 为仿射的情形。
设 $U$ 是 $V$ 中的一个仿射开集；$j(X)\cap U$ 是 $j(X)$ 中的仿射开集，
因而与其同构的 $U_0=\pi(j(X)\cap U)$ 是 $X$ 中的仿射开集；对 $X$ 中每个
仿射开集 $W_0\subset U_0$，$W=\pi^{-1}(W_0)\cap U$ 是 $V$ 中的仿射开集，
因为 $X$ 是概形（\egaref{I.5.5.10-zh}{I, 5.5.10}）；特别地，
$U'=U\cap\pi^{-1}(U_0)$ 是 $V$ 中的仿射开集，并且显然有
$\pi(U')=U_0$ 和 $j(U_0)=j(X)\cap U$。在此情形，按定义
\[
  \Gamma(W_0,\pi_*(\mathcal{O}_V/\mathcal{J}^{n+1}))
  =\Gamma(\pi^{-1}(W_0),\mathcal{O}_V/\mathcal{J}^{n+1})~;
\]
但 $\pi^{-1}(W_0)$ 中每个不属于 $j(W_0)$ 的点都有一个在
$\pi^{-1}(W_0)$ 中的开邻域，它不与 $j(X)$ 相交，且在其中
$\mathcal{O}_V/\mathcal{J}^{n+1}$ 因而为零；故显然，
$\mathcal{O}_V/\mathcal{J}^{n+1}$ 在 $\pi^{-1}(W_0)$ 上的截面与在 $W$
上的截面一一对应。换言之，若 $\pi'$ 是 $\pi$ 在 $U'$ 上的限制，则
$(\mathcal{O}_X|U_0)$-模
\[
  \pi_*(\mathcal{O}_V/\mathcal{J}^{n+1})|U_0
  \quad\hbox{和}\quad
  \pi'_*\bigl((\mathcal{O}_V/\mathcal{J}^{n+1})|U'\bigr)
\]
相同。由于 $U'$ 和 $U_0$ 是仿射的，而且这些 $U_0$ 覆盖 $X$，可知
（\egaref{I.1.6.3-zh}{I, 1.6.3}）
$\pi_*(\mathcal{O}_V/\mathcal{J}^{n+1})$ 是拟凝聚的；对于
$\pi_*(\mathcal{J}/\mathcal{J}^{n+1})$，证明完全相同。
% Locked French authority: ega2-1-fr.tex, 820505 B,
% SHA-256 84EDBE3E83530AF2959B441796337C9DC21EAFCA6A13114A26778760FBF437AC.
% Sequential source scope: lines 16409-16451, part (iii) and closure of
% the proof of Lemma II.8.10.1.

最后，为证明 (iii)，注意 $\mathcal{S}$ 就是
$\mathbf{S}_{\mathcal{O}_X}(\mathcal{L})$；因此由 $\varphi$ 典范地得出
$\mathcal{O}_X$-代数层的同态
\[
  \psi:\mathcal{S}\longrightarrow
  \varprojlim_n\pi_*(\mathcal{O}_V/\mathcal{J}^{n+1})
\]
（\egaref{II.1.7.4-zh}{1.7.4}）；此外，这个同态把
$\mathcal{L}^{\otimes n}$ 映入
$\varprojlim_m\pi_*(\mathcal{J}^n/\mathcal{J}^{m+1})$，所以对于所考虑的
拓扑是连续的，因而确实延拓为同态
\[
  \widehat{\psi}:\widehat{\mathcal{S}}\longrightarrow
  \varprojlim_n\pi_*(\mathcal{O}_V/\mathcal{J}^{n+1}).
\]
为看出这是一个同构，可像 (i) 的证明中那样，化归到
$X=\operatorname{Spec}(A)$ 和 $V=\operatorname{Spec}(B)$ 都是仿射的情形，
其中 $\mathcal{J}=\widetilde{\mathfrak{J}}$，而 $\mathfrak{J}$ 是 $B$ 的
一个理想；与 $\pi$ 相应的是一个单射 $A\to B$，它把 $A$ 等同于 $B$ 的
一个与 $\mathfrak{J}$ 互补的子环；$\mathcal{L}$（相应地，
$\pi_*(\mathcal{O}_V/\mathcal{J}^{n+1})$）是与 $A$-模
$L=\mathfrak{J}/\mathfrak{J}^2$（相应地，$B/\mathfrak{J}^{n+1}$）相伴的
拟凝聚 $\mathcal{O}_X$-模。由于 $\mathcal{J}$ 是一个\emph{可逆}
$\mathcal{O}_V$-模，还可假设 $\mathfrak{J}=Bt$，其中 $t$ 在 $B$ 中不是
零因子。于是由关系 $B=A\oplus Bt$，对每个 $n>0$ 推出
\[
  B=A\oplus At\oplus At^2\oplus\cdots\oplus At^n\oplus Bt^{n+1}
\]
因而有形式幂级数环 $A[[T]]$ 到
$C=\varprojlim B/\mathfrak{J}^{n+1}$ 上的一个典范 $A$-同构，使 $t$ 对应于
$T$。另一方面，有 $L=A\bar t$，其中 $\bar t$ 是 $t$ 模 $Bt^2$ 的类；
按假设，同态 $\varphi$ 把 $\bar t$ 映到一个元素 $t'\in C$，它模
$Ct^2$ 与 $t$ 同余。由此对 $n$ 作归纳得
\[
  A\oplus At'\oplus\cdots\oplus At'{}^n\oplus Ct^{n+1}
  =A\oplus At\oplus\cdots\oplus At^n\oplus Ct^{n+1}
\]
这证明同态 $\widehat{\psi}$ 确实对应于
$\prod_{n\geqslant0}L^{\otimes n}$ 到 $C$ 上的一个同构。
\end{proof}
% Locked French authority: ega2-1-fr.tex, 820505 B,
% SHA-256 84EDBE3E83530AF2959B441796337C9DC21EAFCA6A13114A26778760FBF437AC.
% Sequential source scope: lines 16453-16501, Lemma II.8.10.2 and proof /
% printed pp.185-186.

\begin{lemma}[8.10.2]
\label{II.8.10.2-zh}
在 \egaref{II.8.10.1-zh}{8.10.1} 的假设下，设
$g:X'\to X$ 是一个态射，

\oldpage[II]{186}
设 $V'=V\times_X X'$，又设 $\pi':V'\to X'$、
$g':V'\to V$ 是典范投影，从而有交换图
\[
  \xymatrix{
    V\ar[d]_{\pi} & V'\ar[l]^{g'}\ar[d]^{\pi'}\\
    X & X'\ar[l]_{g}
  }
\]
则 $j'=j\times1_{X'}$ 是 $V'$ 的一个 $X'$-截面，并且是闭浸入；
$\mathcal{J}'=(g')^*(\mathcal{J})\mathcal{O}_{V'}$ 是
$\mathcal{O}_{V'}$ 的拟凝聚理想，它定义了与 $j'$ 相伴的 $V'$ 的闭子预概形。
此外，有
\[
  \pi'_*(\mathcal{O}_{V'}/\mathcal{J}'^{n+1})
  =g^*(\pi_*(\mathcal{O}_V/\mathcal{J}^{n+1})).
\]
最后，$\mathcal{J}'$ 作为 $\mathcal{O}_{V'}$-模典范同构于
$(g')^*(\mathcal{J})$；特别地，若 $\mathcal{J}$ 是可逆
$\mathcal{O}_V$-模，则它也是可逆的。
\end{lemma}

\begin{proof}
$j'$ 是闭浸入这一事实由 \egaref{I.4.3.1-zh}{I, 4.3.1} 得出；由基预概形
扩张的函子性，它是 $V'$ 的一个 $X'$-截面。此外，若 $Z$（相应地，$Z'$）
是与 $j$（相应地，$j'$）相伴的 $V$（相应地，$V'$）的闭子预概形，则有
$Z'=g'{}^{-1}(Z)$（\egaref{I.4.3.1-zh}{I, 4.3.1}），第二项断言于是由
\egaref{I.4.4.5-zh}{I, 4.4.5} 得出。为证明其余断言，像
\egaref{II.8.10.1-zh}{8.10.1} 中那样，可化归到 $X$、$V$ 和 $X'$（因而
$V'$ 也）都是仿射的情形；沿用
\egaref{II.8.10.1-zh}{8.10.1} 的证明中的记号，并设
$X'=\operatorname{Spec}(A')$。于是 $V'=\operatorname{Spec}(B')$，其中
$B'=B\otimes_A A'$，并且 $\mathcal{J}'=\widetilde{\mathfrak{J}'}$，其中
$\mathfrak{J}'=\operatorname{Im}(\mathfrak{J}\otimes_A A')$。因而有
\[
  B'/\mathfrak{J}'^{n+1}=(B/\mathfrak{J}^{n+1})\otimes_A A'~;
\]
此外，因为 $\mathfrak{J}$ 作为 $A$-模是 $B$ 的直和因子，
$\mathfrak{J}\otimes_A A'$ 作为 $A'$-模也是 $B'$ 的直和因子，故典范地
等同于 $\mathfrak{J}'$。
\end{proof}
% Locked French authority: ega2-1-fr.tex, 820505 B,
% SHA-256 84EDBE3E83530AF2959B441796337C9DC21EAFCA6A13114A26778760FBF437AC.
% Sequential source scope: lines 16503-16527, Corollary II.8.10.3 and proof.

\begin{corollary}[8.10.3]
\label{II.8.10.3-zh}
假设 \egaref{II.8.10.1-zh}{8.10.1} 的各项假设成立，并进一步假设
$\pi$ 是有限型的，且 $\mathcal{J}$ 是可逆 $\mathcal{O}_V$-模。则对每个
$x\in X$，在点 $j(x)$ 处，纤维 $\pi^{-1}(x)$ 的局部环是维数为 1 的
正则环（因而是整环），其完备化同构于形式幂级数环 $k(x)[[T]]$
（$T$ 是不定元）；此外，$\pi^{-1}(x)$ 只有一个包含 $j(x)$ 的不可约分支。
\end{corollary}

\begin{proof}
由于 $\pi^{-1}(x)=V\times_X\operatorname{Spec}(k(x))$，由
\egaref{II.8.10.2-zh}{8.10.2} 可归结到 $X$ 是某个域 $K$ 的谱的情形。
由于 $\pi$ 是有限型的（\egaxref{I.6.3.4-zh}{I, 6.3.4, (iv)}），
$\mathcal{O}_{j(x)}$ 此时是诺特局部环，因而对于
$\mathfrak{m}_{j(x)}$-预进拓扑是分离的
（\egaxref{0.7.3.5-zh}{0, 7.3.5}）；由
\egaref{II.8.10.1-zh}{8.10.1, (ii) 和 (iii)} 可知，该环的完备化同构于
$K[[T]]$，从而 $\mathcal{O}_{j(x)}$ 是维数为 1 的正则环
（[1]，第 17-01 页，定理 1）；最后，由于 $\mathcal{O}_{j(x)}$ 是整环，
$j(x)$ 只属于 $V$ 的不可约分支（数目有限）中的一个
（\egaxref{I.5.1.4-zh}{I, 5.1.4}）。
\end{proof}
% Locked French authority: ega2-1-fr.tex, 820505 B,
% SHA-256 84EDBE3E83530AF2959B441796337C9DC21EAFCA6A13114A26778760FBF437AC.
% Sequential source scope: lines 16529-16561, Corollary II.8.10.4 and proof / printed pp.186-187.

\begin{corollary}[8.10.4]
\label{II.8.10.4-zh}
假设 \egaref{II.8.10.1-zh}{8.10.1} 的各项假设成立，并进一步假设
$\mathcal{J}$ 是可逆 $\mathcal{O}_V$-模。令 $W=V-j(X)$；对任意拟凝聚理想
$\mathcal{K}$（作为 $\mathcal{O}_X$ 的理想），置
$\mathcal{K}_V=\pi^*(\mathcal{K})\mathcal{O}_V$ 和
$\mathcal{K}_W=\mathcal{K}_V|W$。那么 $\mathcal{K}_V$ 是
$\mathcal{O}_V$ 的如下最大拟凝聚理想：它在 $W$ 上的限制为
$\mathcal{K}_W$。
\end{corollary}

\begin{proof}
事实上，像在 \egaref{II.8.10.1-zh}{8.10.1} 中一样，可以看出问题在 $X$
和 $V$ 上是局部的；所以可沿用 \egaref{II.8.10.1-zh}{8.10.1} 的证明中的
记号，其中 $\mathfrak{J}=Bt$，而 $t$ 在 $B$ 中不是零因子。此外，
$W=\operatorname{Spec}(B_t)$，且 $\mathcal{K}=\widetilde{\mathfrak{K}}$，
其中 $\mathfrak{K}$ 是 $A$ 的理想；因此
$\pi^*(\mathcal{K})\mathcal{O}_V=(\mathfrak{K}.B)^{\sim}$
（\egaref{I.1.6.9-zh}{I, 1.6.9}），
$\mathcal{K}_W=(\mathfrak{K}.B_t)^{\sim}$；而最大者

\oldpage[II]{187}
——即 $B$ 中典范像在 $B_t$ 内等于 $\mathfrak{K}.B_t$ 的理想中最大者——
是 $\mathfrak{K}.B_t$ 的原像，也就是所有满足下述条件的 $s\in B$ 的集合：
存在整数 $n>0$，使 $t^ns\in\mathfrak{K}.B$。需证明后一关系蕴含
$s\in\mathfrak{K}.B$；等价地，需证明 $t$ 的典范像在
$B/\mathfrak{K}B=(A/\mathfrak{K})\otimes_A B$ 中不是零因子。这由
\egaref{II.8.10.2-zh}{8.10.2} 应用于
$X'=\operatorname{Spec}(A/\mathfrak{K})$ 得出。
\end{proof}
% Locked French authority: ega2-1-fr.tex, 820505 B,
% SHA-256 84EDBE3E83530AF2959B441796337C9DC21EAFCA6A13114A26778760FBF437AC.
% Sequential source scope: lines 16563-16593, Corollary II.8.10.5.

\begin{corollary}[8.10.5]
\label{II.8.10.5-zh}
假设 \egaref{II.8.10.3-zh}{8.10.3} 的诸假设成立；设
$W=V-j(X)$，$x$ 是 $X$ 的一个点，$\mathcal{K}$ 是
$\mathcal{O}_X$ 的一个拟凝聚理想，$z$ 是 $\pi^{-1}(x)$ 中含有 $j(x)$
的不可约分支的泛点（\egaref{II.8.10.3-zh}{8.10.3}）。
\begin{enumerate}
  \item[\rm{(i)}] 设 $g$ 是 $\mathcal{O}_V$ 在 $V$ 上的一个截面，使得
  $g|W$ 是 $\mathcal{K}_W$ 在 $W$ 上的一个截面（记号见
  \egaref{II.8.10.4-zh}{8.10.4}）。则 $g$ 是 $\mathcal{K}_V$ 的一个截面；
  若还有 $g(z)\neq0$，并且对每个整数 $m>0$，以 $g_m^x$ 表示 $g$ 在
  $\Gamma(X,\pi_*(\mathcal{O}_V/\mathcal{J}^{m+1}))$ 中典范像 $g_m$ 在点
  $x$ 处的芽，则存在整数 $m>0$，使 $g_m^x$ 在
  \[
    (\pi_*(\mathcal{O}_V/\mathcal{J}^{m+1}))_x
    \otimes_{\mathcal{O}_x}k(x)
  \]
  中的像满足 $\neq0$。
  \item[\rm{(ii)}] 进一步假设
  \egaref{II.8.10.1-zh}{8.10.1, (iii)} 的条件成立。于是，若存在
  $\mathcal{K}_V$ 在 $V$ 上的截面 $g$，使得 $g(z)\neq0$，则存在整数
  $n\geqslant0$ 和
  $\mathcal{K}.\mathcal{L}^{\otimes n}=
  \mathcal{K}\otimes\mathcal{L}^{\otimes n}\subset\mathcal{L}^{\otimes n}$
  的一个截面 $f$，使得 $f(x)\neq0$。若 $g$ 是 $\mathcal{J}$ 的一个截面，
  则可取 $n>0$。
\end{enumerate}
\end{corollary}
% Locked French authority: ega2-1-fr.tex, 820505 B,
% SHA-256 84EDBE3E83530AF2959B441796337C9DC21EAFCA6A13114A26778760FBF437AC.
% Sequential source scope: lines 16595-16630, proof of Corollary II.8.10.5.

\begin{proof}
\begin{enumerate}
  \item[\rm{(i)}] 按假设，由 $g|W$ 生成的 $\mathcal{O}_W$ 的理想包含在
  $\mathcal{K}_W$ 中，故由 \egaref{II.8.10.4-zh}{8.10.4}，由 $g$ 生成的
  $\mathcal{O}_V$ 的理想包含在 $\mathcal{K}_V$ 中；换言之，$g$ 是
  $\mathcal{K}_V$ 的一个截面。为证明 (i) 的第二项断言，仍可假设 $X$ 和
  $V$ 是仿射的，并沿用 \egaref{II.8.10.1-zh}{8.10.1} 的记号；这时纤维
  $\pi^{-1}(x)$ 是环 $B'=B\otimes_A k(x)$ 的仿射概形，而且 $B'$ 中存在
  非零因子 $t'$，使得 $B'=k(x)\oplus B't'$。由于 $j(x)$ 是 $z$ 的一个
  特化，且 $g(z)\neq0$，必有 $g_{j(x)}\neq0$。但
  $\mathcal{O}_{j(x)}$ 是一个分离局部环（\egaref{II.8.10.3-zh}{8.10.3}），
  因而嵌入其完备化中，所以 $g$ 在该完备化中的像非零。该完备化同构于
  $\varprojlim_n(B'/B't'^{n+1})$（\egaref{II.8.10.3-zh}{8.10.3}）；
  若 $g'=g\otimes1\in B'$，则存在整数 $m$，使得
  $g'\notin B't'^{m+1}$；也就是说，$g'$ 在 $B'/B't'^{m+1}$ 中的像
  $g'_m$ 非零。由于 $g'_m$ 正是 $g_m^x$ 的像，我们的断言得证。
  \item[\rm{(ii)}] 由 \egaref{II.8.10.1-zh}{8.10.1, (iii)}，
  $\pi_*(\mathcal{O}_V/\mathcal{J}^{m+1})$ 同构于诸
  $\mathcal{L}^{\otimes k}$（$0\leqslant k\leqslant m$）的直和；以 $f_k$
  表示 $\mathcal{L}^{\otimes k}$ 在 $X$ 上的截面，它是由此同构对应于
  $g_m$ 的元素
  $\bigoplus_{k=0}^m\Gamma(X,\mathcal{L}^{\otimes k})$ 的分量。按 (i) 选择
  $m$，则由 (i) 存在指标 $k$，使得 $f_k(x)\neq0$。为看出 $f_k$ 是
  $\mathcal{K}\mathcal{L}^{\otimes k}$ 的一个截面，只须如上考虑 $X$ 和
  $V$ 为仿射的情形；这时它立即由 $g\in\mathfrak{K}.B$ 得出（记号见
  \egaref{II.8.10.4-zh}{8.10.4}）。最后一项断言来自以下事实：假设
  $g\in\Gamma(V,\mathcal{J})$ 蕴含 $f_0=0$。
\end{enumerate}
\end{proof}
% Locked French authority: ega2-1-fr.tex, 820505 B,
% SHA-256 84EDBE3E83530AF2959B441796337C9DC21EAFCA6A13114A26778760FBF437AC.
% Sequential source scope: lines 16632-16686, opening of Environment II.8.10.6.

\begin{env}[8.10.6]
\label{II.8.10.6-zh}
\emph{\egaref{II.8.9.4-zh}{8.9.4} 的证明。}问题在 $Y$ 上是局部的
（\egaref{II.4.6.4-zh}{4.6.4}）；由于 $\varepsilon$ 是一个 $Y$-截面，
可把 $C$ 替换为 $\varepsilon(Y)$ 的某个点的一个仿射开邻域 $U$，使得
$\varepsilon(Y)\cap U$ 在 $U$ 中是闭的。换言之，可以假设 $C$ 是仿射的，
而 $Y$ 是 $C$ 的闭子预概形（因而是仿射的），由 $\mathcal{O}_C$ 的一个
拟凝聚理想层 $\mathcal{I}$ 定义。

\oldpage[II]{188}
由于 $p$ 是分离拟紧的，这时 $X$ 是一个拟紧概形，而问题化为证明
$\mathcal{L}$ 是丰沛的（\egaref{II.4.6.4-zh}{4.6.4}）。依照判据
\egaref{II.4.5.2-zh}{4.5.2, a)}，因而必须证明如下命题：对
$\mathcal{O}_X$ 的每个拟凝聚理想 $\mathcal{K}$ 和不属于
$\mathcal{O}_X/\mathcal{K}$ 的支集的每个点 $x\in X$，存在整数 $n>0$ 和
$\mathcal{K}\otimes\mathcal{L}^{\otimes n}$ 在 $X$ 上的一个截面 $f$，
使得 $f(x)\neq0$。

为此，置
\[
  \mathcal{K}_V=\pi^*(\mathcal{K})\mathcal{O}_V
\]
以及
\[
  \mathcal{K}_W=\mathcal{K}_V|W,\qquad\hbox{其中}\qquad W=V-j(X)~;
\]
由于 $q$ 在 $W$ 上的限制是到 $C$ 中的拟紧浸入，由
\egaref{I.9.4.2-zh}{I, 9.4.2}，$\mathcal{K}_W$ 是 $\mathcal{O}_V$ 的一个
拟凝聚理想 $\mathcal{K}'_V$ 在 $W$ 上的限制，其中
\[
  \mathcal{K}'_V=q^*(\mathcal{K}_C)\mathcal{O}_V
\]
而 $\mathcal{K}_C$ 是 $\mathcal{O}_C$ 的一个拟凝聚理想。此外，按假设
$q^{-1}(Y)\subset j(X)$，而 $Y$ 由理想 $\mathcal{I}$ 定义，所以
$q^*(\mathcal{I})\mathcal{O}_V$ 在 $W$ 上的限制与 $\mathcal{O}_V$ 的限制
相同；因此 $\mathcal{K}_W$ 也是
$q^*(\mathcal{I}\mathcal{K}_C)\mathcal{O}_V$ 在 $W$ 上的限制，故可假设
$\mathcal{K}_C\subset\mathcal{I}$，于是
\begin{equation}
\label{II.8.10.6.1-zh}
  \mathcal{K}'_V\subset q^*(\mathcal{I})\mathcal{O}_V\subset\mathcal{J}
\tag{8.10.6.1}
\end{equation}
这里用到了 \egaref{I.4.4.6-zh}{I, 4.4.6} 和
\egaref{II.8.9.4.1-zh}{8.9.4.1} 的交换性。此外，由
\egaref{II.8.10.4-zh}{8.10.4} 得出
\begin{equation}
\label{II.8.10.6.2-zh}
  \mathcal{K}'_V\subset\mathcal{K}_V.
\tag{8.10.6.2}
\end{equation}
% Locked French authority: ega2-1-fr.tex, 820505 B,
% SHA-256 84EDBE3E83530AF2959B441796337C9DC21EAFCA6A13114A26778760FBF437AC.
% Sequential source scope: lines 16688-16704, conclusion of Environment II.8.10.6.

在此情形，由 \egaref{II.8.10.3-zh}{8.10.3}，$j(x)$ 只属于
$\pi^{-1}(x)$ 的一个不可约分支；设 $z$ 为此分支的泛点，并置
$z'=q(z)$。由 \egaref{II.8.10.5-zh}{8.10.5}，若我们证明存在
$\mathcal{K}'_V$ 在 $V$ 上的截面 $g$，使得 $g(z)\neq0$，则证明完成
（这里还用到 \egaref{II.8.10.6.1-zh}{8.10.6.1} 和
\egaref{II.8.10.6.2-zh}{8.10.6.2}）。按假设，$\mathcal{K}$ 在 $x$ 的
一个开邻域上的限制等于 $\mathcal{O}_X$ 的限制；另一方面，由
\egaref{II.8.10.3-zh}{8.10.3} 有 $z\neq j(x)$，故 $z\in W$，因而
$(\mathcal{K}_W)_z=\mathcal{O}_{V,z}$；由定义可知
$(\mathcal{K}_C)_{z'}=\mathcal{O}_{C,z'}$。由于 $C$ 是仿射的，存在
$\mathcal{K}_C$ 在 $C$ 上的一个截面 $g'$，使得 $g'(z')\neq0$；取与
$g'$ 典范对应的 $\mathcal{K}'_V$ 的截面 $g$，便有 $g(z)\neq0$，证明完成。
\end{env}
% Locked French authority: ega2-1-fr.tex, 820505 B,
% SHA-256 84EDBE3E83530AF2959B441796337C9DC21EAFCA6A13114A26778760FBF437AC.
% Sequential source scope: lines 16706-16731, Remark II.8.10.7 /
% printed pp.188-189.

\begin{remark}[8.10.7]
\label{II.8.10.7-zh}
我们尚不知，在 \egaref{II.8.9.4-zh}{8.9.4} 的陈述中，条件 (ii) 是否
多余。无论如何，若不假设存在一个 $Y$-态射 $\pi:V\to X$，使
$\pi\circ j=1_X$，结论就不再成立；下面简要说明如何构造一个反例，其
细节只能留待以后展开。取 $Y=\operatorname{Spec}(k)$，其中 $k$ 是一个域；
取 $C=\operatorname{Spec}(A)$，其中 $A=k[T_1,T_2]$，并令 $Y$-截面
$\varepsilon$ 对应于增广同态 $A\to k$。以 $C'$ 表示由在 $C$ 的闭点
$a=\varepsilon(Y)$ 处暴涨 $C$ 所得的概形；若 $D$ 是 $a$ 在 $C'$ 中的
逆像，则在 $D$ 中取一个闭点 $b$，

\oldpage[II]{189}
并以 $V$ 表示由在 $b$ 处暴涨 $C'$ 所得的概形；$X$ 是 $V$ 的闭子概形，
它是 $a$ 在结构态射 $q:V\to C$ 下的逆像。可以证明，$X$ 是两个不可约分支
$X_1$、$X_2$ 的并，其中 $X_1$ 是 $b$ 在 $V$ 中的逆像。显然，定义 $X$ 的
$\mathcal{O}_V$ 的理想 $\mathcal{J}$ 仍是可逆的；但通过考察
$\mathcal{L}$ 在 $X_1$ 上的逆像的“次数”，可以证明
$j^*(\mathcal{J})=\mathcal{L}$（$j$ 是典范单射 $X\to V$）并不丰沛：若
$\mathcal{L}$ 丰沛，此次数应当 $>0$，而可证明它等于 0（通过初等的
相交计算）。
\end{remark}
% Locked French authority: ega2-1-fr.tex, 820505 B,
% SHA-256 84EDBE3E83530AF2959B441796337C9DC21EAFCA6A13114A26778760FBF437AC.
% Sequential source scope: lines 16733-16796, subsection II.8.11 and Lemma II.8.11.1 with proof.

\subsection{收缩的唯一性}
\label{subsection:II.8.11-zh}

\begin{lemma}[8.11.1]
\label{II.8.11.1-zh}
设 $U$、$V$ 是两个预概形，
$h=(h_0,\lambda):U\to V$ 是满射态射。假设：
\begin{enumerate}
  \item[$1^\circ$]
  $\lambda:\mathcal{O}_V\to h_*(\mathcal{O}_U)
  =(h_0)_*(\mathcal{O}_U)$ 是同构。
  \item[$2^\circ$] $V$ 的底空间与 $U$ 的底空间关于关系
  $R:h_0(x)=h_0(y)$ 的商空间相等同（当态射 $h$ 是开的或闭的时，
  这个条件总是满足；尤其当 $h$ 是紧合的时如此）。
\end{enumerate}
则对每个预概形 $W$，映射
\begin{equation}
\label{II.8.11.1.1-zh}
  \operatorname{Hom}(V,W)\longrightarrow\operatorname{Hom}(U,W)
\tag{8.11.1.1}
\end{equation}
把每个态射 $v=(v_0,\nu)$（从 $V$ 到 $W$）对应到态射
$u=v\circ h=(u_0,\mu)$；它是从 $\operatorname{Hom}(V,W)$ 到满足如下条件的
$u$ 所成集合上的双射：$u_0$ 在每个纤维 $h_0^{-1}(x)$ 上为常值。
\end{lemma}

\begin{proof}
显然，若 $u=v\circ h$，因而 $u_0=v_0\circ h_0$，则 $u_0$ 在每个集合
$h_0^{-1}(x)$ 上为常值。反之，若 $u$ 具有此性质，我们来证明存在唯一的
$v\in\operatorname{Hom}(V,W)$ 使得 $u=v\circ h$。连续映射
$v_0:V\to W$ 满足 $u_0=v_0\circ h_0$，其存在性与唯一性由假设得出，
因为 $h_0$ 可等同于从 $U$ 到 $U/R$ 的典范映射。另一方面，必要时用一个
同构的预概形替换 $V$，可以假设 $\lambda$ 是恒等映射；由假设，$\mu$
于是是一个同态
\[
  \mu:\mathcal{O}_W\longrightarrow(u_0)_*(\mathcal{O}_U)
  =(v_0)_*((h_0)_*(\mathcal{O}_U))
\]
使得相应的同态
$\mu^\sharp:u_0^*(\mathcal{O}_W)\to\mathcal{O}_U$ 在每个纤维上都是局部的。
由于
$(v_0)_*((h_0)_*(\mathcal{O}_U))=(v_0)_*(\mathcal{O}_V)$，必有
$\nu=\mu$，问题因而归结为验证相应同态
$\nu^\sharp:v_0^*(\mathcal{O}_W)\to\mathcal{O}_V$ 在每个纤维上都是局部的。
而每个 $y\in V$ 都具有形式 $h_0(x)$，其中 $x\in U$；置
$z=v_0(y)=u_0(x)$。于是由 \egaxref{0.3.5.5-zh}{0, 3.5.5}，同态
$\mu_x^\sharp$ 分解为
\[
  \mu_x^\sharp:\mathcal{O}_z
  \xrightarrow{\nu_y^\sharp}\mathcal{O}_y
  \xrightarrow{\lambda_x^\sharp}\mathcal{O}_x.
\]
按假设，$\lambda_x^\sharp$ 和 $\mu_x^\sharp$ 是局部同态；因此
$\lambda_x^\sharp$ 把 $\mathcal{O}_y$ 的每个可逆元变为
$\mathcal{O}_x$ 的可逆元。若 $\nu_y^\sharp$ 把 $\mathcal{O}_z$ 的一个
不可逆元变为 $\mathcal{O}_y$ 的可逆元，则 $\mu_x^\sharp$ 会把
$\mathcal{O}_z$ 的这个元变为 $\mathcal{O}_x$ 的可逆元，这与假设矛盾，
故引理得证。
\end{proof}
% Locked French authority: ega2-1-fr.tex, 820505 B,
% SHA-256 84EDBE3E83530AF2959B441796337C9DC21EAFCA6A13114A26778760FBF437AC.
% Sequential source scope: lines 16798-16838, Corollary II.8.11.2, proof, and Remark II.8.11.3 / printed pp.189-190.

\begin{corollary}[8.11.2]
\label{II.8.11.2-zh}
设 $U$ 是整预概形，$V$ 是正规预概形；则任何普遍闭、双有理且纯不可分的
态射 $h:U\to V$ 都是同构。
\end{corollary}

\begin{proof}
若 $h=(h_0,\lambda)$，则由假设可知 $h_0$ 是单射且为闭映射，并且
$h_0(U)$

\oldpage[II]{190}
在 $V$ 中稠密，故 $h_0$ 是从 $U$ 到 $V$ 的一个\emph{同胚}。
为证明该推论，只需看出
$\lambda:\mathcal{O}_V\to(h_0)_*(\mathcal{O}_U)$ 是同构；这时即可应用
\egaref{II.8.11.1-zh}{8.11.1}，它将证明映射
\egaref{II.8.11.1.1-zh}{8.11.1.1} 是双射（因为每个纤维
$h_0^{-1}(x)$ 都只含一个点），因而 $h$ 是同构。问题显然在 $V$ 上是
局部的，所以可假设 $V=\operatorname{Spec}(A)$ 是仿射的，其环为整闭整环
（\egaref{II.8.8.6.1-zh}{8.8.6.1}）；此时 $h$ 对应于
（\egaref{I.2.2.4-zh}{I, 2.2.4}）一个同态
$\varphi:A\to\Gamma(U,\mathcal{O}_U)$，问题归结为证明 $\varphi$ 是同构。
若 $K$ 是 $A$ 的分式域，则按假设 $\Gamma(U,\mathcal{O}_U)$ 的分式域为
$K$，而 $A$ 是 $\Gamma(U,\mathcal{O}_U)$ 的子环，$\varphi$ 是典范单射
（\egaref{I.8.2.7-zh}{I, 8.2.7}）。由于态射 $h$ 满足
\egaref{II.7.3.11-zh}{7.3.11} 的各项假设，
$\Gamma(U,\mathcal{O}_U)$ 是 $A$ 在 $K$ 中整闭包的子环，故由假设它等于
$A$。
\end{proof}

\begin{remark}[8.11.3]
\label{II.8.11.3-zh}
我们将在第三章 \egaref{III.4.4.11-zh}{III, 4.4.11} 看到：当 $V$ 是
\emph{局部诺特}预概形时，任何固有且拟有限的态射 $h:U\to V$（特别是任何
满足 \egaref{II.8.11.2-zh}{8.11.2} 各项假设的态射）必为\emph{有限}态射。
因此，在这种情形下，\egaref{II.8.11.2-zh}{8.11.2} 的结论由
\egaref{II.6.1.15-zh}{6.1.15} 得出。
\end{remark}
% Locked French authority: ega2-1-fr.tex, 820505 B,
% SHA-256 84EDBE3E83530AF2959B441796337C9DC21EAFCA6A13114A26778760FBF437AC.
% Sequential source scope: lines 16840-16883, Environment II.8.11.4 and Lemma II.8.11.5 with proof.

\begin{env}[8.11.4]
\label{II.8.11.4-zh}
现在我们将看到，在格劳尔特判据 \egaref{II.8.9.1-zh}{8.9.1} 中，常常可以
断言预概形 $C$ 和“收缩”$q$ 以本质唯一的方式确定。
\end{env}

\begin{lemma}[8.11.5]
\label{II.8.11.5-zh}
设 $Y$ 是一个预概形，$p:X\to Y$ 是固有态射，$\mathcal{L}$ 是一个
$p$-丰沛的可逆 $\mathcal{O}_X$-模，$C$ 是一个 $Y$-预概形，
$\varepsilon:Y\to C$ 是一个 $Y$-截面，
$q:V=\mathbf{V}(\mathcal{L})\to C$ 是一个 $Y$-态射，并使图表
\egaref{II.8.9.1.1-zh}{8.9.1.1} 交换。进一步假设，若
$p=(p_0,\theta)$，则
$\theta:\mathcal{O}_Y\to p_*(\mathcal{O}_X)$ 是同构。再设
$\mathcal{S}'=\bigoplus_{n\geqslant0}p_*(\mathcal{L}^{\otimes n})$，
$C'=\operatorname{Spec}(\mathcal{S}')$，而
$q':\mathbf{V}(\mathcal{L})\to C'$ 是典范 $Y$-态射
（\egaref{II.8.8.5-zh}{8.8.5}）。则存在唯一的 $Y$-态射 $u:C'\to C$，
使得 $q=u\circ q'$。
\end{lemma}

\begin{proof}
关于 $\theta$ 的假设特别蕴含 $p$ 是满射；由于由
\egaref{II.8.8.4-zh}{8.8.4}，$q'$ 在
$\mathbf{V}(\mathcal{L})-j(X)$ 上的限制是到
$C'-\varepsilon'(Y)$ 上的一个\emph{同构}（$\varepsilon'$ 是 $C'$ 的顶点
截面），仍由 \egaref{II.8.8.4-zh}{8.8.4}，$q'$ 是\emph{固有满射}；此外，
由 \egaref{II.8.8.6-zh}{8.8.6}，若置 $q'=(q'_0,\tau)$，则
$\tau:\mathcal{O}_{C'}\to q'_*(\mathcal{O}_V)$ 是同构。因此可应用引理
\egaref{II.8.11.1-zh}{8.11.1}；若证明 $q$ 在每个纤维
$q'{}^{-1}(z')$ 上为常值（其中 $z'\in C'$），引理即得证。对于
$z'\notin\varepsilon'(Y)$，此条件显然成立。另一方面，若
$z'\in\varepsilon'(Y)$，则存在唯一的 $y\in Y$，使得
$z'=\varepsilon'(y)$；由 \egaref{II.8.8.5.2-zh}{8.8.5.2} 的交换性以及
$q'$ 把 $\mathbf{V}(\mathcal{L})-j(X)$ 映到
$C'-\varepsilon'(Y)$ 中这一事实，有
$q'{}^{-1}(z')=j(p^{-1}(y))$；图表
\egaref{II.8.9.1.1-zh}{8.9.1.1} 的交换性于是建立了断言。
\end{proof}
% Locked French authority: ega2-1-fr.tex, 820505 B,
% SHA-256 84EDBE3E83530AF2959B441796337C9DC21EAFCA6A13114A26778760FBF437AC.
% Sequential source scope: lines 16885-16915, Corollary II.8.11.6 and proof.

\begin{corollary}[8.11.6]
\label{II.8.11.6-zh}
在 \egaref{II.8.11.5-zh}{8.11.5} 的假设下，进一步假设 $q$ 是固有的，且
$q$ 在 $\mathbf{V}(\mathcal{L})-j(X)$ 上的限制是到
$C-\varepsilon(Y)$ 上的一个同构。则态射 $u$ 是普遍闭、满且纯不可分的，
并且它在 $C'-\varepsilon'(Y)$ 上的限制是到 $C-\varepsilon(Y)$ 上的一个
同构。
\end{corollary}

\begin{proof}
由于 $q'$ 是从 $\mathbf{V}(\mathcal{L})-j(X)$ 到
$C'-\varepsilon'(Y)$ 上的一个同构（\egaref{II.8.8.4-zh}{8.8.4}），
最后一项断言立即由 $q=u\circ q'$ 得出。此外，图表

\oldpage[II]{191}
\egaref{II.8.8.5.2-zh}{8.8.5.2} 和
\egaref{II.8.9.1.1-zh}{8.9.1.1} 的交换性表明，$u$ 在 $C'$ 的闭子预概形
$\varepsilon'(Y)$ 上的限制，是到 $C$ 的闭子预概形 $\varepsilon(Y)$ 上的
一个同构；由此立即得出，对每个 $z'\in\varepsilon'(Y)$，若
$z=u(z')$，则 $u$ 定义从 $k(z)$ 到 $k(z')$ 上的一个同构。这些观察证明
$u$ 是双射且纯不可分的；此外，若 $\psi$ 和 $\psi'$ 分别是结构态射
$C\to Y$、$C'\to Y$，则 $\psi'=\psi\circ u$；由于 $\psi'$ 是分离的
（\egaref{II.1.2.4-zh}{1.2.4}），$u$ 也是分离的
（\egaref{I.5.5.1-zh}{I, 5.5.1, (v)}）。另一方面，在引理
\egaref{II.8.11.5-zh}{8.11.5} 的证明中已看到 $q'$ 是满射；由于
$q=u\circ q'$ 是固有的，最终由 \egaref{II.5.4.3-zh}{5.4.3} 和
\egaref{II.5.4.9-zh}{5.4.9} 得出 $u$ 是普遍闭的。
\end{proof}
% Locked French authority: ega2-1-fr.tex, 820505 B,
% SHA-256 84EDBE3E83530AF2959B441796337C9DC21EAFCA6A13114A26778760FBF437AC.
% Sequential source scope: lines 16917-16941, Proposition II.8.11.7 and proof.

\begin{proposition}[8.11.7]
\label{II.8.11.7-zh}
设 $Y$ 是一个预概形，$X$ 是一个整预概形，$p:X\to Y$ 是固有态射，
$\mathcal{L}$ 是一个 $p$-丰沛的可逆 $\mathcal{O}_X$-模，$C$ 是一个正规
$Y$-预概形，$\varepsilon:Y\to C$ 是一个 $Y$-截面，
$q:V=\mathbf{V}(\mathcal{L})\to C$ 是一个 $Y$-态射，并使图表
\egaref{II.8.9.1.1-zh}{8.9.1.1} 交换。进一步假设，若
$p=(p_0,\theta)$，则
$\theta:\mathcal{O}_Y\to p_*(\mathcal{O}_X)$ 是同构。再设
$\mathcal{S}'=\bigoplus_{n\geqslant0}p_*(\mathcal{L}^{\otimes n})$、
$C'=\operatorname{Spec}(\mathcal{S}')$，而
$q':\mathbf{V}(\mathcal{L})\to C'$ 是典范 $Y$-态射
（\egaref{II.8.8.5-zh}{8.8.5}）。则满足 $q=u\circ q'$ 的唯一 $Y$-态射
$u:C'\to C$ 是同构。
\end{proposition}

\begin{proof}
由 \egaref{II.8.8.6-zh}{8.8.6}，$C'$ 是整预概形；由于 $u$ 是底空间
$C'\to C$ 的同胚（由 \egaref{II.8.11.6-zh}{8.11.6}，$u$ 是双射且闭的），
$C$ 是不可约的，因而是整预概形；又因为 $u$ 在 $C'$ 的一个非空开集上的
限制，是到 $C$ 的一个开集上的同构，所以 $u$ 是双有理的。既然假设 $C$
正规，只须应用 \egaref{II.8.11.2-zh}{8.11.2} 即得结论。
\end{proof}
% Locked French authority: ega2-1-fr.tex, 820505 B,
% SHA-256 84EDBE3E83530AF2959B441796337C9DC21EAFCA6A13114A26778760FBF437AC.
% Sequential source scope: lines 16943-16978, Remark II.8.11.8.

\begin{remark}[8.11.8]
\label{II.8.11.8-zh}
\begin{enumerate}
  \item[\rm{(i)}] 注意，$C$ 正规这一假设蕴含 $X$ 也正规。事实上，
  $C'=\operatorname{Spec}(\mathcal{S}')$ 这时是正规的，因为它同构于 $C$；
  并且由 \egaref{II.8.8.6-zh}{8.8.6}，它是整的。由此可知
  $\operatorname{Proj}(\mathcal{S}')$ 是\emph{正规的}。事实上，问题在
  $Y$ 上是局部的；若 $Y$ 是仿射的，$\mathcal{S}'=\widetilde{S'}$，则环
  $S'=\Gamma(C',\mathcal{O}_{C'})$ 是整闭整环
  （\egaref{II.8.8.6.1-zh}{8.8.6.1}），故对每个齐次元
  $f\in S'_+$，分次环 $S'_f$ 是整闭整环
  （[13]，第 I 卷，第 257 和 261 页）；因而其次数 0 项的环
  $S'_{(f)}$ 也是整闭整环，因为 $S'_f$ 与 $S'_{(f)}$ 的分式域的交等于
  $S'_{(f)}$；这就证明了断言（\egaref{II.6.3.4-zh}{6.3.4}）。最后，
  由于 $X$ 同构于 $\operatorname{Proj}(\mathcal{S}')$ 的一个开子预概形
  （\egaref{II.8.8.1-zh}{8.8.1}），$X$ 确为正规的。因此命题
  \egaref{II.8.11.7-zh}{8.11.7} 可表述如下：
  \emph{若 $X$ 是正规整预概形，
  $p=(p_0,\theta):X\to Y$ 是固有态射，且
  $\theta:\mathcal{O}_Y\to p_*(\mathcal{O}_X)$ 是同构，则对每个
  $p$-丰沛的 $\mathcal{O}_X$-模 $\mathcal{L}$，存在唯一一种收缩
  $V=\mathbf{V}(\mathcal{L})$ 的零截面的方式，使所得 $C$ 为正规
  $Y$-概形，而 $q:V\to C$ 为固有 $Y$-态射。}
  \item[\rm{(ii)}] 当 $p$ 是固有的时，假设
  $p_*(\mathcal{O}_X)=\mathcal{O}_Y$ 可视为辅助假设，并不真正限制结果的
  一般性。事实上，若它不成立，只须以 $Y$-概形
  $Y'=\operatorname{Spec}(p_*(\mathcal{O}_X))$ 替换 $Y$，并把 $X$ 视为
  $Y'$-概形。我们将在第三章第 \S~4 节回到这一一般方法。
\end{enumerate}
\end{remark}
% Locked French authority: ega2-1-fr.tex, 820505 B,
% SHA-256 84EDBE3E83530AF2959B441796337C9DC21EAFCA6A13114A26778760FBF437AC.
% Sequential source scope: lines 16979-17050, subsection II.8.12 and Environment II.8.12.1.

\subsection{投影锥上的拟凝聚层}
\label{subsection:II.8.12-zh}

\begin{env}[8.12.1]
\label{II.8.12.1-zh}
重新采用 \egaref{II.8.3.1-zh}{8.3.1} 的假设和记号。设 $\mathcal{M}$ 是
拟凝聚分次 $\mathcal{S}$-模；为避免任何混淆，我们以
$\widetilde{\mathcal{M}}$ 表示相伴的拟凝聚 $\mathcal{O}_C$-模

\oldpage[II]{192}
（它与 $\mathcal{M}$ 相伴，见 \egaref{II.1.4.3-zh}{1.4.3}），这里把
$\mathcal{M}$ 视为\emph{非分次} $\mathcal{S}$-模；并以
$\operatorname{Proj}_0(\mathcal{M})$ 表示拟凝聚 $\mathcal{O}_X$-模，它与
$\mathcal{M}$ 相伴；这次把 $\mathcal{M}$ 视为分次 $\mathcal{S}$-模
（换言之，就是那个 $\mathcal{O}_X$-模 $\widetilde{\mathcal{M}}$；参见
\egaref{II.3.2.2-zh}{3.2.2}）。此外置
\begin{equation}
\label{II.8.12.1.1-zh}
  \mathcal{M}_X=\operatorname{Proj}(\mathcal{M})
  =\bigoplus_{n\in\mathbf{Z}}\operatorname{Proj}_0(\mathcal{M}(n))~;
\tag{8.12.1.1}
\end{equation}
拟凝聚分次 $\mathcal{O}_X$-代数 $\mathcal{S}_X$ 由
\egaref{II.8.6.1.1-zh}{8.6.1.1} 定义；借助典范同态
\egaref{II.3.2.6.1-zh}{3.2.6.1}，
$\operatorname{Proj}(\mathcal{M})$ 具有分次（拟凝聚）
$\mathcal{S}_X$-模结构：
\begin{equation}
\label{II.8.12.1.2-zh}
  \mathcal{O}_X(m)\otimes_{\mathcal{O}_X}
  \operatorname{Proj}_0(\mathcal{M}(n))
  \to \operatorname{Proj}_0(\mathcal{S}(m)\otimes_{\mathcal{S}}\mathcal{M}(n))
  \to \operatorname{Proj}_0(\mathcal{M}(m+n))
\tag{8.12.1.2}
\end{equation}
模公理可借助交换图表 \egaref{II.2.5.11.4-zh}{2.5.11.4} 验证。

若 $Y=\operatorname{Spec}(A)$ 是仿射的，$\mathcal{S}=\widetilde{S}$ 且
$\mathcal{M}=\widetilde{M}$，其中 $S$ 是分次 $A$-代数，$M$ 是分次
$S$-模，则对任意齐次元 $f\in S_+$，有
\begin{equation}
\label{II.8.12.1.3-zh}
  \Gamma(X_f,\operatorname{Proj}(\widetilde{M}))=M_f
\tag{8.12.1.3}
\end{equation}
这是由定义和 \egaref{II.8.2.9.1-zh}{8.2.9.1} 得出的。

现在考虑拟凝聚分次 $\widehat{\mathcal{S}}$-模
\begin{equation}
\label{II.8.12.1.4-zh}
  \widehat{\mathcal{M}}
  =\mathcal{M}\otimes_{\mathcal{S}}\widehat{\mathcal{S}}
\tag{8.12.1.4}
\end{equation}
（$\widehat{\mathcal{S}}$ 由 \egaref{II.8.3.1.1-zh}{8.3.1.1} 定义）；
由此得到拟凝聚分次 $\mathcal{O}_{\widehat{C}}$-模
$\operatorname{Proj}_0(\widehat{\mathcal{M}})$，我们也把它记作
\begin{equation}
\label{II.8.12.1.5-zh}
  \mathcal{M}^{\square}=\operatorname{Proj}_0(\widehat{\mathcal{M}}).
\tag{8.12.1.5}
\end{equation}
显然（\egaref{II.3.2.4-zh}{3.2.4}），$\mathcal{M}^{\square}$ 关于
$\mathcal{M}$ 是一个加性正合函子，并与直和及归纳极限交换。
\end{env}
% Locked French authority: ega2-1-fr.tex, 820505 B,
% SHA-256 84EDBE3E83530AF2959B441796337C9DC21EAFCA6A13114A26778760FBF437AC.
% Sequential source scope: lines 17052-17097, Proposition II.8.12.2, proof, and projective-closure terminology / printed pp.192-193.

\begin{proposition}[8.12.2]
\label{II.8.12.2-zh}
沿用 \egaref{II.8.3.2-zh}{8.3.2} 的记号，存在函子性的典范同构
\begin{equation}
\label{II.8.12.2.1-zh}
  i^*(\mathcal{M}^{\square})\xrightarrow{\sim}\widetilde{\mathcal{M}},
  \qquad
  j^*(\mathcal{M}^{\square})\xrightarrow{\sim}
  \operatorname{Proj}_0(\mathcal{M}).
\tag{8.12.2.1}
\end{equation}
\end{proposition}

\begin{proof}
事实上，$i^*(\mathcal{M}^{\square})$ 典范等同于
$\left(\widehat{\mathcal{M}}/(z-1)\widehat{\mathcal{M}}\right)^{\sim}$，
它位于
$\operatorname{Spec}(\widehat{\mathcal{S}}/(z-1)\widehat{\mathcal{S}})$ 上；
这是由 \egaref{II.3.2.3-zh}{3.2.3} 得出的。于是，
\egaref{II.8.12.2.1-zh}{8.12.2.1} 中第一个典范同构立即由
\egaref{II.1.4.1-zh}{1.4.1} 以及典范同构
$\widehat{\mathcal{M}}/(z-1)\widehat{\mathcal{M}}
\xrightarrow{\sim}\mathcal{M}$ 得出。另一方面，典范浸入
$j:X\to\widehat{C}$ 对应于典范同态
$\widehat{\mathcal{S}}\to\mathcal{S}$，其核为
$z\widehat{\mathcal{S}}$（\egaref{II.8.3.2-zh}{8.3.2}）；
\egaref{II.8.12.2.1-zh}{8.12.2.1} 中第二个同态是典范同态
\egaref{II.3.5.2-zh}{3.5.2, (ii)} 的特殊情形，因为在这里有
$\widehat{\mathcal{M}}\otimes_{\widehat{\mathcal{S}}}\mathcal{S}
=\mathcal{M}$。为验证它是同构，可只考虑
$Y=\operatorname{Spec}(A)$ 为仿射、
$\mathcal{S}=\widetilde{S}$、$\mathcal{M}=\widetilde{M}$ 的情形；参照
\egaref{II.2.8.8-zh}{2.8.8}，对于每个齐次元 $f$，若它属于 $S_+$，前述同态在
$X_f$ 上的限制为同构这一点是立即的。
\end{proof}

\oldpage[II]{193}
由于 \egaref{II.8.12.2.1-zh}{8.12.2.1} 中第一个同构的存在，略微滥用术语，
仍称 $\mathcal{M}^{\square}$ 是 $\mathcal{O}_C$-模
$\widetilde{\mathcal{M}}$ 的\emph{射影闭包}（这里默认为：
$\mathcal{O}_C$-模 $\widetilde{\mathcal{M}}$ 的数据中包括
$\mathcal{S}$-模 $\mathcal{M}$ 的分次）。
% Locked French authority: ega2-1-fr.tex, 820505 B,
% SHA-256 84EDBE3E83530AF2959B441796337C9DC21EAFCA6A13114A26778760FBF437AC.
% Sequential source scope: lines 17099-17124, Environment II.8.12.3.

\begin{env}[8.12.3]
\label{II.8.12.3-zh}
沿用 \egaref{II.8.3.5-zh}{8.3.5} 的记号，有一个函子性的典范同态
\begin{equation}
\label{II.8.12.3.1-zh}
  p^*(\operatorname{Proj}_0(\mathcal{M}))
  \to\mathcal{M}^{\square}|\widehat{E}.
\tag{8.12.3.1}
\end{equation}
事实上，这是 \egaref{II.3.5.6-zh}{3.5.6} 中一般定义的同态
$u^{\sharp}$ 的一个特殊情形。若 $Y=\operatorname{Spec}(A)$ 为仿射的、
$\mathcal{S}=\widetilde{S}$、$\mathcal{M}=\widetilde{M}$，则参照
\egaref{II.2.8.8-zh}{2.8.8} 可见，
\egaref{II.8.12.3.1-zh}{8.12.3.1} 在
$p^{-1}(X_f)=\widehat{C}_f$ 上的限制（其中 $f$ 是 $S_+$ 中的齐次元），
对应于典范同态
\begin{equation}
\label{II.8.12.3.2-zh}
  M_{(f)}\otimes_{S_{(f)}}S_f^{\leq}\to M_f^{\leq}
\tag{8.12.3.2}
\end{equation}
这里用到了 \egaref{II.8.2.3.2-zh}{8.2.3.2} 和
\egaref{II.8.2.5.2-zh}{8.2.5.2}。
\end{env}
% Locked French authority: ega2-1-fr.tex, 820505 B,
% SHA-256 84EDBE3E83530AF2959B441796337C9DC21EAFCA6A13114A26778760FBF437AC.
% Sequential source scope: lines 17126-17160, Environment II.8.12.4.

\begin{env}[8.12.4]
\label{II.8.12.4-zh}
置身于 \egaref{II.8.5.1-zh}{8.5.1} 的假设下，并沿用其中记号。由
\egaref{II.1.5.6-zh}{1.5.6}，对每个拟凝聚分次 $\mathcal{S}$-模
$\mathcal{M}$，一方面有 $\mathcal{O}_{C'}$-模的典范同构
\begin{equation}
\label{II.8.12.4.1-zh}
  \Phi^*(\widetilde{\mathcal{M}})
  \xrightarrow{\sim}
  \left(q^*(\mathcal{M})\otimes_{q^*(\mathcal{S})}\mathcal{S}'\right)^{\sim}
\tag{8.12.4.1}
\end{equation}
；另一方面，\egaref{II.3.5.6-zh}{3.5.6} 蕴含存在一个典范的
$\operatorname{Proj}(\varphi)$-态射
\begin{equation}
\label{II.8.12.4.2-zh}
  \operatorname{Proj}_0(\mathcal{M})\to
  \left(\operatorname{Proj}_0
  (q^*(\mathcal{M})\otimes_{q^*(\mathcal{S})}\mathcal{S}')\right)
  |G(\varphi)
\tag{8.12.4.2}
\end{equation}
以及一个典范的 $\widehat{\Phi}$-态射
\begin{equation}
\label{II.8.12.4.3-zh}
  \operatorname{Proj}_0(\widehat{\mathcal{M}})\to
  \left(\operatorname{Proj}_0
  (q^*(\widehat{\mathcal{M}})
  \otimes_{q^*(\widehat{\mathcal{S}})}\widehat{\mathcal{S}'})\right)
  |G(\widehat{\varphi}).
\tag{8.12.4.3}
\end{equation}
\end{env}
% Locked French authority: ega2-1-fr.tex, 820505 B,
% SHA-256 84EDBE3E83530AF2959B441796337C9DC21EAFCA6A13114A26778760FBF437AC.
% Sequential source scope: lines 17162-17236, Environment II.8.12.5 /
% printed pp.193-194.

\begin{env}[8.12.5]
\label{II.8.12.5-zh}
现在考虑 \egaref{II.8.6.1-zh}{8.6.1} 的情形，并沿用相同的记号；因而在
前文中取 $Y'=X$，其中态射 $q:X\to Y$ 是结构态射，而 $\varphi$ 是典范
$q$-态射（\egaref{II.8.6.1.2-zh}{8.6.1.2}）。于是有典范同构
\begin{equation}
\label{II.8.12.5.1-zh}
  q^*(\mathcal{M})\otimes_{q^*(\mathcal{S})}\mathcal{S}_X^{\geq}
  \xrightarrow{\sim}\mathcal{M}_X^{\geq}
\tag{8.12.5.1}
\end{equation}
其中置
$\mathcal{M}_X^{\geq}=\bigoplus_{n\geq0}\operatorname{Proj}_0(\mathcal{M}(n))$。
事实上，可只考虑 $Y=\operatorname{Spec}(A)$ 为仿射、
$\mathcal{S}=\widetilde{S}$ 且 $\mathcal{M}=\widetilde{M}$ 的情形，并在每个
仿射开集 $X_f$ 中定义同构 \egaref{II.8.12.5.1-zh}{8.12.5.1}（其中
$f$ 是 $S_+$ 中的齐次元），再验证它与改用 $f$ 的齐次倍元相容。
而 \egaref{II.8.12.5.1-zh}{8.12.5.1} 左端在 $X_f$ 上的限制，由
\egaref{II.8.6.2.1-zh}{8.6.2.1} 是
$\widetilde{M}'=
\left((M\otimes_A S_{(f)})
\otimes_{S\otimes_A S_{(f)}}S_f^{\geq}\right)^{\sim}$；由于有
$M\otimes_A S_{(f)}$ 到
$M\otimes_S(S\otimes_A S_{(f)})$ 上的典范同构，便得出
$\widetilde{M}'$ 到
$\left(M\otimes_S S_f^{\geq}\right)^{\sim}$ 上的同构；后者由
\egaref{II.8.2.9.1-zh}{8.2.9.1} 典范同构于
\egaref{II.8.12.5.1-zh}{8.12.5.1} 右端在 $X_f$ 上的限制，并满足所需的
相容条件。

\oldpage[II]{194}
在上述论证中，以 $\widehat{\mathcal{M}}$ 代替 $\mathcal{M}$，以
$\widehat{\mathcal{S}}$ 代替 $\mathcal{S}$，并以
$(\mathcal{S}_X^{\geq})^{\wedge}$ 代替 $\mathcal{S}_X$，同样得到典范同构
\begin{equation}
\label{II.8.12.5.2-zh}
  q^*(\widehat{\mathcal{M}})
  \otimes_{q^*(\widehat{\mathcal{S}})}(\mathcal{S}_X^{\geq})^{\wedge}
  \xrightarrow{\sim}(\mathcal{M}_X^{\geq})^{\wedge}.
\tag{8.12.5.2}
\end{equation}
回想 \egaref{II.8.6.2-zh}{8.6.2}，结构态射
$u:\operatorname{Proj}(\mathcal{S}_X^{\geq})\to X$ 是同构；首先由前文
得出典范的 $u$-同构
\begin{equation}
\label{II.8.12.5.3-zh}
  \operatorname{Proj}_0\mathcal{M}
  \xrightarrow{\sim}\operatorname{Proj}_0(\mathcal{M}_X^{\geq})
\tag{8.12.5.3}
\end{equation}
它是 \egaref{II.8.12.4.2-zh}{8.12.4.2} 的一个特殊情形。事实上，沿用
\egaref{II.8.6.2-zh}{8.6.2} 的证明中的记号，这化为验证典范同态
$M_{(f)}\otimes_{S_{(f)}}(S_f^{\geq})^{(d)}\to(M_f^{\geq})^{(d)}$ 在
$f\in S_d$ 时是同构，而这是立即的。

其次，同构 \egaref{II.8.12.5.2-zh}{8.12.5.2} 使我们可以把
\egaref{II.8.12.4.3-zh}{8.12.4.3} 应用于典范态射
$r=\operatorname{Proj}(\widehat{\alpha}):\widehat{C}_X\to
\widehat{C}$，从而得到典范的 $r$-态射
\begin{equation}
\label{II.8.12.5.4-zh}
  \mathcal{M}^{\square}\to(\mathcal{M}_X^{\geq})^{\square}.
\tag{8.12.5.4}
\end{equation}
现在回想 \egaref{II.8.6.2-zh}{8.6.2}，$r$ 在去顶点锥
$\widehat{E}_X$ 和 $E_X$ 上的限制，分别是到 $\widehat{E}$ 和 $E$ 上的
\emph{同构}。此外：
\end{env}
% Locked French authority: ega2-1-fr.tex, 820505 B,
% SHA-256 84EDBE3E83530AF2959B441796337C9DC21EAFCA6A13114A26778760FBF437AC.
% Sequential source scope: lines 17238-17278, Proposition II.8.12.6 and proof.

\begin{proposition}[8.12.6]
\label{II.8.12.6-zh}
典范 $r$-态射 \egaref{II.8.12.5.4-zh}{8.12.5.4} 在
$\widehat{E}_X$ 和 $E_X$ 上的限制分别是同构
\begin{equation}
\label{II.8.12.6.1-zh}
  \mathcal{M}^{\square}|\widehat{E}
  \xrightarrow{\sim}(\mathcal{M}_X^{\geq})^{\square}|\widehat{E}_X
\tag{8.12.6.1}
\end{equation}
\begin{equation}
\label{II.8.12.6.2-zh}
  \widetilde{\mathcal{M}}|E
  \xrightarrow{\sim}(\mathcal{M}_X^{\geq})^{\sim}|E_X.
\tag{8.12.6.2}
\end{equation}
\end{proposition}

\begin{proof}
如 \egaref{II.8.6.2-zh}{8.6.2} 的证明中那样，化归到 $Y$ 为仿射的情形；
沿用该证明的记号，并化归到定义 \egaref{II.2.8.8-zh}{2.8.8}，须证明
典范同态
\[
  \widehat{M}_{(f)}\otimes_{\widehat{S}_{(f)}}
  (S_f^{\geq})_{(f/1)}^{\wedge}
  \to(M\otimes_S S_f^{\geq})_{(f/1)}^{\wedge}
\]
是同构；但由 \egaref{II.8.2.3.2-zh}{8.2.3.2} 和
\egaref{II.8.2.5.2-zh}{8.2.5.2}，第一项典范地等同于
$M_f^{\leq}\otimes_{S_f^{\leq}}(S_f^{\geq})_{f/1}^{\leq}$，继而由
\egaref{II.8.2.7.2-zh}{8.2.7.2} 等同于 $M_f^{\leq}$；第二项等同于
$(M_f^{\geq})_{f/1}^{\leq}$，继而由
\egaref{II.8.2.9.2-zh}{8.2.9.2} 也等同于 $M_f^{\leq}$。这就证明了
\egaref{II.8.12.6.1-zh}{8.12.6.1}；而公式
\egaref{II.8.12.6.2-zh}{8.12.6.2} 由
\egaref{II.8.12.6.1-zh}{8.12.6.1} 和
\egaref{II.8.12.2.1-zh}{8.12.2.1} 得出。
\end{proof}
% Locked French authority: ega2-1-fr.tex, 820505 B,
% SHA-256 84EDBE3E83530AF2959B441796337C9DC21EAFCA6A13114A26778760FBF437AC.
% Sequential source scope: lines 17280-17317, Corollary II.8.12.7 and Proposition II.8.12.8 with proofs / printed pp.194-195.

\begin{corollary}[8.12.7]
\label{II.8.12.7-zh}
在 \egaref{II.8.6.3-zh}{8.6.3} 的等同下，
$(\mathcal{M}_X^{\geq})^{\square}$ 在 $\widehat{E}_X$ 上的限制等同于
$(\mathcal{M}_X^{\leq})^{\sim}$，而
$(\mathcal{M}_X^{\geq})^{\square}$ 在 $E_X$ 上的限制等同于
$\widetilde{\mathcal{M}}_X$。
\end{corollary}

\begin{proof}
事实上，可化归到仿射情形；此时断言由把
$(M_f^{\geq})_{f/1}^{\leq}$ 等同于 $M_f^{\leq}$，并把
$(M_f^{\geq})_{f/1}$ 等同于 $M_f$ 得出
（\egaref{II.8.2.9.2-zh}{8.2.9.2}）。
\end{proof}

\begin{proposition}[8.12.8]
\label{II.8.12.8-zh}
在 \egaref{II.8.6.4-zh}{8.6.4} 的假设下，典范同态
\egaref{II.8.12.3.1-zh}{8.12.3.1} 是同构。
\end{proposition}

\begin{proof}
考虑到 $\operatorname{Proj}(\mathcal{S}_X^{\geq})\to X$ 是同构
（\egaref{II.8.6.2-zh}{8.6.2}），以及下列

\oldpage[II]{195}
同构 \egaref{II.8.12.5.4-zh}{8.12.5.4} 和
\egaref{II.8.12.6.1-zh}{8.12.6.1}，问题化归为证明典范同态
$p_X^*(\operatorname{Proj}_0(\mathcal{M}_X^{\geq}))
\to(\mathcal{M}_X^{\geq})^{\square}|\widehat{E}_X$ 的相应命题；换言之，
化归到 $\mathcal{S}_1$ 是可逆 $\mathcal{O}_Y$-模且
$\mathcal{S}$ 由 $\mathcal{S}_1$ 生成的情形。沿用
\egaref{II.8.12.3-zh}{8.12.3} 的记号，此时对于 $f\in S_1$，有
$S_f^{\leq}=S_{(f)}[1/f]$，并且典范同态
$M_{(f)}\otimes_{S_{(f)}}S_f^{\leq}\to M_f^{\leq}$ 由
$M_f^{\leq}$ 的定义可知是同构。
\end{proof}
% Locked French authority: ega2-1-fr.tex, 820505 B,
% SHA-256 84EDBE3E83530AF2959B441796337C9DC21EAFCA6A13114A26778760FBF437AC.
% Sequential source scope: lines 17319-17354, Environment II.8.12.9 and Proposition II.8.12.10 with proof.

\begin{env}[8.12.9]
\label{II.8.12.9-zh}
现在考虑拟凝聚 $\mathcal{S}$-模
\[
  \mathcal{M}_{[n]}=\bigoplus_{m\geqslant n}\mathcal{M}_m
\]
以及（沿用 \egaref{II.8.7.2-zh}{8.7.2} 的记号）拟凝聚分次
$\mathcal{S}^{\natural}$-模
\begin{equation}
\label{II.8.12.9.1-zh}
  \mathcal{M}^{\natural}
  =\left(\bigoplus_{n\geqslant0}\mathcal{M}_{[n]}\right)^{\sim}.
\tag{8.12.9.1}
\end{equation}
我们已在 \egaref{II.8.7.3-zh}{8.7.3} 中看到，存在典范的 $C$-同构
$h:C_X\xrightarrow{\sim}\operatorname{Proj}(\mathcal{S}^{\natural})$。
此外：
\end{env}

\begin{proposition}[8.12.10]
\label{II.8.12.10-zh}
存在典范的 $h$-同构
\begin{equation}
\label{II.8.12.10.1-zh}
  \operatorname{Proj}_0(\mathcal{M}^{\natural})
  \xrightarrow{\sim}\widetilde{\mathcal{M}}_X.
\tag{8.12.10.1}
\end{equation}
\end{proposition}

\begin{proof}
仿照 \egaref{II.8.7.3-zh}{8.7.3} 论证，但这次用双同构
\egaref{II.8.2.9.3-zh}{8.2.9.3} 的存在性替代
\egaref{II.8.2.7.3-zh}{8.2.7.3}。细节留给读者。
\end{proof}
% Locked French authority: ega2-1-fr.tex, 820505 B,
% SHA-256 84EDBE3E83530AF2959B441796337C9DC21EAFCA6A13114A26778760FBF437AC.
% Sequential source scope: lines 17356-17411, subsection II.8.13 and Environments II.8.13.1-II.8.13.2.

\subsection{子层与闭子概形的射影闭包}
\label{subsection:II.8.13-zh}

\begin{env}[8.13.1]
\label{II.8.13.1-zh}
假设和记号同 \egaref{II.8.12.1-zh}{8.12.1}。考虑一个拟凝聚子-
$\mathcal{S}$-模 $\mathcal{N}$，它包含于 $\mathcal{M}$，且
\emph{不必是分次的}。因此可考虑与之相伴的拟凝聚
$\mathcal{O}_C$-模 $\widetilde{\mathcal{N}}$；它与 $\mathcal{N}$ 相伴，
并且是 $\mathcal{O}_C$-模 $\widetilde{\mathcal{M}}$ 的子模。另一方面，
我们已经看到（\egaref{II.8.12.2.1-zh}{8.12.2.1}），
$\widetilde{\mathcal{M}}$ 可等同于 $\mathcal{M}^{\square}$ 在 $C$ 上的限制。
由于典范单射 $i:C\to\widehat{C}$ 是仿射态射
（\egaref{II.8.3.2-zh}{8.3.2}），因而\emph{更是}拟紧的，所以
\emph{典范延拓} $(\widetilde{\mathcal{N}})^{-}$，即最大的
子-$\mathcal{O}_{\widehat{C}}$-模，它包含于 $\mathcal{M}^{\square}$ 中，
并诱导 $\widetilde{\mathcal{N}}$ 于 $C$ 上，是拟凝聚
$\mathcal{O}_{\widehat{C}}$-模（\egaxref{I.9.4.2-zh}{I, 9.4.2}）。我们将借助
一个分次 $\widehat{\mathcal{S}}$-模明确描述它。
\end{env}

\begin{env}[8.13.2]
\label{II.8.13.2-zh}
为此，对每个整数 $n\geqslant0$，考虑同态
$\bigoplus_{i\leqslant n}\mathcal{M}_i\to\mathcal{M}$；对每个开集 $U$
（它是 $Y$ 的开集），该同态把族
\[
  (s_i)\in\bigoplus_{i\leqslant n}\Gamma(U,\mathcal{M}_i)
\]
对应到截面 $\sum_i s_i\in\Gamma(U,\mathcal{M})$。以
$\mathcal{N}'_n$ 表示 $\mathcal{N}$ 在此同态下的原像；它是
$\mathcal{O}_Y$-模 $\bigoplus_{i\leqslant n}\mathcal{M}_i$ 的拟凝聚子模。
现在考虑同态
$\bigoplus_{i\leqslant n}\mathcal{M}_i\to
\widehat{\mathcal{M}}=\mathcal{M}[z]$，它把 $(s_i)$ 对应到截面
$\sum_{i\leqslant n}s_i z^{n-i}\in\Gamma(U,\widehat{\mathcal{M}}_n)$；令
$\mathcal{N}_n$ 为 $\mathcal{N}'_n$ 在此同态下的像。立即可验证，
$\overline{\mathcal{N}}=\bigoplus_{n\geqslant0}\mathcal{N}_n$ 是
$\widehat{\mathcal{S}}$-模 $\widehat{\mathcal{M}}$ 的分次拟凝聚子模；称
$\overline{\mathcal{N}}$ 是由 $\mathcal{N}$ 通过\emph{齐次化}、借助
“齐次化变量” $z$ 得到的。应注意，

\oldpage[II]{196}
若 $\mathcal{N}$ 已经是一个子-$\mathcal{S}$-模，而且是\emph{分次}的，
并包含于 $\mathcal{M}$，则 $\overline{\mathcal{N}}$ 可等同于
$\widehat{\mathcal{N}}_n$ 中次数为 $n\geqslant0$ 的各分量之直和，其中
$\widehat{\mathcal{N}}=\mathcal{N}[z]$。
\end{env}
% Locked French authority: ega2-1-fr.tex, 820505 B,
% SHA-256 84EDBE3E83530AF2959B441796337C9DC21EAFCA6A13114A26778760FBF437AC.
% Sequential source scope: lines 17413-17466, Proposition II.8.13.3 and proof.

\begin{proposition}[8.13.3]
\label{II.8.13.3-zh}
$\mathcal{O}_{\widehat{C}}$-模
$\operatorname{Proj}_0(\overline{\mathcal{N}})$ 是
$\widetilde{\mathcal{N}}$ 到 $\widehat{C}$ 上的典范延拓
$(\widetilde{\mathcal{N}})^{-}$。
\end{proposition}

\begin{proof}
由典范延拓的定义（\egaref{I.9.4.1-zh}{I, 9.4.1}），问题在 $Y$ 和
$\widehat{C}$ 上是局部的。因此可先假设 $Y=\operatorname{Spec}(A)$ 为
仿射的，且 $\mathcal{S}=\widetilde{S}$、$\mathcal{M}=\widetilde{M}$、
$\mathcal{N}=\widetilde{N}$，其中 $N$ 是 $M$ 的一个不必分次的 $S$-子模。
此外，由 \egaref{II.8.3.2.6-zh}{8.3.2.6}，$\widehat{C}$ 是诸仿射开集
$\widehat{C}_z=C$ 和
$\widehat{C}_f=\operatorname{Spec}(S_f^{\leq})$（其中 $f$ 是 $S_+$ 中的
齐次元）的并。因此只须证明：$1^\circ$
$\operatorname{Proj}_0(\overline{\mathcal{N}})$ 在 $C$ 上的限制是
$\widetilde{\mathcal{N}}$；$2^\circ$
$\operatorname{Proj}_0(\overline{\mathcal{N}})$ 在每个 $\widehat{C}_f$
上的限制，是 $\mathcal{N}$ 在
$C\cap\widehat{C}_f=\operatorname{Spec}(S_f)$ 上的限制的典范延拓
（\egaref{II.8.3.2.6-zh}{8.3.2.6}）。对于第一点，注意
$\operatorname{Proj}_0(\overline{\mathcal{N}})|C$ 等同于
$(\overline{N}_{(z)})^{\sim}$
（\egaref{II.8.3.2.4-zh}{8.3.2.4}）；而
$\overline{N}_{(z)}$ 典范等同于
（\egaref{II.2.2.5-zh}{2.2.5}）$\overline{N}$ 在
$\widehat{M}/(z-1)\widehat{M}$ 中的像；借助后者到 $M$ 上的典范同构
（\egaref{II.8.2.5-zh}{8.2.5}），这个像由
\egaref{II.8.13.2-zh}{8.13.2} 中 $\overline{N}$ 的定义等同于 $N$。

为证明第二点，注意单射
$i:C\cap\widehat{C}_f\to\widehat{C}_f$ 对应于典范单射
$S_f^{\leq}\to S_f$（\egaref{II.8.3.2.6-zh}{8.3.2.6}）；另一方面，
有 $\Gamma(\widehat{C}_f,\mathcal{M}^{\square})=M_f^{\leq}$、
$\Gamma(\widehat{C}_f,i_*(\widetilde{\mathcal{N}}))=N_f$，以及
（由 \egaref{II.8.12.2.1-zh}{8.12.2.1}）
$\Gamma(\widehat{C}_f,i_*(i^*(\mathcal{M}^{\square})))=M_f$。结合
\egaref{I.9.4.2-zh}{I, 9.4.2}，问题因而化为证明：
$\overline{N}_{(f)}\subset\widehat{M}_{(f)}=M_f^{\leq}$ 典范等同于 $N_f$
在典范单射 $M_f^{\leq}\to M_f$ 下的逆像。事实上，置
$d=\deg(f)>0$，并假设 $M_f$ 的一个元素
$(\sum_{k\leqslant md}x_k)/f^m$（其中 $x_k\in M_k$）具有形式
$y/f^m$，其中 $y\in N$。把 $y$ 与诸 $x_k$ 乘以同一个适当的 $f^h$，
可先假设 $\sum_{k\leqslant md}x_k=y$。但在等同
\egaref{II.8.2.5.2-zh}{8.2.5.2} 中，
$(\sum_{k\leqslant md}x_k)/f^m$ 对应于
$\sum_{k\leqslant md}x_kz^{md-k}/f^m$；这个元素确实属于
$\overline{N}_{(f)}$，因为 $\sum_{k\leqslant md}x_k\in N$；反之显然。
\end{proof}
% Locked French authority: ega2-1-fr.tex, 820505 B,
% SHA-256 84EDBE3E83530AF2959B441796337C9DC21EAFCA6A13114A26778760FBF437AC.
% Sequential source scope: lines 17468-17509, Remark II.8.13.4 closing subsection II.8.13 / printed pp.196-197.

\begin{remark}[8.13.4]
\label{II.8.13.4-zh}
\begin{enumerate}
  \item[\rm{(i)}] \egaref{II.8.13.3-zh}{8.13.3} 最重要的应用情形是
  $\mathcal{M}=\mathcal{S}$；于是 $\widetilde{\mathcal{N}}$ 是一个
  \emph{任意的}拟凝聚理想层 $\mathcal{J}$；它是
  $\mathcal{O}_C$ 的理想层（\egaref{II.1.4.3-zh}{1.4.3}），并与一个
  \emph{闭子预概形} $Z$ 一一对应（该子预概形位于 $C$ 中）。此时，
  典范延拓 $\overline{\mathcal{J}}$（源自 $\mathcal{J}$）是
  $\mathcal{O}_{\widehat{C}}$ 的拟凝聚理想层，它定义了\emph{闭包}
  $\overline{Z}$，也就是 $Z$ 在 $\widehat{C}$ 中的闭包
  （\egaref{I.9.5.10-zh}{I, 9.5.10}）；命题
  \egaref{II.8.13.3-zh}{8.13.3} 给出了一种定义 $\overline{Z}$ 的典范方法：
  使用 $\widehat{\mathcal{S}}=\mathcal{S}[z]$ 中的分次理想。
  \item[\rm{(ii)}] 为简单起见，假设 $Y$ 是仿射的，并沿用
  \egaref{II.8.13.3-zh}{8.13.3} 证明中的记号。对每个非零的 $x\in N$，
  令 $d(x)$ 为各齐次分量 $x_i$ 的次数的最大值；这里所说的是 $x$ 在
  $M$ 中的齐次分量。按定义，
  $\overline{N}$ 是 $\widehat{M}$ 的如下子模：它由 $0$ 以及形如
  $h(x,k)=z^k\sum_{i\leqslant d(x)}x_i z^{d(x)-i}$ 的元素组成（$k$ 为整数
  $\geqslant0$）；因而，作为 $\widehat{S}=S[z]$-模，它由形如
  \[
    h(x,0)=\sum_{i\leqslant d(x)}x_i z^{d(x)-i}.
  \]
  的元素生成。

  \oldpage[II]{197}
  称 $h(x,0)$ 是由 $x$ 借助“齐次化变量”$z$ 进行\emph{齐次化}而得到的。
  但是，由于 $h(x,0)$ 对 $x$ 不具有可加性（因而\emph{更不用说}关于
  $S$ 为线性），\emph{切勿误以为}（即使 $M=S$），各 $h(x,0)$ 会构成
  一组\emph{生成元}来生成分次 $\widehat{S}$-模 $\overline{N}$，只要令
  $x$ 遍历一组\emph{生成元}，而这组生成元属于 $S$-模 $N$。然而，在
  $N$ 是 $S$-模且\emph{秩一自由}的情形下（初等代数几何只考虑这一情形），
  这确实成立：若 $t$ 是 $N$ 的一个基，则 $h(t,0)$ 生成
  $\widehat{S}$-模 $\overline{N}$。
\end{enumerate}
\end{remark}
% Locked French authority: ega2-1-fr.tex, 820505 B,
% SHA-256 84EDBE3E83530AF2959B441796337C9DC21EAFCA6A13114A26778760FBF437AC.
% Sequential source scope: lines 17511-17538, subsection II.8.14 and Environment II.8.14.1.

\subsection{分次 \texorpdfstring{$\mathcal{S}$}{S}-模相伴层的补充}
\label{subsection:II.8.14-zh}

\begin{env}[8.14.1]
\label{II.8.14.1-zh}
设 $Y$ 是一个预概形，$\mathcal{S}$ 是一个\emph{正次数}拟凝聚分次
$\mathcal{O}_Y$-代数，$X=\operatorname{Proj}(\mathcal{S})$，而
$q:X\to Y$ 是结构态射（由 \egaref{II.3.1.3-zh}{3.1.3}，它是分离的）。
沿用 \egaref{II.8.12.1-zh}{8.12.1} 的记号；于是我们定义了关于
$\mathcal{M}$ 的函子 $\mathcal{M}_X=\operatorname{Proj}(\mathcal{M})$，
它从拟凝聚分次 $\mathcal{S}$-模范畴到拟凝聚分次 $\mathcal{S}_X$-模范畴；
此外，由 \egaref{II.3.2.4-zh}{3.2.4}，它显然是与归纳极限交换的
\emph{正合加性}函子。

还要注意，由定义 \egaref{II.8.12.1.1-zh}{8.12.1.1} 立即得出
\begin{equation}
\label{II.8.14.1.1-zh}
  \operatorname{Proj}(\mathcal{M}(n))
  =(\operatorname{Proj}(\mathcal{M}))(n)
  \qquad\text{对每个 }n\in\mathbf{Z}.
\tag{8.14.1.1}
\end{equation}
\end{env}
% Locked French authority: ega2-1-fr.tex, 820505 B,
% SHA-256 84EDBE3E83530AF2959B441796337C9DC21EAFCA6A13114A26778760FBF437AC.
% Sequential source scope: lines 17540-17608, Environment II.8.14.2.

\begin{env}[8.14.2]
\label{II.8.14.2-zh}
我们首先把 \egaref{II.3.2.6-zh}{3.2.6} 中定义的典范同态 $\lambda$ 和
$\mu$ 延拓到形如 $\operatorname{Proj}(\mathcal{M})$ 的
$\mathcal{S}_X$-模。为此，注意任取 $m\in\mathbf{Z}$ 和
$n\in\mathbf{Z}$，由 \egaref{II.2.1.2.1-zh}{2.1.2.1}，都有
$\mathcal{O}_X$-模的典范同态
\begin{equation}
\label{II.8.14.2.1-zh}
  \lambda_{mn}:\operatorname{Proj}_0(\mathcal{M}(m))
  \otimes_{\mathcal{O}_X}\operatorname{Proj}_0(\mathcal{N}(n))
  \to\operatorname{Proj}_0
  ((\mathcal{M}\otimes_{\mathcal{S}}\mathcal{N})(m+n))
\tag{8.14.2.1}
\end{equation}
同样，由 \egaref{II.2.1.2.2-zh}{2.1.2.2}，有
$\mathcal{O}_X$-模的典范同态
\begin{equation}
\label{II.8.14.2.2-zh}
  \mu_{mn}:\operatorname{Proj}_0
  ((\mathcal{H}om_{\mathcal{S}}(\mathcal{M},\mathcal{N}))(n-m))
  \to\mathcal{H}om_{\mathcal{O}_X}
  (\operatorname{Proj}_0(\mathcal{M}(m)),
   \operatorname{Proj}_0(\mathcal{N}(n)))
\tag{8.14.2.2}
\end{equation}
这里 $\mathcal{M}$、$\mathcal{N}$ 是任意两个拟凝聚分次
$\mathcal{S}$-模。由此得出同态
\[
  \mu_k:\operatorname{Proj}_0
  ((\mathcal{H}om_{\mathcal{S}}(\mathcal{M},\mathcal{N}))(k))
  \to
  (\mathcal{H}om_{\mathcal{S}_X}
  (\operatorname{Proj}(\mathcal{M}),\operatorname{Proj}(\mathcal{N})))_k
\]
其定义如下：对每个
$u\in\Gamma(U,\operatorname{Proj}_0
((\mathcal{H}om_{\mathcal{S}}(\mathcal{M},\mathcal{N}))(k)))$，令其对应于
分次 $\mathbf{Z}$-模的 $k$ 次同态 $\mu_k(u)$
$\Gamma(U,\operatorname{Proj}(\mathcal{M}))\to
\Gamma(U,\operatorname{Proj}(\mathcal{N}))$（$U$ 是 $X$ 的开集）；它在每个
$\Gamma(U,\operatorname{Proj}_0(\mathcal{M}(m)))$ 上都与
$\mu_{m,m+k}(u)$ 一致。此外，回到 $\mu_{mn}$ 的定义
（\egaref{II.2.5.12.1-zh}{2.5.12.1}），立即可见 $\mu_k(u)$ 实际上是分次
$\Gamma(U,\mathcal{S}_X)$-模的 $k$ 次同态；并且诸 $\mu_k$ 定义了分次
$\mathcal{S}_X$-模的同态
\begin{equation}
\label{II.8.14.2.3-zh}
  \operatorname{Proj}
  (\mathcal{H}om_{\mathcal{S}}(\mathcal{M},\mathcal{N}))
  \to\mathcal{H}om_{\mathcal{S}_X}
  (\operatorname{Proj}(\mathcal{M}),\operatorname{Proj}(\mathcal{N})).
\tag{8.14.2.3}
\end{equation}

同样，结合结合性图表 \egaref{II.2.5.11.4-zh}{2.5.11.4}，同态
\egaref{II.8.14.2.1-zh}{8.14.2.1} 给出分次 $\mathcal{S}_X$-模的同态
\begin{equation}
\label{II.8.14.2.4-zh}
  \lambda:\operatorname{Proj}(\mathcal{M})
  \otimes_{\mathcal{S}_X}\operatorname{Proj}(\mathcal{N})
  \to\operatorname{Proj}(\mathcal{M}\otimes_{\mathcal{S}}\mathcal{N}).
\tag{8.14.2.4}
\end{equation}
\end{env}
% Locked French authority: ega2-1-fr.tex, 820505 B,
% SHA-256 84EDBE3E83530AF2959B441796337C9DC21EAFCA6A13114A26778760FBF437AC.
% Sequential source scope: lines 17610-17658, Proposition II.8.14.3 and proof.

\oldpage[II]{198}
\begin{proposition}[8.14.3]
\label{II.8.14.3-zh}
同态 \egaref{II.8.14.2.4-zh}{8.14.2.4} 是双射；同态
\egaref{II.8.14.2.3-zh}{8.14.2.3} 也是如此，只要分次
$\mathcal{S}$-模 $\mathcal{M}$ 有限表示
（\egaref{II.3.1.1-zh}{3.1.1}）。
\end{proposition}

\begin{proof}
问题显然在 $X$ 和 $Y$ 上都是局部的；因此可假设
$Y=\operatorname{Spec}(A)$ 是仿射的，$\mathcal{S}=\widetilde{S}$、
$\mathcal{M}=\widetilde{M}$、$\mathcal{N}=\widetilde{N}$，其中 $S$ 是一个
正次数分次 $A$-代数，而 $M$ 和 $N$ 是两个分次 $S$-模。若 $f$ 是
$S_+$ 的齐次元，则同态 \egaref{II.8.14.2.1-zh}{8.14.2.1} 和
\egaref{II.8.14.2.2-zh}{8.14.2.2} 在仿射开集 $D_+(f)$ 上的限制，分别
对应于典范同态 \egaref{II.2.5.11.1-zh}{2.5.11.1} 和
\egaref{II.2.5.12.1-zh}{2.5.12.1}：
\begin{align*}
  M(m)_{(f)}\otimes_{S_{(f)}}N(n)_{(f)}
  &\to(M\otimes_S N)(m+n)_{(f)},\\
  (\operatorname{Hom}_S(M,N))(n-m)_{(f)}
  &\to\operatorname{Hom}_{S_{(f)}}(M(m)_{(f)},N(n)_{(f)}).
\end{align*}
考察这些同态的定义，并计及 \egaref{II.8.2.9.1-zh}{8.2.9.1}，可知
\egaref{II.8.14.2.4-zh}{8.14.2.4} 在 $D_+(f)$ 上的限制对应于典范同态
\[
  M_f\otimes_{S_f}N_f\to(M\otimes_S N)_f
\]
它在 \egaxref{0.1.3.4-zh}{0, 1.3.4} 中定义，并且已知是同构。同样，
\egaref{II.8.14.2.3-zh}{8.14.2.3} 在 $D_+(f)$ 上的限制对应于典范同态
\egaxref{0.1.3.5-zh}{0, 1.3.5}
\[
  (\operatorname{Hom}_S(M,N))_f\to
  \operatorname{Hom}_{S_f}(M_f,N_f)
\]
这里考虑到 $M$ 是有限型的，模 $\operatorname{Hom}_S(M,N)$ 是由
$S$-模的\emph{齐次}同态构成的各子群之直和
（\egaref{II.2.1.2-zh}{2.1.2}），并与\emph{所有}同态 $M\to N$
（作为 $S$-模的同态）所成的集合重合。于是 $M$ 有限表示这一假设蕴含
（\egaxref{0.1.3.5-zh}{0, 1.3.5}），所考虑的典范同态确实是同构。
\end{proof}
% Locked French authority: ega2-1-fr.tex, 820505 B,
% SHA-256 84EDBE3E83530AF2959B441796337C9DC21EAFCA6A13114A26778760FBF437AC.
% Sequential source scope: lines 17660-17675, Proposition II.8.14.4 and proof.

\begin{proposition}[8.14.4]
\label{II.8.14.4-zh}
若 $U$ 是 $X$ 的拟紧开集，则存在整数 $d$，使得对每个整数 $n$，若它为
$d$ 的倍数，则 $\mathcal{O}_X(n)|U$ 是可逆的，其逆为
$\mathcal{O}_X(-n)|U$。
\end{proposition}

\begin{proof}
由于 $q(U)$ 拟紧，它可由有限多个仿射开集 $V_i$ 覆盖；因而每个 $x\in U$
都包含在形如 $D_+(f)$ 的仿射开集中，其中 $f$ 是次数 $>0$ 的齐次元，
属于某个环 $\Gamma(V_i,\mathcal{S})$。由于 $U$ 拟紧，可用有限多个这样的开集
$D_+(f_j)$ 覆盖它；令 $d$ 为各 $f_j$ 的次数的一个公倍数。由
\egaref{II.2.5.17-zh}{2.5.17}，整数 $d$ 满足所求条件。
\end{proof}
% Locked French authority: ega2-1-fr.tex, 820505 B,
% SHA-256 84EDBE3E83530AF2959B441796337C9DC21EAFCA6A13114A26778760FBF437AC.
% Sequential source scope: lines 17677-17723, Environment II.8.14.5.

\begin{env}[8.14.5]
\label{II.8.14.5-zh}
沿用 \egaref{II.8.14.1-zh}{8.14.1} 的假设和记号，我们在
\egaref{II.3.3.2-zh}{3.3.2} 中定义了诸典范的 $\mathcal{O}_Y$-模同态
\begin{equation}
\label{II.8.14.5.1-zh}
  \alpha_n:\mathcal{M}_n\to
  q_*(\operatorname{Proj}_0(\mathcal{M}(n)))
  \qquad(n\in\mathbf{Z})
\tag{8.14.5.1}
\end{equation}
推广 \egaref{II.3.3.1-zh}{3.3.1} 的记号，对每个分次
$\mathcal{S}_X$-模 $\mathcal{F}$，我们置
\begin{equation}
\label{II.8.14.5.2-zh}
  \Gamma_*(\mathcal{F})=
  \bigoplus_{n\in\mathbf{Z}}q_*(\mathcal{F}_n).
\tag{8.14.5.2}
\end{equation}
特别地，
$\Gamma_*(\mathcal{S}_X)=\bigoplus_{n\in\mathbf{Z}}q_*(\mathcal{O}_X(n))$
就是 \egaref{II.3.3.1.2-zh}{3.3.1.2} 中记作
$\Gamma_*(\mathcal{O}_X)$ 的分次 $\mathcal{O}_Y$-代数；显然，
$\Gamma_*(\mathcal{F})$ 是一个分次 $\Gamma_*(\mathcal{S}_X)$-模
（\egaref{0.4.2.2-zh}{0, 4.2.2}）。当在同态

\oldpage[II]{199}
\egaref{II.8.14.5.1-zh}{8.14.5.1} 中取 $\mathcal{M}=\mathcal{S}$ 时，
便得到分次 $\mathcal{O}_Y$-代数同态
\begin{equation}
\label{II.8.14.5.3-zh}
  \alpha:\mathcal{S}\to\Gamma_*(\mathcal{S}_X)
\tag{8.14.5.3}
\end{equation}
这正是 \egaref{II.3.3.2-zh}{3.3.2} 中已经定义的同态，因而它使
$\Gamma_*(\mathcal{F})$ 成为一个分次 $\mathcal{S}$-模；于是诸同态
\egaref{II.8.14.5.1-zh}{8.14.5.1} 定义了一个分次
$\mathcal{S}$-模同态（次数为 $0$）
\begin{equation}
\label{II.8.14.5.4-zh}
  \alpha:\mathcal{M}\to\Gamma_*(\operatorname{Proj}(\mathcal{M})).
\tag{8.14.5.4}
\end{equation}
\end{env}
% Locked French authority: ega2-1-fr.tex, 820505 B,
% SHA-256 84EDBE3E83530AF2959B441796337C9DC21EAFCA6A13114A26778760FBF437AC.
% Sequential source scope: lines 17725-17790, Environment II.8.14.6.

\begin{env}[8.14.6]
\label{II.8.14.6-zh}
一般而言，对拟凝聚分次 $\mathcal{S}_X$-模 $\mathcal{F}$，不能确定分次
$\mathcal{S}$-模 $\Gamma_*(\mathcal{F})$ 必为拟凝聚的。考虑 $X$ 的一个
开集 $X'$，使 $q$ 在 $X'$ 上的限制 $q':X'\to Y$ 是\emph{拟紧}态射。
由于 $q'$ 还是分离的，对每个拟凝聚 $\mathcal{O}_{X'}$-模
$\mathcal{F}'$，$q'_*(\mathcal{F}')$ 是拟凝聚 $\mathcal{O}_Y$-模
（\egaref{I.9.2.2-zh}{I, 9.2.2, b}）。我们置
\begin{equation}
\label{II.8.14.6.1-zh}
  \mathcal{S}_{X'}=\mathcal{S}_X|X'
  =\bigoplus_{n\in\mathbf{Z}}\mathcal{O}_X(n)|X'
\tag{8.14.6.1}
\end{equation}
并且对每个分次 $\mathcal{S}_{X'}$-模 $\mathcal{F}'$，置
\begin{equation}
\label{II.8.14.6.2-zh}
  \Gamma'_*(\mathcal{F}')=
  \bigoplus_{n\in\mathbf{Z}}q'_*(\mathcal{F}'_n).
\tag{8.14.6.2}
\end{equation}
上述评注于是表明，若 $\mathcal{F}'$ 是拟凝聚
$\mathcal{S}_{X'}$-模，则 $\Gamma'_*(\mathcal{F}')$ 是拟凝聚分次
$\mathcal{S}$-模（\egaref{I.9.6.1-zh}{I, 9.6.1}）。

还要注意，典范单射 $j:X'\to X$ 是\emph{拟紧的}，因为
$q'=q\circ j$ 是拟紧的，而 $q$ 是分离的
（\egaref{I.6.6.4-zh}{I, 6.6.4, (v)}）。因此，对每个拟凝聚分次
$\mathcal{S}_{X'}$-模 $\mathcal{F}'$，
$\mathcal{F}=j_*(\mathcal{F}')$ 是拟凝聚分次 $\mathcal{S}_X$-模；由前述
定义得
\begin{equation}
\label{II.8.14.6.3-zh}
  \Gamma'_*(\mathcal{F}')=\Gamma_*(\mathcal{F}).
\tag{8.14.6.3}
\end{equation}
在对 $X'$ 的相同假设下，对每个拟凝聚分次 $\mathcal{S}$-模
$\mathcal{M}$，我们置
\begin{equation}
\label{II.8.14.6.4-zh}
  \operatorname{Proj}'(\mathcal{M})
  =\operatorname{Proj}(\mathcal{M})|X'
\tag{8.14.6.4}
\end{equation}
它是拟凝聚分次 $\mathcal{S}_{X'}$-模。典范同态
\[
  \operatorname{Proj}(\mathcal{M})\to
  j_*(\operatorname{Proj}'(\mathcal{M}))
\]
（\egaxref{0.4.4.3-zh}{0, 4.4.3}）因而给出分次 $\mathcal{S}$-模的典范同态
$\Gamma_*(\operatorname{Proj}(\mathcal{M}))\to
\Gamma'_*(\operatorname{Proj}'(\mathcal{M}))$；再把它与
\egaref{II.8.14.5.4-zh}{8.14.5.4} 复合，就得到拟凝聚分次
$\mathcal{S}$-模的函子性典范同态（次数为 $0$）
\begin{equation}
\label{II.8.14.6.5-zh}
  \alpha':\mathcal{M}\to
  \Gamma'_*(\operatorname{Proj}'(\mathcal{M})).
\tag{8.14.6.5}
\end{equation}
\end{env}
% Locked French authority: ega2-1-fr.tex, 820505 B,
% SHA-256 84EDBE3E83530AF2959B441796337C9DC21EAFCA6A13114A26778760FBF437AC.
% Sequential source scope: lines 17792-17827, Environment II.8.14.7.

\begin{env}[8.14.7]
\label{II.8.14.7-zh}
在 $X'$ 上保留 \egaref{II.8.14.6-zh}{8.14.6} 的假设，并设
$\mathcal{F}'$ 是拟凝聚分次 $\mathcal{S}_{X'}$-模，于是
$\operatorname{Proj}'(\Gamma'_*(\mathcal{F}'))$ 也是拟凝聚分次
$\mathcal{S}_{X'}$-模。我们将定义一个典范函子性同态（次数为 $0$），
它是分次 $\mathcal{S}_{X'}$-模的同态：
\begin{equation}
\label{II.8.14.7.1-zh}
  \beta':\operatorname{Proj}'(\Gamma'_*(\mathcal{F}'))\to\mathcal{F}'.
\tag{8.14.7.1}
\end{equation}

\oldpage[II]{200}
先假设 $Y=\operatorname{Spec}(A)$ 是仿射的，
$\mathcal{S}=\widetilde{S}$，其中 $S$ 是一个正次数分次 $A$-代数；此时
$\Gamma'_*(\mathcal{F}')=\widetilde{M}$，其中
$M=\bigoplus_{n\in\mathbf{Z}}\Gamma(X',\mathcal{F}'_n)$ 是分次 $S$-模。
设 $f\in S_d$ 满足 $D_+(f)\subset X'$；由定义
\egaref{II.2.6.2-zh}{2.6.2}，$\alpha_d(f)$ 在 $D_+(f)$ 上的限制，是
$\mathcal{O}_X(d)$ 在 $D_+(f)$ 上的截面，它对应于元 $f/1$，后者属于
$(S(d))_{(f)}$；该截面因而可逆。所以 $\alpha_d(f^n)$ 对每个 $n>0$
也都是可逆的。由此
立即得到分次模的一个 $S_f$-同态（次数为 $0$）
$\beta_f:M_f\to\Gamma(D_+(f),\mathcal{F}')$：它把每个元
$z/f^n\in M_f$（其中 $z\in M$）对应到截面
$(z|D_+(f))(\alpha_d(f^n)|D_+(f))^{-1}$；这是 $\mathcal{F}'$ 在
$D_+(f)$ 上的截面。此外，有一个与
\egaref{II.2.6.4.1-zh}{2.6.4.1} 对应的交换图表，由此便在此情形定义了
$\beta'$。为过渡到一般情形，必须考虑一个 $A$-代数 $A'$、分次
$A'$-代数 $S'=S\otimes_A A'$，并使用一个与
\egaref{II.2.8.13.2-zh}{2.8.13.2} 类似的交换图表；细节留给读者。
\end{env}
% Locked French authority: ega2-1-fr.tex, 820505 B,
% SHA-256 84EDBE3E83530AF2959B441796337C9DC21EAFCA6A13114A26778760FBF437AC.
% Sequential source scope: lines 17829-17864, Proposition II.8.14.8 and Corollary II.8.14.9 with proofs.

\begin{proposition}[8.14.8]
\label{II.8.14.8-zh}
若 $X'$ 是 $X=\operatorname{Proj}(\mathcal{S})$ 的开集，使得
$q':X'\to Y$ 拟紧，则 \egaref{II.8.14.7-zh}{8.14.7} 中定义的同态
$\beta'$ 是双射。
\end{proposition}

\begin{proof}
显然可只考虑 $Y$ 为仿射的情形；问题归结为证明（沿用
\egaref{II.8.14.7-zh}{8.14.7} 的记号）同态
$\beta_f:M_f\to\Gamma(D_+(f),\mathcal{F}')$ 是同构。将 $f$ 换成它的
某个幂，并不改变 $D_+(f)$ 或 $\beta_f$；由于按假设 $X'$ 是
\emph{拟紧的}，由 \egaref{II.8.14.4-zh}{8.14.4}，总可假设层
$\mathcal{O}_X(d)$ 是\emph{可逆的}。由于 $X'$ 是概形（因为 $q'$ 分离），
此时该命题正是 \egaref{I.9.3.1-zh}{I, 9.3.1}。
\end{proof}

\begin{corollary}[8.14.9]
\label{II.8.14.9-zh}
在 \egaref{II.8.14.8-zh}{8.14.8} 的假设下，每个拟凝聚分次
$\mathcal{S}_{X'}$-模 $\mathcal{F}'$ 都同构于一个形如
$\mathcal{S}_{X'}$-分次模 $\operatorname{Proj}'(\mathcal{M})$ 的模，其中
$\mathcal{M}$ 是拟凝聚分次 $\mathcal{S}$-模。此外，若 $\mathcal{F}'$ 为
有限型，并且假设 $Y$ 是拟紧概形，或其底空间为诺特空间的预概形，则可假设
$\mathcal{M}$ 为有限型。
\end{corollary}

\begin{proof}
由 \egaref{II.8.14.8-zh}{8.14.8} 推出本推论的证明，完全仿照
\egaref{II.3.4.5-zh}{3.4.5} 的证明：后者从
\egaref{II.3.4.4-zh}{3.4.4} 出发；细节留给读者。
\end{proof}
% Locked French authority: ega2-1-fr.tex, 820505 B,
% SHA-256 84EDBE3E83530AF2959B441796337C9DC21EAFCA6A13114A26778760FBF437AC.
% Sequential source scope: lines 17866-17894, Proposition II.8.14.10 and proof.

\begin{proposition}[8.14.10]
\label{II.8.14.10-zh}
在 \egaref{II.8.14.7-zh}{8.14.7} 的假设下，设
$\mathcal{M}$ 是一个拟凝聚分次 $\mathcal{S}$-模，
$\mathcal{F}'$ 是一个拟凝聚分次 $\mathcal{S}_{X'}$-模；复合同态
\begin{equation}
\label{II.8.14.10.1-zh}
  \operatorname{Proj}'(\mathcal{M})
  \xrightarrow{\operatorname{Proj}'(\alpha')}
  \operatorname{Proj}'(\Gamma'_*(\operatorname{Proj}'(\mathcal{M})))
  \xrightarrow{\beta'}\operatorname{Proj}'(\mathcal{M})
\tag{8.14.10.1}
\end{equation}
以及
\begin{equation}
\label{II.8.14.10.2-zh}
  \Gamma'_*(\mathcal{F}')
  \xrightarrow{\alpha'}
  \Gamma'_*(\operatorname{Proj}'(\Gamma'_*(\mathcal{F}')))
  \xrightarrow{\Gamma'_*(\beta')}\Gamma'_*(\mathcal{F}')
\tag{8.14.10.2}
\end{equation}
都是恒等同构。
\end{proposition}

\begin{proof}
问题在 $Y$ 上是局部的，并且验证与 \egaref{II.2.6.5-zh}{2.6.5}
中的验证相同；我们把细节留给读者。
\end{proof}
% Locked French authority: ega2-1-fr.tex, 820505 B,
% SHA-256 84EDBE3E83530AF2959B441796337C9DC21EAFCA6A13114A26778760FBF437AC.
% Sequential source scope: lines 17896-17908, Remark II.8.14.11.

\begin{remark}[8.14.11]
\label{II.8.14.11-zh}
我们将在第三章（\egaref{III.2.3.1-zh}{III, 2.3.1}）中看到：当
$Y$ 是\emph{局部诺特的}，而 $\mathcal{S}$ 是一个\emph{有限型}拟凝聚
分次 $\mathcal{O}_Y$-代数时（此时

\oldpage[II]{201}
可取 $X'=X$），对于每个满足条件 \textbf{(TF)} 的拟凝聚分次
$\mathcal{S}$-模 $\mathcal{M}$，同态 $\alpha$
（\egaref{II.8.14.5.4-zh}{8.14.5.4}）都是
\textbf{(TN)}-\emph{双射}。
\end{remark}
% Locked French authority: ega2-1-fr.tex, 820505 B,
% SHA-256 84EDBE3E83530AF2959B441796337C9DC21EAFCA6A13114A26778760FBF437AC.
% Sequential source scope: lines 17910-17976, Remark II.8.14.12.

\begin{remark}[8.14.12]
\label{II.8.14.12-zh}
在 \egaref{II.8.14.4-zh}{8.14.4} 中描述的情形是下述情形的一个特例。
设 $X$ 是一个环空间，$\mathcal{S}$ 是一个 $\mathcal{O}_X$-分次代数
（其次数可正可负）；假设存在一个整数 $d>0$，使得 $\mathcal{S}_d$ 和
$\mathcal{S}_{-d}$ 都是\emph{可逆的}，并且典范同态
\begin{equation}
\label{II.8.14.12.1-zh}
  \mathcal{S}_d\otimes_{\mathcal{O}_X}\mathcal{S}_{-d}
  \to\mathcal{O}_X
\tag{8.14.12.1}
\end{equation}
是一个\emph{同构}（因而 $\mathcal{S}_{-d}$ 与
$\mathcal{S}_d^{-1}$ 等同）。此时称分次 $\mathcal{O}_X$-代数
$\mathcal{S}$ \emph{是周期的，其周期为 $d$}。这一名称源于以下性质：
\emph{在上述假设下，对于任意分次 $\mathcal{S}$-模 $\mathcal{F}$，
典范同态}
\begin{equation}
\label{II.8.14.12.2-zh}
  \mathcal{S}_d\otimes\mathcal{F}_n\to\mathcal{F}_{n+d}
\tag{8.14.12.2}
\end{equation}
\emph{对每个 $n\in\mathbf{Z}$ 都是同构。}事实上，这个问题在 $X$ 上
是局部的，因而可以假设 $\mathcal{S}_d$ 在 $X$ 上有一个\emph{可逆}
截面 $s$，其逆 $s'$ 是 $\mathcal{S}_{-d}$ 的一个截面。对于任意截面
$z\in\Gamma(U,\mathcal{F}_{n+d})$，令其对应于 $U$ 上
$\mathcal{S}_d\otimes\mathcal{F}_n$ 的截面
$(s|U)\otimes(s'|U)z$，由此得到的同态
$\mathcal{F}_{n+d}\to\mathcal{S}_d\otimes\mathcal{F}_n$ 正是
\egaref{II.8.14.12.2-zh}{8.14.12.2} 的逆同态，这就证明了所断言的结论。
由此，对于每个 $k\in\mathbf{Z}$，得到一个典范同构
\[
  (\mathcal{S}_d)^{\otimes k}\otimes\mathcal{F}_n
  \xrightarrow{\sim}\mathcal{F}_{n+kd}.
\]
因此，\emph{若已知一个分次 $\mathcal{S}$-模 $\mathcal{F}$ 的各
$\mathcal{S}_0$-模 $\mathcal{F}_i$ $(0\leqslant i\leqslant d-1)$
以及下列典范同态，则该分次模即被确定：}
\[
  \mathcal{S}_i\otimes\mathcal{F}_j\to\mathcal{F}_{i+j}
  \qquad\text{对 }0\leqslant i,j\leqslant d-1
\]
（当 $i+j\geqslant d$ 时，约定
$\mathcal{F}_{i+j}=\mathcal{S}_d\otimes_{\mathcal{S}_0}
\mathcal{F}_{i+j-d}$）。当然，为使这些同态确实在
$(\mathcal{S}_d)^{\otimes k}\otimes\mathcal{F}_i$
$(k\in\mathbf{Z},\ 0\leqslant i\leqslant d-1)$ 的直和上定义一个
$\mathcal{S}$-模结构，它们必须满足某些结合性条件；这里不再把这些条件
明确写出。

当 $d=1$ 时（这正是 \egaref{subsection:II.3.3-zh}{3.3} 中考察的情形），
可以说，分次 $\mathcal{S}$-模范畴（相应地，当 $X$ 是预概形且
$\mathcal{S}$ 拟凝聚时，指其中的拟凝聚对象所成范畴）与任意
$\mathcal{S}_0$-模范畴（相应地，其中的拟凝聚对象所成范畴）
\emph{等价}；正是在这个意义下，可以把本条的论述看作
\S~\egaref{section:II.3-zh}{3} 中论述的推广。此外可以看出，在适当的
有限性条件下，本条的结果与 \egaref{II.8.14.11-zh}{8.14.11} 合在一起，
在某种程度上把预概形上的拟凝聚分次代数以及这类代数上的分次模
«~模 \textbf{(TN)}~» 的研究，归结为所考察代数\emph{是周期的}这一特例
（此时，$\mathcal{M}$ 的条件 \textbf{(TN)}
（\egaref{II.3.4.2-zh}{3.4.2}）因而蕴含 $\mathcal{M}=0$）。
\end{remark}
% Locked French authority: ega2-1-fr.tex, 820505 B,
% SHA-256 84EDBE3E83530AF2959B441796337C9DC21EAFCA6A13114A26778760FBF437AC.
% Sequential source scope: lines 17978-18058, Remark II.8.14.13 / printed pp.201-202.

\begin{remark}[8.14.13]
\label{II.8.14.13-zh}
在 \egaref{II.8.14.1-zh}{8.14.1} 的假设下，设 $d$ 是整数且 $>0$；
已经定义了一个典范 $Y$-同构 $h$，它把 $X$ 映到
$X^{(d)}=\operatorname{Proj}(\mathcal{S}^{(d)})$
（\egaref{II.3.1.8-zh}{3.1.8}）。对于每个

\oldpage[II]{202}
拟凝聚分次 $\mathcal{S}$-模 $\mathcal{M}$ 和每个整数 $k$，若其满足
$0\leqslant k\leqslant d-1$，也有一个典范 $h$-同构
（沿用 \egaref{II.3.1.1-zh}{3.1.1} 的记号）
\begin{equation}
\label{II.8.14.13.1-zh}
  (\operatorname{Proj}(\mathcal{M}))^{(d,k)}
  \xleftarrow{\sim}\operatorname{Proj}(\mathcal{M}^{(d,k)}).
\tag{8.14.13.1}
\end{equation}

先假设 $Y=\operatorname{Spec}(A)$ 为仿射，
$\mathcal{S}=\widetilde{S}$、$\mathcal{M}=\widetilde{M}$，其中 $S$ 是按
正次数分次的 $A$-代数，而 $M$ 是分次 $S$-模。已知对每个
$f\in S_e$（$e>0$），$h$ 把 $D_+(f)$ 映到 $D_+(f^d)$，并对应于典范同构
$S_{(f^d)}\xrightarrow{\sim}S_{(f)}$
（\egaref{II.2.2.2-zh}{2.2.2}）。于是，
\egaref{II.8.14.13.1-zh}{8.14.13.1} 在 $D_+(f^d)$ 上的限制对应于典范双同构
$M_{f^d}\xrightarrow{\sim}M_f$ 在 $M_{f^d}$ 中次数与 $k$ 模 $d$ 同余的
各元素上的限制。我们把下列验证留给读者：这些同构与把 $f$ 换为它的齐次
倍数 $fg$ 相容；然后，当把 $S$ 换为一个分次 $A'$-代数
$S'=S\otimes_A A'$ 时，也有类似的相容性，其中 $A'$ 是 $A$-代数。

特别地，由此有一个 $h$-同构
\begin{equation}
\label{II.8.14.13.2-zh}
  (\mathcal{S}^{(d)})_{X^{(d)}}
  \xrightarrow{\sim}(\mathcal{S}_X)^{(d)}
\tag{8.14.13.2}
\end{equation}
它保持两边的乘法结构；借助这个同构，
\egaref{II.8.14.13.1-zh}{8.14.13.1} 成为一个 $h$-双同构，它把分次
$(\mathcal{S}^{(d)})_{X^{(d)}}$-模映到分次
$(\mathcal{S}_X)^{(d)}$-模。同样，有一个 $h$-同构
\begin{equation}
\label{II.8.14.13.3-zh}
  \operatorname{Proj}_0(\mathcal{S}^{(d,k)}(n))
  \xrightarrow{\sim}\mathcal{O}_X(nd+k)
\tag{8.14.13.3}
\end{equation}
它补充了 \egaref{II.3.2.9-zh}{3.2.9, (ii)} 的结果。

由同构 \egaref{II.8.14.13.1-zh}{8.14.13.1} 立即得出分次
$\mathcal{S}^{(d)}$-模的一个同构
\begin{equation}
\label{II.8.14.13.4-zh}
  \Gamma_*^{(d)}(\operatorname{Proj}(\mathcal{M}^{(d,k)}))
  \xrightarrow{\sim}
  \Gamma_*((\operatorname{Proj}(\mathcal{M}))^{(d,k)})
\tag{8.14.13.4}
\end{equation}
其中 $\Gamma_*^{(d)}$ 对应于结构态射
$q^{(d)}:X^{(d)}\to Y$。直接验证可知，典范同态 $\alpha$
（\egaref{II.8.14.5.4-zh}{8.14.5.4}）和类似同态 $\alpha^{(d)}$
（对应于 $X^{(d)}$）使图表
\begin{equation}
\label{II.8.14.13.5-zh}
  \xymatrix{
    & \mathcal{M}^{(d,k)}\ar[dl]_{\alpha^{(d)}}\ar[dr]^{\alpha} & \\
    \Gamma_*^{(d)}(\operatorname{Proj}(\mathcal{M}^{(d,k)}))
      \ar[rr]^{\sim}
      && \Gamma_*((\operatorname{Proj}(\mathcal{M}))^{(d,k)})
  }
\tag{8.14.13.5}
\end{equation}
交换。验证时，假设 $Y$ 为仿射，并计算在 $\alpha^{(d)}$ 和 $\alpha$ 下，
由 $M^{(d,k)}$ 中同一元素所得之像在开集 $D_+(f^d)$ 和 $D_+(f)$ 上的限制
（记号同上）。
\end{remark}
% Locked French authority: ega2-1-fr.tex, 820505 B,
% SHA-256 84EDBE3E83530AF2959B441796337C9DC21EAFCA6A13114A26778760FBF437AC.
% Sequential source scope: lines 18060-18087, Proposition II.8.14.14 and proof.

\begin{proposition}[8.14.14]
\label{II.8.14.14-zh}
设 $Y$ 是一个拟紧预概形，$\mathcal{S}$ 是一个有限型拟凝聚分次
$\mathcal{O}_Y$-代数，$\mathcal{M}$ 是一个满足条件 \textbf{(TF)} 的
拟凝聚分次 $\mathcal{S}$-模；置 $X=\operatorname{Proj}(\mathcal{S})$。
那么 $\mathcal{S}_X$ 是一个周期分次 $\mathcal{O}_X$-代数
（\egaref{II.8.14.12-zh}{8.14.12}），并且 $\mathcal{S}_X$ 存在一个

\oldpage[II]{203}
周期 $d$，使得诸
$(\operatorname{Proj}(\mathcal{M}))^{(d,k)}$
$(0\leqslant k\leqslant d-1)$ 都是有限型
$(\mathcal{S}_X)^{(d)}$-模。
\end{proposition}

\begin{proof}
事实上，\egaref{II.3.1.10-zh}{3.1.10} 证明存在 $d$，使得
$\mathcal{S}^{(d)}$ 由
$\mathcal{S}_d=(\mathcal{S}^{(d)})_1$ 生成，而后者是一个有限型
$\mathcal{S}_0$-模。因而，为证明第一个断言，根据
\egaref{II.8.14.13.2-zh}{8.14.13.2} 可归结到 $d=1$ 的情形；这时命题由
\egaref{II.3.2.7-zh}{3.2.7} 得出。此外，结合
\egaref{II.8.14.13.1-zh}{8.14.13.1}，第二个断言是
\egaref{II.2.1.6-zh}{2.1.6, (iii)} 和
\egaref{II.3.4.3-zh}{3.4.3} 的推论。
\end{proof}
\clearpage
\renewcommand{\bibname}{参考文献\textit{（续）}}

\begin{thebibliography}{26}
\phantomsection
\addcontentsline{toc}{section}{参考文献（续）}

\bibitem[23]{II-23}
\oldpage[II]{205}
C.~Chevalley,
\emph{Introduction to the theory of algebraic functions of one variable}.
Math. Surveys, n\textsuperscript{o}~6, Amer. Math. Soc., New York, 1951.

\bibitem[24]{II-24}
P.~Jaffard,
\emph{Les systèmes d'idéaux}, Paris (Dunod), 1960.

\bibitem[25]{II-25}
M.~Nagata,
\emph{On the derived normal rings of Noetherian integral domains},
Mem. Coll. Sci. Kyoto, sér.~A, t.~XXIX (1955), p.~293--303.

\bibitem[26]{II-26}
M.~Nagata,
\emph{Existence theorems for non projective complete algebraic varieties},
Ill. J. Math., t.~II (1958), p.~490--498.

\end{thebibliography}
\clearpage
\oldpage[II]{207}
\markboth{A. GROTHENDIECK --- 第二章}{A. GROTHENDIECK --- 第二章}
\phantomsection
\addcontentsline{toc}{section}{记号索引}
\section*{记号索引}
\thispagestyle{plain}

\begingroup
\small
\setlength{\parindent}{0pt}
\setlength{\parskip}{2pt}
\newcommand{\egaShProj}{\mathcal{P}\!\operatorname{roj}}

$\mathcal{A}(X),\mathcal{A}(\mathcal{F}),\mathcal{A}(\mathcal{B})$
（$X$ 为 $S$-预概形，$\mathcal{F}$ 为 $\mathcal{O}_X$-模，
$\mathcal{B}$ 为 $\mathcal{O}_X$-代数）：\hyperref[II.1.1.1-zh]{1.1.1}。

$\mathcal{A}(h)$（$h$ 为 $S$-态射）：\hyperref[II.1.1.2-zh]{1.1.2}。

$\Spec(\mathcal{B})$（$\mathcal{B}$ 为 $\mathcal{O}_S$-代数）：\hyperref[II.1.3.1-zh]{1.3.1}。

$\widetilde{\mathcal{M}}$（$\mathcal{M}$ 为 $\mathcal{B}$-模，
$\mathcal{B}$ 为 $\mathcal{O}_S$-代数）：\hyperref[II.1.4.3-zh]{1.4.3}。

$\mathbf{T}(E),\mathbf{S}(E),\mathbf{S}_A(E)$（$E$ 为 $A$-模）：\hyperref[II.1.7.1-zh]{1.7.1}。

$\mathbf{S}(\mathcal{E}),\mathbf{S}_{\mathcal{A}}(\mathcal{E})$
（$\mathcal{E}$ 为 $\mathcal{A}$-模）：\hyperref[II.1.7.4-zh]{1.7.4}。

$\mathbf{V}(\mathcal{E})$（$\mathcal{E}$ 为拟凝聚 $\mathcal{O}_S$-模）：\hyperref[II.1.7.8-zh]{1.7.8}。

$S_n,S_+,M_n,S^{(d)},M^{(d,k)},M^{(d)}$
（$S$ 为分次环，$M$ 为分次 $S$-模）：\hyperref[II.2.1.1-zh]{2.1.1}。

$M(n),S(n)$（$S$ 为分次环，$M$ 为分次 $S$-模）：\hyperref[II.2.1.1-zh]{2.1.1}。

$\Hom_S(M,N)$（$S$ 为分次环，$M,N$ 为分次 $S$-模）：\hyperref[II.2.1.2-zh]{2.1.2}。

$\mathfrak N_+,\mathfrak r_+(\mathfrak J)$
（$\mathfrak J$ 为 $S_+$ 的理想）：\hyperref[II.2.1.10-zh]{2.1.10}。

$S_{(f)},M_{(f)}$（$f$ 为 $S_+$ 的齐次元）：\hyperref[II.2.2.1-zh]{2.2.1}。

$S_{(T)},M_{(T)},S_{(\mathfrak p)},M_{(\mathfrak p)}$
（$T$ 为 $S_+$ 中由齐次元组成的乘法子集，
$\mathfrak p$ 为 $S_+$ 的分次素理想）：\hyperref[II.2.2.7-zh]{2.2.7}。

$\Proj(S)$（$S$ 为分次环）：\hyperref[II.2.3.1-zh]{2.3.1}。

$V_+(E)$（$E$ 为 $S_+$ 的子集）：\hyperref[II.2.3.2-zh]{2.3.2}。

$D_+(f)$（$f$ 为 $S$ 的元素）：\hyperref[II.2.3.3-zh]{2.3.3}。

$\mathfrak j_+(Y)$（$Y$ 为 $\Proj(S)$ 的子集）：\hyperref[II.2.3.10-zh]{2.3.10}。

$\widetilde M$（$M$ 为分次 $S$-模）：\hyperref[II.2.5.2-zh]{2.5.2}。

$\widetilde u$（$u$ 为次数零的同态）：\hyperref[II.2.5.4-zh]{2.5.4}。

$\mathcal{O}_X(n),\mathcal{F}(n)$
（$X=\Proj(S)$，$\mathcal{F}$ 为 $\mathcal{O}_X$-模）：\hyperref[II.2.5.10-zh]{2.5.10}。

$\lambda$（典范同态）：\hyperref[II.2.5.11-zh]{2.5.11}。

$\mu$（典范同态）：\hyperref[II.2.5.12-zh]{2.5.12}。

$X^{(d)}$（$X=\Proj(S)$，$d$ 为整数且 $>0$）：\hyperref[II.2.5.16-zh]{2.5.16}。

$\Gamma_*(\mathcal{F})$（$\mathcal{F}$ 为 $\mathcal{O}_X$-模，
$X=\Proj(S)$）：\hyperref[II.2.6.1-zh]{2.6.1}。

$\alpha_n,\alpha,\alpha_{\mathcal{M}}$（典范同态）：\hyperref[II.2.6.2-zh]{2.6.2}。

$\beta,\beta_{\mathcal{F}}$（典范同态）：\hyperref[II.2.6.4-zh]{2.6.4}。

$G(\vphi),{}^a\vphi$（滥用语言）
（$\vphi$ 为分次环同态）：\hyperref[II.2.8.1-zh]{2.8.1}。

$\Proj(\vphi)$（$\vphi$ 为分次环同态）：\hyperref[II.2.8.2-zh]{2.8.2}。

$\mathcal{S}^{(d)},\mathcal{M}^{(d,k)},\mathcal{M}^{(d)},\mathcal{M}(n)$
（$\mathcal{S}$ 为分次 $\mathcal{O}_Y$-代数，
$\mathcal{M}$ 为分次 $\mathcal{S}$-模）：\hyperref[II.3.1.1-zh]{3.1.1}。

$\Proj(\mathcal{S})$（$\mathcal{S}$ 为拟凝聚分次 $\mathcal{O}_Y$-代数）：\hyperref[II.3.1.3-zh]{3.1.3}。

$X_f$（$X=\Proj(\mathcal{S})$，$f$ 为 $\mathcal{S}_d$ 在 $Y$ 上的截面）：\hyperref[II.3.1.4-zh]{3.1.4}。

$\widetilde{\mathcal{M}}$（$\mathcal{M}$ 为分次 $\mathcal{S}$-模）：\hyperref[II.3.2.2-zh]{3.2.2}。

$\mathcal{O}_X(n),\mathcal{F}(n)$
（$X=\Proj(\mathcal{S})$，$\mathcal{F}$ 为 $\mathcal{O}_X$-模）：\hyperref[II.3.2.5-zh]{3.2.5}。

$\lambda,\mu$（典范同态）：\hyperref[II.3.2.6-zh]{3.2.6}。

$\Gamma_*(\mathcal{F})$（$\mathcal{F}$ 为 $\mathcal{O}_X$-模，
$X=\Proj(\mathcal{S})$）：\hyperref[II.3.3.1-zh]{3.3.1}。

$\alpha_n,\alpha,\alpha_{\mathcal{M}}$（典范同态）：\hyperref[II.3.3.2-zh]{3.3.2}。

\oldpage[II]{208}
\markboth{A. GROTHENDIECK --- 第二章}{A. GROTHENDIECK --- 第二章}
$\beta,\beta_{\mathcal{F}}$（典范同态）：\hyperref[II.3.3.4-zh]{3.3.4}。

$G(\vphi),\Proj(\vphi)$
（$\vphi$ 为分次 $\mathcal{O}_Y$-代数同态）：\hyperref[II.3.5.1-zh]{3.5.1}。

$\Proj(u)$（$u$ 为一个 $g$-态射，它从一个分次 $\mathcal{O}_Y$-代数到一个分次
$\mathcal{O}_{Y'}$-代数）：\hyperref[II.3.5.6-zh]{3.5.6}。

$r_{\mathcal{L},\psi}$（$\mathcal{L}$ 为可逆 $\mathcal{O}_X$-模）：\hyperref[II.3.7.1-zh]{3.7.1}。

$\mathbf{P}(\mathcal{E}),\mathbf{P}(E),\mathbf{P}_Y^n,\mathbf{P}_A^n$：\hyperref[II.4.1.1-zh]{4.1.1}。

$\mathbf{P}(u)$（$u$ 为 $\mathcal{O}_Y$-模的满射同态）：\hyperref[II.4.1.2-zh]{4.1.2}。

$\Hyp_Y(X,\mathcal{E})$（$\mathcal{E}$ 为拟凝聚 $\mathcal{O}_Y$-模）：\hyperref[II.4.2.5-zh]{4.2.5}。

$\varsigma$（塞格雷态射）：\hyperref[II.4.3.1-zh]{4.3.1}。

$\mathcal{F}(n)$（$\mathcal{F}$ 为 $\mathcal{O}_X$-模，$X$ 配备一个
可逆 $\mathcal{O}_X$-模）：\hyperref[II.4.5.1-zh]{4.5.1}。

$\varepsilon$（恒等自同构）：\hyperref[II.4.5.1-zh]{4.5.1}。

$P(u,T),\sigma_i(u)$（$u$ 为有限秩自由模的自同态）：\hyperref[II.6.4.1-zh]{6.4.1}。

$P(u,T),\sigma_i(u),\det u$
（$u$ 为整环或约化诺特环上的有限型模的自同态）：\hyperref[II.6.4.2-zh]{6.4.2} 和 \hyperref[II.6.4.7-zh]{6.4.7}。

$\sigma_i(u),\det u$（$u$ 为局部自由 $\mathcal{A}$-模的自同态）：\hyperref[II.6.4.8-zh]{6.4.8}。

$\sigma_i,\det$（层同态）：\hyperref[II.6.4.8-zh]{6.4.8}、\hyperref[II.6.4.9-zh]{6.4.9} 和 \hyperref[II.6.4.10-zh]{6.4.10}。

$\sigma_i(u),\det u$
（$u$ 为局部整预概形或局部诺特约化预概形上的有限型
$\mathcal{O}_X$-模的自同态）：
\hyperref[II.6.4.9-zh]{6.4.9} 和 \hyperref[II.6.4.10-zh]{6.4.10}。

$N_{\mathcal{B}/\mathcal{A}}(f)$
（$f$ 为 $\mathcal{B}$-代数 $\mathcal{A}$ 的截面）：\hyperref[II.6.5.1-zh]{6.5.1}。

$N_{\mathcal{B}/\mathcal{A}}(\mathcal{L}')$
（$\mathcal{L}'$ 为可逆 $\mathcal{B}$-模）：\hyperref[II.6.5.2-zh]{6.5.2}。

$N_{\mathcal{B}/\mathcal{A}}(h')$
（$h'$ 为可逆 $\mathcal{B}$-模的同态）：\hyperref[II.6.5.3-zh]{6.5.3}。

$N_{X'/X}(\mathcal{L}')$
（$X'$ 在 $X$ 上有限，$\mathcal{L}'$ 为可逆 $\mathcal{O}_{X'}$-模）：\hyperref[II.6.5.5-zh]{6.5.5}。

$N_{X'/X}(h')$
（$h'$ 为可逆 $\mathcal{O}_{X'}$-模的同态）：\hyperref[II.6.5.5-zh]{6.5.5}。

$\mathcal{S}_X$（$X=\Proj(\mathcal{S})$，$\mathcal{S}$ 为
$\mathcal{O}_Y$ 的诸理想之和）：\hyperref[II.8.1.5-zh]{8.1.5}。

$S^\geq,S^\leq,S_f^\geq,S_f^\leq,
M^\geq,M^\leq,M_f^\geq,M_f^\leq$
（$S$ 为分次环，$M$ 为分次 $S$-模，
$f$ 为 $S_+$ 的齐次元）：\hyperref[II.8.2.1-zh]{8.2.1}。

$\widehat S$（$S$ 为分次环）：\hyperref[II.8.2.2-zh]{8.2.2}。

$\widehat M$（$M$ 为分次 $S$-模）：\hyperref[II.8.2.4-zh]{8.2.4}。

$S_{[n]},S^\natural,f^\natural$
（$S$ 为分次环，$f$ 为 $S$ 中的齐次元）：\hyperref[II.8.2.6-zh]{8.2.6}。

$M_{[n]},M^\natural$（$M$ 为分次 $S$-模）：\hyperref[II.8.2.8-zh]{8.2.8}。

$\widehat{\mathcal{S}}$（$\mathcal{S}$ 为分次 $\mathcal{O}_Y$-代数）：\hyperref[II.8.3.1-zh]{3.8.1}。

$G_m$（群概形）：\hyperref[II.8.3.9-zh]{8.3.9}。

$\mathcal{S}_X$（$X=\Proj(\mathcal{S})$，$\mathcal{S}$ 为分次 $\mathcal{O}_Y$-代数）：\hyperref[II.8.6.1-zh]{8.6.1}。

$\mathcal{S}_X,C_X,\widehat C_X,E_X,\widehat E_X$
（$X=\Proj(\mathcal{S})$，$\mathcal{S}$ 为分次 $\mathcal{O}_Y$-代数，
且 $\mathcal{S}_0=\mathcal{O}_Y$）：\hyperref[II.8.6.1-zh]{8.6.1}。

$\mathcal{S}_{[n]},\mathcal{S}^{\natural}$
（$\mathcal{S}$ 为分次 $\mathcal{O}_Y$-代数）：\hyperref[II.8.7.2-zh]{8.7.2}。

$\egaShProj_0(\mathcal{M}),\egaShProj(\mathcal{M}),\mathcal{M}_X,
\widehat{\mathcal{M}},\mathcal{M}^{\square}$
（$\mathcal{M}$ 为分次 $\mathcal{S}$-模）：\hyperref[II.8.12.1-zh]{8.12.1}。

$\mathcal{M}_X^\geq$
（$X=\Proj(\mathcal{S})$，$\mathcal{M}$ 为分次 $\mathcal{S}$-模）：\hyperref[II.8.12.5-zh]{8.12.5}。

$\mathcal{M}_{[n]},\mathcal{M}^{\natural}$
（$\mathcal{M}$ 为分次 $\mathcal{S}$-模）：\hyperref[II.8.12.9-zh]{8.12.9}。

$(\widetilde{\mathcal{N}})^-$
（$\mathcal{N}$ 为某个分次 $\mathcal{S}$-模的子-$\mathcal{S}$-模）：\hyperref[II.8.13.1-zh]{8.13.1}。

$\overline{\mathcal{N}}$
（$\mathcal{N}$ 为某个分次 $\mathcal{S}$-模的子-$\mathcal{S}$-模）：\hyperref[II.8.13.2-zh]{8.13.2}。

$\Gamma_*(\mathcal{F})$
（$X=\Proj(\mathcal{S})$，$\mathcal{F}$ 为分次 $\mathcal{S}_X$-模）：\hyperref[II.8.14.5-zh]{8.14.5}。

$\mathcal{S}_{X'},\Gamma'_*(\mathcal{F}'),\egaShProj'(\mathcal{M}),
\alpha',\beta'$
（$X'$ 为 $X=\Proj(\mathcal{S})$ 的开集，$\mathcal{F}'$ 为分次 $\mathcal{S}_{X'}$-模，
$\mathcal{M}$ 为分次 $\mathcal{S}$-模）：\hyperref[II.8.14.6-zh]{8.14.6}。

\endgroup

\clearpage
\oldpage[II]{209}
\markboth{A. GROTHENDIECK --- 第二章}{A. GROTHENDIECK --- 第二章}
\thispagestyle{plain}
\phantomsection
\addcontentsline{toc}{section}{术语索引}
\section*{术语索引}

\begingroup
\small
\setlength{\parindent}{0pt}
\setlength{\parskip}{2pt}
\newcommand{\egaTerme}[2]{\textbf{#1}：#2。\par}

\egaTerme{仿射（态射）}{\hyperref[II.1.6.1-zh]{1.6.1}}
\egaTerme{在一个预概形上仿射（预概形）}{\hyperref[II.1.2.1-zh]{1.2.1}}
\egaTerme{分次环上的分次代数}{\hyperref[II.2.1.2-zh]{2.1.2}}
\egaTerme{$A$-模的对称代数}{\hyperref[II.1.7.1-zh]{1.7.1}}
\egaTerme{$\mathcal{A}$-模的对称 $\mathcal{A}$-代数}{\hyperref[II.1.7.4-zh]{1.7.4}}
\egaTerme{本质约化的分次 $\mathcal{O}_Y$-代数}{\hyperref[II.3.1.12-zh]{3.1.12}}
\egaTerme{整的分次 $\mathcal{O}_Y$-代数}{\hyperref[II.3.1.12-zh]{3.1.12}}
\egaTerme{丰沛（可逆 $\mathcal{O}_X$-模）}{\hyperref[II.4.5.3-zh]{4.5.3}}
\egaTerme{相对于 $f$ 丰沛、$f$-丰沛、$Y$-丰沛（可逆 $\mathcal{O}_X$-模）}{\hyperref[II.4.6.1-zh]{4.6.1}}
\egaTerme{本质整分次环}{\hyperref[II.2.1.11-zh]{2.1.11}}
\egaTerme{本质约化分次环}{\hyperref[II.2.1.10-zh]{2.1.10}}
\egaTerme{与一个 $\mathcal{B}$-模相伴（层）}{\hyperref[II.1.4.3-zh]{1.4.3}}
\egaTerme{与一个分次 $S$-模相伴（层）}{\hyperref[II.2.5.3-zh]{2.5.3}}
\egaTerme{与一个分次 $\mathcal{S}$-模相伴（层）}{\hyperref[II.3.2.2-zh]{3.2.2}}
\egaTerme{与一个 $\mathcal{O}_S$-代数相伴（$S$-概形）}{\hyperref[II.1.3.1-zh]{1.3.1}}

\egaTerme{条件 (TF)、条件 (TN)}{\hyperref[II.2.7.2-zh]{2.7.2}}
\egaTerme{条件 (TF)、条件 (TN)}{\hyperref[II.3.4.2-zh]{3.4.2}}
\egaTerme{由分次 $\mathcal{O}_Y$-代数定义的仿射锥、射影锥}{\hyperref[II.8.3.1-zh]{8.3.1}}
\egaTerme{由分次 $\mathcal{O}_Y$-代数定义的去顶点仿射锥、去顶点射影锥}{\hyperref[II.8.3.4-zh]{8.3.4}}
\egaTerme{预概形 $\Proj(\mathcal{S})$ 的仿射投影锥、射影投影锥}{\hyperref[II.8.3.1-zh]{8.3.1}}
\egaTerme{域上的代数曲线}{\hyperref[II.7.4.2-zh]{7.4.2}}
\egaTerme{完备代数曲线}{\hyperref[II.7.4.7-zh]{7.4.7}}
\egaTerme{塞尔判据}{\hyperref[II.5.2.1-zh]{5.2.1}}

\egaTerme{整环上的有限型模的自同态的行列式}{\hyperref[II.6.4.2-zh]{6.4.2}}
\egaTerme{约化诺特环上的有限型模的自同态的行列式}{\hyperref[II.6.4.7-zh]{6.4.7}}
\egaTerme{局部自由 $\mathcal{A}$-模的自同态的行列式}{\hyperref[II.6.4.8-zh]{6.4.8}}
\egaTerme{有限型 $\mathcal{O}_X$-模的自同态在局部整预概形 $X$ 上的行列式}{\hyperref[II.6.4.9-zh]{6.4.9}}
\egaTerme{有限型 $\mathcal{O}_X$-模的自同态在局部诺特且约化的预概形 $X$ 上的行列式}{\hyperref[II.6.4.10-zh]{6.4.10}}

\egaTerme{暴涨（预概形）}{\hyperref[II.8.1.3-zh]{8.1.3}}
\egaTerme{整（态射）}{\hyperref[II.6.1.1-zh]{6.1.1}}
\egaTerme{在一个预概形上整（预概形）}{\hyperref[II.6.1.1-zh]{6.1.1}}
\egaTerme{整的（拟凝聚 $\mathcal{O}_S$-代数）}{\hyperref[II.6.1.2-zh]{6.1.2}}
\egaTerme{$\mathcal{A}$-代数的一个截面在 $\mathcal{A}$ 上整}{\hyperref[II.6.3.1-zh]{6.3.1}}

\egaTerme{射影丛的基本层}{\hyperref[II.4.1.1-zh]{4.1.1}}
\egaTerme{环层 $\mathcal{A}$ 在一个 $\mathcal{A}$-代数中的整闭包}{\hyperref[II.6.3.2-zh]{6.3.2}}
\egaTerme{预概形 $X$ 相对于一个 $\mathcal{O}_X$-代数的整闭包}{\hyperref[II.6.3.4-zh]{6.3.4}}
\egaTerme{仿射锥的泛射影闭包}{\hyperref[II.8.3.1-zh]{8.3.1}}

\oldpage[II]{210}
\markboth{A. GROTHENDIECK --- 第二章}{A. GROTHENDIECK --- 第二章}
\egaTerme{射影丛在扩域 $K$ 上有理的几何纤维（扩张自 $k(y)$）}{\hyperref[II.4.2.6-zh]{4.2.6}}
\egaTerme{由 $\mathcal{O}_Y$-模定义的射影丛}{\hyperref[II.4.1.1-zh]{4.1.1}}
\egaTerme{由 $\mathcal{O}_S$-模定义的向量丛}{\hyperref[II.1.7.8-zh]{1.7.8}}
\egaTerme{有限（态射）}{\hyperref[II.6.1.1-zh]{6.1.1}}
\egaTerme{在一个预概形上有限（预概形）}{\hyperref[II.6.1.1-zh]{6.1.1}}
\egaTerme{有限的（拟凝聚 $\mathcal{O}_S$-代数）}{\hyperref[II.6.1.2-zh]{6.1.2}}
\egaTerme{局部自由模的自同态的初等对称函数}{\hyperref[II.6.4.1-zh]{6.4.1}}
\egaTerme{整环上的有限型模的自同态的初等对称函数}{\hyperref[II.6.4.2-zh]{6.4.2}}
\egaTerme{约化诺特环上的有限型模的自同态的初等对称函数}{\hyperref[II.6.4.7-zh]{6.4.7}}
\egaTerme{常秩局部自由 $\mathcal{A}$-模的自同态的初等对称函数}{\hyperref[II.6.4.8-zh]{6.4.8}}
\egaTerme{有限型 $\mathcal{O}_X$-模的自同态在局部整预概形 $X$ 上的初等对称函数}{\hyperref[II.6.4.9-zh]{6.4.9}}
\egaTerme{有限型 $\mathcal{O}_X$-模的自同态在局部诺特且约化的预概形 $X$ 上的初等对称函数}{\hyperref[II.6.4.10-zh]{6.4.10}}
\egaTerme{$\mathcal{A}$-代数的一个截面的初等对称函数}{\hyperref[II.6.5.1-zh]{6.5.1}}

\egaTerme{分次 $\mathcal{S}$-模的一个 $\mathcal{S}$-子模的齐次化}{\hyperref[II.8.13.2-zh]{8.13.2}}
\egaTerme{分次环同态}{\hyperref[II.2.1.2-zh]{2.1.2}}
\egaTerme{分次模的 $k$ 次同态}{\hyperref[II.2.1.2-zh]{2.1.2}}

\egaTerme{$S_+$ 的理想、$S_+$ 的分次素理想（$S$ 为分次环）}{\hyperref[II.2.1.10-zh]{2.1.10}}

\egaTerme{周引理}{\hyperref[II.5.6.1-zh]{5.6.1}}
\egaTerme{射影锥的无穷远轨迹}{\hyperref[II.8.3.3-zh]{8.3.3}}

\egaTerme{典范态射 $G(\mathcal{E})\to\Proj(S)$，$S=\bigoplus_{n\geq0}\Gamma(X,\mathcal{L}^{\otimes n})$}{\hyperref[II.4.5.1-zh]{4.5.1}}
\egaTerme{典范态射 $X\to\Spec(\Gamma(X,\mathcal{O}_X))$}{\hyperref[II.5.1.1-zh]{5.1.1}}
\egaTerme{塞格雷态射}{\hyperref[II.4.3.1-zh]{4.3.1}}

\egaTerme{$S_+$ 的幂零根（$S$ 为分次环）}{\hyperref[II.2.1.10-zh]{2.1.10}}
\egaTerme{$\mathcal{S}_+$ 的幂零根（$\mathcal{S}$ 为拟凝聚分次 $\mathcal{O}_Y$-代数）}{\hyperref[II.3.1.12-zh]{3.1.12}}
\egaTerme{$\mathcal{A}$-代数的一个截面的范数}{\hyperref[II.6.5.1-zh]{6.5.1}}
\egaTerme{可逆 $\mathcal{B}$-模的范数}{\hyperref[II.6.5.2-zh]{6.5.2}}
\egaTerme{可逆 $\mathcal{B}$-模的一个截面的范数}{\hyperref[II.6.5.3-zh]{6.5.3}}
\egaTerme{可逆 $\mathcal{O}_{X'}$-模的范数}{\hyperref[II.6.5.5-zh]{6.5.5}}

\egaTerme{周期的（分次代数）}{\hyperref[II.8.14.12-zh]{8.14.12}}
\egaTerme{整环上的有限型模的自同态的特征多项式}{\hyperref[II.6.4.2-zh]{6.4.2}}
\egaTerme{约化诺特环上的有限型模的自同态的特征多项式}{\hyperref[II.6.4.7-zh]{6.4.7}}
\egaTerme{仿射（射影）锥的顶点预概形}{\hyperref[II.8.3.3-zh]{8.3.3}}
\egaTerme{具有有限表示（分次模）}{\hyperref[II.2.1.1-zh]{2.1.1}}
\egaTerme{具有有限表示（分次 $\mathcal{S}$-模）}{\hyperref[II.3.1.1-zh]{3.1.1}}
\egaTerme{两个分次模的分次张量积}{\hyperref[II.2.1.2-zh]{2.1.2}}
\egaTerme{射影（态射）}{\hyperref[II.5.5.2-zh]{5.5.2}}
\egaTerme{在一个预概形上射影（预概形）}{\hyperref[II.5.5.2-zh]{5.5.2}}
\egaTerme{紧合（态射）}{\hyperref[II.5.4.1-zh]{5.4.1}}
\egaTerme{紧合（子集）}{\hyperref[II.5.4.10-zh]{5.4.10}}
\egaTerme{在一个预概形上紧合（预概形）}{\hyperref[II.5.4.1-zh]{5.4.1}}

\egaTerme{拟仿射（态射）}{\hyperref[II.5.1.1-zh]{5.1.1}}
\egaTerme{在一个预概形上拟仿射（预概形）}{\hyperref[II.5.1.1-zh]{5.1.1}}
\egaTerme{拟仿射（概形）}{\hyperref[II.5.1.1-zh]{5.1.1}}

\oldpage[II]{211}
\markboth{第二章}{第二章}
\egaTerme{拟有限（态射）}{\hyperref[II.6.2.3-zh]{6.2.3}}
\egaTerme{在一个预概形上拟有限（预概形）}{\hyperref[II.6.2.3-zh]{6.2.3}}
\egaTerme{拟射影（态射）}{\hyperref[II.5.3.1-zh]{5.3.1}}
\egaTerme{在一个预概形上拟射影（预概形）}{\hyperref[II.5.3.1-zh]{5.3.1}}

\egaTerme{$S_+$ 的一个理想的根（$S$ 为分次环）}{\hyperref[II.2.1.10-zh]{2.1.10}}
\egaTerme{去顶点射影锥到其无穷远轨迹上的典范回缩}{\hyperref[II.8.3.5-zh]{8.3.5}}

\egaTerme{$\mathcal{O}_X(-1)$ 的典范截面（$X$ 为暴涨预概形）}{\hyperref[II.8.1.9-zh]{8.1.9}}
\egaTerme{仿射锥的零截面}{\hyperref[II.8.3.3-zh]{8.3.3}}
\egaTerme{向量丛的零截面}{\hyperref[II.1.7.9-zh]{1.7.9}}
\egaTerme{仿射（射影）锥的顶点截面}{\hyperref[II.8.3.3-zh]{8.3.3}}
\egaTerme{$\mathcal{O}_S$-代数的谱}{\hyperref[II.1.3.1-zh]{1.3.1}}
\egaTerme{分次环的齐次素谱}{\hyperref[II.2.3.1-zh]{2.3.1}}
\egaTerme{拟凝聚分次 $\mathcal{O}_S$-代数的齐次谱}{\hyperref[II.3.1.3-zh]{3.1.3}}

\egaTerme{(TN)-单射、(TN)-满射、(TN)-双射（分次模同态）}{\hyperref[II.2.7.2-zh]{2.7.2}}
\egaTerme{分次模的 (TN)-同构}{\hyperref[II.2.7.2-zh]{2.7.2}}
\egaTerme{(TN)-单射、(TN)-满射、(TN)-双射（分次环同态）}{\hyperref[II.2.9.1-zh]{2.9.1}}
\egaTerme{分次环的 (TN)-同构}{\hyperref[II.2.9.1-zh]{2.9.1}}
\egaTerme{(TN)-单射、(TN)-满射、(TN)-双射（分次 $\mathcal{S}$-模同态）}{\hyperref[II.3.4.2-zh]{3.4.2}}
\egaTerme{分次 $\mathcal{S}$-模的 (TN)-同构}{\hyperref[II.3.4.2-zh]{3.4.2}}
\egaTerme{(TN)-单射、(TN)-满射、(TN)-双射（分次 $\mathcal{O}_Y$-代数同态）}{\hyperref[II.3.6.1-zh]{3.6.1}}
\egaTerme{分次 $\mathcal{O}_Y$-代数的 (TN)-同构}{\hyperref[II.3.6.1-zh]{3.6.1}}
\egaTerme{$\Proj(S)$ 上的谱拓扑}{\hyperref[II.2.3.3-zh]{2.3.3}}
\egaTerme{相对于 $q$ 极丰沛、相对于 $Y$ 极丰沛（可逆 $\mathcal{O}_X$-模）}{\hyperref[II.4.4.2-zh]{4.4.2}}
\egaTerme{有限型（分次 $\mathcal{S}$-模）}{\hyperref[II.3.1.1-zh]{3.1.1}}

\egaTerme{普遍闭（态射）}{\hyperref[II.5.4.9-zh]{5.4.9}}

\endgroup

\clearpage
\oldpage[II]{212}
\thispagestyle{empty}
\mbox{}
% Locked French authority: ega2-indexes-and-contents-fr.tex, 27454 B,
% SHA-256 2CEC53C734543C30ED2FE2FD01B462E5C51DE4D303CEAA7F8001316BB78EDBB1.
% Sequential source scope: lines 361-461, original table of contents.

\clearpage
\oldpage[II]{213}
\markboth{目录}{目录}
\phantomsection
\addcontentsline{toc}{section}{原目录}
\section*{目录}
\thispagestyle{plain}

\begingroup
\small
\renewcommand{\arraystretch}{1.04}
\begin{longtable}{@{}p{0.09\textwidth}p{0.69\textwidth}r@{}}
\textbf{单元} & \textbf{标题} & \textbf{原页码}\\
\hline
\textbf{II} & \textbf{若干类态射的初步整体研究} & \textbf{5}\\
\hyperref[section:II.1-zh]{\textsection1} & 仿射态射 & 5\\
\hyperref[subsection:II.1.1-zh]{1.1} & $S$-预概形与 $\mathcal{O}_S$-代数层 & 5\\
\hyperref[subsection:II.1.2-zh]{1.2} & 相对仿射的预概形 & 6\\
\hyperref[subsection:II.1.3-zh]{1.3} & 在 $S$ 上与 $\mathcal{O}_S$-代数层相伴的仿射预概形 & 8\\
\hyperref[subsection:II.1.4-zh]{1.4} & 相对于 $S$ 仿射的预概形上的拟凝聚层 & 9\\
\hyperref[II.1.5-zh]{1.5} & 基预概形的换基 & 12\\
\hyperref[II.1.6-zh]{1.6} & 仿射态射 & 14\\
\hyperref[II.1.7-zh]{1.7} & 与模层相伴的向量丛 & 14\\
\hyperref[section:II.2-zh]{\textsection2} & 齐次素谱 & 19\\
\hyperref[subsection:II.2.1-zh]{2.1} & 分次环与分次模概述 & 19\\
\hyperref[subsection:II.2.2-zh]{2.2} & 分次环的分式环 & 23\\
\hyperref[subsection:II.2.3-zh]{2.3} & 分次环的齐次素谱 & 25\\
\hyperref[subsection:II.2.4-zh]{2.4} & $\Proj(S)$ 上的概形结构 & 28\\
\hyperref[subsection:II.2.5-zh]{2.5} & 分次模的伴随层 & 30\\
\hyperref[subsection:II.2.6-zh]{2.6} & 与 $\Proj(S)$ 上一个层相伴的分次 $S$-模 & 36\\
\hyperref[subsection:II.2.7-zh]{2.7} & 有限性条件 & 38\\
\hyperref[subsection:II.2.8-zh]{2.8} & 函子性 & 41\\
\hyperref[subsection:II.2.9-zh]{2.9} & 概形 $\Proj(S)$ 的闭子概形 & 48\\
\hyperref[section:II.3-zh]{\textsection3} & 分次代数层的齐次谱 & 49\\
\hyperref[subsection:II.3.1-zh]{3.1} & 拟凝聚分次 $\mathcal{O}_Y$-代数层的齐次谱 & 49\\
\hyperref[subsection:II.3.2-zh]{3.2} & 与分次 $\mathcal{S}$-模层相伴的 $\Proj(\mathcal{S})$ 上的层 & 54\\
\hyperref[subsection:II.3.3-zh]{3.3} & 与 $\Proj(\mathcal{S})$ 上一个层相伴的分次 $\mathcal{S}$-模 & 56\\
\hyperref[subsection:II.3.4-zh]{3.4} & 有限性条件 & 59\\
\hyperref[subsection:II.3.5-zh]{3.5} & 函子性 & 61\\
\hyperref[subsection:II.3.6-zh]{3.6} & 预概形 $\Proj(\mathcal{S})$ 的闭子概形 & 64\\
\hyperref[subsection:II.3.7-zh]{3.7} & 预概形到齐次谱的态射 & 65\\
\hyperref[subsection:II.3.8-zh]{3.8} & 浸入齐次谱的判别法 & 69\\
\end{longtable}
\oldpage[II]{214}
\markboth{A. GROTHENDIECK --- 第二章}{A. GROTHENDIECK --- 第二章}
\begin{longtable}{@{}p{0.09\textwidth}p{0.69\textwidth}r@{}}
\hyperref[section:II.4-zh]{\textsection4} & 泛射影丛；丰沛层 & 71\\
\hyperref[subsection:II.4.1-zh]{4.1} & 泛射影丛的定义 & 71\\
\hyperref[subsection:II.4.2-zh]{4.2} & 预概形到泛射影丛的态射 & 72\\
\hyperref[subsection:II.4.3-zh]{4.3} & Segre 态射 & 76\\
\hyperref[subsection:II.4.4-zh]{4.4} & 泛射影丛中的浸入；极丰沛层 & 78\\
\hyperref[subsection:II.4.5-zh]{4.5} & 丰沛层 & 83\\
\hyperref[II.4.6-zh]{4.6} & 相对丰沛层 & 89\\
\hyperref[section:II.5-zh]{\textsection5} & 拟仿射态射；拟射影态射；紧合态射；射影态射 & 94\\
\hyperref[subsection:II.5.1-zh]{5.1} & 拟仿射态射 & 94\\
\hyperref[subsection:II.5.2-zh]{5.2} & 塞尔判别法 & 97\\
\hyperref[subsection:II.5.3-zh]{5.3} & 拟射影态射 & 99\\
\hyperref[subsection:II.5.4-zh]{5.4} & 紧合态射和普遍闭态射 & 100\\
\hyperref[subsection:II.5.5-zh]{5.5} & 射影态射 & 103\\
\hyperref[subsection:II.5.6-zh]{5.6} & 周氏引理 & 106\\
\hyperref[section:II.6-zh]{\textsection6} & 整态射与有限态射 & 110\\
\hyperref[subsection:II.6.1-zh]{6.1} & 在另一预概形上整的预概形 & 110\\
\hyperref[subsection:II.6.2-zh]{6.2} & 拟有限态射 & 114\\
\hyperref[subsection:II.6.3-zh]{6.3} & 预概形的整闭包 & 116\\
\hyperref[subsection:II.6.4-zh]{6.4} & $\mathcal{O}_X$-模自同态的行列式 & 120\\
\hyperref[subsection:II.6.5-zh]{6.5} & 可逆层的范数 & 125\\
\hyperref[subsection:II.6.6-zh]{6.6} & 应用：丰沛判据 & 130\\
\hyperref[subsection:II.6.7-zh]{6.7} & Chevalley 定理 & 135\\
\hyperref[section:II.7-zh]{\textsection7} & 赋值判别法 & 138\\
\hyperref[subsection:II.7.1-zh]{7.1} & 赋值环回顾 & 138\\
\hyperref[subsection:II.7.2-zh]{7.2} & 分离性的赋值判别法 & 141\\
\hyperref[subsection:II.7.3-zh]{7.3} & 紧合性的赋值判别法 & 143\\
\hyperref[subsection:II.7.4-zh]{7.4} & 代数曲线与一维函数域 & 148\\
\hyperref[section:II.8-zh]{\textsection8} & 概形的暴涨；投影锥；泛射影闭包 & 152\\
\hyperref[subsection:II.8.1-zh]{8.1} & 预概形的暴涨 & 152\\
\hyperref[subsection:II.8.2-zh]{8.2} & 分次环中局部化的预备结果 & 157\\
\hyperref[subsection:II.8.3-zh]{8.3} & 投影锥 & 162\\
\hyperref[subsection:II.8.4-zh]{8.4} & 向量丛的泛射影闭包 & 168\\
\hyperref[subsection:II.8.5-zh]{8.5} & 函子性 & 169\\
\hyperref[subsection:II.8.6-zh]{8.6} & 去顶点锥的一个典范同构 & 171\\
\end{longtable}
\oldpage[II]{215}
\markboth{目录}{目录}
\begin{longtable}{@{}p{0.09\textwidth}p{0.69\textwidth}r@{}}
\hyperref[subsection:II.8.7-zh]{8.7} & 投影锥的暴涨 & 173\\
\hyperref[subsection:II.8.8-zh]{8.8} & 丰沛层与收缩 & 177\\
\hyperref[subsection:II.8.9-zh]{8.9} & 格劳尔特丰沛性判据：陈述 & 182\\
\hyperref[subsection:II.8.10-zh]{8.10} & 格劳尔特丰沛性判据：证明 & 184\\
\hyperref[subsection:II.8.11-zh]{8.11} & 收缩的唯一性 & 189\\
\hyperref[subsection:II.8.12-zh]{8.12} & 投影锥上的拟凝聚层 & 191\\
\hyperref[subsection:II.8.13-zh]{8.13} & 子层与闭子概形的射影闭包 & 195\\
\hyperref[subsection:II.8.14-zh]{8.14} & 分次 $\mathcal{S}$-模相伴层的补充 & 197\\
\multicolumn{2}{@{}l}{参考文献（续）} & 205\\
\multicolumn{2}{@{}l}{符号索引} & 207\\
\multicolumn{2}{@{}l}{术语索引} & 209\\
\multicolumn{2}{@{}l}{勘误与补遗（第一表）} & 217\\
\end{longtable}

\begin{flushright}
\emph{1960 年 7 月 22 日收到。}
\end{flushright}
\endgroup
% Locked French authority: ega2-errata-addenda-fr.tex, 19746 B,
% SHA-256 EC20D329248B99CF0533CB868DBDF8135D5BDAFA233133814DB67F8CD4F09643.
% Sequential source scope: lines 1-200, Errata et Addenda, Liste 1,
% through the proof of I.1.8.1 and before Environment I.1.8.2.

\clearpage
\oldpage[II]{217}
\markboth{勘误与增补，表 1}{勘误与增补，表 1}
\phantomsection
\addcontentsline{toc}{section}{勘误与增补（表 1）}
\section*{勘误与增补}
\thispagestyle{plain}
\begin{center}
\large （表 1）
\end{center}

\subsection*{第 0 章}

\paragraph{(0, 2.1.3).}
第 21 页倒数第 15 行，将 «~inter-section~» 替换为
«~intersection~»。

\paragraph{(0, 2.1.6).}
第 22 页第 13 行，将
\[
  Z_i\cap C\left(\bigcup_{j\ne i}Z_j\right)
  \quad\hbox{替换为}\quad
  \bigcup_{j\ne i}Z_j.
\]

\paragraph{(0, 3.1.6).}
第 25 页倒数第 17 行，将 $\Gamma(U,\sh{F})$ 替换为
$\Gamma(V,\sh{F})$。

\paragraph{(0, 4.1.1).}
第 36 页第 12 行，将 $\sh{G}\to\sh{F}$ 替换为
$\sh{B}\to\sh{A}$。

\paragraph{(0, 4.2.1).}
第 39 页倒数第 11 行，将 «~faisceaux~» 替换为 «~faisceau~»。

\paragraph{(0, 5.1.2).}
在第 45 页第 14 行之后添加：

如果 $\sh{G}$ 是一个由其在 $Y$ 上的截面生成的 $\sh{O}_Y$-模，
则 $f^*(\sh{G})$ 由其在 $X$ 上的截面生成，因为 $f^*$ 是右正合函子。

\paragraph{(0, 5.2.4).}
第 46 页倒数第 18 行，将 $f_U$ 替换为 $f_U^*$。

\paragraph{(0, 5.3.4).}
在第 47 页倒数第 16 行之前添加：

如果
\[
  \sh{A}\to\sh{B}\to\sh{C}\to\sh{D}\to\sh{E}
\]
是一个 $\sh{O}_X$-模正合列，并且
$\sh{A},\sh{B},\sh{D},\sh{E}$ 都是凝聚的，则 $\sh{C}$ 是凝聚的。

\paragraph{(0, 5.4.1).}
在第 49 页第 2 行之后添加：

假设 $\sh{O}_X$ 是凝聚的，并设 $\sh{F}$ 是一个凝聚
$\sh{O}_X$-模。如果在一点 $x\in X$ 处，$\sh{F}_x$ 是秩为 $n$ 的
自由 $\sh{O}_x$-模，那么存在 $x$ 的一个邻域 $U$，使得
$\sh{F}|U$ 是秩为 $n$ 的局部自由模；事实上，此时 $\sh{F}_x$
同构于 $\sh{O}_x^n$，而该命题由 (0, 5.2.7) 得出。

\paragraph{(0, 6.6.2).}
将第 59 页第 10、11 行替换为：

事实上，这由 (0, 6.4.1, d)) 得出。

\paragraph{(0, 6.7.1).}
在第 59 页倒数第 10 行之前添加：

如果 $f$ 是平坦态射，也称 $X$ \emph{在 $Y$ 上平坦}，或称其
\emph{为 $Y$-平坦的}。

\paragraph{(0, 7.2.9).}
第 65 页倒数第 11 行，将 $M_n$ 替换为 $M_{n-1}$。

\subsection*{第 I 章}

\paragraph{(I, 1.2.1).}
第 83 页第 18 行，将 $K$ 替换为 $k$（共两处）。

\paragraph{(I, 1.2.7).}
第 84 页倒数第 16 行，将 $A$ 替换为 $A'$（共两处），并将
$\mathfrak N$ 替换为 $\mathfrak N'$；倒数第 17 行，将 $X$ 替换为
$X'$。

\paragraph{(I, 1.7.3) 和 (I, 2.2.4).}
这些条目中的陈述可作如下推广（据 J.~Tate 的一项评注）：

\subsection*{1.8. 仿射概形中的局部赋环空间态射}
\label{ega2:addendum:subsection:I.1.8-zh}

\begin{proposition}[1.8.1]
\label{ega2:addendum:I.1.8.1-zh}
设 $(S,\sh{O}_S)$ 是一个仿射概形，$(X,\sh{O}_X)$ 是一个局部赋环空间。
则环同态
\oldpage[II]{218}
\markboth{A. GROTHENDIECK --- 第二章}{A. GROTHENDIECK --- 第二章}
\[
  \Gamma(S,\sh{O}_S)\to\Gamma(X,\sh{O}_X)
\]
的集合与满足下述条件的赋环空间态射
$(\psi,\theta):(X,\sh{O}_X)\to(S,\sh{O}_S)$ 的集合之间存在典范双射：
对每个 $x\in X$，$\theta_x^\sharp:\sh{O}_{\psi(x)}\to\sh{O}_x$
都是局部同态。
\end{proposition}

首先注意，如果 $(X,\sh{O}_X)$、$(S,\sh{O}_S)$ 是任意两个赋环空间，
从 $(X,\sh{O}_X)$ 到 $(S,\sh{O}_S)$ 的一个态射 $(\psi,\theta)$
便典范地定义一个环同态
$\Gamma(\theta):\Gamma(S,\sh{O}_S)\to\Gamma(X,\sh{O}_X)$；由此得到
第一个映射
\begin{equation}
\label{ega2:addendum:I.1.8.1.1-zh}
  \rho:\Hom((X,\sh{O}_X),(S,\sh{O}_S))
  \to\Hom(\Gamma(S,\sh{O}_S),\Gamma(X,\sh{O}_X)).
\tag{1.8.1.1}
\end{equation}
反过来，在命题的假设下，令 $A=\Gamma(S,\sh{O}_S)$，并考虑一个环同态
$\vphi:A\to\Gamma(X,\sh{O}_X)$。对每个 $x\in X$，满足
$\vphi(f)(x)=0$ 的 $f\in A$ 所成集合显然是 $A$ 的一个素理想，因为
$\sh{O}_x/\mathfrak m_x=\kres(x)$ 是一个域；因此它是
$S=\Spec(A)$ 的一个元素，仍记为 ${}^a\vphi(x)$。此外，按照定义
(0, 5.5.2)，对每个 $f\in A$ 有
\[
  {}^a\vphi^{-1}(D(f))=X_f,
\]
这证明了 ${}^a\vphi$ 是一个连续映射 $X\to S$。接着定义一个
$\sh{O}_S$-模同态
\[
  \widetilde{\vphi}:\sh{O}_S\to{}^a\vphi_*(\sh{O}_X)
\]
；对每个 $f\in A$，有 $\Gamma(D(f),\sh{O}_S)=A_f$（I, 1.3.6）；
对每个 $s\in A$，令 $s/f\in A_f$ 对应于元素
\[
  (\vphi(s)|X_f)(\vphi(f)|X_f)^{-1}
  \in\Gamma(X_f,\sh{O}_X)
  =\Gamma(D(f),{}^a\vphi_*(\sh{O}_X)),
\]
立即可验证（将 $D(f)$ 换成 $D(fg)$）这确实定义了一个
$\sh{O}_S$-模同态，从而得到一个赋环空间态射
$({}^a\vphi,\widetilde{\vphi})$。此外，沿用上述记号，并为简便起见令
$y={}^a\vphi(x)$，立即由 (0, 3.7.1) 看出
\[
  \widetilde{\vphi}_x^\sharp(s_y/f_y)
  =(\vphi(s)_x)(\vphi(f)_x)^{-1};
\]
由于关系 $s_y\in\mathfrak m_y$ 按定义等价于
$\vphi(s)_x\in\mathfrak m_x$，可见 $\widetilde{\vphi}_x^\sharp$ 是一个
局部同态 $\sh{O}_y\to\sh{O}_x$，这样便定义了第二个映射
\[
  \sigma:\Hom(\Gamma(S,\sh{O}_S),\Gamma(X,\sh{O}_X))\to\mathfrak L,
\]
其中 $\mathfrak L$ 是满足下述条件的赋环空间态射
$(\psi,\theta):(X,\sh{O}_X)\to(S,\sh{O}_S)$ 所成集合：对每个
$x\in X$，$\theta_x^\sharp$ 都是局部同态。

还需证明 $\sigma$ 与 $\rho$（限制于 $\mathfrak L$）互为逆映射。
$\widetilde{\vphi}$ 的定义立即表明
$\Gamma(\widetilde{\vphi})=\vphi$，因而 $\rho\circ\sigma$ 是恒等映射。
为证明 $\sigma\circ\rho$ 是恒等映射，从一个态射
$(\psi,\theta)\in\mathfrak L$ 出发，并令 $\vphi=\Gamma(\theta)$；
关于 $\theta_x^\sharp$ 的假设使我们能够通过取商，由该同态导出一个单同态
\[
  \theta^x:\kres(\psi(x))\to\kres(x)
\]
使得对任意截面 $f\in A=\Gamma(S,\sh{O}_S)$，都有
$\theta^x(f(\psi(x)))=\vphi(f)(x)$；因此关系 $f(\psi(x))=0$ 等价于
$\vphi(f)(x)=0$，这已经证明 ${}^a\vphi=\psi$。另一方面，由定义可知
交换图
\[
  \xymatrix{
    A\ar[r]^\vphi\ar[d] &
    \Gamma(X,\sh{O}_X)\ar[d]\\
    A_{\psi(x)}\ar[r]^{\theta_x^\sharp} &
    \sh{O}_x
  }
\]
是交换的；将 $\theta_x^\sharp$ 替换为
$\widetilde{\vphi}_x^\sharp$ 后所得的类似图也是交换的，故
$\widetilde{\vphi}_x^\sharp=\theta_x^\sharp$（0, 1.2.4），从而
$\widetilde{\vphi}=\theta$。

% Locked French authority: ega2-errata-addenda-fr.tex, 19746 B,
% SHA-256 EC20D329248B99CF0533CB868DBDF8135D5BDAFA233133814DB67F8CD4F09643.
% Sequential source scope: lines 201-372, Environments I.1.8.2--I.1.8.8.

\begin{env}[1.8.2]
\label{ega2:addendum:I.1.8.2-zh}
当 $(X,\sh{O}_X)$ 和 $(Y,\sh{O}_Y)$ 是局部赋环空间时，我们将考察满足
下述条件的态射 $(\psi,\theta):(X,\sh{O}_X)\to(Y,\sh{O}_Y)$：
对每个 $x\in X$，$\theta_x^\sharp$ 是一个局部同态
$\sh{O}_{\psi(x)}\to\sh{O}_x$。

今后当我们谈到
\oldpage[II]{219}
\markboth{勘误与增补，表 1}{勘误与增补，表 1}
\emph{局部赋环空间态射}时，始终是指这样的态射；以此定义态射后，
局部赋环空间显然构成一个范畴。对于这个范畴中的两个对象 $X,Y$，
$\Hom(X,Y)$ 表示从 $X$ 到 $Y$ 的局部赋环空间态射所成的集合
（即在 (1.8.1) 中记为 $\mathfrak L$ 的集合）；当需要考察从 $X$ 到
$Y$ 的赋环空间态射所成集合时，为避免混淆，将其记为
$\Hom_{\mathrm{an}}(X,Y)$。因此，映射 (1.8.1.1) 写成
\begin{equation}
\label{ega2:addendum:I.1.8.2.1-zh}
  \rho:\Hom_{\mathrm{an}}(X,Y)
  \to\Hom(\Gamma(Y,\sh{O}_Y),\Gamma(X,\sh{O}_X)),
\tag{1.8.2.1}
\end{equation}
而它的限制
\begin{equation}
\label{ega2:addendum:I.1.8.2.2-zh}
  \rho':\Hom(X,Y)
  \to\Hom(\Gamma(Y,\sh{O}_Y),\Gamma(X,\sh{O}_X))
\tag{1.8.2.2}
\end{equation}
在局部赋环空间范畴中是关于 $X$ 和 $Y$ 的函子性映射。
\end{env}

\begin{corollary}[1.8.3]
\label{ega2:addendum:I.1.8.3-zh}
设 $(Y,\sh{O}_Y)$ 是一个局部赋环空间。$Y$ 是仿射概形的充分必要条件是：
对于每个局部赋环空间 $(X,\sh{O}_X)$，映射 (1.8.2.2) 都是双射。
\end{corollary}

命题 (1.8.1) 表明该条件是必要的。反过来，假设该条件成立，并令
$A=\Gamma(Y,\sh{O}_Y)$；由此假设和 (1.8.1) 可知，从局部赋环空间范畴
到集合范畴的函子 $X\mapsto\Hom(X,Y)$ 与
$X\mapsto\Hom(X,\Spec(A))$ 同构。我们知道，这蕴含存在一个典范同构
$X\to\Spec(A)$（参见 0, 8）。

\begin{env}[1.8.4]
\label{ega2:addendum:I.1.8.4-zh}
设 $S=\Spec(A)$ 是一个仿射概形；以 $(S',A')$ 表示如下赋环空间：其底空间
只有一个点，其结构层 $A'$ 是由环 $A$ 在 $S'$ 上定义的（必为简单的）层。
设 $\pi:S\to S'$ 是从 $S$ 到 $S'$ 的唯一映射；另一方面注意到，对于
$S$ 的每个开集 $U$，都有典范映射
$\Gamma(S',A')=\Gamma(S,\sh{O}_S)\to\Gamma(U,\sh{O}_S)$，
因而它定义一个环层的 $\pi$-态射 $\iota:A'\to\sh{O}_S$。这样便典范地
定义了一个赋环空间态射
$i=(\pi,\iota):(S,\sh{O}_S)\to(S',A')$。对于每个 $A$-模 $M$，
用 $M'$ 表示由 $M$ 在 $S'$ 上定义的简单层；它显然是一个 $A'$-模。
显然有 $i_*(\widetilde M)=M'$（I, 1.3.7）。
\end{env}

\begin{lemma}[1.8.5]
\label{ega2:addendum:I.1.8.5-zh}
沿用 (1.8.4) 的记号，对于每个 $A$-模 $M$，函子性的典范
$\sh{O}_S$-同态（0, 4.4.3.3）
\begin{equation}
\label{ega2:addendum:I.1.8.5.1-zh}
  i^*(i_*(\widetilde M))\to\widetilde M
\tag{1.8.5.1}
\end{equation}
是一个同构。
\end{lemma}

事实上，(1.8.5.1) 的两端都是右正合的（函子
$M\mapsto i_*(\widetilde M)$ 显然是正合的），并且与直和可交换；把 $M$
看作某个同态 $A^{(I)}\to A^{(J)}$ 的余核，便归结为在 $M=A$ 的情形下
证明引理，而在该情形下结论是显然的。

\begin{corollary}[1.8.6]
\label{ega2:addendum:I.1.8.6-zh}
设 $(X,\sh{O}_X)$ 是一个赋环空间，$u:X\to S$ 是一个
\oldpage[II]{220}
\markboth{A. GROTHENDIECK --- 第二章}{A. GROTHENDIECK --- 第二章}
赋环空间态射。对于每个 $A$-模 $M$，沿用 (1.8.4) 的记号，有函子性的
典范 $\sh{O}_X$-模同构
\begin{equation}
\label{ega2:addendum:I.1.8.6.1-zh}
  u^*(\widetilde M)\isoto u^*(i^*(M')).
\tag{1.8.6.1}
\end{equation}
\end{corollary}

\begin{corollary}[1.8.7]
\label{ega2:addendum:I.1.8.7-zh}
在 (1.8.6) 的假设下，对于每个 $A$-模 $M$ 和每个 $\sh{O}_X$-模
$\sh{F}$，都有关于 $M$ 和 $\sh{F}$ 为函子性的典范同构
\begin{equation}
\label{ega2:addendum:I.1.8.7.1-zh}
  \Hom_{\sh{O}_S}(\widetilde M,u_*(\sh{F}))
  \isoto\Hom_A(M,\Gamma(X,\sh{F})).
\tag{1.8.7.1}
\end{equation}
\end{corollary}

事实上，由 (0, 4.4.3) 和引理 (1.8.5)，有双函子的典范同构
\[
  \Hom_{\sh{O}_S}(\widetilde M,u_*(\sh{F}))
  \isoto\Hom_{A'}(M',i_*(u_*(\sh{F}))),
\]
而右端显然正是 $\Hom_A(M,\Gamma(X,\sh{F}))$。注意，典范同态
(1.8.7.1) 把每个 $\sh{O}_S$-同态
$h:\widetilde M\to u_*(\sh{F})$（换言之，把每个 $u$-态射
$\widetilde M\to\sh{F}$）映为 $A$-同态
\[
  \Gamma(h):M\to\Gamma(S,u_*(\sh{F}))=\Gamma(X,\sh{F}).
\]

\begin{env}[1.8.8]
\label{ega2:addendum:I.1.8.8-zh}
沿用 (1.8.4) 的记号，显然（0, 4.1.1），每个赋环空间态射
$(\psi,\theta):X\to S'$ 都等价于给定一个环同态
$A\to\Gamma(X,\sh{O}_X)$。因此，还可以把 (1.8.1) 解释为定义了典范双射
\[
  \Hom(X,S)\isoto\Hom(X,S')
\]
（当然，右端指的是赋环空间态射，因为一般而言 $A$ 并不是局部环）。
更一般地，设 $X,Y$ 是两个局部赋环空间，并设 $(Y',A')$ 是如下赋环空间：
其底空间只有一个点，而其环层 $A'$ 是由环
$\Gamma(Y,\sh{O}_Y)$ 定义的简单层。可以把 (1.8.2.1) 解释为映射
\begin{equation}
\label{ega2:addendum:I.1.8.8.1-zh}
  \rho:\Hom_{\mathrm{an}}(X,Y)\to\Hom(X,Y').
\tag{1.8.8.1}
\end{equation}
(1.8.3) 的结果因而可以表述为：在局部赋环空间中，仿射概形恰由下述
性质刻画——$\rho$ 在 $\Hom(X,Y)$ 上的限制
\begin{equation}
\label{ega2:addendum:I.1.8.8.2-zh}
  \rho':\Hom(X,Y)\to\Hom(X,Y')
\tag{1.8.8.2}
\end{equation}
对于每个局部赋环空间 $X$ 都是双射。在后面的一章中，我们将推广这一
定义，从而能把一个仍记为 $\Spec(Z)$ 的局部赋环空间与任意赋环空间
$Z$ 相伴（而不再只限于底空间仅有一个点的赋环空间）；这将成为研究任意
赋环空间上预概形的«~相对~»理论的出发点，并推广第 I 章的结果。
\end{env}
% Locked French authority: ega2-errata-addenda-fr.tex, 19746 B,
% SHA-256 EC20D329248B99CF0533CB868DBDF8135D5BDAFA233133814DB67F8CD4F09643.
% Sequential source scope: lines 374-522, Environment I.1.8.9 through
% the terminal I.10.14.2 erratum on printed p.222.

\begin{env}[1.8.9]
\label{ega2:addendum:I.1.8.9-zh}
可将由一个局部赋环空间 $X$ 和一个 $\sh{O}_X$-模 $\sh{F}$ 组成的偶对
$(X,\sh{F})$ 看作构成一个范畴；该范畴中的一个态射
$(X,\sh{F})\to(Y,\sh{G})$ 是一个偶对 $(u,h)$，其中 $u:X\to Y$
是局部赋环空间的一个态射，
\oldpage[II]{221}
\markboth{勘误与增补，表 1}{勘误与增补，表 1}
而 $h:\sh{G}\to\sh{F}$ 是模的一个 $u$-态射。对于固定的
$(X,\sh{F})$ 和 $(Y,\sh{G})$，这些态射构成一个集合，记为
$\Hom((X,\sh{F}),(Y,\sh{G}))$；映射
$(u,h)\mapsto(\rho'(u),\Gamma(h))$ 是一个典范映射
\begin{equation}
\label{ega2:addendum:I.1.8.9.1-zh}
\begin{split}
  \Hom((X,\sh{F}),(Y,\sh{G}))\to
  \Hom\bigl(
    &(\Gamma(Y,\sh{O}_Y),\Gamma(Y,\sh{G})),\\
    &(\Gamma(X,\sh{O}_X),\Gamma(X,\sh{F}))
  \bigr),
\end{split}
\tag{1.8.9.1}
\end{equation}
它关于 $(X,\sh{F})$ 和 $(Y,\sh{G})$ 是函子性的，右端是与所考察的环和模
对应的双同态集合 (0, 1.0.2)。
\end{env}

\begin{corollary}[1.8.10]
\label{ega2:addendum:I.1.8.10-zh}
设 $Y$ 是一个局部赋环空间，$\sh{G}$ 是一个 $\sh{O}_Y$-模。$Y$ 是仿射概形且
$\sh{G}$ 是拟凝聚 $\sh{O}_Y$-模的充分必要条件是：对于每个由局部赋环空间
$X$ 和 $\sh{O}_X$-模 $\sh{F}$ 组成的偶对 $(X,\sh{F})$，典范映射
(1.8.8.1) 都是双射。
\end{corollary}

证明细节留给读者展开；该论证仿照 (1.8.3) 的论证，并使用 (1.8.1) 和
(1.8.7)。

\begin{remark}[1.8.11]
\label{ega2:addendum:I.1.8.11-zh}
陈述 (I, 1.7.3)、(I, 1.7.4) 和 (I, 2.2.4)，以及定义 (I, 1.6.1)，
都是 (1.8.1) 的特殊情形；同样，(I, 2.2.5) 由 (1.8.7) 得出。命题
(1.8.7) 还蕴含 (I, 1.6.3)（从而也蕴含 (I, 1.6.4)）这一特殊情形：
事实上，如果 $X$ 是仿射的且 $\Gamma(X,\sh{F})=N$，则函子
\[
  M\longmapsto
  \Hom_{\sh{O}_S}(\widetilde M,u_*(\widetilde N))
  \quad\hbox{和}\quad
  M\longmapsto
  \Hom_{\sh{O}_S}(\widetilde M,(N_{[\vphi]})\supertilde)
\]
由 (1.8.7) 和 (I, 1.3.8) 同构（其中
$\vphi:A\to\Gamma(X,\sh{O}_X)$ 对应于 $u$）。最后，(I, 1.6.5)
（从而也有 (I, 1.6.6)）由 (1.8.6) 以及下述事实得出：对于每个
$f\in A$，$A_f$-模
\[
  N'\otimes_{A'}A_f
  \quad\hbox{和}\quad
  (N'\otimes_{A'}A)_f
\]
（采用 (I, 1.6.5) 的记号）典范同构。
\end{remark}

\paragraph{插入在 (I, 3.2.8) 之后：}

\begin{env}[3.2.9]
\label{ega2:addendum:I.3.2.9-zh}
由 (1.8.1)，可以将 (I, 3.2.2) 细化如下：
$Z=\Spec(B\otimes_A C)$ 不仅是 $X=\Spec(B)$ 与 $Y=\Spec(C)$ 在
$S$-预概形范畴中的积，而且也是它们在 $S$ 上的局部赋环空间范畴中的积
（其中 $S$-态射的定义仿照 (I, 2.5.2)）。于是，(I, 3.2.6) 的证明实际上
表明，对于任意两个 $S$-预概形 $X,Y$，预概形 $X\times_S Y$ 不仅是
$X$ 与 $Y$ 在 $S$-预概形范畴中的积，而且也是它们在预概形 $S$ 上的
局部赋环空间范畴中的积。
\end{env}

\paragraph{(I, 4.4.5).}
第 126 页倒数第 14 行，将 $B$ 替换为 $A$；倒数第 13 行，将
«~$A$-代数~» 替换为 «~$B$-代数~»。

\paragraph{(I, 5.3.5).}
第 132 页倒数第 2 行，将 $Y\to S$ 替换为 $g:Y\to S$。

\paragraph{(I, 6.3.2.1).}
第 144 页倒数第 6 行，应为：
\[
  D(g_i)\subset W.
\]

\paragraph{(I, 7.3.8).}
将现有的错误文本替换为以下文本：

设 $X,Y$ 是两个整预概形；由此可知，$\sh{R}(X)$（相应地，
$\sh{R}(Y)$）是拟凝聚 $\sh{O}_X$-模（相应地，$\sh{O}_Y$-模）
(I, 7.3.3)。设 $f:X\to Y$ 是一个支配态射；则存在一个典范的
$\sh{O}_X$-模同态
\oldpage[II]{222}
\markboth{A. GROTHENDIECK --- 第二章}{A. GROTHENDIECK --- 第二章}
\begin{equation}
\label{ega2:addendum:I.7.3.8.1-zh}
  \tau:f^*(\sh{R}(Y))\to\sh{R}(X).
\tag{7.3.8.1}
\end{equation}

先假设 $X=\Spec(A)$ 和 $Y=\Spec(B)$ 是由整环 $A$ 和 $B$ 给出的仿射概形；
于是 $f$ 对应于一个单同态 $B\to A$，它延拓为从 $B$ 的分式域 $L$
到 $A$ 的分式域 $K$ 的一个单同态 $L\to K$。此时，同态 (7.3.8.1)
对应于典范同态
\[
  L\otimes_B A\to K
\]
(I, 1.6.5)。在一般情形下，对于满足 $U\subset X$、$V\subset Y$ 且
$f(U)\subset V$ 的每一对非空仿射开集，按上述方式定义一个同态
$\tau_{U,V}$；立即可见，如果 $U'\subset U$、$V'\subset V$ 且
$f(U')\subset V'$，则 $\tau_{U,V}$ 延拓 $\tau_{U',V'}$，这便证明了
我们的断言。如果 $x,y$ 分别是 $X$ 和 $Y$ 的泛点，则 $f(x)=y$，且
\[
  \bigl(f^*(\sh{R}(Y))\bigr)_x
  =\sh{O}_y\otimes_{\sh{O}_y}\sh{O}_x
  =\sh{O}_x
\]
(0, 4.3.1)，因此 $\tau_x$ 是一个同构。

\paragraph{(I, 9.5.2).}
在第 176 页的最后 3 行和第 177 页的最初 5 行中，将各处的 $X'$ 替换为
$Y'$，并将 $X''$ 替换为 $Y''$。

\paragraph{(I, 10.5.4).}
第 187 页倒数第 11 行，将 (10.3.4) 替换为 (10.3.5)。

\paragraph{(I, 10.6.3).}
第 189 页倒数第 3 行，将 $X$ 替换为 $\mathfrak X$。

\paragraph{(I, 10.14.2).}
第 210 页第 5 行，将 «~凝聚 $\sh{O}_{\mathfrak X}$-模~» 替换为
«~$\sh{O}_{\mathfrak X}$ 的凝聚理想~»；第 7 行，将
«~凝聚 $\sh{O}_{\mathfrak X}$-模~» 替换为
«~$\sh{O}_{\mathfrak X}$ 的凝聚理想~»。
\end{document}
